% main.tex
\documentclass[UTF8,openany,zihao=-4]{ctexbook}

% =========================
% 页面与版式
% =========================
\usepackage[a4paper,margin=2.5cm]{geometry}
\usepackage{setspace}
\usepackage{fancyhdr}
\usepackage{titlesec}
\usepackage{tocloft}

% =========================
% 数学公式
% =========================
\usepackage{amsmath}
\usepackage{amssymb}
\usepackage{mathtools}
\usepackage{bm}

% =========================
% 图形与架构图
% =========================
\usepackage{graphicx}
\usepackage{float}
\usepackage{subcaption}

% TikZ：用于绘制 AI Infra 架构图、并行策略图、网络拓扑图等
\usepackage{tikz}
\usetikzlibrary{
  arrows.meta,
  positioning,
  shapes.geometric,
  shapes.misc,
  shapes.symbols,
  fit,
  calc,
  backgrounds,
  matrix,
  chains,
  decorations.pathreplacing
}

% =========================
% 表格
% =========================
\usepackage{booktabs}
\usepackage{multirow}
\usepackage{array}
\usepackage{longtable}

% =========================
% 代码高亮
% =========================
\usepackage{xcolor}
\usepackage{listings}

\definecolor{codebg}{RGB}{248,248,248}
\definecolor{codegray}{RGB}{90,90,90}
\definecolor{codeblue}{RGB}{40,90,180}
\definecolor{codegreen}{RGB}{40,130,80}
\definecolor{codered}{RGB}{180,60,60}

\lstdefinestyle{aicode}{
  backgroundcolor=\color{codebg},
  basicstyle=\ttfamily\small,
  keywordstyle=\color{codeblue}\bfseries,
  commentstyle=\color{codegreen},
  stringstyle=\color{codered},
  numberstyle=\tiny\color{codegray},
  numbers=left,
  stepnumber=1,
  numbersep=8pt,
  frame=single,
  breaklines=true,
  tabsize=2,
  showstringspaces=false,
  captionpos=b
}

\lstset{style=aicode}

% 如需使用 minted，可替换 listings：
% \usepackage{minted}
% 编译时需加 -shell-escape

% =========================
% 超链接与交叉引用
% =========================
\usepackage[
  colorlinks=true,
  linkcolor=blue,
  citecolor=blue,
  urlcolor=blue
]{hyperref}

\usepackage[nameinlink]{cleveref}

% =========================
% 算法伪代码
% =========================
\usepackage{algorithm}
\usepackage{algorithmic}

% =========================
% 自定义命令
% =========================
\newcommand{\tensor}[1]{\bm{#1}}
\newcommand{\mat}[1]{\bm{#1}}
\newcommand{\vect}[1]{\bm{#1}}
\newcommand{\softmax}{\operatorname{softmax}}
\newcommand{\LayerNorm}{\operatorname{LayerNorm}}
\newcommand{\RMSNorm}{\operatorname{RMSNorm}}
\newcommand{\Attention}{\operatorname{Attention}}

% 图像路径
\graphicspath{{figures/}}

% =========================
% 页眉页脚
% =========================
\pagestyle{fancy}
\fancyhf{}
\fancyhead[LE,RO]{\thepage}
\fancyhead[RE]{\leftmark}
\fancyhead[LO]{\rightmark}

% =========================
% 文档信息
% =========================
\title{从基础到深渊：AI Infrastructure 全栈指南}
\author{作者名}
\date{\today}

\begin{document}

\frontmatter

\maketitle

\chapter*{前言}
\addcontentsline{toc}{chapter}{前言}

当你第一次在命令行敲下 \texttt{torchrun --nproc\_per\_node=8 train.py}，看到八张 GPU 同时开始吐出 loss 曲线时，也许会感到一种强烈的兴奋：大模型训练，终于跑起来了。但兴奋过后，问题很快浮现——显存不够、梯度 NaN、通信效率低下、推理延迟抖动、KV Cache 占满、集群调度混乱……你发现，真正让大模型跑起来并且跑得好的，不只是写好一行 \texttt{model.forward()}，而是一整套从数学、硬件、通信到调度都紧密耦合的系统工程。

这正是 AI Infrastructure 所处的位置。它不是某个单一框架的配置手册，也不是 GPU 编程的入门教程，更不是模型算法的论文综述。AI Infra 是那条把数学公式落到张量算子、把算子落到 CUDA kernel、把 kernel 落到 GPU 硬件、把多 GPU 通信落到网络拓扑、把整个训练和推理流程落到集群调度的长链条。这条链条上的每一环，都可能成为瓶颈；而理解整条链条，才是从"能用"走向"好用"的关键。

\section*{为什么写这本书}

近年来，大语言模型的发展速度惊人。GPT、LLaMA、Qwen、DeepSeek 等模型不断刷新人们对 AI 能力的认知，开源生态也让越来越多的工程师有机会接触、微调、部署甚至训练这些模型。然而，与模型能力的快速发展相比，AI Infrastructure 的知识体系仍然分散在论文、博客、代码仓库和内部文档中，缺乏一份系统性的、从基础到前沿的完整指南。

很多工程师的典型学习路径是这样的：先从 PyTorch 入门，学会写训练脚本；然后遇到显存不够，开始查 ZeRO 和 FSDP 的配置；再遇到训练不稳定，开始搜索 loss spike 和 gradient overflow 的帖子；接着尝试部署推理服务，发现 KV Cache 和 continuous batching 的概念从未系统学过；最终面对集群调度和网络拓扑，才意识到自己对 RDMA 和 NVLink 几乎一无所知。

这条路径的特点是\textbf{碎片化}——每个问题都是单独搜索、单独解决，缺乏全局视角。工程师可能在某个具体问题上积累了经验，却难以把不同层级的知识串联起来。一个熟悉 NCCL 调试的人，可能不理解 FlashAttention 为什么是 IO 优化而非近似算法；一个熟悉 vLLM 配置的人，可能不知道 PagedAttention 的 block table 与操作系统页表之间的类比；一个熟悉 tensor parallel 的人，可能不清楚 pipeline parallel 的 bubble 与 micro-batch 数量的精确关系。

本书的目标，正是填补这条知识链条上的空白，帮助读者建立\textbf{从数学公式到硬件成本的完整心智模型}。

\section*{这本书的核心理念}

贯穿全书的一条主线是：

\begin{center}
\textbf{数学对象 $\rightarrow$ 张量算子 $\rightarrow$ CUDA kernel $\rightarrow$ GPU 硬件 $\rightarrow$ 多 GPU 通信 $\rightarrow$ 集群调度}
\end{center}

我们不孤立地讲任何一个环节。每一章都会追问：这个概念在上面是什么数学形式，在下面是什么硬件成本，在旁边是什么系统 trade-off。

具体来说，本书坚持三个原则：

第一，\textbf{不只讲"怎么做"，更要讲"为什么"。}很多框架文档会告诉你 ZeRO Stage 3 怎么配置，但不会告诉你为什么 optimizer state 的显存是 $16N$ bytes，为什么 reduce-scatter 比 all-reduce 更适合梯度分片，为什么 ZeRO Stage 3 在某些场景下反而比 Stage 2 更慢。本书在每一个关键决策点，都会从数学和系统两个层面解释因果。

第二，\textbf{不回避底层，也不脱离上层。}我们会深入 CUDA 的 grid/block/thread、shared memory 和 Tensor Core，因为不理解这些就无法解释 FlashAttention 的 tiling 策略和 PagedAttention 的访存模式。但我们不会停留在 kernel 层面——每个 kernel 优化都要回到模型结构和系统成本来解释意义。

第三，\textbf{把 trade-off 显式化。}AI Infra 中几乎没有"免费的优化"。降低显存往往增加通信或计算；提高吞吐往往牺牲延迟或稳定性；简化调度往往损失灵活性。本书在每个关键选择处，都会明确写出代价和适用条件，而不是给出一个"最佳配置"就结束。

\section*{本书的结构}

全书分为五篇，外加附录：

\begin{itemize}
\item \textbf{第一篇：基石}——线性代数、微积分、概率论与 CUDA C++ 基础、深度学习基本原理。这是整本书的地基：从 shape 推导显存，从矩阵乘法推导 FLOPs，从链式法则推导反向传播与激活保存，从 softmax 推导数值稳定性，从 GPU 编程模型推导 kernel 性能。
\item \textbf{第二篇：演进}——Transformer 架构深度剖析、大语言模型发展史与主流范式、GPT/LLaMA 与现代 Decoder-only 架构。这是把地基上的数学对象组装成真实模型：Attention、FFN、RoPE、RMSNorm、SwiGLU、GQA、MoE，以及它们对训练和推理系统的具体影响。
\item \textbf{第三篇：分布式训练架构}——数据并行、张量并行、流水线并行、3D 并行策略、ZeRO 与 DeepSpeed。这是回答"大模型放不下、训不动"的核心篇章：从单卡瓶颈出发，逐步推导每种并行策略的通信模式、显存收益和适用边界，最终给出 3D 并行配置的推导方法。
\item \textbf{第四篇：推理架构优化}——KV Cache、PagedAttention 与 vLLM、FlashAttention、量化、Continuous Batching。这是回答"如何让模型高吞吐、低延迟地服务"的核心篇章：从自回归推理的 prefill/decode 两阶段出发，逐步分析 KV Cache 管理、内存碎片、IO 瓶颈、精度压缩和动态调度。
\item \textbf{第五篇：集群网络与硬件调度}——GPU 架构、NVLink/NVSwitch、RDMA/RoCE、K8s AI 调度。这是回到物理世界：理解 GPU 芯片内部结构、节点内互联、节点间网络和集群级调度，让前面的分布式策略和推理优化落到真实硬件上。
\end{itemize}

附录提供符号表、CUDA 环境配置指南、集群部署 checklist 和参考文献。

\section*{读者对象}

本书面向三类读者：

\begin{itemize}
\item \textbf{AI Infra 工程师}：日常工作中需要配置分布式训练、优化推理服务、调试 NCCL 通信、分析 GPU profiling 数据的工程师。本书帮助他们把零散经验整合为系统性理解。
\item \textbf{算法工程师转向 Infra}：已经熟悉模型结构和训练流程，但希望深入理解"为什么 OOM"、"为什么训练发散"、"为什么推理慢"的根本原因，从而更好地与 Infra 团队协作或独立解决系统问题的工程师。
\item \textbf{系统工程师转向 AI}：已经熟悉分布式系统、网络、调度和存储，但需要理解模型计算特征、显存模式和通信需求的工程师，从而把传统系统经验有效迁移到 AI 工作负载上。
\end{itemize}

无论你属于哪一类，本书都假设你具备基本的 Python 编程能力，了解深度学习的基本概念（神经网络、训练、梯度），并有过至少一次使用 GPU 运行模型的经验。不需要你预先精通线性代数、CUDA 或分布式系统——这些正是本书要从 AI Infra 角度重新教给你的。

\section*{如何使用这本书}

本书的章节之间有较强的依赖关系。第一篇是后续所有篇章的基础；第二篇是理解训练和推理篇章中模型结构的前提；第三篇和第四篇相对独立，你可以根据自己的工作重点选择先读训练篇或推理篇；第五篇为第三篇和第四篇提供硬件支撑，建议在读完至少一篇并行策略后再读。

如果你是完全的新手，建议从头到尾顺序阅读。如果你已经有分布式训练经验，可以跳过第一篇的部分章节，直接进入第三篇，但在遇到不理解的数学推导或 kernel 概念时回到第一篇补充。如果你主要做推理服务，可以先读第二篇的架构章节和第四篇，再回到第三篇了解训练侧的并行策略——因为很多推理优化（如 tensor parallel、KV Cache 管理）的训练侧对应知识同样重要。

具体的阅读路线建议，请参阅"本书阅读路线"。

\section*{致谢}

这本书的写作离不开许多人的帮助。感谢所有在开源社区中分享训练配置、推理优化、kernel 实现和调试经验的前辈和同行——你们的实践是本书最重要的素材来源。感谢 NVIDIA、PyTorch、DeepSpeed、vLLM、Megatron-LM 等团队的开源工作和文档——没有这些基础设施，AI Infra 的知识体系无法如此快速积累。感谢每一位在 GPU 集群上熬过深夜调试 loss spike、NCCL timeout 和 OOM 的工程师——你们踩过的坑，本书试图系统化地记录下来，让后来者少走弯路。

最后，感谢我的家人和朋友们在写作期间的耐心与支持。

\begin{flushright}
作者\\
\today
\end{flushright}

\chapter*{本书阅读路线}
\addcontentsline{toc}{chapter}{本书阅读路线}

本书由五篇组成，各篇之间有依赖关系，也有相对独立的模块。不同背景和目标的读者，可以参照以下路线图规划阅读顺序。

\section*{总体结构}

\begin{center}
\begin{tikzpicture}[
    font=\small,
    part/.style={draw, rounded corners=4pt, minimum width=4.2cm, minimum height=1.0cm, align=center, thick},
    arrow/.style={-{Latex[length=3mm]}, thick}
  ]

  \node[part, fill=blue!10] (p1) at (0,4.0) {第一篇：基石\\数学 $\cdot$ CUDA $\cdot$ 深度学习};
  \node[part, fill=green!10] (p2) at (0,2.2) {第二篇：演进\\Transformer $\cdot$ LLM 架构};
  \node[part, fill=orange!12] (p3) at (-3.5,0.2) {第三篇：分布式训练\\DP $\cdot$ TP $\cdot$ PP $\cdot$ 3D $\cdot$ ZeRO};
  \node[part, fill=red!10] (p4) at (3.5,0.2) {第四篇：推理优化\\KV Cache $\cdot$ PagedAttn $\cdot$ FlashAttn $\cdot$ 量化};
  \node[part, fill=purple!10] (p5) at (0,-1.8) {第五篇：硬件与调度\\GPU $\cdot$ NVLink $\cdot$ RDMA $\cdot$ K8s};

  \draw[arrow] (p1) -- (p2);
  \draw[arrow] (p2) -- (p3);
  \draw[arrow] (p2) -- (p4);
  \draw[arrow] (p3) -- (p5);
  \draw[arrow] (p4) -- (p5);

  \node[align=left, text width=8cm] at (0,-3.2) {
    第一篇是全书基础；第二篇连接基础与系统；\\
    第三、四篇相对独立，可按需选择；第五篇提供硬件支撑。
  };
\end{tikzpicture}
\end{center}

\section*{路线一：系统入门路线}

适合\textbf{刚接触 AI Infra} 的读者，从零开始建立完整知识体系。

\begin{enumerate}
\item \textbf{第一篇：基石}（第 1--3 章）——顺序阅读。这是全书的数学和编程基础，后续所有章节都会引用其中的概念。重点关注：shape 与显存估算、矩阵乘法与 FLOPs、链式法则与反向传播、CUDA 线程层次与内存模型、混合精度训练。
\item \textbf{第二篇：演进}（第 4--6 章）——顺序阅读。理解 Transformer 的每个组件如何从数学公式变成系统成本。重点关注：Attention 的 $S^2$ 瓶颈、FFN 的参数量来源、RoPE 与 KV Cache 的关系、GQA 与推理效率。
\item \textbf{第三篇：分布式训练}（第 7--11 章）——顺序阅读。从单卡瓶颈出发，逐步理解为什么需要分布式、每种并行策略解决什么问题。重点关注：Ring AllReduce 的通信量推导、TP 的 column/row splitting、PP 的 bubble 分析、3D 并行配置推导、ZeRO 各 Stage 的显存与通信 trade-off。
\item \textbf{第四篇：推理优化}（第 12--16 章）——顺序阅读。从自回归推理流程出发，逐步理解推理系统的每个优化层次。重点关注：KV Cache 的显存估算、PagedAttention 的内存管理、FlashAttention 的 IO 优化本质、量化对精度和吞吐的影响、Continuous Batching 的调度逻辑。
\item \textbf{第五篇：硬件与调度}（第 17--20 章）——按需阅读。回到物理世界，理解分布式策略和推理优化的硬件基础。
\end{enumerate}

\section*{路线二：训练工程师路线}

适合\textbf{已有模型训练经验}，需要深入理解分布式训练架构的读者。你可以快速回顾第一篇的关键章节，然后集中精力在第二篇和第三篇。

\begin{enumerate}
\item \textbf{快速回顾}：第 1 章的显存/FLOPs/通信量估算、第 2 章的 CUDA 内存模型与 kernel 性能分析、第 3 章的混合精度训练与优化器状态。
\item \textbf{重点阅读}：第 4 章的 Attention 系统成本、第 6 章的 GQA/SwiGLU/MoE 架构对训练的影响。
\item \textbf{核心章节}：第 7--11 章全部。这是训练工程师的核心知识——数据并行、张量并行、流水线并行、3D 并行策略和 ZeRO。建议仔细推导每种策略的通信量和显存收益，不要只记结论。
\item \textbf{补充阅读}：第 12 章的 KV Cache 基础和第 14 章的 FlashAttention，因为训练中同样需要高效 attention kernel。第 17--19 章的硬件与网络知识，帮助你理解并行策略在真实集群上的表现。
\end{enumerate}

\section*{路线三：推理工程师路线}

适合\textbf{主要做推理服务部署与优化}的读者。你可以跳过训练篇的大部分内容，但第二篇的架构知识和第四篇是核心。

\begin{enumerate}
\item \textbf{快速回顾}：第 1 章的显存估算与数值稳定性、第 2 章的 GPU 内存层级与访存模式、第 3 章的混合精度格式。
\item \textbf{重点阅读}：第 4 章的 Attention mask 与 causal 结构、第 6 章的 GQA/RoPE/SwiGLU 架构——这些直接影响 KV Cache 设计、attention kernel 实现和 serving 配置。
\item \textbf{核心章节}：第 12--16 章全部。这是推理工程师的核心知识——KV Cache 机制与显存估算、PagedAttention 与 vLLM 架构、FlashAttention 的 IO 优化原理、量化方法与精度-吞吐权衡、Continuous Batching 与调度策略。
\item \textbf{补充阅读}：第 8 章的张量并行（推理中的 TP 通信与权重切分方式与训练一致）、第 17 章的 GPU 架构（decode 阶段的 memory-bound 特性与 Tensor Core 利用率）、第 18 章的 NVLink/NVSwitch（多卡推理的互联基础）。
\end{enumerate}

\section*{路线四：系统工程师路线}

适合\textbf{已有分布式系统/网络/调度经验}，需要理解 AI 工作负载特征的读者。

\begin{enumerate}
\item \textbf{快速回顾}：第 1 章的显存/计算量/通信量估算公式、第 3 章的优化器状态与混合精度。
\item \textbf{理解负载}：第 4--6 章。重点理解 Transformer 的计算模式（GEMM 密集 vs memory-bound 小算子）、LLM 训练与推理的形态差异、MoE 的 all-to-all 通信特征。
\item \textbf{理解并行}：第 7--11 章。重点不是配置方法，而是通信模式：Ring AllReduce、reduce-scatter、all-gather、all-to-all 的通信量与拓扑关系；TP/PP/DP 的 group 划分与 rank 映射；ZeRO 的分片策略与 checkpoint 写入模式。
\item \textbf{理解推理}：第 12--13 章、第 16 章。重点理解 prefill 与 decode 的差异、KV Cache 的显存与带宽瓶颈、continuous batching 的调度逻辑。
\item \textbf{核心章节}：第 17--20 章。GPU 架构、NVLink/NVSwitch 拓扑、RDMA/RoCE 通信、K8s AI 调度——这是系统工程师最能发挥专业优势的领域。
\end{enumerate}

\section*{各章依赖关系}

\begin{center}
\begin{tikzpicture}[
    font=\small,
    ch/.style={draw, rounded corners=2pt, minimum width=2.0cm, minimum height=0.55cm, align=center, fill=blue!6},
    arrow/.style={-{Latex[length=2mm]}, thick}
  ]

  \node[ch, fill=blue!10] (c1) at (0,6.0) {第1章\\数学基础};
  \node[ch, fill=blue!10] (c2) at (3.5,6.0) {第2章\\CUDA C++};
  \node[ch, fill=blue!10] (c3) at (7.0,6.0) {第3章\\深度学习};

  \node[ch, fill=green!10] (c4) at (1.75,4.5) {第4章\\Transformer};
  \node[ch, fill=green!10] (c5) at (5.25,4.5) {第5章\\LLM 历史};

  \node[ch, fill=green!10] (c6) at (3.5,3.0) {第6章\\LLM 架构};

  \node[ch, fill=orange!12] (c7) at (-1.5,1.5) {第7章\\数据并行};
  \node[ch, fill=orange!12] (c8) at (1.75,1.5) {第8章\\张量并行};
  \node[ch, fill=orange!12] (c9) at (5.0,1.5) {第9章\\流水线并行};

  \node[ch, fill=orange!12] (c10) at (1.75,0.0) {第10章\\3D并行};
  \node[ch, fill=orange!12] (c11) at (5.0,0.0) {第11章\\ZeRO};

  \node[ch, fill=red!10] (c12) at (8.5,3.0) {第12章\\KV Cache};
  \node[ch, fill=red!10] (c13) at (11.5,3.0) {第13章\\PagedAttn};
  \node[ch, fill=red!10] (c14) at (8.5,1.5) {第14章\\FlashAttn};
  \node[ch, fill=red!10] (c15) at (11.5,1.5) {第15章\\量化};
  \node[ch, fill=red!10] (c16) at (10.0,0.0) {第16章\\Cont.Batching};

  \node[ch, fill=purple!10] (c17) at (3.5,-1.5) {第17章\\GPU架构};
  \node[ch, fill=purple!10] (c18) at (6.5,-1.5) {第18章\\NVLink};
  \node[ch, fill=purple!10] (c19) at (9.5,-1.5) {第19章\\RDMA};
  \node[ch, fill=purple!10] (c20) at (12.5,-1.5) {第20章\\K8s调度};

  \draw[arrow] (c1) -- (c2);
  \draw[arrow] (c2) -- (c3);
  \draw[arrow] (c1) -- (c4);
  \draw[arrow] (c3) -- (c4);
  \draw[arrow] (c4) -- (c5);
  \draw[arrow] (c5) -- (c6);
  \draw[arrow] (c4) -- (c6);

  \draw[arrow] (c6) -- (c7);
  \draw[arrow] (c6) -- (c8);
  \draw[arrow] (c6) -- (c9);
  \draw[arrow] (c7) -- (c10);
  \draw[arrow] (c8) -- (c10);
  \draw[arrow] (c9) -- (c10);
  \draw[arrow] (c7) -- (c11);
  \draw[arrow] (c10) -- (c11);

  \draw[arrow] (c6) -- (c12);
  \draw[arrow] (c12) -- (c13);
  \draw[arrow] (c2) -- (c14);
  \draw[arrow] (c12) -- (c14);
  \draw[arrow] (c6) -- (c15);
  \draw[arrow] (c13) -- (c16);
  \draw[arrow] (c14) -- (c16);

  \draw[arrow] (c2) -- (c17);
  \draw[arrow] (c8) -- (c17);
  \draw[arrow] (c10) -- (c18);
  \draw[arrow] (c17) -- (c18);
  \draw[arrow] (c7) -- (c19);
  \draw[arrow] (c18) -- (c19);
  \draw[arrow] (c19) -- (c20);
\end{tikzpicture}
\end{center}

\section*{符号与约定}

本书使用如下符号约定，详细符号表见附录 A：

\begin{itemize}
\item 标量用小写斜体：$x,y,z,n,b,s$；
\item 向量用粗体小写：$\vect{x},\vect{y},\vect{b}$；
\item 矩阵用粗体大写：$\mat{X},\mat{W},\mat{A}$；
\item 高阶张量用粗体大写加特殊标记：$\tensor{X}\in\mathbb{R}^{B\times S\times H}$；
\item $B$ 为 batch size，$S$ 为 sequence length，$H$ 为 hidden size，$V$ 为词表大小，$L$ 为层数；
\item $N$ 为模型参数量，FLOPs 表示浮点运算次数，$1\text{ FLOP}=1$ 次乘法或加法；
\item 显存以 bytes 为单位，带宽以 bytes/s 或 GB/s 为单位；
\item 代码示例使用 Python（PyTorch 风格）或 CUDA C++。
\end{itemize}

\section*{实践建议}

本书不是一本只需要"读"的书。AI Infra 的很多直觉只能通过动手获得。以下建议适用于所有读者：

\begin{itemize}
\item \textbf{运行代码示例}：书中的代码都可以在单卡 GPU 上运行。如果你没有 GPU，可以使用 Google Colab 或云上 GPU 实例。亲手观察 shape、stride、显存和 kernel 时间，比单纯阅读公式有效得多。
\item \textbf{做显存估算练习}：选择一个你熟悉的模型（如 LLaMA-7B），用第 1 章的公式估算训练和推理的显存，然后与实际运行结果对比。差距会帮助你理解哪些因素被简化了。
\item \textbf{用 Nsight Systems 做一次 profiling}：找一个训练或推理任务，用 Nsight Systems 拍一张时间线截图。识别 forward、backward、communication、optimizer step 各占多少时间。这个习惯会让你对系统瓶颈产生直觉。
\item \textbf{尝试不同的并行配置}：如果你有集群条件，尝试用不同的 TP/PP/DP 组合运行同一个模型，观察吞吐、显存和通信时间的变化。
\item \textbf{部署一个推理服务}：用 vLLM 或 TensorRT-LLM 部署一个模型，用压测工具模拟多用户并发请求，观察 KV Cache 使用率、batch 动态变化和延迟分布。
\end{itemize}

AI Infra 是一门实践性极强的学科。希望你在阅读本书的同时，始终有一张 GPU 在旁边运行着代码。

\tableofcontents

\mainmatter

% ==================================================
% 第一篇：基石 Foundations
% ==================================================
\part{基石：从数学、CUDA 到深度学习}

\chapter{线性代数与微积分回顾}
\label{chap:math-foundations}

\section{为什么 AI Infra 工程师仍然需要数学}
\label{sec:why-math-for-ai-infra}

很多工程师第一次接触 AI Infrastructure 时，会误以为这里的核心能力只是“会部署 GPU 集群”“会调 NCCL 环境变量”“会把模型用 vLLM 跑起来”。这些能力当然重要，但它们只是表层现象。真正决定一个工程师能否从“调参执行者”成长为“系统架构师”的，是他能否在如下三个层面之间自由切换：

\begin{itemize}
  \item \textbf{数学层面}：模型的计算究竟是什么？矩阵乘法、归一化、注意力、损失函数和梯度传播在公式上如何定义？
  \item \textbf{系统层面}：这些数学操作如何变成张量算子、计算图、kernel launch、通信 collective 和调度策略？
  \item \textbf{硬件层面}：这些算子最终如何落到 GPU 的 register、shared memory、HBM、NVLink、RDMA 网络与存储系统之上？
\end{itemize}

如果只看系统而不看数学，我们会知道某个训练任务 ``OOM''，却不知道到底是参数、梯度、优化器状态还是激活占用了显存。如果只看数学而不看系统，我们可以写出漂亮的 Transformer 公式，却无法解释为什么同样的模型在不同 batch size、sequence length、tensor parallel size 下吞吐差异巨大。如果只看硬件而不看模型结构，我们会知道 HBM 带宽很贵，却难以判断一个 attention kernel 为什么是 memory-bound，而一个 GEMM kernel 为什么更可能是 compute-bound。

\textbf{AI Infra 的很多问题，本质上是数学对象在工程系统中的生命周期管理问题。}例如：

\begin{itemize}
  \item 一个参数矩阵 $\mat{W}\in\mathbb{R}^{d_{\mathrm{in}}\times d_{\mathrm{out}}}$ 在训练中不只是一个矩阵，它还伴随梯度 $\nabla \mat{W}$、优化器状态 $\mat{m},\mat{v}$、混合精度 master weight、checkpoint 分片与通信 buffer。
  \item 一个隐藏状态张量 $\tensor{X}\in\mathbb{R}^{B\times S\times H}$ 在公式中只是 batch、sequence、hidden 三个维度，但在工程中它有具体的 stride、memory layout、是否 contiguous、是否被切分到多个 GPU。
  \item 一个 softmax 公式看起来简单，但在长上下文 attention 中，若直接物化 $S\times S$ 的 attention matrix，会带来巨大的 HBM 读写压力，这正是 FlashAttention 这类算法诞生的根源。
\end{itemize}

本章不是为了把线性代数、微积分和概率论重新讲成一门数学课，而是从 AI Infrastructure 的视角重新组织这些知识。我们关心的不是“会不会证明定理”，而是：

\begin{itemize}
  \item 如何从 shape 推导显存占用；
  \item 如何从矩阵乘法推导计算量和通信量；
  \item 如何从链式法则理解反向传播与激活保存；
  \item 如何从 softmax 和交叉熵理解语言模型训练的数值稳定性；
  \item 如何从张量 layout 判断算子是否容易被 GPU 高效执行。
\end{itemize}

\section{向量、矩阵与张量}
\label{sec:vectors-matrices-tensors}

\subsection{向量空间与基}
\label{subsec:vector-space-basis}

向量可以被看作一组有序数字，也可以被看作空间中的一个点或方向。在深度学习中，我们更常使用前一种视角：一个 token embedding 是一个向量，一个神经网络隐藏状态是一个向量，一行权重也是一个向量。

设 $\vect{x}\in\mathbb{R}^{d}$，我们通常写作：
\[
  \vect{x}
  =
  \begin{bmatrix}
    x_1 & x_2 & \cdots & x_d
  \end{bmatrix}^{\mathsf{T}}.
\]
如果给定一组基向量 $\{\vect{e}_1,\vect{e}_2,\ldots,\vect{e}_d\}$，则任意向量都可以表示为：
\[
  \vect{x}=\sum_{i=1}^{d}x_i\vect{e}_i.
\]

在语言模型中，一个 token 被 tokenizer 映射为离散 id，然后通过 embedding table 查表得到一个连续向量。若词表大小为 $V$，隐藏维度为 $H$，embedding table 可写为：
\[
  \mat{E}\in\mathbb{R}^{V\times H}.
\]
当 token id 为 $t$ 时，对应 embedding 为：
\[
  \vect{x}=\mat{E}_{t,:}\in\mathbb{R}^{H}.
\]

这个查表操作看似不是矩阵乘法，但也可以等价地看作 one-hot 向量与 embedding matrix 相乘。令 $\vect{o}_t\in\mathbb{R}^{V}$ 是第 $t$ 个位置为 $1$ 的 one-hot 向量，则：
\[
  \vect{x}=\vect{o}_t^{\mathsf{T}}\mat{E}.
\]
工程上不会真的构造 one-hot 向量，因为那会浪费大量内存与计算；框架通常用 gather 操作直接读取 $\mat{E}$ 的第 $t$ 行。这是本书反复出现的一个主题：\textbf{数学表达强调统一性，工程实现强调避免无效计算和无效访存。}

\subsection{矩阵乘法的几何意义}
\label{subsec:matrix-multiplication-geometry}

矩阵乘法是深度学习中最重要的计算模式。一个线性层可以写为：
\[
  \mat{Y}=\mat{X}\mat{W}+\vect{b},
\]
其中：
\[
  \mat{X}\in\mathbb{R}^{B\times d_{\mathrm{in}}},
  \quad
  \mat{W}\in\mathbb{R}^{d_{\mathrm{in}}\times d_{\mathrm{out}}},
  \quad
  \mat{Y}\in\mathbb{R}^{B\times d_{\mathrm{out}}}.
\]

从几何上看，矩阵 $\mat{W}$ 将输入空间中的向量映射到输出空间；从工程上看，矩阵乘法是一个高度规则的三重循环：
\[
  Y_{ij}=\sum_{k=1}^{d_{\mathrm{in}}}X_{ik}W_{kj}.
\]
其乘加操作数量近似为：
\[
  \mathrm{FLOPs}\approx 2B\,d_{\mathrm{in}}\,d_{\mathrm{out}},
\]
这里乘法和加法各算一次浮点操作。

在 GPU 上，矩阵乘法之所以高效，是因为它具有三个重要特征：

\begin{itemize}
  \item \textbf{规则访存}：输入和权重以规则块状方式读取，适合 coalesced memory access。
  \item \textbf{高算术强度}：同一块数据可以被多次复用，单位字节访存对应较多计算。
  \item \textbf{易于分块}：可以将大矩阵切分成 tile，映射到 CUDA block、warp 和 Tensor Core。
\end{itemize}

然而，大模型系统里并不是所有操作都像 GEMM 一样幸运。LayerNorm、softmax、embedding lookup、KV Cache 读取等操作往往具有较低算术强度，容易受内存带宽限制。理解矩阵乘法的规律，有助于我们在后续章节中区分 \textbf{compute-bound} 与 \textbf{memory-bound} 算子。

\begin{table}[H]
  \centering
  \small
  \caption{常见深度学习张量操作的数学视角与系统视角}
  \label{tab:math-vs-system-view}
  \begin{tabular}{p{0.20\textwidth}p{0.34\textwidth}p{0.34\textwidth}}
    \toprule
    \textbf{操作} & \textbf{数学视角} & \textbf{系统视角} \\
    \midrule
    Linear / GEMM
    & 矩阵乘法 $\mat{Y}=\mat{X}\mat{W}$
    & 高吞吐 dense compute，适合 Tensor Core，通常追求高 occupancy 和高数据复用。 \\
    \midrule
    LayerNorm
    & 对 hidden 维求均值和方差
    & reduction + elementwise，访存密集，kernel fusion 很重要。 \\
    \midrule
    Softmax
    & 指数归一化 $\exp(x_i)/\sum_j\exp(x_j)$
    & 需要 max 和 sum reduction，数值稳定性与访存模式都关键。 \\
    \midrule
    Attention
    & $\softmax(\mat{Q}\mat{K}^{\mathsf{T}}/\sqrt{d_k})\mat{V}$
    & 若直接物化 attention matrix，会产生巨大中间张量和 HBM IO。 \\
    \midrule
    Embedding
    & one-hot 乘 embedding matrix
    & 实际为 gather，随机访问较多，缓存局部性差异明显。 \\
    \bottomrule
  \end{tabular}
\end{table}

\subsection{张量的 shape、stride 与 memory layout}
\label{subsec:tensor-shape-stride-layout}

张量是向量和矩阵的推广。一个三维张量可以写为：
\[
  \tensor{X}\in\mathbb{R}^{B\times S\times H},
\]
其中 $B$ 是 batch size，$S$ 是 sequence length，$H$ 是 hidden size。对于 Decoder-only LLM，这样的张量通常表示一批 token 序列在某一层 Transformer block 中的隐藏状态。

\textbf{shape} 描述张量每个维度的大小；\textbf{stride} 描述沿某个维度移动一个单位时，底层存储地址需要跳过多少个元素。设一个三维张量采用 row-major 连续布局，即最后一个维度连续，则：
\[
  \mathrm{shape}=(B,S,H),
\]
对应 stride 为：
\[
  \mathrm{stride}=(S\cdot H,\;H,\;1).
\]
元素 $X_{b,s,h}$ 的线性偏移量为：
\[
  \mathrm{offset}(b,s,h)=b\cdot S H+s\cdot H+h.
\]

这条公式看似基础，却是理解 GPU 性能的入口。因为 GPU 访存效率高度依赖相邻线程是否访问相邻地址。如果一个 kernel 让同一个 warp 中的线程沿 hidden 维访问连续元素，就更容易形成合并访存；如果沿 batch 维访问，则相邻线程之间可能相隔 $S\cdot H$ 个元素，访存效率会显著下降。

\begin{figure}[H]
  \centering
  \begin{tikzpicture}[
      font=\small,
      tensor/.style={draw, rounded corners=2pt, minimum width=2.2cm, minimum height=0.8cm, align=center, fill=blue!6},
      cell/.style={draw, minimum width=0.42cm, minimum height=0.42cm},
      arrow/.style={-{Latex[length=2.5mm]}, thick}
    ]

    \node[align=center] at (0,3.0) {\textbf{逻辑张量视图}\\$\tensor{X}\in\mathbb{R}^{B\times S\times H}$};

    \foreach \b/\xshift in {0/0,1/0.45} {
      \begin{scope}[xshift=\xshift cm,yshift=-\xshift cm]
        \draw[fill=blue!4,draw=blue!60] (-2,0) rectangle (1.2,1.8);
        \foreach \i in {0,...,3} {
          \draw[blue!45] (-2,{0.45*\i}) -- (1.2,{0.45*\i});
        }
        \foreach \j in {0,...,7} {
          \draw[blue!45] ({-2+0.4*\j},0) -- ({-2+0.4*\j},1.8);
        }
      \end{scope}
    }

    \node at (-2.35,1.9) {$B$};
    \node at (-2.25,-0.25) {$S$};
    \node at (1.45,0.75) {$H$};

    \node[tensor] (shape) at (4.2,1.4) {shape\\$(B,S,H)$};
    \node[tensor] (stride) at (4.2,0.2) {stride\\$(SH,H,1)$};
    \draw[arrow] (1.4,1.0) -- (shape.west);
    \draw[arrow] (1.4,0.4) -- (stride.west);

    \node[align=center] at (0,-1.2) {\textbf{连续物理存储：最后一维 $H$ 相邻}};
    \foreach \i in {0,...,15} {
      \node[cell, fill=green!6] (c\i) at ({-3.4+0.42*\i},-2.0) {};
    }
    \node[font=\scriptsize] at (-3.4,-2.55) {$X_{0,0,0}$};
    \node[font=\scriptsize] at (-2.14,-2.55) {$X_{0,0,3}$};
    \node[font=\scriptsize] at (-0.04,-2.55) {$X_{0,1,0}$};
    \node[font=\scriptsize] at (2.06,-2.55) {$\cdots$};

    \draw[decorate,decoration={brace,amplitude=4pt},thick]
      (-3.6,-1.65) -- (-2.0,-1.65)
      node[midway,above=6pt,font=\scriptsize] {同一 token 的 hidden 维连续};

    \node[align=left, text width=4.8cm] at (5.2,-1.5) {
       \textbf{工程含义：}\\
       沿 $H$ 维访问通常更友好；\\
       transpose、view、slice 可能改变 stride；\\
       non-contiguous tensor 可能触发额外 copy。
    };
  \end{tikzpicture}
  \caption{张量 shape、stride 与连续内存布局示意图}
  \label{fig:tensor-shape-stride-layout}
\end{figure}

在 PyTorch 中，很多操作并不立即复制数据，而是改变张量的 metadata。例如 \texttt{view}、\texttt{transpose}、\texttt{permute} 可能改变 shape 或 stride。一个常见性能陷阱是：模型代码中反复对张量做转置，然后传给需要连续内存的 kernel，框架不得不隐式调用 \texttt{contiguous()}，导致额外 HBM 读写。

\begin{lstlisting}[language=Python,caption={观察 PyTorch 张量的 shape、stride 与 contiguous 状态},label={lst:pytorch-shape-stride}]
import torch

B, S, H = 2, 4, 8
x = torch.randn(B, S, H, device="cuda")

print("x.shape      =", x.shape)
print("x.stride     =", x.stride())
print("x.contiguous =", x.is_contiguous())

# Transpose sequence and hidden dimensions.
# This usually changes the stride metadata but does not copy data immediately.
y = x.transpose(1, 2)

print("y.shape      =", y.shape)
print("y.stride     =", y.stride())
print("y.contiguous =", y.is_contiguous())

# Some kernels require contiguous inputs.
# contiguous() creates a new physical layout if needed.
z = y.contiguous()

print("z.shape      =", z.shape)
print("z.stride     =", z.stride())
print("z.contiguous =", z.is_contiguous())
\end{lstlisting}

上面的代码在训练框架中非常常见。对于小模型，它带来的额外 copy 可能不明显；但在长上下文 LLM 中，一个形如 $(B,S,H)$ 的激活张量可能有数百 MB，隐式 \texttt{contiguous()} 会成为不可忽视的开销。

\subsection{batch 维度与广播机制}
\label{subsec:batch-broadcasting}

广播机制允许不同 shape 的张量在不显式复制数据的情况下进行逐元素运算。例如：
\[
  \tensor{X}\in\mathbb{R}^{B\times S\times H},
  \quad
  \vect{b}\in\mathbb{R}^{H},
\]
则：
\[
  Y_{b,s,h}=X_{b,s,h}+b_h.
\]
在逻辑上，$\vect{b}$ 被扩展到 $(B,S,H)$；在工程实现中，理想情况下并不会真的复制 $B\cdot S$ 份 bias，而是在 kernel 中通过索引复用同一段内存。

\begin{figure}[H]
  \centering
  \begin{tikzpicture}[
      font=\small,
      box/.style={draw, rounded corners=2pt, align=center, minimum height=0.9cm},
      arrow/.style={-{Latex[length=2.5mm]}, thick},
      smallcell/.style={draw, minimum width=0.48cm, minimum height=0.42cm}
    ]

    \node[box, fill=blue!6, minimum width=3.0cm] (x) at (-4.4,1.4)
      {$\tensor{X}$\\$(B,S,H)$};
    \node[box, fill=orange!10, minimum width=2.3cm] (b) at (-4.4,-0.1)
      {$\vect{b}$\\$(H)$};

    \node[box, fill=green!8, minimum width=3.2cm] (y) at (3.9,0.65)
      {$\tensor{Y}=\tensor{X}+\vect{b}$\\$(B,S,H)$};

    \draw[arrow] (x.east) -- node[above] {elementwise add} (y.west);
    \draw[arrow] (b.east) -- node[below] {broadcast on $B,S$} (y.west);

    \node[align=center] at (-0.25,2.6) {\textbf{广播不是复制，而是索引规则}};
    \foreach \i in {0,...,5} {
      \node[smallcell, fill=orange!10] (bias\i) at ({-1.5+0.5*\i},1.5) {};
    }
    \node[font=\scriptsize] at (-1.5,1.0) {$b_0$};
    \node[font=\scriptsize] at (1.0,1.0) {$b_{H-1}$};

    \foreach \row/\yy in {0/0.25,1/-0.25,2/-0.75} {
      \foreach \i in {0,...,5} {
        \node[smallcell, fill=blue!5] at ({-1.5+0.5*\i},\yy) {};
      }
    }

    \draw[decorate,decoration={brace,amplitude=4pt},thick]
      (-1.75,0.55) -- (1.25,0.55)
      node[midway,above=6pt,font=\scriptsize] {同一 bias 向量被多个 $(b,s)$ 位置复用};

    \node[align=left, text width=5.2cm] at (0.0,-2.0) {
      在 LayerNorm、RMSNorm、bias add、residual add 中，广播非常常见。\\
      好的 kernel 会把广播、激活函数、归一化等 elementwise 操作融合，减少 HBM 往返。
    };
  \end{tikzpicture}
  \caption{batch 维度与广播机制示意图}
  \label{fig:tensor-broadcasting}
\end{figure}

广播虽然方便，但也可能隐藏开销。若某个操作不支持 stride 为 $0$ 或不规则 stride 的广播输入，框架可能会先把广播结果 materialize 成完整张量。对于大模型而言，这类“看不见的张量”经常是性能问题的根源。

\section{矩阵分解与深度学习中的低秩思想}
\label{sec:matrix-factorization-low-rank}

\subsection{特征值与特征向量}
\label{subsec:eigenvalue-eigenvector}

对于方阵 $\mat{A}\in\mathbb{R}^{n\times n}$，如果存在非零向量 $\vect{v}$ 和标量 $\lambda$，使得：
\[
  \mat{A}\vect{v}=\lambda\vect{v},
\]
则称 $\lambda$ 是 $\mat{A}$ 的特征值，$\vect{v}$ 是对应的特征向量。

直观地说，特征向量是在矩阵变换下方向不变的向量，特征值表示该方向被拉伸或压缩的比例。虽然现代大模型工程中很少直接对巨大的权重矩阵做完整特征分解，但这个概念仍然帮助我们理解很多问题：

\begin{itemize}
  \item 优化中的 Hessian 特征值分布影响训练稳定性；
  \item 梯度下降在不同曲率方向上的收敛速度不同；
  \item 模型压缩和低秩近似依赖“主要方向”与“次要方向”的区分；
  \item LoRA 等参数高效微调方法假设任务增量可以落在较低维子空间中。
\end{itemize}

\subsection{SVD 分解}
\label{subsec:svd}

对于任意矩阵 $\mat{A}\in\mathbb{R}^{m\times n}$，奇异值分解写作：
\[
  \mat{A}=\mat{U}\mat{\Sigma}\mat{V}^{\mathsf{T}},
\]
其中：
\[
  \mat{U}\in\mathbb{R}^{m\times m},\quad
  \mat{\Sigma}\in\mathbb{R}^{m\times n},\quad
  \mat{V}\in\mathbb{R}^{n\times n}.
\]
矩阵 $\mat{\Sigma}$ 的对角线元素 $\sigma_1,\sigma_2,\ldots$ 称为奇异值，通常按从大到小排列：
\[
  \sigma_1\geq\sigma_2\geq\cdots\geq 0.
\]

如果只保留最大的 $r$ 个奇异值，可以得到 rank-$r$ 近似：
\[
  \mat{A}\approx \mat{U}_r\mat{\Sigma}_r\mat{V}_r^{\mathsf{T}}.
\]
这个近似在 Frobenius 范数意义下是最优的：
\[
  \mat{A}_r
  =
  \arg\min_{\operatorname{rank}(\mat{B})\le r}
  \|\mat{A}-\mat{B}\|_F.
\]

工程上，完整 SVD 对大模型权重通常过于昂贵，但“低秩结构”的思想非常重要。它告诉我们：如果一个矩阵的有效信息主要集中在少数方向上，就可能用更少的参数表达它。

\subsection{低秩近似与 LoRA 的直觉}
\label{subsec:low-rank-lora}

设线性层权重为：
\[
  \mat{W}\in\mathbb{R}^{d_{\mathrm{in}}\times d_{\mathrm{out}}}.
\]
一次完整微调会直接更新 $\mat{W}$：
\[
  \mat{W}'=\mat{W}+\Delta\mat{W}.
\]
LoRA 的核心思想是：不直接学习完整的 $\Delta\mat{W}$，而是假设它可以由两个低秩矩阵相乘近似：
\[
  \Delta\mat{W}=\mat{B}\mat{A},
\]
其中：
\[
  \mat{A}\in\mathbb{R}^{r\times d_{\mathrm{out}}},
  \quad
  \mat{B}\in\mathbb{R}^{d_{\mathrm{in}}\times r},
  \quad
  r\ll \min(d_{\mathrm{in}},d_{\mathrm{out}}).
\]
参数量从：
\[
  d_{\mathrm{in}}d_{\mathrm{out}}
\]
下降为：
\[
  r(d_{\mathrm{in}}+d_{\mathrm{out}}).
\]

当 $d_{\mathrm{in}}=d_{\mathrm{out}}=4096$ 且 $r=8$ 时，完整矩阵参数量为：
\[
  4096^2=16,777,216,
\]
而低秩增量参数量为：
\[
  8(4096+4096)=65,536.
\]
两者相差 $256$ 倍。这也是 LoRA 在微调系统中极具工程价值的原因：它不仅减少可训练参数，还减少优化器状态、梯度通信和 checkpoint 大小。

\begin{figure}[H]
  \centering
  \begin{tikzpicture}[
      font=\small,
      matbox/.style={draw, rounded corners=2pt, align=center, fill=blue!6},
      lowrank/.style={draw, rounded corners=2pt, align=center, fill=orange!12},
      arrow/.style={-{Latex[length=2.5mm]}, thick}
    ]

    \node[matbox, minimum width=2.8cm, minimum height=3.0cm] (A) at (-4.2,0)
      {$\Delta\mat{W}$\\[2mm]$d_{\mathrm{in}}\times d_{\mathrm{out}}$\\[2mm]\textbf{完整更新}};
    \node at (-1.75,0) {$\approx$};

    \node[lowrank, minimum width=1.0cm, minimum height=3.0cm] (B) at (-0.4,0)
      {$\mat{B}$\\[2mm]$d_{\mathrm{in}}\times r$};
    \node at (0.65,0) {$\times$};
    \node[lowrank, minimum width=3.0cm, minimum height=1.0cm] (C) at (2.55,0)
      {$\mat{A}$\\[2mm]$r\times d_{\mathrm{out}}$};

    \draw[decorate,decoration={brace,amplitude=5pt},thick]
      (-0.95,-1.75) -- (4.1,-1.75)
      node[midway,below=7pt] {低秩瓶颈 $r\ll d$};

    \node[align=left, text width=4.4cm] at (2.3,2.6) {
       \textbf{Infra 视角：}\\
       可训练参数减少；\\
       Adam 状态减少；\\
       梯度同步减少；\\
       checkpoint 与加载成本减少。
    };

    \draw[arrow] (-4.2,-2.15) -- node[below] {从大矩阵更新转为低秩增量学习} (2.55,-2.15);
  \end{tikzpicture}
  \caption{低秩矩阵分解与 LoRA 增量权重的直觉}
  \label{fig:low-rank-lora}
\end{figure}

\subsection{从矩阵分解理解模型压缩}
\label{subsec:model-compression-low-rank}

低秩近似可以用于模型压缩。设某个线性层权重为 $\mat{W}\in\mathbb{R}^{m\times n}$，如果我们用：
\[
  \mat{W}\approx \mat{U}\mat{V},
  \quad
  \mat{U}\in\mathbb{R}^{m\times r},
  \quad
  \mat{V}\in\mathbb{R}^{r\times n},
\]
代替原矩阵，那么原来的前向传播：
\[
  \vect{y}=\vect{x}\mat{W}
\]
可以变为：
\[
  \vect{y}=(\vect{x}\mat{U})\mat{V}.
\]

参数量由 $mn$ 变为 $r(m+n)$。计算量也类似下降。但这并不意味着低秩分解总是让推理更快，原因包括：

\begin{itemize}
  \item 两个小 GEMM 未必比一个大 GEMM 更容易充分利用 Tensor Core；
  \item rank 太小时精度损失明显，rank 太大时压缩收益不足；
  \item 额外 kernel launch、内存读写和中间激活可能抵消部分收益；
  \item 在 tensor parallel 场景中，低秩分解可能改变通信路径。
\end{itemize}

因此，AI Infra 工程师不能只看参数量，还要看 \textbf{实际 kernel 形状、硬件吞吐、访存模式与通信代价}。一个数学上更小的模型，在系统上不一定更快。

\section{微积分与自动微分}
\label{sec:calculus-autodiff}

\subsection{导数、偏导与梯度}
\label{subsec:derivative-gradient}

训练神经网络的核心是最小化损失函数。设参数为 $\vect{\theta}$，训练目标可以写为：
\[
  \min_{\vect{\theta}}\; \mathcal{L}(\vect{\theta}).
\]
梯度下降更新为：
\[
  \vect{\theta}_{t+1}
  =
  \vect{\theta}_t-\eta\nabla_{\vect{\theta}}\mathcal{L}(\vect{\theta}_t),
\]
其中 $\eta$ 是学习率，$\nabla_{\vect{\theta}}\mathcal{L}$ 是损失对参数的梯度。

对于多元函数 $f:\mathbb{R}^{n}\rightarrow\mathbb{R}$，梯度定义为：
\[
  \nabla_{\vect{x}}f
  =
  \begin{bmatrix}
    \frac{\partial f}{\partial x_1}&
    \frac{\partial f}{\partial x_2}&
    \cdots&
    \frac{\partial f}{\partial x_n}
  \end{bmatrix}^{\mathsf{T}}.
\]
梯度方向是函数增长最快的方向，负梯度方向是局部下降最快的方向。

对 AI Infra 而言，梯度不仅是数学对象，也是系统对象。每个参数张量通常都有对应梯度张量；在数据并行训练中，多个 GPU 上的梯度需要通过 all-reduce 或 reduce-scatter 同步；在 ZeRO 中，梯度可能被切分保存；在混合精度训练中，梯度可能还要经历 loss scaling、unscale、overflow check 等步骤。

\subsection{链式法则}
\label{subsec:chain-rule}

反向传播的数学核心是链式法则。若：
\[
  y=f(u),\quad u=g(x),
\]
则：
\[
  \frac{dy}{dx}
  =
  \frac{dy}{du}
  \frac{du}{dx}.
\]

对于向量函数，链式法则表现为 Jacobian 的乘法。设：
\[
  \vect{y}=f(\vect{u}),\quad \vect{u}=g(\vect{x}),
\]
则：
\[
  \frac{\partial \vect{y}}{\partial \vect{x}}
  =
  \frac{\partial \vect{y}}{\partial \vect{u}}
  \frac{\partial \vect{u}}{\partial \vect{x}}.
\]

在神经网络中，计算图由许多局部操作组成。反向传播并不会显式构造完整 Jacobian，而是从输出损失开始，逐个节点传播 \textbf{vector-Jacobian product}。这也是反向模式自动微分高效的原因：当输出是标量 loss，而参数数量巨大时，反向模式只需一次 backward 就能得到所有参数的梯度。

考虑一个两层 MLP：
\[
  \mat{H}=\operatorname{ReLU}(\mat{X}\mat{W}_1+\vect{b}_1),
\]
\[
  \mat{Y}=\mat{H}\mat{W}_2+\vect{b}_2,
\]
\[
  \mathcal{L}=\ell(\mat{Y},\mat{T}).
\]
反向传播会依次计算：
\[
  \frac{\partial \mathcal{L}}{\partial \mat{W}_2},
  \quad
  \frac{\partial \mathcal{L}}{\partial \mat{H}},
  \quad
  \frac{\partial \mathcal{L}}{\partial \mat{W}_1}.
\]
为了计算这些梯度，前向传播中的某些中间结果必须被保存，例如 $\mat{X}$、$\mat{H}$、ReLU mask 等。这就是激活显存的来源。

\begin{figure}[H]
  \centering
  \begin{tikzpicture}[
      font=\small,
      op/.style={draw, rounded corners=2pt, minimum width=1.55cm, minimum height=0.8cm, align=center, fill=blue!6},
      data/.style={draw, rounded corners=2pt, minimum width=1.35cm, minimum height=0.7cm, align=center, fill=green!7},
      grad/.style={-{Latex[length=2.5mm]}, thick, red!70},
      fwd/.style={-{Latex[length=2.5mm]}, thick, blue!70},
      save/.style={draw, dashed, rounded corners=2pt, fill=yellow!10, inner sep=5pt}
    ]

    \node[data] (x) at (-5.4,0) {$\mat{X}$};
    \node[op] (linear1) at (-3.6,0) {Linear\\$\mat{W}_1$};
    \node[op] (relu) at (-1.7,0) {ReLU};
    \node[data] (h) at (0.0,0) {$\mat{H}$};
    \node[op] (linear2) at (1.8,0) {Linear\\$\mat{W}_2$};
    \node[op] (loss) at (3.8,0) {Loss\\$\mathcal{L}$};

    \node[save, fit=(x)(h), inner sep=8pt] (saved) {};
    \node[font=\scriptsize, align=center] at (-2.6,0.95) {前向保存的激活\\用于 backward};

    \draw[fwd] (x) -- (linear1);
    \draw[fwd] (linear1) -- (relu);
    \draw[fwd] (relu) -- (h);
    \draw[fwd] (h) -- (linear2);
    \draw[fwd] (linear2) -- (loss);

    \draw[grad] (loss.south) .. controls (3.0,-1.4) and (2.1,-1.4) .. (linear2.south);
    \draw[grad] (linear2.south) .. controls (1.0,-1.8) and (0.2,-1.8) .. (h.south);
    \draw[grad] (h.south) .. controls (-0.9,-1.8) and (-1.5,-1.8) .. (relu.south);
    \draw[grad] (relu.south) .. controls (-2.4,-1.4) and (-3.0,-1.4) .. (linear1.south);

    \node[align=left, text width=5.8cm] at (0,-2.75) {
      蓝色箭头表示前向计算，红色箭头表示梯度反传。\\
      反向传播不是免费的：它需要前向激活、额外 kernel、梯度张量和通信同步。
    };
  \end{tikzpicture}
  \caption{计算图中的前向路径、激活保存与反向传播路径}
  \label{fig:autograd-computation-graph}
\end{figure}

\subsection{Jacobian 与 Hessian 的工程意义}
\label{subsec:jacobian-hessian}

Jacobian 描述向量输出对向量输入的一阶导数。若 $f:\mathbb{R}^{n}\rightarrow\mathbb{R}^{m}$，则：
\[
  \mat{J}_{ij}
  =
  \frac{\partial f_i}{\partial x_j},
  \quad
  \mat{J}\in\mathbb{R}^{m\times n}.
\]
Hessian 描述标量函数对向量输入的二阶导数：
\[
  \mat{H}_{ij}
  =
  \frac{\partial^2 f}{\partial x_i\partial x_j},
  \quad
  \mat{H}\in\mathbb{R}^{n\times n}.
\]

对于拥有数十亿参数的 LLM，显式构造 Hessian 几乎不现实，因为它需要 $O(n^2)$ 存储。但 Hessian 的思想仍然广泛存在：

\begin{itemize}
  \item Adam 使用一阶梯度的二阶矩估计，虽然不是 Hessian，但体现了不同参数方向尺度不同的事实；
  \item 二阶量化方法会使用 Hessian 或其近似来估计量化误差对 loss 的影响；
  \item 训练稳定性分析经常关心曲率、sharpness、梯度爆炸与梯度消失；
  \item 一些优化器使用 Hessian-vector product，而不是显式 Hessian。
\end{itemize}

Hessian-vector product 可以写为：
\[
  \mat{H}\vect{v}
  =
  \nabla_{\vect{x}}
  \left(
    \nabla_{\vect{x}}f(\vect{x})^{\mathsf{T}}\vect{v}
  \right).
\]
这个公式说明：通过自动微分，我们可以在不显式构造 Hessian 的情况下计算它与向量的乘积。这种“避免物化大矩阵”的思想，与后续 FlashAttention 避免物化 attention matrix 的系统思想非常相似。

\subsection{前向模式与反向模式自动微分}
\label{subsec:forward-reverse-autodiff}

自动微分主要有两种模式：

\begin{itemize}
  \item \textbf{前向模式}：从输入出发，随着计算图向前传播导数，适合输入维度少、输出维度多的情况。
  \item \textbf{反向模式}：从输出出发，沿计算图反向传播梯度，适合输出维度少、输入维度多的情况。
\end{itemize}

深度学习训练通常最小化一个标量 loss，而参数数量巨大，因此反向模式更合适。

设计算图中某个节点为：
\[
  v_i=f_i(v_{p_1},v_{p_2},\ldots).
\]
反向模式维护 adjoint：
\[
  \bar{v}_i=\frac{\partial \mathcal{L}}{\partial v_i}.
\]
若 $v_i$ 是 $v_j$ 的输入，则反向传播根据局部导数累加：
\[
  \bar{v}_i
  \leftarrow
  \bar{v}_i+
  \bar{v}_j
  \frac{\partial v_j}{\partial v_i}.
\]

\begin{lstlisting}[language=Python,caption={极简反向模式自动微分核心逻辑},label={lst:tiny-reverse-mode-autodiff}]
class Node:
    def __init__(self, value, parents=(), backward_fn=None):
        self.value = value
        self.parents = parents
        self.grad = 0.0
        self.backward_fn = backward_fn

def add(a, b):
    out = Node(a.value + b.value, parents=(a, b))

    def backward():
        # d(a + b) / da = 1
        # d(a + b) / db = 1
        a.grad += out.grad
        b.grad += out.grad

    out.backward_fn = backward
    return out

def mul(a, b):
    out = Node(a.value * b.value, parents=(a, b))

    def backward():
        # d(a * b) / da = b
        # d(a * b) / db = a
        a.grad += out.grad * b.value
        b.grad += out.grad * a.value

    out.backward_fn = backward
    return out

def backward(loss, topo_order):
    # Seed gradient: dL / dL = 1.
    loss.grad = 1.0

    # Reverse topological traversal.
    # Each node pushes its gradient to its parents.
    for node in reversed(topo_order):
        if node.backward_fn is not None:
            node.backward_fn()
\end{lstlisting}

真实框架当然远比这段代码复杂：它需要处理张量 shape、广播梯度归约、in-place 操作、混合精度、分布式通信、动态图生命周期、activation checkpointing、CUDA stream 等问题。但反向模式的核心思想仍然如此：\textbf{前向构图，反向沿图传播梯度。}

\subsection{激活保存、重计算与显存权衡}
\label{subsec:activation-recompute-memory}

在训练中，显存主要由以下部分组成：

\begin{itemize}
  \item 参数 parameters；
  \item 梯度 gradients；
  \item 优化器状态 optimizer states；
  \item 前向激活 activations；
  \item 临时 buffer，例如通信 buffer、workspace、attention 中间结果。
\end{itemize}

对于很深的模型，激活显存可能非常可观。若每层保存一个形如 $(B,S,H)$ 的激活，层数为 $L$，数据类型为 FP16，则仅这一项粗略为：
\[
  M_{\mathrm{act}}
  \approx
  L\cdot B\cdot S\cdot H\cdot 2\ \mathrm{bytes}.
\]
实际情况还要乘上若干常数，因为 attention、MLP、norm、dropout 等都会保存不同中间结果。

Gradient Checkpointing 的思想是：前向时只保存部分节点，反向时重新计算缺失的中间激活，以计算换显存。它体现了 AI Infra 中非常典型的 trade-off：
\[
  \text{显存下降}
  \quad\Longleftrightarrow\quad
  \text{额外计算上升}.
\]
在 GPU 集群中，这种 trade-off 还会进一步影响 pipeline bubble、通信重叠和吞吐稳定性。

\section{概率、熵与交叉熵}
\label{sec:probability-entropy-crossentropy}

\subsection{概率分布与采样}
\label{subsec:distribution-sampling}

语言模型的输出不是一个确定的 token，而是下一个 token 的概率分布。设词表大小为 $V$，模型在位置 $t$ 输出 logits：
\[
  \vect{z}_t\in\mathbb{R}^{V}.
\]
通过 softmax 得到概率：
\[
  p(x_{t+1}=i\mid x_{\le t})
  =
  \frac{\exp(z_{t,i})}{\sum_{j=1}^{V}\exp(z_{t,j})}.
\]

训练时，我们通常最大化真实下一个 token 的概率；推理时，我们可以从该分布中采样，也可以选择概率最大的 token。不同采样策略包括：

\begin{itemize}
  \item \textbf{greedy decoding}：总是选择概率最大的 token；
  \item \textbf{temperature sampling}：用温度调节分布尖锐程度；
  \item \textbf{top-k sampling}：只在概率最高的 $k$ 个 token 中采样；
  \item \textbf{top-p sampling}：只在累计概率达到 $p$ 的最小 token 集合中采样。
\end{itemize}

从 Infra 视角看，训练中的 softmax 通常是大 batch、大矩阵上的密集操作；推理 decode 阶段的 softmax 则可能变成许多小 batch、逐 token 的操作。两者对 kernel fusion 和调度系统的要求不同。

\subsection{信息熵、交叉熵与 KL 散度}
\label{subsec:entropy-crossentropy-kl}

离散分布 $p$ 的信息熵定义为：
\[
  H(p)=-\sum_{i}p_i\log p_i.
\]
熵可以理解为不确定性的度量。分布越均匀，熵越大；分布越集中，熵越小。

给定真实分布 $p$ 与模型分布 $q$，交叉熵定义为：
\[
  H(p,q)=-\sum_i p_i\log q_i.
\]
KL 散度定义为：
\[
  D_{\mathrm{KL}}(p\|q)
  =
  \sum_i p_i\log\frac{p_i}{q_i}.
\]
三者关系为：
\[
  H(p,q)=H(p)+D_{\mathrm{KL}}(p\|q).
\]
当真实分布 $p$ 固定时，最小化交叉熵等价于最小化 KL 散度。

在 next-token prediction 中，真实标签通常是 one-hot。若真实 token id 为 $y$，则：
\[
  p_i=
  \begin{cases}
    1, & i=y,\\
    0, & i\ne y.
  \end{cases}
\]
于是交叉熵变为：
\[
  \mathcal{L}
  =
  -\log q_y.
\]
也就是说，语言模型训练的基本目标就是让真实下一个 token 的概率尽可能高。

\subsection{Softmax 与温度参数}
\label{subsec:softmax-temperature}

带温度参数 $T$ 的 softmax 定义为：
\[
  \softmax(\vect{z}/T)_i
  =
  \frac{\exp(z_i/T)}{\sum_j\exp(z_j/T)}.
\]
当 $T<1$ 时，分布更尖锐，高概率 token 更突出；当 $T>1$ 时，分布更平滑，采样更随机。

直接计算 softmax 可能出现数值溢出。若某个 $z_i$ 很大，则 $\exp(z_i)$ 可能超过浮点数范围。稳定做法是减去最大值：
\[
  m=\max_j z_j,
\]
\[
  \softmax(\vect{z})_i
  =
  \frac{\exp(z_i-m)}{\sum_j\exp(z_j-m)}.
\]
因为 softmax 对整体平移不变：
\[
  \frac{\exp(z_i-c)}{\sum_j\exp(z_j-c)}
  =
  \frac{\exp(z_i)\exp(-c)}{\sum_j\exp(z_j)\exp(-c)}
  =
  \frac{\exp(z_i)}{\sum_j\exp(z_j)}.
\]

交叉熵也常通过 log-sum-exp 稳定计算：
\[
  \log\sum_j\exp(z_j)
  =
  m+\log\sum_j\exp(z_j-m).
\]
若真实类别为 $y$，则：
\[
  \mathcal{L}
  =
  -z_y+\log\sum_j\exp(z_j).
\]

\begin{lstlisting}[language=Python,caption={数值稳定的 softmax 与 cross entropy},label={lst:stable-softmax-cross-entropy}]
import torch

def stable_softmax(logits, dim=-1):
    # Subtracting max does not change softmax results.
    # It prevents exp(logits) from overflowing.
    max_value = torch.max(logits, dim=dim, keepdim=True).values
    shifted = logits - max_value
    exp_value = torch.exp(shifted)
    return exp_value / torch.sum(exp_value, dim=dim, keepdim=True)

def stable_cross_entropy(logits, target):
    # logits: [batch, vocab]
    # target: [batch], each entry is the correct token id.
    max_value = torch.max(logits, dim=-1, keepdim=True).values
    shifted = logits - max_value

    # logsumexp(logits) = max + log(sum(exp(logits - max)))
    logsumexp = max_value.squeeze(-1) + torch.log(torch.sum(torch.exp(shifted), dim=-1))

    # Gather z_y for each sample.
    zy = logits.gather(dim=-1, index=target[:, None]).squeeze(-1)

    # Cross entropy for one-hot target: -z_y + logsumexp(z)
    loss = -zy + logsumexp
    return loss.mean()
\end{lstlisting}

在实际训练中，框架通常会使用 fused cross entropy kernel，将 max reduction、sum reduction、log、gather、loss 计算融合，减少中间张量写回 HBM 的次数。对于词表很大的 LLM，logits 张量形如 $(B,S,V)$，其中 $V$ 可能达到数万甚至更多。此时，cross entropy 的实现效率会影响端到端训练吞吐。

\subsection{语言模型中的 next-token prediction}
\label{subsec:next-token-prediction}

给定 token 序列：
\[
  x_1,x_2,\ldots,x_T,
\]
自回归语言模型将联合概率分解为：
\[
  p(x_1,x_2,\ldots,x_T)
  =
  \prod_{t=1}^{T}p(x_t\mid x_{<t}).
\]
训练目标通常是最小化负对数似然：
\[
  \mathcal{L}
  =
  -\sum_{t=1}^{T}\log p(x_t\mid x_{<t}).
\]
在 batch 训练中，若 batch size 为 $B$、序列长度为 $S$，则平均 loss 可以写为：
\[
  \mathcal{L}
  =
  -\frac{1}{BS}
  \sum_{b=1}^{B}
  \sum_{s=1}^{S}
  \log p(x_{b,s}\mid x_{b,<s}).
\]

\begin{figure}[H]
  \centering
  \begin{tikzpicture}[
      font=\small,
      token/.style={draw, rounded corners=2pt, minimum width=0.95cm, minimum height=0.55cm, fill=blue!6},
      bar/.style={draw, fill=orange!25},
      arrow/.style={-{Latex[length=2.4mm]}, thick}
    ]

    \node[token] (t1) at (-4.8,1.5) {从};
    \node[token] (t2) at (-3.7,1.5) {基础};
    \node[token] (t3) at (-2.6,1.5) {到};
    \node[token] (t4) at (-1.5,1.5) {深渊};

    \draw[arrow] (t1) -- (t2);
    \draw[arrow] (t2) -- (t3);
    \draw[arrow] (t3) -- (t4);

    \node[draw, rounded corners=2pt, fill=green!8, minimum width=2.5cm, minimum height=0.9cm, align=center] (model) at (0.8,1.5)
      {Decoder-only\\LLM};
    \draw[arrow] (t4) -- (model);

    \node[align=center] at (4.0,2.25) {\textbf{下一个 token 概率分布}};
    \draw[->] (2.4,0.2) -- (6.1,0.2) node[right] {token};
    \draw[->] (2.4,0.2) -- (2.4,2.0) node[above] {prob};

    \foreach \x/\h/\name in {2.8/0.55/指南,3.4/1.45/架构,4.0/0.75/训练,4.6/0.35/推理,5.2/0.25/GPU} {
      \draw[bar] (\x,0.2) rectangle ({\x+0.35},{0.2+\h});
      \node[font=\scriptsize, rotate=45, anchor=west] at (\x,0.0) {\name};
    }

    \draw[arrow] (model.east) -- (2.4,1.5);

    \node[align=left, text width=5.0cm] at (-1.5,-1.0) {
       训练时：真实下一个 token 给出监督信号。\\
       推理时：根据概率分布采样或选择 token，并把生成 token 追加到上下文。
    };
  \end{tikzpicture}
  \caption{语言模型 next-token prediction 与 token 概率分布}
  \label{fig:next-token-distribution}
\end{figure}

next-token prediction 的形式非常简单，但它对系统提出了巨大挑战：

\begin{itemize}
  \item 训练时，需要对所有位置并行计算 logits 和 loss；
  \item 推理时，decode 阶段必须逐 token 生成，难以像训练那样完全并行；
  \item 长上下文会使 attention 和 KV Cache 的显存需求快速上升；
  \item 大词表会让 logits 与 cross entropy 占用显著带宽；
  \item 多用户在线请求会引入动态 batch、不同长度序列和复杂调度。
\end{itemize}

\section{从数学公式到系统成本}
\label{sec:from-formula-to-system-cost}

本章最后，我们把前面的数学对象和 AI Infra 中最关心的几个系统指标联系起来：显存、计算量、通信量与数值稳定性。

\subsection{显存估算的基本套路}
\label{subsec:memory-estimation-basics}

假设一个张量 shape 为：
\[
  (d_1,d_2,\ldots,d_k),
\]
元素类型占用 $s$ bytes，则张量显存为：
\[
  M=\left(\prod_{i=1}^{k}d_i\right)s.
\]
例如，隐藏状态：
\[
  \tensor{X}\in\mathbb{R}^{B\times S\times H}
\]
使用 FP16 时显存为：
\[
  M_X=B S H \cdot 2\ \mathrm{bytes}.
\]

对于参数量为 $N$ 的模型，若使用混合精度 Adam 训练，一个粗略估算是：

\begin{itemize}
  \item FP16/BF16 参数：$2N$ bytes；
  \item FP16/BF16 梯度：$2N$ bytes；
  \item FP32 master weight：$4N$ bytes；
  \item Adam 一阶矩和二阶矩：$8N$ bytes。
\end{itemize}

仅参数相关部分就可能达到：
\[
  M_{\mathrm{param+grad+optim}}
  \approx
  16N\ \mathrm{bytes}.
\]
对于 $N=70\mathrm{B}$ 的模型，这个数字约为：
\[
  16\times 70\times 10^9
  =
  1.12\times 10^{12}\ \mathrm{bytes}
  \approx
  1.12\ \mathrm{TB}.
\]
这还没有包括激活、临时 buffer、通信 buffer 和碎片。因此，分布式训练不是“可选优化”，而是大模型训练的必要条件。

\subsection{计算量估算的基本套路}
\label{subsec:flops-estimation-basics}

线性层 $\mat{Y}=\mat{X}\mat{W}$ 的 FLOPs 约为：
\[
  2B d_{\mathrm{in}}d_{\mathrm{out}}.
\]
对于 Transformer，粗略估算训练 FLOPs 时，经常使用参数量、token 数和常数因子进行估算：
\[
  \mathrm{FLOPs}
  \approx
  c\cdot N\cdot T,
\]
其中 $N$ 是参数量，$T$ 是训练 token 数，$c$ 是与前向、反向和模型结构相关的常数。这个公式不是精确 kernel 级模型，但足以帮助架构师判断训练任务需要多少 GPU-days。

当我们进入后续章节，会进一步把 FLOPs 分解到 attention、MLP、embedding、norm 等模块，并分析它们在不同 sequence length 下的主导项。

\subsection{通信量估算的基本套路}
\label{subsec:communication-estimation-basics}

数据并行训练中，每张卡有完整模型副本，每次 backward 后需要同步梯度。若参数量为 $N$，梯度使用 FP16/BF16，则每步需要同步的数据量量级为：
\[
  2N\ \mathrm{bytes}.
\]
Ring AllReduce 的每个 rank 通信量近似为：
\[
  2\frac{P-1}{P}M,
\]
其中 $P$ 是数据并行进程数，$M$ 是待同步消息大小。当 $P$ 很大时，通信量接近 $2M$。

张量并行、流水线并行和 ZeRO 会改变通信模式：有些策略减少显存但增加通信，有些策略减少单卡参数但引入更多 all-gather 或 reduce-scatter。后续分布式训练章节会详细分析这些 trade-off。

\subsection{数值稳定性不是细节}
\label{subsec:numerical-stability-is-not-detail}

AI Infra 工程师经常要处理训练发散、loss spike、NaN、overflow、underflow 等问题。这些问题并不总是模型算法本身的问题，也可能来自系统实现：

\begin{itemize}
  \item FP16 指数范围较小，容易 overflow 或 underflow；
  \item 梯度累积时，如果 loss scaling 不当，可能导致溢出；
  \item softmax 如果不减最大值，可能产生 inf；
  \item reduce 操作的顺序不同，会带来浮点舍入差异；
  \item 分布式训练中，不同 rank 上的数值异常可能通过 collective 扩散。
\end{itemize}

因此，数值稳定性不是“数学洁癖”，而是生产系统可靠性的基础。一个成熟的训练系统通常需要监控：

\begin{itemize}
  \item loss 是否出现 spike；
  \item gradient norm 是否异常；
  \item 是否出现 NaN 或 inf；
  \item loss scaler 是否频繁回退；
  \item 不同 rank 的统计量是否一致；
  \item checkpoint 恢复后 loss 曲线是否连续。
\end{itemize}

\section{本章小结}
\label{sec:chapter01-summary}

本章从 AI Infrastructure 的角度回顾了线性代数、微积分和概率论中最重要的概念。我们没有停留在抽象公式，而是反复追问：这些数学对象在 GPU、训练框架和分布式系统中会变成什么成本。

\begin{itemize}
  \item \textbf{向量、矩阵与张量} 是模型计算的基本对象。shape 决定逻辑结构，stride 和 memory layout 决定实际访存模式。
  \item \textbf{矩阵乘法} 是深度学习吞吐的核心来源。它规则、密集、易分块，因此适合 GPU 和 Tensor Core。
  \item \textbf{低秩思想} 连接了矩阵分解、模型压缩和参数高效微调。LoRA 的工程价值来自参数、梯度、优化器状态和 checkpoint 的系统性减少。
  \item \textbf{链式法则与反向模式自动微分} 是训练的数学基础。反向传播不仅产生梯度，也带来激活保存、重计算和显存压力。
  \item \textbf{softmax、交叉熵和 next-token prediction} 构成语言模型训练目标，但它们在大词表和长序列场景下也会形成显著系统开销。
  \item \textbf{显存、计算量、通信量与数值稳定性} 是 AI Infra 工程师理解模型训练和推理系统的四个基本坐标。
\end{itemize}

从下一章开始，我们将进入 CUDA C++ 基础。届时，本章中的张量、矩阵乘法、stride、访存和并行计算会落到更具体的 GPU 编程模型中：grid、block、thread、warp、shared memory、register 和 global memory。换句话说，从下一章开始，我们将把数学对象真正放到硬件上运行。
\chapter{CUDA C++ 基础}
\label{chap:cuda-cpp}

\section{从 CPU 编程模型到 GPU 编程模型}
\label{sec:cpu-to-gpu-programming-model}

上一章我们从张量、矩阵乘法、自动微分和显存估算出发，建立了 AI Infrastructure 工程师理解模型计算成本的数学视角。本章开始，我们把这些数学对象真正放到硬件上执行。这里的核心问题不再只是：
\[
  \mat{C}=\mat{A}\mat{B}
\]
这条公式如何成立，而是进一步追问：

\begin{itemize}
  \item 这个矩阵乘法由多少个线程共同完成？
  \item 每个线程负责哪些元素？
  \item 数据从 HBM 进入 SM 的路径是什么？
  \item 同一份数据能否在 shared memory 中被多个线程复用？
  \item 为什么同样的 FLOPs，在不同 memory layout 下速度可能相差数倍？
\end{itemize}

CUDA C++ 是 NVIDIA GPU 编程的基础模型。即使现代 AI Infra 工程师更多使用 PyTorch、Triton、XLA、TensorRT、vLLM、DeepSpeed 或 Megatron-LM，也仍然需要理解 CUDA 的底层抽象。原因很简单：\textbf{所有高层框架最终都要落到 kernel、memory hierarchy、stream、event、collective communication 与 device runtime 这些基本机制上。}

\subsection{CPU 的强控制流与 GPU 的强吞吐}
\label{subsec:cpu-control-gpu-throughput}

CPU 和 GPU 都能执行通用计算，但它们的设计目标非常不同。CPU 更强调低延迟、复杂控制流、分支预测、乱序执行和大缓存；GPU 更强调高吞吐、大规模并行、简单控制逻辑和海量线程调度。

从 AI Infra 的角度看，可以把二者理解为两类不同的计算哲学：

\begin{itemize}
  \item \textbf{CPU 像少数几位经验丰富的工程师}：每个核心很强，擅长处理复杂逻辑、系统调用、调度、数据预处理、网络协议栈和控制面任务。
  \item \textbf{GPU 像成千上万个执行统一任务的工人}：单个线程能力有限，但数量巨大，适合执行高度规则、数据并行、算术密集的任务。
\end{itemize}

深度学习中的矩阵乘法、卷积、attention、归一化和逐元素操作，大多具有可并行结构。以向量加法为例：
\[
  c_i=a_i+b_i,\quad i=0,1,\ldots,n-1.
\]
每个位置 $i$ 的计算互不依赖，因此可以让大量 GPU 线程同时处理不同元素。这种模式正是 GPU 擅长的。

\begin{figure}[H]
  \centering
  \begin{tikzpicture}[
      font=\small,
      core/.style={draw, rounded corners=2pt, minimum width=0.75cm, minimum height=0.55cm, align=center},
      cache/.style={draw, rounded corners=2pt, minimum width=3.4cm, minimum height=0.45cm, align=center, fill=yellow!15},
      mem/.style={draw, rounded corners=2pt, minimum width=4.2cm, minimum height=0.65cm, align=center, fill=green!8},
      title/.style={font=\bfseries},
      note/.style={align=left, text width=5.0cm},
      arrow/.style={-{Latex[length=2.3mm]}, thick}
    ]

    \node[title] at (-3.9,3.0) {CPU：低延迟与复杂控制};
    \node[title] at (3.9,3.0) {GPU：高吞吐与海量并行};

    \foreach \i/\x in {1/-5.0,2/-4.2,3/-3.4,4/-2.6} {
      \node[core, fill=blue!8] (cpu\i) at (\x,2.0) {Core};
    }
    \node[cache] (l3) at (-3.8,1.15) {Large Cache};
    \node[mem] (dram) at (-3.8,0.25) {System Memory};

    \foreach \i in {0,...,7} {
      \foreach \j in {0,...,3} {
        \node[core, fill=orange!12, minimum width=0.55cm, minimum height=0.38cm]
          at ({1.6+0.58*\i},{2.35-0.42*\j}) {};
      }
    }
    \node[cache, minimum width=4.9cm] (smem) at (3.9,0.95) {Shared Memory / L1 / Register File};
    \node[mem, minimum width=4.9cm] (hbm) at (3.9,0.05) {HBM};

    \node[note] at (-3.8,-1.2) {
      \textbf{典型任务：}\\
      操作系统、调度器、RPC、数据加载、复杂业务逻辑、少量强依赖计算。
    };

    \node[note] at (3.9,-1.2) {
      \textbf{典型任务：}\\
      GEMM、attention、卷积、归一化、向量化逐元素计算、大规模并行训练。
    };

    \draw[arrow] (-3.8,-3.5) -- node[below] {控制流强} (-1.7,-3.5);
    \draw[arrow] (2.0,-3.5) -- node[below] {数据并行强} (5.8,-3.5);
  \end{tikzpicture}
  \caption{CPU 与 GPU 编程模型的直观对比}
  \label{fig:cpu-vs-gpu-programming-model}
\end{figure}

在实际 AI 系统中，CPU 和 GPU 并不是互相替代的关系。一个训练或推理任务通常包含如下路径：

\begin{itemize}
  \item CPU 负责数据读取、tokenization、batch 组装、调度、RPC、日志与监控；
  \item GPU 负责大规模张量计算；
  \item CPU 通过 CUDA runtime 向 GPU 提交 kernel；
  \item GPU 异步执行 kernel，并通过 stream 和 event 与 CPU 协同；
  \item 多 GPU 之间通过 NVLink、PCIe、InfiniBand 或 RoCE 完成通信。
\end{itemize}

因此，AI Infra 架构师不能只懂 GPU kernel，也不能只懂上层服务。真正的系统性能往往取决于 CPU 调度、GPU 执行、内存拷贝、通信同步和存储 IO 之间是否形成良好的流水线。

\subsection{SIMD、SIMT 与 Warp}
\label{subsec:simd-simt-warp}

CPU 中常见的向量化模型是 SIMD，即 Single Instruction Multiple Data。一个 SIMD 指令可以对多个数据元素执行同一操作。例如，一个 AVX 指令可能同时对多个浮点数做加法。

GPU 使用的模型通常称为 SIMT，即 Single Instruction Multiple Threads。程序员编写的是单个线程的逻辑，但硬件会把多个线程组织成一个执行单位。在 NVIDIA GPU 中，这个基本执行单位称为 \textbf{warp}。一个 warp 通常包含 $32$ 个线程。

设一个 kernel 中有大量线程，每个线程执行如下逻辑：
\[
  c_i=a_i+b_i.
\]
逻辑上，我们可以认为每个线程独立处理一个 $i$；硬件上，GPU 会以 warp 为单位调度这些线程。若同一个 warp 中的线程走相同控制流，执行效率很高；若它们因为条件分支走不同路径，就会发生 \textbf{warp divergence}。

\begin{figure}[H]
  \centering
  \begin{tikzpicture}[
      font=\small,
      thread/.style={draw, minimum width=0.45cm, minimum height=0.42cm, fill=blue!7},
      warpbox/.style={draw, rounded corners=2pt, fill=blue!3},
      arrow/.style={-{Latex[length=2.3mm]}, thick}
    ]

    \node[font=\bfseries] at (0,2.5) {一个 Warp 中的 32 个线程};

    \foreach \i in {0,...,31} {
      \node[thread] (t\i) at ({-7.4+0.48*\i},1.5) {\tiny \i};
    }

    \node[warpbox, fit=(t0)(t31), inner sep=4pt] (warp) {};
    \node[align=center, text width=5.5cm, font=\footnotesize] at (0,0.75) {同一条指令被 32 个线程同时执行，线程 id 不同，处理的数据位置不同};

    \node[draw, rounded corners=2pt, fill=green!8, minimum width=3.2cm, minimum height=0.8cm, align=center] (same) at (-3.4,-0.5)
      {无分歧\\所有线程执行路径 A};
    \node[draw, rounded corners=2pt, fill=red!8, minimum width=3.2cm, minimum height=0.8cm, align=center] (div) at (3.4,-0.5)
      {有分歧\\部分线程路径 A，部分路径 B};

    \draw[arrow] (-1.0,1.1) -- (same.north);
    \draw[arrow] (1.0,1.1) -- (div.north);

    \node[align=left, text width=5.2cm] at (-3.4,-1.75) {
      \textbf{理想情况：}\\
      warp 内线程执行相同指令，硬件吞吐高。
    };

    \node[align=left, text width=5.2cm] at (3.4,-1.75) {
      \textbf{低效情况：}\\
      分支路径需要串行化，部分线程暂时空转。
    };
  \end{tikzpicture}
  \caption{SIMT、Warp 与控制流分歧}
  \label{fig:simt-warp-divergence}
\end{figure}

对于 AI Infra 工程师，warp divergence 最常出现在以下场景：

\begin{itemize}
  \item 处理变长序列时，不同线程面对不同有效长度；
  \item attention mask 或 padding mask 造成条件分支；
  \item 稀疏模型、MoE routing 或 block-sparse attention 中，线程处理的数据形状不规则；
  \item kernel 边界处理时，最后一个 block 的线程不满。
\end{itemize}

需要强调的是，分支本身并不一定糟糕。糟糕的是同一个 warp 内的线程频繁走不同分支，并且每个分支内工作量差异很大。很多高性能 kernel 会通过重排数据、固定 tile shape、使用 mask 计算替代显式分支等方式降低 divergence。

\subsection{为什么 GPU 适合矩阵运算}
\label{subsec:why-gpu-good-at-matmul}

矩阵乘法：
\[
  C_{ij}=\sum_{k=0}^{K-1}A_{ik}B_{kj}
\]
天然具有大量并行性。不同的 $C_{ij}$ 可以由不同线程或线程组计算；同一个 $C_{ij}$ 的求和也可以被进一步并行化。更重要的是，矩阵乘法具有很强的数据复用：

\begin{itemize}
  \item $\mat{A}$ 的一行会被用于计算 $\mat{C}$ 中多个列位置；
  \item $\mat{B}$ 的一列会被用于计算 $\mat{C}$ 中多个行位置；
  \item 分块后，一个 tile 的 $\mat{A}$ 和 $\mat{B}$ 可以在 shared memory 中被多个线程反复使用。
\end{itemize}

算术强度是衡量一个算子是否适合 GPU 的重要指标。它通常定义为：
\[
  I=\frac{\text{FLOPs}}{\text{Bytes moved}}.
\]
若 $I$ 很高，说明每读取一个字节可以完成很多计算，算子更可能是 compute-bound；若 $I$ 很低，说明计算单元经常等待内存，算子更可能是 memory-bound。

对一个理想化的 $N\times N$ 矩阵乘法，FLOPs 约为：
\[
  2N^3.
\]
如果只粗略考虑读取 $\mat{A}$、$\mat{B}$ 和写回 $\mat{C}$ 的数据量，约为：
\[
  3N^2s,
\]
其中 $s$ 是每个元素的字节数。算术强度近似为：
\[
  I\approx \frac{2N^3}{3N^2s}=\frac{2N}{3s}.
\]
随着 $N$ 增大，矩阵乘法的算术强度上升。这解释了为什么大尺寸 GEMM 往往能很好地利用 GPU，而小 batch、小矩阵、零散逐元素操作则更容易受 kernel launch 和访存开销影响。

\section{CUDA 线程层次结构}
\label{sec:cuda-thread-hierarchy}

CUDA 编程模型把并行执行组织成三层：

\begin{itemize}
  \item \textbf{Grid}：一次 kernel launch 产生一个 grid；
  \item \textbf{Block}：grid 由多个 block 组成；
  \item \textbf{Thread}：每个 block 由多个 thread 组成。
\end{itemize}

线程通过内置变量识别自己的位置：

\begin{itemize}
  \item \texttt{threadIdx.x/y/z}：线程在 block 内的索引；
  \item \texttt{blockIdx.x/y/z}：block 在 grid 内的索引；
  \item \texttt{blockDim.x/y/z}：每个 block 的维度；
  \item \texttt{gridDim.x/y/z}：grid 的维度。
\end{itemize}

\subsection{Grid、Block、Thread}
\label{subsec:grid-block-thread}

一个典型一维索引写法是：
\[
  i=\texttt{blockIdx.x}\cdot\texttt{blockDim.x}+\texttt{threadIdx.x}.
\]
这条公式几乎是 CUDA 入门的第一行公式，但它在 AI kernel 中反复出现。无论是向量加法、embedding lookup、LayerNorm 还是简单的 elementwise activation，本质上都需要把逻辑张量元素映射到线程索引。

对于二维矩阵，我们常使用二维 block：
\[
  row=\texttt{blockIdx.y}\cdot\texttt{blockDim.y}+\texttt{threadIdx.y},
\]
\[
  col=\texttt{blockIdx.x}\cdot\texttt{blockDim.x}+\texttt{threadIdx.x}.
\]

\begin{figure}[H]
  \centering
  \begin{tikzpicture}[
      font=\small,
      block/.style={draw, rounded corners=2pt, minimum width=1.35cm, minimum height=0.85cm, fill=blue!6, align=center},
      thread/.style={draw, minimum width=0.35cm, minimum height=0.32cm, fill=orange!12},
      arrow/.style={-{Latex[length=2.3mm]}, thick}
    ]

    \node[font=\bfseries] at (0,3.15) {CUDA 线程层次结构：Grid $\rightarrow$ Block $\rightarrow$ Thread};

    \foreach \r in {0,...,1} {
      \foreach \c in {0,...,3} {
        \node[block] (b\r\c) at ({-3.6+2.4*\c},{1.85-1.2*\r})
          {Block\\$(\c,\r)$};
      }
    }

    \node[draw, rounded corners=4pt, fit=(b00)(b13), inner sep=8pt, label={[font=\small]above left:Grid}] (gridbox) {};

    \node[block, fill=green!8, minimum width=2.1cm, minimum height=1.3cm] (oneblock) at (0,-1.4) {一个 Block};
    \draw[arrow] (b12.south) -- (oneblock.north);

    \foreach \r in {0,...,2} {
      \foreach \c in {0,...,5} {
        \node[thread] at ({-1.1+0.42*\c},{-0.95-0.35*\r}) {};
      }
    }

    \node[align=left, text width=5.2cm] at (4.0,-1.3) {
      \textbf{映射关系：}\\
      Grid 是一次 kernel launch 的全局任务；\\
      Block 是调度到 SM 的基本工作单元；\\
      Thread 是执行标量逻辑的最小编程单位。
    };
  \end{tikzpicture}
  \caption{Grid、Block、Thread 三层结构}
  \label{fig:cuda-grid-block-thread}
\end{figure}

Block 是 CUDA 编程中的关键边界。一个 block 内的线程可以：

\begin{itemize}
  \item 通过 shared memory 共享数据；
  \item 使用 \texttt{\_\_syncthreads()} 做同步；
  \item 协作完成一个 tile 的计算。
\end{itemize}

不同 block 之间通常不能在同一个 kernel 内直接同步。若需要全局同步，常见方式是结束当前 kernel，再 launch 下一个 kernel。这个约束影响了很多算法设计。例如，一个大规模 reduction 往往需要多个阶段：每个 block 先做局部 reduction，再由另一个 kernel 汇总局部结果。

\subsection{Warp 与线程束调度}
\label{subsec:warp-scheduling}

虽然程序员创建的是 thread，但硬件调度的基本单位是 warp。一个 block 会被划分成若干 warp。例如，一个 block 有 $256$ 个线程，则包含：
\[
  \frac{256}{32}=8
\]
个 warp。

GPU 的 SM 会维护多个活跃 warp。当一个 warp 因为 global memory load 等待数据时，调度器可以切换到另一个 ready warp 执行。这样，GPU 用大量线程隐藏内存延迟。

这与 CPU 的思路不同。CPU 通常依靠大缓存、分支预测和乱序执行降低单线程延迟；GPU 则通过同时驻留大量 warp，把某个 warp 的等待时间用另一个 warp 的计算填满。

\begin{figure}[H]
  \centering
  \begin{tikzpicture}[
      font=\small,
      sm/.style={draw, rounded corners=3pt, fill=blue!5, minimum width=9.5cm, minimum height=3.5cm},
      warp/.style={draw, rounded corners=2pt, fill=orange!12, minimum width=1.4cm, minimum height=0.55cm, align=center},
      sched/.style={draw, rounded corners=2pt, fill=green!10, minimum width=2.2cm, minimum height=0.75cm, align=center},
      arrow/.style={-{Latex[length=2.3mm]}, thick}
    ]

    \node[sm] (sm0) at (0,0) {};
    \node[font=\bfseries] at (-4.1,1.45) {Streaming Multiprocessor};

    \node[sched] (sched0) at (-2.9,0.75) {Warp Scheduler};
    \node[sched] (exec0) at (2.8,0.75) {Execution Units};

    \foreach \i/\y in {0/0.2,1/-0.45,2/-1.1} {
      \node[warp] (w\i) at (-2.9,\y) {Warp \i};
      \draw[arrow] (w\i.east) -- (exec0.west);
    }

    \node[warp, fill=red!10] (wait) at (-2.9,-1.75) {Warp 3\\waiting};
    \draw[dashed, -{Latex[length=2.3mm]}, thick] (wait.east) -- (exec0.west);

    \node[align=left, text width=4.2cm] at (0,-2.6) {
       当某个 warp 等待内存时，\\
       调度器选择其他 ready warp 执行，\\
       以隐藏访存延迟。
    };
  \end{tikzpicture}
  \caption{SM 内部的 warp 调度直觉}
  \label{fig:warp-scheduling}
\end{figure}

在写 CUDA kernel 时，我们常看到两类优化目标：

\begin{itemize}
  \item \textbf{减少单个线程的无效工作}：避免冗余计算、减少分支、减少寄存器压力。
  \item \textbf{提高 SM 上可驻留的活跃 warp 数量}：让硬件有足够 warp 用于隐藏延迟。
\end{itemize}

后者就是 occupancy 概念的入口。

\subsection{Occupancy 的基本概念}
\label{subsec:occupancy}

Occupancy 通常指一个 SM 上实际活跃 warp 数量与硬件允许最大活跃 warp 数量的比例：
\[
  \mathrm{Occupancy}
  =
  \frac{\text{Active Warps per SM}}{\text{Maximum Warps per SM}}.
\]
更高 occupancy 可以帮助隐藏延迟，但 \textbf{occupancy 并不是越高越好}。原因是：

\begin{itemize}
  \item 如果 kernel 是纯计算密集型，较低 occupancy 也可能充分利用 Tensor Core；
  \item 如果每个线程使用太多 register，会降低可驻留 warp 数量；
  \item 如果每个 block 使用太多 shared memory，会减少同一 SM 上可同时驻留的 block；
  \item 如果为了提高 occupancy 而降低 tile 复用，可能反而增加 HBM 访问。
\end{itemize}

Occupancy 受多种资源约束影响。可以粗略写成：
\[
  B_{\mathrm{active}}
  =
  \min
  \left(
    B_{\mathrm{thread}},
    B_{\mathrm{reg}},
    B_{\mathrm{smem}},
    B_{\mathrm{arch}}
  \right),
\]
其中：

\begin{itemize}
  \item $B_{\mathrm{thread}}$ 由每个 block 的线程数限制；
  \item $B_{\mathrm{reg}}$ 由每个线程使用的寄存器数量限制；
  \item $B_{\mathrm{smem}}$ 由每个 block 使用的 shared memory 限制；
  \item $B_{\mathrm{arch}}$ 是硬件架构规定的最大 block 数限制。
\end{itemize}

对 AI kernel 来说，occupancy 只是性能分析的一个维度。更完整的判断需要结合：

\begin{itemize}
  \item achieved occupancy；
  \item SM utilization；
  \item memory throughput；
  \item L2 hit rate；
  \item tensor core utilization；
  \item eligible warps per cycle；
  \item stall reason。
\end{itemize}

这也是为什么后续 GPU 架构和 profiling 章节会重新讨论 roofline model 与 Nsight Compute。

\subsection{\texttt{blockDim}、\texttt{gridDim}、\texttt{threadIdx}、\texttt{blockIdx}}
\label{subsec:cuda-built-in-indices}

下面这段代码展示了最典型的一维 CUDA 索引模式。它几乎是所有 CUDA kernel 的起点。

\begin{lstlisting}[language=C++,caption={CUDA 一维线程索引模式},label={lst:cuda-1d-indexing}]
__global__ void kernel_1d(float* out, const float* in, int n) {
  // Global linear index of the current thread.
  int idx = blockIdx.x * blockDim.x + threadIdx.x;

  // Boundary check is required because n may not be a multiple of blockDim.x.
  if (idx < n) {
    out[idx] = in[idx] * 2.0f;
  }
}
\end{lstlisting}

二维索引常用于矩阵、图像或二维 tile：

\begin{lstlisting}[language=C++,caption={CUDA 二维线程索引模式},label={lst:cuda-2d-indexing}]
__global__ void kernel_2d(float* out, const float* in, int rows, int cols) {
  int col = blockIdx.x * blockDim.x + threadIdx.x;
  int row = blockIdx.y * blockDim.y + threadIdx.y;

  if (row < rows && col < cols) {
    int offset = row * cols + col;
    out[offset] = in[offset];
  }
}
\end{lstlisting}

这两个例子看似简单，但它们已经包含了 CUDA kernel 的三个基本要点：

\begin{itemize}
  \item \textbf{线程到数据的映射}：每个线程通过索引找到自己负责的元素。
  \item \textbf{边界检查}：grid 通常会向上取整，最后一部分线程可能越界。
  \item \textbf{内存访问模式}：连续线程最好访问连续地址，以形成合并访存。
\end{itemize}

\section{CUDA 内存模型}
\label{sec:cuda-memory-model}

如果说线程层次结构决定了计算如何并行展开，那么内存模型就决定了这些计算能否高效喂饱 GPU。AI Infra 中很多性能问题不是因为 FLOPs 不够，而是因为数据搬运太慢。

CUDA 中常见的内存层级包括：

\begin{itemize}
  \item \textbf{Global Memory}：容量大，位于显存 HBM，延迟高，带宽高；
  \item \textbf{Shared Memory}：位于 SM 内部，容量小，延迟低，可由同一 block 内线程共享；
  \item \textbf{Register}：每个线程私有，速度最快，但数量有限；
  \item \textbf{Constant Memory}：适合广播读取只读常量；
  \item \textbf{Texture Memory}：适合某些具有空间局部性的只读访问模式；
  \item \textbf{L1/L2 Cache}：硬件管理，用于缓解 global memory 访问成本。
\end{itemize}

\subsection{Global Memory}
\label{subsec:global-memory}

Global memory 是 GPU 上容量最大的显式可访问内存，通常对应 HBM。模型参数、激活、梯度、优化器状态、KV Cache 等绝大多数大张量都存储在 global memory 中。

Global memory 的特点是：

\begin{itemize}
  \item 容量大，能容纳大模型张量；
  \item 延迟高，单次访问成本远高于寄存器和 shared memory；
  \item 带宽高，但需要良好的访问模式才能充分利用；
  \item 所有 block 都可访问，但 block 之间不能依赖 global memory 实现安全同步，除非跨 kernel 或使用特定同步机制。
\end{itemize}

一个低效 kernel 的典型模式是：每次计算都直接从 global memory 读取数据，并且同一份数据被不同线程重复读取。高效 kernel 则尽量把可复用数据先加载到 shared memory 或 register，再在更快的层级中重复使用。

\subsection{Shared Memory}
\label{subsec:shared-memory}

Shared memory 是 CUDA 性能优化中最重要的显式内存层级之一。它位于 SM 内部，被同一个 block 内的所有线程共享。访问 shared memory 通常远快于 global memory，因此适合保存 tile 数据、中间结果或 block 内 reduction 的临时数据。

例如，在 tiled GEMM 中，一个 block 会负责计算 $\mat{C}$ 的一个 tile。线程先把 $\mat{A}$ 和 $\mat{B}$ 的对应 tile 从 global memory 读入 shared memory，然后反复使用：

\[
  \mat{C}_{\mathrm{tile}}
  +=
  \mat{A}_{\mathrm{tile}}
  \mat{B}_{\mathrm{tile}}.
\]

这能显著减少 global memory 访问次数。若没有 shared memory，同一个 $\mat{A}$ 元素可能被多个线程反复从 HBM 读取；有了 shared memory，这个元素只需被加载一次，然后由多个线程复用。

Shared memory 也有两个常见陷阱：

\begin{itemize}
  \item \textbf{容量有限}：每个 block 使用太多 shared memory，会降低 occupancy。
  \item \textbf{bank conflict}：多个线程同时访问落在同一 bank 的地址，可能导致访问串行化。
\end{itemize}

本章先关注 shared memory 的基本复用思想，bank conflict 会在后续 GPU 架构与性能优化章节中进一步展开。

\subsection{Register}
\label{subsec:register}

Register 是每个线程私有的最快存储资源。局部变量、循环累加器、临时标量通常会被编译器放入 register。例如矩阵乘法中，每个线程计算一个 $C_{ij}$ 时，累加变量：

\[
  acc=\sum_k A_{ik}B_{kj}
\]
通常应保存在 register 中。

Register 使用过多会带来两个问题：

\begin{itemize}
  \item 降低一个 SM 上可同时驻留的线程或 warp 数量；
  \item 触发 register spilling，把本应在 register 中的变量溢出到 local memory，而 local memory 实际上位于 global memory，代价很高。
\end{itemize}

因此，高性能 CUDA kernel 常常需要在单线程工作量、register 使用量、occupancy 和 memory reuse 之间做平衡。对于 Transformer 中的 attention、LayerNorm、RMSNorm 等 kernel，这种平衡尤其关键。

\subsection{Constant Memory 与 Texture Memory}
\label{subsec:constant-texture-memory}

Constant memory 适合保存所有线程读取相同地址的只读数据。例如某些固定超参数、查表常量或小型配置数组。如果一个 warp 中所有线程读取同一个 constant memory 地址，硬件可以高效广播。

Texture memory 原本服务于图形渲染，但在通用计算中也可用于某些只读、具有空间局部性的访问模式。现代深度学习框架中，开发者直接使用 texture memory 的机会不如 shared memory 和 register 多，但理解它有助于完整认识 GPU 的存储层级。

对于 AI Infra 工程师，更常见的是通过高层库间接受益。例如 cuDNN、TensorRT 或某些 fused kernel 可能根据访问模式选择合适的缓存路径。大多数情况下，我们需要优先掌握 global memory、shared memory、register、L2 cache 与 HBM 带宽之间的关系。

\subsection{Memory Coalescing}
\label{subsec:memory-coalescing}

Memory coalescing 是 CUDA 性能优化的核心概念之一。若同一个 warp 中的连续线程访问连续内存地址，硬件可以把这些访问合并成更少的内存事务。反之，如果线程访问地址分散，内存事务数量会增加，带宽利用率下降。

考虑一个 FP32 数组 \texttt{x}，线程 $t$ 访问：
\[
  x[t].
\]
这是良好的连续访问。若线程 $t$ 访问：
\[
  x[t\cdot stride],
\]
当 $stride$ 很大时，同一个 warp 内线程访问的地址相距很远，coalescing 变差。

\begin{figure}[H]
  \centering
  \begin{tikzpicture}[
      font=\small,
      cell/.style={draw, minimum width=0.38cm, minimum height=0.38cm},
      thread/.style={draw, rounded corners=1pt, minimum width=0.38cm, minimum height=0.35cm, fill=blue!8},
      arrow/.style={-{Latex[length=2.0mm]}, thick},
      good/.style={draw, rounded corners=2pt, fill=green!7, align=center},
      bad/.style={draw, rounded corners=2pt, fill=red!7, align=center}
    ]

    \node[font=\bfseries] at (0,3.0) {Coalesced vs. Uncoalesced Memory Access};

    \node[good, minimum width=5.6cm, minimum height=0.55cm] at (-3.4,2.25) {连续线程访问连续地址};
    \foreach \i in {0,...,7} {
      \node[thread] (gt\i) at ({-5.0+0.45*\i},1.55) {\tiny T\i};
      \node[cell, fill=green!8] (gm\i) at ({-5.0+0.45*\i},0.65) {};
      \draw[arrow] (gt\i) -- (gm\i);
    }
    \node at (-3.4,0.05) {一次或少量内存事务};

    \node[bad, minimum width=5.6cm, minimum height=0.55cm] at (3.4,2.25) {连续线程访问分散地址};
    \foreach \i/\x in {0/1.2,1/1.8,2/2.7,3/3.5,4/4.4,5/5.2,6/5.9,7/6.5} {
      \node[thread] (bt\i) at ({1.7+0.45*\i},1.55) {\tiny T\i};
      \node[cell, fill=red!8] (bm\i) at (\x,0.65) {};
      \draw[arrow] (bt\i) -- (bm\i);
    }
    \node at (3.4,0.05) {多次内存事务，带宽浪费};

    \node[align=left, text width=10.8cm] at (0,-1.15) {
      在张量计算中，coalescing 与 memory layout 密切相关。若张量最后一维连续，让相邻线程沿最后一维访问，通常更容易获得高带宽。
      这也是上一章讨论 shape、stride 和 contiguous tensor 的工程原因。
    };
  \end{tikzpicture}
  \caption{合并访存与非合并访存示意图}
  \label{fig:memory-coalescing}
\end{figure}

在 LLM 推理中，KV Cache 的读取经常成为 decode 阶段瓶颈。若 KV Cache 布局不能让连续线程访问连续地址，或者请求长度高度不规则，attention kernel 就很难达到理想带宽。后续 PagedAttention 和 FlashAttention 章节会继续围绕这个问题展开。

\subsection{GPU 内存层级图}
\label{subsec:gpu-memory-hierarchy-placeholder}

真实 GPU 的硬件结构非常复杂，不同架构在 SM、L1、L2、Tensor Core、HBM 和互联方面均有差异。本书在后续 GPU 架构章节会专门展开。这里先给出一个抽象内存层级图，帮助建立基本直觉。

\begin{figure}[H]
  \centering
  \begin{tikzpicture}[
      font=\small,
      level/.style={draw, rounded corners=3pt, minimum width=8.5cm, minimum height=0.8cm, align=center},
      arrow/.style={-{Latex[length=2.3mm]}, thick}
    ]

    \node[level, fill=red!10] (reg) at (0,2.4) {Register：线程私有，速度最快，容量最小};
    \node[level, fill=orange!12] (smem) at (0,1.35) {Shared Memory / L1 Cache：Block 内共享，低延迟，容量有限};
    \node[level, fill=yellow!18] (l2) at (0,0.3) {L2 Cache：跨 SM 共享缓存，硬件管理};
    \node[level, fill=green!10] (hbm) at (0,-0.75) {Global Memory / HBM：容量大，带宽高，但访问延迟高};
    \node[level, fill=blue!8] (host) at (0,-1.8) {Host Memory：CPU 内存，经 PCIe / NVLink-C2C 等路径访问};

    \draw[arrow] (reg) -- node[right] {容量增加，延迟增加} (smem);
    \draw[arrow] (smem) -- (l2);
    \draw[arrow] (l2) -- (hbm);
    \draw[arrow] (hbm) -- (host);

    \draw[decorate,decoration={brace,amplitude=5pt},thick]
      (-4.8,2.85) -- (-4.8,1.00)
      node[midway,left=8pt,align=center] {SM 内部\\显式优化重点};

    \draw[decorate,decoration={brace,amplitude=5pt},thick]
      (-4.8,0.55) -- (-4.8,-1.25)
      node[midway,left=8pt,align=center] {芯片/显存\\带宽瓶颈重点};
  \end{tikzpicture}
  \caption{GPU 内存层级的抽象视图}
  \label{fig:gpu-memory-hierarchy-abstract}
\end{figure}

下面保留一个真实硬件图的位置，便于后续排版时替换为官方或自绘硬件结构图。即使图片文件不存在，本段也能正常编译。

\begin{figure}[H]
  \centering
  \IfFileExists{figures/part1/gpu_memory_hierarchy_placeholder.png}{
    \includegraphics[width=0.85\textwidth]{part1/gpu_memory_hierarchy_placeholder.png}
  }{
    \fbox{
      \parbox{0.82\textwidth}{
        \centering
        \vspace{1.2cm}
        \textbf{预留图片位：GPU 内存层级与 SM 结构}\\[2mm]
        此处建议放入一张硬件风格示意图，展示 GPU 芯片中的多个 SM、每个 SM 内部的 register file、shared memory/L1、warp scheduler、Tensor Core，以及芯片级 L2 cache 与 HBM。\\[2mm]
        \textbf{图片生成 Prompt 建议：}\\
        ``Technical diagram of NVIDIA-style GPU memory hierarchy for AI infrastructure textbook, showing multiple streaming multiprocessors, register file, shared memory, L1 cache, L2 cache, HBM stacks, PCIe or NVLink connection, clean vector style, blue and green color palette, labeled components, high readability.''
        \vspace{1.2cm}
      }
    }
  }
  \caption{GPU 内存层级结构示意图预留位}
  \label{fig:gpu-memory-hierarchy-placeholder}
\end{figure}

\section{第一个 CUDA Kernel}
\label{sec:first-cuda-kernel}

本节从向量加法开始，逐步过渡到矩阵乘法。向量加法展示线程索引和 kernel launch；naive GEMM 展示二维线程映射；tiled GEMM 展示 shared memory 复用。这三者构成了理解深度学习 kernel 的最小路径。

\subsection{向量加法}
\label{subsec:cuda-vector-add}

向量加法的数学定义非常简单：
\[
  \vect{c}=\vect{a}+\vect{b},
\]
即：
\[
  c_i=a_i+b_i.
\]
因为每个元素相互独立，所以可以让一个线程处理一个元素。

\begin{lstlisting}[language=C++,caption={CUDA 向量加法 Kernel},label={lst:cuda-vector-add-kernel}]
#include <cuda_runtime.h>
#include <iostream>

__global__ void vector_add_kernel(const float* a,
                                  const float* b,
                                  float* c,
                                  int n) {
  // Each thread computes one global element index.
  int idx = blockIdx.x * blockDim.x + threadIdx.x;

  // The grid size is usually rounded up, so some threads may be out of range.
  if (idx < n) {
    c[idx] = a[idx] + b[idx];
  }
}

int main() {
  int n = 1 << 20;
  size_t bytes = n * sizeof(float);

  float* h_a = new float[n];
  float* h_b = new float[n];
  float* h_c = new float[n];

  for (int i = 0; i < n; ++i) {
    h_a[i] = 1.0f;
    h_b[i] = 2.0f;
  }

  float *d_a, *d_b, *d_c;
  cudaMalloc(&d_a, bytes);
  cudaMalloc(&d_b, bytes);
  cudaMalloc(&d_c, bytes);

  cudaMemcpy(d_a, h_a, bytes, cudaMemcpyHostToDevice);
  cudaMemcpy(d_b, h_b, bytes, cudaMemcpyHostToDevice);

  int threads_per_block = 256;
  int blocks = (n + threads_per_block - 1) / threads_per_block;

  vector_add_kernel<<<blocks, threads_per_block>>>(d_a, d_b, d_c, n);

  cudaMemcpy(h_c, d_c, bytes, cudaMemcpyDeviceToHost);

  std::cout << "h_c[0] = " << h_c[0] << std::endl;

  cudaFree(d_a);
  cudaFree(d_b);
  cudaFree(d_c);
  delete[] h_a;
  delete[] h_b;
  delete[] h_c;

  return 0;
}
\end{lstlisting}

这段代码展示了 CUDA 程序的基本结构：

\begin{itemize}
  \item 在 CPU 侧分配 host memory；
  \item 使用 \texttt{cudaMalloc} 在 GPU 侧分配 device memory；
  \item 使用 \texttt{cudaMemcpy} 把数据从 CPU 拷贝到 GPU；
  \item 使用 \texttt{kernel<<<grid, block>>>} 发起 kernel launch；
  \item 把结果从 GPU 拷贝回 CPU；
  \item 释放 GPU 和 CPU 资源。
\end{itemize}

注意，kernel launch 默认是异步的。CPU 发起 kernel 后不一定等待 GPU 执行完成。上面的 \texttt{cudaMemcpyDeviceToHost} 会隐式等待相关 kernel 完成，因为拷贝结果依赖 kernel 输出。在性能敏感场景中，我们会显式使用 stream、event、异步拷贝和同步控制。

\subsection{矩阵乘法 naive 实现}
\label{subsec:naive-matmul}

接下来考虑矩阵乘法：
\[
  \mat{C}=\mat{A}\mat{B},
\]
其中：
\[
  \mat{A}\in\mathbb{R}^{M\times K},
  \quad
  \mat{B}\in\mathbb{R}^{K\times N},
  \quad
  \mat{C}\in\mathbb{R}^{M\times N}.
\]
每个输出元素为：
\[
  C_{row,col}
  =
  \sum_{k=0}^{K-1}A_{row,k}B_{k,col}.
\]

最直接的 CUDA 实现是：每个线程计算一个 $C_{row,col}$。

\begin{lstlisting}[language=C++,caption={Naive CUDA 矩阵乘法 Kernel},label={lst:cuda-naive-matmul}]
__global__ void matmul_naive_kernel(const float* A,
                                    const float* B,
                                    float* C,
                                    int M,
                                    int N,
                                    int K) {
  // 2D thread mapping:
  // x dimension maps to columns of C
  // y dimension maps to rows of C
  int col = blockIdx.x * blockDim.x + threadIdx.x;
  int row = blockIdx.y * blockDim.y + threadIdx.y;

  if (row < M && col < N) {
    float acc = 0.0f;

    for (int k = 0; k < K; ++k) {
      float a = A[row * K + k];
      float b = B[k * N + col];
      acc += a * b;
    }

    C[row * N + col] = acc;
  }
}
\end{lstlisting}

这个 kernel 的优点是简单，缺点也非常明显：它没有充分复用 global memory 数据。对于相邻的多个线程，它们可能读取相同的 $\mat{A}$ 行或相同的 $\mat{B}$ 列，但每个线程都独立从 global memory 加载，导致大量冗余访存。

可以粗略分析其访存模式。每个 $C_{ij}$ 需要读取 $K$ 个 $A$ 元素和 $K$ 个 $B$ 元素，执行 $K$ 次乘加。若每个线程独立完成，global memory 读取量约为：
\[
  2MNK
\]
个浮点数，而 $\mat{A}$ 和 $\mat{B}$ 中的许多元素本可以被多个输出复用。这就是 tiled matrix multiplication 的动机。

\subsection{使用 shared memory 优化 GEMM}
\label{subsec:shared-memory-gemm}

Tiled GEMM 的核心思想是把 $\mat{A}$、$\mat{B}$、$\mat{C}$ 切成小块。每个 block 负责计算 $\mat{C}$ 的一个 tile。计算过程中，线程协作把 $\mat{A}$ 和 $\mat{B}$ 的对应 tile 读入 shared memory，然后在 shared memory 中复用。

设 tile size 为 $T$。对某个 $\mat{C}$ tile：
\[
  \mat{C}_{tile}
  =
  \sum_{p}
  \mat{A}_{tile,p}
  \mat{B}_{p,tile}.
\]
每轮加载一块 $\mat{A}_{tile,p}$ 和一块 $\mat{B}_{p,tile}$，同步后进行局部乘加。

\begin{figure}[H]
  \centering
  \begin{tikzpicture}[
      font=\small,
      mat/.style={draw, rounded corners=2pt, fill=blue!4},
      tileA/.style={draw, fill=orange!22},
      tileB/.style={draw, fill=green!18},
      tileC/.style={draw, fill=red!15},
      arrow/.style={-{Latex[length=2.3mm]}, thick}
    ]

    \node[font=\bfseries] at (0,3.2) {Tiled Matrix Multiplication};

    \draw[mat] (-5.8,0.2) rectangle (-2.6,2.6);
    \node at (-4.2,2.85) {$\mat{A}\in\mathbb{R}^{M\times K}$};
    \draw[tileA] (-5.8,1.0) rectangle (-4.4,1.8);
    \node at (-5.1,1.4) {$A_t$};

    \draw[mat] (-1.6,0.2) rectangle (1.8,2.6);
    \node at (0.1,2.85) {$\mat{B}\in\mathbb{R}^{K\times N}$};
    \draw[tileB] (-0.4,1.0) rectangle (0.6,1.8);
    \node at (0.1,1.4) {$B_t$};

    \draw[mat] (3.0,0.2) rectangle (5.8,2.6);
    \node at (4.4,2.85) {$\mat{C}\in\mathbb{R}^{M\times N}$};
    \draw[tileC] (3.7,1.0) rectangle (4.7,1.8);
    \node at (4.2,1.4) {$C_t$};

    \node[draw, rounded corners=2pt, fill=yellow!12, align=center, minimum width=2.5cm, minimum height=0.9cm] (shared) at (0,-1.1)
      {Shared Memory\\缓存 $A_t,B_t$};

    \draw[arrow] (-5.1,1.0) -- (-0.8,-0.8);
    \draw[arrow] (0.1,1.0) -- (0.1,-0.65);
    \draw[arrow] (shared.east) -- node[above] {复用后累加} (3.7,1.0);

    \node[align=left, text width=10.5cm] at (0,-2.55) {
      每个 tile 从 global memory 加载一次，随后被 block 内多个线程重复使用。
      这种数据复用显著提高算术强度，是高性能 GEMM 的基本思想。
    };
  \end{tikzpicture}
  \caption{Tiled GEMM 中的 shared memory 数据复用}
  \label{fig:tiled-gemm-shared-memory}
\end{figure}

下面给出一个教学版 shared memory GEMM。它没有使用 Tensor Core，也没有做复杂的向量化和双缓冲，但足以展示 tiled 结构。

\begin{lstlisting}[language=C++,caption={使用 shared memory 的教学版 tiled GEMM},label={lst:cuda-tiled-gemm}]
#define TILE 16

__global__ void matmul_tiled_kernel(const float* A,
                                    const float* B,
                                    float* C,
                                    int M,
                                    int N,
                                    int K) {
  __shared__ float As[TILE][TILE];
  __shared__ float Bs[TILE][TILE];

  int tx = threadIdx.x;
  int ty = threadIdx.y;

  int row = blockIdx.y * TILE + ty;
  int col = blockIdx.x * TILE + tx;

  float acc = 0.0f;

  // Number of tiles along K dimension.
  int num_tiles = (K + TILE - 1) / TILE;

  for (int t = 0; t < num_tiles; ++t) {
    int a_col = t * TILE + tx;
    int b_row = t * TILE + ty;

    // Load A tile into shared memory.
    if (row < M && a_col < K) {
      As[ty][tx] = A[row * K + a_col];
    } else {
      As[ty][tx] = 0.0f;
    }

    // Load B tile into shared memory.
    if (b_row < K && col < N) {
      Bs[ty][tx] = B[b_row * N + col];
    } else {
      Bs[ty][tx] = 0.0f;
    }

    // Make sure all threads finish loading before computation.
    __syncthreads();

    // Compute partial dot product using data in shared memory.
    for (int k = 0; k < TILE; ++k) {
      acc += As[ty][k] * Bs[k][tx];
    }

    // Make sure all threads finish using the current tile before overwrite.
    __syncthreads();
  }

  if (row < M && col < N) {
    C[row * N + col] = acc;
  }
}
\end{lstlisting}

这段代码中有两个 \texttt{\_\_syncthreads()}。第一个同步保证所有线程已经把 tile 加载到 shared memory；第二个同步保证所有线程完成当前 tile 的使用，然后才能进入下一轮覆盖 shared memory。

需要特别注意：

\begin{itemize}
  \item \texttt{\_\_syncthreads()} 是 block 内同步，不是 grid 全局同步；
  \item 所有线程必须以一致方式到达同步点，否则可能死锁；
  \item shared memory tile 的大小会影响 occupancy；
  \item 实际高性能 GEMM 会使用更复杂的 tiling、register blocking、warp-level MMA、Tensor Core、pipeline 和 double buffering。
\end{itemize}

现代深度学习框架中的 GEMM 通常不由工程师手写，而是调用 cuBLAS、cuBLASLt 或 CUTLASS 风格的高性能实现。但理解 tiled GEMM 对后续理解 FlashAttention、fused MLP、blockwise quantization 和 Tensor Core mapping 都非常关键。

\subsection{CUDA kernel 性能分析初步}
\label{subsec:cuda-kernel-performance-basics}

写出正确的 CUDA kernel 只是第一步。对 AI Infra 来说，更重要的是判断它为什么快或为什么慢。最基本的性能分析可以从以下四个指标开始：

\begin{itemize}
  \item \textbf{吞吐}：单位时间处理多少元素、token 或 FLOPs；
  \item \textbf{访存带宽}：实际达到的 HBM bandwidth；
  \item \textbf{算术强度}：每搬运一个字节执行多少计算；
  \item \textbf{kernel launch overhead}：大量小 kernel 是否被 launch 开销主导。
\end{itemize}

对于一个向量加法：
\[
  c_i=a_i+b_i,
\]
每个元素读取 $a_i,b_i$，写回 $c_i$，FP32 情况下数据移动约：
\[
  3\times 4=12\ \mathrm{bytes}.
\]
计算只有一次加法，即 $1$ FLOP。因此算术强度约为：
\[
  I_{\mathrm{vecadd}}=\frac{1}{12}\ \mathrm{FLOPs/byte}.
\]
这显然是 memory-bound 操作。

而 GEMM 的算术强度随矩阵尺寸增加，通常更容易成为 compute-bound。因此，在 GPU 上优化深度学习模型时，我们经常追求两类目标：

\begin{itemize}
  \item 把小而碎的 elementwise 操作融合，减少 HBM 往返和 kernel launch；
  \item 把矩阵乘法组织成适合 Tensor Core 的大块计算，提高 compute utilization。
\end{itemize}

一个简单的 kernel 计时代码如下：

\begin{lstlisting}[language=C++,caption={使用 CUDA Event 统计 kernel 执行时间},label={lst:cuda-event-timing}]
cudaEvent_t start, stop;
cudaEventCreate(&start);
cudaEventCreate(&stop);

cudaEventRecord(start);

vector_add_kernel<<<blocks, threads_per_block>>>(d_a, d_b, d_c, n);

cudaEventRecord(stop);
cudaEventSynchronize(stop);

float milliseconds = 0.0f;
cudaEventElapsedTime(&milliseconds, start, stop);

std::cout << "Kernel time: " << milliseconds << " ms" << std::endl;

cudaEventDestroy(start);
cudaEventDestroy(stop);
\end{lstlisting}

CUDA event 计时适合初步测量 GPU 端 kernel 时间。但生产级性能分析需要使用 Nsight Systems 和 Nsight Compute：

\begin{itemize}
  \item \textbf{Nsight Systems} 更适合观察整体时间线，包括 CPU 调度、CUDA API、kernel launch、memory copy、NCCL 通信等；
  \item \textbf{Nsight Compute} 更适合深入单个 kernel，分析 occupancy、warp stall、memory throughput、L2 hit rate、shared memory conflict 等。
\end{itemize}

\section{CUDA 与深度学习框架}
\label{sec:cuda-and-dl-frameworks}

实际 AI Infra 工作中，我们很少直接把完整模型写成 CUDA C++。更常见的路径是：

\begin{itemize}
  \item 使用 PyTorch、JAX、TensorFlow 等框架描述模型；
  \item 框架把张量操作映射为 operator；
  \item operator 调用 cuBLAS、cuDNN、NCCL、Triton kernel 或自定义 CUDA kernel；
  \item GPU driver 和 runtime 负责调度执行；
  \item profiler 帮助工程师定位瓶颈。
\end{itemize}

\subsection{PyTorch Tensor 与 CUDA Tensor}
\label{subsec:pytorch-tensor-cuda-tensor}

在 PyTorch 中，一个张量既包含数据，也包含元数据。元数据包括：

\begin{itemize}
  \item shape；
  \item stride；
  \item dtype；
  \item device；
  \item storage pointer；
  \item requires grad；
  \item layout。
\end{itemize}

当我们执行：

\begin{lstlisting}[language=Python,caption={PyTorch CUDA Tensor 示例},label={lst:pytorch-cuda-tensor}]
import torch

x = torch.randn(1024, 4096, device="cuda", dtype=torch.float16)
w = torch.randn(4096, 4096, device="cuda", dtype=torch.float16)

y = x @ w

print(y.shape)
print(y.device)
print(y.dtype)
\end{lstlisting}

表面上只有一行 \texttt{x @ w}，但背后发生了多层映射：

\begin{itemize}
  \item PyTorch dispatcher 根据 dtype、device 和 operator 类型选择实现；
  \item 对 CUDA tensor，矩阵乘法通常会调用 cuBLAS 或 cuBLASLt；
  \item cuBLAS 根据矩阵 shape、transpose、dtype 等选择具体 kernel；
  \item kernel 在 GPU 上执行，可能使用 Tensor Core；
  \item 输出 tensor 的 metadata 被返回给 PyTorch。
\end{itemize}

这说明 AI Infra 工程师面对性能问题时，不能只看 Python 层代码。一个简单 operator 背后可能有复杂的 backend 选择、workspace 分配、算法启发式选择和 kernel 实现。

\subsection{cuBLAS、cuDNN、NCCL 的角色}
\label{subsec:cublas-cudnn-nccl}

NVIDIA 深度学习软件栈中有几个核心库：

\begin{itemize}
  \item \textbf{cuBLAS}：提供高性能 BLAS，尤其是 GEMM，是 Transformer 中 Linear、QKV projection、FFN 的核心。
  \item \textbf{cuBLASLt}：提供更灵活的 GEMM 接口，支持 layout、epilogue fusion、算法选择等能力。
  \item \textbf{cuDNN}：提供卷积、归一化、RNN、attention 等深度学习算子。
  \item \textbf{NCCL}：提供多 GPU 通信 collective，例如 all-reduce、all-gather、reduce-scatter、broadcast。
  \item \textbf{CUDA Runtime / Driver}：负责 kernel launch、memory management、stream、event、device 控制等基础能力。
\end{itemize}

\begin{figure}[H]
  \centering
  \begin{tikzpicture}[
      font=\small,
      layer/.style={draw, rounded corners=3pt, minimum width=8.5cm, minimum height=0.72cm, align=center},
      arrow/.style={-{Latex[length=2.3mm]}, thick}
    ]

    \node[layer, fill=blue!8] (app) at (0,3.0) {应用层：训练脚本 / 推理服务 / 调度系统};
    \node[layer, fill=green!8] (fw) at (0,2.05) {框架层：PyTorch / JAX / TensorFlow / vLLM / DeepSpeed};
    \node[layer, fill=yellow!14] (op) at (0,1.1) {Operator 层：Linear / Attention / LayerNorm / Softmax};
    \node[layer, fill=orange!14] (lib) at (0,0.15) {库与 Kernel：cuBLAS / cuDNN / NCCL / Triton / Custom CUDA};
    \node[layer, fill=red!10] (rt) at (0,-0.8) {CUDA Runtime / Driver：Stream / Event / Memory / Launch};
    \node[layer, fill=purple!10] (hw) at (0,-1.75) {GPU Hardware：SM / Tensor Core / L2 / HBM / NVLink};

    \draw[arrow] (app) -- (fw);
    \draw[arrow] (fw) -- (op);
    \draw[arrow] (op) -- (lib);
    \draw[arrow] (lib) -- (rt);
    \draw[arrow] (rt) -- (hw);

    \node[align=left, text width=4.3cm] at (6.5,0.35) {
       \textbf{性能问题定位：}\\
       需要从 Python 代码一路追踪到 kernel、通信和硬件计数器。
    };
  \end{tikzpicture}
  \caption{从深度学习框架到 CUDA kernel 的软件栈}
  \label{fig:dl-framework-cuda-stack}
\end{figure}

在分布式训练中，NCCL 的重要性尤其突出。DDP 的梯度同步、Tensor Parallel 的 all-reduce、ZeRO 的 reduce-scatter 和 all-gather，最终大多依赖 NCCL collective。对于大规模集群，一个训练任务能否扩展到数百甚至数千卡，不仅取决于模型结构，也取决于 NCCL 对拓扑、带宽、延迟和通信重叠的利用。

\subsection{自定义 CUDA Extension}
\label{subsec:custom-cuda-extension}

当框架内置算子无法满足需求时，我们可能需要编写自定义 CUDA extension。常见动机包括：

\begin{itemize}
  \item 将多个小算子融合成一个 kernel，减少 HBM 读写和 kernel launch；
  \item 实现框架暂不支持的新算子；
  \item 针对特定 shape、dtype 或硬件进行优化；
  \item 实现自定义量化、稀疏计算或特殊 attention 模式。
\end{itemize}

下面给出一个极简 PyTorch CUDA extension 的结构。真实工程中还需要处理错误检查、dtype dispatch、contiguous 检查、反向传播绑定和构建脚本。

\begin{lstlisting}[language=C++,caption={PyTorch CUDA Extension 的 C++ 绑定示意},label={lst:pytorch-cuda-extension-cpp}]
#include <torch/extension.h>

// CUDA launcher declaration.
torch::Tensor scale_forward_cuda(torch::Tensor x, float scale);

torch::Tensor scale_forward(torch::Tensor x, float scale) {
  TORCH_CHECK(x.is_cuda(), "x must be a CUDA tensor");
  TORCH_CHECK(x.is_contiguous(), "x must be contiguous");

  return scale_forward_cuda(x, scale);
}

PYBIND11_MODULE(TORCH_EXTENSION_NAME, m) {
  m.def("scale_forward", &scale_forward, "Scale tensor forward on CUDA");
}
\end{lstlisting}

对应的 CUDA kernel 可以写成：

\begin{lstlisting}[language=C++,caption={PyTorch CUDA Extension 的 CUDA Kernel 示意},label={lst:pytorch-cuda-extension-cu}]
#include <torch/extension.h>
#include <cuda.h>
#include <cuda_runtime.h>

template <typename scalar_t>
__global__ void scale_kernel(const scalar_t* x,
                             scalar_t* y,
                             float scale,
                             int64_t n) {
  int idx = blockIdx.x * blockDim.x + threadIdx.x;
  if (idx < n) {
    y[idx] = static_cast<scalar_t>(x[idx] * scale);
  }
}

torch::Tensor scale_forward_cuda(torch::Tensor x, float scale) {
  auto y = torch::empty_like(x);
  int64_t n = x.numel();

  int threads = 256;
  int blocks = (n + threads - 1) / threads;

  AT_DISPATCH_FLOATING_TYPES(x.scalar_type(), "scale_forward_cuda", ([&] {
    scale_kernel<scalar_t><<<blocks, threads>>>(
      x.data_ptr<scalar_t>(),
      y.data_ptr<scalar_t>(),
      scale,
      n
    );
  }));

  return y;
}
\end{lstlisting}

在生产环境中，自定义 CUDA extension 的复杂度往往不在 kernel 本身，而在工程边界：

\begin{itemize}
  \item 如何支持 FP16、BF16、FP32、INT8 等 dtype；
  \item 如何处理非 contiguous tensor；
  \item 如何与 PyTorch autograd 对接；
  \item 如何避免 silent correctness bug；
  \item 如何在不同 GPU 架构上保持性能；
  \item 如何进行单元测试、基准测试和 profiler 分析；
  \item 如何在容器、CI、驱动版本和 CUDA Toolkit 版本之间保持兼容。
\end{itemize}

\subsection{从 kernel 到 operator}
\label{subsec:from-kernel-to-operator}

一个 operator 不一定只对应一个 kernel。例如标准 attention 可以拆成：

\begin{itemize}
  \item QKV projection GEMM；
  \item QK 转置乘法；
  \item mask；
  \item softmax；
  \item dropout；
  \item attention weights 与 V 相乘；
  \item output projection GEMM。
\end{itemize}

如果每一步都是独立 kernel，那么中间张量需要反复写回 HBM，再从 HBM 读出。对于长序列 attention，中间的 attention score 形如：
\[
  (B,\;H_{\mathrm{heads}},\;S,\;S).
\]
当 $S$ 很大时，这个张量会占用巨大显存和带宽。FlashAttention 的关键思想之一，就是把多个操作融合在一个分块 kernel 内，并避免物化完整 attention matrix。

从 kernel 到 operator，可以用如下层次理解：

\begin{table}[H]
  \centering
  \small
  \caption{Kernel、Operator 与 Model Layer 的关系}
  \label{tab:kernel-operator-layer}
  \begin{tabular}{p{0.20\textwidth}p{0.34\textwidth}p{0.34\textwidth}}
    \toprule
    \textbf{层次} & \textbf{例子} & \textbf{工程关注点} \\
    \midrule
    Kernel
    & 一个 CUDA 函数，例如 vector add、LayerNorm kernel、GEMM kernel
    & 线程映射、寄存器、shared memory、访存合并、occupancy。 \\
    \midrule
    Operator
    & PyTorch 中的 \texttt{matmul}、\texttt{layer\_norm}、\texttt{scaled\_dot\_product\_attention}
    & dtype dispatch、layout、autograd、workspace、kernel 选择。 \\
    \midrule
    Model Layer
    & Transformer block 中的 attention、MLP、norm、residual
    & 算子组合、显存占用、激活保存、通信切分、并行策略。 \\
    \midrule
    End-to-end System
    & 训练任务或推理服务
    & 调度、batching、通信、checkpoint、容错、SLA 和成本。 \\
    \bottomrule
  \end{tabular}
\end{table}

很多 AI Infra 优化的本质，是把高层模型结构重新组织为更适合硬件执行的 operator 和 kernel。例如：

\begin{itemize}
  \item 把 bias add、activation、dropout 融合；
  \item 把 LayerNorm 与 residual add 融合；
  \item 使用 fused optimizer 减少参数更新阶段的小 kernel；
  \item 使用 FlashAttention 避免 attention score 物化；
  \item 使用 paged KV Cache 改善推理阶段内存管理；
  \item 使用 Tensor Parallel 把大 GEMM 切分到多个 GPU 上。
\end{itemize}

\section{CUDA 编程中的工程痛点}
\label{sec:cuda-engineering-pain-points}

本章已经介绍了 CUDA 的基本概念，但真实工程远比教材示例复杂。这里提前总结几个常见痛点，帮助读者建立问题雷达。

\subsection{异步执行导致的错误定位困难}
\label{subsec:async-debugging}

CUDA kernel launch 通常是异步的。CPU 提交 kernel 后继续执行，错误可能在后续同步点才暴露。例如某个 kernel 越界访问，报错可能出现在后面的 \texttt{cudaMemcpy} 或另一个无关 API 调用处。

调试时常用方法包括：

\begin{itemize}
  \item 在关键位置调用 \texttt{cudaDeviceSynchronize()} 缩小错误范围；
  \item 检查每个 CUDA API 返回值；
  \item 使用 \texttt{cuda-memcheck} 或 compute sanitizer；
  \item 使用较小输入复现问题；
  \item 为自定义 kernel 编写 CPU reference 结果做对比。
\end{itemize}

在训练框架中，异步错误更隐蔽。一个 CUDA illegal memory access 可能导致后续所有 CUDA 调用失败，最终表现为 PyTorch 抛出一个看似不相关的异常。因此，定位 CUDA 错误时，要始终记住：\textbf{报错位置不一定是出错位置。}

\subsection{小 kernel 与 launch overhead}
\label{subsec:small-kernel-launch-overhead}

GPU 擅长大规模并行计算，但 kernel launch 本身有固定开销。如果模型中存在大量小张量操作，每个操作单独 launch 一个 kernel，那么端到端性能可能被 launch overhead 和 HBM 往返吞噬。

这在 Transformer 中尤其常见。例如一个 block 内有许多 elementwise 操作：

\begin{itemize}
  \item residual add；
  \item bias add；
  \item dropout；
  \item activation；
  \item norm；
  \item mask fill；
  \item scale。
\end{itemize}

如果这些操作拆得太碎，GPU 时间线会出现大量短 kernel，中间夹杂频繁 HBM 读写。优化策略通常是 kernel fusion。以 MLP 为例，可以把 bias add 和 GELU 融合到 GEMM epilogue 中，减少一次读写。

\subsection{非连续张量与隐式拷贝}
\label{subsec:non-contiguous-hidden-copy}

上一章讨论过 shape 与 stride。这里从 CUDA 角度再强调一次：很多高性能 kernel 假设输入 tensor 是 contiguous 或满足特定 stride。如果上层代码传入非连续张量，框架可能隐式插入 \texttt{contiguous()}，导致额外 HBM 拷贝。

例如：

\begin{lstlisting}[language=Python,caption={非连续张量可能触发额外拷贝},label={lst:non-contiguous-copy}]
import torch

x = torch.randn(8, 2048, 4096, device="cuda", dtype=torch.float16)

# Transpose changes stride but usually does not move data immediately.
y = x.transpose(1, 2)

print(y.is_contiguous())  # False

# Some custom CUDA kernels may require contiguous input.
# This line may allocate a new tensor and copy a large amount of data.
z = y.contiguous()
\end{lstlisting}

在大模型训练中，一个激活张量可能非常大。一次隐式 contiguous copy 可能比某些计算本身更昂贵。性能优化时，要特别检查 profiler 中是否出现意料之外的 memory copy 或 layout transform。

\subsection{数值精度与硬件路径}
\label{subsec:numeric-precision-hardware-path}

深度学习 kernel 的性能与 dtype 密切相关。FP32、TF32、FP16、BF16、FP8、INT8、INT4 会走不同硬件路径。以矩阵乘法为例，使用 Tensor Core 的低精度矩阵乘通常远快于普通 FP32 CUDA core 路径。

但精度不是越低越好。工程上必须同时考虑：

\begin{itemize}
  \item 模型精度是否可接受；
  \item accumulation 使用什么精度；
  \item 是否需要 loss scaling；
  \item 是否出现 overflow 或 underflow；
  \item kernel 是否真正使用了 Tensor Core；
  \item shape 是否满足 Tensor Core 友好的对齐要求。
\end{itemize}

例如，某些 GEMM 维度如果不是特定倍数，可能无法充分利用 Tensor Core。模型架构设计和 batch size 选择有时也会考虑硬件对齐，这就是算法与系统共同设计的一个例子。

\section{本章小结}
\label{sec:chapter02-summary}

本章从 CPU 与 GPU 编程模型的差异出发，介绍了 CUDA C++ 的核心概念，并把它们与深度学习系统中的真实问题联系起来。

\begin{itemize}
  \item \textbf{CPU 与 GPU 的设计目标不同}：CPU 强控制流、低延迟；GPU 强吞吐、海量并行。AI 系统需要二者协同。
  \item \textbf{SIMT 与 Warp 是 GPU 执行的基本抽象}：程序员写线程逻辑，硬件以 warp 为单位调度。warp divergence 会降低执行效率。
  \item \textbf{Grid、Block、Thread 构成 CUDA 线程层次结构}：block 是共享 shared memory 和同步的基本边界。
  \item \textbf{Occupancy 是理解 GPU 延迟隐藏的重要概念}：但它不是唯一目标，必须与 register、shared memory、访存带宽和 Tensor Core 利用率一起分析。
  \item \textbf{CUDA 内存层级决定 kernel 性能上限}：global memory 容量大但延迟高；shared memory 适合 block 内复用；register 最快但数量有限。
  \item \textbf{Memory coalescing 对带宽利用至关重要}：连续线程访问连续地址通常更高效，这与张量 layout 和 stride 密切相关。
  \item \textbf{Tiled GEMM 展示了高性能 kernel 的基本思想}：通过 shared memory 复用数据，提高算术强度，减少 HBM 访问。
  \item \textbf{深度学习框架最终仍会落到 CUDA kernel 和通信库}：PyTorch operator、cuBLAS、cuDNN、NCCL、自定义 extension 构成了 AI Infra 的执行底座。
\end{itemize}

下一章将回到深度学习模型本身，系统讲解神经网络基本结构、反向传播、优化器、归一化与混合精度训练。届时，本章中的 kernel、显存、dtype 和硬件执行路径，会与模型训练中的参数、梯度、激活和优化器状态连接起来。
\chapter{深度学习基本原理}
\label{chap:deep-learning-basics}

\section{神经网络的基本结构}
\label{sec:nn-basic-structure}

前两章分别从数学与 CUDA 的角度铺垫了 AI Infrastructure 的底层基础。本章开始，我们回到深度学习模型本身。对于 AI Infra 工程师而言，理解神经网络并不是为了从零发明一个新模型，而是为了回答一系列工程问题：

\begin{itemize}
\item 一个模型的参数量从哪里来？
\item 训练时显存为什么远大于模型权重本身？
\item 为什么反向传播需要保存激活？
\item 为什么优化器状态会占用巨大显存？
\item 为什么大模型训练偏好 AdamW、warmup、cosine decay、LayerNorm、RMSNorm 和混合精度？
\item 为什么同样的数学结构，在 GPU 上的实现方式会决定吞吐和稳定性？
\end{itemize}

本章的目标，是把神经网络的数学结构与工程系统中的真实成本联系起来。我们会从线性层、激活函数和 loss function 出发，推导反向传播，分析优化器，比较归一化方法，并解释混合精度训练为什么成为现代大模型训练的默认配置。

\subsection{线性层}
\label{subsec:linear-layer}

线性层是神经网络最基本的构件。给定输入：
\[
\mat{X}\in\mathbb{R}^{B\times d_{\mathrm{in}}},
\]
权重：
\[
\mat{W}\in\mathbb{R}^{d_{\mathrm{in}}\times d_{\mathrm{out}}},
\]
偏置：
\[
\vect{b}\in\mathbb{R}^{d_{\mathrm{out}}},
\]
线性层输出为：
\[
\mat{Y}=\mat{X}\mat{W}+\vect{b}.
\]

其中 $B$ 是 batch size，$d_{\mathrm{in}}$ 是输入维度，$d_{\mathrm{out}}$ 是输出维度。偏置 $\vect{b}$ 会沿 batch 维广播：
\[
Y_{i,j}=\sum_{k=1}^{d_{\mathrm{in}}}X_{i,k}W_{k,j}+b_j.
\]

在 PyTorch 中，\texttt{nn.Linear(d_in, d_out)} 默认保存的权重 shape 通常是：
\[
(d_{\mathrm{out}},d_{\mathrm{in}}),
\]
而数学书写中我们常使用：
\[
\mat{W}\in\mathbb{R}^{d_{\mathrm{in}}\times d_{\mathrm{out}}}.
\]
这只是约定差异。底层 GEMM 会根据 transpose 参数、stride 和 layout 选择实际计算方式。

从参数量角度看，线性层参数为：
\[
d_{\mathrm{in}}d_{\mathrm{out}}+d_{\mathrm{out}}.
\]
如果 $d_{\mathrm{in}}=d_{\mathrm{out}}=4096$，单个线性层权重参数量约为：
\[
4096^2=16,777,216.
\]
使用 BF16 或 FP16 保存，仅权重就需要约：
\[
16,777,216\times 2\ \mathrm{bytes}\approx 32\ \mathrm{MB}.
\]
在 Transformer 中，attention 的 QKV projection、output projection、FFN 的 up projection 和 down projection 都是线性层。因此，线性层不仅是数学基础，也是参数量、FLOPs 与显存占用的主要来源。

\subsection{激活函数}
\label{subsec:activation-functions}

如果神经网络只由线性层组成，那么无论叠多少层，整体仍然是一个线性变换。设：
\[
\mat{H}_1=\mat{X}\mat{W}_1,\quad
\mat{Y}=\mat{H}_1\mat{W}_2,
\]
则：
\[
\mat{Y}=\mat{X}(\mat{W}_1\mat{W}_2).
\]
这等价于一个更大的线性层。激活函数的作用，就是引入非线性，使神经网络能够拟合复杂函数。

常见激活函数包括 ReLU、Sigmoid、Tanh、GELU 和 SiLU。它们的定义如下：

\[
\operatorname{ReLU}(x)=\max(0,x),
\]
\[
\sigma(x)=\frac{1}{1+\exp(-x)},
\]
\[
\tanh(x)=\frac{\exp(x)-\exp(-x)}{\exp(x)+\exp(-x)}.
\]

GELU 常用于 Transformer 早期架构，其定义可写为：
\[
\operatorname{GELU}(x)=x\Phi(x),
\]
其中 $\Phi(x)$ 是标准正态分布的累积分布函数。工程实现中常用近似：
\[
\operatorname{GELU}(x)
\approx
\frac{1}{2}x
\left(
1+\tanh
\left[
\sqrt{\frac{2}{\pi}}
\left(x+0.044715x^3\right)
\right]
\right).
\]

现代 LLM 中，SwiGLU 这类门控激活非常常见。若输入经过两个线性投影：
\[
\mat{U}=\mat{X}\mat{W}_u,\quad
\mat{G}=\mat{X}\mat{W}_g,
\]
则 SwiGLU 可写作：
\[
\operatorname{SwiGLU}(\mat{X})
==============================

\operatorname{SiLU}(\mat{G})\odot \mat{U},
\]
其中：
\[
\operatorname{SiLU}(x)=x\sigma(x).
\]

从工程上看，激活函数大多属于 elementwise 操作，单独执行时往往是 memory-bound。若一个激活函数只做少量计算，却需要从 HBM 读取输入、写回输出，那么它很可能不能充分利用 GPU 算力。因此，在高性能训练和推理系统中，激活函数常常与 bias add、dropout 或 GEMM epilogue 融合，减少中间张量写回。

\begin{figure}[H]
\centering
\begin{tikzpicture}[
font=\small,
axis/.style={-{Latex[length=2.3mm]}, thick},
curve/.style={thick},
sample/.style={circle, inner sep=1.2pt, fill=black}
]

\node[font=\bfseries] at (0,3.0) {常见激活函数的形状直觉};

\draw[axis] (-5.0,0) -- (-1.1,0) node[right] {$x$};
\draw[axis] (-3.05,-1.0) -- (-3.05,1.8) node[above] {$y$};
\draw[curve, blue!70]
  (-4.8,0) -- (-3.05,0) -- (-1.35,1.35);
\node at (-3.05,-1.35) {ReLU};

\draw[axis] (-0.8,0) -- (3.1,0) node[right] {$x$};
\draw[axis] (1.15,-1.0) -- (1.15,1.8) node[above] {$y$};
\draw[curve, green!60!black]
  plot[smooth] coordinates {
    (-0.55,0.03) (0.0,0.08) (0.55,0.22) (1.15,0.5)
    (1.75,0.78) (2.3,0.92) (2.85,0.97)
  };
\node at (1.15,-1.35) {Sigmoid};

\draw[axis] (3.4,0) -- (7.3,0) node[right] {$x$};
\draw[axis] (5.35,-1.0) -- (5.35,1.8) node[above] {$y$};
\draw[curve, orange!80!black]
  plot[smooth] coordinates {
    (3.65,-0.15) (4.1,-0.12) (4.55,-0.05) (5.0,0.0)
    (5.35,0.0) (5.8,0.2) (6.25,0.62) (6.75,1.18) (7.05,1.55)
  };
\node at (5.35,-1.35) {GELU / SiLU};

\node[align=left, text width=10.6cm] at (1.0,-2.25) {
  激活函数提供非线性表达能力。Infra 视角下，它们通常是逐元素操作，单独执行时算术强度较低，因此常与其他 elementwise 操作或 GEMM epilogue 融合。
};

\end{tikzpicture}
\caption{ReLU、Sigmoid、GELU/SiLU 的函数形状直觉}
\label{fig:activation-functions}
\end{figure}

\subsection{Loss Function}
\label{subsec:loss-function}

训练神经网络时，我们需要一个目标函数衡量模型输出与真实标签之间的差异。这个目标函数称为损失函数。优化器通过最小化 loss 来更新参数。

对于分类任务，常用交叉熵损失。设模型输出 logits：
\[
\vect{z}\in\mathbb{R}^{C},
\]
经过 softmax 得到类别概率：
\[
p_i=\frac{\exp(z_i)}{\sum_{j=1}^{C}\exp(z_j)}.
\]
若真实类别为 $y$，则交叉熵损失为：
\[
\mathcal{L}=-\log p_y.
\]
展开可得：
\[
\mathcal{L}
===========

-z_y+\log\sum_{j=1}^{C}\exp(z_j).
\]

对于语言模型，类别数 $C$ 就是词表大小 $V$。每个位置都要预测下一个 token，因此 batch 中的 logits 常见 shape 为：
\[
(B,S,V).
\]
若 $B=8$、$S=2048$、$V=50000$，则 logits 元素个数为：
\[
8\times 2048\times 50000=819,200,000.
\]
使用 BF16 保存也需要超过 $1.6$ GB 显存。实际训练中，cross entropy 往往需要 fused kernel 或 vocabulary parallelism 来降低显存和通信压力。

对于回归任务，常用均方误差：
\[
\mathcal{L}_{\mathrm{MSE}}
==========================

\frac{1}{N}\sum_{i=1}^{N}(\hat{y}_i-y_i)^2.
\]
大模型预训练主要使用 next-token cross entropy，但强化学习、奖励模型、偏好优化、多模态对齐等任务还会使用更多损失函数。作为 Infra 工程师，我们需要重点关注 loss 的两个方面：

\begin{itemize}
\item \textbf{数值稳定性}：是否容易出现 overflow、underflow、NaN 或 inf。
\item \textbf{系统代价}：是否需要物化巨大中间张量，是否能被 fuse，是否会引入跨 GPU 通信。
\end{itemize}

\subsection{模型参数与中间激活}
\label{subsec:parameters-and-activations}

神经网络训练时，显存并不仅仅存储模型参数。考虑一个简单的 MLP：
\[
\mat{H}_1=\phi(\mat{X}\mat{W}_1+\vect{b}_1),
\]
\[
\mat{H}_2=\phi(\mat{H}_1\mat{W}_2+\vect{b}_2),
\]
\[
\mat{Y}=\mat{H}_2\mat{W}_3+\vect{b}_3.
\]
训练时除了参数 $\mat{W}_1,\mat{W}_2,\mat{W}_3$，还需要保存中间激活 $\mat{X},\mat{H}_1,\mat{H}_2$，以便反向传播计算梯度。

\begin{figure}[H]
\centering
\begin{tikzpicture}[
font=\small,
layer/.style={draw, rounded corners=3pt, minimum width=1.55cm, minimum height=0.9cm, align=center, fill=blue!6},
act/.style={draw, rounded corners=3pt, minimum width=1.4cm, minimum height=0.75cm, align=center, fill=green!8},
param/.style={draw, rounded corners=2pt, fill=orange!12, align=center, minimum width=1.2cm, minimum height=0.5cm},
arrow/.style={-{Latex[length=2.3mm]}, thick}
]

\node[act] (x) at (-5.2,0) {$\mat{X}$};
\node[layer] (l1) at (-3.6,0) {Linear\\$W_1,b_1$};
\node[layer] (a1) at (-1.9,0) {Activation\\$\phi$};
\node[act] (h1) at (-0.4,0) {$\mat{H}_1$};
\node[layer] (l2) at (1.2,0) {Linear\\$W_2,b_2$};
\node[layer] (a2) at (2.9,0) {Activation\\$\phi$};
\node[act] (h2) at (4.4,0) {$\mat{H}_2$};
\node[layer] (l3) at (6.0,0) {Linear\\$W_3,b_3$};
\node[act] (y) at (7.45,0) {$\mat{Y}$};

\draw[arrow] (x) -- (l1);
\draw[arrow] (l1) -- (a1);
\draw[arrow] (a1) -- (h1);
\draw[arrow] (h1) -- (l2);
\draw[arrow] (l2) -- (a2);
\draw[arrow] (a2) -- (h2);
\draw[arrow] (h2) -- (l3);
\draw[arrow] (l3) -- (y);

\node[param] at (-3.6,1.2) {参数};
\node[param] at (1.2,1.2) {参数};
\node[param] at (6.0,1.2) {参数};

\draw[dashed, rounded corners=3pt, fill=yellow!8]
  (-5.85,-1.0) rectangle (5.05,-0.45);
\node at (-0.4,-0.72) {训练时需要保存的中间激活，用于 backward};

\node[align=left, text width=10.8cm] at (1.0,-2.0) {
  前向传播得到输出；反向传播需要用前向中的输入和中间结果计算梯度。
  因此训练显存通常远高于仅加载模型权重所需的显存。
};

\end{tikzpicture}
\caption{多层感知机的前向传播与中间激活}
\label{fig:mlp-forward-activations}
\end{figure}

参数是模型的长期状态；激活是一次 forward/backward 中的临时状态。对于推理，很多激活可以在层间流动后立即释放；对于训练，激活需要保留到对应 backward 使用完毕。这个差异是训练显存远高于推理显存的重要原因之一。

\section{反向传播}
\label{sec:backpropagation}

\subsection{从链式法则到计算图}
\label{subsec:chain-rule-to-computation-graph}

反向传播是链式法则在计算图上的系统化应用。假设一个简单标量函数：
\[
y=(xw+b)^2.
\]
可以拆成：
\[
u=xw,\quad
v=u+b,\quad
y=v^2.
\]
根据链式法则：
\[
\frac{\partial y}{\partial w}
=============================

\frac{\partial y}{\partial v}
\frac{\partial v}{\partial u}
\frac{\partial u}{\partial w}.
\]
逐项计算：
\[
\frac{\partial y}{\partial v}=2v,\quad
\frac{\partial v}{\partial u}=1,\quad
\frac{\partial u}{\partial w}=x.
\]
因此：
\[
\frac{\partial y}{\partial w}=2v x.
\]

深度学习框架会把复杂模型拆成许多基本算子，每个算子记录自己的前向输入、输出和 backward 规则。反向传播从 loss 开始，按照计算图的反向拓扑顺序，把梯度传播回所有需要更新的参数。

\begin{figure}[H]
\centering
\begin{tikzpicture}[
font=\small,
nodebox/.style={draw, rounded corners=2pt, minimum width=1.1cm, minimum height=0.65cm, align=center, fill=blue!6},
op/.style={draw, circle, minimum size=0.72cm, fill=orange!12},
fwd/.style={-{Latex[length=2.2mm]}, thick, blue!70},
bwd/.style={-{Latex[length=2.2mm]}, thick, red!70}
]

\node[nodebox] (x) at (-4.8,1.1) {$x$};
\node[nodebox] (w) at (-4.8,-0.1) {$w$};
\node[op] (mul) at (-3.0,0.5) {$\times$};
\node[nodebox] (u) at (-1.7,0.5) {$u$};

\node[nodebox] (b) at (-1.7,-0.9) {$b$};
\node[op] (add) at (-0.2,0.0) {$+$};
\node[nodebox] (v) at (1.1,0.0) {$v$};
\node[op] (sq) at (2.5,0.0) {$^2$};
\node[nodebox] (y) at (3.8,0.0) {$y$};

\draw[fwd] (x) -- (mul);
\draw[fwd] (w) -- (mul);
\draw[fwd] (mul) -- (u);
\draw[fwd] (u) -- (add);
\draw[fwd] (b) -- (add);
\draw[fwd] (add) -- (v);
\draw[fwd] (v) -- (sq);
\draw[fwd] (sq) -- (y);

\draw[bwd] (y.south) .. controls (3.2,-1.4) and (2.5,-1.4) .. (sq.south);
\draw[bwd] (sq.south) .. controls (2.1,-1.65) and (1.3,-1.65) .. (v.south);
\draw[bwd] (v.south) .. controls (0.7,-1.75) and (-0.2,-1.75) .. (add.south);
\draw[bwd] (add.south) .. controls (-0.8,-1.9) and (-1.5,-1.9) .. (u.south);
\draw[bwd] (u.south) .. controls (-2.2,-1.8) and (-3.0,-1.4) .. (mul.south);
\draw[bwd] (mul.south) .. controls (-3.6,-1.2) and (-4.4,-0.8) .. (w.south);

\node[align=left, text width=10.5cm] at (-0.4,-3.0) {
  蓝色路径是 forward，红色路径是 backward。每个局部算子只需要知道自己的局部导数，
  整体梯度由链式法则沿图传播得到。
};

\end{tikzpicture}
\caption{从链式法则到计算图反向传播}
\label{fig:chain-rule-computation-graph}
\end{figure}

\subsection{梯度在图中的流动}
\label{subsec:gradient-flow-in-graph}

考虑一个线性层：
\[
\mat{Y}=\mat{X}\mat{W}+\vect{b}.
\]
设上游梯度为：
\[
\mat{G}_Y=\frac{\partial \mathcal{L}}{\partial \mat{Y}}.
\]
则对应梯度为：
\[
\frac{\partial \mathcal{L}}{\partial \mat{X}}
=============================================

\mat{G}_Y\mat{W}^{\mathsf{T}},
\]
\[
\frac{\partial \mathcal{L}}{\partial \mat{W}}
=============================================

\mat{X}^{\mathsf{T}}\mat{G}_Y,
\]
\[
\frac{\partial \mathcal{L}}{\partial \vect{b}}
==============================================

\sum_{i=1}^{B}\mat{G}_{Y,i,:}.
\]

这三条公式非常重要，因为它们揭示了训练中 GEMM 的来源。一个线性层 forward 是 GEMM，backward 中对输入和权重的梯度计算也都是 GEMM。也就是说，训练一个线性层通常包含：

\begin{itemize}
\item forward GEMM：$\mat{X}\mat{W}$；
\item backward data GEMM：$\mat{G}_Y\mat{W}^{\mathsf{T}}$；
\item backward weight GEMM：$\mat{X}^{\mathsf{T}}\mat{G}_Y$。
\end{itemize}

因此，训练 FLOPs 通常显著高于单次推理 forward。粗略地说，一个线性层训练中的计算量可近似为 forward 的 $3$ 倍左右，还不包括优化器更新、归一化、激活函数和通信。

下面给出一个手动推导线性层 backward 的 PyTorch 风格示例：

\begin{lstlisting}[language=Python,caption={线性层 forward 与 backward 公式的张量实现},label={lst:linear-backward-manual}]
import torch

B, Din, Dout = 8, 1024, 4096

X = torch.randn(B, Din, device="cuda", dtype=torch.float32)
W = torch.randn(Din, Dout, device="cuda", dtype=torch.float32)
b = torch.randn(Dout, device="cuda", dtype=torch.float32)

# Forward

Y = X @ W + b

# Suppose this is the upstream gradient dL/dY.

dY = torch.randn_like(Y)

# Backward formulas

dX = dY @ W.T          # dL/dX
dW = X.T @ dY          # dL/dW
db = dY.sum(dim=0)     # dL/db

print(dX.shape, dW.shape, db.shape)
\end{lstlisting}

在真实训练框架中，这些梯度计算由 autograd 自动生成或调用高性能 backend 实现。AI Infra 工程师需要知道的是：\textbf{反向传播不是抽象魔法，而是一组具体的张量算子、GEMM、reduction、elementwise kernel 与通信操作。}

\subsection{激活保存与显存占用}
\label{subsec:activation-memory-consumption}

反向传播需要前向中间结果。以线性层为例，计算：
\[
\frac{\partial \mathcal{L}}{\partial \mat{W}}
=============================================

\mat{X}^{\mathsf{T}}\mat{G}_Y
\]
需要保存输入 $\mat{X}$；计算某些激活函数的 backward 也需要保存前向输出或 mask。例如 ReLU backward 需要知道前向时哪些位置大于零：
\[
\frac{\partial \operatorname{ReLU}(x)}{\partial x}
==================================================

\begin{cases}
1, & x>0,\
0, & x\le 0.
\end{cases}
\]

对于一个 $L$ 层网络，若每层保存激活：
\[
\tensor{A}*{\ell}\in\mathbb{R}^{B\times S\times H},
\]
则仅隐藏状态激活显存粗略为：
\[
M*{\mathrm{hidden}}
\approx
LBSH\cdot s,
\]
其中 $s$ 是每个元素的字节数。实际显存会更高，因为 attention、MLP、dropout、norm 和 residual 都可能保存额外中间结果。

\begin{figure}[H]
\centering
\begin{tikzpicture}[
font=\small,
layer/.style={draw, rounded corners=2pt, minimum width=1.05cm, minimum height=0.65cm, align=center, fill=blue!6},
act/.style={draw, rounded corners=2pt, minimum width=0.95cm, minimum height=0.5cm, align=center, fill=yellow!16},
arrow/.style={-{Latex[length=2.1mm]}, thick},
bwd/.style={-{Latex[length=2.1mm]}, thick, red!70}
]

\node[font=\bfseries] at (0,2.8) {前向保存激活，反向逐层消费};

\foreach \i/\x in {1/-4.8,2/-3.2,3/-1.6,4/0.0,5/1.6,6/3.2,7/4.8} {
  \node[layer] (l\i) at (\x,0.8) {Layer \i};
  \node[act] (a\i) at (\x,-0.15) {$A_\i$};
  \draw[dashed] (l\i) -- (a\i);
}

\foreach \i/\j in {1/2,2/3,3/4,4/5,5/6,6/7} {
  \draw[arrow] (l\i.east) -- (l\j.west);
}

\draw[bwd] (4.8,-0.85) -- (3.2,-0.85);
\draw[bwd] (3.2,-0.85) -- (1.6,-0.85);
\draw[bwd] (1.6,-0.85) -- (0.0,-0.85);
\draw[bwd] (0.0,-0.85) -- (-1.6,-0.85);
\draw[bwd] (-1.6,-0.85) -- (-3.2,-0.85);
\draw[bwd] (-3.2,-0.85) -- (-4.8,-0.85);

\node[align=center] at (0,-1.35) {Backward 从最后一层开始，逐层读取对应 forward 激活并释放};

\node[draw, rounded corners=3pt, fill=green!8, align=left, text width=10.4cm] at (0,-2.55) {
  \textbf{显存痛点：} 层数越深、batch 越大、序列越长、hidden size 越大，激活显存越高。
  在大模型训练中，激活显存常常与参数、梯度、优化器状态一起成为限制 batch size 和 sequence length 的关键因素。
};

\end{tikzpicture}
\caption{前向激活保存与反向传播中的激活消费}
\label{fig:activation-save-backward-consume}
\end{figure}

激活显存的峰值通常发生在 forward 结束、backward 开始之前。此时大部分层的激活都还没有被释放。之后 backward 从后向前执行，每处理完一层，该层对应的激活就可以释放。这个生命周期是设计 activation checkpointing、pipeline parallelism 和 ZeRO activation partitioning 的基础。

\subsection{Gradient Checkpointing 的动机}
\label{subsec:gradient-checkpointing}

Gradient Checkpointing，也常称为 activation checkpointing，其核心思想是：
\[
\text{少存激活，多做重计算。}
\]
在 forward 阶段，不保存所有中间激活，只保存少数 checkpoint。到了 backward 阶段，如果需要某段中间激活，就从最近的 checkpoint 重新执行一段 forward 来恢复。

假设一个网络有 $L$ 层。普通训练保存每层激活，显存约为：
\[
O(L).
\]
采用 checkpointing 后，可以把激活显存降到更低，例如在某些理想策略下接近：
\[
O(\sqrt{L}),
\]
代价是额外 forward recomputation，计算量上升。

\begin{figure}[H]
\centering
\begin{tikzpicture}[
font=\small,
layer/.style={draw, rounded corners=2pt, minimum width=0.85cm, minimum height=0.55cm, align=center, fill=blue!6},
ckpt/.style={draw, rounded corners=2pt, minimum width=0.85cm, minimum height=0.55cm, align=center, fill=orange!20},
arrow/.style={-{Latex[length=2.1mm]}, thick}
]

\node[font=\bfseries] at (0,2.8) {Gradient Checkpointing：保存少数节点，反向时重计算};

\foreach \i/\x in {1/-5.4,2/-4.2,3/-3.0,4/-1.8,5/-0.6,6/0.6,7/1.8,8/3.0,9/4.2,10/5.4} {
  \node[layer] (l\i) at (\x,1.0) {$L_\i$};
}
\foreach \i/\j in {1/2,2/3,3/4,4/5,5/6,6/7,7/8,8/9,9/10} {
  \draw[arrow] (l\i.east) -- (l\j.west);
}

\node[ckpt] (c1) at (-5.4,-0.2) {$C_1$};
\node[ckpt] (c5) at (-0.6,-0.2) {$C_5$};
\node[ckpt] (c10) at (5.4,-0.2) {$C_{10}$};

\draw[dashed] (l1) -- (c1);
\draw[dashed] (l5) -- (c5);
\draw[dashed] (l10) -- (c10);

\draw[decorate,decoration={brace,amplitude=5pt},thick]
  (-5.85,-0.75) -- (-0.15,-0.75)
  node[midway,below=7pt] {反向时从 $C_1$ 重算 $L_2$ 到 $L_5$};

\draw[decorate,decoration={brace,amplitude=5pt},thick]
  (-0.15,-1.35) -- (5.85,-1.35)
  node[midway,below=7pt] {反向时从 $C_5$ 重算 $L_6$ 到 $L_{10}$};

\node[align=left, text width=10.4cm] at (0,-2.65) {
  Checkpointing 降低激活显存峰值，但增加计算量。它尤其适合显存受限但计算资源相对充足的场景，例如长上下文训练或大 batch 训练。
};

\end{tikzpicture}
\caption{Gradient Checkpointing 的保存与重计算机制}
\label{fig:gradient-checkpointing}
\end{figure}

PyTorch 中可以使用 \texttt{torch.utils.checkpoint} 实现 checkpointing：

\begin{lstlisting}[language=Python,caption={PyTorch 中使用 activation checkpointing},label={lst:pytorch-gradient-checkpointing}]
import torch
import torch.nn as nn
from torch.utils.checkpoint import checkpoint

class Block(nn.Module):
def **init**(self, hidden):
super().**init**()
self.net = nn.Sequential(
nn.Linear(hidden, 4 * hidden),
nn.GELU(),
nn.Linear(4 * hidden, hidden),
)

def forward(self, x):
    return x + self.net(x)

class Model(nn.Module):
def **init**(self, hidden, num_layers):
super().**init**()
self.layers = nn.ModuleList([Block(hidden) for _ in range(num_layers)])

def forward(self, x):
    for layer in self.layers:
        # Checkpoint this layer.
        # During backward, PyTorch will recompute layer(x).
        x = checkpoint(layer, x, use_reentrant=False)
    return x

\end{lstlisting}

在工程实践中，checkpointing 的收益取决于模型结构和硬件瓶颈。如果训练已经是 compute-bound，额外重计算会明显降低吞吐；如果训练受显存限制无法提高 batch size 或 sequence length，checkpointing 可能带来整体收益。成熟训练系统通常会把 checkpointing 与 ZeRO、tensor parallel、pipeline parallel、sequence parallel 等策略组合使用。

\section{优化器}
\label{sec:optimizers}

优化器负责根据梯度更新模型参数。对于 AI Infra 工程师，优化器不仅是数学算法，也是一组额外显存状态和 kernel。Adam 这类优化器会为每个参数保存额外状态，因此在大模型训练中，优化器状态往往比参数本身更占显存。

\subsection{SGD}
\label{subsec:sgd}

最基本的优化器是随机梯度下降。设参数为 $\vect{\theta}$，loss 为 $\mathcal{L}$，学习率为 $\eta$，则更新为：
\[
\vect{\theta}_{t+1}
===================

## \vect{\theta}_{t}

\eta
\nabla_{\vect{\theta}}\mathcal{L}(\vect{\theta}_{t}).
\]

SGD 的优点是简单、状态少。每个参数只需要保存自身和梯度，不需要额外的一阶矩、二阶矩。其缺点是对学习率敏感，在不同曲率方向上收敛速度差异很大。

可以从几何上理解。若 loss landscape 在某个方向很陡，在另一个方向很平，SGD 会在陡峭方向来回震荡，同时在平坦方向缓慢前进。Momentum 和 Adam 都可以看作对这一问题的改进。

\begin{lstlisting}[language=Python,caption={SGD 参数更新的最小实现},label={lst:sgd-minimal}]
import torch

@torch.no_grad()
def sgd_step(params, lr):
for p in params:
if p.grad is None:
continue
p -= lr * p.grad
p.grad = None
\end{lstlisting}

在真实训练中，SGD 仍然常用于某些视觉模型或小规模任务，但大语言模型预训练通常更偏好 AdamW，因为它对不同参数尺度和稀疏梯度更稳健。

\subsection{Momentum}
\label{subsec:momentum}

Momentum 为梯度引入速度项。更新公式为：
\[
\vect{v}_{t+1}
==============

\mu\vect{v}*t+
\nabla*{\vect{\theta}}\mathcal{L}(\vect{\theta}*t),
\]
\[
\vect{\theta}*{t+1}
===================

\vect{\theta}*t-\eta\vect{v}*{t+1}.
\]
其中 $\mu$ 是动量系数，通常取 $0.9$ 左右。

Momentum 的直觉是：如果多个 step 的梯度方向大体一致，就加速前进；如果梯度方向来回震荡，就相互抵消。它像物理系统中的惯性，让优化轨迹更平滑。

从系统角度看，Momentum 为每个参数引入一个额外状态 $\vect{v}$。如果参数量为 $N$，并以 FP32 保存 momentum，则额外显存为：
\[
4N\ \mathrm{bytes}.
\]
当 $N$ 达到数十亿甚至数百亿时，这项显存非常可观。

\subsection{Adam 与 AdamW}
\label{subsec:adam-adamw}

Adam 是大模型训练中最常见的优化器之一。它同时维护梯度的一阶矩和二阶矩估计：
\[
\vect{m}*t=\beta_1\vect{m}*{t-1}+(1-\beta_1)\vect{g}_t,
\]
\[
\vect{v}*t=\beta_2\vect{v}*{t-1}+(1-\beta_2)\vect{g}_t^2,
\]
其中：
\[
\vect{g}*t=\nabla*{\vect{\theta}}\mathcal{L}(\vect{\theta}_t).
\]
为了修正初始时刻的偏差，Adam 使用 bias correction：
\[
\hat{\vect{m}}_t=\frac{\vect{m}_t}{1-\beta_1^t},
\]
\[
\hat{\vect{v}}_t=\frac{\vect{v}*t}{1-\beta_2^t}.
\]
参数更新为：
\[
\vect{\theta}*{t+1}
===================

## \vect{\theta}_t

\eta
\frac{\hat{\vect{m}}_t}{\sqrt{\hat{\vect{v}}_t}+\epsilon}.
\]

AdamW 与 Adam 的关键差异在于 weight decay 的处理。传统 Adam 中，L2 正则项会混入梯度；AdamW 使用 decoupled weight decay：
\[
\vect{\theta}_{t+1}
===================

## \vect{\theta}_t

\eta
\frac{\hat{\vect{m}}_t}{\sqrt{\hat{\vect{v}}_t}+\epsilon}
---------------------------------------------------------

\eta\lambda\vect{\theta}_t.
\]
这种形式通常更适合 Transformer 和 LLM 训练。

\begin{lstlisting}[language=Python,caption={AdamW 更新逻辑的简化实现},label={lst:adamw-simplified}]
import math
import torch

@torch.no_grad()
def adamw_step(params, state, lr, beta1=0.9, beta2=0.95, eps=1e-8, weight_decay=0.1):
for p in params:
if p.grad is None:
continue

    g = p.grad

    if p not in state:
        state[p] = {
            "step": 0,
            "m": torch.zeros_like(p, dtype=torch.float32),
            "v": torch.zeros_like(p, dtype=torch.float32),
        }

    s = state[p]
    s["step"] += 1
    t = s["step"]

    m = s["m"]
    v = s["v"]

    # First and second moment estimates.
    m.mul_(beta1).add_(g, alpha=1.0 - beta1)
    v.mul_(beta2).addcmul_(g, g, value=1.0 - beta2)

    # Bias correction.
    m_hat = m / (1.0 - beta1 ** t)
    v_hat = v / (1.0 - beta2 ** t)

    # Decoupled weight decay.
    p.mul_(1.0 - lr * weight_decay)

    # Adam update.
    p.addcdiv_(m_hat, torch.sqrt(v_hat).add_(eps), value=-lr)

    p.grad = None

\end{lstlisting}

AdamW 的显存代价非常高。若模型参数量为 $N$，使用 BF16 参数和梯度，同时保存 FP32 master weight、FP32 一阶矩和 FP32 二阶矩，则粗略显存为：

\[
M
\approx
2N
+
2N
+
4N
+
4N
+
4N
==

16N\ \mathrm{bytes}.
\]

这与第 1 章的估算一致。对于 $70$B 参数模型，优化器相关状态是 ZeRO、FSDP 和分片 checkpoint 必须存在的重要原因。

\begin{table}[H]
\centering
\small
\caption{常见优化器的状态与工程代价}
\label{tab:optimizer-state-cost}
\begin{tabular}{p{0.18\textwidth}p{0.25\textwidth}p{0.24\textwidth}p{0.24\textwidth}}
\toprule
\textbf{优化器} & \textbf{核心状态} & \textbf{优点} & \textbf{工程代价} \
\midrule
SGD
& 无额外状态或很少状态
& 简单，显存低
& 对学习率敏感，收敛可能慢。 \
\midrule
Momentum
& 一阶速度 $\vect{v}$
& 缓解震荡，加速一致方向
& 每个参数多一个状态张量。 \
\midrule
Adam
& 一阶矩 $\vect{m}$，二阶矩 $\vect{v}$
& 自适应学习率，训练稳定
& 优化器状态显存大。 \
\midrule
AdamW
& 一阶矩、二阶矩、解耦 weight decay
& Transformer/LLM 常用
& 大模型训练中必须配合状态分片。 \
\bottomrule
\end{tabular}
\end{table}

\subsection{Learning Rate Schedule}
\label{subsec:learning-rate-schedule}

学习率不是一个固定超参数，而是训练过程中的时间函数。记第 $t$ 步学习率为 $\eta_t$，则优化器更新依赖：
\[
\vect{\theta}_{t+1}
===================

## \vect{\theta}_{t}

\eta_t
\cdot
\operatorname{UpdateDirection}_t.
\]

常见学习率调度包括：

\begin{itemize}
\item \textbf{constant}：保持固定学习率；
\item \textbf{step decay}：每隔一定步数衰减；
\item \textbf{linear decay}：线性下降；
\item \textbf{cosine decay}：按余弦曲线平滑下降；
\item \textbf{warmup + decay}：先从小学习率升高，再逐步下降。
\end{itemize}

大模型训练通常使用 warmup。原因是训练初期参数尚未稳定，梯度统计也不稳定，如果一开始使用过大学习率，容易产生 loss spike 或直接发散。warmup 让优化器先以较小步长探索，随后进入主训练阶段。

\subsection{Warmup、Cosine Decay 与大模型训练稳定性}
\label{subsec:warmup-cosine-stability}

一种常见 schedule 是 linear warmup + cosine decay。设 warmup 步数为 $T_w$，总训练步数为 $T$，峰值学习率为 $\eta_{\max}$，最小学习率为 $\eta_{\min}$。当 $t<T_w$：
\[
\eta_t=\eta_{\max}\frac{t}{T_w}.
\]
当 $t\ge T_w$：
\[
\eta_t
======

\eta_{\min}
+
\frac{1}{2}(\eta_{\max}-\eta_{\min})
\left[
1+
\cos
\left(
\pi
\frac{t-T_w}{T-T_w}
\right)
\right].
\]

\begin{figure}[H]
\centering
\begin{tikzpicture}[
font=\small,
axis/.style={-{Latex[length=2.4mm]}, thick},
curve/.style={thick},
labelnode/.style={align=center}
]

\node[font=\bfseries] at (0,3.0) {常见 Learning Rate Schedule};

\draw[axis] (-5.8,-1.2) -- (5.8,-1.2) node[right] {step};
\draw[axis] (-5.8,-1.2) -- (-5.8,2.2) node[above] {lr};

% Constant
\draw[curve, blue!70] (-5.3,1.55) -- (5.3,1.55);
\node[blue!70] at (4.2,1.82) {constant};

% Linear decay
\draw[curve, green!60!black] (-5.3,1.8) -- (5.3,-0.55);
\node[green!60!black] at (3.8,-0.2) {linear decay};

% Warmup + cosine hand drawn
\draw[curve, orange!90!black]
  (-5.3,-0.85) -- (-3.6,1.75)
  .. controls (-2.4,1.65) and (-1.2,1.35) .. (0.0,0.85)
  .. controls (1.5,0.25) and (3.4,-0.35) .. (5.3,-0.7);
\node[orange!90!black] at (-1.0,1.85) {warmup + cosine};

\draw[dashed] (-3.6,-1.2) -- (-3.6,1.75);
\node[align=center] at (-3.6,-1.6) {warmup\\end};

\node[align=left, text width=10.6cm] at (0,-2.55) {
  大模型训练中，warmup 有助于避免初期不稳定；cosine decay 让后期学习率平滑降低，
  常用于提升收敛稳定性和最终 loss。
};

\end{tikzpicture}
\caption{Constant、Linear Decay 与 Warmup+Cosine Decay 学习率曲线}
\label{fig:learning-rate-schedules}
\end{figure}

学习率 schedule 不只是算法问题，也是系统问题。大规模训练通常运行数天到数周，如果因为学习率设置不当在中后期发散，损失的不只是实验时间，还有大量 GPU 成本。因此成熟训练平台会记录并监控：

\begin{itemize}
\item 当前 step 与 token 数；
\item 当前学习率；
\item loss 曲线；
\item gradient norm；
\item optimizer state 是否出现 NaN；
\item loss scaler 是否频繁回退；
\item checkpoint 恢复后 schedule 是否正确延续。
\end{itemize}

\section{正则化与归一化}
\label{sec:regularization-normalization}

\subsection{Dropout}
\label{subsec:dropout}

Dropout 是一种正则化方法。训练时，它以概率 $p$ 将某些激活置零：
\[
\tilde{h}_i=
\begin{cases}
0, & \text{with probability }p,\
\frac{h_i}{1-p}, & \text{with probability }1-p.
\end{cases}
\]
其中 $\frac{1}{1-p}$ 是为了保持期望不变：
\[
\mathbb{E}[\tilde{h}_i]=h_i.
\]

Dropout 的直觉是防止模型过度依赖某些特定神经元，从而提升泛化能力。它在传统神经网络和早期 Transformer 中很常见。但在大规模语言模型预训练中，dropout 的使用会更加谨慎，原因包括：

\begin{itemize}
\item 大规模数据本身具有正则化效果；
\item dropout 会引入随机性，影响训练吞吐和确定性复现；
\item dropout mask 需要额外保存或重建，影响 backward；
\item 某些大模型训练 recipe 会减少甚至移除 dropout。
\end{itemize}

从 kernel 角度看，dropout 需要随机数生成、mask 应用和缩放，通常属于 memory-bound elementwise 操作。高性能实现会把 dropout 与 bias、activation、residual 等操作融合。

\subsection{BatchNorm}
\label{subsec:batchnorm}

BatchNorm 对 batch 维统计均值和方差。对于某个 channel，给定 mini-batch 中的值：
\[
x_1,x_2,\ldots,x_m,
\]
均值和方差为：
\[
\mu_B=\frac{1}{m}\sum_{i=1}^{m}x_i,
\]
\[
\sigma_B^2=\frac{1}{m}\sum_{i=1}^{m}(x_i-\mu_B)^2.
\]
归一化输出为：
\[
\hat{x}_i=\frac{x_i-\mu_B}{\sqrt{\sigma_B^2+\epsilon}},
\]
再经过可学习仿射变换：
\[
y_i=\gamma\hat{x}_i+\beta.
\]

BatchNorm 在 CNN 中非常成功，但在 LLM 中不常用。主要原因是语言模型训练中的 batch、sequence 和并行切分更复杂，跨 batch 统计会引入同步和不稳定性。尤其在分布式训练中，如果 BatchNorm 需要跨设备统计均值方差，就会引入额外通信。

\subsection{LayerNorm}
\label{subsec:layernorm}

LayerNorm 不依赖 batch 维统计，而是对每个样本内部的 hidden 维归一化。给定：
\[
\vect{x}\in\mathbb{R}^{H},
\]
均值为：
\[
\mu=\frac{1}{H}\sum_{i=1}^{H}x_i,
\]
方差为：
\[
\sigma^2=\frac{1}{H}\sum_{i=1}^{H}(x_i-\mu)^2.
\]
输出为：
\[
y_i=\gamma_i\frac{x_i-\mu}{\sqrt{\sigma^2+\epsilon}}+\beta_i.
\]

对于张量：
\[
\tensor{X}\in\mathbb{R}^{B\times S\times H},
\]
LayerNorm 通常对最后一维 $H$ 做归一化，即每个 $(b,s)$ 位置独立计算均值和方差。这非常适合 Transformer，因为它不依赖 batch size，也不需要跨样本同步统计。

\subsection{RMSNorm}
\label{subsec:rmsnorm}

RMSNorm 是 LayerNorm 的变体。它不减均值，只使用均方根进行归一化：
\[
\operatorname{RMS}(\vect{x})
============================

\sqrt{
\frac{1}{H}\sum_{i=1}^{H}x_i^2+\epsilon
}.
\]
输出为：
\[
y_i=\gamma_i\frac{x_i}{\operatorname{RMS}(\vect{x})}.
\]

与 LayerNorm 相比，RMSNorm 少计算均值和减均值操作，形式更简单。许多现代 Decoder-only LLM 使用 RMSNorm，因为它在训练稳定性和计算效率之间取得了良好平衡。

\begin{figure}[H]
\centering
\begin{tikzpicture}[
font=\small,
cell/.style={draw, minimum width=0.42cm, minimum height=0.38cm},
box/.style={draw, rounded corners=2pt, fill=blue!4},
bn/.style={draw, rounded corners=2pt, fill=orange!12},
ln/.style={draw, rounded corners=2pt, fill=green!10},
rms/.style={draw, rounded corners=2pt, fill=purple!10}
]

\node[font=\bfseries] at (0,3.0) {BatchNorm、LayerNorm 与 RMSNorm 的统计维度};

\node at (-4.3,2.2) {$\tensor{X}\in\mathbb{R}^{B\times S\times H}$};

\foreach \r in {0,...,3} {
  \foreach \c in {0,...,7} {
    \node[cell, fill=blue!5] at ({-5.8+0.42*\c},{1.55-0.38*\r}) {};
  }
}
\node at (-6.35,0.95) {$B/S$};
\node at (-4.25,-0.25) {$H$};

\node[bn, minimum width=3.3cm, minimum height=0.6cm, align=center] at (0.2,1.65)
  {BatchNorm：跨 batch 统计};
\draw[decorate,decoration={brace,amplitude=4pt},thick]
  (-1.6,1.25) -- (-1.6,0.25)
  node[midway,left=6pt,font=\scriptsize] {batch};

\node[ln, minimum width=3.3cm, minimum height=0.6cm, align=center] at (0.2,0.55)
  {LayerNorm：对 hidden 维统计 $\mu,\sigma^2$};
\draw[decorate,decoration={brace,amplitude=4pt},thick]
  (-1.35,0.15) -- (1.75,0.15)
  node[midway,below=6pt,font=\scriptsize] {hidden};

\node[rms, minimum width=3.8cm, minimum height=0.6cm, align=center] at (0.2,-0.65)
  {RMSNorm：对 hidden 维统计 RMS，不减均值};

\node[align=left, text width=5.0cm] at (4.0,0.55) {
  LLM 中更常用 LayerNorm/RMSNorm，\\
  因为它们不依赖 batch 统计，\\
  更适合变长序列、分布式训练和小 micro-batch。
};

\end{tikzpicture}
\caption{BatchNorm、LayerNorm 与 RMSNorm 的统计维度差异}
\label{fig:normalization-axes}
\end{figure}

\subsection{为什么 LLM 更偏好 LayerNorm / RMSNorm}
\label{subsec:why-llm-prefers-layernorm-rmsnorm}

LLM 更偏好 LayerNorm 或 RMSNorm，主要有以下原因：

\begin{itemize}
\item \textbf{不依赖 batch size}：大模型训练常因显存限制使用很小的 micro-batch，BatchNorm 的统计会不稳定。
\item \textbf{不需要跨设备同步统计}：在数据并行、张量并行和流水线并行中，跨设备 BatchNorm 会增加通信复杂度。
\item \textbf{适合序列建模}：每个 token 的 hidden state 可以独立归一化，不依赖其他样本。
\item \textbf{实现简单且易融合}：LayerNorm/RMSNorm 常与 residual add、dropout 等操作融合。
\item \textbf{训练稳定性好}：Pre-Norm Transformer 中，归一化放在 attention 和 FFN 之前，有助于深层网络训练。
\end{itemize}

LayerNorm 和 RMSNorm 的 kernel 通常包括 reduction 和 elementwise 操作。对于每个 token 的 hidden 向量，需要计算均值、方差或 RMS，再进行归一化和缩放。它们通常比 GEMM 更 memory-bound，因此 fused norm kernel 和高效 reduction 对大模型训练吞吐非常重要。

\begin{lstlisting}[language=Python,caption={RMSNorm 的 PyTorch 实现},label={lst:rmsnorm-pytorch}]
import torch
import torch.nn as nn

class RMSNorm(nn.Module):
def **init**(self, hidden_size, eps=1e-6):
super().**init**()
self.weight = nn.Parameter(torch.ones(hidden_size))
self.eps = eps

def forward(self, x):
    # x: [batch, seq, hidden]
    # Compute RMS along the hidden dimension.
    variance = x.pow(2).mean(dim=-1, keepdim=True)
    x_norm = x * torch.rsqrt(variance + self.eps)
    return self.weight * x_norm

\end{lstlisting}

这段代码表达清晰，但在大模型生产训练中，通常会使用 fused RMSNorm kernel。原因是 \texttt{pow}、\texttt{mean}、\texttt{rsqrt}、乘法等如果分别对应多个 kernel，会产生多次 HBM 读写。融合后可以在一个 kernel 中完成 reduction 和归一化，减少显存带宽压力。

\section{混合精度训练}
\label{sec:mixed-precision-training}

\subsection{FP32、FP16、BF16、FP8}
\label{subsec:fp32-fp16-bf16-fp8}

深度学习训练最早大量使用 FP32。FP32 精度高、动态范围较大，但显存和带宽开销也大。为了提升训练吞吐并降低显存占用，现代 GPU 训练广泛使用低精度格式。

常见格式包括：

\begin{itemize}
\item \textbf{FP32}：传统单精度浮点，稳定但成本高。
\item \textbf{FP16}：半精度浮点，显存减半，Tensor Core 支持好，但指数位少，容易 overflow/underflow。
\item \textbf{BF16}：Brain Floating Point，指数位与 FP32 相同，动态范围大，尾数位较少，训练稳定性通常优于 FP16。
\item \textbf{TF32}：NVIDIA Tensor Core 上用于加速 FP32 GEMM 的格式，对用户通常表现为 FP32 输入、低精度乘法、高精度累加。
\item \textbf{FP8}：更低精度格式，用于进一步提升吞吐和降低显存，但需要更复杂的 scaling 和稳定性控制。
\end{itemize}

浮点数通常可抽象为：
\[
(-1)^{s}\times 2^{e-bias}\times (1.f),
\]
其中 $s$ 是符号位，$e$ 是指数，$f$ 是尾数。指数位决定动态范围，尾数位决定有效精度。FP16 的指数位较少，所以训练中更容易出现数值下溢；BF16 保留与 FP32 相同数量级的指数范围，因此更适合大模型训练。

\begin{table}[H]
\centering
\small
\caption{常见数值格式的直观对比}
\label{tab:numeric-format-comparison}
\begin{tabular}{p{0.16\textwidth}p{0.18\textwidth}p{0.22\textwidth}p{0.32\textwidth}}
\toprule
\textbf{格式} & \textbf{位宽} & \textbf{特点} & \textbf{大模型训练中的意义} \
\midrule
FP32
& 32 bit
& 精度高、范围大
& 稳定但显存和带宽成本高，常用于 master weight 或累加。 \
\midrule
FP16
& 16 bit
& 尾数较多，指数较少
& Tensor Core 高效，但需要 loss scaling 防止梯度下溢。 \
\midrule
BF16
& 16 bit
& 指数范围接近 FP32
& 大模型训练常用，稳定性通常优于 FP16。 \
\midrule
FP8
& 8 bit
& 显存和带宽更低
& 需要 scaling recipe，适合特定硬件和成熟训练栈。 \
\bottomrule
\end{tabular}
\end{table}

下面预留一张位布局图片。实际排版时可以放入更精确的 FP32、FP16、BF16、FP8 bit layout 对比图。

\begin{figure}[H]
\centering
\IfFileExists{figures/part1/numeric_formats_bit_layout.png}{
\includegraphics[width=0.88\textwidth]{part1/numeric_formats_bit_layout.png}
}{
\fbox{
\parbox{0.84\textwidth}{
\centering
\vspace{1.1cm}
\textbf{预留图片位：FP32、FP16、BF16、FP8 数值格式位布局}[2mm]
此处建议放入一张横向对比图，展示 sign、exponent、mantissa 在不同浮点格式中的位数分布。
重点突出 BF16 与 FP32 拥有相近指数范围，而 FP16 指数位更少。[2mm]
\textbf{图片生成 Prompt 建议：}\
``Clean educational diagram comparing floating point bit layouts for FP32, FP16, BF16 and FP8, showing sign bit, exponent bits and mantissa bits as colored blocks, technical textbook style, high readability, white background.''
\vspace{1.1cm}
}
}
}
\caption{不同浮点格式 bit layout 对比图预留位}
\label{fig:numeric-format-bit-layout-placeholder}
\end{figure}

\subsection{Loss Scaling}
\label{subsec:loss-scaling}

FP16 的一个主要问题是梯度下溢。训练中很多梯度值非常小，如果直接用 FP16 表示，可能被舍入为 $0$。Loss scaling 的思想是：在 backward 前把 loss 放大一个系数 $S$：
\[
\tilde{\mathcal{L}}=S\mathcal{L}.
\]
于是梯度也被放大：
\[
\nabla_{\theta}\tilde{\mathcal{L}}
==================================

S\nabla_{\theta}\mathcal{L}.
\]
在 optimizer step 前，再把梯度除以 $S$：
\[
\nabla_{\theta}\mathcal{L}
==========================

\frac{1}{S}
\nabla_{\theta}\tilde{\mathcal{L}}.
\]

这样做不会改变数学上的更新方向，却可以让小梯度落入 FP16 可表示范围。若放大系数过大，可能导致 overflow。因此混合精度训练通常使用 dynamic loss scaling：如果发现 inf 或 NaN，就降低 scale；如果一段时间内稳定，就尝试提高 scale。

\begin{figure}[H]
\centering
\begin{tikzpicture}[
font=\small,
step/.style={draw, rounded corners=3pt, minimum width=2.1cm, minimum height=0.8cm, align=center, fill=blue!6},
check/.style={draw, diamond, aspect=2.2, align=center, fill=orange!12},
arrow/.style={-{Latex[length=2.3mm]}, thick}
]

\node[step] (loss) at (-5.2,0) {Compute\\loss};
\node[step] (scale) at (-2.7,0) {Scale loss\\$\tilde{L}=S L$};
\node[step] (backward) at (-0.2,0) {Backward\\scaled grad};
\node[step] (unscale) at (2.3,0) {Unscale\\grad / $S$};
\node[check] (check) at (4.9,0) {inf/NaN?};

\draw[arrow] (loss) -- (scale);
\draw[arrow] (scale) -- (backward);
\draw[arrow] (backward) -- (unscale);
\draw[arrow] (unscale) -- (check);

\node[step, fill=green!8] (stepok) at (4.9,-1.65) {Optimizer\\step};
\node[step, fill=red!8] (skip) at (4.9,1.65) {Skip step\\reduce $S$};

\draw[arrow] (check) -- node[right] {no} (stepok);
\draw[arrow] (check) -- node[right] {yes} (skip);

\node[align=left, text width=10.5cm] at (-0.1,-2.85) {
  Loss scaling 常用于 FP16 训练。BF16 因指数范围更大，通常不需要像 FP16 那样依赖 loss scaling，
  但仍然需要监控 NaN、inf 和梯度范数。
};

\end{tikzpicture}
\caption{混合精度训练中的 Dynamic Loss Scaling 流程}
\label{fig:dynamic-loss-scaling}
\end{figure}

PyTorch AMP 中，loss scaling 通常通过 \texttt{GradScaler} 完成：

\begin{lstlisting}[language=Python,caption={PyTorch AMP 与 GradScaler 训练循环},label={lst:pytorch-amp-gradscaler}]
import torch
import torch.nn as nn

model = nn.Linear(4096, 4096).cuda()
optimizer = torch.optim.AdamW(model.parameters(), lr=1e-4)

scaler = torch.cuda.amp.GradScaler()

for x, target in dataloader:
x = x.cuda()
target = target.cuda()

optimizer.zero_grad(set_to_none=True)

# autocast chooses lower precision kernels when safe.
with torch.cuda.amp.autocast(dtype=torch.float16):
    logits = model(x)
    loss = nn.functional.cross_entropy(logits, target)

# Scale loss before backward.
scaler.scale(loss).backward()

# Unscale gradients, check inf/NaN, and step if safe.
scaler.step(optimizer)

# Update scale factor dynamically.
scaler.update()

\end{lstlisting}

对于 BF16 训练，很多系统会使用 \texttt{autocast(dtype=torch.bfloat16)}，并不使用 \texttt{GradScaler}。不过具体策略仍取决于硬件、框架版本、模型结构和训练 recipe。

\subsection{Tensor Core 与矩阵乘法吞吐}
\label{subsec:tensor-core-matmul-throughput}

Tensor Core 是现代 NVIDIA GPU 上专门加速矩阵乘法累加的硬件单元。它可以高效执行小矩阵 tile 的乘加，例如：
\[
\mat{D}=\mat{A}\mat{B}+\mat{C}.
\]
深度学习中的 Linear、Attention projection、FFN、卷积等大量操作都可以映射到 Tensor Core。

从编程层面看，用户可能只是写：
\[
\mat{Y}=\mat{X}\mat{W}.
\]
但底层库会把大矩阵切分成许多 tile，每个 tile 通过 Tensor Core 执行 MMA 操作，再写回输出矩阵。

\begin{figure}[H]
\centering
\begin{tikzpicture}[
font=\small,
tile/.style={draw, rounded corners=2pt, minimum width=1.25cm, minimum height=1.0cm, align=center},
arrow/.style={-{Latex[length=2.3mm]}, thick}
]

\node[font=\bfseries] at (0,3.0) {Tensor Core 中的 Tile-based Matrix Multiply-Accumulate};

\node[tile, fill=orange!15] (a) at (-3.6,1.1) {$A_{tile}$\\FP16/BF16};
\node[tile, fill=green!15] (b) at (-1.5,1.1) {$B_{tile}$\\FP16/BF16};
\node[tile, fill=blue!10] (c) at (0.6,1.1) {$C_{tile}$\\FP32 acc};

\node[draw, rounded corners=3pt, fill=red!10, minimum width=2.7cm, minimum height=1.05cm, align=center] (tc) at (2.9,1.1)
  {Tensor Core\\MMA};

\node[tile, fill=purple!10] (d) at (5.4,1.1) {$D_{tile}$\\output};

\draw[arrow] (a) -- (tc);
\draw[arrow] (b) -- (tc);
\draw[arrow] (c) -- (tc);
\draw[arrow] (tc) -- (d);

\node[align=center] at (2.9,0.1) {$D=A\times B+C$};

\node[align=left, text width=10.2cm] at (1.0,-1.35) {
  混合精度矩阵乘通常使用低精度输入、高精度累加。
  要真正利用 Tensor Core，shape、dtype、layout 和对齐都需要满足底层库或 kernel 的要求。
};

\end{tikzpicture}
\caption{Tensor Core 的 tile-based MMA 计算直觉}
\label{fig:tensor-core-mma}
\end{figure}

Tensor Core 的存在解释了为什么混合精度训练不只是节省显存，也能显著提升计算吞吐。但实际能否获得加速，还取决于：

\begin{itemize}
\item GEMM 尺寸是否足够大；
\item $M,N,K$ 是否满足硬件友好的对齐；
\item 输入是否 contiguous 或满足高效 stride；
\item 是否使用支持 Tensor Core 的 dtype；
\item cuBLAS/cuBLASLt 是否选中了合适算法；
\item 是否有过多小 kernel 打断时间线。
\end{itemize}

\subsection{精度、速度与稳定性的取舍}
\label{subsec:precision-speed-stability-tradeoff}

混合精度训练的核心不是简单地把所有张量从 FP32 改成 FP16 或 BF16，而是在不同位置选择合适精度。常见策略包括：

\begin{itemize}
\item 参数 forward/backward 使用 BF16 或 FP16；
\item 梯度通信使用 BF16 或 FP16；
\item 优化器状态使用 FP32；
\item 某些 reduction 或 norm 使用 FP32 累加；
\item loss、softmax、logsumexp 等敏感操作使用更稳定实现；
\item FP8 训练中为不同张量维护 scale。
\end{itemize}

\begin{figure}[H]
\centering
\begin{tikzpicture}[
font=\small,
box/.style={draw, rounded corners=3pt, minimum width=2.15cm, minimum height=0.8cm, align=center},
arrow/.style={-{Latex[length=2.3mm]}, thick}
]

\node[box, fill=blue!8] (param) at (-5.0,0.8) {BF16/FP16\\Parameters};
\node[box, fill=green!8] (fwd) at (-2.4,0.8) {Forward\\Tensor Core};
\node[box, fill=green!8] (bwd) at (0.2,0.8) {Backward\\Tensor Core};
\node[box, fill=orange!12] (grad) at (2.8,0.8) {Gradients\\BF16/FP16};
\node[box, fill=red!10] (optim) at (5.4,0.8) {Optimizer\\FP32 states};

\draw[arrow] (param) -- (fwd);
\draw[arrow] (fwd) -- (bwd);
\draw[arrow] (bwd) -- (grad);
\draw[arrow] (grad) -- (optim);
\draw[arrow] (optim.south) .. controls (2.0,-1.2) and (-2.8,-1.2) .. (param.south);

\node[box, fill=yellow!16, minimum width=3.0cm] at (-2.4,-0.55) {Sensitive ops\\FP32 accumulation};
\node[box, fill=purple!10, minimum width=3.1cm] at (2.8,-0.55) {Communication\\compressed dtype};

\node[align=left, text width=10.7cm] at (0,-2.0) {
  混合精度是一套系统策略：哪些张量低精度保存，哪些运算高精度累加，
  哪些通信可以压缩，哪些统计必须保持稳定，都需要根据模型和硬件共同决定。
};

\end{tikzpicture}
\caption{混合精度训练中的精度分配策略}
\label{fig:mixed-precision-strategy}
\end{figure}

混合精度训练中的常见故障包括：

\begin{itemize}
\item loss 突然变成 NaN；
\item gradient norm 爆炸；
\item FP16 loss scaler 频繁下降；
\item 某些 rank 出现 inf 后通过 all-reduce 扩散；
\item checkpoint 恢复后 optimizer state dtype 不匹配；
\item 自定义 CUDA kernel 未正确处理 BF16/FP16 累加精度。
\end{itemize}

因此，训练系统需要在日志和监控中暴露足够信息。仅仅看到 loss 变坏是不够的，还需要知道当前 dtype、loss scale、梯度范数、参数范数、是否发生 overflow、哪个 rank 先出问题、最近一次 checkpoint 是否可恢复。

\section{端到端训练循环：从公式到工程}
\label{sec:end-to-end-training-loop}

为了把本章概念串起来，我们写出一个最小化训练循环。它包含 forward、loss、backward、optimizer step、zero grad 和混合精度上下文。

\begin{lstlisting}[language=Python,caption={现代 PyTorch 训练循环骨架},label={lst:modern-pytorch-training-loop}]
import torch
import torch.nn as nn

model = MyModel().cuda()
optimizer = torch.optim.AdamW(
model.parameters(),
lr=3e-4,
betas=(0.9, 0.95),
weight_decay=0.1,
)

use_bf16 = True

for step, batch in enumerate(train_loader):
input_ids = batch["input_ids"].cuda(non_blocking=True)
labels = batch["labels"].cuda(non_blocking=True)

optimizer.zero_grad(set_to_none=True)

# Mixed precision forward.
amp_dtype = torch.bfloat16 if use_bf16 else torch.float16
with torch.cuda.amp.autocast(dtype=amp_dtype):
    logits = model(input_ids)
    loss = nn.functional.cross_entropy(
        logits.view(-1, logits.size(-1)),
        labels.view(-1),
    )

# Backward builds gradients for all parameters.
loss.backward()

# Optional: gradient clipping for stability.
torch.nn.utils.clip_grad_norm_(model.parameters(), max_norm=1.0)

# AdamW updates parameters and optimizer states.
optimizer.step()

if step % 100 == 0:
    print(f"step={step}, loss={loss.item():.4f}")

\end{lstlisting}

这段代码在 Python 层看起来很短，但从系统角度展开，会出现大量操作：

\begin{itemize}
\item data loader 将数据从存储读入 CPU 内存；
\item token 或 batch 被拷贝到 GPU；
\item forward 触发 embedding、GEMM、attention、norm、MLP 等 kernel；
\item loss 触发 cross entropy kernel；
\item backward 触发反向 GEMM、reduction、activation backward、norm backward；
\item DDP 或其他并行策略触发通信；
\item optimizer step 更新参数、梯度和状态；
\item 日志、checkpoint、监控系统记录训练状态。
\end{itemize}

\begin{figure}[H]
\centering
\begin{tikzpicture}[
font=\small,
step/.style={draw, rounded corners=3pt, minimum width=1.8cm, minimum height=0.72cm, align=center, fill=blue!6},
sys/.style={draw, rounded corners=3pt, minimum width=1.8cm, minimum height=0.72cm, align=center, fill=orange!12},
arrow/.style={-{Latex[length=2.3mm]}, thick}
]

\node[font=\bfseries] at (0,3.0) {一次训练 step 的系统展开};

\node[sys] (data) at (-5.4,1.4) {Data\\Load};
\node[step] (fwd) at (-3.2,1.4) {Forward};
\node[step] (loss) at (-1.0,1.4) {Loss};
\node[step] (bwd) at (1.2,1.4) {Backward};
\node[sys] (comm) at (3.4,1.4) {Gradient\\Comm};
\node[step] (opt) at (5.6,1.4) {Optimizer\\Step};

\draw[arrow] (data) -- (fwd);
\draw[arrow] (fwd) -- (loss);
\draw[arrow] (loss) -- (bwd);
\draw[arrow] (bwd) -- (comm);
\draw[arrow] (comm) -- (opt);

\node[sys, fill=green!8] (ckpt) at (-2.2,0.0) {Checkpoint};
\node[sys, fill=green!8] (log) at (0.0,0.0) {Logging};
\node[sys, fill=green!8] (monitor) at (2.2,0.0) {Monitoring};

\draw[dashed, -{Latex[length=2.1mm]}] (opt.south) .. controls (3.0,0.1) .. (monitor.east);
\draw[dashed, -{Latex[length=2.1mm]}] (loss.south) -- (log.north);
\draw[dashed, -{Latex[length=2.1mm]}] (opt.south) .. controls (0.5,-0.4) .. (ckpt.east);

\node[align=left, text width=10.4cm] at (0,-1.65) {
  一个训练 step 不只是 forward/backward。数据、通信、优化器、checkpoint 和监控共同决定端到端吞吐与可靠性。
};

\end{tikzpicture}
\caption{训练 step 的端到端系统视角}
\label{fig:training-step-system-view}
\end{figure}

从这里开始，我们已经具备理解后续章节的基本语言：Transformer block 是由这些基础组件堆叠而成；分布式训练是把这些参数、梯度、激活和优化器状态切分到多个 GPU；推理优化则是在不做 backward 的条件下，最大化 forward 和 decode 的效率。

\section{本章小结}
\label{sec:chapter03-summary}

本章系统回顾了深度学习基本原理，并始终从 AI Infrastructure 的视角解释其工程代价。

\begin{itemize}
\item \textbf{线性层} 是神经网络和 Transformer 的主要参数与计算来源。Forward 和 backward 都会触发 GEMM。
\item \textbf{激活函数} 引入非线性，但常属于 memory-bound elementwise 操作，生产系统中常与其他操作融合。
\item \textbf{Loss function} 不仅定义训练目标，也可能形成显存和带宽瓶颈，尤其是大词表语言模型中的 cross entropy。
\item \textbf{反向传播} 是链式法则在计算图上的应用，需要保存前向激活，因此训练显存远大于推理。
\item \textbf{Gradient Checkpointing} 通过重计算减少激活显存，是长上下文和大模型训练中的关键技术。
\item \textbf{SGD、Momentum、AdamW} 体现了不同优化思想，同时也带来不同优化器状态显存成本。
\item \textbf{Learning Rate Schedule} 对大模型训练稳定性至关重要，warmup 和 cosine decay 是常见组合。
\item \textbf{LayerNorm 与 RMSNorm} 更适合 LLM，因为它们不依赖 batch 统计，避免跨设备同步，并适配变长序列。
\item \textbf{混合精度训练} 通过低精度计算和高精度累加平衡速度、显存与稳定性，是现代 GPU 训练的默认路径。
\end{itemize}

到此，第一篇的基础部分已经完成：第 1 章建立数学对象和显存计算意识，第 2 章解释 CUDA 和 GPU 执行模型，第 3 章把神经网络训练过程拆解为参数、激活、梯度、优化器状态和数值精度。下一章将进入 Transformer 架构深度剖析，我们会把这些基础组件组合成现代大语言模型最核心的计算结构。


% ==================================================
% 第二篇：演进 Models
% ==================================================
\part{演进：Transformer 与大语言模型架构}

\chapter{Transformer 架构深度剖析}
\label{chap:transformer}

\section{从 RNN 到 Attention}
\label{sec:from-rnn-to-attention}

前一篇我们已经建立了三条基础线索：第一，神经网络本质上是张量运算图；第二，GPU 更擅长规则、密集、可并行的矩阵计算；第三，训练系统的成本主要来自参数、激活、梯度、优化器状态与通信。本章开始，我们进入现代大语言模型最核心的架构：\textbf{Transformer}。

Transformer 之所以重要，不只是因为它在自然语言处理任务上表现出色，更因为它天然适配现代硬件。传统 RNN 按时间步串行递推，而 Transformer 把序列中所有 token 的交互组织为大规模矩阵乘法和归一化操作。这一结构上的差异，直接决定了它能否被 GPU、TPU 和分布式训练系统高效扩展。

\subsection{RNN 的序列瓶颈}
\label{subsec:rnn-sequential-bottleneck}

循环神经网络 RNN 的基本形式可以写为：
\[
\vect{h}*t=f(\vect{h}*{t-1},\vect{x}_t),
\]
其中 $\vect{x}_t$ 是第 $t$ 个 token 的输入表示，$\vect{h}_t$ 是第 $t$ 个位置的隐藏状态。展开后：
\[
\vect{h}_1=f(\vect{h}_0,\vect{x}_1),
\]
\[
\vect{h}_2=f(\vect{h}_1,\vect{x}_2),
\]
\[
\cdots
\]
\[
\vect{h}*T=f(\vect{h}*{T-1},\vect{x}_T).
\]

这个递推结构有一个根本问题：\textbf{第 $t$ 步必须等待第 $t-1$ 步完成}。即使每一步内部可以使用矩阵乘法，时间维度本身仍然串行。这对 GPU 很不友好，因为 GPU 需要大量并行工作来隐藏访存延迟和提高吞吐。

从系统角度看，RNN 的瓶颈可以总结为：

\begin{itemize}
\item \textbf{时间步串行依赖}：无法把整个 sequence 的所有位置完全并行计算。
\item \textbf{长距离依赖难以传播}：梯度需要穿过许多时间步，容易出现梯度消失或梯度爆炸。
\item \textbf{硬件利用率受限}：每个时间步的计算规模相对有限，难以形成超大 GEMM。
\item \textbf{分布式扩展困难}：沿序列维切分会遇到强依赖，流水线化收益有限。
\end{itemize}

\begin{figure}[H]
\centering
\begin{tikzpicture}[
font=\small,
token/.style={draw, rounded corners=2pt, minimum width=0.85cm, minimum height=0.55cm, fill=blue!6, align=center},
state/.style={draw, circle, minimum size=0.72cm, fill=orange!14},
block/.style={draw, rounded corners=3pt, minimum width=1.2cm, minimum height=0.75cm, fill=green!8, align=center},
arrow/.style={-{Latex[length=2.3mm]}, thick},
faintarrow/.style={-{Latex[length=2.0mm]}, thick, gray!60}
]

\node[font=\bfseries] at (0,3.0) {RNN 串行计算 vs Attention 并行计算};

% RNN part
\node[font=\bfseries] at (-3.7,2.35) {RNN：时间步递推};
\foreach \i/\x in {1/-5.6,2/-4.6,3/-3.6,4/-2.6,5/-1.6} {
  \node[token] (rx\i) at (\x,1.45) {$x_\i$};
  \node[state] (rh\i) at (\x,0.35) {$h_\i$};
  \draw[arrow] (rx\i) -- (rh\i);
}
\foreach \i/\j in {1/2,2/3,3/4,4/5} {
  \draw[arrow] (rh\i) -- (rh\j);
}
\node[align=center, text width=4.6cm] at (-3.6,-0.75) {
  后一个 hidden state 依赖前一个，\\
  序列维度难以完全并行。
};

% Attention part
\node[font=\bfseries] at (3.7,2.35) {Self-Attention：全位置并行};
\foreach \i/\x in {1/1.4,2/2.4,3/3.4,4/4.4,5/5.4} {
  \node[token] (ax\i) at (\x,1.45) {$x_\i$};
  \node[block] (ay\i) at (\x,0.0) {$y_\i$};
  \draw[arrow] (ax\i) -- (ay\i);
}

\foreach \i in {1,...,5} {
  \foreach \j in {1,...,5} {
    \draw[faintarrow] (ax\i.south) .. controls ({(1.4+\i*0.2)},0.95) and ({(1.4+\j*0.2)},0.55) .. (ay\j.north);
  }
}

\node[align=center, text width=4.8cm] at (3.5,-0.75) {
  所有 token 之间的相关性可组织成矩阵乘法，\\
  更适合 GPU 并行执行。
};

\end{tikzpicture}
\caption{RNN 的串行依赖与 Self-Attention 的并行结构对比}
\label{fig:rnn-vs-attention-parallelism}
\end{figure}

RNN 的问题并不意味着它没有价值。在语音、时间序列、小模型或流式低延迟场景中，RNN 仍有其位置。但大模型时代的核心工作负载是海量 token、巨大模型和分布式训练，架构必须最大化矩阵乘法比例，并尽量减少不可并行的依赖。Transformer 正是在这个方向上取得突破。

\subsection{Seq2Seq 与 Encoder-Decoder}
\label{subsec:seq2seq-encoder-decoder}

在 Transformer 出现之前，机器翻译等序列到序列任务常使用 Seq2Seq 架构。它由 encoder 和 decoder 组成：

\[
\text{Encoder}: (x_1,x_2,\ldots,x_S)\rightarrow \vect{c},
\]
\[
\text{Decoder}: \vect{c}\rightarrow (y_1,y_2,\ldots,y_T).
\]

早期 Seq2Seq 常把整个输入序列压缩成一个固定长度上下文向量 $\vect{c}$。这会造成信息瓶颈：长句子的全部语义必须塞进一个向量中，decoder 在生成每个输出 token 时都只能依赖同一个压缩表示。

Attention 机制最早的重要动机之一，就是让 decoder 在每个生成位置动态查看 encoder 的不同位置，而不是依赖单个固定向量。也就是说，decoder 不再只看一个 $\vect{c}$，而是根据当前状态对 encoder 所有隐藏状态加权求和。

若 encoder 输出为：
\[
\mat{H}=
\begin{bmatrix}
\vect{h}_1^{\mathsf{T}}\
\vect{h}_2^{\mathsf{T}}\
\cdots\
\vect{h}_S^{\mathsf{T}}
\end{bmatrix}
\in\mathbb{R}^{S\times d},
\]
decoder 某一步的查询向量为 $\vect{q}$，则可以计算权重：
\[
\alpha_i=
\frac{\exp(\operatorname{score}(\vect{q},\vect{h}*i))}
{\sum*{j=1}^{S}\exp(\operatorname{score}(\vect{q},\vect{h}*j))},
\]
再得到上下文向量：
\[
\vect{c}=\sum*{i=1}^{S}\alpha_i\vect{h}_i.
\]

这个思想非常朴素：\textbf{不同输出位置应该关注输入序列中不同位置的信息。}后来的 Self-Attention 则进一步把这个机制用于同一个序列内部，使每个 token 都可以关注其他 token。

\subsection{Attention 的核心思想}
\label{subsec:attention-core-idea}

Attention 可以用一句话概括：\textbf{根据相似度对一组 value 加权求和。}

给定一个 query $\vect{q}$、一组 key ${\vect{k}_i}$ 和对应 value ${\vect{v}_i}$，attention 的输出为：
\[
\operatorname{Attention}(\vect{q},\mat{K},\mat{V})
==================================================

\sum_i \alpha_i\vect{v}_i,
\]
其中权重：
\[
\alpha_i=
\softmax_i(\operatorname{score}(\vect{q},\vect{k}_i)).
\]

在深度学习中，最常见的 score 函数是点积：
\[
\operatorname{score}(\vect{q},\vect{k}_i)=\vect{q}^{\mathsf{T}}\vect{k}_i.
\]

点积越大，说明 query 与 key 越相似，对应 value 权重越高。这个机制可以被理解为一种可微的内容寻址：query 是查询条件，key 是地址，value 是内容。

从系统角度看，attention 的高效之处在于它可以批量化成矩阵运算。对于多个 query 组成的矩阵：
\[
\mat{Q}\in\mathbb{R}^{S_q\times d_k},
\]
key 矩阵：
\[
\mat{K}\in\mathbb{R}^{S_k\times d_k},
\]
value 矩阵：
\[
\mat{V}\in\mathbb{R}^{S_k\times d_v},
\]
attention 可以写为：
\[
\Attention(\mat{Q},\mat{K},\mat{V})
===================================

\softmax
\left(
\frac{\mat{Q}\mat{K}^{\mathsf{T}}}{\sqrt{d_k}}
\right)
\mat{V}.
\]

这条公式是 Transformer 的核心。它由两个主要矩阵乘法构成：

\begin{itemize}
\item $\mat{Q}\mat{K}^{\mathsf{T}}$：计算所有 query 与所有 key 的相似度；
\item $\mat{P}\mat{V}$：用 attention probability 对 value 加权求和。
\end{itemize}

其中 $\mat{P}$ 表示 softmax 后的 attention probability：
\[
\mat{P}=
\softmax
\left(
\frac{\mat{Q}\mat{K}^{\mathsf{T}}}{\sqrt{d_k}}
\right).
\]

\subsection{Self-Attention 的并行优势}
\label{subsec:self-attention-parallel-advantage}

Self-Attention 指 query、key、value 都来自同一个输入序列。给定输入：
\[
\mat{X}\in\mathbb{R}^{S\times H},
\]
通过三个线性投影得到：
\[
\mat{Q}=\mat{X}\mat{W}_Q,\quad
\mat{K}=\mat{X}\mat{W}_K,\quad
\mat{V}=\mat{X}\mat{W}_V.
\]
随后：
\[
\mat{O}
=======

\softmax
\left(
\frac{\mat{Q}\mat{K}^{\mathsf{T}}}{\sqrt{d_k}}
\right)\mat{V}.
\]

Self-Attention 的并行优势在于：所有 token 的 $\mat{Q},\mat{K},\mat{V}$ 可以一次性通过 GEMM 得到；所有 token 两两之间的相似度也可以通过一次 batched GEMM 或专门 attention kernel 计算。这与 RNN 的时间步递推形成鲜明对比。

当然，Self-Attention 也带来代价。标准 attention 的 score matrix shape 为：
\[
S\times S.
\]
当序列长度 $S$ 增大时，计算量和显存都按 $O(S^2)$ 增长。这也是长上下文模型、FlashAttention、PagedAttention 和各种稀疏 attention 研究的重要动机。

\begin{table}[H]
\centering
\small
\caption{RNN 与 Self-Attention 的系统视角对比}
\label{tab:rnn-self-attention-system-view}
\begin{tabular}{p{0.22\textwidth}p{0.34\textwidth}p{0.34\textwidth}}
\toprule
\textbf{维度} & \textbf{RNN} & \textbf{Self-Attention} \
\midrule
序列依赖
& 时间步递推，强串行依赖
& 所有位置可并行计算相关性 \
\midrule
长距离建模
& 依赖 hidden state 逐步传递
& 任意两个位置可直接交互 \
\midrule
GPU 友好性
& 时间维难以完全并行，GEMM 尺寸较受限
& 大量 GEMM 和 batched matrix multiply，适合 GPU \
\midrule
复杂度
& 通常随序列长度线性递推
& 标准 attention 随序列长度二次增长 \
\midrule
大模型扩展
& 扩展性受串行依赖制约
& 更适合大规模训练，但需要解决 $S^2$ attention 代价 \
\bottomrule
\end{tabular}
\end{table}

\section{Scaled Dot-Product Attention}
\label{sec:scaled-dot-product-attention}

Scaled Dot-Product Attention 是 Transformer 中最核心的算子。它的形式非常简洁：
\[
\Attention(\mat{Q},\mat{K},\mat{V})
===================================

\softmax
\left(
\frac{\mat{Q}\mat{K}^{\mathsf{T}}}{\sqrt{d_k}}
\right)
\mat{V}.
\]
但这条公式背后包含许多工程细节：矩阵 shape、mask、softmax 数值稳定性、中间张量显存、kernel fusion、训练与推理差异。我们逐一展开。

\subsection{Query、Key、Value}
\label{subsec:qkv}

在 attention 中，Query、Key、Value 分别承担不同角色：

\begin{itemize}
\item \textbf{Query}：当前位置想要寻找什么信息；
\item \textbf{Key}：每个位置提供什么索引或匹配特征；
\item \textbf{Value}：真正被加权汇聚的信息内容。
\end{itemize}

对于输入：
\[
\mat{X}\in\mathbb{R}^{B\times S\times H},
\]
通常通过线性层得到：
\[
\mat{Q}=\mat{X}\mat{W}_Q,
\quad
\mat{K}=\mat{X}\mat{W}_K,
\quad
\mat{V}=\mat{X}\mat{W}_V.
\]
其中：
\[
\mat{W}_Q,\mat{W}_K,\mat{W}*V\in\mathbb{R}^{H\times H}
\]
是最常见的写法。多头注意力中会把 $H$ 拆成多个 head，每个 head 的维度为：
\[
d_h=\frac{H}{n*{\mathrm{heads}}}.
\]

在工程实现中，QKV projection 常常被合并成一个大线性层：
\[
[\mat{Q},\mat{K},\mat{V}]
=========================

\mat{X}\mat{W}*{QKV},
\]
其中：
\[
\mat{W}*{QKV}\in\mathbb{R}^{H\times 3H}.
\]
这样可以把三次 GEMM 合并成一次更大的 GEMM，提高 GPU 利用率，并减少 kernel launch 和读取输入 $\mat{X}$ 的次数。

\begin{figure}[H]
\centering
\begin{tikzpicture}[
font=\small,
tensor/.style={draw, rounded corners=3pt, minimum width=1.5cm, minimum height=0.75cm, align=center, fill=blue!6},
proj/.style={draw, rounded corners=3pt, minimum width=1.25cm, minimum height=0.65cm, align=center, fill=orange!12},
mat/.style={draw, rounded corners=3pt, minimum width=1.25cm, minimum height=0.65cm, align=center, fill=green!10},
arrow/.style={-{Latex[length=2.3mm]}, thick}
]

\node[tensor] (x) at (-5.2,0) {$\mat{X}$\\$B,S,H$};

\node[proj] (wq) at (-3.3,1.35) {$W_Q$};
\node[proj] (wk) at (-3.3,0) {$W_K$};
\node[proj] (wv) at (-3.3,-1.35) {$W_V$};

\node[mat] (q) at (-1.5,1.35) {$Q$};
\node[mat] (k) at (-1.5,0) {$K$};
\node[mat] (v) at (-1.5,-1.35) {$V$};

\draw[arrow] (x) -- (wq);
\draw[arrow] (x) -- (wk);
\draw[arrow] (x) -- (wv);
\draw[arrow] (wq) -- (q);
\draw[arrow] (wk) -- (k);
\draw[arrow] (wv) -- (v);

\node[mat, minimum width=2.3cm, minimum height=0.9cm, fill=red!10] (score) at (1.4,0.7)
  {$QK^{\mathsf{T}}/\sqrt{d_k}$};
\node[mat, minimum width=1.65cm, fill=yellow!15] (softmax) at (3.7,0.7)
  {Softmax};
\node[mat, minimum width=1.65cm, fill=purple!10] (out) at (5.8,0)
  {$O=PV$};

\draw[arrow] (q) -- (score);
\draw[arrow] (k) -- (score);
\draw[arrow] (score) -- (softmax);
\draw[arrow] (softmax) -- (out);
\draw[arrow] (v.east) .. controls (2.2,-1.4) and (4.4,-1.25) .. (out.south);

\node[align=left, text width=9.8cm] at (0,-2.55) {
  Q、K、V 可以来自三个独立线性层，也常被合并为一个 QKV projection。
  后者能形成更大的 GEMM，是训练和推理实现中的常见优化。
};

\end{tikzpicture}
\caption{Q/K/V 投影与 Scaled Dot-Product Attention 数据流}
\label{fig:qkv-attention-dataflow}
\end{figure}

\subsection{$QK^{\mathsf{T}}$ 的语义}
\label{subsec:qk-dot-product-semantics}

矩阵：
\[
\mat{S}=\mat{Q}\mat{K}^{\mathsf{T}}
\]
中的元素为：
\[
S_{ij}=\vect{q}_i^{\mathsf{T}}\vect{k}*j.
\]
它表示第 $i$ 个 query 与第 $j$ 个 key 的匹配程度。若 $S*{ij}$ 大，则第 $i$ 个 token 更关注第 $j$ 个 token。

对于一个长度为 $S$ 的序列，$\mat{S}$ 是一个 $S\times S$ 矩阵。第 $i$ 行表示 token $i$ 对所有 token 的注意力打分。经过 softmax 后：
\[
P_{ij}
======

\frac{\exp(S_{ij})}{\sum_{j'=1}^{S}\exp(S_{ij'})}.
\]
每一行的概率和为 $1$：
\[
\sum_{j=1}^{S}P_{ij}=1.
\]

从语言理解角度看，attention matrix 可以捕捉依赖关系。例如在句子“模型把输入映射到输出，因为它学习了参数”中，“它”可能需要关注“模型”。Self-Attention 允许这种长距离连接在单层中直接发生。

从系统角度看，$QK^{\mathsf{T}}$ 是一个重要瓶颈。它的计算量约为：
\[
2B n_h S^2 d_h,
\]
其中 $n_h$ 是 head 数，$d_h$ 是每个 head 的维度。显式 attention score 的元素数为：
\[
B n_h S^2.
\]
当 $S$ 很大时，$S^2$ 项会迅速成为显存和带宽压力来源。

\subsection{scale 因子 $\sqrt{d_k}$}
\label{subsec:attention-scale-factor}

为什么 attention 要除以 $\sqrt{d_k}$？考虑 query 和 key 的元素独立同分布，均值为 $0$，方差为 $1$。点积：
\[
\vect{q}^{\mathsf{T}}\vect{k}
=============================

\sum_{i=1}^{d_k}q_i k_i.
\]
若 $q_i,k_i$ 独立，则 $q_i k_i$ 的方差约为 $1$，因此和的方差约为：
\[
\operatorname{Var}(\vect{q}^{\mathsf{T}}\vect{k})\approx d_k.
\]
当 $d_k$ 较大时，点积值的尺度会变大，使 softmax 输入过大。softmax 在输入尺度很大时会变得极其尖锐，梯度可能变小，不利于训练。

除以 $\sqrt{d_k}$ 后：
\[
\operatorname{Var}
\left(
\frac{\vect{q}^{\mathsf{T}}\vect{k}}{\sqrt{d_k}}
\right)
\approx 1.
\]
这样可以让 softmax 的输入尺度更稳定。

这也是深度学习中常见的尺度控制思想：初始化、归一化、residual scaling、attention scaling 都是在控制信号和梯度的数值范围。对于大模型训练，数值尺度不仅影响收敛，也影响混合精度稳定性。

\subsection{mask 机制}
\label{subsec:attention-mask}

Attention 中常见的 mask 有两类：

\begin{itemize}
\item \textbf{padding mask}：避免模型关注 padding token；
\item \textbf{causal mask}：在自回归语言模型中，避免当前位置看到未来 token。
\end{itemize}

Decoder-only LLM 使用 causal mask。对于序列位置 $i$，只能关注 $j\le i$ 的位置。也就是说：
\[
S_{ij}=
\begin{cases}
S_{ij}, & j\le i,\
-\infty, & j>i.
\end{cases}
\]
softmax 后，未来位置概率为 $0$。

\begin{figure}[H]
\centering
\begin{tikzpicture}[
font=\small,
cell/.style={draw, minimum width=0.42cm, minimum height=0.42cm},
allowed/.style={fill=green!16},
blocked/.style={fill=red!12},
axislabel/.style={font=\scriptsize}
]

\node[font=\bfseries] at (0,3.0) {Causal Mask：当前位置不能看到未来 token};

\foreach \i in {0,...,7} {
  \node[axislabel] at ({-1.7+0.45*\i},2.25) {$k_{\i}$};
  \node[axislabel] at (-2.25,{1.8-0.45*\i}) {$q_{\i}$};
}

\foreach \r in {0,...,7} {
  \foreach \c in {0,...,7} {
    \ifnum\c>\r
      \node[cell, blocked] at ({-1.7+0.45*\c},{1.8-0.45*\r}) {};
    \else
      \node[cell, allowed] at ({-1.7+0.45*\c},{1.8-0.45*\r}) {};
    \fi
  }
}

\node[draw, rounded corners=2pt, fill=green!16, minimum width=0.7cm, minimum height=0.4cm] at (2.6,1.0) {};
\node[align=left] at (4.0,1.0) {可关注历史与当前位置};

\node[draw, rounded corners=2pt, fill=red!12, minimum width=0.7cm, minimum height=0.4cm] at (2.6,0.25) {};
\node[align=left] at (4.0,0.25) {未来位置被 mask};

\node[align=left, text width=9.8cm] at (1.0,-2.15) {
  Causal mask 保证自回归语言模型训练与推理一致。训练时虽然所有位置并行计算，
  但每个位置的预测只能依赖它之前的 token。
};

\end{tikzpicture}
\caption{Decoder-only 模型中的 causal mask}
\label{fig:causal-mask}
\end{figure}

工程实现中，mask 不一定真的构造一个 $S\times S$ 的完整矩阵。高性能 attention kernel 可以根据行列索引动态判断是否越过 causal 边界，从而避免额外显存开销。FlashAttention、xFormers、Triton attention kernel 等实现都会特别关注 mask 的高效处理。

\subsection{attention score 到 weighted sum}
\label{subsec:attention-score-to-weighted-sum}

计算完 attention probability：
\[
\mat{P}=\softmax(\mat{S}),
\]
输出为：
\[
\mat{O}=\mat{P}\mat{V}.
\]
其中：
\[
\vect{o}*i=\sum*{j=1}^{S}P_{ij}\vect{v}_j.
\]
也就是说，每个输出 token 是所有 value 的加权平均。

下面给出 PyTorch 风格的 Scaled Dot-Product Attention 实现。这个实现用于教学，强调公式对应关系；生产系统应使用高性能 fused attention kernel。

\begin{lstlisting}[language=Python,caption={教学版 Scaled Dot-Product Attention},label={lst:scaled-dot-product-attention}]
import math
import torch

def scaled_dot_product_attention(q, k, v, causal=True):
"""
q, k, v: [batch, heads, seq, head_dim]
return: [batch, heads, seq, head_dim]
"""
B, H, S, D = q.shape

# QK^T: [B, H, S, D] @ [B, H, D, S] -> [B, H, S, S]
scores = torch.matmul(q, k.transpose(-2, -1)) / math.sqrt(D)

if causal:
    # Create an upper-triangular mask.
    # True means positions to mask out.
    mask = torch.triu(
        torch.ones(S, S, device=q.device, dtype=torch.bool),
        diagonal=1,
    )
    scores = scores.masked_fill(mask, float("-inf"))

# Softmax along key dimension.
probs = torch.softmax(scores, dim=-1)

# Weighted sum: [B, H, S, S] @ [B, H, S, D] -> [B, H, S, D]
out = torch.matmul(probs, v)
return out

\end{lstlisting}

这段代码清楚但低效，原因包括：

\begin{itemize}
\item 显式物化了 $\mat{S}\in\mathbb{R}^{B\times H\times S\times S}$；
\item 显式物化了 mask；
\item softmax、mask、matmul 被拆成多个 kernel；
\item 中间 attention probability 需要写回 HBM；
\item 当 $S$ 很长时，显存占用和 HBM IO 都很高。
\end{itemize}

后续 FlashAttention 章节会展示如何通过 tiling、online softmax 和 kernel fusion 避免物化完整 attention matrix。

\section{Multi-Head Attention}
\label{sec:multi-head-attention}

单头 attention 可以让每个 token 根据一个相似度空间关注其他 token。但语言中的关系非常复杂：有些 head 可能关注语法依赖，有些 head 可能关注实体指代，有些 head 可能关注位置模式。Multi-Head Attention 通过多个子空间并行建模不同关系。

\subsection{为什么需要多头}
\label{subsec:why-multi-head}

设隐藏维度为 $H$，head 数为 $n_h$，每个 head 维度为：
\[
d_h=\frac{H}{n_h}.
\]
Multi-Head Attention 不是用一个 $H$ 维 attention，而是把隐藏状态投影到多个 head，每个 head 独立计算 attention：
\[
\operatorname{head}_i
=====================

\Attention
\left(
\mat{Q}_i,\mat{K}_i,\mat{V}_i
\right).
\]
然后拼接所有 head：
\[
\mat{O}
=======

\operatorname{Concat}
\left(
\operatorname{head}*1,\ldots,\operatorname{head}*{n_h}
\right)
\mat{W}_O.
\]

多头的好处包括：

\begin{itemize}
\item \textbf{表达多个关系子空间}：不同 head 可以学习不同关注模式；
\item \textbf{保持计算规模可控}：每个 head 的维度较小，点积尺度稳定；
\item \textbf{适配并行计算}：head 维可以与 batch 维合并，形成 batched GEMM；
\item \textbf{便于模型并行}：不同 head 可以在 tensor parallel 中被切分到不同 GPU。
\end{itemize}

\subsection{Head 维度切分}
\label{subsec:head-dimension-split}

输入张量：
\[
\mat{X}\in\mathbb{R}^{B\times S\times H}
\]
经过 QKV projection 后，通常得到：
\[
\mat{Q},\mat{K},\mat{V}\in\mathbb{R}^{B\times S\times H}.
\]
随后 reshape 为：
\[
\mathbb{R}^{B\times S\times n_h\times d_h}.
\]
为了方便 attention kernel 访问，常进一步 permute 为：
\[
\mathbb{R}^{B\times n_h\times S\times d_h}.
\]

这个 reshape/transpose 在数学上很简单，但工程上非常关键。若 transpose 后张量不 contiguous，传给某些 kernel 可能触发额外 copy；若 memory layout 不适合连续访问，attention 的 HBM 带宽利用率会下降。现代 attention kernel 通常会对 QKV layout 有严格要求，例如 packed QKV、separate Q/K/V、BSHD 或 BSHD-like layout 等。

\begin{figure}[H]
\centering
\begin{tikzpicture}[
font=\small,
box/.style={draw, rounded corners=3pt, align=center, fill=blue!6, minimum height=0.75cm},
head/.style={draw, rounded corners=2pt, align=center, fill=orange!12, minimum width=1.0cm, minimum height=0.55cm},
arrow/.style={-{Latex[length=2.3mm]}, thick}
]

\node[box, minimum width=2.4cm] (x) at (-5.3,0) {$X$\\$B,S,H$};
\node[box, minimum width=2.8cm, fill=green!8] (qkv) at (-2.2,0) {QKV Projection\\$B,S,3H$};
\node[box, minimum width=2.8cm, fill=yellow!14] (reshape) at (1.2,0) {Reshape\\$B,S,n_h,d_h$};
\node[box, minimum width=2.8cm, fill=purple!10] (permute) at (4.6,0) {Permute\\$B,n_h,S,d_h$};

\draw[arrow] (x) -- (qkv);
\draw[arrow] (qkv) -- (reshape);
\draw[arrow] (reshape) -- (permute);

\foreach \i/\xpos in {1/-1.1,2/-0.25,3/0.6,4/1.45,5/2.3} {
  \node[head] at (\xpos,-1.35) {head \i};
}

\draw[decorate,decoration={brace,amplitude=5pt},thick]
  (-1.65,-1.75) -- (2.85,-1.75)
  node[midway,below=7pt] {$H=n_h d_h$};

\node[align=left, text width=10.0cm] at (0,-2.75) {
  Head 维度切分在数学上是 reshape，在系统上涉及 stride、layout、contiguous 与 attention kernel 的输入约定。
};

\end{tikzpicture}
\caption{Multi-Head Attention 中 hidden 维到 head 维的切分}
\label{fig:head-dimension-split}
\end{figure}

\subsection{concat 与 output projection}
\label{subsec:concat-output-projection}

每个 head 独立得到输出：
\[
\mat{O}*i\in\mathbb{R}^{B\times S\times d_h}.
\]
将所有 head 沿 hidden 维拼接：
\[
\mat{O}*{\mathrm{concat}}
=========================

\operatorname{Concat}
(\mat{O}*1,\ldots,\mat{O}*{n_h})
\in\mathbb{R}^{B\times S\times H}.
\]
随后通过 output projection：
\[
\mat{Y}=\mat{O}_{\mathrm{concat}}\mat{W}_O.
\]

为什么还需要 output projection？因为多个 head 的结果只是简单拼接，$\mat{W}_O$ 可以重新混合不同 head 的信息，把它们投影回 residual stream 所在的 hidden 空间。

从参数量看，如果 Q、K、V、O 都是 $H\times H$ 矩阵，则 attention block 的主要参数量约为：
\[
4H^2.
\]
这还不包括 LayerNorm 和 bias。FFN 通常参数更多，后文会分析。

\begin{figure}[H]
\centering
\begin{tikzpicture}[
font=\small,
head/.style={draw, rounded corners=3pt, minimum width=1.4cm, minimum height=0.75cm, align=center, fill=orange!12},
op/.style={draw, rounded corners=3pt, minimum width=1.9cm, minimum height=0.8cm, align=center, fill=blue!6},
arrow/.style={-{Latex[length=2.3mm]}, thick}
]

\node[op] (input) at (-5.2,0) {$X$};

\foreach \i/\y in {1/1.6,2/0.55,3/-0.55,4/-1.6} {
  \node[head] (h\i) at (-2.6,\y) {Head \i\\Attention};
  \draw[arrow] (input.east) .. controls (-4.0,\y) .. (h\i.west);
}

\node[op, fill=green!8, minimum width=2.1cm, minimum height=1.0cm] (concat) at (0.3,0) {Concat\\heads};
\node[op, fill=purple!10, minimum width=2.1cm] (wo) at (3.0,0) {Output\\Projection $W_O$};
\node[op, fill=yellow!14] (out) at (5.5,0) {$Y$};

\foreach \i in {1,...,4} {
  \draw[arrow] (h\i.east) -- (concat.west);
}

\draw[arrow] (concat) -- (wo);
\draw[arrow] (wo) -- (out);

\node[align=left, text width=9.8cm] at (0,-2.75) {
  多个 head 独立计算注意力，拼接后通过 $W_O$ 重新混合。实际实现中，所有 head 常在同一个 fused attention kernel 中并行处理。
};

\end{tikzpicture}
\caption{Multi-Head Attention 的 head 并行、concat 与 output projection}
\label{fig:multi-head-attention-structure}
\end{figure}

\subsection{MHA、MQA、GQA 对比}
\label{subsec:mha-mqa-gqa}

标准 Multi-Head Attention 中，每个 query head 都有自己的 key head 和 value head。若 query head 数为 $n_q$，则：
\[
n_k=n_v=n_q.
\]
这称为 MHA。

在推理 decode 阶段，KV Cache 显存占用与 key/value head 数成正比。为了降低 KV Cache 成本，现代 LLM 常使用 MQA 或 GQA。

\textbf{Multi-Query Attention, MQA} 让所有 query heads 共享同一组 key/value：
\[
n_k=n_v=1,\quad n_q>1.
\]
它显著降低 KV Cache 显存和带宽，但可能影响模型表达能力。

\textbf{Grouped-Query Attention, GQA} 是 MHA 与 MQA 的折中。多个 query heads 分成若干组，每组共享一组 key/value：
\[
1<n_k=n_v<n_q.
\]
例如 $n_q=32$，$n_{kv}=8$，则每 $4$ 个 query head 共享一组 K/V。

\begin{figure}[H]
\centering
\begin{tikzpicture}[
font=\small,
q/.style={draw, rounded corners=2pt, fill=blue!10, minimum width=0.48cm, minimum height=0.42cm},
kv/.style={draw, rounded corners=2pt, fill=orange!16, minimum width=0.48cm, minimum height=0.42cm},
title/.style={font=\bfseries},
arrow/.style={-{Latex[length=1.9mm]}, thick}
]

\node[title] at (-4.5,2.5) {MHA};
\node[title] at (0,2.5) {MQA};
\node[title] at (4.5,2.5) {GQA};

% MHA
\foreach \i in {0,...,5} {
  \node[q] (mq\i) at ({-5.8+0.45*\i},1.6) {};
  \node[kv] (mk\i) at ({-5.8+0.45*\i},0.8) {};
  \node[kv] (mv\i) at ({-5.8+0.45*\i},0.1) {};
  \draw[arrow] (mq\i) -- (mk\i);
  \draw[arrow] (mq\i) -- (mv\i);
}
\node[align=center] at (-4.55,-0.7) {每个 Q head\\独享 K/V};

% MQA
\foreach \i in {0,...,5} {
  \node[q] (qq\i) at ({-1.3+0.45*\i},1.6) {};
}
\node[kv, minimum width=2.8cm] (qk) at (0,0.8) {shared K};
\node[kv, minimum width=2.8cm] (qv) at (0,0.1) {shared V};
\foreach \i in {0,...,5} {
  \draw[arrow] (qq\i) -- (qk.north);
  \draw[arrow] (qq\i) -- (qv.north);
}
\node[align=center] at (0,-0.7) {所有 Q head\\共享一组 K/V};

% GQA
\foreach \i in {0,...,7} {
  \node[q] (gq\i) at ({3.0+0.38*\i},1.6) {};
}
\node[kv, minimum width=1.3cm] (gk1) at (3.55,0.8) {K1};
\node[kv, minimum width=1.3cm] (gk2) at (5.1,0.8) {K2};
\node[kv, minimum width=1.3cm] (gv1) at (3.55,0.1) {V1};
\node[kv, minimum width=1.3cm] (gv2) at (5.1,0.1) {V2};
\foreach \i in {0,...,3} {
  \draw[arrow] (gq\i) -- (gk1.north);
  \draw[arrow] (gq\i) -- (gv1.north);
}
\foreach \i in {4,...,7} {
  \draw[arrow] (gq\i) -- (gk2.north);
  \draw[arrow] (gq\i) -- (gv2.north);
}
\node[align=center] at (4.35,-0.7) {每组 Q heads\\共享 K/V};

\node[align=left, text width=10.4cm] at (0,-2.1) {
  MQA/GQA 主要面向推理效率：减少 KV Cache 大小和 decode 阶段的内存带宽压力。
  这是模型结构与 serving infra 共同设计的典型案例。
};

\end{tikzpicture}
\caption{MHA、MQA 与 GQA 的 KV 共享方式对比}
\label{fig:mha-mqa-gqa-comparison}
\end{figure}

KV Cache 的详细机制会在推理篇专门讨论。这里需要提前建立直觉：\textbf{训练阶段的主要瓶颈常是计算与激活，推理 decode 阶段的瓶颈常是 KV Cache 带宽与调度。}GQA 的流行正是因为它能在不显著损害质量的情况下，降低推理系统成本。

\section{Feed-Forward Network}
\label{sec:feed-forward-network}

Transformer block 除了 attention，还有一个逐 token 的前馈网络 Feed-Forward Network，简称 FFN 或 MLP。很多初学者以为 attention 是 Transformer 中参数最多的部分，但在典型 LLM 中，FFN 往往贡献了更大的参数量和计算量。

\subsection{FFN 的参数量来源}
\label{subsec:ffn-parameter-source}

标准 Transformer FFN 可以写为：
\[
\operatorname{FFN}(\vect{x})
============================

\phi(\vect{x}\mat{W}_1+\vect{b}_1)\mat{W}_2+\vect{b}*2.
\]
其中：
\[
\mat{W}*1\in\mathbb{R}^{H\times H*{\mathrm{ffn}}},
\quad
\mat{W}*2\in\mathbb{R}^{H*{\mathrm{ffn}}\times H}.
\]
通常：
\[
H*{\mathrm{ffn}}\approx 4H.
\]
因此 FFN 参数量约为：
\[
H\cdot 4H+4H\cdot H=8H^2.
\]

而标准 attention 的 Q、K、V、O projection 参数量约为：
\[
4H^2.
\]
这说明 FFN 参数量通常约为 attention projection 的两倍。对于 Decoder-only LLM，绝大部分参数分布在每层的 attention 和 FFN 中，而 FFN 往往是更大的参数来源。

\subsection{GELU、SwiGLU 与门控结构}
\label{subsec:gelu-swiglu-gated-ffn}

早期 Transformer 常使用 GELU：
\[
\operatorname{FFN}(\mat{X})
===========================

\operatorname{GELU}(\mat{X}\mat{W}_1)\mat{W}_2.
\]
现代 LLM 中常见 SwiGLU：
\[
\operatorname{SwiGLU}(\mat{X})
==============================

\left[
\operatorname{SiLU}(\mat{X}\mat{W}_g)
\odot
(\mat{X}\mat{W}_u)
\right]\mat{W}_d.
\]
这里有三个 projection：

\[
\mat{W}*g\in\mathbb{R}^{H\times H*{\mathrm{mid}}},
\quad
\mat{W}*u\in\mathbb{R}^{H\times H*{\mathrm{mid}}},
\quad
\mat{W}*d\in\mathbb{R}^{H*{\mathrm{mid}}\times H}.
\]

门控结构的直觉是：一条分支生成内容，另一条分支生成门控权重，二者逐元素相乘。它能提升表达能力，但也改变参数量和计算量。因此使用 SwiGLU 时，$H_{\mathrm{mid}}$ 往往不是简单的 $4H$，而会调整为较小值，以保持整体参数量可控。

\begin{figure}[H]
\centering
\begin{tikzpicture}[
font=\small,
box/.style={draw, rounded corners=3pt, minimum width=1.55cm, minimum height=0.7cm, align=center},
arrow/.style={-{Latex[length=2.3mm]}, thick},
mult/.style={draw, circle, minimum size=0.65cm, fill=yellow!18}
]

\node[box, fill=blue!6] (x) at (-5.2,0) {$X$};

\node[box, fill=orange!12] (wg) at (-3.2,1.0) {$W_g$};
\node[box, fill=orange!12] (wu) at (-3.2,-1.0) {$W_u$};

\node[box, fill=green!10] (silu) at (-1.1,1.0) {SiLU};
\node[box, fill=green!10] (up) at (-1.1,-1.0) {Up\\features};

\node[mult] (mul) at (1.1,0) {$\odot$};
\node[box, fill=purple!10] (wd) at (3.1,0) {$W_d$};
\node[box, fill=blue!6] (out) at (5.0,0) {$Y$};

\draw[arrow] (x) -- (wg);
\draw[arrow] (x) -- (wu);
\draw[arrow] (wg) -- (silu);
\draw[arrow] (wu) -- (up);
\draw[arrow] (silu) -- (mul);
\draw[arrow] (up) -- (mul);
\draw[arrow] (mul) -- (wd);
\draw[arrow] (wd) -- (out);

\node[align=left, text width=10.0cm] at (0,-2.35) {
  SwiGLU 使用门控分支调制内容分支。工程上，$W_g$ 和 $W_u$ 常可合并为一个更大的 projection，
  再在 hidden 维切分，减少输入读取和 kernel launch。
};

\end{tikzpicture}
\caption{SwiGLU 门控 FFN 结构}
\label{fig:swiglu-ffn}
\end{figure}

\subsection{FFN 为什么常成为参数量大户}
\label{subsec:why-ffn-parameter-heavy}

设隐藏维度 $H=4096$，标准 FFN 中间维度 $H_{\mathrm{ffn}}=4H=16384$。两层线性层参数量约为：
\[
4096\times 16384+16384\times 4096
=================================

134,217,728.
\]
这是一层 FFN 的参数量，不含 bias。若模型有 $80$ 层，仅 FFN 部分就会达到约：
\[
80\times 134,217,728
\approx 10.7\mathrm{B}
\]
参数。

FFN 的计算也很重。对于 batch 和 sequence 合并后的 token 数：
\[
N_{\mathrm{tok}}=B S,
\]
FFN 两个 GEMM 的 FLOPs 约为：
\[
2N_{\mathrm{tok}}H H_{\mathrm{ffn}}
+
2N_{\mathrm{tok}}H_{\mathrm{ffn}}H
==================================

4N_{\mathrm{tok}}H H_{\mathrm{ffn}}.
\]
当 $H_{\mathrm{ffn}}=4H$ 时：
\[
\mathrm{FLOPs}*{\mathrm{FFN}}\approx 16N*{\mathrm{tok}}H^2.
\]

因此，在 GPU profiling 中，Transformer block 的时间往往大量消耗在 FFN GEMM 上。好消息是 FFN GEMM 规则、密集、算术强度高，通常能很好利用 Tensor Core。坏消息是参数和激活体积大，在 tensor parallel 和 ZeRO 中会带来显著通信与显存压力。

\section{Positional Encoding}
\label{sec:positional-encoding}

Self-Attention 本身对输入 token 的排列没有天然顺序感。若没有位置编码，模型只知道 token 内容，不知道它们出现在第几个位置。为了建模序列顺序，需要向 token embedding 注入位置信息。

\subsection{绝对位置编码}
\label{subsec:absolute-positional-encoding}

最直接的方法是为每个位置学习一个位置向量：
\[
\mat{P}\in\mathbb{R}^{S_{\max}\times H}.
\]
输入 token embedding 为：
\[
\mat{E}\in\mathbb{R}^{S\times H}.
\]
加入位置编码后：
\[
\mat{X}=\mat{E}+\mat{P}_{1:S}.
\]

绝对位置编码简单有效，但存在外推问题：如果训练时最大长度为 $S_{\max}$，推理时想处理更长序列，模型没有学过更远位置的 embedding。即便可以插值或扩展，也不一定稳定。

早期 Transformer 也使用正弦位置编码：
\[
PE(pos,2i)=\sin\left(\frac{pos}{10000^{2i/H}}\right),
\]
\[
PE(pos,2i+1)=\cos\left(\frac{pos}{10000^{2i/H}}\right).
\]
它不需要学习参数，并具有一定外推结构。但现代 Decoder-only LLM 更常见的是 RoPE。

\subsection{相对位置编码}
\label{subsec:relative-positional-encoding}

绝对位置编码告诉模型“当前位置是第几个 token”；相对位置编码更关注“两个 token 之间相隔多远”。在 attention 中，位置关系主要影响 query 与 key 的匹配，因此可以把相对位置信息加入 attention score：

\[
S_{ij}
======

\vect{q}_i^{\mathsf{T}}\vect{k}*j
+
b*{i-j}.
\]

其中 $b_{i-j}$ 是与相对距离有关的 bias。相对位置编码适合表达局部关系、距离衰减和长上下文结构。ALiBi 等方法也是沿着类似方向设计的。

从 Infra 角度看，相对位置 bias 的实现会影响 attention kernel。如果 bias 需要显式构造 $S\times S$ 矩阵，会增加显存与带宽；更高效的方法是在 kernel 内根据 $(i,j)$ 动态计算 bias。

\subsection{RoPE}
\label{subsec:rope}

RoPE，即 Rotary Position Embedding，旋转位置编码，是现代 LLM 中非常常见的位置编码方式。它不是把位置向量加到 token embedding 上，而是对 query 和 key 的向量分量按位置进行旋转。

对于二维向量：
\[
\begin{bmatrix}
x_{2i}\
x_{2i+1}
\end{bmatrix},
\]
位置 $pos$ 对应旋转角度 $\theta_{pos,i}$，旋转后：
\[
\begin{bmatrix}
x'*{2i}\
x'*{2i+1}
\end{bmatrix}
=============

\begin{bmatrix}
\cos\theta_{pos,i} & -\sin\theta_{pos,i}\
\sin\theta_{pos,i} & \cos\theta_{pos,i}
\end{bmatrix}
\begin{bmatrix}
x_{2i}\
x_{2i+1}
\end{bmatrix}.
\]

RoPE 的关键性质是，旋转后的 query 与 key 点积可以表达相对位置信息。不同位置的旋转角度差会进入点积，从而让 attention score 感知相对距离。

\begin{figure}[H]
\centering
\begin{tikzpicture}[
font=\small,
axis/.style={-{Latex[length=2.1mm]}, thick},
vec/.style={-{Latex[length=2.5mm]}, thick, blue!75},
vec2/.style={-{Latex[length=2.5mm]}, thick, orange!85!black},
arrow/.style={-{Latex[length=2.3mm]}, thick}
]

\node[font=\bfseries] at (0,3.0) {RoPE：按位置旋转 Query / Key 的二维子空间};

\draw[axis] (-4.2,0) -- (-1.0,0) node[right] {$x_{2i}$};
\draw[axis] (-2.6,-1.4) -- (-2.6,1.6) node[above] {$x_{2i+1}$};
\draw[vec] (-2.6,0) -- (-1.45,0.8) node[above] {$\vect{x}$};

\draw[axis] (1.0,0) -- (4.2,0) node[right] {$x'_{2i}$};
\draw[axis] (2.6,-1.4) -- (2.6,1.6) node[above] {$x'_{2i+1}$};
\draw[vec2] (2.6,0) -- (3.15,1.25) node[above] {$R_{\theta(pos)}\vect{x}$};

\draw[arrow] (-0.6,0.4) -- node[above] {position-dependent rotation} (0.6,0.4);

\draw[->, thick, gray!70] (3.25,0.0) arc (0:65:0.65);
\node at (3.55,0.45) {$\theta$};

\node[align=left, text width=10.2cm] at (0,-2.15) {
  RoPE 将位置信息注入 Q/K，而不是直接加到 hidden state 上。
  它让 attention score 自然包含相对位置关系，是 LLaMA 等现代 Decoder-only LLM 的常见选择。
};

\end{tikzpicture}
\caption{RoPE 旋转位置编码的二维几何直觉}
\label{fig:rope-rotation}
\end{figure}

下面给出一个简化的 RoPE 实现：

\begin{lstlisting}[language=Python,caption={简化版 RoPE 应用逻辑},label={lst:rope-implementation}]
import torch

def rotate_half(x):
# x: [..., dim]
x_even = x[..., 0::2]
x_odd = x[..., 1::2]
# [x0, x1] -> [-x1, x0]
out = torch.stack((-x_odd, x_even), dim=-1)
return out.flatten(-2)

def apply_rope(q, k, cos, sin):
"""
q, k: [batch, heads, seq, head_dim]
cos, sin: [seq, head_dim], broadcastable to q/k
"""
cos = cos[None, None, :, :]
sin = sin[None, None, :, :]

q_rot = q * cos + rotate_half(q) * sin
k_rot = k * cos + rotate_half(k) * sin
return q_rot, k_rot

\end{lstlisting}

RoPE 的工程细节包括：

\begin{itemize}
\item cos/sin cache 的预计算与设备放置；
\item 长上下文扩展时的 base 调整或 scaling；
\item Q/K layout 与 head dimension 对齐；
\item 是否在 fused attention kernel 中融合 RoPE；
\item prefill 与 decode 阶段位置索引管理。
\end{itemize}

\subsection{长上下文外推问题}
\label{subsec:long-context-extrapolation}

长上下文是现代 LLM 的重要需求。位置编码直接影响模型能否处理超过训练长度的序列。若模型训练时最大长度为 $4096$，推理时希望处理 $32768$ 或更长上下文，需要解决外推问题。

常见策略包括：

\begin{itemize}
\item 对绝对位置 embedding 进行插值或扩展；
\item 对 RoPE 的频率或 base 做 scaling；
\item 使用相对位置 bias 或 ALiBi 类方法；
\item 在长上下文数据上继续训练；
\item 结合 sliding window、attention sink、memory token 或 retrieval。
\end{itemize}

但长上下文不仅是模型问题，也是 Infra 问题。标准 attention 的 $O(S^2)$ 成本会迅速上升；推理阶段 KV Cache 的显存随 $S$ 线性增长：
\[
M_{\mathrm{KV}}
\propto
L\cdot B\cdot S\cdot n_{kv}\cdot d_h.
\]
因此，长上下文模型需要算法、kernel、显存管理和调度系统共同优化。

\section{Transformer Block 的现代变体}
\label{sec:modern-transformer-block-variants}

Transformer block 是 attention、FFN、归一化和 residual connection 的组合。不同模型族在细节上有所差异，例如 Post-Norm 与 Pre-Norm、LayerNorm 与 RMSNorm、GELU 与 SwiGLU、MHA 与 GQA、绝对位置编码与 RoPE 等。

\subsection{Post-Norm 与 Pre-Norm}
\label{subsec:postnorm-prenorm}

原始 Transformer 常使用 Post-Norm：
\[
\mat{Y}=\LayerNorm(\mat{X}+\operatorname{Sublayer}(\mat{X})).
\]
也就是说，先 residual add，再做 LayerNorm。

现代大语言模型更常使用 Pre-Norm：
\[
\mat{Y}=\mat{X}+\operatorname{Sublayer}(\LayerNorm(\mat{X})).
\]
归一化被放在 attention 或 FFN 之前。

Pre-Norm 的优势是深层网络训练更稳定。因为 residual stream 可以更直接地跨层传递，梯度也更容易沿 residual 路径流动。Post-Norm 在层数较深时更容易出现训练不稳定，需要更精细的初始化和学习率控制。

\begin{figure}[H]
\centering
\begin{tikzpicture}[
font=\small,
box/.style={draw, rounded corners=3pt, minimum width=1.55cm, minimum height=0.7cm, align=center},
arrow/.style={-{Latex[length=2.3mm]}, thick},
plus/.style={draw, circle, minimum size=0.55cm, fill=yellow!18}
]

\node[font=\bfseries] at (0,3.0) {Post-Norm 与 Pre-Norm 对比};

% PostNorm
\node[font=\bfseries] at (-3.8,2.25) {Post-Norm};
\node[box, fill=blue!6] (px) at (-5.5,1.2) {$X$};
\node[box, fill=orange!12] (psub) at (-3.7,1.2) {Sublayer};
\node[plus] (pplus) at (-2.2,1.2) {$+$};
\node[box, fill=green!10] (pln) at (-0.8,1.2) {Norm};
\node[box, fill=blue!6] (py) at (0.6,1.2) {$Y$};

\draw[arrow] (px) -- (psub);
\draw[arrow] (psub) -- (pplus);
\draw[arrow] (pplus) -- (pln);
\draw[arrow] (pln) -- (py);
\draw[arrow] (px.south) .. controls (-4.0,0.35) and (-2.2,0.35) .. (pplus.south);

% PreNorm
\node[font=\bfseries] at (-3.8,-0.1) {Pre-Norm};
\node[box, fill=blue!6] (x) at (-5.5,-1.1) {$X$};
\node[box, fill=green!10] (ln) at (-4.0,-1.1) {Norm};
\node[box, fill=orange!12] (sub) at (-2.4,-1.1) {Sublayer};
\node[plus] (plus) at (-0.8,-1.1) {$+$};
\node[box, fill=blue!6] (y) at (0.6,-1.1) {$Y$};

\draw[arrow] (x) -- (ln);
\draw[arrow] (ln) -- (sub);
\draw[arrow] (sub) -- (plus);
\draw[arrow] (plus) -- (y);
\draw[arrow] (x.south) .. controls (-3.3,-2.0) and (-0.8,-2.0) .. (plus.south);

\node[align=left, text width=5.4cm] at (4.0,0.1) {
  现代 LLM 多采用 Pre-Norm。\\
  它让 residual stream 更稳定，\\
  有利于深层网络中的梯度传播。
};

\end{tikzpicture}
\caption{Post-Norm 与 Pre-Norm Transformer 子层结构}
\label{fig:postnorm-prenorm}
\end{figure}

\subsection{LayerNorm 与 RMSNorm}
\label{subsec:transformer-layernorm-rmsnorm}

LayerNorm 对 hidden 维计算均值和方差：
\[
\operatorname{LN}(\vect{x})
===========================

\gamma
\frac{\vect{x}-\mu}{\sqrt{\sigma^2+\epsilon}}
+
\beta.
\]
RMSNorm 则只计算均方根：
\[
\RMSNorm(\vect{x})
==================

\gamma
\frac{\vect{x}}{\sqrt{\frac{1}{H}\sum_i x_i^2+\epsilon}}.
\]

许多现代 LLM 使用 RMSNorm，原因包括：

\begin{itemize}
\item 计算略简单，不需要减均值；
\item 不依赖 batch 统计，适合小 micro-batch；
\item 在 Decoder-only 架构中训练表现稳定；
\item 更容易写出高效 fused kernel。
\end{itemize}

归一化层的参数量很小，只有：
\[
O(H)
\]
级别。但它们对训练稳定性和 kernel 时间线有重要影响。Norm 通常是 memory-bound 操作，若与 residual add 或 dropout 拆成多个 kernel，会造成额外 HBM 往返。因此生产系统中经常使用 fused RMSNorm、fused LayerNorm 或 fused residual-norm。

\subsection{Residual Connection}
\label{subsec:residual-connection}

Residual connection 是深层 Transformer 能够训练的关键机制。基本形式为：
\[
\mat{Y}=\mat{X}+F(\mat{X}).
\]
它允许模型学习相对于输入的增量，而不是每层都完全重写表示。梯度也可以沿着恒等路径传播，缓解深层网络中的梯度衰减。

在 Transformer block 中，通常有两个 residual 分支：

\[
\mat{X}'=\mat{X}+\operatorname{Attention}(\operatorname{Norm}(\mat{X})),
\]
\[
\mat{Y}=\mat{X}'+\operatorname{FFN}(\operatorname{Norm}(\mat{X}')).
\]

Residual stream 可以被理解为贯穿所有层的信息高速公路。Attention 和 FFN 每层向这个高速公路写入增量信息。许多可解释性研究也会从 residual stream 的角度分析模型内部表征。

从系统角度看，residual add 是简单 elementwise 操作，计算量很小但需要读写完整激活张量。若单独作为 kernel 执行，可能成为时间线上的小 kernel。高性能实现常把 residual add 与 norm 融合。

\subsection{Decoder-only 架构}
\label{subsec:decoder-only-architecture}

现代大语言模型大多采用 Decoder-only Transformer。它的核心特点是：

\begin{itemize}
\item 使用 causal self-attention；
\item 训练目标是 next-token prediction；
\item 输入和输出共享或不共享 embedding/lm head；
\item 推理时逐 token 生成，并使用 KV Cache；
\item 更适合统一处理文本补全、对话、代码生成和指令跟随任务。
\end{itemize}

一个 Pre-Norm Decoder Block 可以写作：

\[
\mat{U}_{\ell}
==============

\mat{X}*{\ell}
+
\operatorname{Attention}
\left(
\operatorname{Norm}*1(\mat{X}*{\ell})
\right),
\]
\[
\mat{X}*{\ell+1}
================

\mat{U}_{\ell}
+
\operatorname{FFN}
\left(
\operatorname{Norm}*2(\mat{U}*{\ell})
\right).
\]

整个模型则是 embedding、多个 decoder block 和 lm head 的堆叠：

\[
\mat{X}*0=\operatorname{Embedding}(\text{input ids})+\operatorname{PositionEncoding},
\]
\[
\mat{X}*{\ell+1}=\operatorname{DecoderBlock}*{\ell}(\mat{X}*{\ell}),
\quad
\ell=0,\ldots,L-1,
\]
\[
\mat{Z}=\mat{X}*{L}\mat{W}*{\mathrm{vocab}}.
\]

其中 logits：
\[
\mat{Z}\in\mathbb{R}^{B\times S\times V}
\]
用于计算 cross entropy 或进行采样。

\begin{figure}[H]
\centering
\begin{tikzpicture}[
font=\small,
box/.style={draw, rounded corners=3pt, minimum width=2.15cm, minimum height=0.72cm, align=center},
arrow/.style={-{Latex[length=2.3mm]}, thick},
plus/.style={draw, circle, minimum size=0.55cm, fill=yellow!18}
]

\node[font=\bfseries] at (0,4.0) {Pre-Norm Transformer Decoder Block};

\node[box, fill=blue!6] (x) at (-4.9,2.4) {$X_\ell$};

\node[box, fill=green!10] (norm1) at (-2.9,2.4) {RMSNorm\\or LayerNorm};
\node[box, fill=orange!12] (attn) at (-0.5,2.4) {Causal\\Self-Attention};
\node[plus] (plus1) at (1.45,2.4) {$+$};
\node[box, fill=blue!6] (u) at (3.15,2.4) {$U_\ell$};

\draw[arrow] (x) -- (norm1);
\draw[arrow] (norm1) -- (attn);
\draw[arrow] (attn) -- (plus1);
\draw[arrow] (plus1) -- (u);
\draw[arrow] (x.south) .. controls (-2.0,1.45) and (1.45,1.45) .. (plus1.south);

\node[box, fill=green!10] (norm2) at (-1.6,0.3) {RMSNorm\\or LayerNorm};
\node[box, fill=purple!10] (ffn) at (0.9,0.3) {FFN / SwiGLU};
\node[plus] (plus2) at (3.0,0.3) {$+$};
\node[box, fill=blue!6] (y) at (5.0,0.3) {$X_{\ell+1}$};

\draw[arrow] (u.south) .. controls (3.15,1.15) and (-1.6,1.15) .. (norm2.north);
\draw[arrow] (norm2) -- (ffn);
\draw[arrow] (ffn) -- (plus2);
\draw[arrow] (plus2) -- (y);
\draw[arrow] (u.south) .. controls (4.0,1.2) and (3.0,1.15) .. (plus2.north);

\node[align=left, text width=10.4cm] at (0,-1.55) {
  现代 Decoder-only LLM 通常使用 Pre-Norm、causal self-attention、SwiGLU FFN、RoPE 以及 RMSNorm。
  每个 block 都围绕 residual stream 添加 attention 和 FFN 两类增量。
};

\end{tikzpicture}
\caption{现代 Pre-Norm Transformer Decoder Block}
\label{fig:pre-norm-transformer-decoder-block}
\end{figure}

下面给出一个教学版 Decoder Block 的 PyTorch 实现。它省略了高性能 fused kernel、KV Cache、tensor parallel、dropout、FlashAttention 等细节，但结构上对应现代 LLM block。

\begin{lstlisting}[language=Python,caption={教学版 Transformer Decoder Block},label={lst:decoder-block-pytorch}]
import math
import torch
import torch.nn as nn
import torch.nn.functional as F

class RMSNorm(nn.Module):
def **init**(self, hidden_size, eps=1e-6):
super().**init**()
self.weight = nn.Parameter(torch.ones(hidden_size))
self.eps = eps

def forward(self, x):
    variance = x.pow(2).mean(dim=-1, keepdim=True)
    x = x * torch.rsqrt(variance + self.eps)
    return self.weight * x

class MultiHeadAttention(nn.Module):
def **init**(self, hidden_size, num_heads):
super().**init**()
assert hidden_size % num_heads == 0
self.hidden_size = hidden_size
self.num_heads = num_heads
self.head_dim = hidden_size // num_heads

    # Fused QKV projection.
    self.qkv = nn.Linear(hidden_size, 3 * hidden_size, bias=False)
    self.out_proj = nn.Linear(hidden_size, hidden_size, bias=False)

def forward(self, x):
    B, S, H = x.shape

    qkv = self.qkv(x)  # [B, S, 3H]
    q, k, v = qkv.chunk(3, dim=-1)

    # [B, S, H] -> [B, num_heads, S, head_dim]
    q = q.view(B, S, self.num_heads, self.head_dim).transpose(1, 2)
    k = k.view(B, S, self.num_heads, self.head_dim).transpose(1, 2)
    v = v.view(B, S, self.num_heads, self.head_dim).transpose(1, 2)

    scores = q @ k.transpose(-2, -1)
    scores = scores / math.sqrt(self.head_dim)

    causal_mask = torch.triu(
        torch.ones(S, S, device=x.device, dtype=torch.bool),
        diagonal=1,
    )
    scores = scores.masked_fill(causal_mask, float("-inf"))

    probs = F.softmax(scores, dim=-1)
    out = probs @ v  # [B, heads, S, head_dim]

    out = out.transpose(1, 2).contiguous().view(B, S, H)
    return self.out_proj(out)

class SwiGLU(nn.Module):
def **init**(self, hidden_size, intermediate_size):
super().**init**()
self.gate_proj = nn.Linear(hidden_size, intermediate_size, bias=False)
self.up_proj = nn.Linear(hidden_size, intermediate_size, bias=False)
self.down_proj = nn.Linear(intermediate_size, hidden_size, bias=False)

def forward(self, x):
    gate = F.silu(self.gate_proj(x))
    up = self.up_proj(x)
    return self.down_proj(gate * up)

class DecoderBlock(nn.Module):
def **init**(self, hidden_size, num_heads, intermediate_size):
super().**init**()
self.norm1 = RMSNorm(hidden_size)
self.attn = MultiHeadAttention(hidden_size, num_heads)
self.norm2 = RMSNorm(hidden_size)
self.mlp = SwiGLU(hidden_size, intermediate_size)

def forward(self, x):
    # Pre-Norm attention residual.
    x = x + self.attn(self.norm1(x))

    # Pre-Norm FFN residual.
    x = x + self.mlp(self.norm2(x))
    return x

\end{lstlisting}

这段代码适合理解结构，但不适合直接作为高性能 LLM 实现。它的低效点包括：

\begin{itemize}
\item 显式物化 attention score；
\item causal mask 每次 forward 都构造；
\item softmax 和 matmul 没有融合；
\item transpose 后调用 \texttt{contiguous()} 可能产生额外拷贝；
\item 没有 KV Cache，无法高效 decode；
\item 没有 tensor parallel、sequence parallel 或 fused RMSNorm；
\item 没有使用 FlashAttention 或 PyTorch SDPA backend。
\end{itemize}

实际大模型框架会用高度优化的 attention kernel、fused norm、fused MLP、并行线性层和分布式通信策略替换这些朴素实现。

\section{从 Transformer 公式到系统成本}
\label{sec:transformer-formula-to-system-cost}

Transformer 的优雅之处在于结构规整；它的困难之处在于规模巨大。对 AI Infra 工程师来说，理解 Transformer 不能停留在 block 图，而要能估算每一部分的参数量、显存、计算量和通信代价。

\subsection{单层参数量估算}
\label{subsec:transformer-layer-parameter-estimation}

设 hidden size 为 $H$，FFN 中间维度为 $H_{\mathrm{ffn}}$。一个标准 Decoder Block 的主要参数包括：

\begin{itemize}
\item QKV projection：$3H^2$；
\item output projection：$H^2$；
\item FFN up/down projection：$2HH_{\mathrm{ffn}}$；
\item Norm 参数：约 $O(H)$，相对较小。
\end{itemize}

因此单层参数量近似：
\[
N_{\mathrm{layer}}
\approx
4H^2+2HH_{\mathrm{ffn}}.
\]
若 $H_{\mathrm{ffn}}=4H$：
\[
N_{\mathrm{layer}}
\approx
4H^2+8H^2=12H^2.
\]

若使用 SwiGLU，参数量为：
\[
N_{\mathrm{FFN}}
\approx
3HH_{\mathrm{mid}}.
\]
为了与标准 FFN 参数量接近，常将 $H_{\mathrm{mid}}$ 设置为小于 $4H$ 的值。

\subsection{Attention 与 FFN 的 FLOPs}
\label{subsec:attention-ffn-flops}

对于 batch token 数：
\[
N_{\mathrm{tok}}=B S,
\]
attention projection FLOPs 约为：
\[
\mathrm{FLOPs}*{QKV+O}
\approx
2N*{\mathrm{tok}}\cdot 4H^2
===========================

8BSH^2.
\]
attention score 和 value 加权的 FLOPs 约为：
\[
\mathrm{FLOPs}_{attn}
\approx
4B n_h S^2 d_h
==============

4BS^2H.
\]
FFN FLOPs 对标准 $H_{\mathrm{ffn}}=4H$ 约为：
\[
\mathrm{FLOPs}_{FFN}
\approx
16BSH^2.
\]

这说明当 $S$ 不是特别大时，FFN 和 projection GEMM 往往占主导；当 $S$ 变得很大时，$S^2H$ 的 attention 项会迅速上升。这也是长上下文训练特别昂贵的根本原因。

\begin{figure}[H]
\centering
\begin{tikzpicture}[
font=\small,
axis/.style={-{Latex[length=2.3mm]}, thick},
curve/.style={thick},
label/.style={align=center}
]

\node[font=\bfseries] at (0,3.0) {Attention 与 FFN 计算量随序列长度的增长趋势};

\draw[axis] (-5.2,-1.2) -- (5.4,-1.2) node[right] {$S$};
\draw[axis] (-5.2,-1.2) -- (-5.2,2.2) node[above] {relative cost};

\draw[curve, blue!75]
  plot[smooth] coordinates {
    (-4.8,-0.7) (-3.2,-0.2) (-1.6,0.3) (0.0,0.8) (1.6,1.3) (3.2,1.8)
  };
\node[blue!75] at (3.7,1.9) {FFN / Projection\\$\propto S$};

\draw[curve, orange!90!black]
  plot[smooth] coordinates {
    (-4.8,-0.9) (-3.2,-0.8) (-1.6,-0.45) (0.0,0.1) (1.6,0.95) (3.2,2.05)
  };
\node[orange!90!black] at (3.8,0.85) {Attention scores\\$\propto S^2$};

\draw[dashed] (0.9,-1.2) -- (0.9,1.0);
\node[align=center] at (0.9,-1.65) {长上下文\\拐点};

\node[align=left, text width=10.2cm] at (0,-2.55) {
  Transformer 的瓶颈随序列长度变化。短到中等长度时，FFN 与 projection GEMM 常占主要时间；
  长上下文时，attention 的 $S^2$ 项和 KV Cache 压力变得突出。
};

\end{tikzpicture}
\caption{Attention 与 FFN 计算复杂度增长趋势}
\label{fig:attention-ffn-cost-growth}
\end{figure}

\subsection{训练与推理的差异}
\label{subsec:transformer-training-inference-difference}

Transformer 训练和推理有显著差异：

\begin{itemize}
\item \textbf{训练}：给定完整序列，所有位置并行计算；需要 backward、保存激活、优化器状态和梯度通信。
\item \textbf{推理 prefill}：给定 prompt，类似训练 forward，处理完整上下文；不需要 backward。
\item \textbf{推理 decode}：每次只生成一个新 token；使用 KV Cache 避免重复计算历史 K/V。
\end{itemize}

训练阶段的 attention 可以一次处理 $S\times S$ 的关系；decode 阶段每一步只计算新 token 对历史所有 token 的 attention：
\[
\mat{q}*{new}\mat{K}*{cache}^{\mathsf{T}}.
\]
这时计算量随历史长度线性增长，但由于 batch 和矩阵形状较小，往往更 memory-bound。KV Cache 的读取带宽和请求调度成为推理系统关键。

\begin{table}[H]
\centering
\small
\caption{Transformer 训练、Prefill 与 Decode 的系统差异}
\label{tab:training-prefill-decode}
\begin{tabular}{p{0.18\textwidth}p{0.25\textwidth}p{0.25\textwidth}p{0.22\textwidth}}
\toprule
\textbf{阶段} & \textbf{输入形态} & \textbf{主要瓶颈} & \textbf{典型优化} \
\midrule
训练
& 大 batch，完整序列
& GEMM、attention、激活显存、通信
& FlashAttention、并行训练、checkpointing、fused kernel \
\midrule
Prefill
& prompt 完整序列
& 大矩阵计算与 attention IO
& FlashAttention、chunked prefill、continuous batching \
\midrule
Decode
& 每步一个新 token
& KV Cache 带宽、调度、小 batch kernel
& PagedAttention、GQA/MQA、continuous batching、量化 \
\bottomrule
\end{tabular}
\end{table}

\section{本章小结}
\label{sec:chapter04-summary}

本章系统剖析了 Transformer 架构，从 RNN 的串行瓶颈讲到 Self-Attention 的并行优势，再到现代 Decoder-only LLM 的具体 block 结构。

\begin{itemize}
\item \textbf{RNN 的核心瓶颈是时间步串行依赖}，难以充分利用 GPU 的大规模并行能力。
\item \textbf{Attention 的本质是根据相似度对 value 加权求和}，矩阵形式为 $\softmax(QK^{\mathsf{T}}/\sqrt{d_k})V$。
\item \textbf{Scaled Dot-Product Attention 通过 $\sqrt{d_k}$ 控制数值尺度}，并通过 mask 机制支持 padding 和自回归训练。
\item \textbf{Multi-Head Attention 在多个子空间中并行建模 token 关系}，同时天然适配 batched GEMM 和 tensor parallel。
\item \textbf{MHA、MQA、GQA 的差异主要体现在 K/V head 共享方式}，这对推理阶段 KV Cache 显存和带宽有重要影响。
\item \textbf{FFN 往往是 Transformer 参数量大户}，尤其当中间维度接近 $4H$ 时，其参数量和 FLOPs 都非常显著。
\item \textbf{位置编码让 attention 感知序列顺序}，RoPE 通过旋转 Q/K 注入相对位置信息，是现代 LLM 的常见选择。
\item \textbf{现代 Decoder-only Block 通常采用 Pre-Norm、RMSNorm、Causal Self-Attention、SwiGLU 和 Residual Connection}。
\item \textbf{Transformer 的系统瓶颈随阶段变化}：训练关注激活、GEMM、attention 和通信；推理 decode 关注 KV Cache、带宽和调度。
\end{itemize}

下一章将继续从模型演进角度出发，梳理大语言模型的发展史与主流范式：从 N-gram、神经语言模型、GPT 系列，到 instruction tuning、RLHF、开源 LLM 生态和 scaling law。那一章会帮助我们理解为什么今天的 AI Infra 需要服务如此巨大的模型、数据和训练任务。

\chapter{大语言模型发展史与主流范式}
\label{chap:llm-history}

\section{语言模型的任务定义}
\label{sec:language-modeling-task-definition}

上一章我们从结构层面拆解了 Transformer：Self-Attention、Multi-Head Attention、FFN、位置编码、Pre-Norm、RMSNorm、RoPE、GQA 等组件共同构成了现代 Decoder-only LLM 的主体。但仅仅理解结构还不够。一个模型为什么要做 next-token prediction？为什么 GPT 系列把“预训练 + 提示”变成主流范式？为什么后来的模型又需要 instruction tuning、RLHF、DPO 或其他偏好优化？为什么开源模型生态对 AI Infrastructure 产生如此大的推动？

本章要回答的是：\textbf{大语言模型作为一种工程系统，是如何从统计语言模型、神经语言模型、Transformer 预训练，逐步演化到今天的基础模型、聊天模型、推理模型和多模态模型的。}

对于 AI Infra 工程师而言，了解模型发展史不是为了背诵年份和名称，而是为了理解工作负载的变化：

\begin{itemize}
\item 从 N-gram 到神经语言模型，计算从查表统计转向矩阵运算；
\item 从小模型 fine-tuning 到大模型 prompting，服务方式从离线训练转向在线推理；
\item 从预训练模型到 Chat Model，训练流水线从单阶段变成多阶段；
\item 从 dense model 到 MoE，系统瓶颈从单纯矩阵计算扩展到路由、负载均衡和 all-to-all 通信；
\item 从封闭模型到开源模型，部署形态从 API 调用扩展到私有化训练、微调、量化和本地推理。
\end{itemize}

本章会沿着模型范式的演进路线展开，并在每个阶段指出它对 AI Infra 的具体影响。

\subsection{N-gram 到神经语言模型}
\label{subsec:ngram-to-neural-lm}

语言模型的基本任务是估计一个 token 序列出现的概率。给定序列：
\[
x_1,x_2,\ldots,x_T,
\]
语言模型希望建模联合概率：
\[
p(x_1,x_2,\ldots,x_T).
\]
根据概率链式法则：
\[
p(x_1,x_2,\ldots,x_T)
=====================

\prod_{t=1}^{T}
p(x_t\mid x_1,x_2,\ldots,x_{t-1}).
\]

N-gram 语言模型做了一个强假设：当前 token 只依赖前面有限的 $n-1$ 个 token：
\[
p(x_t\mid x_{<t})
\approx
p(x_t\mid x_{t-n+1},\ldots,x_{t-1}).
\]
例如 trigram 模型中：
\[
p(x_t\mid x_{<t})
\approx
p(x_t\mid x_{t-2},x_{t-1}).
\]

N-gram 的优点是简单、可解释、训练成本低；缺点也很明显：

\begin{itemize}
\item \textbf{上下文窗口短}：无法有效建模长距离依赖；
\item \textbf{稀疏性严重}：许多 token 组合在训练集中从未出现；
\item \textbf{泛化能力弱}：无法理解语义相似性，例如“猫”和“狗”在统计表中是完全不同的符号；
\item \textbf{存储成本高}：高阶 N-gram 表会急剧膨胀。
\end{itemize}

神经语言模型的重要突破，是用连续向量表示 token。每个 token 不再只是离散 id，而是 embedding table 中的一行：
\[
\mat{E}\in\mathbb{R}^{V\times H}.
\]
token id $x_t$ 对应向量：
\[
\vect{e}*t=\mat{E}*{x_t,:}.
\]
这样，语义相近的 token 可以在向量空间中接近，模型能够更好地泛化。

从系统角度看，这一步把语言建模从离散计数表转化为矩阵乘法、非线性网络和梯度下降。也就是说，语言模型开始成为 GPU 友好的工作负载。Embedding lookup、矩阵乘法、softmax 和反向传播，逐渐取代了传统统计表查询。

\subsection{自回归语言建模}
\label{subsec:autoregressive-language-modeling}

自回归语言模型是现代 Decoder-only LLM 的核心任务。它按从左到右的方式分解序列概率：
\[
p(x_1,x_2,\ldots,x_T)
=====================

\prod_{t=1}^{T}
p(x_t\mid x_{<t}).
\]
训练目标是最大化真实序列的对数似然：
\[
\max_{\theta}
\sum_{t=1}^{T}
\log p_{\theta}(x_t\mid x_{<t}).
\]
等价地，最小化负对数似然：
\[
\mathcal{L}(\theta)
===================

*

\sum_{t=1}^{T}
\log p_{\theta}(x_t\mid x_{<t}).
\]

在 Transformer Decoder 中，训练时可以并行计算所有位置的 logits。给定输入：
\[
x_1,x_2,\ldots,x_{T-1},
\]
模型预测：
\[
x_2,x_3,\ldots,x_T.
\]
这常被称为 shift-one-token 训练：

\[
\text{input}=[x_1,x_2,\ldots,x_{T-1}],
\]
\[
\text{label}=[x_2,x_3,\ldots,x_T].
\]

\begin{figure}[H]
\centering
\begin{tikzpicture}[
font=\small,
token/.style={draw, rounded corners=2pt, minimum width=0.85cm, minimum height=0.5cm, align=center, fill=blue!6},
label/.style={draw, rounded corners=2pt, minimum width=0.85cm, minimum height=0.5cm, align=center, fill=orange!14},
arrow/.style={-{Latex[length=2.2mm]}, thick}
]

\node[font=\bfseries] at (0,2.8) {自回归语言建模：输入与标签右移一位};

\foreach \i/\txt/\x in {1/从/-4.8,2/基础/-3.6,3/到/-2.4,4/深渊/-1.2,5/：/0.0,6/AI/1.2,7/Infra/2.4} {
  \node[token] (in\i) at (\x,1.2) {\txt};
}

\foreach \i/\txt/\x in {1/基础/-4.8,2/到/-3.6,3/深渊/-2.4,4/：/-1.2,5/AI/0.0,6/Infra/1.2,7/指南/2.4} {
  \node[label] (lab\i) at (\x,0.0) {\txt};
  \draw[arrow] (in\i) -- (lab\i);
}

\node[align=right] at (-5.95,1.2) {input};
\node[align=right] at (-5.95,0.0) {label};

\node[align=left, text width=9.8cm] at (0,-1.35) {
  训练时所有位置的预测可并行计算，但 causal mask 保证每个位置只能使用它之前的上下文。
  推理时则必须逐 token 生成，并依赖 KV Cache 降低重复计算。
};

\end{tikzpicture}
\caption{自回归语言建模中的输入与标签对齐方式}
\label{fig:autoregressive-shifted-label}
\end{figure}

自回归建模对 AI Infra 有两个重要影响。第一，训练阶段可以高效并行，因为完整序列一次进入 GPU；第二，推理阶段天然是逐 token 递推，因为生成的 token 会成为下一步输入。这正是 LLM serving 与普通分类模型 serving 的根本差异。

\subsection{Masked Language Modeling}
\label{subsec:masked-language-modeling}

Masked Language Modeling，简称 MLM，是 BERT 等 Encoder-only 模型采用的经典任务。它不是从左到右预测下一个 token，而是在输入中随机 mask 一部分 token，让模型根据双向上下文恢复它们。

给定序列：
\[
x_1,x_2,\ldots,x_T,
\]
随机选择 mask 集合 $\mathcal{M}$，将对应位置替换为特殊 token \texttt{[MASK]}。训练目标是：
\[
\mathcal{L}_{\mathrm{MLM}}
==========================

*

\sum_{t\in\mathcal{M}}
\log p_{\theta}(x_t\mid x_{\setminus \mathcal{M}}).
\]

MLM 的优点是能够使用双向上下文，适合理解类任务，例如文本分类、句子匹配、命名实体识别等。缺点是它不直接对应生成式推理过程。对于需要逐 token 生成的任务，自回归语言模型更自然。

从系统角度看，Encoder-only 模型和 Decoder-only 模型的推理形态不同：

\begin{itemize}
\item Encoder-only 模型通常一次性编码完整输入，输出分类或 token-level 结果；
\item Decoder-only 模型需要自回归生成，decode 阶段不断读取和扩展 KV Cache；
\item Encoder-Decoder 模型在翻译、摘要等任务中仍常见，但大规模通用聊天模型更偏向 Decoder-only。
\end{itemize}

\subsection{Causal LM 与 Chat LM}
\label{subsec:causal-lm-and-chat-lm}

Causal Language Model 是基础的自回归模型。它只学习：
\[
p(x_t\mid x_{<t}).
\]
它可以生成文本、补全代码、续写故事，但并不天然理解“用户指令”。例如用户输入“请总结下面这段话”，基础 LM 可能继续写类似训练语料中的文本，而不是稳定地执行总结任务。

Chat LM 则是在 Causal LM 的基础上，经过指令微调、偏好优化和安全对齐，使模型更倾向于遵循用户意图。它的输入通常被组织成对话格式：

\[
\texttt{<system>} \quad \text{系统指令}
\]
\[
\texttt{<user>} \quad \text{用户问题}
\]
\[
\texttt{<assistant>} \quad \text{模型回答}
\]

从数学上看，Chat LM 仍然是自回归模型，仍然预测下一个 token。但数据分布和训练目标发生了变化：模型不再只是拟合互联网文本，而是拟合“在给定指令和对话上下文下生成有帮助回答”的分布。

这一区分对 Infra 很重要。Base Model 和 Chat Model 的服务模式不同：

\begin{itemize}
\item Base Model 常用于续写、评测、微调和作为后续训练起点；
\item Chat Model 常用于在线对话服务，需要安全策略、系统提示、工具调用、上下文管理和多轮对话缓存；
\item Chat serving 中 prompt 模板、role token、stop token、采样策略都会影响实际输出。
\end{itemize}

\section{GPT 系列架构演进}
\label{sec:gpt-family-evolution}

GPT 系列推动了现代 Decoder-only LLM 的主流范式。这里的重点不是逐一复述每个模型细节，而是理解范式变化：从预训练加微调，到规模扩展，再到 few-shot prompting、instruction tuning 和 RLHF。

\subsection{GPT-1：预训练与微调}
\label{subsec:gpt1-pretrain-finetune}

GPT-1 的重要意义在于确立了“先在大规模无标注文本上预训练，再在下游任务上微调”的路线。这个范式可以写作：

\[
\theta_{\mathrm{pretrain}}
==========================

\arg\min_{\theta}
\mathcal{L}*{\mathrm{LM}}(\theta;\mathcal{D}*{\mathrm{text}}),
\]
然后：
\[
\theta_{\mathrm{task}}
======================

\arg\min_{\theta}
\mathcal{L}*{\mathrm{task}}(\theta;\mathcal{D}*{\mathrm{task}}).
\]

预训练阶段学习通用语言表示，微调阶段适配具体任务。这个思想极大提高了样本效率，因为下游任务不再需要从随机初始化开始训练。

在 Infra 层面，这一阶段的模型规模还远小于今天的 LLM，但已经显露出预训练的工程特点：

\begin{itemize}
\item 需要大规模文本数据清洗和 tokenization；
\item 训练任务以长时间 GPU 计算为主；
\item checkpoint 和 experiment tracking 开始变得重要；
\item 下游任务微调需要可复用模型权重。
\end{itemize}

\subsection{GPT-2：规模扩展}
\label{subsec:gpt2-scaling}

GPT-2 进一步强调规模扩展。模型参数量、训练数据量和训练计算量都显著增加。它展示了一个重要趋势：在相同架构下，扩大模型和数据可以带来更强的生成能力。

这对 AI Infra 的含义非常直接：模型能力开始越来越依赖大规模计算。训练系统不再只是“能跑起来”，而要关注：

\begin{itemize}
\item 多 GPU 数据并行；
\item 高吞吐数据加载；
\item 稳定 checkpoint；
\item 长时间训练容错；
\item tokenizer 与数据管线的可复现性；
\item 日志与监控体系。
\end{itemize}

GPT-2 时代已经能看到一个趋势：语言模型训练不只是机器学习问题，也开始成为大规模系统工程问题。

\subsection{GPT-3：In-Context Learning}
\label{subsec:gpt3-in-context-learning}

GPT-3 的关键影响，是让人们看到大模型在无需参数更新的情况下，通过 prompt 中的少量示例完成任务。这被称为 In-Context Learning。

传统 fine-tuning 的形式是：
\[
\theta'
=======

\theta-\eta\nabla_{\theta}\mathcal{L}_{\mathrm{task}},
\]
即通过梯度更新参数适配任务。而 in-context learning 的形式是：
\[
y
=

f_{\theta}
(
\text{instruction},
\text{examples},
x
),
\]
参数 $\theta$ 不变，任务信息通过上下文提供。

这带来了使用方式的巨大变化。用户不一定需要训练新模型，而可以通过 prompt engineering、few-shot examples、chain-of-thought 提示、工具调用说明等方式引导模型行为。对于 Infra 来说，这意味着：

\begin{itemize}
\item 在线推理负载迅速增长；
\item prompt 长度成为成本因素；
\item 上下文窗口越长，prefill 成本越高；
\item 多租户服务需要处理高度动态的请求长度；
\item 缓存、批处理、限流和计费粒度需要围绕 token 设计。
\end{itemize}

\begin{figure}[H]
\centering
\begin{tikzpicture}[
font=\small,
block/.style={draw, rounded corners=3pt, minimum width=2.35cm, minimum height=0.78cm, align=center, fill=blue!6},
example/.style={draw, rounded corners=2pt, minimum width=2.2cm, minimum height=0.55cm, align=center, fill=orange!12},
arrow/.style={-{Latex[length=2.2mm]}, thick}
]

\node[font=\bfseries] at (0,3.0) {In-Context Learning：不更新参数，通过上下文描述任务};

\node[example] (inst) at (-4.4,1.55) {Instruction};
\node[example] (ex1) at (-2.0,1.55) {Example 1};
\node[example] (ex2) at (0.4,1.55) {Example 2};
\node[example] (query) at (2.8,1.55) {Query};

\draw[arrow] (inst) -- (ex1);
\draw[arrow] (ex1) -- (ex2);
\draw[arrow] (ex2) -- (query);

\node[block, fill=green!8] (model) at (0,0.2) {Frozen LLM\\parameters unchanged};

\draw[arrow] (query.south) .. controls (2.4,0.65) and (1.4,0.2) .. (model.east);
\draw[arrow] (inst.south) .. controls (-3.2,0.65) and (-1.6,0.2) .. (model.west);

\node[block, fill=purple!10] (out) at (0,-1.25) {Generated Answer};
\draw[arrow] (model) -- (out);

\node[align=left, text width=10.4cm] at (0,-2.55) {
  In-context learning 把任务适配从训练时参数更新转移到推理时上下文构造。
  这使 prompt 长度、prefill 计算和上下文管理成为 serving infra 的核心问题。
};

\end{tikzpicture}
\caption{In-Context Learning 的任务适配方式}
\label{fig:in-context-learning}
\end{figure}

\subsection{Instruction Tuning}
\label{subsec:instruction-tuning}

基础语言模型虽然强大，但它的目标只是预测下一个 token。若训练语料中出现问答格式，它可能学会回答；若没有明确格式，它可能只是续写。Instruction Tuning 的目标，是让模型更稳定地遵循自然语言指令。

监督指令微调，简称 SFT，可以写作：
\[
\mathcal{D}_{\mathrm{SFT}}
==========================

{
(p_i,r_i)
}*{i=1}^{N},
\]
其中 $p_i$ 是指令或对话 prompt，$r_i$ 是高质量参考回答。训练目标仍然是交叉熵：
\[
\mathcal{L}*{\mathrm{SFT}}
==========================

*

\sum_{i}
\sum_{t}
\log p_{\theta}(r_{i,t}\mid p_i,r_{i,<t}).
\]

Instruction Tuning 的关键不是改变模型架构，而是改变数据分布。模型学习到：当输入是一条用户指令时，应该输出满足指令的回答，而不是随意续写。

下面给出一个简化的 SFT 数据构造示意：

\begin{lstlisting}[language=Python,caption={监督指令微调 SFT 数据格式示意},label={lst:sft-data-format}]
def build_chat_sample(system_prompt, user_prompt, assistant_answer, tokenizer):
text = (
"<system>\n" + system_prompt + "\n"
"<user>\n" + user_prompt + "\n"
"<assistant>\n" + assistant_answer
)

tokens = tokenizer.encode(text)

# In SFT, loss is often computed only on assistant tokens.
# Tokens belonging to system/user prompt can be masked out with label = -100.
labels = build_labels_masking_non_assistant_tokens(tokens)

return {
    "input_ids": tokens[:-1],
    "labels": labels[1:],
}

\end{lstlisting}

SFT 对 Infra 的影响包括：

\begin{itemize}
\item 训练数据从网页文本转向高质量指令/对话样本；
\item 训练序列通常具有 role token、模板和 loss mask；
\item 数据质量比数据规模更敏感；
\item 微调可使用 LoRA、QLoRA、FSDP、ZeRO 等节省资源；
\item 多任务混合比例会影响模型行为。
\end{itemize}

\subsection{RLHF 与对齐训练}
\label{subsec:rlhf-alignment-training}

RLHF，即 Reinforcement Learning from Human Feedback，是 Chat Model 发展的关键技术之一。它的基本流程通常包含三个阶段：

\begin{enumerate}
\item \textbf{SFT}：用人工示范回答监督微调基础模型；
\item \textbf{Reward Model 训练}：收集人类对多个回答的偏好排序，训练奖励模型；
\item \textbf{RL 优化}：使用 PPO 等强化学习方法，让策略模型生成更符合偏好的回答。
\end{enumerate}

偏好数据通常形式为：
\[
(p, r^+, r^-),
\]
其中 $p$ 是 prompt，$r^+$ 是更受偏好的回答，$r^-$ 是较差回答。Reward Model 训练目标可写为：
\[
\mathcal{L}_{\mathrm{RM}}
=========================

*

\log
\sigma
\left(
R_{\phi}(p,r^+)-R_{\phi}(p,r^-)
\right).
\]
它鼓励奖励模型给更好回答更高分。

RL 阶段希望最大化奖励，同时避免模型偏离 SFT 模型太远。目标可以抽象为：
\[
\max_{\theta}
\mathbb{E}*{r\sim \pi*{\theta}(\cdot\mid p)}
\left[
R_{\phi}(p,r)
-
\beta
D_{\mathrm{KL}}
\left(
\pi_{\theta}(\cdot\mid p)
|
\pi_{\mathrm{ref}}(\cdot\mid p)
\right)
\right].
\]

这里的 KL 惩罚用于限制策略模型不要为了追求 reward 而产生过度偏移。RLHF 在工程上比普通 SFT 复杂得多，因为它涉及多个模型、在线生成、奖励打分、KL 控制和强化学习更新。

\begin{figure}[H]
\centering
\begin{tikzpicture}[
font=\small,
stage/.style={draw, rounded corners=3pt, minimum width=2.35cm, minimum height=0.85cm, align=center, fill=blue!6},
data/.style={draw, rounded corners=2pt, minimum width=2.0cm, minimum height=0.58cm, align=center, fill=orange!12},
model/.style={draw, rounded corners=3pt, minimum width=2.2cm, minimum height=0.72cm, align=center, fill=green!8},
arrow/.style={-{Latex[length=2.3mm]}, thick}
]

\node[font=\bfseries] at (0,3.5) {预训练、SFT、RLHF 三阶段流程};

\node[data] (raw) at (-5.4,2.2) {Raw Text\\Corpus};
\node[stage] (pretrain) at (-3.0,2.2) {Pretraining\\next-token loss};
\node[model] (base) at (-0.6,2.2) {Base LM};

\draw[arrow] (raw) -- (pretrain);
\draw[arrow] (pretrain) -- (base);

\node[data] (demo) at (-5.4,0.75) {Human\\Demos};
\node[stage] (sft) at (-3.0,0.75) {SFT\\CE loss};
\node[model] (sftmodel) at (-0.6,0.75) {SFT Model};

\draw[arrow] (demo) -- (sft);
\draw[arrow] (base) -- (sft);
\draw[arrow] (sft) -- (sftmodel);

\node[data] (pref) at (-5.4,-0.8) {Preference\\Ranking};
\node[stage] (rmtrain) at (-3.0,-0.8) {Reward Model\\Training};
\node[model] (rm) at (-0.6,-0.8) {Reward\\Model};

\draw[arrow] (pref) -- (rmtrain);
\draw[arrow] (rmtrain) -- (rm);

\node[stage, fill=purple!10, minimum width=2.5cm] (rl) at (2.0,0.0) {RLHF / PPO\\KL-regularized};
\node[model, fill=yellow!16] (chat) at (4.9,0.0) {Chat Model};

\draw[arrow] (sftmodel) -- (rl);
\draw[arrow] (rm) -- (rl);
\draw[arrow] (rl) -- (chat);

\node[align=left, text width=10.3cm] at (0,-2.35) {
  RLHF 把训练流水线从单一 next-token 预训练扩展为多模型、多阶段系统。
  它对数据平台、训练调度、在线生成、评测和安全审核都提出更高要求。
};

\end{tikzpicture}
\caption{预训练、监督指令微调与 RLHF 的典型流程}
\label{fig:pretrain-sft-rlhf-pipeline}
\end{figure}

简化的偏好模型训练伪代码如下：

\begin{lstlisting}[language=Python,caption={Reward Model 偏好训练的简化伪代码},label={lst:reward-model-training}]
import torch
import torch.nn.functional as F

def reward_model_loss(reward_model, batch):
"""
batch contains:
prompt_ids
chosen_response_ids
rejected_response_ids
"""
chosen_input = concat_prompt_response(
batch["prompt_ids"],
batch["chosen_response_ids"],
)
rejected_input = concat_prompt_response(
batch["prompt_ids"],
batch["rejected_response_ids"],
)

r_chosen = reward_model(chosen_input)      # [batch]
r_rejected = reward_model(rejected_input)  # [batch]

# Bradley-Terry style pairwise preference loss:
# maximize log sigmoid(r_chosen - r_rejected)
loss = -F.logsigmoid(r_chosen - r_rejected).mean()
return loss

\end{lstlisting}

RLHF 的 Infra 难点包括：

\begin{itemize}
\item 同时维护 policy model、reference model、reward model、value model；
\item 生成阶段与训练阶段交替，GPU 利用率不如纯预训练稳定；
\item PPO batch 中样本长度差异大，padding 和 KV Cache 管理复杂；
\item reward hacking、KL 爆炸和训练不稳定需要实时监控；
\item 人类偏好数据采集、审核、版本管理和质量控制成为数据平台问题。
\end{itemize}

近年来，DPO 等直接偏好优化方法试图简化 RLHF 流程，避免显式训练和在线使用 reward model 或 PPO。无论采用哪种方法，对齐训练都说明了一个事实：LLM 的训练已经不再是单一 loss 的最优化，而是数据、偏好、安全、系统吞吐和评测闭环的综合工程。

\section{开源 LLM 生态}
\label{sec:open-source-llm-ecosystem}

开源模型生态极大改变了 AI Infra 的实践方式。过去，只有少数组织能够训练和部署大模型；随着 LLaMA、Mistral、Mixtral、Qwen、DeepSeek 等模型族出现，越来越多团队可以在本地或私有云中进行推理部署、领域微调、量化压缩和服务优化。

这里的“开源”需要谨慎理解。有些模型开放权重但不开放训练数据；有些允许商业使用但有附加限制；有些开放代码和权重但不满足传统开源定义。对 AI Infra 工程师而言，核心影响在于：\textbf{模型权重可获得后，训练与推理系统的优化空间被显著打开。}

\subsection{LLaMA 系列}
\label{subsec:llama-family}

LLaMA 系列推动了开源 LLM 生态的快速发展。它采用现代 Decoder-only Transformer 设计，并在后续版本中逐渐形成如下常见特征：

\begin{itemize}
\item Pre-Norm 结构；
\item RMSNorm；
\item RoPE；
\item SwiGLU；
\item GQA 或改进的 attention 变体；
\item 面向 base model 与 chat model 的不同训练阶段。
\end{itemize}

LLaMA 对 Infra 的意义非常大。模型权重可下载后，工程团队可以真正进行端到端实验：

\begin{itemize}
\item 使用 vLLM、TensorRT-LLM、TGI、SGLang 等 serving 框架部署；
\item 使用 LoRA、QLoRA、FSDP、DeepSpeed 进行微调；
\item 使用 GPTQ、AWQ、SmoothQuant、FP8/INT8/INT4 等量化方法降低成本；
\item 在私有数据上构建企业级 RAG、Agent、代码助手和领域模型；
\item 通过 profiler 分析真实 kernel、KV Cache、batching 和通信瓶颈。
\end{itemize}

换句话说，LLaMA 系列让大模型从“只能调用 API 的黑盒”变成“可以被系统工程师拆解、优化和部署的基础设施组件”。

\subsection{Mistral / Mixtral}
\label{subsec:mistral-mixtral}

Mistral 7B 的重要特点之一，是在相对较小参数规模下追求高效率。它采用 GQA 以提升推理效率，并使用滑动窗口注意力来降低长序列成本。滑动窗口注意力的基本思想是：每个 token 不必关注所有历史 token，而只关注最近窗口内的 token。

若窗口大小为 $W$，标准 attention 对每个位置关注 $S$ 个历史位置，复杂度约为：
\[
O(S^2).
\]
滑动窗口 attention 则约为：
\[
O(SW),
\]
当 $W\ll S$ 时显著降低成本。

Mixtral 则代表了 Sparse Mixture-of-Experts，简称 MoE 的方向。它不是让每个 token 经过同一个 FFN，而是设置多个 expert，并由 router 为每个 token 选择少数 expert。若每层有 $E$ 个 expert，每个 token 选择 top-$k$ 个 expert，则每个 token 的激活参数只是全部 expert 的一小部分。

MoE 的直觉是：
\[
\text{总参数量大}
\quad
\text{但每 token 计算量相对小}.
\]
这使模型可以拥有更高容量，同时控制单 token FLOPs。

\begin{figure}[H]
\centering
\begin{tikzpicture}[
font=\small,
token/.style={draw, rounded corners=2pt, minimum width=0.8cm, minimum height=0.48cm, align=center, fill=blue!6},
router/.style={draw, rounded corners=3pt, minimum width=1.7cm, minimum height=0.7cm, align=center, fill=orange!14},
expert/.style={draw, rounded corners=3pt, minimum width=1.35cm, minimum height=0.65cm, align=center, fill=green!10},
arrow/.style={-{Latex[length=2.2mm]}, thick},
faint/.style={-{Latex[length=2.0mm]}, thick, gray!45}
]

\node[font=\bfseries] at (0,3.0) {MoE：Router 为每个 token 选择少数 Expert};

\foreach \i/\x in {1/-4.8,2/-3.8,3/-2.8,4/-1.8} {
  \node[token] (t\i) at (\x,1.3) {$t_\i$};
}

\node[router] (router) at (-3.3,0.1) {Router\\Top-$k$};

\foreach \i in {1,...,4} {
  \draw[arrow] (t\i) -- (router);
}

\foreach \i/\x in {1/0.3,2/1.7,3/3.1,4/4.5} {
  \node[expert] (e\i) at (\x,0.85) {Expert \i};
}

\draw[arrow] (router) -- (e1);
\draw[arrow] (router) -- (e3);
\draw[faint] (router) -- (e2);
\draw[faint] (router) -- (e4);

\node[expert, fill=purple!10, minimum width=1.8cm] (combine) at (2.4,-0.75) {Combine};
\draw[arrow] (e1) -- (combine);
\draw[arrow] (e3) -- (combine);

\node[align=left, text width=10.3cm] at (0,-2.2) {
  MoE 提高模型容量，但引入路由、负载均衡和 all-to-all 通信。
  对训练集群而言，MoE 不只是模型结构变化，更是通信模式和调度问题。
};

\end{tikzpicture}
\caption{Mixture-of-Experts 中的 Router 与 Expert 分发}
\label{fig:moe-router-experts}
\end{figure}

MoE 对 AI Infra 的挑战包括：

\begin{itemize}
\item token 在不同 expert 之间动态分发，导致负载不均；
\item expert parallel 通常需要 all-to-all 通信；
\item 每个 expert 的 batch 形状动态变化，不利于稳定 GEMM；
\item 推理时需要优化路由、缓存和专家并行部署；
\item 多机训练中，网络拓扑对 MoE 性能影响非常大。
\end{itemize}

\subsection{Qwen、DeepSeek 等模型族}
\label{subsec:qwen-deepseek-and-others}

Qwen、DeepSeek 等模型族进一步推动了开源 LLM 在多语言、代码、数学、推理、长上下文和 MoE 方向的发展。不同模型族的公开重点不同：有的强调通用能力，有的强调代码，有的强调数学推理，有的强调训练效率，有的强调长上下文或开源部署生态。

对 Infra 工程师而言，模型族差异需要转化为系统决策：

\begin{itemize}
\item 模型是 dense 还是 MoE；
\item 使用 MHA、GQA 还是 MLA 等 attention 变体；
\item 上下文长度是多少；
\item tokenizer 词表大小和特殊 token 规则是什么；
\item 是否支持 FP8、INT8、INT4 量化；
\item 是否有官方推理框架或推荐部署参数；
\item license 是否允许目标业务使用。
\end{itemize}

同样是“70B 模型”，不同架构的部署成本可能完全不同。GQA 会影响 KV Cache；MoE 会影响专家并行；长上下文会影响 prefill 和 cache；大词表会影响 logits 和 cross entropy；多模态输入会引入视觉 encoder、投影层和不同的 batch 组织方式。

\subsection{开源模型对 AI Infra 的推动}
\label{subsec:open-models-push-ai-infra}

开源模型生态使 AI Infra 从少数大厂内部能力，变成广泛工程实践。它推动了以下方向：

\begin{itemize}
\item \textbf{推理框架竞争}：vLLM、TensorRT-LLM、TGI、llama.cpp、SGLang 等不断优化吞吐与延迟；
\item \textbf{量化方法普及}：INT8、INT4、AWQ、GPTQ、GGUF 等让普通 GPU 甚至 CPU 也能运行模型；
\item \textbf{微调工具成熟}：LoRA、QLoRA、PEFT、Axolotl、LLaMA-Factory 等降低领域适配门槛；
\item \textbf{评测体系扩展}：从通用 benchmark 扩展到业务场景、对齐、安全、幻觉和成本评估；
\item \textbf{私有化部署需求增长}：企业希望在自己的数据和合规环境中运行模型；
\item \textbf{硬件适配多样化}：NVIDIA GPU、AMD GPU、国产加速卡、Apple Silicon、CPU 后端都有部署需求。
\end{itemize}

这意味着 AI Infra 工程师需要具备跨层能力：读懂模型结构，理解 kernel 和显存，掌握服务调度，能做量化和并行配置，还要能处理数据、评测和安全上线。

\section{模型规模定律}
\label{sec:scaling-laws}

大语言模型的发展并不是随机堆参数。过去几年，一个核心经验逐渐清晰：模型能力与参数量、数据量和计算量之间存在可预测的 scaling 关系。Scaling Law 并不保证每一次扩展都成功，但它为训练规划提供了重要指导。

\subsection{参数量、数据量与计算量}
\label{subsec:params-data-compute}

预训练一个自回归 Transformer 的主要资源可以概括为三要素：

\begin{itemize}
\item \textbf{参数量 $N$}：模型容量；
\item \textbf{训练 token 数 $D$}：数据规模；
\item \textbf{训练计算量 $C$}：通常与 $N$ 和 $D$ 的乘积相关。
\end{itemize}

粗略地，dense Transformer 的训练 FLOPs 可估为：
\[
C\approx 6ND.
\]
这里的常数 $6$ 来自 forward、backward 和梯度计算的近似开销，具体值会随架构、激活重计算、并行策略和实现细节变化。

在给定计算预算 $C$ 的条件下，训练更大模型意味着可训练 token 数减少；训练更多 token 意味着模型参数量可能需要变小。早期一些大模型更强调参数量扩展，而后来的 compute-optimal 思路强调：模型不能只大，还要用足够多的数据训练。

\begin{figure}[H]
\centering
\begin{tikzpicture}[
font=\small,
vertex/.style={circle, draw, minimum size=1.05cm, align=center},
arrow/.style={-{Latex[length=2.3mm]}, thick},
label/.style={align=center, text width=3.1cm}
]

\node[font=\bfseries] at (0,3.2) {大模型训练的三角关系：参数、数据与算力};

\node[vertex, fill=blue!8] (params) at (-3.7,0.5) {$N$\\参数量};
\node[vertex, fill=green!10] (data) at (3.7,0.5) {$D$\\数据量};
\node[vertex, fill=orange!14] (compute) at (0,-2.2) {$C$\\计算量};

\draw[arrow] (params) -- node[above] {容量与泛化} (data);
\draw[arrow] (data) -- node[right] {训练 token} (compute);
\draw[arrow] (compute) -- node[left] {GPU-days} (params);

\node[label] at (-4.4,-1.45) {
  模型太小：\\
  容量不足
};
\node[label] at (4.4,-1.45) {
  数据太少：\\
  undertrained
};
\node[label] at (0,1.95) {
  算力预算决定\\
  可行组合
};

\end{tikzpicture}
\caption{参数量、数据量与计算量的三角权衡}
\label{fig:params-data-compute-triangle}
\end{figure}

\subsection{Chinchilla Scaling Law}
\label{subsec:chinchilla-scaling-law}

Chinchilla Scaling Law 的核心启发是：在固定计算预算下，许多早期大模型参数很多但训练 token 相对不足。更 compute-optimal 的策略是同时增加模型大小和训练数据量。

直观地说，若模型参数翻倍，训练 token 数也应该相应增加，而不是只让模型变大。否则模型可能没有被充分训练，推理成本却很高。

可以把训练不足的模型理解为：
\[
\text{参数容量很大}
\quad
\text{但数据没有充分填充这些容量}.
\]
这种模型可能在训练曲线上还没有充分收敛，却已经因为参数大而带来高推理成本。

对 AI Infra 来说，Chinchilla 风格的结论意味着：

\begin{itemize}
\item 训练计划不能只问“训练多大模型”，还要问“训练多少 token”；
\item 数据管线、去重、质量过滤和混合比例与算力同等重要；
\item compute-optimal 模型可能参数更小但训练更充分，推理成本更低；
\item 训练集群规划要围绕 token budget、step time 和 MFU 计算；
\item checkpoint、日志、恢复和评测要支持长时间大 token 训练。
\end{itemize}

\subsection{compute-optimal training}
\label{subsec:compute-optimal-training}

设训练计算预算为 $C$，参数量为 $N$，训练 token 数为 $D$。粗略有：
\[
C\propto ND.
\]
compute-optimal training 的目标是在给定 $C$ 下选择合适的 $N$ 和 $D$，使最终 loss 最低。

如果 $N$ 太大、$D$ 太小，模型 undertrained；如果 $N$ 太小、$D$ 太大，模型容量不足。最优点在二者之间。

\begin{figure}[H]
\centering
\begin{tikzpicture}[
font=\small,
axis/.style={-{Latex[length=2.3mm]}, thick},
curve/.style={thick},
point/.style={circle, fill=red!70, inner sep=2pt}
]

\node[font=\bfseries] at (0,3.0) {Compute-Optimal Training 的直觉};

\draw[axis] (-5.4,-1.2) -- (5.4,-1.2) node[right] {model size $N$};
\draw[axis] (-5.4,-1.2) -- (-5.4,2.2) node[above] {final loss};

\draw[curve, blue!75]
  plot[smooth] coordinates {
    (-5.0,1.7) (-4.0,1.05) (-3.0,0.45) (-2.0,0.05)
    (-1.0,-0.15) (0.0,-0.25) (1.0,-0.18) (2.0,0.02)
    (3.0,0.45) (4.0,1.05) (5.0,1.65)
  };

\node[point] at (0.0,-0.25) {};
\node[align=center] at (0.0,-0.72) {compute-optimal\\region};

\node[align=center] at (-3.7,1.75) {模型太小\\容量不足};
\node[align=center] at (3.9,1.75) {模型太大\\数据不足};

\node[align=left, text width=10.2cm] at (0,-2.35) {
  在固定计算预算下，并不是参数越多越好。参数量、训练 token 数和数据质量需要共同规划。
  这也是大模型训练从“堆规模”走向“系统化配比”的关键。
};

\end{tikzpicture}
\caption{给定计算预算下的 compute-optimal 训练直觉}
\label{fig:compute-optimal-training}
\end{figure}

在真实项目中，compute-optimal 不只是公式问题，还要考虑：

\begin{itemize}
\item 数据是否足够高质量；
\item 是否存在重复数据污染 benchmark；
\item 训练集群是否能稳定运行足够长时间；
\item 模型规模是否满足推理成本约束；
\item 下游任务是否更需要知识、推理、代码、多语言或长上下文；
\item 是否计划后续 SFT、RLHF、DPO 或领域继续预训练。
\end{itemize}

\subsection{为什么 Infra 决定模型边界}
\label{subsec:why-infra-defines-model-boundary}

模型能力看起来由算法和数据决定，但在大规模训练中，Infra 决定可行边界。一个团队能训练多大模型，取决于：

\begin{itemize}
\item 有多少 GPU；
\item 单卡 HBM 容量和带宽；
\item 节点内 NVLink/NVSwitch 拓扑；
\item 节点间 RDMA 网络带宽与稳定性；
\item 存储系统能否持续喂数据；
\item 分布式训练框架是否能高效扩展；
\item checkpoint 是否能在故障后恢复；
\item 训练监控是否能及时发现 loss spike、NaN 和 straggler；
\item 调度系统是否能提供长时间稳定资源。
\end{itemize}

同样的模型设计，在不同 Infra 上可能有完全不同结果。一个理论上可行的训练配置，如果网络通信无法支撑 tensor parallel，或者存储吞吐无法支撑数据加载，或者 checkpoint 写入拖慢训练，就无法成为稳定产品。

\begin{table}[H]
\centering
\small
\caption{模型规模扩展对 Infra 的压力}
\label{tab:model-scaling-infra-pressure}
\begin{tabular}{p{0.22\textwidth}p{0.34\textwidth}p{0.34\textwidth}}
\toprule
\textbf{扩展方向} & \textbf{直接收益} & \textbf{Infra 压力} \
\midrule
增大参数量
& 提升模型容量和潜在能力
& 显存、优化器状态、模型并行、checkpoint 体积增加。 \
\midrule
增加训练 token
& 更充分训练，降低 undertraining
& 数据管线、训练时长、稳定性、存储吞吐要求上升。 \
\midrule
增加上下文长度
& 支持长文档、多轮对话和复杂任务
& Attention 成本、KV Cache、激活显存和调度复杂度上升。 \
\midrule
使用 MoE
& 增加总参数容量，控制 active FLOPs
& Router、负载均衡、all-to-all 通信和 expert placement 更复杂。 \
\midrule
多模态扩展
& 支持图像、音频、视频等输入
& 数据处理、encoder 计算、跨模态对齐和 serving pipeline 更复杂。 \
\bottomrule
\end{tabular}
\end{table}

因此，本书后续几篇会围绕这些 Infra 边界展开：分布式训练解决大模型放不下和训不动的问题；推理架构解决高吞吐、低延迟、低成本服务的问题；硬件与网络章节解决多 GPU、多节点和集群调度的问题。

\section{从模型范式到训练流水线}
\label{sec:from-model-paradigm-to-training-pipeline}

为了把本章内容落到工程上，我们可以把一个现代 Chat LLM 的构建过程抽象为如下流水线：

\begin{enumerate}
\item 数据收集与清洗；
\item tokenizer 训练或选择；
\item base model 预训练；
\item 中间 checkpoint 评测；
\item 指令微调 SFT；
\item 偏好数据收集；
\item reward model 或直接偏好优化；
\item 安全对齐与红队测试；
\item 量化、蒸馏或格式转换；
\item serving 部署与在线监控。
\end{enumerate}

\begin{figure}[H]
\centering
\begin{tikzpicture}[
font=\small,
stage/.style={draw, rounded corners=3pt, minimum width=1.8cm, minimum height=0.68cm, align=center, fill=blue!6},
sys/.style={draw, rounded corners=3pt, minimum width=1.8cm, minimum height=0.68cm, align=center, fill=orange!12},
arrow/.style={-{Latex[length=2.1mm]}, thick}
]

\node[font=\bfseries] at (0,3.1) {现代 Chat LLM 构建流水线};

\node[sys] (data) at (-5.2,1.6) {Data\\Curation};
\node[sys] (tok) at (-3.0,1.6) {Tokenizer};
\node[stage] (pre) at (-0.8,1.6) {Pretrain};
\node[sys] (eval1) at (1.4,1.6) {Eval\\Checkpoint};
\node[stage] (sft) at (3.6,1.6) {SFT};

\draw[arrow] (data) -- (tok);
\draw[arrow] (tok) -- (pre);
\draw[arrow] (pre) -- (eval1);
\draw[arrow] (eval1) -- (sft);

\node[sys] (pref) at (-5.2,0.0) {Preference\\Data};
\node[stage] (align) at (-3.0,0.0) {RLHF / DPO};
\node[sys] (safe) at (-0.8,0.0) {Safety\\Eval};
\node[sys] (quant) at (1.4,0.0) {Quantize\\Convert};
\node[stage] (serve) at (3.6,0.0) {Serving};

\draw[arrow] (sft.south) .. controls (3.6,0.75) and (-3.0,0.75) .. (align.north);
\draw[arrow] (pref) -- (align);
\draw[arrow] (align) -- (safe);
\draw[arrow] (safe) -- (quant);
\draw[arrow] (quant) -- (serve);

\node[align=left, text width=10.5cm] at (0,-1.6) {
  每个阶段都有不同 Infra 需求：预训练重吞吐和稳定性，SFT 重数据质量，
  RLHF/DPO 重多模型训练与偏好数据，Serving 重延迟、吞吐、KV Cache 和成本。
};

\end{tikzpicture}
\caption{现代 Chat LLM 的端到端构建流水线}
\label{fig:modern-chat-llm-pipeline}
\end{figure}

下面给出一个极简训练阶段配置伪代码，用来说明不同阶段的资源和目标差异：

\begin{lstlisting}[language=Python,caption={不同 LLM 训练阶段的配置抽象},label={lst:llm-training-stages-config}]
training_plan = [
{
"stage": "pretraining",
"objective": "next_token_prediction",
"data": "large_scale_web_code_math_corpus",
"parallelism": ["data_parallel", "tensor_parallel", "pipeline_parallel"],
"optimizer": "AdamW",
"main_bottleneck": ["GPU_flops", "activation_memory", "network_collectives"],
},
{
"stage": "supervised_finetuning",
"objective": "instruction_following",
"data": "high_quality_prompt_response_pairs",
"parallelism": ["data_parallel", "fsdp_or_zero"],
"optimizer": "AdamW_or_LoRA_optimizer",
"main_bottleneck": ["data_quality", "sequence_packing", "overfitting"],
},
{
"stage": "preference_optimization",
"objective": "align_with_human_preference",
"data": "chosen_rejected_pairs",
"parallelism": ["data_parallel", "model_sharding"],
"optimizer": "PPO_DPO_or_related_method",
"main_bottleneck": ["multi_model_memory", "generation_throughput", "stability"],
},
{
"stage": "serving",
"objective": "low_latency_high_throughput_generation",
"data": "online_requests",
"parallelism": ["tensor_parallel", "pipeline_parallel", "continuous_batching"],
"optimizer": None,
"main_bottleneck": ["KV_cache", "scheduler", "memory_bandwidth", "tail_latency"],
},
\]
\end{lstlisting}

这段伪代码强调：模型生命周期中的不同阶段，对 Infra 的瓶颈完全不同。预训练关心 GPU 利用率和多节点通信；SFT 关心数据格式、loss mask 和样本质量；偏好优化关心多模型显存和生成吞吐；serving 则关心 KV Cache、batching、延迟和成本。

\section{本章小结}
\label{sec:chapter05-summary}

本章从语言模型任务定义出发，梳理了大语言模型的主要发展脉络和训练范式。

\begin{itemize}
\item \textbf{N-gram 模型} 使用有限上下文和离散计数，简单但无法有效建模长距离依赖和语义泛化。
\item \textbf{神经语言模型} 将 token 映射到连续向量空间，把语言建模转化为可由 GPU 加速的张量计算。
\item \textbf{自回归语言建模} 通过 next-token prediction 学习序列概率，是 Decoder-only LLM 的核心任务。
\item \textbf{MLM 与 Causal LM} 代表不同建模范式：前者适合理解，后者适合生成。
\item \textbf{GPT 系列} 推动了预训练、规模扩展、in-context learning、instruction tuning 和 RLHF 等关键范式。
\item \textbf{Instruction Tuning} 通过高质量指令数据改变模型输出分布，让基础模型更稳定地遵循用户意图。
\item \textbf{RLHF 与偏好优化} 把训练扩展为多模型、多阶段系统，连接了人类偏好、奖励建模和策略优化。
\item \textbf{开源模型生态} 让模型权重成为可部署、可微调、可量化、可优化的基础设施对象，推动了 serving、量化和微调工具链成熟。
\item \textbf{Scaling Law} 告诉我们模型能力与参数量、数据量和计算量有关；compute-optimal training 要在给定算力下平衡模型规模与训练 token。
\item \textbf{Infra 决定模型边界}：GPU、网络、存储、并行框架、checkpoint、调度和监控共同决定一个模型能否被训练和服务。
\end{itemize}

下一章将进入具体架构对比：GPT、LLaMA 与现代 Decoder-only 架构。我们会把本章的历史范式落到模型结构细节上，系统比较 token embedding、causal attention、RMSNorm、RoPE、SwiGLU、GQA、MoE 以及 lm head 等组件，并分析这些结构如何改变训练和推理基础设施。

\chapter{GPT、LLaMA 与现代 Decoder-only 架构}
\label{chap:gpt-llama-arch}

\section{Decoder-only LLM 总览}
\label{sec:decoder-only-llm-overview}

前两章分别回答了两个问题：Transformer 为什么适合大规模并行训练，大语言模型范式又是如何从语言建模演进到 Chat LM 的。本章进一步落到具体架构：现代 Decoder-only LLM 到底由哪些模块组成？GPT 风格模型与 LLaMA 风格模型有何差异？为什么 RMSNorm、RoPE、SwiGLU、GQA、MoE 这些组件会成为现代模型设计中的高频词？这些架构选择又如何影响训练和推理基础设施？

一个 Decoder-only LLM 可以被抽象为如下函数：
\[
f_{\theta}:
(x_1,x_2,\ldots,x_S)
\rightarrow
p(x_{S+1}\mid x_{\le S}).
\]
它接收一段 token 序列，输出下一个 token 的概率分布。训练时，它对序列中每个位置并行做 next-token prediction；推理时，它不断把新生成的 token 追加到上下文中，逐步生成完整回答。

从结构上看，现代 Decoder-only LLM 通常包含四个主要部分：

\begin{itemize}
\item \textbf{Token Embedding}：把离散 token id 映射为连续向量；
\item \textbf{Transformer Decoder Blocks}：堆叠多层 causal self-attention 与 FFN；
\item \textbf{Final Norm}：在输出 logits 前对 hidden state 做归一化；
\item \textbf{LM Head}：把 hidden state 投影到词表维度，得到 logits。
\end{itemize}

其整体计算可以写作：
\[
\mat{X}*0=\operatorname{Embed}(\mat{T})+\operatorname{PosEnc},
\]
\[
\mat{X}*{\ell+1}
================

\operatorname{DecoderBlock}*{\ell}(\mat{X}*{\ell}),
\quad
\ell=0,1,\ldots,L-1,
\]
\[
\mat{H}=\operatorname{Norm}(\mat{X}*L),
\]
\[
\mat{Z}=\mat{H}\mat{W}*{\mathrm{lm}},
\]
其中：
\[
\mat{Z}\in\mathbb{R}^{B\times S\times V}
\]
是 logits，$V$ 是词表大小。

\subsection{token embedding}
\label{subsec:token-embedding}

Tokenizer 将文本切分为 token，并把每个 token 映射为整数 id。设词表大小为 $V$，隐藏维度为 $H$，embedding table 为：
\[
\mat{E}\in\mathbb{R}^{V\times H}.
\]
输入 token ids：
\[
\mat{T}\in\mathbb{N}^{B\times S}
\]
经过查表得到：
\[
\tensor{X}*{b,s,:}=\mat{E}*{T_{b,s},:}.
\]
因此：
\[
\tensor{X}\in\mathbb{R}^{B\times S\times H}.
\]

Embedding 查表在数学上可以看作 one-hot 向量乘矩阵：
\[
\vect{x}_{b,s}
==============

\vect{o}*{T*{b,s}}^{\mathsf{T}}\mat{E}.
\]
但工程上不会真的构造 one-hot，因为这样会产生巨大稀疏矩阵。框架通常使用 gather 操作直接读取 embedding table 的对应行。

从 Infra 角度看，embedding 层有几个值得注意的特点：

\begin{itemize}
\item \textbf{参数量与词表大小成正比}：embedding 参数量为 $VH$；
\item \textbf{访问模式偏随机}：不同 token id 对应 embedding table 中不同位置，缓存局部性不如 GEMM；
\item \textbf{可与 LM Head 共享权重}：某些模型使用 tied embedding，即 $\mat{W}_{\mathrm{lm}}=\mat{E}^{\mathsf{T}}$；
\item \textbf{分布式训练中可做 vocabulary parallel}：当 $V$ 很大时，可将词表维切分到多个 GPU。
\end{itemize}

例如，当 $V=128000$，$H=8192$ 时，仅 embedding 参数量就是：
\[
128000\times 8192\approx 1.05\times 10^9.
\]
这约为 $1$B 参数。对于大词表多语言模型，embedding 和 lm head 的系统成本不可忽视。

\subsection{transformer blocks}
\label{subsec:transformer-blocks}

Decoder-only 模型的主体是堆叠的 Transformer blocks。一个现代 Pre-Norm block 通常写为：
\[
\mat{U}_{\ell}
==============

\mat{X}*{\ell}
+
\operatorname{CausalSelfAttention}
\left(
\operatorname{Norm}*1(\mat{X}*{\ell})
\right),
\]
\[
\mat{X}*{\ell+1}
================

\mat{U}_{\ell}
+
\operatorname{FFN}
\left(
\operatorname{Norm}*2(\mat{U}*{\ell})
\right).
\]

这里的 Norm 可以是 LayerNorm 或 RMSNorm；Attention 可以是 MHA、MQA、GQA 或更复杂的注意力变体；FFN 可以是 GELU MLP、SwiGLU MLP 或 MoE FFN。不同模型族的差异，往往就体现在这些选择上。

Transformer blocks 的工程成本主要来自：

\begin{itemize}
\item Attention 中的 QKV projection、attention score、softmax、weighted sum 和 output projection；
\item FFN 中的 up/gate/down projection；
\item Norm、residual add、activation 等 memory-bound 小算子；
\item 训练时保存的激活；
\item 分布式训练中的 tensor parallel、sequence parallel、pipeline parallel 和数据并行通信。
\end{itemize}

\subsection{lm head}
\label{subsec:lm-head}

LM Head 将最终 hidden state 投影到词表维度：
\[
\mat{Z}_{b,s,:}
===============

\mat{H}*{b,s,:}\mat{W}*{\mathrm{lm}},
\]
其中：
\[
\mat{W}_{\mathrm{lm}}\in\mathbb{R}^{H\times V}.
\]
logits：
\[
\mat{Z}\in\mathbb{R}^{B\times S\times V}
\]
再经过 softmax 得到 token 概率。

LM Head 的参数量为：
\[
HV.
\]
如果使用 tied embedding，则：
\[
\mat{W}_{\mathrm{lm}}=\mat{E}^{\mathsf{T}},
\]
可以减少参数量，并在某些情况下提升泛化。

但 LM Head 的主要问题不只是参数量，而是 logits 张量很大。训练中，如果 $B=8$、$S=4096$、$V=128000$，logits 元素数为：
\[
8\times 4096\times 128000\approx 4.19\times 10^9.
\]
即使使用 BF16，也超过 $8$GB。实际系统通常会使用 fused cross entropy、vocabulary parallelism、logits chunking 或 selective logits 计算来降低显存峰值。

\subsection{logits 与采样}
\label{subsec:logits-and-sampling}

推理时，我们通常只关心最后一个位置的 logits：
\[
\vect{z}=\mat{Z}_{B,S,:}\in\mathbb{R}^{V}.
\]
通过 softmax 得到概率：
\[
p_i=\frac{\exp(z_i/T)}{\sum_j\exp(z_j/T)},
\]
其中 $T$ 是 temperature。之后可以使用 greedy、top-$k$、top-$p$、min-$p$、typical sampling 等策略采样下一个 token。

采样策略不是模型结构的一部分，但它强烈影响服务行为。低 temperature 更确定，高 temperature 更随机；top-$p$ 限制候选 token 集合，减少长尾低概率 token 的影响；repetition penalty 可以缓解重复输出。

从 serving infra 角度看，采样阶段通常不是主要 FLOPs 来源，但会影响：

\begin{itemize}
\item 每个请求生成长度；
\item stop condition；
\item batch 内请求退出时机；
\item tail latency；
\item continuous batching 的动态调度。
\end{itemize}

\begin{figure}[H]
\centering
\begin{tikzpicture}[
font=\small,
box/.style={draw, rounded corners=3pt, minimum width=2.05cm, minimum height=0.75cm, align=center},
block/.style={draw, rounded corners=3pt, minimum width=2.2cm, minimum height=0.68cm, align=center, fill=blue!6},
arrow/.style={-{Latex[length=2.3mm]}, thick}
]

\node[font=\bfseries] at (0,4.0) {Decoder-only LLM 整体结构};

\node[box, fill=orange!12] (ids) at (-5.1,2.2) {Token IDs\\$B\times S$};
\node[box, fill=green!10] (emb) at (-2.6,2.2) {Embedding\\$B,S,H$};

\node[block] (b1) at (0.0,2.2) {Decoder\\Block 1};
\node[block] (b2) at (0.0,1.25) {Decoder\\Block 2};
\node[block] (dots) at (0.0,0.3) {$\cdots$};
\node[block] (bL) at (0.0,-0.65) {Decoder\\Block $L$};

\node[box, fill=purple!10] (norm) at (2.7,-0.65) {Final Norm};
\node[box, fill=yellow!16] (lm) at (5.1,-0.65) {LM Head\\$H\rightarrow V$};
\node[box, fill=red!10] (logits) at (5.1,-1.75) {Logits\\$B,S,V$};

\draw[arrow] (ids) -- (emb);
\draw[arrow] (emb) -- (b1);
\draw[arrow] (b1) -- (b2);
\draw[arrow] (b2) -- (dots);
\draw[arrow] (dots) -- (bL);
\draw[arrow] (bL) -- (norm);
\draw[arrow] (norm) -- (lm);
\draw[arrow] (lm) -- (logits);

\node[align=left, text width=10.7cm] at (0,-3.05) {
  Decoder-only LLM 的主体是重复堆叠的 causal decoder block。
  训练时对所有位置计算 logits 与 loss；推理 decode 阶段通常只取最后位置 logits，并将采样 token 追加到上下文。
};

\end{tikzpicture}
\caption{Decoder-only LLM 的端到端结构}
\label{fig:decoder-only-llm-overall}
\end{figure}

\section{GPT 架构细节}
\label{sec:gpt-architecture-details}

GPT 风格架构可以看作现代 Decoder-only LLM 的起点。虽然不同版本细节不同，但核心思想相对稳定：token embedding、位置 embedding、堆叠 causal self-attention block、最终 lm head，以及 next-token prediction 训练目标。

\subsection{learned positional embedding}
\label{subsec:gpt-learned-positional-embedding}

早期 GPT 模型使用 learned absolute positional embedding。设最大上下文长度为 $S_{\max}$，位置 embedding 表为：
\[
\mat{P}\in\mathbb{R}^{S_{\max}\times H}.
\]
第 $s$ 个位置的输入为：
\[
\vect{x}*s=\mat{E}*{t_s,:}+\mat{P}_{s,:}.
\]

这种方法简单直接，模型可以学习每个绝对位置的表示。但它有明显限制：如果推理时序列长度超过训练时最大位置，模型没有对应的位置 embedding。即便可以通过插值或扩展，也不一定保持稳定。

从系统角度看，learned positional embedding 成本很小，参数量为：
\[
S_{\max}H.
\]
相比 Transformer block 中的 $O(H^2)$ 参数，它通常不是主要参数来源。但它影响上下文长度扩展能力。后来的 LLaMA 风格模型更多使用 RoPE，因为它更自然地将相对位置信息注入 Q/K，并在长上下文扩展中有更丰富的调整方式。

\subsection{causal self-attention}
\label{subsec:gpt-causal-self-attention}

GPT 的核心是 causal self-attention：
\[
\mat{O}
=======

\softmax
\left(
\frac{\mat{Q}\mat{K}^{\mathsf{T}}+\mat{M}}{\sqrt{d_h}}
\right)
\mat{V},
\]
其中 mask $\mat{M}$ 满足：
\[
M_{ij}=
\begin{cases}
0, & j\le i,\
-\infty, & j>i.
\end{cases}
\]
它保证第 $i$ 个位置不能看到未来 token。

在训练中，causal mask 允许所有位置并行计算，同时保持自回归约束。这个性质非常关键：如果训练也必须逐 token 递推，那么大模型训练成本会高得难以接受。

在推理中，causal self-attention 以另一种方式工作。假设已经生成了 $S$ 个 token，并缓存了历史 K/V：
\[
\mat{K}*{cache}\in\mathbb{R}^{S\times d_h},
\quad
\mat{V}*{cache}\in\mathbb{R}^{S\times d_h}.
\]
生成第 $S+1$ 个 token 时，只需要计算新 token 的 query：
\[
\vect{q}*{new},
\]
并与历史 key 做 attention：
\[
\vect{o}*{new}
==============

\softmax
\left(
\frac{\vect{q}*{new}\mat{K}*{cache}^{\mathsf{T}}}{\sqrt{d_h}}
\right)
\mat{V}_{cache}.
\]
这就是 KV Cache 的基础。它避免了每一步重复计算历史 token 的 K/V，但引入了显存和带宽压力。

\subsection{residual stream}
\label{subsec:gpt-residual-stream}

GPT 架构中，residual connection 把每一层的输出加回主干 hidden state：
\[
\mat{X}_{\ell+1}
================

\mat{X}*{\ell}
+
F*{\ell}(\mat{X}_{\ell}).
\]
这个贯穿层间的主干表示常被称为 residual stream。Attention 和 MLP 每层都向 residual stream 写入增量信息。

从训练稳定性角度看，residual stream 提供了梯度传播的捷径。若没有 residual connection，深层网络会更难训练。若有 residual connection，梯度可以沿恒等路径传播：
\[
\frac{\partial \mat{X}*{\ell+1}}{\partial \mat{X}*{\ell}}
=========================================================

\mat{I}
+
\frac{\partial F_{\ell}}{\partial \mat{X}_{\ell}}.
\]
其中 $\mat{I}$ 保证了一个直接通路。

从系统角度看，residual add 本身是 elementwise 操作，计算量不大，但会读写完整激活张量：
\[
B\times S\times H.
\]
因此，高性能实现中常将 residual add 与 LayerNorm/RMSNorm 融合，减少 HBM 往返。

\subsection{GPT 风格 block 的工程实现}
\label{subsec:gpt-style-block-engineering}

GPT 风格 block 可抽象为：
\[
\mat{X}'=\mat{X}+\operatorname{Attention}(\operatorname{Norm}(\mat{X})),
\]
\[
\mat{Y}=\mat{X}'+\operatorname{MLP}(\operatorname{Norm}(\mat{X}')).
\]

下面给出一个教学版 GPT block。它使用 LayerNorm、MHA 和 GELU MLP，接近早期 GPT 风格。为突出结构，代码没有使用 FlashAttention、KV Cache 或 fused kernel。

\begin{lstlisting}[language=Python,caption={教学版 GPT 风格 Decoder Block},label={lst:gpt-style-block}]
import math
import torch
import torch.nn as nn
import torch.nn.functional as F

class GPTAttention(nn.Module):
def **init**(self, hidden_size, num_heads):
super().**init**()
assert hidden_size % num_heads == 0
self.hidden_size = hidden_size
self.num_heads = num_heads
self.head_dim = hidden_size // num_heads

    self.qkv = nn.Linear(hidden_size, 3 * hidden_size)
    self.out_proj = nn.Linear(hidden_size, hidden_size)

def forward(self, x):
    B, S, H = x.shape

    qkv = self.qkv(x)
    q, k, v = qkv.chunk(3, dim=-1)

    q = q.view(B, S, self.num_heads, self.head_dim).transpose(1, 2)
    k = k.view(B, S, self.num_heads, self.head_dim).transpose(1, 2)
    v = v.view(B, S, self.num_heads, self.head_dim).transpose(1, 2)

    scores = q @ k.transpose(-2, -1)
    scores = scores / math.sqrt(self.head_dim)

    mask = torch.triu(torch.ones(S, S, device=x.device, dtype=torch.bool), diagonal=1)
    scores = scores.masked_fill(mask, float("-inf"))

    probs = F.softmax(scores, dim=-1)
    out = probs @ v

    out = out.transpose(1, 2).contiguous().view(B, S, H)
    return self.out_proj(out)

class GPTMLP(nn.Module):
def **init**(self, hidden_size, intermediate_size):
super().**init**()
self.fc1 = nn.Linear(hidden_size, intermediate_size)
self.fc2 = nn.Linear(intermediate_size, hidden_size)

def forward(self, x):
    return self.fc2(F.gelu(self.fc1(x)))

class GPTBlock(nn.Module):
def **init**(self, hidden_size, num_heads, intermediate_size):
super().**init**()
self.ln1 = nn.LayerNorm(hidden_size)
self.attn = GPTAttention(hidden_size, num_heads)
self.ln2 = nn.LayerNorm(hidden_size)
self.mlp = GPTMLP(hidden_size, intermediate_size)

def forward(self, x):
    x = x + self.attn(self.ln1(x))
    x = x + self.mlp(self.ln2(x))
    return x

\end{lstlisting}

这段代码的工程问题非常典型：

\begin{itemize}
\item 每次 forward 构造 causal mask；
\item 显式物化 attention score；
\item softmax 和两个 matmul 没有融合；
\item transpose 后调用 \texttt{contiguous()}，可能引入额外拷贝；
\item LayerNorm、residual add、GELU 分散为多个小 kernel；
\item 没有 KV Cache，无法高效自回归 decode。
\end{itemize}

真实 GPT 推理和训练框架会使用更高效实现。例如使用 PyTorch SDPA、FlashAttention、fused bias-GELU、fused LayerNorm、tensor parallel linear，以及更适合推理的 KV Cache layout。

\section{LLaMA 架构细节}
\label{sec:llama-architecture-details}

LLaMA 风格架构代表了现代开源 Decoder-only LLM 的一类重要设计。它继承了 GPT 的 causal decoder 主体，但在归一化、位置编码、FFN 和 attention head 共享上采用了更现代的选择：

\begin{itemize}
\item 使用 RMSNorm 而非标准 LayerNorm；
\item 使用 RoPE 而非 learned absolute positional embedding；
\item 使用 SwiGLU 或类似门控 FFN；
\item 在一些版本中使用 GQA，以降低推理 KV Cache 成本；
\item 通常采用 Pre-Norm 结构。
\end{itemize}

这些变化看似是模型细节，但对训练稳定性、推理显存、kernel 实现和分布式并行都有重要影响。

\subsection{RMSNorm}
\label{subsec:llama-rmsnorm}

RMSNorm 定义为：
\[
\RMSNorm(\vect{x})
==================

\gamma
\frac{\vect{x}}
{\sqrt{
\frac{1}{H}\sum_{i=1}^{H}x_i^2+\epsilon
}}.
\]
它与 LayerNorm 的区别在于不减均值。LayerNorm 为：
\[
\LayerNorm(\vect{x})
====================

\gamma
\frac{\vect{x}-\mu}
{\sqrt{\sigma^2+\epsilon}}
+
\beta.
\]

RMSNorm 的参数通常只有缩放项 $\gamma$，没有 bias。相比 LayerNorm，它计算更简单，少了一次均值计算和减均值步骤。对于每个 token 的 hidden 向量，它需要：

\begin{itemize}
\item 计算平方和 reduction；
\item 求均值；
\item rsqrt；
\item 逐元素缩放。
\end{itemize}

虽然 RMSNorm 参数量很小，但它位于每一层的 attention 和 FFN 前面，调用频繁。若实现不佳，容易在 GPU 时间线上形成大量 memory-bound 小 kernel。因此生产系统常使用 fused RMSNorm kernel。

\begin{lstlisting}[language=Python,caption={RMSNorm 的简洁实现},label={lst:llama-rmsnorm}]
import torch
import torch.nn as nn

class RMSNorm(nn.Module):
def **init**(self, hidden_size, eps=1e-6):
super().**init**()
self.weight = nn.Parameter(torch.ones(hidden_size))
self.eps = eps

def forward(self, x):
    # x: [batch, seq, hidden]
    # Keep reduction in higher precision if needed in real systems.
    variance = x.pow(2).mean(dim=-1, keepdim=True)
    x = x * torch.rsqrt(variance + self.eps)
    return self.weight * x

\end{lstlisting}

实际实现中，RMSNorm 常会考虑：

\begin{itemize}
\item 输入 dtype 为 BF16/FP16 时，是否使用 FP32 accumulation；
\item hidden size 是否满足向量化加载；
\item 是否融合 residual add；
\item 是否支持 tensor parallel 下的 sequence parallel；
\item backward 是否保存 RMS 或重计算。
\end{itemize}

\subsection{RoPE}
\label{subsec:llama-rope}

RoPE 将位置信息注入到 Query 和 Key 中。对于 head 内的向量维度，每两个维度作为一个二维平面进行旋转：
\[
\begin{bmatrix}
x'*{2i}\
x'*{2i+1}
\end{bmatrix}
=============

\begin{bmatrix}
\cos\theta_{pos,i} & -\sin\theta_{pos,i}\
\sin\theta_{pos,i} & \cos\theta_{pos,i}
\end{bmatrix}
\begin{bmatrix}
x_{2i}\
x_{2i+1}
\end{bmatrix}.
\]

在 attention 中，RoPE 通常应用于：
\[
\mat{Q},\mat{K},
\]
而不应用于 $\mat{V}$。原因是位置关系主要通过 query-key 点积影响 attention score：
\[
\mat{S}=\mat{Q}\mat{K}^{\mathsf{T}}.
\]

RoPE 对推理有一个重要工程细节：decode 阶段每一步新 token 的 position id 必须正确。对于多轮对话、prefix cache、paged KV Cache、speculative decoding 或 chunked prefill，position id 管理会变得复杂。如果位置错位，即使模型权重正确，输出也会异常。

下面是一个带 position ids 的 RoPE 简化实现：

\begin{lstlisting}[language=Python,caption={带 position ids 的 RoPE 应用示意},label={lst:rope-with-position-ids}]
def rotate_half(x):
x_even = x[..., 0::2]
x_odd = x[..., 1::2]
return torch.stack((-x_odd, x_even), dim=-1).flatten(-2)

def apply_rope_with_position_ids(q, k, cos_cache, sin_cache, position_ids):
"""
q, k: [batch, heads, seq, head_dim]
cos_cache, sin_cache: [max_seq, head_dim]
position_ids: [batch, seq]
"""
# Gather cos/sin according to actual positions.
cos = cos_cache[position_ids]  # [batch, seq, head_dim]
sin = sin_cache[position_ids]

cos = cos[:, None, :, :]      # [batch, 1, seq, head_dim]
sin = sin[:, None, :, :]

q = q * cos + rotate_half(q) * sin
k = k * cos + rotate_half(k) * sin
return q, k

\end{lstlisting}

RoPE 的长上下文扩展通常涉及 base、频率或位置缩放策略。无论采用哪种策略，Infra 系统都要保证训练、微调、推理、KV Cache 和 batch 调度对 position ids 的理解一致。

\subsection{SwiGLU}
\label{subsec:llama-swiglu}

LLaMA 风格模型通常使用 SwiGLU FFN。其形式为：
\[
\operatorname{FFN}(\mat{X})
===========================

\left[
\operatorname{SiLU}(\mat{X}\mat{W}_g)
\odot
(\mat{X}\mat{W}_u)
\right]\mat{W}_d.
\]

其中 gate projection 和 up projection 可以合并为一个 fused projection：
\[
[\mat{G},\mat{U}]
=================

\mat{X}\mat{W}_{gu},
\]
然后：
\[
\mat{Y}
=======

\left[
\operatorname{SiLU}(\mat{G})
\odot
\mat{U}
\right]\mat{W}_d.
\]

这样做可以把两次读取输入 $\mat{X}$ 的 GEMM 合并为一次更大的 GEMM，提高硬件效率。

\begin{figure}[H]
\centering
\begin{tikzpicture}[
font=\small,
box/.style={draw, rounded corners=3pt, minimum width=1.6cm, minimum height=0.7cm, align=center},
arrow/.style={-{Latex[length=2.3mm]}, thick},
mult/.style={draw, circle, minimum size=0.6cm, fill=yellow!18}
]

\node[font=\bfseries] at (0,3.1) {LLaMA 风格 SwiGLU FFN};

\node[box, fill=blue!6] (x) at (-5.4,0.6) {$X$};

\node[box, fill=orange!12, minimum width=2.3cm] (fused) at (-3.2,0.6) {Fused\\Gate+Up Proj};
\node[box, fill=green!10] (splitg) at (-1.0,1.35) {$G$};
\node[box, fill=green!10] (splitu) at (-1.0,-0.15) {$U$};

\node[box, fill=purple!10] (silu) at (0.8,1.35) {SiLU};
\node[mult] (mul) at (2.2,0.6) {$\odot$};
\node[box, fill=orange!12] (down) at (3.8,0.6) {Down\\Proj};
\node[box, fill=blue!6] (out) at (5.5,0.6) {$Y$};

\draw[arrow] (x) -- (fused);
\draw[arrow] (fused) -- (splitg);
\draw[arrow] (fused) -- (splitu);
\draw[arrow] (splitg) -- (silu);
\draw[arrow] (silu) -- (mul);
\draw[arrow] (splitu) -- (mul);
\draw[arrow] (mul) -- (down);
\draw[arrow] (down) -- (out);

\node[align=left, text width=10.5cm] at (0,-1.6) {
  Gate 与 Up projection 常合并为一个大 GEMM，然后在 hidden 维切分。
  这能减少 kernel launch 与输入读取，是模型结构和硬件效率共同设计的例子。
};

\end{tikzpicture}
\caption{SwiGLU 中 gate/up projection 的融合实现直觉}
\label{fig:llama-swiglu-fused-proj}
\end{figure}

SwiGLU 的参数量取决于中间维度 $H_{\mathrm{mid}}$。其主要参数为：
\[
2HH_{\mathrm{mid}}+H_{\mathrm{mid}}H=3HH_{\mathrm{mid}}.
\]
为了与标准 FFN 的 $8H^2$ 参数量接近，通常会选择：
\[
H_{\mathrm{mid}}\approx \frac{8}{3}H,
\]
再按硬件友好的倍数对齐。例如为了 Tensor Core、tensor parallel 和内存对齐，实际中间维度常会被 round 到某个整数倍。

\subsection{GQA}
\label{subsec:llama-gqa}

Grouped-Query Attention，简称 GQA，是现代 LLM 推理优化的重要结构。设 query head 数为 $n_q$，key/value head 数为 $n_{kv}$，若：
\[
1<n_{kv}<n_q,
\]
则每组 query heads 共享一组 key/value。

KV Cache 显存与 $n_{kv}$ 成正比。对于 $L$ 层模型，batch size 为 $B$，上下文长度为 $S$，head dim 为 $d_h$，每个元素 $s_{\mathrm{dtype}}$ bytes，KV Cache 粗略为：
\[
M_{\mathrm{KV}}
===============

2\cdot L\cdot B\cdot S\cdot n_{kv}\cdot d_h\cdot s_{\mathrm{dtype}}.
\]
其中前面的 $2$ 分别表示 K 和 V。

如果从 MHA 的 $n_{kv}=n_q$ 改为 GQA 的 $n_{kv}=n_q/g$，则 KV Cache 大小约减少 $g$ 倍。对于推理 decode 阶段，这不仅减少显存，还减少读取 KV Cache 的带宽。

\begin{figure}[H]
\centering
\begin{tikzpicture}[
font=\small,
q/.style={draw, rounded corners=2pt, fill=blue!10, minimum width=0.42cm, minimum height=0.38cm},
kv/.style={draw, rounded corners=2pt, fill=orange!16, minimum width=0.9cm, minimum height=0.45cm, align=center},
arrow/.style={-{Latex[length=1.8mm]}, thick}
]

\node[font=\bfseries] at (0,3.0) {GQA：多个 Query Heads 共享一组 K/V};

\foreach \i in {0,...,7} {
  \node[q] (q\i) at ({-3.2+0.55*\i},1.55) {$Q$};
}

\node[kv] (kv1) at (-2.35,0.55) {K/V 1};
\node[kv] (kv2) at (-0.15,0.55) {K/V 2};
\node[kv] (kv3) at (2.05,0.55) {K/V 3};

\foreach \i in {0,1,2} {
  \draw[arrow] (q\i) -- (kv1);
}
\foreach \i in {3,4,5} {
  \draw[arrow] (q\i) -- (kv2);
}
\foreach \i in {6,7} {
  \draw[arrow] (q\i) -- (kv3);
}

\node[align=left, text width=9.8cm] at (0,-1.25) {
  Query heads 保持较高数量以维持表达能力，K/V heads 减少以降低 KV Cache。
  GQA 是服务系统友好的模型结构选择，尤其适合长上下文和大 batch decode。
};

\end{tikzpicture}
\caption{Grouped-Query Attention 中 Query 与 K/V head 的映射}
\label{fig:gqa-query-kv-sharing}
\end{figure}

GQA 也带来实现细节。attention kernel 需要把多个 query heads 映射到同一个 kv head：
\[
h_{kv}=\left\lfloor \frac{h_q}{g}\right\rfloor,
\]
其中：
\[
g=\frac{n_q}{n_{kv}}.
\]
如果 kernel layout 设计不当，会造成不连续访存或额外 broadcast。高性能 serving 框架通常会对 GQA/MQA 做专门优化。

\subsection{LLaMA 与 GPT 的差异}
\label{subsec:llama-vs-gpt}

GPT 风格和 LLaMA 风格都属于 Decoder-only Transformer，但现代 LLaMA 风格在若干组件上做了不同选择。下面给出一个概览表。

\begin{table}[H]
\centering
\small
\caption{GPT 风格与 LLaMA 风格 Decoder-only 架构对比}
\label{tab:gpt-vs-llama}
\begin{tabular}{p{0.22\textwidth}p{0.32\textwidth}p{0.34\textwidth}}
\toprule
\textbf{组件} & \textbf{GPT 风格常见设计} & \textbf{LLaMA 风格常见设计} \
\midrule
位置编码
& Learned absolute positional embedding
& RoPE，将位置信息注入 Q/K \
\midrule
归一化
& LayerNorm
& RMSNorm \
\midrule
FFN
& GELU MLP
& SwiGLU 门控 FFN \
\midrule
Attention
& MHA 为主
& MHA/GQA，现代版本常使用 GQA \
\midrule
Bias
& 一些实现使用 bias
& 常见实现中许多线性层不使用 bias \
\midrule
推理效率
& 依赖具体实现
& GQA、RoPE cache、RMSNorm 等更贴近现代 serving 优化 \
\bottomrule
\end{tabular}
\end{table}

\begin{figure}[H]
\centering
\begin{tikzpicture}[
font=\small,
box/.style={draw, rounded corners=3pt, minimum width=1.8cm, minimum height=0.62cm, align=center},
arrow/.style={-{Latex[length=2.1mm]}, thick}
]

\node[font=\bfseries] at (0,3.5) {GPT Block vs LLaMA Block};

% GPT
\node[font=\bfseries] at (-3.5,2.75) {GPT 风格};
\node[box, fill=blue!6] (gx) at (-3.5,2.0) {$X$};
\node[box, fill=green!10] (gln1) at (-3.5,1.2) {LayerNorm};
\node[box, fill=orange!12] (gattn) at (-3.5,0.4) {MHA\\Causal Attn};
\node[box, fill=green!10] (gln2) at (-3.5,-0.4) {LayerNorm};
\node[box, fill=purple!10] (gmlp) at (-3.5,-1.2) {GELU MLP};

\draw[arrow] (gx) -- (gln1);
\draw[arrow] (gln1) -- (gattn);
\draw[arrow] (gattn) -- (gln2);
\draw[arrow] (gln2) -- (gmlp);

\node[align=center, text width=3.2cm] at (-3.5,-2.25) {
  learned position embedding\\
  LayerNorm + GELU
};

% LLaMA
\node[font=\bfseries] at (3.5,2.75) {LLaMA 风格};
\node[box, fill=blue!6] (lx) at (3.5,2.0) {$X$};
\node[box, fill=green!10] (lrms1) at (3.5,1.2) {RMSNorm};
\node[box, fill=orange!12] (lattn) at (3.5,0.4) {GQA/MHA\\RoPE Attn};
\node[box, fill=green!10] (lrms2) at (3.5,-0.4) {RMSNorm};
\node[box, fill=purple!10] (lmlp) at (3.5,-1.2) {SwiGLU MLP};

\draw[arrow] (lx) -- (lrms1);
\draw[arrow] (lrms1) -- (lattn);
\draw[arrow] (lattn) -- (lrms2);
\draw[arrow] (lrms2) -- (lmlp);

\node[align=center, text width=3.4cm] at (3.5,-2.25) {
  RoPE\\
  RMSNorm + SwiGLU\\
  GQA for serving
};

\draw[dashed, thick] (0,-2.65) -- (0,2.85);

\node[align=left, text width=10.4cm] at (0,-3.35) {
  二者都属于 Decoder-only Transformer。差异不在整体范式，而在位置编码、归一化、FFN 和 K/V head 共享等细节。
};

\end{tikzpicture}
\caption{GPT Block 与 LLaMA Block 的结构差异}
\label{fig:gpt-vs-llama-block}
\end{figure}

\section{MoE 架构简介}
\label{sec:moe-architecture-intro}

Dense Transformer 中，每个 token 都经过相同的 FFN 参数。MoE，即 Mixture-of-Experts，则把 FFN 替换为多个专家网络，并由 router 为每个 token 选择少数专家。它的目标是提升模型总容量，同时控制每个 token 的实际计算量。

\subsection{Expert、Router 与 Top-k Routing}
\label{subsec:expert-router-topk}

设一层 MoE 有 $E$ 个 experts：
\[
{F_1,F_2,\ldots,F_E}.
\]
每个 expert 通常是一个 FFN。给定 token hidden state $\vect{x}$，router 计算每个 expert 的得分：
\[
\vect{r}=\vect{x}\mat{W}_r,
\quad
\vect{r}\in\mathbb{R}^{E}.
\]
然后选择 top-$k$ 个 expert：
\[
\mathcal{E}_k(\vect{x})=\operatorname{TopK}(\vect{r}).
\]
输出为：
\[
\operatorname{MoE}(\vect{x})
============================

\sum_{e\in \mathcal{E}_k(\vect{x})}
\alpha_e F_e(\vect{x}),
\]
其中 $\alpha_e$ 是 router softmax 后的权重。

Top-1 routing 只选择一个 expert，计算最省但容量利用可能不平衡；Top-2 routing 选择两个 expert，通常更稳定但计算和通信更多。

\begin{figure}[H]
\centering
\begin{tikzpicture}[
font=\small,
token/.style={draw, rounded corners=2pt, minimum width=0.9cm, minimum height=0.5cm, align=center, fill=blue!6},
router/.style={draw, rounded corners=3pt, minimum width=1.65cm, minimum height=0.75cm, align=center, fill=orange!14},
expert/.style={draw, rounded corners=3pt, minimum width=1.45cm, minimum height=0.65cm, align=center, fill=green!10},
arrow/.style={-{Latex[length=2.2mm]}, thick},
faint/.style={-{Latex[length=2.0mm]}, thick, gray!45}
]

\node[font=\bfseries] at (0,3.2) {MoE Router 与 Expert 分发};

\node[token] (x) at (-5.2,0.8) {$x$};
\node[router] (r) at (-3.2,0.8) {Router\\Top-$k$};
\draw[arrow] (x) -- (r);

\foreach \i/\y in {1/2.0,2/1.1,3/0.2,4/-0.7,5/-1.6} {
  \node[expert] (e\i) at (-0.6,\y) {Expert \i};
}

\draw[arrow] (r) -- node[above,sloped,font=\scriptsize] {selected} (e1);
\draw[arrow] (r) -- node[above,sloped,font=\scriptsize] {selected} (e3);
\draw[faint] (r) -- (e2);
\draw[faint] (r) -- (e4);
\draw[faint] (r) -- (e5);

\node[router, fill=purple!10] (combine) at (2.0,0.8) {Weighted\\Combine};
\draw[arrow] (e1.east) -- (combine);
\draw[arrow] (e3.east) -- (combine);

\node[token, fill=yellow!16] (out) at (4.3,0.8) {$y$};
\draw[arrow] (combine) -- (out);

\node[align=left, text width=10.4cm] at (0,-2.55) {
  MoE 的核心是条件计算：每个 token 只激活少数 experts。
  这让模型拥有很大的总参数量，但每个 token 的计算量只与 top-$k$ experts 有关。
};

\end{tikzpicture}
\caption{MoE 中 Router、Expert 与 Top-$k$ 路由}
\label{fig:moe-router-topk}
\end{figure}

\subsection{稀疏激活与计算效率}
\label{subsec:sparse-activation-compute-efficiency}

Dense FFN 中，每个 token 都经过同一个大 FFN：
\[
y=F(x).
\]
MoE 中，每个 token 只经过少数专家：
\[
y=\sum_{e\in\mathcal{E}_k(x)}\alpha_e F_e(x).
\]

如果有 $E$ 个 experts，每个 expert 参数量为 $P_e$，则总 expert 参数量为：
\[
P_{\mathrm{total}}=E P_e.
\]
但每个 token 激活的参数量约为：
\[
P_{\mathrm{active}}=kP_e.
\]
当 $k\ll E$ 时：
\[
P_{\mathrm{active}}\ll P_{\mathrm{total}}.
\]

这就是 MoE 的吸引力：模型容量可以很大，但每 token FLOPs 不必等比例增长。对于同等计算预算，MoE 可能拥有更强的知识容量或任务分化能力。

但 MoE 并非免费。Router 需要学习负载均衡，否则所有 token 都可能挤向少数 experts。常见做法是加入 auxiliary load balancing loss。若某些 experts 过载，而另一些空闲，就会出现：

\begin{itemize}
\item 计算负载不均；
\item GPU 间通信不均；
\item token dropping 或 padding waste；
\item 训练不稳定；
\item 推理 tail latency 上升。
\end{itemize}

\subsection{MoE 的通信挑战}
\label{subsec:moe-communication-challenge}

在单 GPU 上，MoE 可以看作动态选择不同 FFN。但在多 GPU 训练中，experts 往往分布在不同 GPU 上。token 需要根据 router 结果发送到对应 expert 所在 GPU，expert 计算后再把结果发回原位置。这通常需要 all-to-all 通信。

MoE 通信流程可以抽象为：

\begin{enumerate}
\item 每个 GPU 上的 token 经过 router，得到目标 expert；
\item 根据 expert placement，将 token dispatch 到对应 GPU；
\item 每个 GPU 对收到的 token 按 expert 分组并执行 FFN；
\item 将 expert 输出 combine 回原 token 顺序；
\item 继续后续 Transformer block。
\end{enumerate}

\begin{figure}[H]
\centering
\begin{tikzpicture}[
font=\small,
gpu/.style={draw, rounded corners=4pt, minimum width=2.2cm, minimum height=1.2cm, align=center, fill=blue!6},
expert/.style={draw, rounded corners=2pt, minimum width=0.8cm, minimum height=0.35cm, align=center, fill=green!12},
arrow/.style={-{Latex[length=2.0mm]}, thick}
]

\node[font=\bfseries] at (0,3.1) {MoE Expert Parallel 中的 All-to-All 通信};

\foreach \i/\x in {0/-4.5,1/-1.5,2/1.5,3/4.5} {
  \node[gpu] (g\i) at (\x,1.2) {GPU \i};
  \node[expert] at (\x,1.55) {Expert};
  \node[expert] at (\x,0.95) {Expert};
}

\foreach \i/\x in {0/-4.5,1/-1.5,2/1.5,3/4.5} {
  \node[gpu, fill=orange!10] (o\i) at (\x,-1.1) {Output\\tokens};
}

\foreach \i in {0,...,3} {
  \foreach \j in {0,...,3} {
    \draw[arrow, opacity=0.35] (g\i.south) -- (o\j.north);
  }
}

\node[align=left, text width=10.4cm] at (0,-2.55) {
  MoE 的 token dispatch 和 combine 通常产生 all-to-all 通信。
  在跨节点训练中，网络拓扑、带宽、expert placement 和负载均衡会直接决定吞吐。
};

\end{tikzpicture}
\caption{MoE Expert Parallel 中的 token dispatch 与 all-to-all 通信}
\label{fig:moe-all-to-all}
\end{figure}

与 dense Transformer 中常见的 all-reduce、all-gather、reduce-scatter 不同，MoE 的 all-to-all 更依赖动态路由结果，通信形状也更不稳定。这使得 MoE 对网络和调度系统更敏感。

\subsection{MoE 对训练与推理 Infra 的影响}
\label{subsec:moe-impact-on-training-serving-infra}

MoE 对训练 Infra 的影响包括：

\begin{itemize}
\item \textbf{expert parallel}：experts 分布到不同 GPU，需要 token dispatch；
\item \textbf{all-to-all 通信}：通信模式更复杂，对网络带宽和延迟敏感；
\item \textbf{负载均衡}：router 不均衡会造成部分 GPU 成为 straggler；
\item \textbf{动态 shape}：每个 expert 收到的 token 数变化，GEMM shape 不稳定；
\item \textbf{checkpoint 复杂}：expert 参数分布式保存和恢复需要清晰映射。
\end{itemize}

MoE 对推理 Infra 的影响也很明显：

\begin{itemize}
\item 单请求只激活部分 expert，但多请求 batch 中专家分布可能不均；
\item expert placement 会影响跨 GPU 通信；
\item 动态路由使 latency 更难预测；
\item 缓存和 batching 策略要考虑专家负载；
\item 量化时不同 experts 的统计分布可能不同。
\end{itemize}

MoE 的本质是用系统复杂度换模型容量。对于资源充足、网络强、工程栈成熟的团队，它可能带来很高的性价比；对于网络较弱或 serving 调度不成熟的环境，MoE 可能带来难以控制的尾延迟和通信瓶颈。

\section{从架构选择到系统成本}
\label{sec:architecture-choice-to-system-cost}

本章讨论的每个组件都不仅是模型设计，也是系统设计。下面把 GPT、LLaMA 和 MoE 中的关键选择映射到系统成本。

\subsection{位置编码与上下文扩展}
\label{subsec:position-encoding-system-cost}

Learned absolute position embedding 简单，但长上下文扩展受限。RoPE 更适合现代 LLM，但需要管理 cos/sin cache、position ids 和长上下文 scaling。对于 serving 系统，位置编码影响：

\begin{itemize}
\item prefill 与 decode 的 position id 一致性；
\item prefix cache 复用时的位置偏移；
\item speculative decoding 中草稿 token 的位置管理；
\item sliding window 或 chunked attention 的相对位置处理；
\item 长上下文外推策略是否与训练配置一致。
\end{itemize}

\subsection{GQA 与 KV Cache 成本}
\label{subsec:gqa-kv-cache-system-cost}

GQA 的系统收益可以直接从 KV Cache 公式看出：
\[
M_{\mathrm{KV}}
===============

2LBSn_{kv}d_hs_{\mathrm{dtype}}.
\]
在其他参数不变时，将 $n_{kv}$ 从 $n_q$ 降到 $n_q/4$，KV Cache 约减少四倍。对于长上下文和高并发 serving，这可能决定单卡能承载多少请求。

但 GQA 也要求 attention kernel 支持 query head 到 kv head 的映射。如果实现不佳，节省的显存可能被不规则访存和低效 kernel 抵消。因此现代 serving 框架通常会对 MHA/MQA/GQA 分别优化。

\subsection{SwiGLU 与 FFN GEMM}
\label{subsec:swiglu-ffn-system-cost}

SwiGLU 提升表达能力，但引入 gate/up/down 三个 projection。高效实现通常把 gate 和 up 合并为一个 GEMM：
\[
[G,U]=XW_{gu}.
\]
这使 GEMM 更大、更适合 Tensor Core，也减少一次输入读取。随后 SiLU、elementwise multiply 和 down projection 可以进一步通过 fusion 优化。

对于 tensor parallel，FFN 的 up/gate projection 通常做 column parallel，down projection 通常做 row parallel。这会影响 all-reduce 或 reduce-scatter 的位置。后续张量并行章节会专门推导。

\subsection{RMSNorm 与 fused kernel}
\label{subsec:rmsnorm-fused-kernel-cost}

RMSNorm 是 memory-bound 操作。若每一层有两个 RMSNorm，模型有 $L$ 层，则每个 forward 至少调用 $2L$ 次 norm。训练还需要 backward。若 norm、residual add、dropout 分散为多个 kernel，会产生大量小 kernel 和 HBM 读写。

Fused RMSNorm 或 fused add-RMSNorm 可以把：
\[
Y=\RMSNorm(X+R)
\]
放进一个 kernel 中，减少一次或多次显存读写。这类优化看似细节，但在大模型端到端训练和推理中非常重要。

\subsection{MoE 与网络拓扑}
\label{subsec:moe-network-topology-cost}

MoE 将 FFN 从 dense compute 变成 sparse conditional compute，同时把系统瓶颈从单纯 GEMM 扩展到 all-to-all 通信。Expert placement 应尽量利用高带宽拓扑：

\begin{itemize}
\item 节点内优先使用 NVLink/NVSwitch；
\item 跨节点需要高性能 RDMA；
\item expert parallel group 应与物理拓扑对齐；
\item router 负载均衡影响每个 GPU 的实际 token 数；
\item straggler 会拖慢整个训练 step。
\end{itemize}

这说明模型架构选择不能与硬件拓扑割裂。一个在论文中高效的 MoE 结构，如果部署在网络较弱的集群上，可能无法达到预期吞吐。

\begin{table}[H]
\centering
\small
\caption{现代 Decoder-only 架构组件与系统影响}
\label{tab:decoder-component-system-impact}
\begin{tabular}{p{0.22\textwidth}p{0.34\textwidth}p{0.34\textwidth}}
\toprule
\textbf{组件} & \textbf{模型侧作用} & \textbf{Infra 侧影响} \
\midrule
RoPE
& 注入相对位置信息，支持更灵活位置建模
& position id、cos/sin cache、长上下文 scaling、KV Cache 对齐。 \
\midrule
RMSNorm
& 稳定训练，简化归一化
& memory-bound，小 kernel 多，适合 fused implementation。 \
\midrule
SwiGLU
& 提升 FFN 表达能力
& gate/up/down GEMM，适合 projection fusion 与 tensor parallel。 \
\midrule
GQA
& 减少 K/V heads，保持较多 query heads
& KV Cache 显存和带宽下降，attention kernel 需支持 head mapping。 \
\midrule
MoE
& 提升总参数容量，稀疏激活
& all-to-all、负载均衡、expert placement、动态 batch shape。 \
\midrule
Tied Embedding
& 共享输入输出词表权重
& 减少参数，但影响 vocabulary parallel 和 checkpoint layout。 \
\bottomrule
\end{tabular}
\end{table}

\section{端到端现代 Decoder-only 模型实现骨架}
\label{sec:modern-decoder-only-code-skeleton}

下面给出一个教学版现代 LLaMA 风格 Decoder-only 模型骨架。它包含 token embedding、RMSNorm、RoPE、GQA、SwiGLU 和 lm head。代码强调结构关系，不追求生产性能。

\begin{lstlisting}[language=Python,caption={教学版 LLaMA 风格 Decoder-only 模型骨架},label={lst:llama-style-model-skeleton}]
import math
import torch
import torch.nn as nn
import torch.nn.functional as F

class RMSNorm(nn.Module):
def **init**(self, hidden_size, eps=1e-6):
super().**init**()
self.weight = nn.Parameter(torch.ones(hidden_size))
self.eps = eps

def forward(self, x):
    var = x.pow(2).mean(dim=-1, keepdim=True)
    return self.weight * x * torch.rsqrt(var + self.eps)

def rotate_half(x):
x_even = x[..., 0::2]
x_odd = x[..., 1::2]
return torch.stack((-x_odd, x_even), dim=-1).flatten(-2)

def apply_rope(q, k, cos, sin):
# q, k: [B, heads, S, D]
cos = cos[None, None, :, :]
sin = sin[None, None, :, :]
q = q * cos + rotate_half(q) * sin
k = k * cos + rotate_half(k) * sin
return q, k

class GQAAttention(nn.Module):
def **init**(self, hidden_size, num_q_heads, num_kv_heads):
super().**init**()
assert hidden_size % num_q_heads == 0
assert num_q_heads % num_kv_heads == 0

    self.hidden_size = hidden_size
    self.num_q_heads = num_q_heads
    self.num_kv_heads = num_kv_heads
    self.head_dim = hidden_size // num_q_heads
    self.group_size = num_q_heads // num_kv_heads

    self.q_proj = nn.Linear(hidden_size, num_q_heads * self.head_dim, bias=False)
    self.k_proj = nn.Linear(hidden_size, num_kv_heads * self.head_dim, bias=False)
    self.v_proj = nn.Linear(hidden_size, num_kv_heads * self.head_dim, bias=False)
    self.o_proj = nn.Linear(hidden_size, hidden_size, bias=False)

def forward(self, x, cos, sin):
    B, S, H = x.shape

    q = self.q_proj(x)
    k = self.k_proj(x)
    v = self.v_proj(x)

    q = q.view(B, S, self.num_q_heads, self.head_dim).transpose(1, 2)
    k = k.view(B, S, self.num_kv_heads, self.head_dim).transpose(1, 2)
    v = v.view(B, S, self.num_kv_heads, self.head_dim).transpose(1, 2)

    q, k = apply_rope(q, k, cos, sin)

    # Repeat K/V heads so each query head has a corresponding K/V head.
    # Production kernels avoid explicit repeat to save memory bandwidth.
    k = k.repeat_interleave(self.group_size, dim=1)
    v = v.repeat_interleave(self.group_size, dim=1)

    scores = q @ k.transpose(-2, -1)
    scores = scores / math.sqrt(self.head_dim)

    causal_mask = torch.triu(
        torch.ones(S, S, device=x.device, dtype=torch.bool),
        diagonal=1,
    )
    scores = scores.masked_fill(causal_mask, float("-inf"))

    probs = F.softmax(scores, dim=-1)
    out = probs @ v

    out = out.transpose(1, 2).contiguous().view(B, S, H)
    return self.o_proj(out)

class SwiGLUMLP(nn.Module):
def **init**(self, hidden_size, intermediate_size):
super().**init**()
self.gate_proj = nn.Linear(hidden_size, intermediate_size, bias=False)
self.up_proj = nn.Linear(hidden_size, intermediate_size, bias=False)
self.down_proj = nn.Linear(intermediate_size, hidden_size, bias=False)

def forward(self, x):
    return self.down_proj(F.silu(self.gate_proj(x)) * self.up_proj(x))

class DecoderBlock(nn.Module):
def **init**(self, hidden_size, num_q_heads, num_kv_heads, intermediate_size):
super().**init**()
self.input_norm = RMSNorm(hidden_size)
self.attn = GQAAttention(hidden_size, num_q_heads, num_kv_heads)
self.post_attn_norm = RMSNorm(hidden_size)
self.mlp = SwiGLUMLP(hidden_size, intermediate_size)

def forward(self, x, cos, sin):
    x = x + self.attn(self.input_norm(x), cos, sin)
    x = x + self.mlp(self.post_attn_norm(x))
    return x

class LlamaStyleModel(nn.Module):
def **init**(self, vocab_size, hidden_size, num_layers,
num_q_heads, num_kv_heads, intermediate_size):
super().**init**()
self.embed_tokens = nn.Embedding(vocab_size, hidden_size)
self.layers = nn.ModuleList([
DecoderBlock(hidden_size, num_q_heads, num_kv_heads, intermediate_size)
for _ in range(num_layers)
])
self.norm = RMSNorm(hidden_size)
self.lm_head = nn.Linear(hidden_size, vocab_size, bias=False)

def forward(self, input_ids, cos, sin):
    x = self.embed_tokens(input_ids)

    for layer in self.layers:
        x = layer(x, cos, sin)

    x = self.norm(x)
    logits = self.lm_head(x)
    return logits

\end{lstlisting}

这段代码中有意保留了一些低效实现，方便读者看到生产系统需要优化什么：

\begin{itemize}
\item GQA 中显式 \texttt{repeat_interleave} 会浪费显存和带宽；
\item attention score 被显式物化；
\item RoPE、mask、softmax、matmul 没有融合；
\item 没有 KV Cache；
\item 没有 fused RMSNorm；
\item 没有 tensor parallel；
\item logits 没有做 vocabulary parallel 或 fused cross entropy。
\end{itemize}

在真实训练和推理框架中，这些低效点会分别由 FlashAttention、PagedAttention、fused norm、parallel linear、KV Cache manager、vocab parallel cross entropy 等组件解决。也就是说，现代 AI Infra 的大量工作，正是把这个“教学上清晰的模型”重写成“硬件上高效的模型”。

\section{本章小结}
\label{sec:chapter06-summary}

本章系统剖析了 GPT、LLaMA 与现代 Decoder-only LLM 架构，并把模型组件与系统成本联系起来。

\begin{itemize}
\item \textbf{Decoder-only LLM} 由 token embedding、堆叠 decoder blocks、final norm 和 lm head 组成，本质上执行 next-token prediction。
\item \textbf{Embedding 与 LM Head} 的参数量与词表大小直接相关，大词表模型需要关注 vocabulary parallel、logits 显存和 fused cross entropy。
\item \textbf{GPT 风格架构} 使用 causal self-attention、residual stream、LayerNorm/GELU 等组件，是现代 Decoder-only 模型的基础范式。
\item \textbf{LLaMA 风格架构} 更常使用 RMSNorm、RoPE、SwiGLU 和 GQA，体现了现代模型对训练稳定性与推理效率的共同优化。
\item \textbf{RMSNorm} 计算更简单，但仍然是 memory-bound 高频小算子，适合 fused kernel。
\item \textbf{RoPE} 将位置信息注入 Q/K，带来更灵活的位置建模，同时要求推理系统正确管理 position ids 和 cos/sin cache。
\item \textbf{SwiGLU} 通过门控 FFN 提升表达能力，工程上常将 gate/up projection 融合为大 GEMM。
\item \textbf{GQA} 通过减少 K/V heads 降低 KV Cache 显存和带宽，是现代 LLM serving 友好的重要架构选择。
\item \textbf{MoE} 用稀疏激活提升总模型容量，但引入 router、expert parallel、all-to-all 通信和负载均衡挑战。
\item \textbf{模型结构就是系统设计}：每个架构选择都会影响显存、FLOPs、通信、kernel 实现、checkpoint 和 serving 调度。
\end{itemize}

至此，第二篇完成了从 Transformer 原理、LLM 发展范式到现代 Decoder-only 架构的主线。下一篇将进入分布式训练架构：从数据并行开始，分析为什么单卡无法支撑大模型训练，DDP 如何同步梯度，Ring AllReduce 如何工作，以及通信与计算如何重叠。


% ==================================================
% 第三篇：分布式训练架构 Training Infra
% ==================================================
\part{分布式训练架构：从单卡到千卡集群}

\chapter{数据并行 DP / DDP}
\label{chap:data-parallelism}

\section{为什么需要分布式训练}
\label{sec:why-distributed-training}

前两篇我们已经从数学、CUDA、深度学习原理、Transformer 架构和现代 Decoder-only LLM 结构出发，建立了大模型的基本计算图。现在开始进入第三篇：\textbf{分布式训练架构}。

单卡训练适合理解算法，但无法支撑现代大模型的训练规模。原因并不只是“模型太大”，而是模型参数、优化器状态、激活、训练数据、训练时长和容错要求共同超过了单机单卡的能力边界。一个 AI Infrastructure 工程师必须能回答：

\begin{itemize}
\item 模型参数能否放进单卡显存？
\item batch size 能否满足训练稳定性？
\item 单卡训练需要多长时间？
\item 多卡训练时梯度如何同步？
\item 通信量如何估算？
\item 什么时候数据并行足够，什么时候必须引入张量并行或流水线并行？
\item 为什么 DDP 的性能可能被 bucket、通信重叠、拓扑和小参数层影响？
\end{itemize}

本章聚焦最基础也最常用的并行方式：\textbf{数据并行}。数据并行的思想非常直观：每张 GPU 保存一份完整模型副本，不同 GPU 处理不同 mini-batch 分片，然后在 backward 后同步梯度，使所有副本保持一致。

\subsection{模型规模、数据规模与训练时间}
\label{subsec:model-data-training-time}

训练一个模型的总计算量大致与参数量和训练 token 数成正比。对于 dense Transformer，自回归预训练 FLOPs 可粗略估算为：
\[
C \approx 6ND,
\]
其中 $N$ 是模型参数量，$D$ 是训练 token 数。这个估算忽略了具体架构、激活重计算、通信、kernel 效率和实现细节，但足以说明规模问题。

设训练任务总计算量为 $C$，单张 GPU 的有效算力为 $F_{\mathrm{eff}}$，单卡训练时间近似为：
\[
T_1 \approx \frac{C}{F_{\mathrm{eff}}}.
\]
如果使用 $P$ 张 GPU，并且扩展效率为 $\eta(P)$，则训练时间近似为：
\[
T_P \approx \frac{C}{P\eta(P)F_{\mathrm{eff}}}.
\]

这里的 $\eta(P)$ 非常重要。理想线性扩展时：
\[
\eta(P)=1.
\]
但真实训练中，通信、同步、数据加载、straggler、checkpoint 和调度开销都会降低效率：
\[
0<\eta(P)<1.
\]

因此，分布式训练的目标不是简单地“把 GPU 数量堆上去”，而是在给定硬件、网络和模型结构下，让每张 GPU 尽可能多地做有效计算，尽可能少地等待通信和同步。

\subsection{单卡瓶颈}
\label{subsec:single-gpu-bottlenecks}

单卡训练常见瓶颈包括四类。

第一类是 \textbf{显存瓶颈}。训练显存大致包括：

\begin{itemize}
\item 参数；
\item 梯度；
\item 优化器状态；
\item 前向激活；
\item 临时 buffer；
\item CUDA runtime、allocator 和通信 workspace。
\end{itemize}

对于 AdamW，如果参数使用 BF16，梯度使用 BF16，同时保留 FP32 master weight、FP32 一阶矩和 FP32 二阶矩，参数相关训练状态粗略为：
\[
M_{\mathrm{param\ states}}\approx 16N\ \mathrm{bytes}.
\]
一个 $7$B 参数模型，仅这些状态就约为：
\[
16\times 7\times 10^9\ \mathrm{bytes}
\approx 112\ \mathrm{GB}.
\]
这已经超过许多单卡显存，更不用说激活和临时 buffer。

第二类是 \textbf{吞吐瓶颈}。即使模型能勉强放下，单卡训练可能需要过长时间。大模型预训练通常需要消耗数万亿 token，单卡训练时间不可接受。

第三类是 \textbf{batch size 瓶颈}。训练稳定性往往需要足够大的 global batch size。若单卡 micro-batch 太小，训练噪声、归一化统计、梯度累积和吞吐都会受到影响。

第四类是 \textbf{可靠性瓶颈}。训练时间越长，硬件故障、驱动问题、节点重启、网络抖动和存储错误的概率越高。分布式训练系统不仅要快，还要能保存 checkpoint、恢复作业，并在故障后保持一致性。

\subsection{并行训练的基本分类}
\label{subsec:parallel-training-taxonomy}

大模型训练中的并行策略可以粗略分为三类：

\begin{itemize}
\item \textbf{数据并行 Data Parallelism, DP}：每张 GPU 保存完整模型副本，不同 GPU 处理不同数据分片，同步梯度。
\item \textbf{张量并行 Tensor Parallelism, TP}：把一个算子的张量维度切到多张 GPU 上，例如把 Linear 层的权重按列或按行切分。
\item \textbf{流水线并行 Pipeline Parallelism, PP}：把模型不同层切到不同 GPU 或节点上，使用 micro-batch 流水线执行。
\end{itemize}

数据并行最容易理解，也最通用。只要单张 GPU 能放下完整模型及训练状态，通常可以优先使用 DDP 或 FSDP/ZeRO 的数据并行变体。张量并行和流水线并行则用于模型太大、单卡放不下，或者需要更大规模扩展的场景。

\begin{figure}[H]
\centering
\begin{tikzpicture}[
font=\small,
gpu/.style={draw, rounded corners=3pt, minimum width=1.35cm, minimum height=0.75cm, align=center, fill=blue!6},
data/.style={draw, rounded corners=2pt, minimum width=1.15cm, minimum height=0.45cm, align=center, fill=orange!14},
layer/.style={draw, rounded corners=2pt, minimum width=0.95cm, minimum height=0.48cm, align=center, fill=green!10},
arrow/.style={-{Latex[length=2.2mm]}, thick}
]

\node[font=\bfseries] at (0,3.4) {三类基础并行策略：DP、TP、PP};

% DP
\node[font=\bfseries] at (-4.7,2.65) {Data Parallel};
\foreach \i/\y in {0/1.75,1/0.85,2/-0.05} {
  \node[data] (d\i) at (-6.0,\y) {batch};
  \node[gpu] (dp\i) at (-4.35,\y) {Full\\Model};
  \draw[arrow] (d\i) -- (dp\i);
}
\draw[<->, thick, red!70] (-4.35,0.35) -- (-4.35,0.45);
\node[align=center, text width=3.2cm] at (-4.8,-1.05) {模型复制\\数据切分\\梯度同步};

% TP
\node[font=\bfseries] at (0,2.65) {Tensor Parallel};
\node[data] (td) at (-1.2,1.3) {batch};
\node[gpu] (tp0) at (0.1,1.75) {$W_0$\\shard};
\node[gpu] (tp1) at (0.1,0.85) {$W_1$\\shard};
\node[gpu] (tp2) at (0.1,-0.05) {$W_2$\\shard};
\draw[arrow] (td) -- (tp0);
\draw[arrow] (td) -- (tp1);
\draw[arrow] (td) -- (tp2);
\node[align=center, text width=3.4cm] at (0,-1.05) {单层权重切分\\算子内部通信};

% PP
\node[font=\bfseries] at (4.7,2.65) {Pipeline Parallel};
\foreach \i/\x in {1/3.1,2/4.25,3/5.4,4/6.55} {
  \node[layer] (l\i) at (\x,1.3) {L\i};
}
\node[gpu] (pp0) at (3.65,0.2) {GPU0};
\node[gpu] (pp1) at (5.95,0.2) {GPU1};
\draw[arrow] (l1) -- (l2);
\draw[arrow] (l2) -- (l3);
\draw[arrow] (l3) -- (l4);
\draw[dashed] (3.0,0.75) rectangle (4.65,1.85);
\draw[dashed] (5.0,0.75) rectangle (6.95,1.85);
\draw[arrow] (pp0) -- (pp1);
\node[align=center, text width=3.5cm] at (4.9,-1.05) {按层切分\\micro-batch 流水线};

\node[align=left, text width=10.6cm] at (0,-2.4) {
  数据并行是最基础的扩展方式；张量并行解决单层算子过大问题；流水线并行解决模型层数太多、整体放不下的问题。
  现代大模型训练常组合三者形成 3D 并行。
};

\end{tikzpicture}
\caption{数据并行、张量并行与流水线并行的基本区别}
\label{fig:parallelism-overview-dp-tp-pp}
\end{figure}

本章只深入数据并行。后续章节会分别讨论张量并行、流水线并行和 3D 并行。

\section{Data Parallel 基本原理}
\label{sec:data-parallel-basics}

\subsection{replica}
\label{subsec:data-parallel-replica}

数据并行中，每张 GPU 上都有一份完整模型副本，称为 replica。假设有 $P$ 张 GPU，每张 GPU 的模型参数初始相同：
\[
\theta^{(0)}_0=\theta^{(0)}*1=\cdots=\theta^{(0)}*{P-1}.
\]
第 $p$ 张 GPU 处理自己的数据分片 $\mathcal{B}_p$，计算本地 loss：
\[
\mathcal{L}_p(\theta)
\frac{1}{|\mathcal{B}*p|}
\sum*{x\in\mathcal{B}_p}
\ell(x;\theta).
\]
本地 backward 得到本地梯度：
\[
\vect{g}_p
\nabla_{\theta}\mathcal{L}_p(\theta).
\]

为了让所有 replica 保持一致，需要对梯度做平均：
\[
\bar{\vect{g}}
\frac{1}{P}
\sum_{p=0}^{P-1}
\vect{g}_p.
\]
然后每张 GPU 使用相同的平均梯度更新参数：
\[
\theta^{(t+1)}_p
## \theta^{(t)}_p
\eta \bar{\vect{g}},
\quad
p=0,\ldots,P-1.
\]
只要初始参数相同、梯度同步一致、优化器状态更新一致，各 replica 就会保持相同参数。

这也是数据并行的核心不变量：

\[
\boxed{
\theta_0^{(t)}=\theta_1^{(t)}=\cdots=\theta_{P-1}^{(t)}
}
\]

\subsection{mini-batch 切分}
\label{subsec:minibatch-sharding}

设 global batch size 为 $B_{\mathrm{global}}$，GPU 数量为 $P$，每张 GPU 的 local batch size 为 $B_{\mathrm{local}}$。若不使用梯度累积，则：
\[
B_{\mathrm{global}}
P\cdot B_{\mathrm{local}}.
\]

如果使用 gradient accumulation，设累积步数为 $A$，则：
\[
B_{\mathrm{global}}
P\cdot B_{\mathrm{local}}\cdot A.
\]

这个公式非常重要。很多训练配置表面上写的是 micro-batch size，但真正影响优化器更新的是 global batch size。学习率 schedule、warmup、梯度噪声、loss 曲线和收敛稳定性都与 global batch size 相关。

在 PyTorch DDP 中，每个进程通常绑定一张 GPU。使用 \texttt{DistributedSampler} 时，数据集会按 rank 切分，每个 rank 读取不同样本。一个常见错误是忘记设置 sampler 的 epoch，导致不同 epoch 中 shuffle 不正确：

\begin{lstlisting}[language=Python,caption={DistributedSampler 中设置 epoch 的重要性},label={lst:distributed-sampler-set-epoch}]
from torch.utils.data.distributed import DistributedSampler

sampler = DistributedSampler(dataset, shuffle=True)
loader = DataLoader(dataset, sampler=sampler, batch_size=local_batch_size)

for epoch in range(num_epochs):
# Important: ensure each epoch uses a different shuffle order
# while keeping all ranks aligned.
sampler.set_epoch(epoch)

for batch in loader:
    train_step(batch)

\end{lstlisting}

如果多个 rank 读到相同样本，表面上训练仍然能跑，但有效数据吞吐会下降；如果不同 rank 的 step 数不一致，可能导致 collective communication hang，因为有些 rank 进入 all-reduce，而另一些 rank 已经退出。

\subsection{梯度同步}
\label{subsec:gradient-synchronization}

数据并行的关键通信操作是梯度同步。假设模型有参数：
\[
\theta={\theta_1,\theta_2,\ldots,\theta_K}.
\]
每个 rank backward 后得到本地梯度：
\[
\nabla\theta_k^{(p)}.
\]
同步后希望每个 rank 都得到：
\[
\nabla\theta_k
\frac{1}{P}
\sum_{p=0}^{P-1}
\nabla\theta_k^{(p)}.
\]

这通常通过 all-reduce 完成。All-reduce 的语义是：所有 rank 输入一个 tensor，执行规约操作，例如 sum，然后把规约结果返回给所有 rank。若先做 sum，再除以 world size，就得到平均梯度。

\begin{figure}[H]
\centering
\begin{tikzpicture}[
font=\small,
gpu/.style={draw, rounded corners=3pt, minimum width=1.65cm, minimum height=0.78cm, align=center, fill=blue!6},
grad/.style={draw, rounded corners=2pt, minimum width=1.35cm, minimum height=0.52cm, align=center, fill=orange!14},
arrow/.style={-{Latex[length=2.2mm]}, thick}
]

\node[font=\bfseries] at (0,3.2) {数据并行中的本地梯度与 AllReduce 同步};

\foreach \i/\x in {0/-5.1,1/-2.2,2/0.7,3/3.6} {
  \node[gpu] (gpu\i) at (\x,1.6) {GPU \i\\Model};
  \node[grad] (grad\i) at (\x,0.65) {$g_\i$};
  \draw[arrow] (gpu\i) -- node[right,font=\scriptsize] {backward} (grad\i);
}

\node[draw, rounded corners=4pt, fill=green!10, minimum width=3.1cm, minimum height=0.85cm, align=center] (ar) at (-0.75,-0.45)
  {AllReduce\\$\bar{g}=\frac{1}{P}\sum_p g_p$};

\foreach \i in {0,1,2,3} {
  \draw[arrow] (grad\i) -- (ar);
}

\foreach \i/\x in {0/-5.1,1/-2.2,2/0.7,3/3.6} {
  \node[grad, fill=purple!10] (avg\i) at (\x,-1.65) {$\bar{g}$};
  \draw[arrow] (ar) -- (avg\i);
  \draw[arrow] (avg\i) -- node[right,font=\scriptsize] {optimizer step} (\x,-2.45);
}

\node[align=left, text width=10.4cm] at (-0.75,-3.05) {
  每张 GPU 先独立计算本地梯度，再通过 all-reduce 得到平均梯度。
  之后所有 replica 使用相同梯度更新，因此参数继续保持一致。
};

\end{tikzpicture}
\caption{数据并行中的梯度同步}
\label{fig:data-parallel-gradient-sync}
\end{figure}

需要注意，实际实现中不一定等所有参数 backward 完成后才一次性同步所有梯度。PyTorch DDP 会把参数梯度组织成 bucket，当某个 bucket 中的梯度都 ready 后，就可以立即发起 all-reduce，从而与后续 backward 计算重叠。

\subsection{all-reduce}
\label{subsec:all-reduce}

All-reduce 可以理解为两个逻辑阶段：

\begin{enumerate}
\item \textbf{Reduce}：把所有 rank 的 tensor 做求和；
\item \textbf{Broadcast}：把求和结果发回所有 rank。
\end{enumerate}

但高性能实现通常不会真的先集中到一个 rank 再广播，因为这会产生中心瓶颈。Ring AllReduce、Tree AllReduce、CollNet、NVLink/NVSwitch 拓扑感知算法等都用于提高通信效率。

All-reduce 的时间常用带宽-延迟模型估算：
\[
T_{\mathrm{comm}}
\approx
\alpha\cdot n_{\mathrm{steps}}
+
\frac{\mathrm{bytes}}{\beta_{\mathrm{eff}}},
\]
其中：

\begin{itemize}
\item $\alpha$ 是每轮通信的延迟成本；
\item $n_{\mathrm{steps}}$ 是通信步数；
\item $\beta_{\mathrm{eff}}$ 是有效带宽；
\item $\mathrm{bytes}$ 是每个 rank 需要发送或接收的数据量。
\end{itemize}

对于大梯度 tensor，带宽项占主导；对于很多小 tensor，延迟项和 kernel launch/API 开销更明显。因此 DDP 使用 bucket 把多个小梯度合并为较大通信单元，减少小消息开销。

\section{PyTorch DDP}
\label{sec:pytorch-ddp}

PyTorch DistributedDataParallel，简称 DDP，是最常用的数据并行实现之一。DDP 的基本模式是：每个进程绑定一张 GPU，加载同样的模型，处理不同数据分片，并在 backward 时自动同步梯度。

\subsection{process group}
\label{subsec:process-group}

Process group 是分布式通信的基本上下文。它定义哪些进程参与 collective communication。例如，一个 8 卡单机训练中，通常有 8 个进程组成一个默认 process group。

初始化 process group 的典型代码如下：

\begin{lstlisting}[language=Python,caption={初始化 PyTorch Process Group},label={lst:init-process-group}]
import os
import torch
import torch.distributed as dist

def init_distributed():
# Environment variables are usually set by torchrun.
rank = int(os.environ["RANK"])
world_size = int(os.environ["WORLD_SIZE"])
local_rank = int(os.environ["LOCAL_RANK"])

torch.cuda.set_device(local_rank)

dist.init_process_group(
    backend="nccl",
    init_method="env://",
    rank=rank,
    world_size=world_size,
)

return rank, world_size, local_rank

\end{lstlisting}

在 GPU 训练中，backend 通常使用 NCCL。CPU 通信或调试场景可使用其他 backend，但大模型训练基本依赖 NCCL 进行高性能 GPU collective。

Process group 不一定只有一个。后续张量并行、流水线并行和 3D 并行中，会创建多个不同 group：

\begin{itemize}
\item data parallel group；
\item tensor parallel group；
\item pipeline parallel group；
\item expert parallel group；
\item sequence parallel group。
\end{itemize}

每个 group 代表一组 rank 的通信边界。理解 group 是理解分布式训练框架的基础。

\subsection{rank、world size、local rank}
\label{subsec:rank-world-size-local-rank}

分布式训练中有三个常见概念：

\begin{itemize}
\item \textbf{rank}：全局进程编号，范围为 $0,\ldots,P-1$；
\item \textbf{world size}：全局进程总数 $P$；
\item \textbf{local rank}：当前节点内的进程编号，通常用于选择本地 GPU。
\end{itemize}

例如有 2 台机器，每台 8 张 GPU，共 16 个进程。则：

\[
\mathrm{world\ size}=16.
\]
第一台机器上的 local rank 为 $0,\ldots,7$，第二台机器上的 local rank 也为 $0,\ldots,7$，但全局 rank 可能为 $8,\ldots,15$。

\begin{figure}[H]
\centering
\begin{tikzpicture}[
font=\small,
nodebox/.style={draw, rounded corners=4pt, minimum width=4.3cm, minimum height=2.1cm, fill=blue!3},
gpu/.style={draw, rounded corners=2pt, minimum width=0.72cm, minimum height=0.48cm, align=center, fill=orange!12}
]

\node[font=\bfseries] at (0,3.0) {rank、world size 与 local rank};

\node[nodebox, label={[font=\bfseries]above:Node 0}] (n0) at (-3.0,1.0) {};
\node[nodebox, label={[font=\bfseries]above:Node 1}] (n1) at (3.0,1.0) {};

\foreach \i/\x in {0/-4.2,1/-3.4,2/-2.6,3/-1.8} {
  \node[gpu] at (\x,1.35) {G\i};
  \node[font=\scriptsize] at (\x,0.75) {rank \i};
  \node[font=\scriptsize] at (\x,0.35) {local \i};
}

\foreach \i/\x/\r in {0/1.8/4,1/2.6/5,2/3.4/6,3/4.2/7} {
  \node[gpu] at (\x,1.35) {G\i};
  \node[font=\scriptsize] at (\x,0.75) {rank \r};
  \node[font=\scriptsize] at (\x,0.35) {local \i};
}

\node[align=left, text width=10.2cm] at (0,-1.05) {
  local rank 用于绑定本机 GPU；rank 用于全局通信排序；world size 决定 collective 的参与进程数。
  多节点训练问题中，rank mapping 与物理拓扑会影响通信性能。
};

\end{tikzpicture}
\caption{两节点八进程示例中的 rank 与 local rank}
\label{fig:rank-world-size-local-rank}
\end{figure}

一个常见启动方式是：

\begin{lstlisting}[language=bash,caption={使用 torchrun 启动单机 8 卡 DDP},label={lst:torchrun-single-node}]
torchrun 
--nproc_per_node=8 
train_ddp.py 
--config config.yaml
\end{lstlisting}

多机启动时，还需要指定节点数、节点 rank、master 地址和端口：

\begin{lstlisting}[language=bash,caption={使用 torchrun 启动多机 DDP 的示意},label={lst:torchrun-multi-node}]
torchrun 
--nnodes=2 
--nproc_per_node=8 
--node_rank=0 
--master_addr=10.0.0.1 
--master_port=29500 
train_ddp.py
\end{lstlisting}

第二台机器将 \texttt{--node_rank} 设置为 \texttt{1}，其余 master 配置保持一致。

\subsection{gradient bucket}
\label{subsec:gradient-bucket}

如果模型有很多参数 tensor，逐个参数 all-reduce 会产生大量小消息，延迟开销很高。DDP 的做法是把多个参数梯度打包成 bucket。一个 bucket 满足条件后，就作为一个较大的 tensor 发起 all-reduce。

在 backward 中，梯度不是同时产生的，而是按照计算图从后往前逐步 ready。DDP 会在 autograd hook 中捕获参数梯度 ready 事件。当某个 bucket 里的所有梯度都 ready 后，DDP 可以立即发起该 bucket 的 all-reduce。

\begin{figure}[H]
\centering
\begin{tikzpicture}[
font=\small,
param/.style={draw, rounded corners=2pt, minimum width=0.82cm, minimum height=0.45cm, align=center, fill=blue!6},
bucket/.style={draw, rounded corners=3pt, minimum width=2.8cm, minimum height=0.7cm, align=center, fill=orange!12},
arrow/.style={-{Latex[length=2.0mm]}, thick}
]

\node[font=\bfseries] at (0,3.0) {Gradient Bucket：将多个小梯度合并通信};

\foreach \i/\x in {1/-5.0,2/-4.1,3/-3.2,4/-2.3,5/-1.4,6/-0.5,7/0.4,8/1.3,9/2.2,10/3.1} {
  \node[param] (p\i) at (\x,1.55) {$g_{\i}$};
}

\node[bucket] (b0) at (-3.65,0.3) {Bucket 0};
\node[bucket] (b1) at (-0.95,0.3) {Bucket 1};
\node[bucket] (b2) at (1.75,0.3) {Bucket 2};

\foreach \i in {1,2,3} {
  \draw[arrow] (p\i) -- (b0);
}
\foreach \i in {4,5,6} {
  \draw[arrow] (p\i) -- (b1);
}
\foreach \i in {7,8,9,10} {
  \draw[arrow] (p\i) -- (b2);
}

\node[draw, rounded corners=3pt, fill=green!10, minimum width=2.6cm, minimum height=0.65cm, align=center] (ar) at (-0.95,-1.05) {AllReduce buckets};

\draw[arrow] (b0) -- (ar);
\draw[arrow] (b1) -- (ar);
\draw[arrow] (b2) -- (ar);

\node[align=left, text width=10.2cm] at (-0.4,-2.15) {
  Bucket 太小会导致大量小通信；bucket 太大则可能延迟通信启动，减少与 backward 的重叠机会。
  DDP 性能调优中，bucket 大小是一个重要参数。
};

\end{tikzpicture}
\caption{DDP 中的梯度 bucket 机制}
\label{fig:ddp-gradient-bucket}
\end{figure}

Bucket 大小会影响性能：

\begin{itemize}
\item bucket 太小：通信次数多，延迟开销大；
\item bucket 太大：必须等待更多梯度 ready，通信启动晚，重叠减少；
\item 参数顺序不合理：ready 时间相差大的梯度被放在同一 bucket，导致等待；
\item 模型中存在很多小层：更容易受到 bucket 和延迟影响。
\end{itemize}

DDP 的默认 bucket 策略对大多数模型有效，但在大规模模型或特殊结构中，仍可能需要调优。

\subsection{overlap communication and computation}
\label{subsec:overlap-communication-computation}

DDP 的关键优化是通信与计算重叠。Backward 从最后一层开始向前计算梯度。假设后几层梯度已经 ready，DDP 可以立即发起这些梯度 bucket 的 all-reduce，同时 GPU 继续计算更前面层的 backward。

理想情况下，通信时间被隐藏在 backward 计算时间中。若总 backward 计算时间为 $T_{\mathrm{bwd}}$，总通信时间为 $T_{\mathrm{comm}}$，可重叠部分为 $T_{\mathrm{overlap}}$，则 step 时间中的 backward+communication 近似为：
\[
T_{\mathrm{bwd+comm}}
\approx
T_{\mathrm{bwd}}
+
T_{\mathrm{comm}}
-----------------
T_{\mathrm{overlap}}.
\]
若通信完全被隐藏，则：
\[
T_{\mathrm{overlap}}\approx T_{\mathrm{comm}}.
\]
若无法重叠，则：
\[
T_{\mathrm{overlap}}\approx 0.
\]

\begin{figure}[H]
\centering
\begin{tikzpicture}[
font=\small,
comp/.style={draw, rounded corners=2pt, minimum height=0.5cm, fill=blue!12},
comm/.style={draw, rounded corners=2pt, minimum height=0.38cm, fill=orange!16},
axis/.style={-{Latex[length=2.1mm]}, thick}
]

\node[font=\bfseries] at (0,3.1) {DDP 中 backward 与 gradient all-reduce 的重叠};

\draw[axis] (-5.6,2.1) -- (5.6,2.1) node[right] {time};

\node[comp, minimum width=1.45cm] (b4) at (-4.4,1.35) {Bwd L4};
\node[comp, minimum width=1.45cm] (b3) at (-2.85,1.35) {Bwd L3};
\node[comp, minimum width=1.45cm] (b2) at (-1.3,1.35) {Bwd L2};
\node[comp, minimum width=1.45cm] (b1) at (0.25,1.35) {Bwd L1};

\node[comm, minimum width=2.1cm] (c4) at (-2.7,0.55) {AllReduce bucket 4};
\node[comm, minimum width=2.1cm] (c3) at (-1.0,0.05) {AllReduce bucket 3};
\node[comm, minimum width=2.1cm] (c2) at (0.7,-0.45) {AllReduce bucket 2};
\node[comm, minimum width=2.1cm] (c1) at (2.4,-0.95) {AllReduce bucket 1};

\draw[dashed] (-3.65,1.1) -- (-3.65,0.75);
\draw[dashed] (-2.1,1.1) -- (-2.1,0.25);
\draw[dashed] (-0.55,1.1) -- (-0.55,-0.25);
\draw[dashed] (1.0,1.1) -- (1.0,-0.75);

\node[align=left, text width=10.4cm] at (0,-2.15) {
  后层梯度先 ready，对应 bucket 可以先开始 all-reduce。
  如果网络通信能与前面层 backward 重叠，DDP 的 step time 会显著降低。
};

\end{tikzpicture}
\caption{DDP 中通信与计算重叠的时间线}
\label{fig:ddp-overlap-timeline}
\end{figure}

通信重叠并不总是有效。常见失败原因包括：

\begin{itemize}
\item bucket 太大，导致 all-reduce 启动太晚；
\item backward 计算太短，无法隐藏通信；
\item 网络带宽不足，通信尾部暴露；
\item 参数依赖或 autograd 图结构导致梯度 ready 顺序不理想；
\item GPU 上计算 kernel 与通信 kernel 资源竞争；
\item 小模型或小 batch 下，通信延迟占比过高。
\end{itemize}

下面给出一个最小 DDP 训练脚本骨架：

\begin{lstlisting}[language=Python,caption={PyTorch DDP 训练循环骨架},label={lst:pytorch-ddp-training-loop}]
import os
import torch
import torch.distributed as dist
import torch.nn as nn
from torch.nn.parallel import DistributedDataParallel as DDP
from torch.utils.data import DataLoader
from torch.utils.data.distributed import DistributedSampler

def setup_distributed():
rank = int(os.environ["RANK"])
world_size = int(os.environ["WORLD_SIZE"])
local_rank = int(os.environ["LOCAL_RANK"])

torch.cuda.set_device(local_rank)
dist.init_process_group(backend="nccl", init_method="env://")
return rank, world_size, local_rank

def train():
rank, world_size, local_rank = setup_distributed()

dataset = build_dataset()
sampler = DistributedSampler(
    dataset,
    num_replicas=world_size,
    rank=rank,
    shuffle=True,
)
loader = DataLoader(
    dataset,
    batch_size=local_batch_size,
    sampler=sampler,
    num_workers=4,
    pin_memory=True,
)

model = MyModel().cuda(local_rank)

# Each process owns one GPU.
model = DDP(
    model,
    device_ids=[local_rank],
    output_device=local_rank,
    bucket_cap_mb=25,
    find_unused_parameters=False,
)

optimizer = torch.optim.AdamW(model.parameters(), lr=3e-4)

for epoch in range(num_epochs):
    sampler.set_epoch(epoch)

    for batch in loader:
        batch = move_to_cuda(batch, local_rank)

        optimizer.zero_grad(set_to_none=True)

        output = model(batch["input_ids"])
        loss = compute_loss(output, batch["labels"])

        loss.backward()
        optimizer.step()

        if rank == 0:
            print("loss:", loss.item())

dist.destroy_process_group()

if **name** == "**main**":
train()
\end{lstlisting}

实际工程中还需要处理 AMP、gradient clipping、checkpoint、resume、logging、evaluation、异常退出和多节点启动等问题。

\section{通信复杂度分析}
\label{sec:communication-complexity-analysis}

数据并行的计算很容易理解：每张 GPU 处理自己的 local batch。真正决定扩展性的，是梯度同步通信。假设模型梯度总大小为 $G$ bytes，GPU 数为 $P$。每个训练 step 至少需要让所有 rank 获得平均梯度，因此通信不可避免。

\subsection{Ring AllReduce}
\label{subsec:ring-allreduce}

Ring AllReduce 是常见 all-reduce 算法之一。它把 $P$ 个 rank 组织成一个环，并把梯度 tensor 切成 $P$ 个 chunk。算法分为两个阶段：

\begin{enumerate}
\item \textbf{Reduce-Scatter}：每个 rank 逐步把 chunk 沿环传递并累加，最终每个 rank 得到一个 reduced chunk；
\item \textbf{All-Gather}：每个 rank 再把 reduced chunk 沿环传播，最终所有 rank 得到完整规约结果。
\end{enumerate}

每个阶段有 $P-1$ 步，每步每个 rank 发送和接收一个大小为 $G/P$ 的 chunk。因此每个 rank 的发送数据量近似为：
\[
V_{\mathrm{send}}
# 2(P-1)\frac{G}{P}
\frac{2(P-1)}{P}G.
\]
接收量相同。

通信时间粗略为：
\[
T_{\mathrm{ring}}
\approx
2(P-1)\alpha
+
\frac{2(P-1)}{P}\frac{G}{\beta},
\]
其中 $\alpha$ 是每步延迟，$\beta$ 是有效链路带宽。

\begin{figure}[H]
\centering
\begin{tikzpicture}[
font=\small,
rank/.style={draw, circle, minimum size=0.9cm, align=center, fill=blue!8},
arrow/.style={-{Latex[length=2.2mm]}, thick},
chunk/.style={draw, rounded corners=2pt, minimum width=0.55cm, minimum height=0.35cm, fill=orange!14}
]

\node[font=\bfseries] at (0,3.2) {Ring AllReduce：环形传递 chunk};

\node[rank] (r0) at (0,2.0) {R0};
\node[rank] (r1) at (2.7,0.9) {R1};
\node[rank] (r2) at (1.7,-1.7) {R2};
\node[rank] (r3) at (-1.7,-1.7) {R3};
\node[rank] (r4) at (-2.7,0.9) {R4};

\draw[arrow] (r0) -- (r1);
\draw[arrow] (r1) -- (r2);
\draw[arrow] (r2) -- (r3);
\draw[arrow] (r3) -- (r4);
\draw[arrow] (r4) -- (r0);

\foreach \i/\x in {0/-0.85,1/-0.28,2/0.28,3/0.85} {
  \node[chunk] at (\x,0.25) {\scriptsize c\i};
}

\node[align=center] at (0,-0.45) {gradient tensor\\split into chunks};

\node[align=left, text width=10.4cm] at (0,-3.0) {
  Ring AllReduce 带宽利用率高，适合大 tensor 通信。
  但通信步数随 $P$ 增长，延迟项在大规模或小消息场景中会变得明显。
};

\end{tikzpicture}
\caption{Ring AllReduce 的环形通信结构}
\label{fig:ring-allreduce}
\end{figure}

Ring AllReduce 的优势是带宽效率高，所有 rank 同时发送和接收，避免中心节点瓶颈。缺点是延迟步数为 $O(P)$。当 GPU 数很大、梯度 bucket 较小或跨节点链路延迟高时，延迟项会影响扩展效率。

\subsection{Tree AllReduce}
\label{subsec:tree-allreduce}

Tree AllReduce 使用树结构做 reduce 和 broadcast。其延迟步数通常为：
\[
O(\log P).
\]
因此在小消息或大规模 rank 数场景中，tree-based 算法可能更有优势。

可以把 Tree AllReduce 理解为：

\begin{enumerate}
\item 叶子 rank 把数据发送给父节点；
\item 父节点累加子节点数据并向上发送；
\item 根节点得到总和；
\item 根节点再沿树向下广播结果。
\end{enumerate}

Tree 的缺点是某些链路或节点可能承载更多数据，带宽利用可能不如 ring 均衡。现代通信库通常会根据消息大小、拓扑和硬件选择不同算法，而不是固定使用一种。

\begin{figure}[H]
\centering
\begin{tikzpicture}[
font=\small,
rank/.style={draw, circle, minimum size=0.78cm, align=center, fill=green!10},
arrow/.style={-{Latex[length=2.1mm]}, thick}
]

\node[font=\bfseries] at (0,3.0) {Tree AllReduce 的 Reduce 与 Broadcast};

\node[rank] (r0) at (0,2.0) {R0};
\node[rank] (r1) at (-2.0,0.9) {R1};
\node[rank] (r2) at (2.0,0.9) {R2};
\node[rank] (r3) at (-3.0,-0.3) {R3};
\node[rank] (r4) at (-1.0,-0.3) {R4};
\node[rank] (r5) at (1.0,-0.3) {R5};
\node[rank] (r6) at (3.0,-0.3) {R6};

\draw[arrow, blue!70] (r3) -- (r1);
\draw[arrow, blue!70] (r4) -- (r1);
\draw[arrow, blue!70] (r5) -- (r2);
\draw[arrow, blue!70] (r6) -- (r2);
\draw[arrow, blue!70] (r1) -- (r0);
\draw[arrow, blue!70] (r2) -- (r0);

\draw[arrow, orange!80!black, dashed] (r0) .. controls (-1.1,1.8) .. (r1);
\draw[arrow, orange!80!black, dashed] (r0) .. controls (1.1,1.8) .. (r2);
\draw[arrow, orange!80!black, dashed] (r1) .. controls (-2.3,0.3) .. (r3);
\draw[arrow, orange!80!black, dashed] (r1) .. controls (-1.5,0.3) .. (r4);
\draw[arrow, orange!80!black, dashed] (r2) .. controls (1.5,0.3) .. (r5);
\draw[arrow, orange!80!black, dashed] (r2) .. controls (2.3,0.3) .. (r6);

\node[align=left, text width=10.2cm] at (0,-1.85) {
  蓝色实线表示 reduce 方向，橙色虚线表示 broadcast 方向。
  Tree AllReduce 延迟步数较低，但带宽利用和链路负载取决于树形拓扑。
};

\end{tikzpicture}
\caption{Tree AllReduce 的树形通信模式}
\label{fig:tree-allreduce}
\end{figure}

\subsection{带宽、延迟与扩展性}
\label{subsec:bandwidth-latency-scalability}

分布式训练的扩展效率可写为：
\[
\eta(P)
\frac{T_1}{P T_P}.
\]
如果 $P$ 张 GPU 的 step time 为：
\[
T_P
T_{\mathrm{compute}}
+
T_{\mathrm{comm}}
-----------------
T_{\mathrm{overlap}}
+
T_{\mathrm{other}},
\]
那么扩展效率取决于通信和其他开销占比。

当模型较大、计算量充足时，communication 可以被 backward 部分隐藏，DDP 扩展效率较好。当模型较小或 batch 较小时，计算时间不足以覆盖通信，效率下降。

几个关键比值：

\[
\rho_{\mathrm{comm}}
\frac{T_{\mathrm{comm}}-T_{\mathrm{overlap}}}{T_P},
\]
表示暴露通信占 step time 的比例。若 $\rho_{\mathrm{comm}}$ 很高，说明 GPU 大量时间在等待通信。

\[
\rho_{\mathrm{data}}
\frac{T_{\mathrm{data}}}{T_P},
\]
表示数据加载占比。若数据加载慢，即使通信优化很好，GPU 仍会空等。

\[
\rho_{\mathrm{checkpoint}}
\frac{T_{\mathrm{checkpoint}}}{N_{\mathrm{steps\ per\ ckpt}}T_P},
\]
表示摊销到每步的 checkpoint 成本。大规模训练中，checkpoint 体积可能达到 TB 级，不能忽略。

\begin{table}[H]
\centering
\small
\caption{影响 DDP 扩展效率的常见因素}
\label{tab:ddp-scaling-factors}
\begin{tabular}{p{0.22\textwidth}p{0.34\textwidth}p{0.34\textwidth}}
\toprule
\textbf{因素} & \textbf{表现} & \textbf{优化方向} \
\midrule
网络带宽不足
& all-reduce 暴露时间长，GPU 等待通信
& 使用更高速互联，优化 bucket，拓扑感知 rank mapping。 \
\midrule
通信延迟高
& 小 bucket、小参数层扩展差
& 合并小梯度，调整 bucket size，减少同步频率。 \
\midrule
batch 太小
& 计算时间短，无法隐藏通信
& 增大 local batch 或使用 gradient accumulation。 \
\midrule
数据加载慢
& GPU utilization 周期性下降
& 提升 dataloader、缓存、预取、存储吞吐。 \
\midrule
straggler
& 某些 rank 慢，所有 rank 等待
& 检查硬件健康、数据倾斜、网络拥塞、节点温度。 \
\midrule
checkpoint 慢
& 周期性训练停顿
& 异步 checkpoint、分片保存、并行写入、降低频率。 \
\bottomrule
\end{tabular}
\end{table}

\subsection{NCCL 的角色}
\label{subsec:nccl-role}

NCCL 是 NVIDIA GPU 集群中最常用的 collective communication 库之一。PyTorch DDP 使用 NCCL backend 时，all-reduce、broadcast、reduce-scatter、all-gather 等操作通常由 NCCL 执行。

NCCL 的作用包括：

\begin{itemize}
\item 根据 GPU 拓扑选择通信路径；
\item 利用 NVLink、NVSwitch、PCIe、InfiniBand 或 RoCE；
\item 实现 ring、tree 等 collective 算法；
\item 支持 GPU Direct P2P 和 GPUDirect RDMA；
\item 提供高性能 GPU buffer 之间通信。
\end{itemize}

AI Infra 工程中，NCCL 问题非常常见。例如：

\begin{itemize}
\item 某个 rank 卡在 all-reduce；
\item 多机通信带宽远低于预期；
\item 单机 8 卡快，多机扩展后效率骤降；
\item rank mapping 与 NVLink/NIC 拓扑不匹配；
\item 防火墙、网卡、驱动、容器权限或环境变量配置错误；
\item 某个节点硬件异常导致 collective hang。
\end{itemize}

在定位 NCCL 问题时，常用思路是先脱离训练框架，使用最小通信测试确认集群基础通信能力，再回到训练任务分析 bucket、overlap 和拓扑。

\section{数据并行的工程实践与常见陷阱}
\label{sec:data-parallel-engineering-practices}

\subsection{global batch size 与学习率}
\label{subsec:global-batch-size-and-lr}

数据并行改变了 global batch size。如果保持 local batch 不变，GPU 数从 $P$ 增加到 $2P$，global batch size 也会翻倍：
\[
B_{\mathrm{global}}=P B_{\mathrm{local}} A.
\]
这会改变优化过程。常见经验是随着 batch size 增大适当调整学习率，但具体策略取决于模型、数据和优化器。

在大模型训练中，global batch 通常以 token 数衡量：
\[
B_{\mathrm{tokens}}
P\cdot B_{\mathrm{local}}\cdot A\cdot S.
\]
其中 $S$ 是序列长度。许多训练配置会写成每 step 多少 tokens，而不是多少 samples。

如果恢复训练时 world size、gradient accumulation 或 sequence length 变化，需要特别小心学习率 schedule、optimizer state 和 consumed tokens 的一致性。否则虽然程序能跑，但训练曲线可能出现异常。

\subsection{梯度累积与 no_sync}
\label{subsec:gradient-accumulation-no-sync}

当显存不足以使用较大 local batch 时，可以使用 gradient accumulation。假设累积 $A$ 步，每步只做 forward/backward，不立即 optimizer step，累积到第 $A$ 步再更新参数。

在 DDP 中，如果每个 micro-step 都同步梯度，会产生不必要通信。可以使用 \texttt{no_sync()} 在前 $A-1$ 个 micro-step 跳过 all-reduce，只在最后一步同步。

\begin{lstlisting}[language=Python,caption={DDP 中使用 no_sync 做梯度累积},label={lst:ddp-no-sync-gradient-accumulation}]
accum_steps = 4

for step, batch in enumerate(loader):
micro_step = step % accum_steps

# Avoid gradient synchronization for intermediate accumulation steps.
if micro_step < accum_steps - 1:
    context = model.no_sync()
else:
    context = nullcontext()

with context:
    output = model(batch["input_ids"])
    loss = compute_loss(output, batch["labels"])
    loss = loss / accum_steps
    loss.backward()

if micro_step == accum_steps - 1:
    optimizer.step()
    optimizer.zero_grad(set_to_none=True)

\end{lstlisting}

这里有两个关键点：

\begin{itemize}
\item loss 需要除以 \texttt{accum_steps}，否则梯度尺度会变大；
\item 只有最后一个 micro-step 才同步梯度并执行 optimizer step。
\end{itemize}

如果忘记使用 \texttt{no_sync()}，训练结果可能仍然正确，但通信量增加 $A$ 倍，吞吐明显下降。

\subsection{unused parameters 与动态图}
\label{subsec:unused-parameters-and-dynamic-graph}

DDP 默认假设所有参数都会在每次 forward/backward 中参与计算。如果某些参数没有产生梯度，可能导致同步等待异常。可以设置：
\[
\texttt{find_unused_parameters=True}.
\]
但这会引入额外图遍历开销，不应随意开启。

unused parameters 常见于：

\begin{itemize}
\item 条件分支模型；
\item 多任务模型中某些 head 不参与当前 batch；
\item MoE 中部分 expert 未被选中；
\item 冻结部分参数但未正确设置 \texttt{requires_grad=False}；
\item 返回值中某些 loss 没有连接到所有参数。
\end{itemize}

对于大模型训练，动态图和 unused parameters 会影响 DDP bucket 构建和通信重叠。若模型结构固定，尽量保持每步计算图一致；若必须使用条件计算，需要理解通信框架如何处理未使用参数。

\subsection{checkpoint 与 rank 0 陷阱}
\label{subsec:checkpoint-rank0-pitfall}

在普通 DDP 中，每张 GPU 保存完整模型参数，理论上只需 rank 0 保存 checkpoint。但这有几个注意点：

\begin{itemize}
\item optimizer state 在每个 rank 上也相同，可以 rank 0 保存；
\item scheduler、global step、consumed samples/tokens 必须保存；
\item 如果使用 AMP，需要保存 scaler；
\item 多节点训练时 rank 0 写 checkpoint 可能成为单点 IO 瓶颈；
\item 如果结合 ZeRO/FSDP，checkpoint 可能是分片的，不能简单只保存 rank 0。
\end{itemize}

一个简化 checkpoint 写法如下：

\begin{lstlisting}[language=Python,caption={DDP 中仅 rank 0 保存 checkpoint 的简化示意},label={lst:ddp-rank0-checkpoint}]
def save_checkpoint(model, optimizer, scheduler, step, rank, path):
if rank != 0:
return

# DDP wraps the original model under .module.
state = {
    "model": model.module.state_dict(),
    "optimizer": optimizer.state_dict(),
    "scheduler": scheduler.state_dict(),
    "step": step,
}

torch.save(state, path)

\end{lstlisting}

恢复时，所有 rank 都需要加载相同 checkpoint，或者 rank 0 加载后 broadcast 参数。简单单机实验可以每个 rank 从共享文件系统读取；大规模训练中，这可能造成存储风暴，需要分布式 checkpoint 方案。

\subsection{DDP 适用边界}
\label{subsec:ddp-applicability-boundary}

DDP 的前提是每张 GPU 能保存完整模型、梯度和优化器状态。如果模型太大，DDP 就不够了。此时需要：

\begin{itemize}
\item \textbf{ZeRO / FSDP}：在数据并行维度切分 optimizer state、gradient 和 parameter；
\item \textbf{Tensor Parallel}：切分单层矩阵；
\item \textbf{Pipeline Parallel}：按层切分模型；
\item \textbf{Activation Checkpointing}：降低激活显存；
\item \textbf{Offload}：把部分状态移到 CPU 或 NVMe。
\end{itemize}

可以把 DDP 的适用条件写成显存不等式：
\[
M_{\mathrm{params}}
+
M_{\mathrm{grads}}
+
M_{\mathrm{optimizer}}
+
M_{\mathrm{activations}}
+
M_{\mathrm{buffers}}
\le
M_{\mathrm{HBM}}.
\]
若这个不等式不成立，就必须引入更复杂的并行或分片技术。

\section{本章小结}
\label{sec:chapter07-summary}

本章系统介绍了数据并行和 PyTorch DDP 的基本原理、通信机制与工程实践。

\begin{itemize}
\item \textbf{分布式训练的动机} 来自显存、吞吐、batch size、训练时长和可靠性等多重瓶颈。
\item \textbf{数据并行} 的核心思想是每张 GPU 保存完整模型副本，处理不同数据分片，并同步平均梯度。
\item \textbf{global batch size} 满足 $B_{\mathrm{global}}=P\cdot B_{\mathrm{local}}\cdot A$，其中 $P$ 是 GPU 数，$A$ 是梯度累积步数。
\item \textbf{梯度同步} 通常通过 all-reduce 完成，所有 rank 得到相同平均梯度后执行相同 optimizer step。
\item \textbf{PyTorch DDP} 使用 process group 管理通信，使用 rank/world size/local rank 描述分布式进程。
\item \textbf{gradient bucket} 将多个梯度合并通信，减少小消息开销，并允许通信与 backward 计算重叠。
\item \textbf{Ring AllReduce} 带宽效率高，每个 rank 通信量约为 $\frac{2(P-1)}{P}G$，但延迟步数随 $P$ 增长。
\item \textbf{Tree AllReduce} 具有较低延迟步数，适合某些小消息或大规模场景，但带宽利用取决于树形拓扑。
\item \textbf{NCCL} 是 GPU collective communication 的核心库，负责利用 NVLink、NVSwitch、PCIe、InfiniBand 或 RoCE 执行高性能通信。
\item \textbf{DDP 的边界} 是单卡必须能保存完整模型和训练状态；当模型太大时，需要 ZeRO、FSDP、TP、PP 或 3D 并行。
\end{itemize}

下一章将进入张量并行。我们会分析为什么数据并行在模型太大时不够，如何把 Linear 层按列或按行切分到多张 GPU，Megatron-LM 如何在 Transformer 内部组织 tensor parallel group，以及前向和反向传播中 all-reduce、all-gather、reduce-scatter 分别出现在什么位置。

\chapter{张量并行 TP：Megatron-LM 的核心思想}
\label{chap:tensor-parallelism}

\section{为什么数据并行不够}
\label{sec:why-data-parallel-not-enough}

上一章讨论了数据并行。数据并行的核心假设是：每张 GPU 都能保存一份完整模型副本、梯度、优化器状态和前向激活。只要这个假设成立，DDP 就是一种简单、稳定、通用的扩展方式。

但大模型训练很快会撞到这个假设的边界。对于几十亿、几百亿甚至更大规模的模型，单张 GPU 不再能容纳完整训练状态。此时仅仅增加数据并行副本并不能解决问题，因为每个副本仍然需要放下一整份模型。

张量并行 Tensor Parallelism，简称 TP，解决的是另一个问题：\textbf{把单个算子的张量切分到多张 GPU 上，让多张 GPU 共同完成一个 Transformer block 内部的矩阵计算。}

如果说数据并行是：
[
\text{复制模型，切分数据},
]
那么张量并行就是：
[
\text{切分模型内部张量，共同计算一个算子}.
]

\subsection{单卡显存放不下模型}
\label{subsec:model-does-not-fit-single-gpu}

训练显存主要由以下部分组成：

[
M_{\mathrm{train}}
==================

M_{\mathrm{params}}
+
M_{\mathrm{grads}}
+
M_{\mathrm{optimizer}}
+
M_{\mathrm{activations}}
+
M_{\mathrm{buffers}}.
]

对于 AdamW，若参数使用 BF16，梯度使用 BF16，优化器状态和 master weight 使用 FP32，则参数相关状态粗略为：
[
M_{\mathrm{param\ states}}\approx 16N\ \mathrm{bytes},
]
其中 $N$ 是参数量。

如果 $N=70\mathrm{B}$，则仅参数、梯度和优化器相关状态就约为：
[
16\times 70\times 10^9
======================

1.12\times 10^{12}\ \mathrm{bytes}
\approx 1.12\ \mathrm{TB}.
]
这远远超过单张 GPU 的 HBM 容量。

数据并行无法降低每张 GPU 上的模型状态大小。使用 $P$ 张 GPU 做 DDP，每张 GPU 仍然保存完整模型：
[
M_{\mathrm{per\ GPU}}^{\mathrm{DDP}}
\approx
M_{\mathrm{train}}.
]
这意味着，如果单卡放不下模型，DDP 的 world size 再大也无济于事。

张量并行则直接切分模型参数。若某个大矩阵参数被切到 $P_{\mathrm{TP}}$ 张 GPU 上，理想情况下该矩阵在每张 GPU 上的参数显存可下降为：
[
\frac{1}{P_{\mathrm{TP}}}.
]
当然，代价是前向和反向传播中需要引入额外通信。

\subsection{参数、梯度、优化器状态}
\label{subsec:param-grad-optimizer-states}

大模型训练中，一个参数张量并不只是一个权重矩阵。以 AdamW 为例，一个参数 $\theta$ 通常关联：

\begin{itemize}
\item 参数本身 $\theta$；
\item 梯度 $\nabla \theta$；
\item FP32 master weight；
\item 一阶矩 $m$；
\item 二阶矩 $v$。
\end{itemize}

如果参数矩阵：
[
\mat{W}\in\mathbb{R}^{H\times 4H}
]
被切成 $P_{\mathrm{TP}}$ 份，则不仅 $\mat{W}$ 被切分，与之对应的梯度和优化器状态也会随同切分：

[
\mat{W}
=======

[\mat{W}*0,\mat{W}*1,\ldots,\mat{W}*{P*{\mathrm{TP}}-1}],
]
[
\nabla\mat{W}
=============

[\nabla\mat{W}*0,\nabla\mat{W}*1,\ldots,\nabla\mat{W}*{P*{\mathrm{TP}}-1}].
]

每个 rank 只保存本地 shard：
[
\mat{W}_p,\quad
\nabla\mat{W}_p,\quad
m_p,\quad
v_p.
]

这与 ZeRO/FSDP 的切分维度不同。ZeRO 在数据并行维度上切分 optimizer state、gradient 或 parameter；张量并行则通常在模型内部矩阵维度上切分参数，使多张 GPU 共同执行同一个层。

\begin{table}[H]
\centering
\small
\caption{DDP、ZeRO/FSDP 与 Tensor Parallelism 的切分对象对比}
\label{tab:ddp-zero-tp-comparison}
\begin{tabular}{p{0.2\textwidth}p{0.28\textwidth}p{0.42\textwidth}}
\toprule
\textbf{策略} & \textbf{切分对象} & \textbf{核心作用} \
\midrule
DDP
& 数据 batch
& 每卡完整模型副本，处理不同数据分片，同步梯度。 \
\midrule
ZeRO / FSDP
& 参数、梯度、优化器状态
& 在数据并行 rank 之间切分训练状态，降低每卡显存。 \
\midrule
Tensor Parallelism
& 层内矩阵或 attention heads
& 把单个算子拆到多张 GPU 上共同计算，使单层模型参数可扩展。 \
\midrule
Pipeline Parallelism
& 模型层
& 把不同层放到不同 GPU 上，解决整体层堆叠过深、模型太大的问题。 \
\bottomrule
\end{tabular}
\end{table}

\subsection{模型并行的必要性}
\label{subsec:necessity-of-model-parallelism}

当模型规模继续增大，仅依靠 DDP 不够；仅依靠 ZeRO/FSDP 也不一定够。原因包括：

\begin{itemize}
\item 单个 Linear 层的权重矩阵可能很大；
\item 单层 GEMM 的计算量巨大，单卡执行时间过长；
\item attention heads 数量多，可以天然切分；
\item FFN 中间维度大，适合按 hidden 维切分；
\item 训练吞吐需要更多 GPU 共同参与单个样本或单个 micro-batch 的计算。
\end{itemize}

模型并行可以粗略分为两类：

\begin{itemize}
\item \textbf{层内并行}：切分单层算子，例如 tensor parallelism；
\item \textbf{层间并行}：切分模型层，例如 pipeline parallelism。
\end{itemize}

张量并行属于层内并行。它最适合 Transformer，因为 Transformer 中绝大部分计算都集中在规则的矩阵乘法上：

[
\mat{Y}=\mat{X}\mat{W}.
]

只要能合理切分 $\mat{W}$ 或 $\mat{X}$，并在必要位置插入通信，就可以让多张 GPU 协作完成一个 Linear 层。

\begin{figure}[H]
\centering
\begin{tikzpicture}[
font=\small,
gpu/.style={draw, rounded corners=3pt, minimum width=1.55cm, minimum height=0.75cm, align=center, fill=blue!6},
tensor/.style={draw, rounded corners=2pt, minimum width=1.2cm, minimum height=0.55cm, align=center, fill=orange!14},
arrow/.style={-{Latex[length=2.2mm]}, thick}
]

```
\node[font=\bfseries] at (0,3.2) {数据并行与张量并行的直观区别};

% Data parallel
\node[font=\bfseries] at (-3.8,2.5) {Data Parallel};
\foreach \i/\y in {0/1.5,1/0.55,2/-0.4} {
  \node[tensor] (d\i) at (-5.3,\y) {data \i};
  \node[gpu] (dp\i) at (-3.5,\y) {Full\\Model};
  \draw[arrow] (d\i) -- (dp\i);
}
\node[align=center, text width=3.7cm] at (-4.3,-1.35) {
  每张 GPU 有完整模型\\
  只切分 batch
};

% Tensor parallel
\node[font=\bfseries] at (3.3,2.5) {Tensor Parallel};
\node[tensor] (x) at (1.0,0.6) {same input};
\foreach \i/\y in {0/1.55,1/0.55,2/-0.45} {
  \node[gpu] (tp\i) at (3.2,\y) {Weight\\Shard \i};
  \draw[arrow] (x) -- (tp\i);
}
\node[tensor, fill=green!10] (combine) at (5.4,0.55) {combine};
\foreach \i in {0,1,2} {
  \draw[arrow] (tp\i) -- (combine);
}

\node[align=center, text width=3.8cm] at (3.5,-1.35) {
  多张 GPU 共同计算\\
  同一个层或算子
};

\draw[dashed, thick] (0,-1.9) -- (0,2.7);
```

\end{tikzpicture}
\caption{数据并行复制模型，张量并行切分模型内部张量}
\label{fig:dp-vs-tp}
\end{figure}

\section{张量并行基本原理}
\label{sec:tensor-parallel-basics}

张量并行的核心是对 Linear 层进行切分。因为 Transformer 的 attention projection、output projection、FFN up projection、down projection 和 lm head 都可以视作 Linear 层，理解 Linear 的切分就能理解大部分 TP 设计。

考虑线性层：
[
\mat{Y}=\mat{X}\mat{W},
]
其中：
[
\mat{X}\in\mathbb{R}^{B S\times H_{\mathrm{in}}},
\quad
\mat{W}\in\mathbb{R}^{H_{\mathrm{in}}\times H_{\mathrm{out}}},
\quad
\mat{Y}\in\mathbb{R}^{B S\times H_{\mathrm{out}}}.
]

为了简化记号，我们把 batch 和 sequence 合并成 token 维：
[
N_{\mathrm{tok}}=B S.
]
则：
[
\mat{X}\in\mathbb{R}^{N_{\mathrm{tok}}\times H_{\mathrm{in}}}.
]

张量并行常见切法有两种：

\begin{itemize}
\item \textbf{Column Parallel Linear}：按输出维切分 $\mat{W}$；
\item \textbf{Row Parallel Linear}：按输入维切分 $\mat{W}$。
\end{itemize}

这两个切法在 Megatron-LM 中非常关键。一个 Transformer block 的 attention 和 FFN 会交替使用这两种并行线性层，使通信次数尽可能少。

\subsection{按列切分 Linear Layer}
\label{subsec:column-parallel-linear}

Column Parallel Linear 按输出维切分权重矩阵。将：
[
\mat{W}\in\mathbb{R}^{H_{\mathrm{in}}\times H_{\mathrm{out}}}
]
按列切成 $P$ 份：
[
\mat{W}
=======

[\mat{W}*0,\mat{W}*1,\ldots,\mat{W}*{P-1}],
]
其中：
[
\mat{W}*p\in
\mathbb{R}^{H*{\mathrm{in}}\times \frac{H*{\mathrm{out}}}{P}}.
]

每个 TP rank 持有一份 $\mat{W}_p$，输入 $\mat{X}$ 在所有 rank 上相同。每个 rank 本地计算：
[
\mat{Y}_p=\mat{X}\mat{W}_p.
]
最终完整输出为：
[
\mat{Y}
=======

[\mat{Y}_0,\mat{Y}*1,\ldots,\mat{Y}*{P-1}].
]

如果后续算子可以直接消费分片输出，则不需要立刻 all-gather。若后续算子需要完整 $\mat{Y}$，则需要在输出维上 all-gather。

\begin{figure}[H]
\centering
\begin{tikzpicture}[
font=\small,
matrix/.style={draw, rounded corners=2pt, minimum width=1.6cm, minimum height=1.15cm, align=center},
shard/.style={draw, rounded corners=2pt, minimum width=0.9cm, minimum height=1.15cm, align=center},
gpu/.style={draw, rounded corners=3pt, minimum width=1.35cm, minimum height=0.58cm, align=center, fill=blue!6},
arrow/.style={-{Latex[length=2.2mm]}, thick}
]

```
\node[font=\bfseries] at (0,3.25) {Column Parallel Linear：按输出维切分权重};

\node[matrix, fill=orange!12] (x) at (-5.2,0.8) {$X$\\$N,H_{in}$};

\foreach \i/\xpos/\col in {0/-2.4/red!12,1/-1.35/green!12,2/-0.3/purple!12,3/0.75/yellow!18} {
  \node[shard, fill=\col] (w\i) at (\xpos,0.8) {$W_\i$};
  \node[gpu] (g\i) at (\xpos,-0.55) {rank \i};
  \draw[dashed] (w\i) -- (g\i);
  \draw[arrow] (x.east) -- (w\i.west);
}

\foreach \i/\xpos/\col in {0/2.1/red!12,1/3.05/green!12,2/4.0/purple!12,3/4.95/yellow!18} {
  \node[shard, fill=\col, minimum height=0.95cm] (y\i) at (\xpos,0.8) {$Y_\i$};
  \draw[arrow] (w\i.east) -- (y\i.west);
}

\draw[decorate,decoration={brace,amplitude=5pt},thick]
  (1.65,1.55) -- (5.4,1.55)
  node[midway,above=7pt] {$Y=[Y_0,Y_1,Y_2,Y_3]$};

\node[align=left, text width=10.5cm] at (0,-2.0) {
  Column Parallel Linear 让每个 rank 计算一部分输出通道。
  如果下一步能继续使用 sharded output，可以避免立即 all-gather。
};
```

\end{tikzpicture}
\caption{Column Parallel Linear 的张量切分方式}
\label{fig:column-parallel-linear}
\end{figure}

Column Parallel Linear 的本地参数量为：
[
\frac{H_{\mathrm{in}}H_{\mathrm{out}}}{P}.
]
本地计算量为：
[
2N_{\mathrm{tok}}H_{\mathrm{in}}\frac{H_{\mathrm{out}}}{P}.
]
因此参数和计算都按 TP size 近似缩小 $P$ 倍。

反向传播中，设上游梯度分片为：
[
\mat{G}_{Y,p}
=============

\frac{\partial \mathcal{L}}{\partial \mat{Y}_p}.
]
本地权重梯度：
[
\frac{\partial \mathcal{L}}{\partial \mat{W}_p}
===============================================

\mat{X}^{\mathsf{T}}\mat{G}*{Y,p}.
]
输入梯度的本地贡献为：
[
\mat{G}*{X,p}
=============

\mat{G}_{Y,p}\mat{W}_p^{\mathsf{T}}.
]
完整输入梯度为所有 rank 的求和：
[
\mat{G}_X
=========

\sum_{p=0}^{P-1}
\mat{G}_{X,p}.
]
因此 Column Parallel Linear 在 backward 中通常需要对输入梯度做 all-reduce，除非上游/下游结构可以把这个通信延后或与其他通信融合。

\subsection{按行切分 Linear Layer}
\label{subsec:row-parallel-linear}

Row Parallel Linear 按输入维切分权重矩阵。将：
[
\mat{W}\in\mathbb{R}^{H_{\mathrm{in}}\times H_{\mathrm{out}}}
]
按行切成 $P$ 份：
[
\mat{W}
=======

\begin{bmatrix}
\mat{W}*0\
\mat{W}*1\
\vdots\
\mat{W}*{P-1}
\end{bmatrix},
]
其中：
[
\mat{W}*p\in
\mathbb{R}^{\frac{H*{\mathrm{in}}}{P}\times H*{\mathrm{out}}}.
]

相应地，输入也按 hidden 维切分：
[
\mat{X}
=======

[\mat{X}*0,\mat{X}*1,\ldots,\mat{X}*{P-1}],
]
其中：
[
\mat{X}*p\in
\mathbb{R}^{N*{\mathrm{tok}}\times \frac{H*{\mathrm{in}}}{P}}.
]

每个 rank 本地计算：
[
\mat{Y}_p=\mat{X}_p\mat{W}_p.
]
完整输出为：
[
\mat{Y}
=======

\sum_{p=0}^{P-1}
\mat{Y}_p.
]
因此 Row Parallel Linear 在 forward 中通常需要 all-reduce 求和。

\begin{figure}[H]
\centering
\begin{tikzpicture}[
font=\small,
shard/.style={draw, rounded corners=2pt, minimum width=1.1cm, minimum height=0.52cm, align=center},
mat/.style={draw, rounded corners=2pt, minimum width=1.45cm, minimum height=1.45cm, align=center},
gpu/.style={draw, rounded corners=3pt, minimum width=1.2cm, minimum height=0.5cm, align=center, fill=blue!6},
arrow/.style={-{Latex[length=2.1mm]}, thick}
]

```
\node[font=\bfseries] at (0,3.25) {Row Parallel Linear：按输入维切分权重};

\foreach \i/\y/\col in {0/1.55/red!12,1/0.9/green!12,2/0.25/purple!12,3/-0.4/yellow!18} {
  \node[shard, fill=\col] (x\i) at (-4.8,\y) {$X_\i$};
  \node[shard, fill=\col, minimum width=1.35cm] (w\i) at (-2.4,\y) {$W_\i$};
  \node[gpu] (g\i) at (-0.55,\y) {rank \i};
  \node[shard, fill=\col, minimum width=1.25cm] (yp\i) at (1.25,\y) {$Y_\i$};
  \draw[arrow] (x\i) -- (w\i);
  \draw[arrow] (w\i) -- (g\i);
  \draw[arrow] (g\i) -- (yp\i);
}

\node[mat, fill=orange!10] (sum) at (3.65,0.58) {$Y=\sum_pY_p$};
\foreach \i in {0,1,2,3} {
  \draw[arrow] (yp\i) -- (sum);
}

\node[align=left, text width=10.5cm] at (0,-2.0) {
  Row Parallel Linear 要对各 rank 的局部输出求和，因此 forward 中通常需要 all-reduce。
  它常接在 Column Parallel Linear 后面，用来消费按 hidden 维切分的中间激活。
};
```

\end{tikzpicture}
\caption{Row Parallel Linear 的张量切分与输出求和}
\label{fig:row-parallel-linear}
\end{figure}

Row Parallel Linear 的反向传播中，设完整输出梯度为：
[
\mat{G}_Y
=========

\frac{\partial \mathcal{L}}{\partial \mat{Y}}.
]
由于：
[
\mat{Y}=\sum_p \mat{X}_p\mat{W}_p,
]
每个 rank 本地计算：
[
\frac{\partial \mathcal{L}}{\partial \mat{W}_p}
===============================================

\mat{X}_p^{\mathsf{T}}\mat{G}_Y,
]
[
\frac{\partial \mathcal{L}}{\partial \mat{X}_p}
===============================================

\mat{G}_Y\mat{W}_p^{\mathsf{T}}.
]
此时每个 rank 得到的是输入梯度的一个 shard，不需要在 backward 中对 $\mat{G}_X$ 做 all-reduce，除非后续结构要求完整输入梯度。

\subsection{前向传播中的通信}
\label{subsec:tp-forward-communication}

Column Parallel Linear 和 Row Parallel Linear 的通信位置不同。

对于 Column Parallel Linear：

[
\mat{Y}_p=\mat{X}\mat{W}_p.
]

如果输出保持分片：
[
\mat{Y}=[\mat{Y}*0,\ldots,\mat{Y}*{P-1}],
]
则 forward 不需要通信。如果需要完整输出，则需要 all-gather：

[
\operatorname{AllGather}(\mat{Y}_p)\rightarrow \mat{Y}.
]

对于 Row Parallel Linear：

[
\mat{Y}_p=\mat{X}_p\mat{W}_p,
]
[
\mat{Y}=\sum_p\mat{Y}_p.
]
因此 forward 需要 all-reduce：

[
\operatorname{AllReduceSum}(\mat{Y}_p)\rightarrow \mat{Y}.
]

Megatron-LM 的关键设计之一，就是把 Column Parallel Linear 和 Row Parallel Linear 成对安排。例如 FFN 中：

[
\mat{H}=\phi(\mat{X}\mat{W}_1),
]
[
\mat{Y}=\mat{H}\mat{W}_2.
]

将 $\mat{W}_1$ 做 Column Parallel，得到分片 $\mat{H}_p$；然后将 $\mat{W}_2$ 做 Row Parallel，直接消费 $\mat{H}_p$。这样中间激活 $\mat{H}$ 不需要 all-gather，只在第二个线性层输出处做一次 all-reduce。

[
\mat{H}*p=\phi(\mat{X}\mat{W}*{1,p}),
]
[
\mat{Y}=
\sum_p
\mat{H}*p\mat{W}*{2,p}.
]

这个组合非常优雅，因为它把两个大矩阵乘法之间的中间通信消掉了。

\begin{figure}[H]
\centering
\begin{tikzpicture}[
font=\small,
box/.style={draw, rounded corners=3pt, minimum width=1.75cm, minimum height=0.68cm, align=center},
shard/.style={draw, rounded corners=2pt, minimum width=1.0cm, minimum height=0.5cm, align=center},
arrow/.style={-{Latex[length=2.2mm]}, thick}
]

```
\node[font=\bfseries] at (0,3.15) {Column Parallel + Row Parallel：避免中间 all-gather};

\node[box, fill=blue!6] (x) at (-5.2,0.8) {$X$\\replicated};
\node[box, fill=orange!12] (col) at (-2.9,0.8) {Column\\Parallel};
\node[shard, fill=green!10] (h0) at (-0.7,1.45) {$H_0$};
\node[shard, fill=green!10] (h1) at (-0.7,0.75) {$H_1$};
\node[shard, fill=green!10] (h2) at (-0.7,0.05) {$H_2$};

\node[box, fill=purple!10] (row) at (1.6,0.8) {Row\\Parallel};
\node[box, fill=yellow!16] (ar) at (3.8,0.8) {AllReduce\\sum};
\node[box, fill=blue!6] (y) at (5.8,0.8) {$Y$\\replicated};

\draw[arrow] (x) -- (col);
\draw[arrow] (col) -- (h0);
\draw[arrow] (col) -- (h1);
\draw[arrow] (col) -- (h2);
\draw[arrow] (h0) -- (row);
\draw[arrow] (h1) -- (row);
\draw[arrow] (h2) -- (row);
\draw[arrow] (row) -- (ar);
\draw[arrow] (ar) -- (y);

\node[align=left, text width=10.4cm] at (0,-1.4) {
  第一个线性层按列切分，输出保持 sharded；第二个线性层按行切分，直接消费 sharded hidden。
  最终只需要一次 all-reduce 得到完整输出。
};
```

\end{tikzpicture}
\caption{Column Parallel 与 Row Parallel 组合减少中间通信}
\label{fig:column-row-parallel-pair}
\end{figure}

\subsection{反向传播中的通信}
\label{subsec:tp-backward-communication}

张量并行的 backward 通信与 forward 通信互为镜像。理解这一点有助于判断通信是否可被隐藏。

对于 Column Parallel Linear：

[
\mat{Y}*p=\mat{X}\mat{W}*p.
]
forward 如果不 all-gather，则每个 rank 只持有 $\mat{Y}*p$。Backward 中，每个 rank 根据 $\mat{G}*{Y,p}$ 计算本地 $\mat{G}*{X,p}$：
[
\mat{G}*{X,p}
=============

\mat{G}_{Y,p}\mat{W}_p^{\mathsf{T}}.
]
如果上一层需要完整 $\mat{G}_X$，则：
[
\mat{G}*X=\sum_p\mat{G}*{X,p}.
]
这需要 all-reduce。

对于 Row Parallel Linear：

[
\mat{Y}=\sum_p \mat{X}_p\mat{W}_p.
]
forward 已经通过 all-reduce 得到完整 $\mat{Y}$。Backward 中，$\mat{G}*Y$ 在所有 rank 上相同，每个 rank 计算自己的 $\mat{G}*{X,p}$ 和 $\nabla \mat{W}*p$。因此 backward 不需要对 $\mat{G}*{X,p}$ 求和。

这种互补性让 Megatron-LM 能够通过特定的通信算子安排，控制每个 Transformer block 中的通信次数。

\begin{table}[H]
\centering
\small
\caption{Column Parallel 与 Row Parallel Linear 的通信位置}
\label{tab:column-row-parallel-communication}
\begin{tabular}{p{0.24\textwidth}p{0.32\textwidth}p{0.32\textwidth}}
\toprule
\textbf{并行线性层} & \textbf{Forward 通信} & \textbf{Backward 通信} \
\midrule
Column Parallel
& 若输出保持 sharded，无需通信；若需要完整输出，all-gather
& 输入梯度通常需要 all-reduce 求和。 \
\midrule
Row Parallel
& 局部输出需要 all-reduce 求和
& 输入梯度天然是 sharded，通常无需 all-reduce。 \
\bottomrule
\end{tabular}
\end{table}

下面给出教学版的 Column Parallel Linear 和 Row Parallel Linear。代码用于说明通信语义，不是生产级高性能实现。

\begin{lstlisting}[language=Python,caption={教学版 Column Parallel Linear 与 Row Parallel Linear},label={lst:column-row-parallel-linear}]
import torch
import torch.nn as nn
import torch.distributed as dist

class ColumnParallelLinear(nn.Module):
def **init**(self, in_features, out_features, tp_group, gather_output=False):
super().**init**()
self.tp_group = tp_group
self.tp_size = dist.get_world_size(tp_group)
self.gather_output = gather_output

```
    assert out_features % self.tp_size == 0
    self.out_per_rank = out_features // self.tp_size

    self.weight = nn.Parameter(
        torch.empty(in_features, self.out_per_rank)
    )
    nn.init.xavier_normal_(self.weight)

def forward(self, x):
    # x: [tokens, in_features], replicated across TP ranks
    y_local = x @ self.weight  # [tokens, out_per_rank]

    if not self.gather_output:
        return y_local

    # Gather output shards along hidden dimension.
    y_parts = [torch.empty_like(y_local) for _ in range(self.tp_size)]
    dist.all_gather(y_parts, y_local, group=self.tp_group)
    y = torch.cat(y_parts, dim=-1)
    return y
```

class RowParallelLinear(nn.Module):
def **init**(self, in_features, out_features, tp_group):
super().**init**()
self.tp_group = tp_group
self.tp_size = dist.get_world_size(tp_group)

```
    assert in_features % self.tp_size == 0
    self.in_per_rank = in_features // self.tp_size

    self.weight = nn.Parameter(
        torch.empty(self.in_per_rank, out_features)
    )
    nn.init.xavier_normal_(self.weight)

def forward(self, x_local):
    # x_local: [tokens, in_per_rank], sharded across TP ranks
    y_local = x_local @ self.weight  # [tokens, out_features]

    # Sum partial outputs across TP ranks.
    dist.all_reduce(y_local, op=dist.ReduceOp.SUM, group=self.tp_group)
    return y_local
```

\end{lstlisting}

生产系统会进一步优化：

\begin{itemize}
\item 使用 fused bias、activation 和 GEMM epilogue；
\item 避免不必要的 all-gather；
\item 使用 reduce-scatter 替代 all-reduce 的一部分；
\item 与 sequence parallel 结合降低激活显存；
\item 使用自定义 autograd function 精确控制 forward/backward 通信。
\end{itemize}

\section{Transformer 中的 TP}
\label{sec:tp-in-transformer}

Transformer 之所以适合张量并行，是因为它的大部分参数和计算都来自 Linear 层。一个 Decoder Block 主要包括：

\begin{itemize}
\item Attention QKV projection；
\item Attention output projection；
\item FFN up/gate projection；
\item FFN down projection；
\item Norm 和 residual 等小算子。
\end{itemize}

张量并行的设计目标是：在尽量少通信的前提下，把这些大矩阵乘法切到多个 TP rank 上。

\subsection{Attention QKV projection 切分}
\label{subsec:tp-attention-qkv-projection}

输入：
[
\mat{X}\in\mathbb{R}^{N_{\mathrm{tok}}\times H}.
]
QKV projection 通常写作：
[
[\mat{Q},\mat{K},\mat{V}]
=========================

\mat{X}\mat{W}*{QKV},
]
其中：
[
\mat{W}*{QKV}\in\mathbb{R}^{H\times 3H}.
]

Megatron 风格 TP 通常对 QKV projection 使用 Column Parallel Linear。也就是说，按输出维切分：
[
\mat{W}_{QKV}
=============

[
\mat{W}*{QKV,0},
\mat{W}*{QKV,1},
\ldots,
\mat{W}_{QKV,P-1}
].
]
每个 rank 计算一部分 Q/K/V heads。

如果模型有 $n_h$ 个 attention heads，TP size 为 $P$，则每个 rank 持有：
[
\frac{n_h}{P}
]
个 heads。每个 rank 本地计算这些 heads 的 attention，不需要与其他 rank 交换 attention score。原因是每个 head 的 attention 是独立的：

[
\operatorname{head}_i
=====================

\Attention(\mat{Q}_i,\mat{K}_i,\mat{V}_i).
]

只要每个 head 的 Q、K、V 都在同一个 rank 上，attention 计算就是本地的。

\begin{figure}[H]
\centering
\begin{tikzpicture}[
font=\small,
box/.style={draw, rounded corners=3pt, minimum width=1.7cm, minimum height=0.68cm, align=center},
head/.style={draw, rounded corners=2pt, minimum width=1.05cm, minimum height=0.5cm, align=center},
arrow/.style={-{Latex[length=2.1mm]}, thick}
]

```
\node[font=\bfseries] at (0,3.2) {Attention QKV Projection 的 Head 维切分};

\node[box, fill=blue!6] (x) at (-5.1,0.8) {$X$\\replicated};
\node[box, fill=orange!12] (qkv) at (-2.9,0.8) {Column Parallel\\QKV Projection};

\foreach \i/\xpos/\col in {0/-0.8/red!12,1/0.55/green!12,2/1.9/purple!12,3/3.25/yellow!18} {
  \node[head, fill=\col] (h\i) at (\xpos,1.25) {heads\\rank \i};
  \node[box, fill=\col, minimum width=1.2cm, minimum height=0.58cm] (attn\i) at (\xpos,0.25) {Attn};
  \draw[arrow] (h\i) -- (attn\i);
}

\node[box, fill=purple!8, minimum width=2.0cm] (out) at (5.2,0.75) {sharded\\attention out};

\draw[arrow] (x) -- (qkv);
\draw[arrow] (qkv) -- (h0);
\draw[arrow] (qkv) -- (h1);
\draw[arrow] (qkv) -- (h2);
\draw[arrow] (qkv) -- (h3);

\foreach \i in {0,1,2,3} {
  \draw[arrow] (attn\i) -- (out);
}

\node[align=left, text width=10.4cm] at (0,-1.6) {
  QKV projection 按 head 切分后，每个 rank 只计算自己负责的 heads。
  Attention score 与 softmax 都在本地完成，无需跨 TP rank 同步 attention matrix。
};
```

\end{tikzpicture}
\caption{Transformer Attention 中 QKV projection 的张量并行切分}
\label{fig:tp-qkv-head-split}
\end{figure}

对于 GQA/MQA，切分会稍微复杂，因为 query heads 和 kv heads 数量不同。要求 TP size 与 $n_q$、$n_{kv}$ 的整除关系合理，否则需要更复杂的 head mapping 或跨 rank 共享 K/V。实际训练框架通常会限制或自动处理这些约束。

\subsection{Attention output projection 切分}
\label{subsec:tp-attention-output-projection}

Attention 计算后，每个 rank 得到自己负责 heads 的输出：
[
\mat{O}*p\in
\mathbb{R}^{N*{\mathrm{tok}}\times \frac{H}{P}}.
]
Multi-Head Attention 的 output projection 为：
[
\mat{Y}=\mat{O}\mat{W}_O,
]
其中：
[
\mat{W}_O\in\mathbb{R}^{H\times H}.
]

因为 $\mat{O}$ 已经按 hidden/head 维切分，最自然的做法是对 $\mat{W}*O$ 做 Row Parallel Linear：
[
\mat{W}*O=
\begin{bmatrix}
\mat{W}*{O,0}\
\mat{W}*{O,1}\
\vdots\
\mat{W}_{O,P-1}
\end{bmatrix}.
]
每个 rank 计算：
[
\mat{Y}_p=\mat{O}*p\mat{W}*{O,p}.
]
最终：
[
\mat{Y}=\sum_p\mat{Y}_p.
]
这需要一次 all-reduce。

也就是说，attention 子层的典型通信模式是：

[
\text{QKV Column Parallel}
\rightarrow
\text{local attention}
\rightarrow
\text{Output Row Parallel}
\rightarrow
\text{AllReduce}.
]

这个 all-reduce 发生在 attention output projection 后，用于恢复 replicated residual stream。

\subsection{FFN up projection 与 down projection 切分}
\label{subsec:tp-ffn-projection-split}

FFN 的结构与 attention 类似，也可以使用 Column Parallel + Row Parallel 的组合。以标准 MLP 为例：

[
\mat{H}=\phi(\mat{X}\mat{W}_1),
]
[
\mat{Y}=\mat{H}\mat{W}_2.
]

其中：
[
\mat{W}*1\in\mathbb{R}^{H\times H*{\mathrm{ffn}}},
\quad
\mat{W}*2\in\mathbb{R}^{H*{\mathrm{ffn}}\times H}.
]

对 $\mat{W}*1$ 使用 Column Parallel：
[
\mat{W}*1=[\mat{W}*{1,0},\ldots,\mat{W}*{1,P-1}],
]
每个 rank 得到：
[
\mat{H}*p=\phi(\mat{X}\mat{W}*{1,p}).
]

再对 $\mat{W}*2$ 使用 Row Parallel：
[
\mat{W}*2=
\begin{bmatrix}
\mat{W}*{2,0}\
\vdots\
\mat{W}*{2,P-1}
\end{bmatrix}.
]
每个 rank 计算：
[
\mat{Y}_p=\mat{H}*p\mat{W}*{2,p},
]
然后 all-reduce：
[
\mat{Y}=\sum_p\mat{Y}_p.
]

对于 SwiGLU：

[
\operatorname{FFN}(\mat{X})
===========================

[\operatorname{SiLU}(\mat{X}\mat{W}_g)\odot(\mat{X}\mat{W}_u)]\mat{W}_d.
]

通常 $\mat{W}_g$ 和 $\mat{W}_u$ 都做 Column Parallel，且使用相同切分方式。每个 rank 本地拥有 gate 和 up 的对应 shard：

[
\mat{G}*p=\mat{X}\mat{W}*{g,p},
]
[
\mat{U}*p=\mat{X}\mat{W}*{u,p},
]
[
\mat{H}_p=\operatorname{SiLU}(\mat{G}_p)\odot \mat{U}_p.
]

然后 down projection 使用 Row Parallel：

[
\mat{Y}=\sum_p \mat{H}*p\mat{W}*{d,p}.
]

\begin{figure}[H]
\centering
\begin{tikzpicture}[
font=\small,
box/.style={draw, rounded corners=3pt, minimum width=1.7cm, minimum height=0.65cm, align=center},
shard/.style={draw, rounded corners=2pt, minimum width=1.1cm, minimum height=0.5cm, align=center},
arrow/.style={-{Latex[length=2.1mm]}, thick}
]

```
\node[font=\bfseries] at (0,3.25) {FFN 中的 Tensor Parallel：Up/Gate Column Parallel，Down Row Parallel};

\node[box, fill=blue!6] (x) at (-5.2,1.0) {$X$};

\node[box, fill=orange!12] (gateup) at (-3.0,1.0) {Gate/Up\\Column Parallel};

\foreach \i/\y/\col in {0/1.75/red!12,1/1.0/green!12,2/0.25/purple!12} {
  \node[shard, fill=\col] (h\i) at (-0.9,\y) {$H_\i$};
}

\node[box, fill=purple!10] (down) at (1.3,1.0) {Down\\Row Parallel};
\node[box, fill=yellow!16] (ar) at (3.4,1.0) {AllReduce};
\node[box, fill=blue!6] (y) at (5.3,1.0) {$Y$};

\draw[arrow] (x) -- (gateup);
\draw[arrow] (gateup) -- (h0);
\draw[arrow] (gateup) -- (h1);
\draw[arrow] (gateup) -- (h2);
\draw[arrow] (h0) -- (down);
\draw[arrow] (h1) -- (down);
\draw[arrow] (h2) -- (down);
\draw[arrow] (down) -- (ar);
\draw[arrow] (ar) -- (y);

\node[align=left, text width=10.4cm] at (0,-1.25) {
  FFN 中间维度通常很大，天然适合按 hidden 维切分。
  Column Parallel 产生 sharded intermediate，Row Parallel 直接消费该分片，只在输出处 all-reduce。
};
```

\end{tikzpicture}
\caption{Transformer FFN 中的张量并行切分}
\label{fig:tp-ffn-split}
\end{figure}

\subsection{TP 中的 all-gather 与 reduce-scatter}
\label{subsec:tp-allgather-reducescatter}

基础 TP 常使用 all-reduce 恢复 replicated hidden state。但 all-reduce 可以分解为：

[
\operatorname{AllReduce}
========================

\operatorname{ReduceScatter}
+
\operatorname{AllGather}.
]

在 sequence parallel 中，系统会利用这一点，把某些激活沿 sequence 维切分，让每个 rank 只保存部分 token 的激活。这样可以降低 activation memory。

例如，在普通 TP 中，attention output projection 后做 all-reduce，所有 rank 得到完整：
[
\mat{Y}\in\mathbb{R}^{B\times S\times H}.
]

在 sequence parallel 中，可以改为 reduce-scatter，让每个 rank 只得到 sequence 维的一部分：
[
\mat{Y}_p\in
\mathbb{R}^{B\times \frac{S}{P}\times H}.
]

在后续需要完整 sequence 的位置，再使用 all-gather 恢复。这样通信总量可能类似，但激活显存下降，并且一些 norm、dropout、residual 等操作可以在 sequence shard 上本地执行。

\begin{figure}[H]
\centering
\begin{tikzpicture}[
font=\small,
box/.style={draw, rounded corners=3pt, minimum width=2.0cm, minimum height=0.68cm, align=center},
shard/.style={draw, rounded corners=2pt, minimum width=0.9cm, minimum height=0.5cm, align=center},
arrow/.style={-{Latex[length=2.1mm]}, thick}
]

```
\node[font=\bfseries] at (0,3.2) {AllReduce 与 ReduceScatter + AllGather};

\node[box, fill=blue!6] (local) at (-5.1,1.0) {partial\\outputs};
\node[box, fill=orange!12] (ar) at (-2.9,1.0) {AllReduce};
\node[box, fill=green!10] (full) at (-0.6,1.0) {replicated\\full tensor};

\draw[arrow] (local) -- (ar);
\draw[arrow] (ar) -- (full);

\node[box, fill=blue!6] (local2) at (-5.1,-0.65) {partial\\outputs};
\node[box, fill=orange!12] (rs) at (-2.9,-0.65) {ReduceScatter};
\node[shard, fill=green!10] (s0) at (-0.9,-0.25) {seq 0};
\node[shard, fill=green!10] (s1) at (-0.9,-0.95) {seq 1};
\node[box, fill=purple!10] (ag) at (1.4,-0.65) {AllGather\\when needed};
\node[box, fill=green!10] (full2) at (3.8,-0.65) {full tensor};

\draw[arrow] (local2) -- (rs);
\draw[arrow] (rs) -- (s0);
\draw[arrow] (rs) -- (s1);
\draw[arrow] (s0) -- (ag);
\draw[arrow] (s1) -- (ag);
\draw[arrow] (ag) -- (full2);

\node[align=left, text width=10.5cm] at (0,-2.2) {
  Sequence Parallel 利用 reduce-scatter 保持 sequence 维 sharded，降低激活显存。
  后续在需要完整输入的算子前，再通过 all-gather 恢复。
};
```

\end{tikzpicture}
\caption{AllReduce 与 ReduceScatter/AllGather 的关系}
\label{fig:allreduce-reducescatter-allgather}
\end{figure}

张量并行中的通信选择不是纯数学问题，而是系统设计问题。需要综合考虑：

\begin{itemize}
\item 通信量；
\item 激活显存；
\item kernel 是否支持 sharded 输入；
\item 是否能与计算重叠；
\item NCCL collective 的性能；
\item rank mapping 与物理拓扑。
\end{itemize}

\section{Megatron-LM 并行设计}
\label{sec:megatron-lm-parallel-design}

Megatron-LM 的核心贡献之一，是系统化地把 Transformer 内部的 Linear 层按列/行切分，并组织 tensor model parallel group，使多张 GPU 共同训练一个大 Transformer 模型。

在 Megatron 风格设计中，一个训练作业通常包含多种并行维度：

\begin{itemize}
\item tensor model parallel；
\item pipeline model parallel；
\item data parallel；
\item sequence parallel；
\item context parallel 或 expert parallel 等扩展维度。
\end{itemize}

本章重点关注 tensor model parallel。

\subsection{tensor model parallel group}
\label{subsec:tensor-model-parallel-group}

Tensor model parallel group 是参与同一个层内张量切分的一组 rank。假设总 GPU 数为：
[
W
]
即 world size。TP size 为：
[
P_{\mathrm{TP}}.
]
数据并行 size 为：
[
P_{\mathrm{DP}}.
]
如果暂不考虑 pipeline parallel，则：
[
W=P_{\mathrm{TP}}\cdot P_{\mathrm{DP}}.
]

例如 $W=8$，$P_{\mathrm{TP}}=2$，则 $P_{\mathrm{DP}}=4$。可以组织为：

[
\text{TP groups}:
(0,1),(2,3),(4,5),(6,7),
]
[
\text{DP groups}:
(0,2,4,6),(1,3,5,7).
]

同一个 TP group 内的 rank 共同保存一个模型副本的不同 tensor shards；不同 DP group 之间保存不同 replica，并在数据并行维度同步梯度。

\begin{figure}[H]
\centering
\begin{tikzpicture}[
font=\small,
rank/.style={draw, circle, minimum size=0.72cm, align=center, fill=blue!8},
tpbox/.style={draw, rounded corners=4pt, dashed, thick, fill=orange!5},
dpbox/.style={draw, rounded corners=4pt, dotted, thick, fill=green!5},
arrow/.style={-{Latex[length=2.0mm]}, thick}
]

```
\node[font=\bfseries] at (0,3.2) {TP Group 与 DP Group 的正交关系};

\foreach \r/\x/\y in {
  0/-4.2/1.4,1/-3.1/1.4,
  2/-4.2/0.2,3/-3.1/0.2,
  4/0.6/1.4,5/1.7/1.4,
  6/0.6/0.2,7/1.7/0.2
} {
  \node[rank] (r\r) at (\x,\y) {\r};
}

\node[tpbox, fit=(r0)(r1), label={[font=\scriptsize]above:TP group 0}] {};
\node[tpbox, fit=(r2)(r3), label={[font=\scriptsize]below:TP group 1}] {};
\node[tpbox, fit=(r4)(r5), label={[font=\scriptsize]above:TP group 2}] {};
\node[tpbox, fit=(r6)(r7), label={[font=\scriptsize]below:TP group 3}] {};

\draw[green!60!black, thick, dotted] (r0) -- (r2) -- (r6) -- (r4) -- (r0);
\draw[green!60!black, thick, dotted] (r1) -- (r3) -- (r7) -- (r5) -- (r1);

\node[align=left, text width=10.2cm] at (0,-1.55) {
  横向虚线框表示 TP group：组内 rank 共同切分同一个模型副本。
  绿色点线表示 DP group：组内 rank 对应相同 tensor shard 的不同数据并行副本，需要同步梯度。
};
```

\end{tikzpicture}
\caption{Tensor Parallel Group 与 Data Parallel Group 的组织方式}
\label{fig:tp-dp-groups}
\end{figure}

这种 group 组织方式非常关键。若 rank mapping 不合理，TP 通信可能跨节点发生，性能会很差。通常希望 TP group 尽量放在同一节点内，利用 NVLink/NVSwitch；DP group 可以跨节点，因为数据并行梯度同步通信量较大但频率和模式更规则，可以利用 RDMA 网络。

\subsection{vocabulary parallel embedding}
\label{subsec:vocabulary-parallel-embedding}

LLM 的 embedding 和 lm head 与词表大小 $V$ 成正比。当 $V$ 很大时，embedding table 和 output projection 会很大：

[
\mat{E}\in\mathbb{R}^{V\times H},
]
[
\mat{W}_{\mathrm{lm}}\in\mathbb{R}^{H\times V}.
]

Vocabulary parallelism 将词表维切分到 TP ranks 上。假设 TP size 为 $P$，每个 rank 保存一段词表：

[
V_p=\frac{V}{P}.
]
Embedding table 被切为：
[
\mat{E}
=======

\begin{bmatrix}
\mat{E}_0\
\mat{E}*1\
\vdots\
\mat{E}*{P-1}
\end{bmatrix},
]
其中：
[
\mat{E}_p\in\mathbb{R}^{V_p\times H}.
]

对于输入 token id，每个 rank 只处理属于自己词表范围的 token，其余位置输出零向量。然后对所有 rank 的 embedding 输出做 all-reduce sum，得到完整 embedding。

输出 lm head 也可以按词表维切分。每个 rank 计算一部分 logits：
[
\mat{Z}*p=\mat{H}\mat{W}*{\mathrm{lm},p},
]
其中：
[
\mat{Z}*p\in\mathbb{R}^{N*{\mathrm{tok}}\times V_p}.
]
训练 cross entropy 时，可以使用 vocabulary parallel cross entropy，避免显式 all-gather 完整 logits。

\begin{figure}[H]
\centering
\begin{tikzpicture}[
font=\small,
vocab/.style={draw, rounded corners=3pt, minimum width=1.15cm, minimum height=1.0cm, align=center},
box/.style={draw, rounded corners=3pt, minimum width=1.8cm, minimum height=0.65cm, align=center},
arrow/.style={-{Latex[length=2.0mm]}, thick}
]

```
\node[font=\bfseries] at (0,3.1) {Vocabulary Parallel Embedding / LM Head};

\foreach \i/\x/\col in {0/-3.3/red!12,1/-2.1/green!12,2/-0.9/purple!12,3/0.3/yellow!18} {
  \node[vocab, fill=\col] (v\i) at (\x,1.1) {$V_\i$\\rank \i};
}

\draw[decorate,decoration={brace,amplitude=5pt},thick]
  (-3.9,1.85) -- (0.9,1.85)
  node[midway,above=7pt] {vocabulary dimension split};

\node[box, fill=blue!6] (embed) at (3.0,1.1) {Embedding\\AllReduce sum};
\foreach \i in {0,1,2,3} {
  \draw[arrow] (v\i) -- (embed);
}

\node[box, fill=orange!12] (hidden) at (-2.0,-0.75) {Hidden\\$N,H$};
\node[box, fill=green!10] (logits) at (1.0,-0.75) {Local logits\\$N,V/P$};
\node[box, fill=purple!10] (loss) at (4.0,-0.75) {Vocab Parallel\\Cross Entropy};

\draw[arrow] (hidden) -- (logits);
\draw[arrow] (logits) -- (loss);

\node[align=left, text width=10.2cm] at (0,-2.15) {
  词表并行避免每个 rank 保存完整 embedding/lm head。
  对大词表模型，vocabulary parallel cross entropy 还能避免物化完整 logits。
};
```

\end{tikzpicture}
\caption{词表维度上的张量并行}
\label{fig:vocabulary-parallelism}
\end{figure}

词表并行对训练大词表模型非常重要。否则 lm head logits：
[
B\times S\times V
]
可能成为显存峰值和通信瓶颈。

\subsection{sequence parallelism}
\label{subsec:sequence-parallelism}

Tensor Parallelism 通常切分 hidden 维，但某些操作不适合 hidden 维切分，例如 LayerNorm、RMSNorm、Dropout、Residual Add 等。它们通常需要完整 hidden 向量，但可以沿 sequence 维切分。

Sequence Parallelism 的思想是：在 TP group 内，将激活沿 sequence 维分片，使每个 rank 只保存部分 token 的完整 hidden：

[
\tensor{X}\in\mathbb{R}^{B\times S\times H}
\rightarrow
\tensor{X}_p\in
\mathbb{R}^{B\times \frac{S}{P}\times H}.
]

这样，LayerNorm/RMSNorm 可以在本地执行，因为每个 token 的 hidden 维是完整的。相比所有 rank 都保存完整 sequence，激活显存降低约 $P$ 倍。

但是，当进入需要 hidden 维切分的 Column Parallel Linear 时，需要 all-gather sequence shards，让每个 rank 获得完整 sequence 的输入；当 Row Parallel Linear 输出后，可以用 reduce-scatter 再回到 sequence-sharded 状态。

\begin{figure}[H]
\centering
\begin{tikzpicture}[
font=\small,
tensor/.style={draw, rounded corners=3pt, minimum width=2.0cm, minimum height=0.75cm, align=center},
shard/.style={draw, rounded corners=2pt, minimum width=1.25cm, minimum height=0.5cm, align=center},
arrow/.style={-{Latex[length=2.1mm]}, thick}
]

```
\node[font=\bfseries] at (0,3.25) {Sequence Parallelism：沿序列维切分激活};

\node[tensor, fill=blue!6] (full) at (-5.0,1.1) {$X$\\$B,S,H$};
\node[tensor, fill=orange!12] (rs) at (-2.7,1.1) {ReduceScatter\\sequence dim};

\foreach \i/\y/\col in {0/1.8/red!12,1/1.1/green!12,2/0.4/purple!12} {
  \node[shard, fill=\col] (s\i) at (-0.55,\y) {$B,S/P,H$};
}

\node[tensor, fill=green!10] (norm) at (1.9,1.1) {Norm / Dropout\\local};
\node[tensor, fill=purple!10] (ag) at (4.5,1.1) {AllGather\\when needed};

\draw[arrow] (full) -- (rs);
\draw[arrow] (rs) -- (s0);
\draw[arrow] (rs) -- (s1);
\draw[arrow] (rs) -- (s2);
\draw[arrow] (s0) -- (norm);
\draw[arrow] (s1) -- (norm);
\draw[arrow] (s2) -- (norm);
\draw[arrow] (norm) -- (ag);

\node[align=left, text width=10.4cm] at (0,-1.45) {
  Sequence Parallel 将非 TP 线性层之间的激活按 sequence 维切分。
  这样 LayerNorm/RMSNorm 等需要完整 hidden 的操作仍可本地执行，同时降低激活显存。
};
```

\end{tikzpicture}
\caption{Sequence Parallelism 的激活切分方式}
\label{fig:sequence-parallelism}
\end{figure}

Sequence Parallelism 的收益主要体现在训练显存，尤其是长序列、大 hidden 和大 TP size 场景。但它也引入更多 all-gather/reduce-scatter，对通信实现提出更高要求。

\subsection{activation memory 优化}
\label{subsec:activation-memory-optimization-in-tp}

张量并行降低了参数和部分激活的显存，但训练中的 activation memory 仍可能很大。常见优化包括：

\begin{itemize}
\item \textbf{Activation Checkpointing}：少存激活，backward 时重算；
\item \textbf{Sequence Parallelism}：沿 sequence 维切分非 tensor-parallel 区间的激活；
\item \textbf{Selective Recomputation}：只重算 attention softmax、MLP activation 等高显存中间结果；
\item \textbf{FlashAttention}：避免物化完整 attention matrix；
\item \textbf{Fused Kernels}：减少中间张量写回；
\item \textbf{Mixed Precision}：使用 BF16/FP16/FP8 降低激活字节数。
\end{itemize}

一个 Transformer block 的训练显存与 TP 的关系并不是简单的 $1/P$。参数和部分线性层激活会随 TP 切分下降，但 residual stream、norm 输入、dropout mask、attention 中间结果和临时 buffer 不一定都按 TP size 下降。因此工程上需要 profiler 或显存估算工具逐项分析。

可以把 TP 后的每卡显存粗略写为：
[
M_{\mathrm{per\ GPU}}
\approx
\frac{M_{\mathrm{tp\ shardable}}}{P_{\mathrm{TP}}}
+
M_{\mathrm{replicated}}
+
M_{\mathrm{comm\ buffers}}.
]
其中：

\begin{itemize}
\item $M_{\mathrm{tp\ shardable}}$ 是可随 TP 切分的参数、梯度、优化器状态和部分激活；
\item $M_{\mathrm{replicated}}$ 是仍然复制的状态；
\item $M_{\mathrm{comm\ buffers}}$ 是 all-gather、all-reduce、reduce-scatter 等通信需要的临时 buffer。
\end{itemize}

TP size 不是越大越好。随着 TP size 增大：

\begin{itemize}
\item 每卡参数和计算下降；
\item 通信频率和相对开销上升；
\item GEMM 变小，Tensor Core 利用率可能下降；
\item rank mapping 更依赖高速互联；
\item 某些 head 数、hidden size、intermediate size 必须能整除 TP size。
\end{itemize}

\section{张量并行的性能权衡}
\label{sec:tensor-parallel-performance-tradeoff}

\subsection{计算量下降与通信量上升}
\label{subsec:compute-down-communication-up}

张量并行将一个大 GEMM 切到 $P$ 个 rank 上。理想情况下，每个 rank 的 GEMM 计算量约下降为：
[
\frac{1}{P}.
]
但每个 Transformer block 通常会引入若干次 TP collective communication。例如在基础 Megatron 风格中，每层 attention output projection 后有一次 all-reduce，FFN down projection 后有一次 all-reduce。因此每层至少有两次 TP 通信。

如果单层计算时间为：
[
T_{\mathrm{compute}}(P)
\approx
\frac{T_{\mathrm{compute}}(1)}{P\eta_{\mathrm{gemm}}(P)},
]
通信时间为：
[
T_{\mathrm{comm}}(P),
]
则 TP 后单层时间近似为：
[
T_{\mathrm{layer}}(P)
\approx
T_{\mathrm{compute}}(P)+T_{\mathrm{comm}}(P).
]

当 $P$ 较小时，计算下降带来收益；当 $P$ 过大时，GEMM 变小、通信占比上升，收益可能下降甚至反转。

\begin{figure}[H]
\centering
\begin{tikzpicture}[
font=\small,
axis/.style={-{Latex[length=2.2mm]}, thick},
curve/.style={thick}
]

```
\node[font=\bfseries] at (0,3.0) {TP Size 增大时计算与通信的权衡};

\draw[axis] (-5.3,-1.2) -- (5.3,-1.2) node[right] {TP size};
\draw[axis] (-5.3,-1.2) -- (-5.3,2.2) node[above] {relative time};

\draw[curve, blue!75]
  plot[smooth] coordinates {
    (-4.8,1.8) (-3.2,1.05) (-1.6,0.55) (0,0.25) (1.6,0.12) (3.2,0.08) (4.8,0.06)
  };
\node[blue!75] at (3.2,0.45) {compute per rank};

\draw[curve, orange!90!black]
  plot[smooth] coordinates {
    (-4.8,-0.75) (-3.2,-0.6) (-1.6,-0.35) (0,0.05) (1.6,0.55) (3.2,1.15) (4.8,1.9)
  };
\node[orange!90!black] at (3.2,1.6) {communication overhead};

\draw[curve, purple!80]
  plot[smooth] coordinates {
    (-4.8,1.2) (-3.2,0.55) (-1.6,0.18) (0,0.0) (1.6,0.18) (3.2,0.7) (4.8,1.45)
  };
\node[purple!80] at (0.9,-0.35) {total layer time};

\node[align=left, text width=10.2cm] at (0,-2.35) {
  TP size 过小无法解决显存和计算瓶颈；TP size 过大则通信和小 GEMM 开销上升。
  合适的 TP size 通常受节点内高速互联、head 数、hidden size 和模型规模共同约束。
};
```

\end{tikzpicture}
\caption{张量并行中计算收益与通信开销的权衡}
\label{fig:tp-compute-communication-tradeoff}
\end{figure}

\subsection{节点内优先原则}
\label{subsec:intra-node-first-principle}

TP 通信发生在每个 Transformer block 内部，频率很高。因此 TP group 通常应尽量放在同一节点内，利用 NVLink/NVSwitch 等高速互联。

常见经验是：

\begin{itemize}
\item 单节点 8 GPU 且有 NVSwitch 时，TP size 可选 2、4 或 8；
\item 尽量避免 TP group 跨低带宽 PCIe 或跨节点网络；
\item 跨节点通信更适合 DP 或 PP 维度；
\item rank mapping 应与物理拓扑对齐；
\item 使用 NCCL tests 验证实际 all-reduce、all-gather、reduce-scatter 带宽。
\end{itemize}

如果 TP group 跨节点，层内通信会频繁走 RDMA 网络，step time 可能显著上升。除非模型极大、单节点 TP 不够，否则通常优先保持 TP 在节点内。

\subsection{TP 与 batch size}
\label{subsec:tp-and-batch-size}

TP 切分模型后，每个 rank 的 GEMM 维度会变小。如果 micro-batch 或 sequence length 也较小，GEMM 可能无法充分利用 Tensor Core。矩阵乘法形状大致为：
[
(B S)\times H
\quad
\text{乘}
\quad
H\times \frac{H_{\mathrm{out}}}{P}.
]
当 $B S$ 很小或 $H_{\mathrm{out}}/P$ 太小，GEMM 变成小矩阵，吞吐下降。

因此 TP 配置要与 micro-batch、sequence length、hidden size 和 FFN intermediate size 一起考虑。不能只看显存是否足够，还要看 GEMM shape 是否高效。

\subsection{TP 与数值一致性}
\label{subsec:tp-numerical-consistency}

张量并行改变了规约顺序。浮点数加法不满足严格结合律：
[
(a+b)+c \ne a+(b+c)
]
在有限精度下可能产生微小差异。因此 TP size 不同、all-reduce 算法不同、rank 数不同，训练结果可能不完全 bitwise identical。

通常这类差异很小，不影响模型收敛。但在调试时需要注意：

\begin{itemize}
\item 不同 TP size 下 loss 曲线可能有微小差异；
\item mixed precision 和不同规约顺序会放大数值差异；
\item checkpoint 在不同 TP size 之间转换需要正确重组权重；
\item 单元测试应使用合理 tolerance，而不是强制逐 bit 相等。
\end{itemize}

\section{本章小结}
\label{sec:chapter08-summary}

本章系统介绍了张量并行的动机、基本原理、Transformer 内部切分方式以及 Megatron-LM 风格的并行设计。

\begin{itemize}
\item \textbf{数据并行的边界} 是每张 GPU 必须保存完整模型副本；当模型状态单卡放不下时，需要模型并行。
\item \textbf{张量并行} 将单个 Linear 层或 attention heads 切分到多个 GPU 上，让多张 GPU 共同计算一个 Transformer block。
\item \textbf{Column Parallel Linear} 按输出维切分权重，forward 可输出 sharded tensor，backward 通常需要对输入梯度 all-reduce。
\item \textbf{Row Parallel Linear} 按输入维切分权重，forward 需要对局部输出 all-reduce 求和，backward 通常自然得到 sharded 输入梯度。
\item \textbf{Column + Row 的组合} 可以避免中间 all-gather，只在输出处通信，是 Megatron-LM 张量并行的核心模式。
\item \textbf{Attention 中的 QKV projection} 通常按 head 做 Column Parallel，每个 rank 本地计算自己负责的 heads。
\item \textbf{Attention output projection} 通常做 Row Parallel，并在输出处 all-reduce 恢复 replicated residual stream。
\item \textbf{FFN 中的 up/gate projection} 通常做 Column Parallel，down projection 做 Row Parallel，适合大中间维度的 Transformer MLP。
\item \textbf{Tensor model parallel group} 定义了哪些 rank 共同切分一个模型副本；它与 data parallel group 正交。
\item \textbf{Vocabulary parallelism} 将词表维切分到 TP ranks 上，降低 embedding/lm head 显存，并可配合 parallel cross entropy。
\item \textbf{Sequence Parallelism} 沿 sequence 维切分激活，降低 LayerNorm、Dropout、Residual 等区间的 activation memory。
\item \textbf{TP size 不是越大越好}：计算量下降的同时通信开销上升，GEMM shape 变小，拓扑和 rank mapping 更关键。
\end{itemize}

下一章将进入流水线并行 Pipeline Parallelism。我们会分析为什么仅靠张量并行仍无法解决所有问题，如何把模型按层切成多个 pipeline stages，GPipe、1F1B 和 interleaved schedule 如何减少 pipeline bubble，以及 micro-batch 数量如何影响吞吐、显存和调度复杂度。

\chapter{流水线并行 PP}
\label{chap:pipeline-parallelism}

\section{Pipeline Parallelism 的动机}
\label{sec:pipeline-parallelism-motivation}

上一章我们讨论了张量并行。张量并行把一个 Transformer block 内部的 Linear、Attention heads、FFN 中间维度切到多张 GPU 上，使多张 GPU 共同完成一层内部的矩阵计算。但张量并行也有边界：当模型层数很多、总参数量巨大时，即使每一层内部已经被切分，单个 GPU 仍可能无法承担全部层的参数、优化器状态和激活。

这时需要另一类并行方式：\textbf{流水线并行 Pipeline Parallelism, PP}。它的核心思想是：

[
\text{把模型按层切分到不同设备上，让不同 micro-batch 像流水线一样穿过这些 stage。}
]

如果说数据并行是“复制模型、切分数据”，张量并行是“切分层内矩阵”，那么流水线并行就是“切分模型层”。它特别适用于深层 Transformer，因为 Transformer 是由许多相似 block 堆叠而成的结构：

[
\mat{X}_{\ell+1}
================

\operatorname{DecoderBlock}*{\ell}(\mat{X}*{\ell}),
\quad
\ell=0,1,\ldots,L-1.
]

当 $L$ 很大时，可以把这些 block 分给多个 pipeline stages。例如 $L=80$，pipeline size 为 $P=8$，每个 stage 大约负责 $10$ 层。

\subsection{层级切分}
\label{subsec:layer-partitioning}

流水线并行最常见的切分方式是按层切分。设模型有 $L$ 层，pipeline stage 数为 $P$。若均匀切分，每个 stage 拥有：
[
\frac{L}{P}
]
层。第 $p$ 个 stage 负责的层集合可表示为：
[
\mathcal{L}_p
=============

\left{
\ell
\mid
p\frac{L}{P}
\le
\ell
<
(p+1)\frac{L}{P}
\right}.
]

如果 $L$ 不能被 $P$ 整除，则需要不均匀切分。实际系统中，层切分并不总是简单平均，因为不同层的计算量可能不同。例如：

\begin{itemize}
\item embedding 层和 lm head 参数量大，但计算形态不同；
\item 某些模型的前几层或后几层结构特殊；
\item MoE 层与 dense 层成本不同；
\item 使用 recomputation、sequence parallel 或 tensor parallel 后，不同 stage 显存和时间可能不均；
\item 最后一层可能包含 final norm 和 vocab parallel cross entropy。
\end{itemize}

流水线并行要求各 stage 时间尽量均衡。如果某个 stage 比其他 stage 慢，它就会成为整个流水线的瓶颈。设第 $p$ 个 stage 的单个 micro-batch 处理时间为 $t_p$，则流水线稳定阶段的节拍约由最慢 stage 决定：

[
t_{\mathrm{cycle}}
==================

\max_{p} t_p.
]

因此，Pipeline Parallelism 的第一项工程任务就是\textbf{合理切分层，使每个 stage 的计算时间和显存压力尽量接近}。

\subsection{stage}
\label{subsec:pipeline-stage}

Pipeline stage 是流水线并行中的基本计算单元。一个 stage 通常运行在一张 GPU 或一组 GPU 上。如果结合张量并行，一个 pipeline stage 可能不是单张 GPU，而是一个 TP group。例如：

[
\text{PP stage}
===============

\text{一组 TP ranks 共同负责若干层}.
]

在纯流水线并行中，stage 之间传递的是激活张量。前向传播时：

[
\mat{A}_{p+1}
=============

f_p(\mat{A}_p),
]
其中 $f_p$ 表示第 $p$ 个 stage 负责的若干层。第 $p$ 个 stage 计算完成后，将输出激活发送给第 $p+1$ 个 stage。

反向传播时，梯度沿相反方向传播。若第 $p+1$ 个 stage 传回激活梯度：
[
\frac{\partial \mathcal{L}}{\partial \mat{A}*{p+1}},
]
则第 $p$ 个 stage 根据本地保存的前向激活和参数，计算：
[
\frac{\partial \mathcal{L}}{\partial \mat{A}*{p}}
]
并发送给前一个 stage。

因此，流水线并行引入两类点对点通信：

\begin{itemize}
\item forward activation send/recv；
\item backward activation gradient send/recv。
\end{itemize}

这与数据并行和张量并行不同。DDP 主要是 collective all-reduce；TP 常用 all-reduce、all-gather、reduce-scatter；PP 主要是相邻 stage 之间的 point-to-point communication。

\begin{figure}[H]
\centering
\begin{tikzpicture}[
font=\small,
stage/.style={draw, rounded corners=4pt, minimum width=2.1cm, minimum height=1.0cm, align=center, fill=blue!6},
act/.style={draw, rounded corners=2pt, minimum width=1.0cm, minimum height=0.45cm, align=center, fill=orange!14},
arrow/.style={-{Latex[length=2.3mm]}, thick},
back/.style={-{Latex[length=2.3mm]}, thick, red!70}
]

```
\node[font=\bfseries] at (0,3.1) {模型按层切分为多个 Pipeline Stages};

\node[stage] (s0) at (-4.8,0.9) {Stage 0\\Layers 0--9};
\node[stage] (s1) at (-1.6,0.9) {Stage 1\\Layers 10--19};
\node[stage] (s2) at (1.6,0.9) {Stage 2\\Layers 20--29};
\node[stage] (s3) at (4.8,0.9) {Stage 3\\Layers 30--39};

\node[act] (a0) at (-6.4,0.9) {$A_0$};
\node[act] (a1) at (-3.2,0.9) {$A_1$};
\node[act] (a2) at (0.0,0.9) {$A_2$};
\node[act] (a3) at (3.2,0.9) {$A_3$};
\node[act] (a4) at (6.4,0.9) {$A_4$};

\draw[arrow] (a0) -- (s0);
\draw[arrow] (s0) -- (a1);
\draw[arrow] (a1) -- (s1);
\draw[arrow] (s1) -- (a2);
\draw[arrow] (a2) -- (s2);
\draw[arrow] (s2) -- (a3);
\draw[arrow] (a3) -- (s3);
\draw[arrow] (s3) -- (a4);

\draw[back] (a4.south) .. controls (5.6,-0.6) and (4.8,-0.6) .. (s3.south);
\draw[back] (s3.south) .. controls (4.0,-0.95) and (3.2,-0.95) .. (a3.south);
\draw[back] (a3.south) .. controls (2.4,-1.1) and (1.6,-1.1) .. (s2.south);
\draw[back] (s2.south) .. controls (0.8,-1.25) and (0.0,-1.25) .. (a2.south);
\draw[back] (a2.south) .. controls (-0.8,-1.4) and (-1.6,-1.4) .. (s1.south);
\draw[back] (s1.south) .. controls (-2.4,-1.55) and (-3.2,-1.55) .. (a1.south);
\draw[back] (a1.south) .. controls (-4.0,-1.7) and (-4.8,-1.7) .. (s0.south);

\node[align=left, text width=10.7cm] at (0,-2.55) {
  前向传播沿 stage 从左到右传递激活；反向传播沿相反方向传递激活梯度。
  PP 的通信主要发生在相邻 stage 之间，是点对点 send/recv。
};
```

\end{tikzpicture}
\caption{模型按层切分为多个 pipeline stages}
\label{fig:pipeline-stage-layer-partition}
\end{figure}

\subsection{micro-batch}
\label{subsec:micro-batch}

如果只把一个 batch 直接送入流水线，会发生严重空闲。考虑 $P=4$ 个 stages，一个 batch 需要依次经过 Stage 0、Stage 1、Stage 2、Stage 3。若没有切分 micro-batch，则同一时刻只有一个 stage 在工作，其他 stage 都在等待。

为了解决这个问题，流水线并行会把一个 mini-batch 切成多个 micro-batches。设一个训练 step 的 batch 被切成 $M$ 个 micro-batches：

[
\mathcal{B}
===========

{\mathcal{B}_0,\mathcal{B}*1,\ldots,\mathcal{B}*{M-1}}.
]

每个 micro-batch 都会依次经过所有 stages。这样，当 Stage 1 正在处理 micro-batch 0 时，Stage 0 可以处理 micro-batch 1；当 Stage 2 处理 micro-batch 0 时，Stage 1 处理 micro-batch 1，Stage 0 处理 micro-batch 2。多个 micro-batches 填满流水线后，各 stage 可以同时工作。

Global batch size 与 pipeline micro-batch 的关系为：

[
B_{\mathrm{global}}
===================

B_{\mathrm{micro}}
\cdot
M
\cdot
P_{\mathrm{DP}},
]
其中 $B_{\mathrm{micro}}$ 是每个 micro-batch 的 batch size，$M$ 是每个 step 的 micro-batch 数，也常被称为 gradient accumulation steps，$P_{\mathrm{DP}}$ 是数据并行副本数。

需要注意，流水线并行中的 micro-batch 与数据并行中的 local batch、梯度累积密切相关。一个 optimizer step 通常要等待所有 micro-batches 的 forward 和 backward 都完成，然后聚合梯度并更新参数。

\subsection{pipeline bubble}
\label{subsec:pipeline-bubble}

流水线并行的核心问题是 pipeline bubble，即流水线无法完全填满时产生的空闲时间。即使使用 micro-batch，开始阶段需要 warmup，结束阶段需要 drain/cooldown。

以 $P$ 个 stages、$M$ 个 micro-batches 为例。在最简单的流水线中，至少需要 $P-1$ 个时间片填满流水线，也需要 $P-1$ 个时间片排空流水线。因此 bubble 与 $P$ 有关，而可用计算量与 $M$ 有关。粗略地，micro-batch 数越多，bubble 占比越低。

对于简单 GPipe all-forward-all-backward 调度，流水线效率可粗略写作：

[
\eta_{\mathrm{pipe}}
\approx
\frac{M}{M+P-1}.
]

这个公式表达了一个重要直觉：

[
M\gg P
\quad
\Rightarrow
\quad
\eta_{\mathrm{pipe}}\rightarrow 1.
]

也就是说，micro-batch 数越多，流水线越容易被填满，bubble 占比越低。但 $M$ 变大也会带来代价：

\begin{itemize}
\item 每个 optimizer step 的 micro-batch 数增加，step 内调度更复杂；
\item GPipe 中需要保存更多 micro-batch 的激活，显存上升；
\item micro-batch 太小会导致 GEMM shape 变小，计算效率下降；
\item pipeline flush 时间和通信调度变复杂；
\item 与数据并行、张量并行组合时，global batch size 可能过大。
\end{itemize}

\begin{figure}[H]
\centering
\begin{tikzpicture}[
font=\small,
cell/.style={draw, minimum width=0.7cm, minimum height=0.45cm, align=center},
work/.style={fill=blue!14},
bubble/.style={fill=gray!18},
axis/.style={-{Latex[length=2.2mm]}, thick}
]

```
\node[font=\bfseries] at (0,3.1) {Pipeline Bubble：warmup 与 cooldown 中的空闲};

\foreach \t in {0,...,7} {
  \node[font=\scriptsize] at ({-3.5+0.75*\t},2.35) {t\t};
}

\foreach \s/\y in {0/1.7,1/1.15,2/0.6,3/0.05} {
  \node[font=\scriptsize] at (-4.35,\y) {Stage \s};
}

% Stage 0
\foreach \t/\m in {0/0,1/1,2/2,3/3} {
  \node[cell, work] at ({-3.5+0.75*\t},1.7) {$F\m$};
}
\foreach \t in {4,...,7} {
  \node[cell, bubble] at ({-3.5+0.75*\t},1.7) {};
}

% Stage 1
\node[cell, bubble] at (-3.5,1.15) {};
\foreach \t/\m in {1/0,2/1,3/2,4/3} {
  \node[cell, work] at ({-3.5+0.75*\t},1.15) {$F\m$};
}
\foreach \t in {5,...,7} {
  \node[cell, bubble] at ({-3.5+0.75*\t},1.15) {};
}

% Stage 2
\foreach \t in {0,1} {
  \node[cell, bubble] at ({-3.5+0.75*\t},0.6) {};
}
\foreach \t/\m in {2/0,3/1,4/2,5/3} {
  \node[cell, work] at ({-3.5+0.75*\t},0.6) {$F\m$};
}
\foreach \t in {6,7} {
  \node[cell, bubble] at ({-3.5+0.75*\t},0.6) {};
}

% Stage 3
\foreach \t in {0,1,2} {
  \node[cell, bubble] at ({-3.5+0.75*\t},0.05) {};
}
\foreach \t/\m in {3/0,4/1,5/2,6/3} {
  \node[cell, work] at ({-3.5+0.75*\t},0.05) {$F\m$};
}
\node[cell, bubble] at ({-3.5+0.75*7},0.05) {};

\node[draw, rounded corners=2pt, fill=blue!14, minimum width=0.55cm, minimum height=0.35cm] at (3.6,1.3) {};
\node[align=left] at (4.8,1.3) {有效计算};
\node[draw, rounded corners=2pt, fill=gray!18, minimum width=0.55cm, minimum height=0.35cm] at (3.6,0.7) {};
\node[align=left] at (4.8,0.7) {bubble 空闲};

\node[align=left, text width=10.3cm] at (0,-1.45) {
  流水线开始时，后面的 stage 还没有收到输入；结束时，前面的 stage 已经没有新 micro-batch 可处理。
  这些空闲时间就是 pipeline bubble。
};
```

\end{tikzpicture}
\caption{Pipeline warmup/cooldown 造成的 bubble}
\label{fig:pipeline-bubble}
\end{figure}

Pipeline Parallelism 的本质就是在模型切分带来的显存收益与 bubble、通信、调度复杂度之间做权衡。

\section{GPipe}
\label{sec:gpipe}

GPipe 是早期经典的流水线并行调度方式。它的基本策略是：对所有 micro-batches 先执行 forward，等所有 forward 完成后，再按反向顺序执行 backward。这种方式也被称为 all-forward-all-backward。

\subsection{micro-batch 调度}
\label{subsec:gpipe-microbatch-scheduling}

假设有 $P=4$ 个 pipeline stages，$M=4$ 个 micro-batches。GPipe forward 阶段中，micro-batch 0 先进入 Stage 0，下一时间片进入 Stage 1，同时 micro-batch 1 进入 Stage 0。随着时间推进，所有 micro-batches 依次穿过所有 stages。

Forward 完成后，最后一个 stage 开始对最后一个 micro-batch 做 backward，然后梯度逐 stage 向前传播。Backward 阶段与 forward 类似，但方向相反。

GPipe 的优点是概念简单、易于实现。缺点是需要在 backward 前保存所有 micro-batches 的前向激活，因此显存压力较大。

\subsection{all-forward-all-backward}
\label{subsec:all-forward-all-backward}

GPipe 的调度可以写作：

[
F_{0,0},F_{0,1},\ldots,F_{0,P-1},
]
表示 micro-batch 0 从 Stage 0 到 Stage $P-1$ 的 forward；对所有 micro-batches 重复。之后再执行 backward：

[
B_{M-1,P-1},B_{M-1,P-2},\ldots,B_{M-1,0}.
]

更直观地看，它先填充并完成所有 forward，然后再排空 backward。这意味着在 forward 结束时，各 stage 保存了大量 micro-batch 激活，直到对应 backward 使用完毕。

\begin{figure}[H]
\centering
\begin{tikzpicture}[
font=\scriptsize,
cell/.style={draw, minimum width=0.62cm, minimum height=0.42cm, align=center},
fwd/.style={fill=blue!14},
bwd/.style={fill=red!14},
bubble/.style={fill=gray!18}
]

```
\node[font=\bfseries\small] at (0,3.0) {GPipe All-Forward-All-Backward 时间线};

\foreach \t in {0,...,13} {
  \node at ({-4.6+0.65*\t},2.35) {\t};
}

\foreach \s/\y in {0/1.75,1/1.2,2/0.65,3/0.1} {
  \node[font=\scriptsize] at (-5.35,\y) {S\s};
}

% initialize all cells grey
\foreach \s/\y in {0/1.75,1/1.2,2/0.65,3/0.1} {
  \foreach \t in {0,...,13} {
    \node[cell, bubble] at ({-4.6+0.65*\t},\y) {};
  }
}

% Forward schedule M=4, P=4
\foreach \m in {0,...,3} {
  \foreach \s in {0,...,3} {
    \pgfmathtruncatemacro{\t}{\m+\s}
    \pgfmathsetmacro{\y}{1.75-0.55*\s}
    \node[cell, fwd] at ({-4.6+0.65*\t},\y) {$F_{\m}$};
  }
}

% Backward schedule after forward drain, start t=7
\foreach \m/\idx in {3/0,2/1,1/2,0/3} {
  \foreach \s in {3,2,1,0} {
    \pgfmathtruncatemacro{\stageoffset}{3-\s}
    \pgfmathtruncatemacro{\t}{7+\idx+\stageoffset}
    \pgfmathsetmacro{\y}{1.75-0.55*\s}
    \node[cell, bwd] at ({-4.6+0.65*\t},\y) {$B_{\m}$};
  }
}

\node[draw, rounded corners=2pt, fill=blue!14, minimum width=0.5cm, minimum height=0.32cm] at (3.9,1.5) {};
\node[anchor=west] at (4.3,1.5) {Forward};
\node[draw, rounded corners=2pt, fill=red!14, minimum width=0.5cm, minimum height=0.32cm] at (3.9,1.0) {};
\node[anchor=west] at (4.3,1.0) {Backward};
\node[draw, rounded corners=2pt, fill=gray!18, minimum width=0.5cm, minimum height=0.32cm] at (3.9,0.5) {};
\node[anchor=west] at (4.3,0.5) {Bubble};

\node[align=left, text width=10.4cm, font=\small] at (0,-1.35) {
  GPipe 先完成所有 micro-batch 的 forward，再执行所有 backward。
  调度简单，但 forward 激活需要保存到 backward 阶段，显存压力较高。
};
```

\end{tikzpicture}
\caption{GPipe 的 all-forward-all-backward 调度}
\label{fig:gpipe-timeline}
\end{figure}

\subsection{bubble 分析}
\label{subsec:gpipe-bubble-analysis}

对于 $P$ 个 stages、$M$ 个 micro-batches，GPipe forward 阶段需要：
[
M+P-1
]
个时间片完成。Backward 阶段也需要：
[
M+P-1
]
个时间片。因此总时间片约为：
[
T_{\mathrm{GPipe}}
==================

2(M+P-1).
]

真正做有效计算的 forward 和 backward 各有 $MP$ 个 stage-microbatch 任务，总有效任务数为：
[
2MP.
]
如果理想情况下每个时间片每个 stage 都能工作，总容量为：
[
2P(M+P-1).
]
因此效率可近似写为：
[
\eta_{\mathrm{GPipe}}
=====================

# \frac{2MP}{2P(M+P-1)}

\frac{M}{M+P-1}.
]

这就是前面提到的 bubble 公式。

当 $M=1$ 时：
[
\eta_{\mathrm{GPipe}}
=====================

\frac{1}{P}.
]
流水线完全没有被填满，效率很低。

当 $M\gg P$ 时：
[
\eta_{\mathrm{GPipe}}
\approx 1.
]
bubble 占比下降。

但实际中不能无限增大 $M$，因为 micro-batch 数越多，GPipe 需要保存的激活越多。若每个 micro-batch 在每个 stage 保存激活大小为 $A_{\mathrm{stage}}$，则每个 stage 需要保存的激活粗略与 $M$ 成正比：
[
M_{\mathrm{act}}^{\mathrm{GPipe}}
\propto
M A_{\mathrm{stage}}.
]

\subsection{激活保存问题}
\label{subsec:gpipe-activation-saving}

GPipe 的显存问题来自 all-forward-all-backward。所有 micro-batches 的 forward 先完成，backward 尚未开始，因此 stage 必须保存每个 micro-batch 的必要激活。

对于某个 stage，若它处理 $L_p$ 层，每层激活大小约为：
[
B_{\mathrm{micro}}S H s_{\mathrm{dtype}},
]
那么保存 $M$ 个 micro-batches 的激活显存大致与：
[
M L_p B_{\mathrm{micro}} S H s_{\mathrm{dtype}}
]
成正比。实际还包括 attention softmax、MLP 中间激活、dropout mask、norm 中间结果等。

为了缓解这个问题，GPipe 通常与 activation checkpointing 搭配。即每个 stage 内部不保存全部激活，而是在 backward 时重算部分 forward。这样可以显著降低显存，但增加计算量。

下面给出一个简化的 GPipe 调度生成示意代码：

\begin{lstlisting}[language=Python,caption={生成 GPipe all-forward-all-backward 调度的简化示意},label={lst:gpipe-schedule-generator}]
def build_gpipe_schedule(num_stages, num_microbatches):
schedule = []

```
# Forward phase
for t in range(num_microbatches + num_stages - 1):
    tasks = []
    for stage in range(num_stages):
        microbatch = t - stage
        if 0 <= microbatch < num_microbatches:
            tasks.append((stage, "F", microbatch))
    schedule.append(tasks)

# Backward phase
for t in range(num_microbatches + num_stages - 1):
    tasks = []
    for stage in reversed(range(num_stages)):
        # Run backward in reverse stage order.
        microbatch = (num_microbatches - 1) - (t - (num_stages - 1 - stage))
        if 0 <= microbatch < num_microbatches:
            tasks.append((stage, "B", microbatch))
    schedule.append(tasks)

return schedule
```

\end{lstlisting}

这个调度易懂，但不是现代大模型训练中最常见的高效方式。为了降低激活显存并提升 overlap，1F1B 调度更常用。

\section{1F1B Pipeline Schedule}
\label{sec:one-forward-one-backward}

1F1B 是 one-forward-one-backward 的缩写。它的核心思想是：流水线填满后，每个 stage 尽量交替执行一个 forward 和一个 backward，而不是先把所有 forward 做完再做所有 backward。

相比 GPipe，1F1B 的主要优点是减少激活驻留时间。一个 micro-batch 的 forward 完成后，不必等所有 micro-batches 都 forward 完成才 backward。随着 backward 更早开始，部分激活可以更早释放。

1F1B 通常分为三个阶段：

\begin{itemize}
\item \textbf{warmup}：前几个 micro-batches 进入流水线，尚无 backward 可执行；
\item \textbf{steady state}：每个 stage 交替执行 forward 和 backward；
\item \textbf{cooldown}：没有新的 forward，只剩 backward 排空流水线。
\end{itemize}

\subsection{warmup}
\label{subsec:onef1b-warmup}

在 1F1B 中，不同 stage 的 warmup 长度不同。靠前的 stage 需要先产生足够多的 forward activation，让后面的 stage 能够开始 backward。靠后的 stage 较早开始 backward。

对于 $P$ 个 stages，第 $p$ 个 stage 的 warmup micro-batch 数通常与它距离最后 stage 的距离有关。越靠前的 stage，需要更多 warmup forward；越靠后的 stage，warmup 更短。

直观地说，Stage 0 必须先不断发送 micro-batches，直到最后 stage 完成某个 micro-batch 的 forward 并开始 backward，反向梯度再逐 stage 传回来。只有此时 Stage 0 才有 backward 可以执行。

\subsection{steady state}
\label{subsec:onef1b-steady-state}

进入 steady state 后，每个 stage 大致执行：

[
F, B, F, B, F, B,\ldots
]

或者根据实现细节交错执行。这样做有两个收益：

\begin{itemize}
\item backward 更早开始，激活更早释放；
\item stage 的计算更均匀，减少 GPipe 中 forward 与 backward 完全分离带来的激活峰值。
\end{itemize}

1F1B 并不会消除 bubble，但它能显著改善显存。对于大模型训练，这通常比 GPipe 更实用。

\subsection{cooldown}
\label{subsec:onef1b-cooldown}

当所有 micro-batches 都已经完成 forward 后，流水线进入 cooldown。此时不再有新的 forward 输入，剩余任务是把所有 micro-batches 的 backward 执行完。靠后的 stage 较早完成，靠前的 stage 最后完成 backward。

Cooldown 与 warmup 一样会产生 bubble，但当 micro-batch 数足够大时，bubble 占比可以被摊薄。

\begin{figure}[H]
\centering
\begin{tikzpicture}[
font=\scriptsize,
cell/.style={draw, minimum width=0.62cm, minimum height=0.42cm, align=center},
fwd/.style={fill=blue!14},
bwd/.style={fill=red!14},
bubble/.style={fill=gray!18}
]

```
\node[font=\bfseries\small] at (0,3.0) {1F1B Pipeline 时间线示意};

\foreach \t in {0,...,13} {
  \node at ({-4.6+0.65*\t},2.35) {\t};
}

\foreach \s/\y in {0/1.75,1/1.2,2/0.65,3/0.1} {
  \node[font=\scriptsize] at (-5.35,\y) {S\s};
}

\foreach \s/\y in {0/1.75,1/1.2,2/0.65,3/0.1} {
  \foreach \t in {0,...,13} {
    \node[cell, bubble] at ({-4.6+0.65*\t},\y) {};
  }
}

% A hand-crafted 1F1B-like timeline for P=4, M=6
% Stage 0
\foreach \t/\op/\m/\style in {
  0/F/0/fwd,1/F/1/fwd,2/F/2/fwd,3/F/3/fwd,4/B/0/bwd,5/F/4/fwd,6/B/1/bwd,7/F/5/fwd,8/B/2/bwd,9/B/3/bwd,10/B/4/bwd,11/B/5/bwd
} {
  \node[cell, \style] at ({-4.6+0.65*\t},1.75) {$\op_{\m}$};
}

% Stage 1
\foreach \t/\op/\m/\style in {
  1/F/0/fwd,2/F/1/fwd,3/F/2/fwd,4/B/0/bwd,5/F/3/fwd,6/B/1/bwd,7/F/4/fwd,8/B/2/bwd,9/F/5/fwd,10/B/3/bwd,11/B/4/bwd,12/B/5/bwd
} {
  \node[cell, \style] at ({-4.6+0.65*\t},1.2) {$\op_{\m}$};
}

% Stage 2
\foreach \t/\op/\m/\style in {
  2/F/0/fwd,3/F/1/fwd,4/B/0/bwd,5/F/2/fwd,6/B/1/bwd,7/F/3/fwd,8/B/2/bwd,9/F/4/fwd,10/B/3/bwd,11/F/5/fwd,12/B/4/bwd,13/B/5/bwd
} {
  \node[cell, \style] at ({-4.6+0.65*\t},0.65) {$\op_{\m}$};
}

% Stage 3
\foreach \t/\op/\m/\style in {
  3/F/0/fwd,4/B/0/bwd,5/F/1/fwd,6/B/1/bwd,7/F/2/fwd,8/B/2/bwd,9/F/3/fwd,10/B/3/bwd,11/F/4/fwd,12/B/4/bwd,13/F/5/fwd
} {
  \node[cell, \style] at ({-4.6+0.65*\t},0.1) {$\op_{\m}$};
}

\node[draw, rounded corners=2pt, fill=blue!14, minimum width=0.5cm, minimum height=0.32cm] at (3.9,1.5) {};
\node[anchor=west] at (4.3,1.5) {Forward};
\node[draw, rounded corners=2pt, fill=red!14, minimum width=0.5cm, minimum height=0.32cm] at (3.9,1.0) {};
\node[anchor=west] at (4.3,1.0) {Backward};
\node[draw, rounded corners=2pt, fill=gray!18, minimum width=0.5cm, minimum height=0.32cm] at (3.9,0.5) {};
\node[anchor=west] at (4.3,0.5) {Bubble};

\node[align=left, text width=10.4cm, font=\small] at (0,-1.4) {
  1F1B 在 warmup 后交替执行 forward 与 backward。
  相比 GPipe，它能让激活更早释放，因此更适合显存紧张的大模型训练。
};
```

\end{tikzpicture}
\caption{1F1B pipeline schedule 的 warmup、steady state 与 cooldown}
\label{fig:onef1b-timeline}
\end{figure}

\subsection{与 GPipe 的对比}
\label{subsec:gpipe-vs-onef1b}

GPipe 和 1F1B 都使用 micro-batch 填充流水线，但它们的显存和调度特征不同。

\begin{table}[H]
\centering
\small
\caption{GPipe 与 1F1B Pipeline Schedule 对比}
\label{tab:gpipe-vs-1f1b}
\begin{tabular}{p{0.2\textwidth}p{0.34\textwidth}p{0.34\textwidth}}
\toprule
\textbf{维度} & \textbf{GPipe} & \textbf{1F1B} \
\midrule
调度方式
& 所有 forward 完成后再执行所有 backward
& warmup 后 forward/backward 交替执行 \
\midrule
实现复杂度
& 较简单
& 较复杂，需要更精细的 send/recv 调度 \
\midrule
激活显存
& 较高，需要保存所有 micro-batch 激活
& 较低，部分 micro-batch 可更早 backward 并释放激活 \
\midrule
pipeline bubble
& 存在，micro-batch 数越多占比越低
& 也存在，但 steady state 更适合长时间运行 \
\midrule
大模型适用性
& 概念清晰，适合教学和部分场景
& 更常用于大规模模型训练 \
\bottomrule
\end{tabular}
\end{table}

1F1B 的一个关键点是梯度累积。由于一个 optimizer step 包含多个 micro-batches，每个 stage 在完成所有 micro-batches 的 backward 之前不能更新参数。否则不同 micro-batch 会看到不同版本的参数，破坏同步语义。

因此，1F1B 中通常存在一个参数版本一致性要求：

[
\theta_{\mathrm{used\ by\ all\ microbatches}}
=============================================

\theta^{(t)}.
]
直到本 step 所有 micro-batches backward 完成后，才执行：

[
\theta^{(t+1)}
==============

\operatorname{OptimizerStep}
\left(
\theta^{(t)},\nabla\theta
\right).
]

\section{Interleaved Pipeline}
\label{sec:interleaved-pipeline}

随着模型越来越大，单纯把连续层分配给每个 physical stage 可能仍然存在 bubble 和负载不均问题。Interleaved Pipeline 使用 virtual pipeline stage，把每个物理设备上的层再切成多个虚拟段，交错执行，以减少 bubble 并改善负载均衡。

\subsection{virtual pipeline stage}
\label{subsec:virtual-pipeline-stage}

假设有 $P$ 个 physical pipeline stages，每个 stage 原本负责连续的一段层。如果引入 virtual pipeline size $V$，则每个 physical stage 被拆成 $V$ 个 virtual stages。总 virtual stages 数为：
[
P_{\mathrm{virtual}}=P\cdot V.
]

例如模型有 $32$ 层，physical PP size 为 $4$，每个 stage 原本负责 $8$ 层：

[
\text{Stage 0}: L0-L7,
\quad
\text{Stage 1}: L8-L15,
]
[
\text{Stage 2}: L16-L23,
\quad
\text{Stage 3}: L24-L31.
]

若 virtual pipeline size 为 $2$，则每个 physical stage 可以拥有两段非连续或交错的层：

[
\text{GPU0}: L0-L3,\ L16-L19,
]
[
\text{GPU1}: L4-L7,\ L20-L23,
]
[
\text{GPU2}: L8-L11,\ L24-L27,
]
[
\text{GPU3}: L12-L15,\ L28-L31.
]

这种布局使同一物理 GPU 在流水线中扮演多个 virtual stage，从而缩短每个 virtual stage 的计算时间，使流水线粒度更细。

\begin{figure}[H]
\centering
\begin{tikzpicture}[
font=\small,
gpu/.style={draw, rounded corners=4pt, minimum width=2.1cm, minimum height=1.35cm, align=center, fill=blue!6},
vstage/.style={draw, rounded corners=2pt, minimum width=1.55cm, minimum height=0.42cm, align=center, fill=orange!14},
arrow/.style={-{Latex[length=2.2mm]}, thick}
]

```
\node[font=\bfseries] at (0,3.2) {Interleaved Pipeline：一个物理 GPU 包含多个 Virtual Stages};

\foreach \i/\x in {0/-4.8,1/-1.6,2/1.6,3/4.8} {
  \node[gpu] (g\i) at (\x,0.9) {GPU \i};
  \node[vstage] (a\i) at (\x,1.25) {Virtual Stage A};
  \node[vstage] (b\i) at (\x,0.55) {Virtual Stage B};
}

\draw[arrow] (a0) -- (a1);
\draw[arrow] (a1) -- (a2);
\draw[arrow] (a2) -- (a3);
\draw[arrow] (a3.south) .. controls (3.6,-0.55) and (-3.6,-0.55) .. (b0.south);
\draw[arrow] (b0) -- (b1);
\draw[arrow] (b1) -- (b2);
\draw[arrow] (b2) -- (b3);

\node[align=left, text width=10.5cm] at (0,-1.65) {
  每个物理 GPU 负责多个 virtual stages。模型层在 virtual pipeline 中更细粒度流动，
  可以降低 bubble，但调度、通信和激活管理更复杂。
};
```

\end{tikzpicture}
\caption{Interleaved Pipeline 中的 virtual pipeline stage}
\label{fig:interleaved-virtual-pipeline-stage}
\end{figure}

\subsection{减少 bubble}
\label{subsec:interleaved-reduce-bubble}

Interleaving 的核心收益是减少 pipeline bubble。直觉上，virtual stage 数增多后，流水线更细，单个 stage 的处理时间更短，填充和排空的粒度更细。对于相同 physical PP size，interleaved schedule 可以提高 pipeline 利用率。

但收益不是免费的。Interleaving 会增加：

\begin{itemize}
\item stage 间激活传递次数；
\item send/recv 调度复杂度；
\item 每个 GPU 上多个 virtual stage 的激活管理；
\item layer 到 device 的映射复杂度；
\item checkpoint 保存和恢复的映射复杂度。
\end{itemize}

如果 virtual stage 过多，单个 virtual stage 的计算太小，通信和调度开销可能抵消 bubble 收益。因此 virtual pipeline size 也需要调参。

\subsection{调度复杂度}
\label{subsec:interleaved-scheduling-complexity}

Interleaved schedule 需要同时考虑：

\begin{itemize}
\item physical stage；
\item virtual stage；
\item micro-batch id；
\item forward/backward 方向；
\item 激活 send/recv；
\item 梯度 send/recv；
\item 与 tensor parallel communication 的重叠；
\item 与 data parallel gradient reduction 的关系。
\end{itemize}

一个训练 step 中，某个 GPU 可能在不同时间片执行不同 virtual stage 的 forward 或 backward。它还要保存每个 micro-batch 对应 virtual stage 的激活。若实现不当，很容易出现死锁、显存泄漏、激活错配或梯度错配。

下面给出一个抽象的 interleaved task 表示：

\begin{lstlisting}[language=Python,caption={Interleaved Pipeline 中任务描述的抽象表示},label={lst:interleaved-task-abstraction}]
from dataclasses import dataclass

@dataclass
class PipelineTask:
physical_stage: int
virtual_stage: int
microbatch: int
direction: str  # "forward" or "backward"

def run_task(task: PipelineTask):
if task.direction == "forward":
recv_activation_if_needed(task)
output = run_forward_layers(task)
send_activation_if_needed(task, output)
save_activation_for_backward(task, output)
else:
recv_grad_if_needed(task)
grad_input = run_backward_layers(task)
send_grad_if_needed(task, grad_input)
release_saved_activation(task)
\end{lstlisting}

真实系统会比这复杂得多。尤其是当 pipeline parallel 与 tensor parallel、data parallel、sequence parallel、activation checkpointing 同时存在时，每个任务可能同时涉及 point-to-point communication 和 collective communication。

\subsection{Megatron PP 实践}
\label{subsec:megatron-pipeline-practice}

Megatron-LM 的 3D 并行通常组合 TP、PP 和 DP。对于一个 rank，它同时属于：

[
\text{tensor model parallel group},
]
[
\text{pipeline model parallel group},
]
[
\text{data parallel group}.
]

Pipeline rank 决定它负责哪些层；Tensor parallel rank 决定它负责这些层中哪些 tensor shards；Data parallel rank 决定它属于哪个 replica。

如果 world size 为：
[
W=P_{\mathrm{TP}}\cdot P_{\mathrm{PP}}\cdot P_{\mathrm{DP}},
]
则总 rank 会被组织成三维网格。后续 3D 并行章节会详细讨论 rank mapping。这里先建立直觉：Pipeline Parallelism 不会替代 TP 或 DP，而是与它们组合使用。

在 Megatron 风格训练中，流水线并行实践需要考虑：

\begin{itemize}
\item 每个 pipeline stage 的层数是否均衡；
\item embedding 和 lm head 放在哪些 stage；
\item 是否使用 interleaved schedule；
\item micro-batch 数是否足够填满流水线；
\item activation checkpointing 的粒度；
\item stage 间通信是否跨节点；
\item 与 TP collective 是否发生资源竞争；
\item 与 DP gradient synchronization 的重叠方式。
\end{itemize}

\section{流水线并行的工程权衡}
\label{sec:pipeline-parallel-engineering-tradeoffs}

\subsection{显存收益}
\label{subsec:pp-memory-benefit}

流水线并行按层切分模型，因此每个 stage 只保存部分层的参数、梯度和优化器状态。若模型层均匀切分到 $P_{\mathrm{PP}}$ 个 stages，理想情况下每个 stage 的层相关参数显存约下降为：

[
\frac{1}{P_{\mathrm{PP}}}.
]

但并非所有显存都严格按 PP size 缩小。例如：

\begin{itemize}
\item embedding 和 lm head 可能集中在首尾 stage；
\item pipeline boundary 需要保存发送和接收 buffer；
\item micro-batch 数增加会影响激活数量；
\item GPipe 会保存大量 forward 激活；
\item 1F1B 降低激活驻留，但仍需要保存未 backward 的 micro-batch；
\item checkpointing 会用额外计算换显存。
\end{itemize}

可以粗略写作：

[
M_{\mathrm{per\ stage}}
\approx
\frac{M_{\mathrm{layer\ states}}}{P_{\mathrm{PP}}}
+
M_{\mathrm{stage\ activations}}
+
M_{\mathrm{pipeline\ buffers}}
+
M_{\mathrm{imbalanced}}.
]

其中 $M_{\mathrm{imbalanced}}$ 表示由不均匀层切分、embedding/lm head 或特殊层造成的额外显存不平衡。

\subsection{通信成本}
\label{subsec:pp-communication-cost}

PP 的通信主要是 stage boundary 的 activation 和 activation gradient。设 stage 间传递的激活 shape 为：
[
\mathbb{R}^{B_{\mathrm{micro}}\times S\times H}.
]
若使用 BF16，每个元素 2 bytes，则一次 forward 边界通信大小约为：
[
B_{\mathrm{micro}} S H \cdot 2\ \mathrm{bytes}.
]
Backward 还需要传回同样大小的梯度。因此每个 micro-batch 每条边界的双向通信约为：
[
2B_{\mathrm{micro}} S H s_{\mathrm{dtype}}.
]

对于 $P_{\mathrm{PP}}$ 个 stages，有 $P_{\mathrm{PP}}-1$ 条边界。每个 step 有 $M$ 个 micro-batches。因此 stage boundary 通信总量与：
[
M(P_{\mathrm{PP}}-1)B_{\mathrm{micro}}S Hs_{\mathrm{dtype}}
]
成正比。

与 TP 不同，PP 通信通常是相邻 stages 的 point-to-point，通信量与 hidden activation 大小有关，而不是与参数量直接相关。长序列和大 hidden 会增加 PP 通信压力。

\subsection{负载均衡}
\label{subsec:pp-load-balance}

Pipeline 的吞吐由最慢 stage 决定。若某个 stage 时间为 $t_{\max}$，其他 stages 更快，则更快的 stages 会等待它。负载不均来自：

\begin{itemize}
\item 层数切分不均；
\item embedding/lm head 集中在首尾 stage；
\item MoE 层造成某些 stage 更慢；
\item 激活 checkpointing 只在部分 stage 开启；
\item 硬件性能差异或节点温度降频；
\item stage 间通信跨越不同带宽链路。
\end{itemize}

负载均衡不仅要看理论 FLOPs，还要看 profiler 中的实际 stage time。某个 stage 即使 FLOPs 不高，也可能因为通信、内存带宽或小 kernel 过多而变慢。

\subsection{调试难度}
\label{subsec:pp-debugging-difficulty}

Pipeline Parallelism 的调试比 DDP 和单纯 TP 更复杂。常见问题包括：

\begin{itemize}
\item send/recv 不匹配导致死锁；
\item 某个 rank 提前退出导致其他 rank hang；
\item micro-batch id 与 activation buffer 错配；
\item backward gradient 发送给错误 stage；
\item 不同 stage 参数版本不一致；
\item pipeline flush 时遗漏某些 backward；
\item checkpoint 保存的 layer 到 stage 映射错误；
\item stage 负载不均导致吞吐低；
\item 日志只在 rank 0 打印，难以看到中间 stage 错误。
\end{itemize}

因此，工程实践中通常需要：

\begin{itemize}
\item 为每个 rank/stage 打印可控调试日志；
\item 在小模型、小 micro-batch 上验证 schedule；
\item 使用张量 shape assertion；
\item 在通信前后记录 micro-batch id；
\item 使用 profiler 查看每个 stage timeline；
\item 将调度逻辑与模型层逻辑解耦；
\item 提供 deterministic 的 pipeline schedule 测试。
\end{itemize}

\section{一个极简流水线模拟器}
\label{sec:minimal-pipeline-simulator}

为了帮助理解 pipeline bubble，可以写一个极简模拟器。它不执行真实模型，只生成每个时间片每个 stage 的任务。

\begin{lstlisting}[language=Python,caption={极简 GPipe 时间线模拟器},label={lst:minimal-gpipe-simulator}]
def gpipe_timeline(num_stages, num_microbatches):
total_forward_steps = num_microbatches + num_stages - 1
total_backward_steps = num_microbatches + num_stages - 1

```
timeline = []

# Forward
for t in range(total_forward_steps):
    row = []
    for s in range(num_stages):
        m = t - s
        if 0 <= m < num_microbatches:
            row.append(f"F{m}")
        else:
            row.append("--")
    timeline.append(row)

# Backward
for t in range(total_backward_steps):
    row = []
    for s in range(num_stages):
        # backward flows from last stage to first
        m = num_microbatches - 1 - (t - (num_stages - 1 - s))
        if 0 <= m < num_microbatches:
            row.append(f"B{m}")
        else:
            row.append("--")
    timeline.append(row)

return timeline
```

def print_timeline(timeline):
num_stages = len(timeline[0])
print("time | " + " ".join([f"S{s:>3}" for s in range(num_stages)]))
print("-" * (7 + 5 * num_stages))
for t, row in enumerate(timeline):
print(f"{t:>4} | " + " ".join([f"{x:>4}" for x in row]))

timeline = gpipe_timeline(num_stages=4, num_microbatches=6)
print_timeline(timeline)
\end{lstlisting}

对于 1F1B，调度生成会更复杂，需要根据 stage、micro-batch、warmup 数、剩余 backward 数决定下一步执行 forward 还是 backward。真实框架通常不会用一个简单循环描述完整调度，而是使用状态机或预生成 schedule table。

下面给出一个更抽象的 1F1B 状态机伪代码：

\begin{lstlisting}[language=Python,caption={1F1B 调度的状态机式伪代码},label={lst:onef1b-state-machine}]
class PipelineStageState:
def **init**(self, stage_id, num_stages, num_microbatches):
self.stage_id = stage_id
self.num_stages = num_stages
self.num_microbatches = num_microbatches

```
    self.next_forward = 0
    self.next_backward = 0
    self.saved_activations = {}

def has_forward_to_run(self):
    return self.next_forward < self.num_microbatches

def has_backward_to_run(self):
    # In a real implementation, this depends on receiving grad from next stage.
    return self.next_backward < self.next_forward - (self.num_stages - 1 - self.stage_id)

def step(self):
    if self.in_warmup_phase():
        return self.run_forward()
    elif self.has_backward_to_run():
        # In steady state, prefer backward to release activation.
        return self.run_backward()
    elif self.has_forward_to_run():
        return self.run_forward()
    else:
        return self.idle_or_cooldown()

def run_forward(self):
    mb = self.next_forward
    self.next_forward += 1
    return ("F", mb)

def run_backward(self):
    mb = self.next_backward
    self.next_backward += 1
    return ("B", mb)
```

\end{lstlisting}

这段伪代码不是可直接用于训练的实现，但它说明了 1F1B 的核心：调度器要知道每个 stage 当前能否 forward、能否 backward、是否已经收到必要输入或梯度，并尽可能优先执行能释放激活的 backward。

\section{本章小结}
\label{sec:chapter09-summary}

本章系统介绍了流水线并行 Pipeline Parallelism 的动机、基本概念、GPipe、1F1B 和 Interleaved Pipeline。

\begin{itemize}
\item \textbf{流水线并行的核心思想} 是把模型按层切成多个 pipeline stages，让不同 micro-batch 像流水线一样穿过这些 stages。
\item \textbf{Stage} 是流水线中的基本计算单元，负责一段连续或虚拟切分的模型层；stage 间传递 forward activation 和 backward activation gradient。
\item \textbf{Micro-batch} 用于填满流水线。Global batch size 满足 $B_{\mathrm{global}}=B_{\mathrm{micro}}\cdot M\cdot P_{\mathrm{DP}}$。
\item \textbf{Pipeline bubble} 是 warmup 和 cooldown 中的空闲时间。简单情况下，流水线效率近似为 $\eta_{\mathrm{pipe}}\approx \frac{M}{M+P-1}$。
\item \textbf{GPipe} 使用 all-forward-all-backward 调度，概念简单，但需要保存所有 micro-batch 的 forward 激活，显存压力较大。
\item \textbf{1F1B} 在 warmup 后交替执行 forward 和 backward，能让激活更早释放，更适合大模型训练。
\item \textbf{Interleaved Pipeline} 引入 virtual pipeline stage，让每个物理 GPU 承担多个虚拟段，以减少 bubble，但调度和通信更复杂。
\item \textbf{PP 的显存收益} 来自按层切分模型状态，但激活、pipeline buffer、embedding/lm head 和负载不均会影响真实显存。
\item \textbf{PP 的通信成本} 主要来自相邻 stage 之间传递激活和激活梯度，通信量与 micro-batch、sequence length 和 hidden size 相关。
\item \textbf{PP 的工程难点} 包括 stage 负载均衡、send/recv 死锁、micro-batch 调度、参数版本一致性、checkpoint 映射和 profiler 分析。
\end{itemize}

下一章将进入 3D 并行策略。我们会把数据并行、张量并行和流水线并行组合起来，分析 DP、TP、PP 三者如何正交组织 rank group，如何根据 GPU 数量反推并行配置，以及如何在计算、显存和通信之间做系统级权衡。

\chapter{3D 并行策略}
\label{chap:3d-parallelism}

\section{DP、TP、PP 的组合}
\label{sec:dp-tp-pp-composition}

前面三章分别讨论了数据并行、张量并行和流水线并行。到这一章，我们把三者组合起来，进入大模型训练中最常见的系统形态：\textbf{3D 并行}。

所谓 3D 并行，并不是指某个单独的新算法，而是把三类并行维度组织成一个三维并行网格：

[
\text{Data Parallelism}
\times
\text{Tensor Parallelism}
\times
\text{Pipeline Parallelism}.
]

如果用符号表示：

[
W
=

P_{\mathrm{DP}}
\cdot
P_{\mathrm{TP}}
\cdot
P_{\mathrm{PP}},
]

其中：

\begin{itemize}
\item $W$ 是总 GPU 数，也就是 world size；
\item $P_{\mathrm{DP}}$ 是数据并行 size；
\item $P_{\mathrm{TP}}$ 是张量并行 size；
\item $P_{\mathrm{PP}}$ 是流水线并行 size。
\end{itemize}

3D 并行的目标是把一个超大 Transformer 训练任务拆解到许多 GPU 上，同时控制三个核心成本：

\begin{itemize}
\item \textbf{显存}：模型参数、梯度、优化器状态、激活和通信 buffer 能否放下；
\item \textbf{计算}：每张 GPU 是否有足够大的 GEMM 和稳定工作量；
\item \textbf{通信}：DP、TP、PP 三类通信是否能被网络拓扑和调度承受。
\end{itemize}

如果只使用一种并行方式，往往会遇到边界。只用 DDP，单卡必须放下完整模型；只用 TP，层内通信过于频繁，TP size 不能无限增大；只用 PP，pipeline bubble 和 stage 负载均衡会限制效率。因此，大模型训练通常需要三者组合。

\subsection{三类并行的正交性}
\label{subsec:orthogonality-of-parallelism}

DP、TP、PP 的“正交性”是理解 3D 并行的关键。

\textbf{数据并行}切分 batch。每个数据并行副本看到不同数据分片，但副本之间逻辑上训练同一个模型。每个 DP replica 内部可以包含多个 TP ranks 和多个 PP stages。DP 维度的通信主要是梯度同步，常见操作包括 all-reduce、reduce-scatter 或 ZeRO/FSDP 中的参数/梯度分片通信。

[
\text{DP}: \quad
\mathcal{B}
===========

\mathcal{B}*0
\cup
\mathcal{B}*1
\cup
\cdots
\cup
\mathcal{B}*{P*{\mathrm{DP}}-1}.
]

\textbf{张量并行}切分每一层内部的大矩阵。一个 Transformer block 内的 QKV projection、output projection、FFN up/gate/down projection 会被切到 TP group 内的多张 GPU 上。TP 维度通信发生在每层内部，频率很高，因此通常要求 TP group 位于同一节点内或高速互联域内。

[
\mat{W}
=======

[\mat{W}*0,\mat{W}*1,\ldots,\mat{W}*{P*{\mathrm{TP}}-1}]
\quad
\text{or}
\quad
\mat{W}
=======

\begin{bmatrix}
\mat{W}*0\
\mat{W}*1\
\vdots\
\mat{W}*{P*{\mathrm{TP}}-1}
\end{bmatrix}.
]

\textbf{流水线并行}切分模型层。不同 PP stages 负责不同层段，micro-batches 像流水线一样依次穿过这些 stages。PP 维度通信主要是相邻 stage 之间发送 forward activation 和 backward activation gradient。

[
{0,1,\ldots,L-1}
================

\mathcal{L}*0
\cup
\mathcal{L}*1
\cup
\cdots
\cup
\mathcal{L}*{P*{\mathrm{PP}}-1}.
]

三者正交的意思是：同一个全局 rank 可以同时属于一个 DP group、一个 TP group 和一个 PP group。它在 DP 维度上代表某个数据副本；在 TP 维度上代表某个 tensor shard；在 PP 维度上代表某个 layer stage。

可以把一个 rank 的身份写成三元组：

[
r
\leftrightarrow
(d,t,p),
]

其中：

\begin{itemize}
\item $d\in[0,P_{\mathrm{DP}})$ 表示 data parallel index；
\item $t\in[0,P_{\mathrm{TP}})$ 表示 tensor parallel index；
\item $p\in[0,P_{\mathrm{PP}})$ 表示 pipeline parallel index。
\end{itemize}

一个完整的全局 rank 映射函数可以写为：

[
r
=

((d\cdot P_{\mathrm{PP}})+p)\cdot P_{\mathrm{TP}}+t.
]

也可以使用其他排列方式，例如让 TP 维最快变化，PP 维次之，DP 维最慢变化。不同排列方式会影响 rank mapping 和物理拓扑匹配，这一点稍后会详细讨论。

\begin{figure}[H]
\centering
\begin{tikzpicture}[
font=\small,
cube/.style={draw, rounded corners=3pt, minimum width=1.1cm, minimum height=0.65cm, align=center, fill=blue!6},
arrow/.style={-{Latex[length=2.2mm]}, thick},
group/.style={draw, rounded corners=4pt, dashed, thick}
]

```
\node[font=\bfseries] at (0,4.0) {3D 并行：DP、TP、PP 三个正交维度};

% Draw a 2 x 2 x 3-ish grid projected
\foreach \d/\dx in {0/0,1/0.55} {
  \foreach \p/\py in {0/0,1/0.9,2/1.8} {
    \foreach \t/\tx in {0/0,1/1.25} {
      \pgfmathsetmacro{\x}{-3.6+\tx+\dx}
      \pgfmathsetmacro{\y}{-0.5+\py+\dx*0.35}
      \node[cube] (r\d\t\p) at (\x,\y) {$d\d$\\$t\t,p\p$};
    }
  }
}

\draw[arrow] (-4.4,-1.2) -- (-2.3,-1.2) node[midway,below] {TP};
\draw[arrow] (-4.55,-0.85) -- (-4.55,2.1) node[midway,left] {PP};
\draw[arrow] (-4.25,-0.95) -- (-3.55,-0.55) node[midway,below right] {DP};

\node[group, fit=(r000)(r010)] {};
\node[align=center, font=\scriptsize] at (-2.95,0.1) {同一层段内\\TP group};

\draw[green!60!black, thick, dotted] (r000) -- (r100);
\draw[green!60!black, thick, dotted] (r010) -- (r110);
\node[align=center, font=\scriptsize, green!50!black] at (-1.2,-0.25) {DP replicas};

\draw[red!70, thick, dashed] (r000) -- (r001) -- (r002);
\node[align=center, font=\scriptsize, red!70] at (-4.85,1.4) {PP stages};

\node[align=left, text width=10.8cm] at (1.8,-2.0) {
  每个 rank 同时有 DP、TP、PP 三个坐标。TP group 负责层内张量切分；
  PP group 负责层间流水线；DP group 负责不同数据副本之间的梯度同步。
};
```

\end{tikzpicture}
\caption{3D 并行 rank 拓扑示意图}
\label{fig:3d-parallel-rank-topology}
\end{figure}

\subsection{group 划分}
\label{subsec:3d-parallel-group-partition}

在实际分布式训练框架中，3D 并行不是抽象概念，而是具体的 process groups。每个 rank 需要加入多个通信组。最少包括：

\begin{itemize}
\item \textbf{Tensor Parallel Group}：组内做层内 all-reduce、all-gather、reduce-scatter；
\item \textbf{Pipeline Parallel Group}：组内相邻 rank 做 send/recv；
\item \textbf{Data Parallel Group}：组内同步相同模型 shard 的梯度或状态；
\item \textbf{Embedding Group}：首尾 stage 之间共享或同步 embedding/lm head 时可能需要；
\item \textbf{Sequence Parallel Group}：通常与 TP group 相同或相关；
\item \textbf{Expert Parallel Group}：MoE 训练中用于 expert dispatch 和 all-to-all。
\end{itemize}

假设：

[
P_{\mathrm{TP}}=2,\quad
P_{\mathrm{PP}}=4,\quad
P_{\mathrm{DP}}=2.
]

则总 GPU 数为：

[
W=2\times 4\times 2=16.
]

如果使用 rank 映射：

[
r=((dP_{\mathrm{PP}})+p)P_{\mathrm{TP}}+t,
]

则全局 rank 表如下：

[
\begin{array}{c|c|c|c}
d & p & t & r\
\hline
0 & 0 & 0 & 0\
0 & 0 & 1 & 1\
0 & 1 & 0 & 2\
0 & 1 & 1 & 3\
0 & 2 & 0 & 4\
0 & 2 & 1 & 5\
0 & 3 & 0 & 6\
0 & 3 & 1 & 7\
1 & 0 & 0 & 8\
1 & 0 & 1 & 9\
1 & 1 & 0 & 10\
1 & 1 & 1 & 11\
1 & 2 & 0 & 12\
1 & 2 & 1 & 13\
1 & 3 & 0 & 14\
1 & 3 & 1 & 15
\end{array}
]

此时 TP groups 为：

[
(0,1),(2,3),(4,5),(6,7),
]
[
(8,9),(10,11),(12,13),(14,15).
]

PP groups 为：

[
(0,2,4,6),\quad (1,3,5,7),
]
[
(8,10,12,14),\quad (9,11,13,15).
]

DP groups 为：

[
(0,8),(1,9),(2,10),(3,11),
]
[
(4,12),(5,13),(6,14),(7,15).
]

注意 DP group 中的 rank 必须代表\textbf{相同的 TP shard 和相同的 PP stage}，只是数据副本不同。因此 rank $0$ 和 rank $8$ 都是 $t=0,p=0$，只是 $d$ 不同。它们之间同步的是同一 shard 的梯度。

下面给出一个简单的 group 生成代码。它用于帮助理解，真实框架会考虑更多并行维度和拓扑策略。

\begin{lstlisting}[language=Python,caption={3D 并行中 TP/PP/DP groups 的生成示意},label={lst:3d-parallel-group-generation}]
def rank_from_coords(d, p, t, dp, pp, tp):
# Layout: DP outer, PP middle, TP inner.
return ((d * pp) + p) * tp + t

def build_3d_parallel_groups(dp, pp, tp):
world_size = dp * pp * tp

```
tp_groups = []
for d in range(dp):
    for p in range(pp):
        group = [
            rank_from_coords(d, p, t, dp, pp, tp)
            for t in range(tp)
        ]
        tp_groups.append(group)

pp_groups = []
for d in range(dp):
    for t in range(tp):
        group = [
            rank_from_coords(d, p, t, dp, pp, tp)
            for p in range(pp)
        ]
        pp_groups.append(group)

dp_groups = []
for p in range(pp):
    for t in range(tp):
        group = [
            rank_from_coords(d, p, t, dp, pp, tp)
            for d in range(dp)
        ]
        dp_groups.append(group)

assert sorted(sum(tp_groups, [])) == list(range(world_size))
return tp_groups, pp_groups, dp_groups
```

tp_groups, pp_groups, dp_groups = build_3d_parallel_groups(
dp=2,
pp=4,
tp=2,
)

print("TP groups:", tp_groups)
print("PP groups:", pp_groups)
print("DP groups:", dp_groups)
\end{lstlisting}

在 PyTorch 中，这些 group 通常会通过 \texttt{torch.distributed.new_group()} 创建。训练框架会保存每个 rank 所属 group，并在不同模块中调用对应 group 的 collective 或 point-to-point 通信。

\subsection{rank mapping}
\label{subsec:rank-mapping}

Rank mapping 是把逻辑三维坐标 $(d,p,t)$ 映射到物理 GPU 的过程。它极其重要，因为不同并行维度的通信频率和通信模式不同。

一般经验是：

\begin{itemize}
\item \textbf{TP group 优先放在节点内高速互联}，因为 TP 通信发生在每层内部，频率最高；
\item \textbf{PP 相邻 stage 尽量放在高速链路上}，因为每个 micro-batch 都要传递激活和梯度；
\item \textbf{DP 可以跨节点}，因为 DP 通信虽然数据量大，但通常可与 backward 重叠，且模式规则；
\item \textbf{同一节点内的 GPU 拓扑要考虑 NVLink/NVSwitch/PCIe 层级}；
\item \textbf{跨节点要考虑 NIC 数量、rail、RDMA 路径和机架拓扑}。
\end{itemize}

设一个节点有 $G_{\mathrm{node}}=8$ 张 GPU。如果选择：

[
P_{\mathrm{TP}}=8,
]

则一个 TP group 可以完整放在单节点内。这样每层 TP all-reduce 都走节点内 NVLink/NVSwitch，通常性能较好。如果选择：

[
P_{\mathrm{TP}}=16,
]

则一个 TP group 跨两个节点，每层 TP 通信都要走跨节点网络，性能可能明显下降。除非模型极大必须这么做，否则通常避免跨节点 TP。

PP group 是否跨节点要看 stage 数和节点布局。如果每个 stage 是一个 TP group，而 TP group 占满一个节点，那么 PP 相邻 stage 自然跨节点。此时 PP 通信是激活传递，其大小与：

[
B_{\mathrm{micro}}S H
]

成正比。对于长序列和大 hidden，PP 通信也不可忽略。因此有些训练会把相邻 PP stages 安排在网络拓扑相近的节点上。

\begin{figure}[H]
\centering
\begin{tikzpicture}[
font=\small,
nodebox/.style={draw, rounded corners=5pt, minimum width=5.0cm, minimum height=2.3cm, fill=blue!3},
gpu/.style={draw, rounded corners=2pt, minimum width=0.7cm, minimum height=0.46cm, align=center, fill=orange!12},
arrow/.style={-{Latex[length=2.0mm]}, thick},
fast/.style={thick, green!60!black},
slow/.style={thick, red!70, dashed}
]

```
\node[font=\bfseries] at (0,3.5) {Rank Mapping：让高频通信走高速链路};

\node[nodebox, label={[font=\bfseries]above:Node 0}] (n0) at (-3.1,1.1) {};
\node[nodebox, label={[font=\bfseries]above:Node 1}] (n1) at (3.1,1.1) {};

\foreach \i/\x/\y in {0/-4.6/1.55,1/-3.8/1.55,2/-3.0/1.55,3/-2.2/1.55,4/-4.6/0.75,5/-3.8/0.75,6/-3.0/0.75,7/-2.2/0.75} {
  \node[gpu] (g\i) at (\x,\y) {\i};
}

\foreach \i/\x/\y/\r in {0/1.6/1.55/8,1/2.4/1.55/9,2/3.2/1.55/10,3/4.0/1.55/11,4/1.6/0.75/12,5/2.4/0.75/13,6/3.2/0.75/14,7/4.0/0.75/15} {
  \node[gpu] (h\i) at (\x,\y) {\r};
}

\draw[fast] (g0) -- (g1);
\draw[fast] (g1) -- (g2);
\draw[fast] (g2) -- (g3);
\draw[fast] (g4) -- (g5);
\draw[fast] (g5) -- (g6);
\draw[fast] (g6) -- (g7);

\draw[fast] (h0) -- (h1);
\draw[fast] (h1) -- (h2);
\draw[fast] (h2) -- (h3);
\draw[fast] (h4) -- (h5);
\draw[fast] (h5) -- (h6);
\draw[fast] (h6) -- (h7);

\draw[slow] (g3) -- node[above,font=\scriptsize] {cross-node} (h0);

\node[align=left, text width=10.8cm] at (0,-1.2) {
  TP 通信通常最高频，应尽量限制在节点内高速互联范围。
  如果 TP group 跨节点，每层都会付出跨节点通信成本，扩展效率可能急剧下降。
};
```

\end{tikzpicture}
\caption{rank mapping 对通信路径的影响}
\label{fig:rank-mapping-topology}
\end{figure}

\subsection{为什么 rank mapping 影响性能}
\label{subsec:why-rank-mapping-affects-performance}

通信性能不是只由总 GPU 数决定，而是由\textbf{哪个 rank 和哪个 rank 通信、通信走哪条链路、链路是否拥塞}决定。

TP 通信常发生在每层内部。一个 80 层模型，每层 attention output 和 FFN down projection 至少各一次 TP all-reduce，则一个 forward/backward step 中 TP collective 次数非常多。如果 TP group 跨节点，每层都要跨网络同步，通信尾部可能无法隐藏。

PP 通信发生在 stage 边界。若 micro-batch 数 $M$ 很大，stage 间 send/recv 次数为 $O(MP_{\mathrm{PP}})$。如果相邻 PP stages 被放在跨机架或拥塞链路上，pipeline bubble 和等待时间会上升。

DP 通信通常发生在 backward 梯度同步。DDP bucket 可以与 backward 重叠，ZeRO/FSDP 也能通过 reduce-scatter/all-gather 改变通信形态。DP 跨节点是常见选择，但仍需要高带宽网络。

因此，一个经验性 rank mapping 原则是：

[
\text{通信频率越高、延迟越敏感的 group，越应该放在更近的拓扑范围内。}
]

按优先级粗略排序：

[
\text{TP}

>

\text{PP adjacent stages}

>

\text{DP}.
]

但这不是绝对。若 PP activation 很大、micro-batch 很多，PP 也可能成为主要瓶颈。若 ZeRO-3/FSDP 频繁 all-gather 参数，DP 维通信也会非常重。真正配置时需要结合 profiler、NCCL tests 和训练 step breakdown。

\section{计算、显存与通信的三角权衡}
\label{sec:compute-memory-communication-triangle}

3D 并行配置的本质，是在计算、显存与通信之间做权衡。一个配置可能显存很省，但通信太重；另一个配置计算效率高，但单卡显存放不下；还有一个配置通信少，但 pipeline bubble 大。

我们可以把三者画成一个三角形：

[
\text{Compute}
\leftrightarrow
\text{Memory}
\leftrightarrow
\text{Communication}.
]

\begin{figure}[H]
\centering
\begin{tikzpicture}[
font=\small,
vertex/.style={circle, draw, minimum size=1.25cm, align=center},
arrow/.style={-{Latex[length=2.2mm]}, thick}
]

```
\node[font=\bfseries] at (0,3.2) {3D 并行的 Compute-Memory-Communication 三角权衡};

\node[vertex, fill=blue!8] (compute) at (0,1.7) {Compute\\效率};
\node[vertex, fill=green!10] (memory) at (-4.0,-1.2) {Memory\\显存};
\node[vertex, fill=orange!14] (comm) at (4.0,-1.2) {Communication\\通信};

\draw[arrow] (memory) -- node[left, align=center] {切分模型\\降低显存} (compute);
\draw[arrow] (compute) -- node[right, align=center] {GEMM 变小\\通信增加} (comm);
\draw[arrow] (comm) -- node[below, align=center] {通信 buffer\\拓扑限制} (memory);

\node[align=left, text width=10.5cm] at (0,-2.7) {
  DP、TP、PP 的配置不是单目标优化。TP 增大可降低层内显存但增加高频通信；
  PP 增大可降低层数相关显存但增加 bubble 和 activation send/recv；
  DP 增大可提升数据吞吐但增加梯度同步压力。
};
```

\end{tikzpicture}
\caption{计算、显存与通信的三角权衡}
\label{fig:compute-memory-communication-triangle}
\end{figure}

\subsection{显存占用估算}
\label{subsec:memory-estimation}

训练显存可以拆成五部分：

[
M_{\mathrm{total}}
==================

M_{\mathrm{params}}
+
M_{\mathrm{grads}}
+
M_{\mathrm{optimizer}}
+
M_{\mathrm{activations}}
+
M_{\mathrm{buffers}}.
]

对于一个 dense Transformer，如果暂时忽略 embedding/lm head 和不均匀层切分，模型层参数在 TP 和 PP 下的每卡占用近似为：

[
M_{\mathrm{params,perGPU}}
\approx
\frac{N_{\mathrm{params}}\cdot s_{\mathrm{param}}}
{P_{\mathrm{TP}}P_{\mathrm{PP}}}.
]

这里的解释是：TP 切分层内矩阵，PP 切分层数。DP 不减少每个副本中的参数占用，除非使用 ZeRO/FSDP。

梯度与参数 shard 对应，若不使用 ZeRO/FSDP：

[
M_{\mathrm{grads,perGPU}}
\approx
\frac{N_{\mathrm{params}}\cdot s_{\mathrm{grad}}}
{P_{\mathrm{TP}}P_{\mathrm{PP}}}.
]

AdamW 优化器状态通常包括 master weight、一阶矩和二阶矩。如果它们随 TP/PP shard 保存，则：

[
M_{\mathrm{optimizer,perGPU}}
\approx
\frac{N_{\mathrm{params}}\cdot s_{\mathrm{optim}}}
{P_{\mathrm{TP}}P_{\mathrm{PP}}}.
]

如果使用 ZeRO-1 在 DP 维切分 optimizer state，则优化器状态进一步除以 $P_{\mathrm{DP}}$：

[
M_{\mathrm{optimizer,perGPU}}^{\mathrm{ZeRO1}}
\approx
\frac{N_{\mathrm{params}}\cdot s_{\mathrm{optim}}}
{P_{\mathrm{TP}}P_{\mathrm{PP}}P_{\mathrm{DP}}}.
]

如果使用 ZeRO-2，梯度也可在 DP 维切分；ZeRO-3 则参数也可在 DP 维切分。但 ZeRO-3 与 TP/PP 组合后通信更加复杂，后续 ZeRO 章节会专门讨论。

激活显存的估算更复杂。对一个 Transformer 层，隐藏状态激活基础大小为：

[
B_{\mathrm{micro}}S H s_{\mathrm{act}}.
]

若每个 PP stage 有 $L/P_{\mathrm{PP}}$ 层，则每个 stage 的激活显存粗略与：

[
M_{\mathrm{act,stage}}
\propto
N_{\mathrm{live\ microbatches}}
\cdot
\frac{L}{P_{\mathrm{PP}}}
\cdot
B_{\mathrm{micro}}S Hs_{\mathrm{act}}
]

相关。其中 $N_{\mathrm{live\ microbatches}}$ 与 pipeline schedule 有关。GPipe 中它可能接近 $M$；1F1B 中通常显著更低。

TP 对激活显存的影响取决于是否使用 sequence parallel。若仅使用 TP，一些中间激活沿 hidden 维切分，但 residual stream 可能仍 replicated。使用 sequence parallel 后，部分激活可进一步除以 $P_{\mathrm{TP}}$。

因此，真实估算可以写成：

[
M_{\mathrm{act}}
================

M_{\mathrm{act,sharded}}
+
M_{\mathrm{act,replicated}}
+
M_{\mathrm{act,temp}}.
]

如果粗略引入一个激活切分因子 $q_{\mathrm{act}}$：

[
M_{\mathrm{act,perGPU}}
\approx
\frac{
c_{\mathrm{act}}
\cdot
N_{\mathrm{live\ microbatches}}
\cdot
\frac{L}{P_{\mathrm{PP}}}
\cdot
B_{\mathrm{micro}}S Hs_{\mathrm{act}}
}
{q_{\mathrm{act}}},
]

其中 $c_{\mathrm{act}}$ 表示每层实际保存激活相对 hidden state 的倍数，$q_{\mathrm{act}}$ 可能为 $1$、$P_{\mathrm{TP}}$ 或介于两者之间。

下面给出一个简化显存估算代码，用于配置初筛。

\begin{lstlisting}[language=Python,caption={3D 并行训练显存的粗略估算器},label={lst:3d-memory-estimator}]
def estimate_3d_memory_gb(
num_params,
num_layers,
hidden_size,
seq_len,
micro_batch_size,
tp_size,
pp_size,
dp_size,
live_microbatches,
param_bytes=2,      # BF16 parameters
grad_bytes=2,       # BF16 gradients
master_bytes=4,     # FP32 master weights
adam_m_bytes=4,     # FP32 Adam m
adam_v_bytes=4,     # FP32 Adam v
act_bytes=2,        # BF16 activations
activation_factor=8.0,
activation_shard_factor=1.0,
zero_stage=0,
buffer_gb=4.0,
):
"""
A rough estimator for screening configurations.
It is not a replacement for real profiler-based measurement.
"""

```
# Parameters are sharded by TP and PP.
param_shard = num_params / (tp_size * pp_size)

param_mem = param_shard * param_bytes
grad_mem = param_shard * grad_bytes

optim_bytes = master_bytes + adam_m_bytes + adam_v_bytes
optim_mem = param_shard * optim_bytes

# ZeRO effects along DP dimension.
if zero_stage >= 1:
    optim_mem /= dp_size
if zero_stage >= 2:
    grad_mem /= dp_size
if zero_stage >= 3:
    param_mem /= dp_size

layers_per_stage = num_layers / pp_size
hidden_act = micro_batch_size * seq_len * hidden_size * act_bytes

act_mem = (
    live_microbatches
    * layers_per_stage
    * hidden_act
    * activation_factor
    / activation_shard_factor
)

total = param_mem + grad_mem + optim_mem + act_mem + buffer_gb * (1024 ** 3)

return {
    "params_gb": param_mem / (1024 ** 3),
    "grads_gb": grad_mem / (1024 ** 3),
    "optimizer_gb": optim_mem / (1024 ** 3),
    "activations_gb": act_mem / (1024 ** 3),
    "buffers_gb": buffer_gb,
    "total_gb": total / (1024 ** 3),
}
```

memory = estimate_3d_memory_gb(
num_params=70e9,
num_layers=80,
hidden_size=8192,
seq_len=4096,
micro_batch_size=1,
tp_size=8,
pp_size=4,
dp_size=8,
live_microbatches=4,
activation_shard_factor=8,  # assume sequence parallel roughly helps
zero_stage=1,
)

for k, v in memory.items():
print(k, round(v, 2))
\end{lstlisting}

这段代码只能用于直觉估算。真实显存还会受到 allocator fragmentation、kernel workspace、CUDA graph、NCCL buffer、FlashAttention workspace、recompute 策略、模型具体结构、MoE、embedding/lm head、vocab parallel cross entropy 等影响。

\subsection{通信量估算}
\label{subsec:communication-estimation}

3D 并行中的通信可以分为三类：

[
C_{\mathrm{comm}}
=================

C_{\mathrm{DP}}
+
C_{\mathrm{TP}}
+
C_{\mathrm{PP}}.
]

\textbf{DP 通信}主要是梯度同步或 ZeRO/FSDP 的状态同步。如果不使用 ZeRO，DDP 每 step 需要同步参数梯度。假设每个 DP replica 内部经过 TP/PP 后，每个 rank 持有的参数 shard 大小为：

[
G_{\mathrm{rank}}
=================

\frac{N_{\mathrm{params}}s_{\mathrm{grad}}}
{P_{\mathrm{TP}}P_{\mathrm{PP}}}.
]

普通 all-reduce 的每 rank 通信量可粗略按 Ring AllReduce 估算：

[
V_{\mathrm{DP}}
\approx
\frac{2(P_{\mathrm{DP}}-1)}{P_{\mathrm{DP}}}
G_{\mathrm{rank}}.
]

若使用 ZeRO-2 或 ZeRO-3，通信形态会变成 reduce-scatter、all-gather 和参数预取，不能简单用 DDP all-reduce 代替。

\textbf{TP 通信}发生在每层内部。Megatron 风格每层至少有 attention output projection 后一次 all-reduce、FFN down projection 后一次 all-reduce。每次 all-reduce 的张量大小大致为：

[
B_{\mathrm{micro}}S Hs_{\mathrm{act}}.
]

如果每个 stage 有 $L/P_{\mathrm{PP}}$ 层，每个 micro-batch 在本 stage 触发约 $n_{\mathrm{tp_collectives_per_layer}}$ 次 TP collective，则每个 step TP 通信量与：

[
M
\cdot
\frac{L}{P_{\mathrm{PP}}}
\cdot
n_{\mathrm{tp}}
\cdot
B_{\mathrm{micro}}S Hs_{\mathrm{act}}
]

相关。再乘以 all-reduce 的通信放大系数：

[
\frac{2(P_{\mathrm{TP}}-1)}{P_{\mathrm{TP}}}.
]

因此：

[
V_{\mathrm{TP}}
\approx
M
\cdot
\frac{L}{P_{\mathrm{PP}}}
\cdot
n_{\mathrm{tp}}
\cdot
\frac{2(P_{\mathrm{TP}}-1)}{P_{\mathrm{TP}}}
\cdot
B_{\mathrm{micro}}S Hs_{\mathrm{act}}.
]

\textbf{PP 通信}是 stage 边界传递 activation 和 activation gradient。每条边界每个 micro-batch forward 发送一次激活，backward 发送一次梯度：

[
V_{\mathrm{PP,boundary}}
\approx
2M B_{\mathrm{micro}}S Hs_{\mathrm{act}}.
]

每个 rank 只和相邻 stage 通信，不一定经过所有边界。因此每 rank PP 通信取决于它是否在中间 stage、首 stage 还是尾 stage。中间 stage 既接收又发送，边界通信更多；首尾 stage 少一侧。

通信估算的重点不是得到一个完美数字，而是判断哪个维度最可能成为瓶颈：

[
\text{TP 高频，偏低延迟和高带宽节点内通信；}
]
[
\text{PP 与 micro-batch 和激活大小相关，偏点对点通信；}
]
[
\text{DP 与参数/梯度大小相关，偏大块 collective 通信。}
]

\begin{table}[H]
\centering
\small
\caption{DP、TP、PP 通信模式对比}
\label{tab:dp-tp-pp-communication-comparison}
\begin{tabular}{p{0.16\textwidth}p{0.25\textwidth}p{0.24\textwidth}p{0.25\textwidth}}
\toprule
\textbf{并行维度} & \textbf{通信对象} & \textbf{常见操作} & \textbf{性能敏感点} \
\midrule
DP
& 梯度、参数 shard、优化器状态
& all-reduce, reduce-scatter, all-gather
& 跨节点带宽、bucket overlap、ZeRO stage。 \
\midrule
TP
& 层内激活、partial output、sequence shard
& all-reduce, all-gather, reduce-scatter
& 频率高，强依赖节点内 NVLink/NVSwitch。 \
\midrule
PP
& stage 边界激活与激活梯度
& send/recv
& micro-batch 数、hidden/seq 大小、stage 拓扑。 \
\midrule
MoE EP
& token dispatch 和 expert output
& all-to-all
& 负载均衡、网络拥塞、expert placement。 \
\bottomrule
\end{tabular}
\end{table}

\subsection{计算效率估算}
\label{subsec:compute-efficiency-estimation}

大模型训练常用 MFU，Model FLOPs Utilization，衡量有效模型计算占硬件峰值计算的比例：

[
\mathrm{MFU}
============

\frac{\mathrm{Model\ FLOPs\ per\ step}}
{\mathrm{Step\ time}\cdot \mathrm{Peak\ FLOPs\ of\ GPUs}}.
]

如果模型每个 token 的训练 FLOPs 粗略为 $6N$，每 step 训练 token 数为：

[
T_{\mathrm{tokens/step}}
========================

B_{\mathrm{global}}S,
]

则每 step 模型 FLOPs 为：

[
F_{\mathrm{step}}
\approx
6N B_{\mathrm{global}}S.
]

若使用 $W$ 张 GPU，每张 GPU 峰值为 $F_{\mathrm{peak}}$，step time 为 $T_{\mathrm{step}}$，则：

[
\mathrm{MFU}
\approx
\frac{6N B_{\mathrm{global}}S}
{T_{\mathrm{step}} W F_{\mathrm{peak}}}.
]

MFU 下降可能来自：

\begin{itemize}
\item TP size 太大导致 GEMM 太小；
\item PP bubble 占比高；
\item DP 通信暴露；
\item 数据加载阻塞；
\item checkpoint 阻塞；
\item kernel fusion 不足，小 kernel 太多；
\item activation checkpointing 重计算太多；
\item straggler 或节点性能不一致；
\item 通信与计算资源争抢。
\end{itemize}

Pipeline bubble 对计算效率的影响可以单独估算。若使用 $M$ 个 micro-batches、$P_{\mathrm{PP}}$ 个 stages，简单近似：

[
\eta_{\mathrm{pipe}}
\approx
\frac{M}{M+P_{\mathrm{PP}}-1}.
]

例如 $P_{\mathrm{PP}}=8$，$M=8$：

[
\eta_{\mathrm{pipe}}
====================

# \frac{8}{8+8-1}

\frac{8}{15}
\approx 53.3%.
]

这说明 micro-batch 数与 PP size 相当时，bubble 仍然很大。若 $M=64$：

[
\eta_{\mathrm{pipe}}
====================

# \frac{64}{64+8-1}

\frac{64}{71}
\approx 90.1%.
]

但 $M$ 变大意味着 gradient accumulation 变大，global batch size 也变大，可能影响训练动力学。因此 pipeline efficiency 不能孤立优化。

\begin{lstlisting}[language=Python,caption={MFU 与 Pipeline Efficiency 的简化估算},label={lst:mfu-pipeline-efficiency-estimator}]
def estimate_mfu(
num_params,
global_batch_size,
seq_len,
step_time_sec,
num_gpus,
peak_flops_per_gpu,
):
# Dense Transformer training FLOPs rough estimate.
flops_per_step = 6 * num_params * global_batch_size * seq_len
peak_flops_per_step = step_time_sec * num_gpus * peak_flops_per_gpu
return flops_per_step / peak_flops_per_step

def pipeline_efficiency(num_microbatches, pp_size):
return num_microbatches / (num_microbatches + pp_size - 1)

for m in [4, 8, 16, 32, 64]:
print(m, round(pipeline_efficiency(m, pp_size=8), 3))
\end{lstlisting}

\subsection{不同规模模型的并行策略选择}
\label{subsec:parallel-strategy-for-different-model-sizes}

并行策略没有唯一最优答案。下面给出一组工程直觉，用于初始配置选择。

\textbf{小模型，例如 1B 到 7B。}如果单卡可以放下模型和训练状态，优先使用 DDP 或 FSDP。TP/PP 可能降低单卡计算效率，除非 batch 或序列足够大。常见策略：

[
P_{\mathrm{TP}}=1,\quad
P_{\mathrm{PP}}=1,\quad
P_{\mathrm{DP}}=W.
]

如果显存不足，可以使用 ZeRO/FSDP、activation checkpointing 或 LoRA/QLoRA 等方法。

\textbf{中等模型，例如 7B 到 30B。}可能需要 TP 或 ZeRO。若节点内 8 卡互联较好，可尝试：

[
P_{\mathrm{TP}}=2\ \text{or}\ 4,
\quad
P_{\mathrm{PP}}=1,
\quad
P_{\mathrm{DP}}=\frac{W}{P_{\mathrm{TP}}}.
]

如果层数较多、显存仍不足，可引入 PP：

[
P_{\mathrm{TP}}=4,
\quad
P_{\mathrm{PP}}=2,
\quad
P_{\mathrm{DP}}=\frac{W}{8}.
]

\textbf{大模型，例如 30B 到 100B。}通常需要 TP+PP+DP。TP 放在节点内，PP 跨节点或跨节点组，DP 复制多个 3D replica。典型方向：

[
P_{\mathrm{TP}}=4\ \text{or}\ 8,
\quad
P_{\mathrm{PP}}=4\ \text{or}\ 8,
\quad
P_{\mathrm{DP}}=\frac{W}{P_{\mathrm{TP}}P_{\mathrm{PP}}}.
]

\textbf{超大模型，例如 100B 以上。}通常需要 3D 并行加 ZeRO、activation checkpointing、sequence parallel、distributed checkpoint 和拓扑感知调度。TP 可能取节点内最大高速互联规模，例如 8；PP 取决于层数和显存；DP size 由总 GPU 数和 global batch 需求决定。

\begin{table}[H]
\centering
\small
\caption{不同模型规模下的并行策略直觉}
\label{tab:parallel-strategy-by-model-size}
\begin{tabular}{p{0.18\textwidth}p{0.25\textwidth}p{0.24\textwidth}p{0.23\textwidth}}
\toprule
\textbf{模型规模} & \textbf{优先策略} & \textbf{常见配置方向} & \textbf{主要风险} \
\midrule
1B--7B
& DDP / FSDP
& TP=1, PP=1，扩大 DP
& 显存、数据加载、小 batch 通信。 \
\midrule
7B--30B
& TP 或 ZeRO/FSDP
& TP=2/4，PP 可为 1/2
& TP 通信、GEMM 变小、ZeRO 通信。 \
\midrule
30B--100B
& TP+PP+DP
& TP=4/8，PP=2/4/8
& pipeline bubble、stage 均衡、rank mapping。 \
\midrule
100B+
& 3D + ZeRO + checkpointing
& TP 节点内，PP 跨节点，DP 多副本
& 通信、checkpoint、故障恢复、长时间稳定性。 \
\bottomrule
\end{tabular}
\end{table}

这些只是初始猜测。真实项目必须通过 profile 和试跑验证。一个配置是否好，不看理论是否优雅，而看：

\begin{itemize}
\item 单卡显存是否有安全余量；
\item MFU 是否合理；
\item step time 是否稳定；
\item pipeline bubble 是否可接受；
\item NCCL 时间占比是否过高；
\item checkpoint 是否能在可接受时间内完成；
\item 故障恢复是否可靠；
\item global batch 和学习率是否满足训练 recipe。
\end{itemize}

\section{大模型训练配置推导}
\label{sec:large-model-training-config-derivation}

3D 并行配置推导是 AI Infra 工程师的核心能力之一。给定模型规模、GPU 数量、单卡显存、节点拓扑和目标 global batch，需要反推出 TP、PP、DP、micro-batch、gradient accumulation 等参数。

\subsection{global batch size}
\label{subsec:global-batch-size}

在 3D 并行中，global batch size 由数据并行、副本内 micro-batch 和梯度累积共同决定：

[
B_{\mathrm{global}}
===================

B_{\mathrm{micro}}
\cdot
M
\cdot
P_{\mathrm{DP}}.
]

其中：

\begin{itemize}
\item $B_{\mathrm{micro}}$ 是每个 pipeline micro-batch 的 batch size；
\item $M$ 是每个 optimizer step 的 micro-batch 数；
\item $P_{\mathrm{DP}}$ 是数据并行副本数。
\end{itemize}

如果按 token 数计算：

[
T_{\mathrm{global}}
===================

B_{\mathrm{global}}
\cdot
S
=

B_{\mathrm{micro}}
\cdot
M
\cdot
P_{\mathrm{DP}}
\cdot
S.
]

在语言模型训练中，常用 tokens per step 描述 batch，因为序列长度直接影响计算量。若使用 sequence packing，实际每个 sample 的有效 token 可能不同，此时需要按 non-padding token 数统计。

\subsection{micro batch size}
\label{subsec:micro-batch-size}

Micro-batch size 是 3D 并行中非常敏感的参数。它影响：

\begin{itemize}
\item 单个 micro-batch 激活显存；
\item GEMM shape 和 Tensor Core 利用率；
\item pipeline bubble 摊薄方式；
\item gradient accumulation steps；
\item global batch size；
\item data loader 和 sequence packing。
\end{itemize}

通常显存越紧，$B_{\mathrm{micro}}$ 越小。很多大模型训练中，micro-batch size 甚至可能为 $1$。这并不意味着 global batch size 小，因为可以通过更大的 $M$ 和 $P_{\mathrm{DP}}$ 累积到目标 global batch。

但是，$B_{\mathrm{micro}}=1$ 也可能带来问题：

\begin{itemize}
\item GEMM 的 $M$ 维，即 token 数 $B_{\mathrm{micro}}S$，可能较小；
\item pipeline micro-batch 数需要增加，调度开销上升；
\item 如果序列长度也不大，GPU 利用率会下降；
\item Batch 维相关的数据增强或采样统计变弱。
\end{itemize}

因此配置时通常先从显存允许的最大 micro-batch 试起，再根据 global batch 需求调 gradient accumulation。

\subsection{gradient accumulation}
\label{subsec:gradient-accumulation-in-3d}

在 PP 中，micro-batch 数 $M$ 既是 gradient accumulation 的一部分，也是 pipeline 填充效率的关键。为了达到目标 global batch：

[
M
=

\frac{B_{\mathrm{global}}}
{B_{\mathrm{micro}}P_{\mathrm{DP}}}.
]

同时为了降低 pipeline bubble，希望：

[
M \gg P_{\mathrm{PP}}.
]

这两个条件可能冲突。如果 $P_{\mathrm{DP}}$ 很大，给定 global batch 下 $M$ 可能变小，导致 pipeline bubble 大。反过来，如果为了提高 pipeline efficiency 增大 $M$，global batch 可能过大，需要调整学习率或训练 recipe。

例如：

[
B_{\mathrm{global}}=1024,
\quad
B_{\mathrm{micro}}=1,
\quad
P_{\mathrm{DP}}=64.
]

则：

[
M=\frac{1024}{1\times 64}=16.
]

若：

[
P_{\mathrm{PP}}=16,
]

则 pipeline efficiency 近似为：

[
\eta_{\mathrm{pipe}}
====================

# \frac{16}{16+16-1}

\frac{16}{31}
\approx 51.6%.
]

这很低。可以考虑：

\begin{itemize}
\item 减小 DP size，增大 $M$；
\item 减小 PP size；
\item 增大 global batch；
\item 增大 micro-batch；
\item 使用 interleaved pipeline；
\item 调整并行切分，使 PP bubble 降低。
\end{itemize}

\subsection{TP size、PP size、DP size}
\label{subsec:tp-pp-dp-size}

选择 TP、PP、DP size 时，可以按以下顺序思考。

\textbf{第一步：选择 TP size。}

TP size 主要受以下因素约束：

\begin{itemize}
\item hidden size 是否能整除 TP size；
\item attention heads 是否能整除 TP size；
\item GQA 中 kv heads 是否能合理切分；
\item TP group 是否能放进节点内高速互联；
\item 单层矩阵是否太大，需要更多 TP；
\item TP 后 GEMM 是否仍然足够大。
\end{itemize}

常见策略是 TP size 选择节点内 GPU 数的因子，例如 2、4、8。对于 8 卡 NVSwitch 节点，TP=8 是常见上限。若 TP 跨节点，需要非常谨慎。

\textbf{第二步：选择 PP size。}

PP size 主要受以下因素约束：

\begin{itemize}
\item 模型层数 $L$ 是否能较均匀切分；
\item 每个 stage 的参数和激活显存是否放得下；
\item PP bubble 是否可接受；
\item micro-batch 数是否足够大；
\item stage 间通信是否会跨慢链路；
\item embedding/lm head 是否造成首尾 stage 不均。
\end{itemize}

一般来说，PP size 越大，每个 stage 层数越少，显存越省；但 bubble 和调度复杂度越高。PP size 也不宜超过层数的合理切分粒度。

\textbf{第三步：由总 GPU 数反推 DP size。}

当 $P_{\mathrm{TP}}$ 和 $P_{\mathrm{PP}}$ 确定后：

[
P_{\mathrm{DP}}
===============

\frac{W}{P_{\mathrm{TP}}P_{\mathrm{PP}}}.
]

DP size 决定可并行处理多少数据副本，也影响 global batch 和梯度同步规模。若 DP size 太小，global batch 可能不足或训练吞吐不足；若 DP size 太大，在固定 global batch 下 micro-batch 数可能太小，pipeline bubble 上升。

\subsection{从 GPU 数量反推并行配置}
\label{subsec:derive-parallel-config-from-gpu-count}

下面给出一个配置推导流程。假设：

\begin{itemize}
\item 总 GPU 数 $W=256$；
\item 每节点 8 张 GPU；
\item 模型有 $L=80$ 层；
\item hidden size $H=8192$；
\item attention heads $n_h=64$；
\item 目标 global batch size $B_{\mathrm{global}}=1024$；
\item 序列长度 $S=4096$；
\item 初始设 micro-batch size $B_{\mathrm{micro}}=1$。
\end{itemize}

\textbf{候选配置 A：}

[
P_{\mathrm{TP}}=8,\quad
P_{\mathrm{PP}}=4.
]

则：

[
P_{\mathrm{DP}}
===============

# \frac{256}{8\times 4}

8.

]

Gradient accumulation micro-batch 数为：

[
M
=

# \frac{1024}{1\times 8}

128.

]

Pipeline efficiency 近似为：

[
\eta_{\mathrm{pipe}}
====================

# \frac{128}{128+4-1}

\frac{128}{131}
\approx 97.7%.
]

这个配置 pipeline bubble 很低。TP=8 可放在每个节点内；PP=4 意味着一个 DP replica 需要：

[
8\times 4=32
]

张 GPU，也就是 4 个节点。总共 256 GPU 可形成 8 个 DP replicas。

\textbf{候选配置 B：}

[
P_{\mathrm{TP}}=4,\quad
P_{\mathrm{PP}}=8.
]

则：

[
P_{\mathrm{DP}}
===============

# \frac{256}{4\times 8}

8.

]

Micro-batch 数仍为：

[
M=128.
]

Pipeline efficiency：

[
\eta_{\mathrm{pipe}}
====================

# \frac{128}{128+8-1}

\frac{128}{135}
\approx 94.8%.
]

这个配置 TP 通信更少，单个 TP group 只用 4 张 GPU；PP 更深，每个 stage 层数约 $10$。如果显存允许，这可能比 TP=8 更好，因为 GEMM 更大、TP 通信更少。但 PP stage 间通信和调度稍多。

\textbf{候选配置 C：}

[
P_{\mathrm{TP}}=8,\quad
P_{\mathrm{PP}}=8.
]

则：

[
P_{\mathrm{DP}}
===============

# \frac{256}{8\times 8}

4.

]

Micro-batch 数：

[
M
=

# \frac{1024}{1\times 4}

256.

]

Pipeline efficiency：

[
\eta_{\mathrm{pipe}}
====================

# \frac{256}{256+8-1}

\frac{256}{263}
\approx 97.3%.
]

这个配置显存更省，但每个 DP replica 需要 64 GPU，只有 4 个 DP replicas。它适合模型更大、显存更紧的情况，但 TP 和 PP 都较高，通信和调度复杂度更大。

下面给出一个枚举候选配置的脚本。

\begin{lstlisting}[language=Python,caption={从 GPU 数量枚举 3D 并行候选配置},label={lst:enumerate-3d-configs}]
def divisors(n):
return [x for x in range(1, n + 1) if n % x == 0]

def enumerate_3d_configs(
world_size,
gpus_per_node,
num_layers,
num_heads,
hidden_size,
global_batch_size,
micro_batch_size,
preferred_tp_values=(1, 2, 4, 8),
):
configs = []

```
for tp in preferred_tp_values:
    if world_size % tp != 0:
        continue
    if num_heads % tp != 0:
        continue
    if hidden_size % tp != 0:
        continue

    # Prefer TP within a node.
    if tp > gpus_per_node:
        continue

    remaining = world_size // tp

    for pp in divisors(remaining):
        if num_layers % pp != 0:
            continue

        dp = remaining // pp
        if dp <= 0:
            continue

        denom = micro_batch_size * dp
        if global_batch_size % denom != 0:
            continue

        microbatches = global_batch_size // denom
        pipe_eff = microbatches / (microbatches + pp - 1)

        configs.append({
            "tp": tp,
            "pp": pp,
            "dp": dp,
            "microbatches": microbatches,
            "layers_per_stage": num_layers // pp,
            "pipeline_efficiency": pipe_eff,
            "gpus_per_replica": tp * pp,
        })

# Sort by high pipeline efficiency, then smaller TP/PP complexity.
configs.sort(
    key=lambda c: (
        -c["pipeline_efficiency"],
        c["tp"] * c["pp"],
        c["pp"],
    )
)
return configs
```

configs = enumerate_3d_configs(
world_size=256,
gpus_per_node=8,
num_layers=80,
num_heads=64,
hidden_size=8192,
global_batch_size=1024,
micro_batch_size=1,
)

for cfg in configs[:10]:
print(cfg)
\end{lstlisting}

这个脚本只做基本整除和 pipeline efficiency 检查。真实系统还需要把显存估算、通信估算、节点拓扑、ZeRO stage、checkpoint 策略和训练 recipe 一起纳入搜索。

\section{故障恢复与 checkpoint}
\label{sec:fault-tolerance-and-checkpoint}

大模型训练通常持续数天、数周甚至更久。随着 GPU 数量增加，故障概率也会增加。假设单个节点在训练期间发生故障的概率为 $p$，节点数为 $N_{\mathrm{nodes}}$，整个作业无故障概率近似为：

[
(1-p)^{N_{\mathrm{nodes}}}.
]

当节点数很大时，即使单节点故障概率很低，整个作业也很可能遇到故障。因此，大规模训练系统必须具备 checkpoint 和恢复能力。

Checkpoint 不只是保存模型权重。一个完整训练 checkpoint 通常包括：

\begin{itemize}
\item 模型参数；
\item 优化器状态；
\item 梯度缩放器状态；
\item 学习率 scheduler 状态；
\item global step；
\item consumed samples 或 consumed tokens；
\item tokenizer 和数据状态；
\item RNG states；
\item 并行配置和 rank mapping；
\item ZeRO/FSDP/TP/PP shard metadata；
\item dataloader 或 dataset cursor；
\item 训练配置文件与代码版本信息。
\end{itemize}

如果少保存其中某些状态，训练可能“能恢复”，但不是严格恢复。比如 scheduler step 错了，学习率会跳；consumed tokens 错了，数据会重复或跳过；RNG state 错了，dropout 和数据 shuffle 不一致；rank mapping metadata 错了，sharded 权重可能加载到错误 rank。

\subsection{distributed checkpoint}
\label{subsec:distributed-checkpoint}

单机小模型常由 rank 0 保存完整 checkpoint。但 3D 并行训练中，这通常不可行。原因包括：

\begin{itemize}
\item 单个 rank 不一定拥有完整模型；
\item TP/PP/ZeRO/FSDP 后参数和优化器状态是分片的；
\item checkpoint 体积可能达到 TB 级；
\item rank 0 聚合完整状态会造成显存和内存峰值；
\item 单点写入会造成 IO 瓶颈；
\item 恢复时重新广播完整状态代价巨大。
\end{itemize}

Distributed checkpoint 的思想是：每个 rank 或每组 rank 保存自己负责的 shard，形成一个分布式 checkpoint 目录。目录中包含：

\begin{itemize}
\item 每个 shard 的数据文件；
\item 描述全局张量如何切分的 metadata；
\item rank 到 shard 的映射；
\item tensor 名称、shape、dtype、offset；
\item 训练状态元数据。
\end{itemize}

恢复时，每个 rank 根据当前并行配置和 metadata 读取自己需要的 shard。如果恢复配置与保存配置相同，读取较简单；如果 TP/PP/DP size 改变，则需要做 resharding。

\begin{figure}[H]
\centering
\begin{tikzpicture}[
font=\small,
rank/.style={draw, rounded corners=3pt, minimum width=1.45cm, minimum height=0.65cm, align=center, fill=blue!6},
file/.style={draw, rounded corners=2pt, minimum width=1.35cm, minimum height=0.5cm, align=center, fill=orange!14},
meta/.style={draw, rounded corners=3pt, minimum width=2.2cm, minimum height=0.65cm, align=center, fill=green!10},
arrow/.style={-{Latex[length=2.1mm]}, thick}
]

```
\node[font=\bfseries] at (0,3.35) {分布式 Checkpoint：每个 rank 保存自己的 shard};

\foreach \i/\x in {0/-4.8,1/-2.4,2/0,3/2.4,4/4.8} {
  \node[rank] (r\i) at (\x,1.55) {Rank \i};
  \node[file] (f\i) at (\x,0.45) {shard \i};
  \draw[arrow] (r\i) -- (f\i);
}

\node[meta] (meta) at (0,-0.8) {metadata.json\\tensor layout};
\foreach \i in {0,1,2,3,4} {
  \draw[arrow] (f\i) -- (meta);
}

\node[align=left, text width=10.6cm] at (0,-2.25) {
  分布式 checkpoint 不聚合到单个 rank，而是保存 shard 与 metadata。
  metadata 描述每个全局张量如何被 TP、PP、DP、ZeRO 等维度切分。
};
```

\end{tikzpicture}
\caption{分布式 checkpoint 保存流程}
\label{fig:distributed-checkpoint-flow}
\end{figure}

\subsection{sharded checkpoint}
\label{subsec:sharded-checkpoint}

Sharded checkpoint 是分布式训练中的常见格式。以一个线性层权重为例：

[
\mat{W}\in\mathbb{R}^{H_{\mathrm{in}}\times H_{\mathrm{out}}}.
]

如果它在 TP 维做 column parallel，rank $t$ 保存：

[
\mat{W}_t
=========

\mat{W}*{:,\ t\frac{H*{\mathrm{out}}}{P_{\mathrm{TP}}}:(t+1)\frac{H_{\mathrm{out}}}{P_{\mathrm{TP}}}}.
]

Checkpoint metadata 需要记录：

\begin{itemize}
\item tensor name；
\item global shape；
\item shard shape；
\item shard offset；
\item dtype；
\item parallel dimension；
\item owning rank；
\item file path。
\end{itemize}

一个简化 metadata 可以写成：

\begin{lstlisting}[language=Python,caption={Sharded checkpoint metadata 的简化示意},label={lst:sharded-checkpoint-metadata}]
metadata = {
"global_step": 120000,
"consumed_tokens": 503316480000,
"parallel_config": {
"tp_size": 8,
"pp_size": 4,
"dp_size": 8,
"zero_stage": 1,
},
"tensors": {
"layers.0.attention.qkv.weight": [
{
"rank": 0,
"global_shape": [8192, 24576],
"local_shape": [8192, 3072],
"offset": [0, 0],
"dtype": "bfloat16",
"file": "rank_00000_model.pt",
},
{
"rank": 1,
"global_shape": [8192, 24576],
"local_shape": [8192, 3072],
"offset": [0, 3072],
"dtype": "bfloat16",
"file": "rank_00001_model.pt",
},
]
}
}
\end{lstlisting}

如果后续想把 TP size 从 8 改为 4，需要把原来的 8 个 column shards 重新拼接再切成 4 个 shards。这就是 checkpoint resharding。它要求 metadata 足够完整，否则无法安全转换。

Sharded checkpoint 的优点是：

\begin{itemize}
\item 保存时无需聚合完整模型；
\item 每个 rank 并行写入，IO 吞吐更高；
\item 与 ZeRO/FSDP/TP/PP 的分片状态自然匹配；
\item 大模型 checkpoint 可扩展到 TB 级。
\end{itemize}

缺点是：

\begin{itemize}
\item 文件数量多；
\item metadata 更复杂；
\item resharding 困难；
\item 部分 shard 损坏会影响整体恢复；
\item 需要与对象存储或并行文件系统良好配合。
\end{itemize}

\subsection{optimizer state checkpoint}
\label{subsec:optimizer-state-checkpoint}

仅保存模型参数不足以恢复训练。AdamW 的优化器状态包括：

[
m_t,\quad v_t,\quad \theta_{\mathrm{master}}.
]

如果不保存这些状态，只用模型权重恢复，等价于重新初始化优化器。训练可能继续下降，但 loss 曲线、收敛路径和最终结果都会改变。

优化器状态 checkpoint 的挑战是体积大。对于 AdamW，优化器状态通常比参数本身大得多。若模型参数为 $N$，优化器状态可能为：

[
12N\ \mathrm{bytes}
]

甚至更多，取决于是否保存 master weight、dtype 和 ZeRO stage。对于 70B 模型，优化器状态可以达到数百 GB 甚至更高。

使用 ZeRO 或 FSDP 时，优化器状态通常在 DP 维切分。Checkpoint 时每个 DP rank 只保存自己负责的 optimizer shard。恢复时必须保证：

[
\text{optimizer shard}
\leftrightarrow
\text{parameter shard}
]

正确对应。如果参数 resharding 了，优化器状态也要同步 reshard，否则会出现状态错配。

下面给出一个简化的 checkpoint 保存流程。

\begin{lstlisting}[language=Python,caption={保存模型与优化器 shard 的简化伪代码},label={lst:save-sharded-checkpoint-pseudocode}]
def save_sharded_checkpoint(
model,
optimizer,
scheduler,
step,
consumed_tokens,
rank,
checkpoint_dir,
parallel_context,
):
# Each rank saves only local shards.
model_state = collect_local_model_shards(model)
optim_state = collect_local_optimizer_shards(optimizer)

```
rank_state = {
    "model": model_state,
    "optimizer": optim_state,
    "scheduler": scheduler.state_dict(),
    "step": step,
    "consumed_tokens": consumed_tokens,
    "rng_state": get_rng_state(),
    "parallel_context": parallel_context.local_metadata(),
}

path = f"{checkpoint_dir}/rank_{rank:05d}.pt"
torch.save(rank_state, path)

# Rank 0 writes global metadata after all ranks finish writing shards.
torch.distributed.barrier()

if rank == 0:
    metadata = build_global_checkpoint_metadata(
        model=model,
        optimizer=optimizer,
        step=step,
        consumed_tokens=consumed_tokens,
        parallel_context=parallel_context,
    )
    write_json(metadata, f"{checkpoint_dir}/metadata.json")

torch.distributed.barrier()
```

\end{lstlisting}

真实 checkpoint 系统还需要处理原子性。不能让训练在 metadata 写完但部分 shard 没写完时误认为 checkpoint 可用。常见做法是：

\begin{itemize}
\item 写入临时目录；
\item 所有 rank 完成后写 success marker；
\item 最后原子 rename 或更新 latest 指针；
\item 恢复时只读取带完整 marker 的 checkpoint。
\end{itemize}

\subsection{checkpoint 与训练吞吐的冲突}
\label{subsec:checkpoint-vs-throughput}

Checkpoint 会中断或拖慢训练。设每隔 $K$ steps 保存一次 checkpoint，每次 checkpoint 花费 $T_{\mathrm{ckpt}}$，训练 step time 为 $T_{\mathrm{step}}$，则摊销开销为：

[
\rho_{\mathrm{ckpt}}
====================

\frac{T_{\mathrm{ckpt}}}{KT_{\mathrm{step}}}.
]

如果每 1000 steps 保存一次，每步 2 秒，checkpoint 花 10 分钟：

[
\rho_{\mathrm{ckpt}}
====================

# \frac{600}{1000\times 2}

30%.
]

这非常高。为了降低影响，可以：

\begin{itemize}
\item 降低 checkpoint 频率；
\item 使用异步 checkpoint；
\item 使用分布式并行写入；
\item 只保存必要状态，减少临时 checkpoint 体积；
\item 周期性保存 full checkpoint，期间保存 lightweight checkpoint；
\item 使用高吞吐并行文件系统或对象存储；
\item 将 checkpoint 写入与训练计算 overlap。
\end{itemize}

但 checkpoint 频率不能无限降低。故障后需要回退到最近 checkpoint。如果 checkpoint 间隔太长，一次故障会损失大量 GPU 时间。设平均故障间隔为 MTBF，checkpoint 间隔为 $K T_{\mathrm{step}}$，每次故障平均损失约为半个 checkpoint 间隔：

[
T_{\mathrm{lost}}
\approx
\frac{1}{2}K T_{\mathrm{step}}.
]

因此 checkpoint 策略需要在吞吐开销和故障损失之间平衡。

\begin{figure}[H]
\centering
\begin{tikzpicture}[
font=\small,
axis/.style={-{Latex[length=2.1mm]}, thick},
train/.style={draw, rounded corners=2pt, minimum width=1.0cm, minimum height=0.45cm, align=center, fill=blue!12},
ckpt/.style={draw, rounded corners=2pt, minimum width=0.75cm, minimum height=0.45cm, align=center, fill=orange!16},
fail/.style={draw, star, star points=5, minimum size=0.8cm, fill=red!20}
]

```
\node[font=\bfseries] at (0,3.1) {Checkpoint 间隔与故障回退成本};

\draw[axis] (-5.5,1.2) -- (5.5,1.2) node[right] {time};

\foreach \i/\x in {0/-4.5,1/-3.3,2/-2.1,3/-0.9,4/0.3,5/1.5,6/2.7,7/3.9} {
  \node[train] at (\x,1.2) {train};
}

\foreach \x in {-5.1,-1.5,2.1} {
  \node[ckpt] at (\x,2.0) {ckpt};
  \draw[dashed] (\x,1.55) -- (\x,1.2);
}

\node[fail] at (4.2,1.95) {};
\node[align=center] at (4.2,2.65) {failure};

\draw[red!70, thick, <->] (2.1,0.45) -- (4.2,0.45);
\node[red!70, align=center] at (3.15,0.05) {lost work\\since last checkpoint};

\node[align=left, text width=10.5cm] at (0,-1.35) {
  checkpoint 太频繁会降低吞吐；太稀疏会在故障后损失更多训练时间。
  大规模训练必须根据故障率、checkpoint 时长和训练成本选择合适间隔。
};
```

\end{tikzpicture}
\caption{checkpoint 间隔与故障恢复损失}
\label{fig:checkpoint-interval-failure-cost}
\end{figure}

\section{3D 并行配置实战检查清单}
\label{sec:3d-parallel-config-checklist}

为了把本章内容落到工程实践，下面给出一个 3D 并行配置检查清单。

\subsection{模型结构检查}
\label{subsec:model-structure-checklist}

配置前需要确认：

\begin{itemize}
\item hidden size 是否能整除 TP size；
\item attention heads 是否能整除 TP size；
\item kv heads 是否能整除 TP size 或支持对应 GQA 切分；
\item FFN intermediate size 是否能整除 TP size；
\item vocabulary size 是否需要 padding 到 TP 友好的倍数；
\item layer 数是否能被 PP size 合理切分；
\item embedding/lm head 是否共享权重；
\item 是否有 MoE 层，expert 数是否能与 expert parallel 匹配；
\item 是否使用 sequence parallel；
\item 是否使用 activation checkpointing。
\end{itemize}

\subsection{硬件拓扑检查}
\label{subsec:hardware-topology-checklist}

需要确认：

\begin{itemize}
\item 每节点 GPU 数；
\item 节点内 NVLink/NVSwitch/PCIe 拓扑；
\item 每节点 NIC 数；
\item GPU 到 NIC 的亲和性；
\item 跨节点网络带宽；
\item 是否多 rail；
\item rank mapping 是否让 TP group 尽量节点内；
\item PP 相邻 stage 是否落在合理网络距离；
\item DP group 是否均匀分布；
\item NCCL 是否识别正确拓扑。
\end{itemize}

\subsection{训练配置检查}
\label{subsec:training-config-checklist}

需要确认：

\begin{itemize}
\item $W=P_{\mathrm{DP}}P_{\mathrm{TP}}P_{\mathrm{PP}}$；
\item $B_{\mathrm{global}}=B_{\mathrm{micro}}MP_{\mathrm{DP}}$；
\item $M$ 是否足够大以降低 pipeline bubble；
\item micro-batch size 是否能放入显存；
\item global batch 是否符合训练 recipe；
\item 学习率是否按 global batch 调整；
\item warmup steps 是否按 token 或 step 正确设置；
\item gradient clipping 是否在正确维度上执行；
\item mixed precision 与 loss scaling 是否配置正确；
\item activation checkpointing 粒度是否合理。
\end{itemize}

\subsection{运行时检查}
\label{subsec:runtime-checklist}

训练启动后，需要观察：

\begin{itemize}
\item 每张 GPU 显存峰值是否有余量；
\item 每个 rank step time 是否一致；
\item pipeline stage 是否负载均衡；
\item TP collective 是否占比过高；
\item DP gradient sync 是否被 backward 隐藏；
\item PP send/recv 是否出现等待；
\item dataloader 是否阻塞；
\item loss 是否稳定下降；
\item gradient norm 是否异常；
\item 是否出现 NaN/inf；
\item checkpoint 是否成功写入并可恢复；
\item tokens per second 和 MFU 是否达到预期。
\end{itemize}

\section{本章小结}
\label{sec:chapter10-summary}

本章把数据并行、张量并行和流水线并行组合起来，系统介绍了 3D 并行策略。

\begin{itemize}
\item \textbf{3D 并行} 将 DP、TP、PP 三个维度正交组合，总 GPU 数满足 $W=P_{\mathrm{DP}}P_{\mathrm{TP}}P_{\mathrm{PP}}$。
\item \textbf{DP 切分数据}，不同副本处理不同 batch 分片，并同步梯度或分片状态。
\item \textbf{TP 切分层内张量}，例如 attention heads、QKV projection、FFN 中间维度，通信频率最高，通常应限制在节点内高速互联范围。
\item \textbf{PP 切分模型层}，不同 stages 负责不同层段，通过 micro-batches 填充流水线，通信主要是相邻 stage 的 activation send/recv。
\item \textbf{group 划分} 是 3D 并行的实现基础。每个 rank 同时属于 TP group、PP group 和 DP group。
\item \textbf{rank mapping} 直接影响通信路径。高频通信的 group 应优先映射到更近、更快的物理拓扑。
\item \textbf{显存估算} 需要同时考虑参数、梯度、优化器状态、激活和 buffer，并区分哪些状态被 TP/PP/DP/ZeRO 切分。
\item \textbf{通信估算} 需要分开看 DP、TP、PP：DP 主要同步梯度，TP 高频层内 collective，PP 传递 stage 边界激活。
\item \textbf{计算效率} 受 GEMM shape、pipeline bubble、通信重叠、数据加载、checkpoint 和 straggler 影响，MFU 是重要观察指标。
\item \textbf{训练配置推导} 要满足 global batch、micro-batch、gradient accumulation、TP/PP/DP 整除关系和硬件拓扑约束。
\item \textbf{Distributed checkpoint} 是大规模训练可靠性的核心。checkpoint 需要保存模型、优化器、scheduler、RNG、consumed tokens、parallel metadata 和 shard layout。
\item \textbf{Checkpoint 策略} 需要在训练吞吐和故障回退成本之间平衡。频率太高会拖慢训练，频率太低会在故障后损失大量工作。
\end{itemize}

至此，我们已经建立了大模型训练最核心的并行框架：DP 解决数据扩展，TP 解决层内矩阵过大，PP 解决层数和整体模型过大，3D 并行则把三者组合为可扩展的训练系统。下一章将进入 ZeRO 优化器与 DeepSpeed，进一步分析如何在数据并行维度切分 optimizer state、gradient 和 parameter，从而显著降低每卡显存占用，并理解 ZeRO Stage 1、Stage 2、Stage 3 的通信与显存权衡。

\chapter{ZeRO 优化器与 DeepSpeed}
\label{chap:zero-deepspeed}

\section{ZeRO 的问题背景}
\label{sec:zero-background}

前面几章讨论了 DP、TP、PP 与 3D 并行。数据并行通过复制模型来提升吞吐；张量并行切分层内矩阵；流水线并行切分模型层；3D 并行将三者组合，解决超大模型的训练扩展问题。

但即使使用 3D 并行，我们仍然会遇到一个非常现实的问题：\textbf{训练状态太大}。

在大模型训练中，显存不只被模型参数占用。对于每个参数，训练系统还要保存梯度、优化器状态、可能的 FP32 master weight、通信 buffer、激活和临时 workspace。尤其使用 AdamW 时，优化器状态通常比参数本身大得多。

ZeRO，全称 Zero Redundancy Optimizer，其核心目标是消除数据并行中的冗余训练状态。普通 DDP 中，每个 data parallel rank 都保存完整参数、完整梯度和完整优化器状态。ZeRO 则把这些状态在数据并行 rank 之间切分，让每个 rank 只保存一部分，从而显著降低每卡显存。

从一句话理解：

[
\text{DDP: 每个 DP rank 都保存完整训练状态。}
]
[
\text{ZeRO: 在 DP ranks 之间切分训练状态，减少冗余。}
]

ZeRO 并不替代数据并行。相反，它是在数据并行维度上的显存优化。ZeRO 也可以与 TP、PP、Sequence Parallel、Activation Checkpointing、CPU/NVMe Offload 等技术组合使用。

\subsection{Adam 优化器状态膨胀}
\label{subsec:adam-state-explosion}

先从显存账本开始。假设模型参数量为 $N$。若使用 BF16 参数训练，参数本身占用：

[
M_{\mathrm{param}}=2N\ \mathrm{bytes}.
]

反向传播后，每个参数有对应梯度。如果梯度也使用 BF16 保存：

[
M_{\mathrm{grad}}=2N\ \mathrm{bytes}.
]

AdamW 通常维护两个 FP32 状态：

[
m_t,\quad v_t,
]
分别是一阶矩和二阶矩。它们各占：

[
4N\ \mathrm{bytes}.
]

很多混合精度训练还会保留 FP32 master weight，用于更稳定的优化器更新：

[
M_{\mathrm{master}}=4N\ \mathrm{bytes}.
]

因此，仅参数相关训练状态就可能达到：

[
M_{\mathrm{train\ states}}
==========================

2N
+
2N
+
4N
+
4N
+
4N
==

16N\ \mathrm{bytes}.
]

这就是常说的 AdamW 训练状态约为每参数 16 bytes。注意这个数字还没有包括 activation、临时 buffer、通信 buffer、CUDA allocator fragment、dropout mask、attention 中间结果、checkpoint buffer 等。

如果模型参数量为 $7$B，则仅训练状态约为：

[
16\times 7\times 10^9
=====================

112\times 10^9\ \mathrm{bytes}
\approx 112\ \mathrm{GB}.
]

如果模型参数量为 $70$B：

[
16\times 70\times 10^9
======================

1120\ \mathrm{GB}
\approx 1.12\ \mathrm{TB}.
]

这说明，大模型训练中的显存压力并不只是“参数太大”，而是“参数相关状态成倍膨胀”。

\begin{table}[H]
\centering
\small
\caption{混合精度 AdamW 训练中常见的参数相关状态}
\label{tab:adamw-state-memory}
\begin{tabular}{p{0.24\textwidth}p{0.24\textwidth}p{0.18\textwidth}p{0.24\textwidth}}
\toprule
\textbf{状态} & \textbf{用途} & \textbf{常见 dtype} & \textbf{每参数字节数} \
\midrule
参数 $\theta$
& 前向和反向计算
& BF16 / FP16
& 2 bytes \
\midrule
梯度 $\nabla\theta$
& 反向传播结果
& BF16 / FP16 / FP32
& 常见 2 或 4 bytes \
\midrule
FP32 master weight
& 优化器高精度更新
& FP32
& 4 bytes \
\midrule
Adam 一阶矩 $m$
& 动量估计
& FP32
& 4 bytes \
\midrule
Adam 二阶矩 $v$
& 方差估计
& FP32
& 4 bytes \
\bottomrule
\end{tabular}
\end{table}

如果使用 SGD，优化器状态可能更少；如果使用 Adafactor、8-bit optimizer 或其他低精度优化器，状态也可能下降。但 AdamW 仍是大模型训练中最常见的基线之一，因此理解 AdamW 状态膨胀非常重要。

\subsection{DDP 中的冗余}
\label{subsec:ddp-redundancy}

普通 DDP 中，每个 data parallel rank 保存完整模型副本。假设 DP size 为 $P_{\mathrm{DP}}$，每个 rank 的参数状态为：

[
\theta^{(r)},\quad
r=0,\ldots,P_{\mathrm{DP}}-1.
]

DDP 的不变量是所有 rank 的参数相同：

[
\theta^{(0)}=\theta^{(1)}=\cdots=\theta^{(P_{\mathrm{DP}}-1)}.
]

每个 rank 处理不同数据分片，得到本地梯度：

[
g^{(r)}.
]

然后通过 all-reduce 得到平均梯度：

[
\bar{g}
=======

\frac{1}{P_{\mathrm{DP}}}
\sum_{r=0}^{P_{\mathrm{DP}}-1}
g^{(r)}.
]

所有 rank 使用相同梯度更新参数，因此继续保持一致。

这个设计简单可靠，但存在显存冗余。对于同一个参数张量，所有 DP ranks 都保存：

\begin{itemize}
\item 完整参数；
\item 完整梯度；
\item 完整优化器状态；
\item 完整 FP32 master weight。
\end{itemize}

从数据并行角度看，这些状态是冗余的。因为所有 DP ranks 上的参数最终应该相同，优化器状态也逻辑相同。ZeRO 的核心问题是：\textbf{既然 DP ranks 之间有冗余，为什么不把这些状态分片保存？}

\begin{figure}[H]
\centering
\begin{tikzpicture}[
font=\small,
rank/.style={draw, rounded corners=4pt, minimum width=2.0cm, minimum height=1.2cm, align=center, fill=blue!6},
state/.style={draw, rounded corners=2pt, minimum width=1.55cm, minimum height=0.35cm, align=center, fill=orange!14},
arrow/.style={-{Latex[length=2.2mm]}, thick}
]

```
\node[font=\bfseries] at (0,3.25) {DDP 中每个 rank 保存完整训练状态};

\foreach \i/\x in {0/-4.5,1/-1.5,2/1.5,3/4.5} {
  \node[rank] (r\i) at (\x,1.1) {DP Rank \i};
  \node[state] at (\x,1.45) {params};
  \node[state] at (\x,1.05) {grads};
  \node[state] at (\x,0.65) {optimizer};
}

\draw[<->, thick, red!70] (-3.5,-0.15) -- (3.5,-0.15)
  node[midway,below=4pt] {gradient all-reduce};

\node[align=left, text width=10.5cm] at (0,-1.6) {
  DDP 的简单性来自完整复制：每个 rank 都有完整模型和完整优化器状态。
  但这种复制使显存随 DP size 线性重复，而不是被 DP size 分摊。
};
```

\end{tikzpicture}
\caption{DDP 中参数、梯度和优化器状态的冗余复制}
\label{fig:ddp-redundant-states}
\end{figure}

\subsection{ZeRO 的核心思想}
\label{subsec:zero-core-idea}

ZeRO 将训练状态分为三类：

\begin{itemize}
\item Optimizer States；
\item Gradients；
\item Parameters。
\end{itemize}

然后分阶段将这些状态在 data parallel ranks 之间切分。

[
\text{ZeRO Stage 1: partition optimizer states}
]
[
\text{ZeRO Stage 2: partition optimizer states + gradients}
]
[
\text{ZeRO Stage 3: partition optimizer states + gradients + parameters}
]

如果 DP size 为 $P_{\mathrm{DP}}$，被切分的状态理想情况下每卡下降为：

[
\frac{1}{P_{\mathrm{DP}}}.
]

这种思想与 TP/PP 不同。TP/PP 是模型并行，切的是层内矩阵或模型层；ZeRO 是数据并行维度上的状态分片，切的是 DP replica 之间原本冗余的训练状态。

可以把 ZeRO 看作对 DDP 的显存账本重写：

[
M_{\mathrm{DDP}}
================

M_{\theta}
+
M_g
+
M_{\mathrm{optim}},
]

[
M_{\mathrm{ZeRO1}}
==================

M_{\theta}
+
M_g
+
\frac{M_{\mathrm{optim}}}{P_{\mathrm{DP}}},
]

[
M_{\mathrm{ZeRO2}}
==================

M_{\theta}
+
\frac{M_g}{P_{\mathrm{DP}}}
+
\frac{M_{\mathrm{optim}}}{P_{\mathrm{DP}}},
]

[
M_{\mathrm{ZeRO3}}
==================

\frac{M_{\theta}}{P_{\mathrm{DP}}}
+
\frac{M_g}{P_{\mathrm{DP}}}
+
\frac{M_{\mathrm{optim}}}{P_{\mathrm{DP}}}
+
M_{\mathrm{gather\ buffers}}.
]

这里 $M_{\mathrm{gather\ buffers}}$ 是 Stage 3 中为了计算某层而临时 all-gather 参数所需的显存。Stage 3 显存最低，但通信和实现复杂度最高。

\section{ZeRO Stage 1}
\label{sec:zero-stage1}

ZeRO Stage 1 切分优化器状态，但参数和梯度仍然在每个 DP rank 上完整保存。它是 ZeRO 中最容易理解、通信变化相对较小的一档。

\subsection{partition optimizer states}
\label{subsec:zero1-partition-optimizer-states}

假设模型参数为：

[
\theta=
{\theta_0,\theta_1,\ldots,\theta_{K-1}}.
]

在 DDP 中，每个 rank 保存所有参数的 Adam 状态：

[
{m_i,v_i,\theta_{master,i}}_{i=0}^{K-1}.
]

ZeRO Stage 1 将这些 optimizer states 按参数分片给不同 DP ranks。若 $P_{\mathrm{DP}}=4$，则 rank 0 可能负责参数 0 到 25% 的 optimizer states，rank 1 负责 25% 到 50%，以此类推。

每个 rank 仍然保存完整参数 $\theta$ 和完整梯度 $g$，但只保存自己负责参数分片的优化器状态：

[
\mathcal{O}^{(r)}
=================

{m_i,v_i,\theta_{master,i}\mid i\in\mathcal{I}_r}.
]

其中 $\mathcal{I}_r$ 是 rank $r$ 负责的参数索引集合。

优化器 step 时，每个 rank 只更新自己负责的参数 shard：

[
\theta_{\mathcal{I}*r}
\leftarrow
\operatorname{AdamWUpdate}
(
\theta*{\mathcal{I}*r},
g*{\mathcal{I}*r},
m*{\mathcal{I}*r},
v*{\mathcal{I}_r}
).
]

更新完成后，需要让所有 DP ranks 得到完整、最新的参数。因为 Stage 1 中参数仍然是 replicated 的，所以更新后的参数 shards 需要同步到所有 ranks。这个同步可以通过 all-gather、broadcast 或等价的通信实现。

\begin{figure}[H]
\centering
\begin{tikzpicture}[
font=\small,
rank/.style={draw, rounded corners=4pt, minimum width=2.0cm, minimum height=1.4cm, align=center, fill=blue!6},
full/.style={draw, rounded corners=2pt, minimum width=1.55cm, minimum height=0.35cm, align=center, fill=orange!14},
shard/.style={draw, rounded corners=2pt, minimum width=1.55cm, minimum height=0.35cm, align=center, fill=green!12}
]

```
\node[font=\bfseries] at (0,3.35) {ZeRO Stage 1：切分优化器状态};

\foreach \i/\x in {0/-4.5,1/-1.5,2/1.5,3/4.5} {
  \node[rank] at (\x,1.2) {Rank \i};
  \node[full] at (\x,1.65) {full params};
  \node[full] at (\x,1.2) {full grads};
  \node[shard] at (\x,0.75) {optim shard};
}

\node[align=left, text width=10.5cm] at (0,-0.95) {
  Stage 1 只切分 optimizer states。参数和梯度仍然完整复制，因此实现相对简单。
  显存收益主要来自 Adam 状态和 master weight 的分片。
};
```

\end{tikzpicture}
\caption{ZeRO Stage 1 的状态分片方式}
\label{fig:zero-stage1-states}
\end{figure}

\subsection{显存节省}
\label{subsec:zero1-memory-saving}

设参数使用 BF16，梯度使用 BF16，FP32 master weight、Adam $m$、Adam $v$ 各 4 bytes。DDP 每参数训练状态约为：

[
2+2+4+4+4=16\ \mathrm{bytes}.
]

ZeRO Stage 1 中，参数和梯度仍完整保存：

[
M_{\theta}=2N,
\quad
M_g=2N.
]

优化器相关状态为：

[
M_{\mathrm{optim}}=12N,
]

并被 $P_{\mathrm{DP}}$ 分片：

[
M_{\mathrm{optim,perGPU}}
=========================

\frac{12N}{P_{\mathrm{DP}}}.
]

因此每卡参数相关状态约为：

[
M_{\mathrm{ZeRO1}}
==================

4N+\frac{12N}{P_{\mathrm{DP}}}.
]

当 $P_{\mathrm{DP}}=8$ 时：

[
M_{\mathrm{ZeRO1}}
==================

# 4N+\frac{12N}{8}

5.5N\ \mathrm{bytes}.
]

相比 DDP 的 $16N$：

[
\frac{5.5}{16}
\approx 34.4%.
]

这已经是非常明显的显存节省。

\subsection{通信特征}
\label{subsec:zero1-communication}

Stage 1 的通信主要包括两部分：

\begin{itemize}
\item backward 后仍然需要梯度同步；
\item optimizer step 后需要同步更新后的参数。
\end{itemize}

由于参数和梯度在每个 rank 上完整存在，Stage 1 与 DDP 在训练流程上比较接近。它的主要复杂度在 optimizer step：每个 rank 只更新自己负责的参数分片，因此更新结果需要传播给所有 rank。

Stage 1 的优点是显存收益大、实现相对容易；缺点是梯度仍然完整保存，参数也完整复制，显存节省有上限。

\section{ZeRO Stage 2}
\label{sec:zero-stage2}

ZeRO Stage 2 在 Stage 1 的基础上进一步切分梯度。它保存完整参数，但优化器状态和梯度都在 DP ranks 之间分片。

\subsection{partition gradients}
\label{subsec:zero2-partition-gradients}

DDP 的梯度同步通常是 all-reduce。每个 rank 输入完整梯度，输出完整平均梯度。ZeRO Stage 2 的目标是：既然每个参数分片的优化器状态只在某个 rank 上，那么该参数分片的梯度也只需要被那个 rank 保存和用于更新。

因此 Stage 2 用 reduce-scatter 替代完整 all-reduce。对完整梯度 $g$，按参数维切分为：

[
g=[g_0,g_1,\ldots,g_{P_{\mathrm{DP}}-1}].
]

每个 rank 本地计算完整梯度贡献 $g^{(r)}$，然后 reduce-scatter：

[
\bar{g}_r
=========

\frac{1}{P_{\mathrm{DP}}}
\sum_{q=0}^{P_{\mathrm{DP}}-1}
g_r^{(q)}.
]

最终 rank $r$ 只获得自己负责的平均梯度 shard $\bar{g}_r$，而不是完整平均梯度。

这与 DDP all-reduce 的区别是：

[
\operatorname{AllReduce}
========================

\operatorname{ReduceScatter}
+
\operatorname{AllGather}.
]

DDP 做完整 all-reduce，相当于 reduce-scatter 后再 all-gather，让所有 rank 拿到完整梯度。ZeRO Stage 2 只需要 reduce-scatter，因为每个 rank 只需要自己负责的梯度分片来更新对应参数。

\begin{figure}[H]
\centering
\begin{tikzpicture}[
font=\small,
rank/.style={draw, rounded corners=4pt, minimum width=2.0cm, minimum height=1.35cm, align=center, fill=blue!6},
full/.style={draw, rounded corners=2pt, minimum width=1.55cm, minimum height=0.35cm, align=center, fill=orange!14},
shard/.style={draw, rounded corners=2pt, minimum width=1.55cm, minimum height=0.35cm, align=center, fill=green!12}
]

```
\node[font=\bfseries] at (0,3.35) {ZeRO Stage 2：切分优化器状态与梯度};

\foreach \i/\x in {0/-4.5,1/-1.5,2/1.5,3/4.5} {
  \node[rank] at (\x,1.2) {Rank \i};
  \node[full] at (\x,1.65) {full params};
  \node[shard] at (\x,1.2) {grad shard};
  \node[shard] at (\x,0.75) {optim shard};
}

\node[align=left, text width=10.5cm] at (0,-0.95) {
  Stage 2 使用 reduce-scatter 聚合梯度，每个 rank 只保留自己负责的梯度 shard。
  参数仍然完整复制，因此 forward/backward 不需要频繁 all-gather 参数。
};
```

\end{tikzpicture}
\caption{ZeRO Stage 2 的状态分片方式}
\label{fig:zero-stage2-states}
\end{figure}

\subsection{reduce-scatter}
\label{subsec:zero2-reduce-scatter}

Reduce-scatter 的语义是：所有 ranks 输入相同 shape 的 tensor，先按元素规约求和，再把结果分片散发给不同 ranks。

对于 $P$ 个 ranks，每个 rank 输入：

[
\mat{G}^{(r)}\in\mathbb{R}^{N}.
]

Reduce-scatter 后，rank $r$ 得到：

[
\mat{Y}^{(r)}
=============

\left[
\sum_{q=0}^{P-1}\mat{G}^{(q)}
\right]_{\mathrm{shard}\ r}.
]

如果需要平均梯度，再除以 $P$。

与 all-reduce 相比，reduce-scatter 的输出大小只有原来的：

[
\frac{1}{P}.
]

这直接降低了每 rank 保存梯度的显存。

\begin{figure}[H]
\centering
\begin{tikzpicture}[
font=\small,
chunk/.style={draw, rounded corners=2pt, minimum width=0.65cm, minimum height=0.38cm, align=center},
rank/.style={draw, rounded corners=3pt, minimum width=1.5cm, minimum height=0.55cm, align=center, fill=blue!6},
op/.style={draw, rounded corners=4pt, minimum width=2.1cm, minimum height=0.7cm, align=center, fill=orange!14},
arrow/.style={-{Latex[length=2.1mm]}, thick}
]

```
\node[font=\bfseries] at (0,3.25) {Reduce-Scatter：规约后只保留本 rank 的 shard};

\foreach \i/\x in {0/-4.5,1/-1.5,2/1.5,3/4.5} {
  \node[rank] (r\i) at (\x,1.65) {Rank \i};
  \foreach \j/\dx/\col in {0/-0.45/red!12,1/-0.15/green!12,2/0.15/purple!12,3/0.45/yellow!18} {
    \node[chunk, fill=\col] at (\x+\dx,1.05) {};
  }
}

\node[op] (rs) at (0,0.0) {ReduceScatter\\sum + shard};

\foreach \i in {0,1,2,3} {
  \draw[arrow] (r\i) -- (rs);
}

\foreach \i/\x/\col in {0/-4.5/red!12,1/-1.5/green!12,2/1.5/purple!12,3/4.5/yellow!18} {
  \node[rank] (out\i) at (\x,-1.25) {Rank \i};
  \node[chunk, fill=\col, minimum width=1.1cm] at (\x,-1.85) {reduced shard};
  \draw[arrow] (rs) -- (out\i);
}

\node[align=left, text width=10.5cm] at (0,-2.75) {
  Stage 2 不需要每个 rank 保存完整平均梯度。
  Reduce-scatter 直接把规约后的梯度分片交给对应 optimizer shard 的 owner rank。
};
```

\end{tikzpicture}
\caption{ZeRO Stage 2 中 reduce-scatter 梯度聚合}
\label{fig:zero2-reduce-scatter}
\end{figure}

\subsection{显存节省}
\label{subsec:zero2-memory-saving}

Stage 2 中，参数完整保存，梯度和优化器状态分片。若仍使用 BF16 参数和梯度、FP32 master/m/v，则：

[
M_{\theta}=2N.
]

梯度分片：

[
M_g=\frac{2N}{P_{\mathrm{DP}}}.
]

优化器状态分片：

[
M_{\mathrm{optim}}=\frac{12N}{P_{\mathrm{DP}}}.
]

所以：

[
M_{\mathrm{ZeRO2}}
==================

2N+\frac{14N}{P_{\mathrm{DP}}}.
]

当 $P_{\mathrm{DP}}=8$：

[
M_{\mathrm{ZeRO2}}
==================

# 2N+\frac{14N}{8}

3.75N\ \mathrm{bytes}.
]

相比 DDP 的 $16N$：

[
\frac{3.75}{16}
\approx 23.4%.
]

Stage 2 比 Stage 1 进一步节省了梯度显存，但参数仍然完整复制，因此显存下限仍有参数本身的 $2N$。

\subsection{通信特征}
\label{subsec:zero2-communication}

Stage 2 的 backward 梯度同步使用 reduce-scatter，而不是 all-reduce。通信量与 DDP all-reduce 同量级，但输出显存更小，并且不需要完整梯度 all-gather。

不过，由于参数仍然完整复制，optimizer step 后仍然需要让所有 ranks 获得一致的参数。这通常涉及将更新后的参数 shard 同步回所有 ranks。

Stage 2 的优点是：

\begin{itemize}
\item 显存比 Stage 1 更低；
\item forward/backward 不需要按层 all-gather 参数；
\item 比 Stage 3 通信模式简单；
\item 适合许多中等到大型模型训练。
\end{itemize}

Stage 2 的缺点是：

\begin{itemize}
\item 参数仍然完整复制，模型非常大时仍可能放不下；
\item optimizer step 后参数同步仍有成本；
\item 与 TP/PP 组合时，梯度分片 owner 必须与参数 shard 映射一致。
\end{itemize}

\section{ZeRO Stage 3}
\label{sec:zero-stage3}

ZeRO Stage 3 是最彻底的一档。它将优化器状态、梯度和参数都在 DP ranks 之间分片。每个 rank 只保存一部分参数。需要某层参数参与 forward/backward 时，再临时 all-gather 该层参数，用完后释放。

Stage 3 的显存最低，但通信最复杂，也最依赖实现质量。

\subsection{partition parameters}
\label{subsec:zero3-partition-parameters}

Stage 3 中，每个参数张量也被切分到 DP ranks 上。设参数向量化后为：

[
\theta=[\theta_0,\theta_1,\ldots,\theta_{P_{\mathrm{DP}}-1}].
]

Rank $r$ 只常驻保存：

[
\theta_r.
]

这意味着在任意时刻，单个 rank 不再拥有完整模型参数。为了执行某一层的 forward，需要把该层参数从各 DP ranks all-gather 到当前计算 rank 组中。

对于某层权重 $\mat{W}^{(\ell)}$：

[
\mat{W}^{(\ell)}
================

\operatorname{AllGather}
(
\mat{W}^{(\ell)}*0,
\mat{W}^{(\ell)}*1,
\ldots,
\mat{W}^{(\ell)}*{P*{\mathrm{DP}}-1}
).
]

Forward 使用完整 $\mat{W}^{(\ell)}$ 计算。Backward 也需要对应参数用于梯度计算。用完后，可以释放完整参数，只保留本 rank 的 shard。

\begin{figure}[H]
\centering
\begin{tikzpicture}[
font=\small,
rank/.style={draw, rounded corners=4pt, minimum width=2.0cm, minimum height=1.45cm, align=center, fill=blue!6},
shard/.style={draw, rounded corners=2pt, minimum width=1.5cm, minimum height=0.35cm, align=center, fill=green!12},
temp/.style={draw, rounded corners=2pt, minimum width=1.5cm, minimum height=0.35cm, align=center, fill=orange!14},
op/.style={draw, rounded corners=4pt, minimum width=2.4cm, minimum height=0.65cm, align=center, fill=purple!10},
arrow/.style={-{Latex[length=2.1mm]}, thick}
]

```
\node[font=\bfseries] at (0,3.45) {ZeRO Stage 3：参数也被分片，计算前临时 AllGather};

\foreach \i/\x in {0/-4.5,1/-1.5,2/1.5,3/4.5} {
  \node[rank] (r\i) at (\x,1.55) {Rank \i};
  \node[shard] at (\x,1.95) {param shard};
  \node[shard] at (\x,1.55) {grad shard};
  \node[shard] at (\x,1.15) {optim shard};
}

\node[op] (ag) at (0,0.15) {AllGather layer params};

\foreach \i in {0,1,2,3} {
  \draw[arrow] (r\i) -- (ag);
}

\foreach \i/\x in {0/-4.5,1/-1.5,2/1.5,3/4.5} {
  \node[temp] at (\x,-0.95) {full layer temp};
  \draw[arrow] (ag) -- (\x,-0.7);
}

\node[align=left, text width=10.5cm] at (0,-2.2) {
  Stage 3 常驻状态都是 sharded 的。某层计算前临时聚合完整参数，用完后释放。
  因此显存最低，但 forward/backward 中会出现频繁参数 all-gather。
};
```

\end{tikzpicture}
\caption{ZeRO Stage 3 的参数分片与临时 all-gather}
\label{fig:zero-stage3-param-gather}
\end{figure}

\subsection{all-gather 参数}
\label{subsec:zero3-allgather-parameters}

Stage 3 的关键性能问题是参数 all-gather。如果每一层都在使用前同步等待参数聚合，GPU 会大量空等。高性能实现必须做：

\begin{itemize}
\item \textbf{bucketed all-gather}：把多个小参数合并成较大 bucket；
\item \textbf{prefetch}：在当前层计算时提前聚合后续层参数；
\item \textbf{release}：某层参数不再需要时尽快释放完整副本；
\item \textbf{reuse distance 控制}：对即将再次使用的参数暂时保留，避免反复 gather；
\item \textbf{overlap communication and computation}：尽量隐藏参数 all-gather 时间。
\end{itemize}

Stage 3 的 forward 可以抽象为：

[
\text{prefetch params of layer }\ell+1,
]
[
\text{gather params of layer }\ell,
]
[
\text{compute layer }\ell,
]
[
\text{release params of layer }\ell-k.
]

Backward 也类似，但方向相反。由于 backward 需要参数计算梯度，参数释放时机要特别小心。

下面给出一个简化的 ZeRO-3 forward 伪代码。

\begin{lstlisting}[language=Python,caption={ZeRO Stage 3 forward 中参数 gather/prefetch/release 的伪代码},label={lst:zero3-forward-pseudocode}]
def zero3_forward(layers, input_tensor, param_manager):
x = input_tensor

```
for layer_id, layer in enumerate(layers):
    # Start prefetching future layer parameters if possible.
    if layer_id + 1 < len(layers):
        param_manager.prefetch(layer_id + 1)

    # Ensure current layer parameters are available.
    param_manager.all_gather(layer_id)

    # Compute with full temporary parameters.
    x = layer.forward(x)

    # Release parameters that will not be reused soon.
    param_manager.release_if_safe(layer_id)

return x
```

\end{lstlisting}

真实实现要处理 autograd、backward、参数生命周期、gradient hooks、bucket、CUDA stream、通信 group、CPU/NVMe offload 和异常恢复，复杂度远高于这段伪代码。

\subsection{显存节省}
\label{subsec:zero3-memory-saving}

Stage 3 中，参数、梯度、优化器状态都被 DP size 分片。理想情况下，每参数常驻显存约为：

[
M_{\mathrm{ZeRO3,resident}}
===========================

\frac{16N}{P_{\mathrm{DP}}}.
]

但这只是常驻状态。计算时需要临时聚合某些层的参数，因此要加上 gather buffer：

[
M_{\mathrm{ZeRO3}}
==================

\frac{16N}{P_{\mathrm{DP}}}
+
M_{\mathrm{param\ gather}}
+
M_{\mathrm{activations}}
+
M_{\mathrm{buffers}}.
]

如果 all-gather bucket 太大，$M_{\mathrm{param\ gather}}$ 会很高；bucket 太小，通信效率下降。因此 Stage 3 的性能高度依赖 bucket 配置和参数生命周期管理。

当 $P_{\mathrm{DP}}=8$ 时，常驻参数相关状态理论上为：

[
\frac{16N}{8}=2N\ \mathrm{bytes}.
]

这非常低，甚至接近原始 BF16 参数大小。但 Stage 3 的通信代价也最高。

\subsection{通信特征}
\label{subsec:zero3-communication}

Stage 3 的通信包括：

\begin{itemize}
\item forward 前参数 all-gather；
\item backward 前或 backward 中参数 all-gather；
\item backward 后梯度 reduce-scatter；
\item optimizer step 只更新本地 shard；
\item 可能的参数 prefetch 和释放；
\item 若使用 offload，还包含 GPU 与 CPU/NVMe 之间的数据搬运。
\end{itemize}

Stage 3 的通信量可能很大，但它允许训练远超单卡显存限制的模型。它非常适合以下场景：

\begin{itemize}
\item 模型太大，Stage 1/2 仍放不下；
\item DP size 足够大，可以有效分摊参数；
\item 网络带宽较好；
\item 实现支持通信计算重叠；
\item 可以接受一定训练吞吐下降换取显存可行性。
\end{itemize}

Stage 3 不适合盲目开启。对于能用 Stage 1/2 或 FSDP 更简单配置跑起来的模型，Stage 3 可能因为参数 all-gather 造成额外开销。工程上常见做法是：先用较低 stage 满足显存，再逐步提高 stage，而不是一开始就默认使用 Stage 3。

\section{DeepSpeed 实践}
\label{sec:deepspeed-practice}

DeepSpeed 是围绕大模型训练和推理构建的一套系统工具，其中 ZeRO 是其最著名的训练优化之一。本节不介绍安装细节，而关注 DeepSpeed 配置与训练系统设计思想。

DeepSpeed 典型使用方式是：用户定义 PyTorch 模型、优化器和数据，然后通过 DeepSpeed engine 接管分布式训练、混合精度、ZeRO 分片、梯度累积、通信和 checkpoint。

\subsection{DeepSpeed engine}
\label{subsec:deepspeed-engine}

普通 PyTorch 训练循环大致是：

\begin{lstlisting}[language=Python,caption={普通 PyTorch 训练循环的抽象形式},label={lst:normal-pytorch-loop}]
for batch in dataloader:
optimizer.zero_grad()
output = model(batch)
loss = loss_fn(output, batch["labels"])
loss.backward()
optimizer.step()
\end{lstlisting}

DeepSpeed 会把模型、优化器、scheduler 包装到 engine 中。训练循环变为：

\begin{lstlisting}[language=Python,caption={DeepSpeed engine 训练循环的抽象形式},label={lst:deepspeed-engine-loop}]
import deepspeed

model_engine, optimizer, _, scheduler = deepspeed.initialize(
model=model,
model_parameters=model.parameters(),
config=ds_config,
)

for batch in dataloader:
output = model_engine(batch)
loss = compute_loss(output, batch["labels"])

```
model_engine.backward(loss)
model_engine.step()
```

\end{lstlisting}

这看起来只是把 \texttt{loss.backward()} 换成 \texttt{model_engine.backward(loss)}，把 \texttt{optimizer.step()} 换成 \texttt{model_engine.step()}。但背后发生了大量事情：

\begin{itemize}
\item 按配置启用 ZeRO Stage；
\item 管理参数、梯度、优化器状态分片；
\item 处理 gradient accumulation；
\item 执行 mixed precision；
\item 管理通信 bucket；
\item 在 backward hook 中触发 reduce-scatter；
\item Stage 3 中进行参数 all-gather/prefetch/release；
\item 管理 checkpoint 保存与加载；
\item 与分布式 process group 协作。
\end{itemize}

DeepSpeed engine 的价值在于把这些复杂逻辑封装起来，让用户仍能以相对接近 PyTorch 的方式写训练代码。

\subsection{配置文件}
\label{subsec:deepspeed-config}

DeepSpeed 的核心行为由 JSON 配置控制。下面是一个简化 ZeRO Stage 2 配置：

\begin{lstlisting}[language=json,caption={DeepSpeed ZeRO Stage 2 配置示例},label={lst:deepspeed-zero2-config}]
{
"train_batch_size": 1024,
"train_micro_batch_size_per_gpu": 1,
"gradient_accumulation_steps": 128,

"bf16": {
"enabled": true
},

"optimizer": {
"type": "AdamW",
"params": {
"lr": 3e-4,
"betas": [0.9, 0.95],
"eps": 1e-8,
"weight_decay": 0.1
}
},

"zero_optimization": {
"stage": 2,
"reduce_scatter": true,
"allgather_partitions": true,
"allgather_bucket_size": 500000000,
"reduce_bucket_size": 500000000,
"overlap_comm": true,
"contiguous_gradients": true
},

"gradient_clipping": 1.0,
"steps_per_print": 100
}
\end{lstlisting}

几个字段需要特别理解：

\begin{itemize}
\item \texttt{train_batch_size}：全局 batch size；
\item \texttt{train_micro_batch_size_per_gpu}：每 GPU micro-batch；
\item \texttt{gradient_accumulation_steps}：梯度累积步数；
\item \texttt{bf16.enabled}：使用 BF16 混合精度；
\item \texttt{zero_optimization.stage}：ZeRO stage；
\item \texttt{reduce_bucket_size}：reduce-scatter 或 all-reduce bucket 大小；
\item \texttt{allgather_bucket_size}：all-gather bucket 大小；
\item \texttt{overlap_comm}：尝试通信计算重叠；
\item \texttt{contiguous_gradients}：将梯度放到连续 buffer，减少碎片和小通信。
\end{itemize}

DeepSpeed 配置不是越激进越好。Bucket 太大可能增加显存峰值和延迟；bucket 太小可能通信效率差。开启 overlap 后，也可能因通信和计算争抢 GPU/NIC 资源而收益不稳定。实际需要结合 profiler 调优。

\subsection{ZeRO Offload}
\label{subsec:zero-offload}

ZeRO Offload 将部分训练状态从 GPU 显存 offload 到 CPU 内存，甚至 NVMe。它的动机是：GPU HBM 很快但容量有限；CPU 内存容量大但带宽和延迟较差；NVMe 容量更大但更慢。

可以把存储层级粗略排序为：

[
\text{HBM}

>

\text{CPU DRAM}

>

\text{NVMe SSD}.
]

其中 $>$ 表示带宽更高、延迟更低。Offload 的目标是用更慢但更大的存储扩展可训练模型规模。

ZeRO Offload 常见对象包括：

\begin{itemize}
\item optimizer states；
\item gradients；
\item parameters；
\item temporary buffers。
\end{itemize}

一个简化 Stage 2 + optimizer offload 配置：

\begin{lstlisting}[language=json,caption={DeepSpeed ZeRO Offload 配置示例},label={lst:deepspeed-zero-offload-config}]
{
"bf16": {
"enabled": true
},

"zero_optimization": {
"stage": 2,
"offload_optimizer": {
"device": "cpu",
"pin_memory": true
},
"contiguous_gradients": true,
"overlap_comm": true,
"reduce_scatter": true
},

"optimizer": {
"type": "AdamW",
"params": {
"lr": 3e-4,
"betas": [0.9, 0.95],
"eps": 1e-8,
"weight_decay": 0.1
}
}
}
\end{lstlisting}

Offload 的收益是降低 HBM 占用；代价是增加 PCIe/NVLink-C2C/CPU 内存通路的数据搬运。如果 offload 数据无法与计算重叠，训练会明显变慢。因此 Offload 更适合显存极其紧张、愿意用吞吐换可行性的场景。

\begin{figure}[H]
\centering
\begin{tikzpicture}[
font=\small,
level/.style={draw, rounded corners=4pt, minimum width=3.2cm, minimum height=0.75cm, align=center},
arrow/.style={-{Latex[length=2.2mm]}, thick}
]

```
\node[font=\bfseries] at (0,3.3) {ZeRO Offload 的存储层级};

\node[level, fill=blue!8] (hbm) at (0,2.05) {GPU HBM\\fast, small};
\node[level, fill=green!10] (dram) at (0,0.85) {CPU DRAM\\larger, slower};
\node[level, fill=orange!14] (nvme) at (0,-0.35) {NVMe SSD\\very large, much slower};

\draw[arrow] (hbm) -- node[right] {offload / prefetch} (dram);
\draw[arrow] (dram) -- node[right] {spill / load} (nvme);

\node[align=left, text width=10.5cm] at (0,-1.95) {
  Offload 使用更大但更慢的存储扩展可训练模型规模。
  关键是预取、异步搬运和计算重叠，否则 PCIe/NVMe 会成为训练瓶颈。
};
```

\end{tikzpicture}
\caption{ZeRO Offload 中 GPU、CPU 与 NVMe 的层级关系}
\label{fig:zero-offload-memory-hierarchy}
\end{figure}

\subsection{ZeRO-Infinity}
\label{subsec:zero-infinity}

ZeRO-Infinity 可以理解为 ZeRO-Offload 思想的进一步扩展：不仅使用 CPU 内存，也使用 NVMe 来支持更大模型训练。它需要更精细的内存调度，包括：

\begin{itemize}
\item 分层参数管理；
\item CPU/NVMe 异步预取；
\item 计算与数据搬运重叠；
\item buffer 池管理；
\item 带宽感知调度；
\item 分片 checkpoint 与状态恢复。
\end{itemize}

从系统角度看，ZeRO-Infinity 把大模型训练从“GPU 显存管理”扩展为“分层存储系统调度”。如果数据在需要时还没有从 NVMe 加载到 CPU，再从 CPU 传到 GPU，GPU 就会空等。要让它高效，需要提前预测未来层需要哪些参数，并提前搬运。

这种思想与操作系统中的 paging 有相似之处，但难度更高，因为深度学习训练对吞吐和同步非常敏感。一次 page miss 就可能拖慢整个 step。

\subsection{配置调优项}
\label{subsec:deepspeed-tuning-knobs}

DeepSpeed ZeRO 常见调优项包括：

\begin{itemize}
\item \texttt{stage}：选择 ZeRO stage；
\item \texttt{reduce_bucket_size}：梯度 reduce bucket；
\item \texttt{allgather_bucket_size}：参数或状态 all-gather bucket；
\item \texttt{overlap_comm}：是否重叠通信与计算；
\item \texttt{contiguous_gradients}：是否使用连续梯度 buffer；
\item \texttt{offload_optimizer.device}：optimizer offload 到 CPU 或 NVMe；
\item \texttt{offload_param.device}：Stage 3 参数 offload；
\item \texttt{prefetch_bucket_size}：Stage 3 参数预取大小；
\item \texttt{param_persistence_threshold}：小参数是否持久保留；
\item \texttt{max_live_parameters}：控制同时驻留参数数量；
\item \texttt{max_reuse_distance}：控制参数复用距离。
\end{itemize}

这些参数本质上控制三个东西：

[
\text{显存峰值},
\quad
\text{通信效率},
\quad
\text{计算等待时间}.
]

例如增大 bucket size 通常提升通信带宽利用率，但也增加显存峰值和等待时间；减小 bucket size 可以更早发起通信，但小消息开销上升。参数持久化可以减少重复 all-gather，但会增加常驻显存。Offload 可以降低 HBM，但可能引入 PCIe/NVMe 瓶颈。

因此 ZeRO 调优不是简单套模板，而是通过 profiler 观察瓶颈后做取舍。

\section{ZeRO 与 3D 并行的组合}
\label{sec:zero-with-3d-parallelism}

ZeRO 经常与 TP、PP、DP 组合。理解组合方式非常关键，否则容易误以为 ZeRO 会在所有 GPU 之间切分状态。实际上，ZeRO 主要沿数据并行维度切分状态。

\subsection{ZeRO 沿 DP 维切分}
\label{subsec:zero-shards-along-dp}

在 3D 并行中，一个全局 rank 有坐标：

[
(d,t,p).
]

ZeRO 的分片通常发生在相同 $(t,p)$、不同 $d$ 的 DP group 内。也就是说，rank：

[
(0,t,p),(1,t,p),\ldots,(P_{\mathrm{DP}}-1,t,p)
]

共同切分同一个 TP shard、同一个 PP stage 的训练状态。

这是因为不同 TP rank 保存的是不同 tensor shard，不同 PP rank 保存的是不同 layers。只有 DP rank 之间才是同一模型 shard 的副本。ZeRO 消除的是这些 DP 副本之间的冗余。

[
\text{ZeRO group}
=================

\text{DP group for fixed }(t,p).
]

如果把 ZeRO 错误地跨 TP 或 PP 维度切分，就会破坏模型并行语义。

\begin{figure}[H]
\centering
\begin{tikzpicture}[
font=\small,
rank/.style={draw, circle, minimum size=0.72cm, align=center, fill=blue!8},
zline/.style={thick, green!60!black, dotted},
tpbox/.style={draw, rounded corners=4pt, dashed, thick, fill=orange!5}
]

```
\node[font=\bfseries] at (0,3.2) {ZeRO 在固定 TP/PP 坐标的 DP group 内切分};

\foreach \d/\dx in {0/0,1/0.55} {
  \foreach \p/\py in {0/0,1/0.9,2/1.8} {
    \foreach \t/\tx in {0/0,1/1.15} {
      \pgfmathsetmacro{\x}{-3.5+\tx+\dx*3.2}
      \pgfmathsetmacro{\y}{-0.4+\py}
      \node[rank] (r\d\t\p) at (\x,\y) {$d\d$};
    }
  }
}

\draw[zline] (r000) -- (r100);
\draw[zline] (r010) -- (r110);
\draw[zline] (r001) -- (r101);
\draw[zline] (r011) -- (r111);
\draw[zline] (r002) -- (r102);
\draw[zline] (r012) -- (r112);

\node[tpbox, fit=(r000)(r010), label={[font=\scriptsize]above:TP group at d0,p0}] {};
\node[tpbox, fit=(r100)(r110), label={[font=\scriptsize]above:TP group at d1,p0}] {};

\node[align=left, text width=10.5cm] at (0,-1.75) {
  绿色点线表示 ZeRO/DP 分片 group：固定 TP rank 和 PP stage，只沿 DP replica 方向切分训练状态。
  TP 和 PP 维度仍然保持各自的模型并行语义。
};
```

\end{tikzpicture}
\caption{3D 并行中 ZeRO 沿 DP 维切分状态}
\label{fig:zero-along-dp-dimension}
\end{figure}

\subsection{与 Tensor Parallel 的关系}
\label{subsec:zero-with-tensor-parallel}

TP 已经把某个权重矩阵切成多个 tensor shards。例如 Column Parallel Linear 中：

[
\mat{W}
=======

[\mat{W}*0,\mat{W}*1,\ldots,\mat{W}*{P*{\mathrm{TP}}-1}].
]

TP rank $t$ 只保存 $\mat{W}_t$。如果再使用 ZeRO，则 $\mat{W}_t$ 的训练状态会在 DP ranks 之间切分。换句话说，ZeRO 不是切完整 $\mat{W}$，而是切 TP 后的本地 shard。

因此每卡参数相关常驻状态可粗略写为：

[
M_{\mathrm{perGPU}}
\approx
\frac{M_{\mathrm{states}}}
{P_{\mathrm{TP}}P_{\mathrm{PP}}q_{\mathrm{ZeRO}}},
]

其中 $q_{\mathrm{ZeRO}}$ 取决于 ZeRO stage 和被切分的状态：

[
q_{\mathrm{ZeRO}}
=================

\begin{cases}
1, & \text{未被 ZeRO 切分的状态},\
P_{\mathrm{DP}}, & \text{被 ZeRO 切分的状态}.
\end{cases}
]

TP 和 ZeRO 都能降低单卡显存，但通信位置不同：

\begin{itemize}
\item TP 通信发生在层内，每层高频 collective；
\item ZeRO 通信发生在 DP 状态同步、梯度 reduce-scatter、参数 all-gather 等位置；
\item 两者可能同时占用 NCCL 和网络资源；
\item 若 TP 和 ZeRO group 拓扑映射不佳，会出现通信竞争。
\end{itemize}

\subsection{与 Pipeline Parallel 的关系}
\label{subsec:zero-with-pipeline-parallel}

PP 将层分配到不同 stages。每个 stage 只保存自己负责的层。ZeRO 在每个 stage 内部的 DP replicas 之间切分该 stage 的状态。

如果第 $p$ 个 stage 负责参数量为 $N_p$，ZeRO Stage 2 下该 stage 每个 rank 参数相关显存约为：

[
M_p
===

2N_p
+
\frac{14N_p}{P_{\mathrm{DP}}}.
]

若 stage 切分不均，ZeRO 不能自动解决 stage 间负载不均。它只能在同一个 stage 的 DP replicas 之间分摊状态。若 Stage 0 因 embedding 很大而显存高，Stage 0 的 DP ranks 会同时面临高显存；其他 stages 不能帮助它保存这些状态，除非改变 PP layer partition 或使用更复杂的分片策略。

因此 PP + ZeRO 配置中，仍然需要关注：

\begin{itemize}
\item 每个 stage 的参数量；
\item 每个 stage 的激活显存；
\item 首尾 stage 的 embedding/lm head；
\item pipeline schedule 造成的 live micro-batches；
\item ZeRO partition 是否在每个 stage 内均衡。
\end{itemize}

\subsection{与 FSDP 的概念对照}
\label{subsec:zero-vs-fsdp}

PyTorch FSDP，Fully Sharded Data Parallel，与 ZeRO Stage 3 在思想上非常接近：都在数据并行维度切分参数、梯度和优化器状态，并在计算需要时 all-gather 参数，用完释放。

可以粗略对应：

[
\text{FSDP full shard}
\approx
\text{ZeRO Stage 3}.
]

但二者在工程实现、API、参数 flatten、wrap 策略、prefetch、checkpoint、生态集成和调优项上有所不同。

FSDP 常通过 wrapping module 来决定分片粒度。例如一个 Transformer block 可以被 wrap 成一个 FSDP unit。Forward 进入该 block 前 all-gather 参数，退出后释放。Backward 时再次 all-gather 或复用参数，并 reduce-scatter 梯度。

ZeRO Stage 3 通常由 DeepSpeed engine 管理参数生命周期和 bucket。两者共同的核心问题都是：

\begin{itemize}
\item 参数何时 all-gather；
\item 参数何时释放；
\item 梯度如何 reduce-scatter；
\item 如何与 activation checkpointing 交互；
\item 如何保存 sharded checkpoint；
\item 如何避免频繁小通信；
\item 如何控制显存峰值。
\end{itemize}

\begin{table}[H]
\centering
\small
\caption{ZeRO 与 FSDP 的概念对照}
\label{tab:zero-fsdp-comparison}
\begin{tabular}{p{0.22\textwidth}p{0.33\textwidth}p{0.33\textwidth}}
\toprule
\textbf{维度} & \textbf{ZeRO / DeepSpeed} & \textbf{PyTorch FSDP} \
\midrule
核心思想
& 按 ZeRO stage 切分 optimizer、gradient、parameter
& 按 FSDP unit 全分片参数、梯度和优化器状态 \
\midrule
Stage 3 对应
& 参数也分片，计算前 all-gather
& Full shard 模式中参数计算前 all-gather \
\midrule
配置方式
& JSON 配置 + DeepSpeed engine
& PyTorch API + wrapping policy \
\midrule
参数粒度
& bucket 和参数管理器控制
& FSDP wrap unit 控制 \
\midrule
Checkpoint
& DeepSpeed sharded checkpoint
& full / sharded / local state dict 等模式 \
\midrule
调优重点
& bucket、offload、prefetch、overlap
& wrap 粒度、prefetch、reshard、state dict 策略 \
\bottomrule
\end{tabular}
\end{table}

\section{选择 ZeRO Stage 的工程决策}
\label{sec:choosing-zero-stage}

ZeRO Stage 的选择不是越高越好，而是根据显存、通信、模型规模和工程复杂度权衡。

\subsection{Stage 选择的基本原则}
\label{subsec:zero-stage-choice-principles}

一个实用原则是：\textbf{选择满足显存要求的最低 ZeRO stage。}

如果 DDP 能放下并且吞吐好，不一定需要 ZeRO。若 DDP 放不下或显存余量太小，可以先尝试 Stage 1。若梯度显存也成为问题，尝试 Stage 2。若参数本身也放不下，才使用 Stage 3。

[
\text{DDP}
\rightarrow
\text{ZeRO1}
\rightarrow
\text{ZeRO2}
\rightarrow
\text{ZeRO3}
]

这个顺序的原因是：stage 越高，显存越省，但通信和系统复杂度越高。

\begin{table}[H]
\centering
\small
\caption{ZeRO 各 Stage 的显存与通信取舍}
\label{tab:zero-stage-tradeoffs}
\begin{tabular}{p{0.16\textwidth}p{0.25\textwidth}p{0.25\textwidth}p{0.24\textwidth}}
\toprule
\textbf{模式} & \textbf{切分状态} & \textbf{显存收益} & \textbf{主要代价} \
\midrule
DDP
& 无，完整复制
& 最低
& 显存冗余高，但简单稳定 \
\midrule
ZeRO-1
& optimizer states
& 中等到较高
& optimizer step 后参数同步 \
\midrule
ZeRO-2
& optimizer states + gradients
& 更高
& reduce-scatter 梯度，参数仍复制 \
\midrule
ZeRO-3
& optimizer states + gradients + parameters
& 最高
& 参数 all-gather/prefetch/release，通信复杂 \
\bottomrule
\end{tabular}
\end{table}

\subsection{何时使用 Stage 1}
\label{subsec:when-to-use-zero1}

Stage 1 适合以下场景：

\begin{itemize}
\item 参数和梯度能放下，但 AdamW optimizer states 太大；
\item 希望尽量保持接近 DDP 的训练行为；
\item 通信网络不想承担 Stage 3 的频繁参数 all-gather；
\item 模型规模中等，显存压力主要来自 optimizer；
\item 需要相对简单的 checkpoint 和调试流程。
\end{itemize}

Stage 1 常常是一个性价比很高的选择。因为 AdamW 的优化器状态占比很大，切分 optimizer states 能带来明显显存收益，而不会像 Stage 3 那样改变 forward 中的参数可用性。

\subsection{何时使用 Stage 2}
\label{subsec:when-to-use-zero2}

Stage 2 适合以下场景：

\begin{itemize}
\item Stage 1 后梯度仍占用过多显存；
\item 参数本身仍能完整复制；
\item 希望避免 Stage 3 的参数 all-gather；
\item DP size 较大，梯度和 optimizer 分片收益明显；
\item 通信库对 reduce-scatter 性能较好。
\end{itemize}

Stage 2 是许多训练任务中很常用的折中点。它显存比 Stage 1 更省，但 forward/backward 不需要像 Stage 3 那样频繁聚合参数，因此吞吐可能更好。

\subsection{何时使用 Stage 3}
\label{subsec:when-to-use-zero3}

Stage 3 适合以下场景：

\begin{itemize}
\item 参数本身太大，无法在每个 DP rank 完整复制；
\item Stage 2 仍然 OOM；
\item DP size 足够大，可以有效分摊参数；
\item 网络和实现足够好，可以承受参数 all-gather；
\item 训练目标是尽可能扩展模型规模，而不是极致吞吐；
\item 可以接受更复杂的 checkpoint 和调试。
\end{itemize}

Stage 3 的调优重点是：

\begin{itemize}
\item 参数 bucket 大小；
\item prefetch 距离；
\item 参数释放策略；
\item 是否持久化小参数；
\item 与 activation checkpointing 的交互；
\item 是否 offload；
\item sharded checkpoint 是否可靠。
\end{itemize}

\subsection{常见故障与排查}
\label{subsec:zero-common-failures}

ZeRO 训练中常见问题包括：

\begin{itemize}
\item \textbf{OOM}：bucket 太大、live parameters 太多、activation 未 checkpoint、offload 未生效；
\item \textbf{吞吐很低}：参数 all-gather 暴露、offload 阻塞、bucket 太小、通信拓扑差；
\item \textbf{hang}：某些 rank 未进入 collective、参数生命周期错误、数据 loader step 数不一致；
\item \textbf{loss 异常}：梯度累积配置错误、global batch 不一致、参数更新同步错误；
\item \textbf{checkpoint 无法恢复}：shard metadata 不匹配、并行配置变化未 reshard、optimizer state 丢失；
\item \textbf{显存碎片}：大量小参数、小 bucket、频繁 gather/release 造成 allocator 压力。
\end{itemize}

排查顺序通常是：

\begin{enumerate}
\item 先用小模型、小 batch、小数据集验证逻辑；
\item 关闭复杂优化，逐步开启 ZeRO stage、offload、overlap；
\item 打印每个 rank 的并行坐标和 group；
\item 检查 global batch 与 gradient accumulation；
\item 使用 profiler 查看 all-gather/reduce-scatter 时间；
\item 检查显存峰值与 bucket 配置；
\item 验证 checkpoint 保存和恢复；
\item 扩展到多节点前先做 NCCL tests。
\end{enumerate}

\section{本章小结}
\label{sec:chapter11-summary}

本章系统介绍了 ZeRO 优化器、DeepSpeed 实践以及 ZeRO 与 3D 并行的组合方式。

\begin{itemize}
\item \textbf{AdamW 训练状态膨胀} 是大模型显存压力的重要来源。BF16 参数、BF16 梯度、FP32 master weight、Adam 一阶矩和二阶矩合计常约为每参数 16 bytes。
\item \textbf{普通 DDP 存在显存冗余}：每个 data parallel rank 都保存完整参数、完整梯度和完整优化器状态。
\item \textbf{ZeRO 的核心思想} 是在数据并行 ranks 之间切分训练状态，消除 DP 维度上的冗余。
\item \textbf{ZeRO Stage 1} 切分 optimizer states，参数和梯度仍完整复制，显存收益主要来自 Adam 状态分片。
\item \textbf{ZeRO Stage 2} 进一步切分 gradients，通常使用 reduce-scatter 聚合梯度，每个 rank 只保存自己负责的梯度 shard。
\item \textbf{ZeRO Stage 3} 切分 parameters、gradients 和 optimizer states，计算前临时 all-gather 参数，用完释放，显存最低但通信最复杂。
\item \textbf{DeepSpeed engine} 封装了 ZeRO 分片、混合精度、梯度累积、通信 bucket、overlap、offload 和 checkpoint 等复杂逻辑。
\item \textbf{ZeRO Offload} 将 optimizer states、parameters 或 gradients 移到 CPU/NVMe，以吞吐换显存容量。
\item \textbf{ZeRO-Infinity} 将训练扩展为分层存储调度问题，需要 CPU/NVMe 预取、异步搬运和计算重叠。
\item \textbf{ZeRO 与 3D 并行组合时，通常沿 DP 维切分状态}。固定 TP rank 和 PP stage 的 DP group 共同切分同一个模型 shard 的训练状态。
\item \textbf{ZeRO 与 FSDP 在概念上相近}，尤其 FSDP full shard 与 ZeRO Stage 3 都采用参数分片、按需 all-gather 和梯度 reduce-scatter。
\item \textbf{选择 ZeRO stage 的原则} 是使用满足显存要求的最低 stage。Stage 越高，显存越省，但通信和调试复杂度越高。
\end{itemize}

至此，第三篇完成了分布式训练架构的核心主线：DDP 解决数据并行，TP 解决层内矩阵切分，PP 解决模型层切分，3D 并行组合多维扩展，ZeRO/DeepSpeed 则进一步消除数据并行维度上的训练状态冗余。下一篇将进入推理架构：我们会从训练和推理的差异讲起，分析 prefill、decode、KV Cache、batching、attention kernel、调度器和 serving 系统如何共同决定 LLM 在线服务的吞吐、延迟和成本。


% ==================================================
% 第四篇：推理架构优化 Serving Infra
% ==================================================
\part{推理架构优化：让大模型高吞吐、低延迟地服务}

\chapter{KV Cache 机制}
\label{chap:kv-cache}

\section{自回归推理流程}
\label{sec:autoregressive-inference-flow}

前面几篇主要围绕训练展开：从 Transformer 结构到分布式训练，再到 3D 并行和 ZeRO。现在开始进入第四篇：\textbf{推理架构优化 Serving Infra}。训练系统关心的是如何在大量 GPU 上稳定、高吞吐地完成参数更新；推理系统关心的是如何在在线请求中以低延迟、低成本、高并发生成 token。

大语言模型推理和普通深度学习推理有本质差异。一个图像分类模型通常一次 forward 得到一个分类结果；一个推荐模型通常一次 forward 得到候选排序；而自回归语言模型必须不断生成下一个 token，并把这个 token 追加到上下文，再生成下一个 token。也就是说，LLM 推理不是一次计算，而是一串强依赖的迭代过程：

\[
x_{t+1}\sim p_{\theta}(x_{t+1}\mid x_{\le t}),
\]
\[
x_{t+2}\sim p_{\theta}(x_{t+2}\mid x_{\le t+1}),
\]
\[
\cdots
\]

对于 Decoder-only LLM，每一步生成都依赖前面已经生成的全部 token。这个依赖使得推理阶段天然分成两个部分：

\begin{itemize}
\item \textbf{Prefill 阶段}：处理用户输入 prompt，计算 prompt 中所有 token 的 hidden states，并建立 KV Cache；
\item \textbf{Decode 阶段}：每次只处理一个新 token，读取历史 KV Cache，生成下一个 token，并追加新的 K/V。
\end{itemize}

KV Cache 正是 LLM 推理性能优化的核心。没有 KV Cache，decode 每一步都要重复计算历史 token 的 Key 和 Value；有了 KV Cache，系统只需要为新 token 计算新的 K/V，同时复用历史 K/V。这个优化看似简单，却决定了 LLM serving 的显存结构、调度策略、batch 管理、长上下文能力和吞吐上限。

\subsection{prefill 阶段}
\label{subsec:prefill-stage}

Prefill 阶段处理的是用户输入 prompt。假设 prompt token 序列为：

\[
x_1,x_2,\ldots,x_S.
\]

模型需要一次性计算这些 token 的 hidden states，并得到最后一个位置的 logits，用于采样第一个输出 token：

\[
p(x_{S+1}\mid x_1,\ldots,x_S).
\]

在 prefill 中，输入序列长度为 $S$，模型可以像训练 forward 一样并行处理整个 prompt。每一层 attention 需要计算：

\[
\mat{Q}\in\mathbb{R}^{B\times n_q\times S\times d_h},
\]
\[
\mat{K}\in\mathbb{R}^{B\times n_{kv}\times S\times d_h},
\]
\[
\mat{V}\in\mathbb{R}^{B\times n_{kv}\times S\times d_h}.
\]

对于标准 causal attention，prefill 阶段每个 token 只能关注它之前的位置，因此仍需要 causal mask。但与训练一样，所有位置可以在同一次 forward 中并行计算。Attention 的 score matrix 大小约为：

\[
B\times n_q\times S\times S.
\]

因此 prefill 的计算复杂度包含 $S^2$ 项。序列越长，prefill 越昂贵。尤其当 prompt 很长时，TTFT，即 time to first token，往往主要由 prefill 决定。

Prefill 阶段完成后，系统会把每一层的 Key 和 Value 保存下来：

\[
\mat{K}^{(\ell)}*{1:S},
\quad
\mat{V}^{(\ell)}*{1:S},
\quad
\ell=0,\ldots,L-1.
\]

这些保存下来的张量就是 KV Cache。它们会在后续 decode 中被反复读取。

从系统角度看，prefill 阶段具有以下特点：

\begin{itemize}
\item \textbf{输入长度可能很长}：prompt 可以是几百、几千甚至更多 token；
\item \textbf{计算密度较高}：QKV projection、FFN 和 attention 都能形成较大 GEMM；
\item \textbf{更接近训练 forward}：一次处理完整序列；
\item \textbf{会建立 KV Cache}：为每层、每个请求保存历史 K/V；
\item \textbf{决定首 token 延迟}：用户体感中的“模型多久开始回答”主要受 prefill 影响。
\end{itemize}

\subsection{decode 阶段}
\label{subsec:decode-stage}

Decode 阶段从模型生成第一个输出 token 后开始。假设已经有上下文：

\[
x_1,x_2,\ldots,x_t.
\]

下一步只需要处理最新 token $x_t$，生成：

\[
x_{t+1}.
\]

在第 $\ell$ 层，对最新 token 计算 query、key、value：

\[
\vect{q}_t^{(\ell)},
\quad
\vect{k}_t^{(\ell)},
\quad
\vect{v}_t^{(\ell)}.
\]

然后把新的 $\vect{k}_t,\vect{v}_t$ 追加到该层的 KV Cache 中：

\[
\mat{K}_{cache}^{(\ell)}[t]
\vect{k}*t^{(\ell)},
\]
\[
\mat{V}*{cache}^{(\ell)}[t]
\vect{v}_t^{(\ell)}.
\]

Attention 输出为：

\[
\vect{o}_t^{(\ell)}
\softmax
\left(
\frac{
\vect{q}*t^{(\ell)}
\left(\mat{K}*{cache}^{(\ell)}[1:t]\right)^{\mathsf{T}}
}
{\sqrt{d_h}}
\right)
\mat{V}_{cache}^{(\ell)}[1:t].
\]

这里非常关键：decode 每一步只为当前 token 计算 query，但需要读取历史所有 Key 和 Value。因此 decode 的每一步计算量随当前上下文长度 $t$ 线性增长，而不是像 prefill 那样一次计算 $S\times S$ 的矩阵。

Decode 阶段的系统特征与 prefill 不同：

\begin{itemize}
\item \textbf{每步只处理一个新 token}；
\item \textbf{强自回归依赖}：下一步必须等待当前 token 采样完成；
\item \textbf{GEMM 尺寸小}：batch 内每个请求通常只增加一个 token；
\item \textbf{频繁读取 KV Cache}：每一步、每一层都读取历史 K/V；
\item \textbf{通常 memory-bound}：吞吐受 HBM 带宽和 cache layout 影响很大；
\item \textbf{决定持续生成速度}：用户体感中的 tokens per second 主要受 decode 影响。
\end{itemize}

这也是 LLM serving 的难点：prefill 像一个大 batch forward，decode 像大量小步迭代；两者在算子形状、显存访问、调度目标上完全不同。

\subsection{token-by-token generation}
\label{subsec:token-by-token-generation}

完整生成流程可以写成如下循环：

\[
\text{prompt}
\xrightarrow{\text{prefill}}
\text{first logits}
\xrightarrow{\text{sample}}
x_{S+1}
\xrightarrow{\text{decode}}
x_{S+2}
\xrightarrow{\text{decode}}
\cdots
\]

每次 decode 后，模型得到下一个 token 的 logits：

\[
\vect{z}_t\in\mathbb{R}^{V},
\]
其中 $V$ 是词表大小。随后经过采样策略得到 token：

\[
x_{t+1}
\operatorname{Sample}
(
\vect{z}_t;
T,\ top\text{-}k,\ top\text{-}p,\ \ldots
).
\]

生成循环通常会在以下条件之一满足时结束：

\begin{itemize}
\item 生成了 EOS token；
\item 命中 stop words 或 stop sequence；
\item 达到最大生成长度；
\item 请求被用户取消；
\item server 触发超时或资源限制；
\item 安全策略或工具调用逻辑中断生成。
\end{itemize}

下面给出一个极简自回归推理伪代码。它强调 prefill、decode 和 KV Cache 的关系，不代表高性能实现。

\begin{lstlisting}[language=Python,caption={自回归推理中 prefill、decode 与 KV Cache 的极简伪代码},label={lst:autoregressive-generation-kv-cache}]
def generate(model, input_ids, max_new_tokens, sampler):
"""
input_ids: [batch, prompt_len]
"""
# Prefill: process the whole prompt and build KV cache.
logits, kv_cache = model.forward_prefill(input_ids)

next_token = sampler(logits[:, -1, :])
output_tokens = [next_token]

# Decode: one token per iteration.
for _ in range(max_new_tokens - 1):
    logits, kv_cache = model.forward_decode(
        input_ids=next_token[:, None],
        kv_cache=kv_cache,
    )

    next_token = sampler(logits[:, -1, :])
    output_tokens.append(next_token)

    if should_stop(next_token):
        break

return concat_tokens(output_tokens)

\end{lstlisting}

真实 serving 系统会比这复杂得多。它需要把多个请求放到同一个 batch 中，支持请求动态加入和退出，管理不同长度的 KV Cache，处理 prefix cache、分页分配、并行采样、流式输出、超时、取消、优先级、SLA 和多 GPU 通信。

\subsection{为什么推理不同于训练}
\label{subsec:why-inference-differs-from-training}

训练与推理都执行 Transformer forward，但系统形态差异很大。训练通常有固定 batch、固定 sequence length、完整 forward+backward、规则的张量形状和长时间稳定运行。推理则面向在线请求，输入长度和输出长度高度动态，且 decode 阶段存在严格 token 依赖。

训练中，一次 step 可能处理：

\[
B\times S
\]
个 token，并对所有 token 计算 loss。推理 prefill 处理 prompt，而 decode 每步通常只处理每个请求的一个 token。

训练主要瓶颈包括：

\begin{itemize}
\item 大 GEMM 计算吞吐；
\item 激活显存；
\item backward 与 optimizer；
\item DP/TP/PP 通信；
\item checkpoint 和数据加载。
\end{itemize}

推理主要瓶颈包括：

\begin{itemize}
\item KV Cache 显存；
\item decode 阶段 HBM 带宽；
\item 请求调度；
\item batch 内长度不均；
\item TTFT 与 TPOT；
\item cache 碎片；
\item 多租户和 SLA。
\end{itemize}

\begin{table}[H]
\centering
\small
\caption{LLM 训练与推理的系统差异}
\label{tab:training-vs-inference}
\begin{tabular}{p{0.22\textwidth}p{0.34\textwidth}p{0.34\textwidth}}
\toprule
\textbf{维度} & \textbf{训练} & \textbf{推理} \
\midrule
输入形态
& 大 batch、固定或近似固定序列长度
& 请求长度动态，prompt 和 output 长度差异大 \
\midrule
计算流程
& forward + backward + optimizer step
& prefill + 多轮 decode \
\midrule
并行性
& 一个 step 内大量 token 可并行
& decode token-by-token，强依赖上一 token \
\midrule
主要显存
& 参数、梯度、优化器状态、激活
& 参数、KV Cache、临时 attention buffer \
\midrule
主要瓶颈
& GEMM、激活、分布式通信
& HBM 带宽、KV Cache 管理、调度 \
\midrule
关键指标
& tokens/s、MFU、loss、收敛稳定性
& TTFT、TPOT、吞吐、尾延迟、成本/token \
\bottomrule
\end{tabular}
\end{table}

\begin{figure}[H]
\centering
\begin{tikzpicture}[
font=\small,
box/.style={draw, rounded corners=3pt, minimum width=1.8cm, minimum height=0.7cm, align=center},
token/.style={draw, rounded corners=2pt, minimum width=0.55cm, minimum height=0.4cm, align=center},
arrow/.style={-{Latex[length=2.2mm]}, thick}
]

\node[font=\bfseries] at (0,3.4) {Prefill 与 Decode 阶段对比};

% Prefill
\node[font=\bfseries] at (-3.8,2.55) {Prefill};
\foreach \i/\x in {1/-5.2,2/-4.55,3/-3.9,4/-3.25,5/-2.6} {
  \node[token, fill=blue!10] (p\i) at (\x,1.8) {$x_{\i}$};
}
\node[box, fill=orange!14] (pf) at (-3.9,0.75) {Full prompt\\forward};
\node[box, fill=green!10] (pkv) at (-3.9,-0.25) {Build KV Cache};
\node[box, fill=purple!10] (plogit) at (-3.9,-1.25) {First logits};

\foreach \i in {1,2,3,4,5} {
  \draw[arrow] (p\i) -- (pf);
}
\draw[arrow] (pf) -- (pkv);
\draw[arrow] (pkv) -- (plogit);

% Decode
\node[font=\bfseries] at (3.8,2.55) {Decode};
\node[token, fill=blue!10] (dnew) at (2.25,1.8) {$x_t$};
\node[box, fill=orange!14] (df) at (3.8,1.25) {One-token\\forward};
\node[box, fill=green!10] (dkv) at (5.35,1.8) {Read old KV\\Append new KV};
\node[box, fill=purple!10] (dlogit) at (3.8,0.25) {Next logits};
\node[token, fill=yellow!18] (sample) at (3.8,-0.75) {$x_{t+1}$};
\node[box, fill=red!8] (loop) at (3.8,-1.6) {repeat};

\draw[arrow] (dnew) -- (df);
\draw[arrow] (dkv) -- (df);
\draw[arrow] (df) -- (dlogit);
\draw[arrow] (dlogit) -- (sample);
\draw[arrow] (sample) -- (loop);
\draw[arrow] (loop.west) .. controls (1.7,-1.1) and (1.7,1.6) .. (dnew.west);

\draw[dashed, thick] (0,-2.0) -- (0,2.8);

\node[align=left, text width=10.6cm] at (0,-2.75) {
  Prefill 一次处理完整 prompt，计算密度高；Decode 每次只处理一个新 token，
  但每层都要读取不断增长的 KV Cache，因此更容易受显存带宽和 cache 管理影响。
};

\end{tikzpicture}
\caption{Prefill 与 Decode 阶段的计算流程对比}
\label{fig:prefill-vs-decode}
\end{figure}

\section{KV Cache 的基本原理}
\label{sec:kv-cache-basics}

KV Cache 的核心问题是：为什么 Key 和 Value 可以缓存？为什么 Query 不缓存？缓存的 shape 是什么？显存如何估算？这些问题决定了后续所有推理优化。

\subsection{Key / Value 为什么可以缓存}
\label{subsec:why-cache-key-value}

考虑 Decoder-only Transformer 的第 $\ell$ 层。对某个 token 的 hidden state $\vect{h}_i^{(\ell)}$，线性投影得到：

\[
\vect{q}_i^{(\ell)}=\vect{h}_i^{(\ell)}\mat{W}_Q^{(\ell)},
\]
\[
\vect{k}_i^{(\ell)}=\vect{h}_i^{(\ell)}\mat{W}_K^{(\ell)},
\]
\[
\vect{v}_i^{(\ell)}=\vect{h}_i^{(\ell)}\mat{W}_V^{(\ell)}.
\]

在自回归 decode 中，当已经生成了 token $x_1,\ldots,x_t$ 后，未来生成 $x_{t+1}$ 不会改变过去 token 的 hidden states。原因是 causal mask 保证过去位置不能看未来位置；模型参数在推理中也不更新。因此，对于历史位置 $i\le t$：

\[
\vect{k}_i^{(\ell)},\quad
\vect{v}_i^{(\ell)}
\]

在后续 decode 步中保持不变。

当生成第 $t+1$ 个 token 时，需要计算的是新 token 的 query：

\[
\vect{q}_{t+1}^{(\ell)}.
\]

它要和历史所有 key 做点积：

\[
\vect{q}*{t+1}^{(\ell)}
\left[
\vect{k}*{1}^{(\ell)},
\vect{k}*{2}^{(\ell)},
\ldots,
\vect{k}*{t+1}^{(\ell)}
\right]^{\mathsf{T}}.
\]

历史 key/value 不变，因此可以缓存。若不缓存，就必须在每个 decode step 重新对整个上下文做 forward，重复计算所有历史 token 的 K/V。这会把 decode 成本从每步近似 $O(t)$ 变成重复的 $O(t^2)$ 或更高，完全不可接受。

为什么不缓存 Query？因为每一步只需要当前新 token 的 query。历史 token 的 query 不会再用于未来 token 的 attention。未来 token 关注历史 token 时，需要的是未来 token 的 query 与历史 token 的 key，而不是历史 token 的 query。因此 Query 没有长期缓存价值。

可以总结为：

\[
\text{Key: 历史 token 的可匹配地址，需要被未来 query 查询。}
\]
\[
\text{Value: 历史 token 的内容，需要被未来 attention 加权汇聚。}
\]
\[
\text{Query: 当前 token 的查询向量，只在当前步使用。}
\]

\subsection{cache shape 推导}
\label{subsec:kv-cache-shape-derivation}

设模型有：

\begin{itemize}
\item 层数 $L$；
\item batch size $B$；
\item 当前最大上下文长度 $S$；
\item query head 数 $n_q$；
\item key/value head 数 $n_{kv}$；
\item head dimension $d_h$；
\item hidden size $H=n_qd_h$；
\item 每个元素字节数 $s_{\mathrm{dtype}}$。
\end{itemize}

对于每一层，K Cache 的 shape 通常是：

\[
\mat{K}*{cache}^{(\ell)}
\in
\mathbb{R}^{B\times n*{kv}\times S\times d_h}.
\]

V Cache 同理：

\[
\mat{V}*{cache}^{(\ell)}
\in
\mathbb{R}^{B\times n*{kv}\times S\times d_h}.
\]

所有层合并后：

\[
\mat{K}*{cache}
\in
\mathbb{R}^{L\times B\times n*{kv}\times S\times d_h},
\]
\[
\mat{V}*{cache}
\in
\mathbb{R}^{L\times B\times n*{kv}\times S\times d_h}.
\]

有些实现会使用不同 layout，例如：

\[
[B,L,n_{kv},S,d_h],
\]
\[
[L,B,S,n_{kv},d_h],
\]
\[
[L,B,n_{kv},d_h,S].
\]

layout 的选择会影响 attention kernel 的访存连续性、append 新 token 的写入效率、batch 内请求重排成本以及分页 KV Cache 的实现方式。

一个简单连续 KV Cache layout 如下图所示。

\begin{figure}[H]
\centering
\begin{tikzpicture}[
font=\small,
cell/.style={draw, minimum width=0.45cm, minimum height=0.38cm},
label/.style={font=\scriptsize},
box/.style={draw, rounded corners=3pt, minimum width=2.0cm, minimum height=0.65cm, align=center},
arrow/.style={-{Latex[length=2.1mm]}, thick}
]

\node[font=\bfseries] at (0,3.25) {KV Cache 张量布局示意};

\node[box, fill=blue!6] (layer) at (-4.7,1.9) {Layer $\ell$};
\node[box, fill=orange!12] (kv) at (-2.4,1.9) {K Cache / V Cache};
\draw[arrow] (layer) -- (kv);

\foreach \b/\y in {0/1.15,1/0.35,2/-0.45} {
  \node[label] at (-5.1,\y) {batch \b};
  \foreach \h/\xbase in {0/-3.95,1/-2.85,2/-1.75} {
    \node[label] at (\xbase,1.65) {head \h};
    \foreach \s in {0,...,4} {
      \node[cell, fill=green!10] at ({\xbase+0.18*\s},{\y}) {};
    }
  }
}

\draw[decorate,decoration={brace,amplitude=5pt},thick]
  (-4.05,-0.95) -- (-3.15,-0.95)
  node[midway,below=6pt] {$S$ positions};

\draw[decorate,decoration={brace,amplitude=5pt},thick]
  (-4.3,1.35) -- (-1.25,1.35)
  node[midway,above=6pt] {$n_{kv}$ heads};

\node[box, fill=purple!10, minimum width=2.4cm] (shape) at (2.6,0.4)
  {$[L,B,n_{kv},S,d_h]$};

\draw[arrow] (-1.1,0.4) -- (shape);

\node[align=left, text width=10.5cm] at (0,-2.05) {
  KV Cache 需要为每一层保存 K 和 V。每个请求的 cache 会随着生成长度增长。
  实际系统会根据 attention kernel 的访存模式选择具体内存 layout。
};

\end{tikzpicture}
\caption{KV Cache 的基本张量布局}
\label{fig:kv-cache-tensor-layout}
\end{figure}

\subsection{batch、head、seq\_len、head\_dim}
\label{subsec:batch-head-seqlen-headdim}

KV Cache shape 中每个维度都有系统含义。

\textbf{Batch 维 $B$} 对应同一时刻正在服务的请求数。与训练不同，推理中的 batch 是动态的。请求可能随时加入或退出，且每个请求的上下文长度不同。因此 KV Cache 的 batch 维并不总是一个固定矩形张量那样简单。传统实现可能为每个 batch slot 预留最大长度；更高效的实现会使用 block-based 或 paged allocation。

\textbf{Head 维 $n_{kv}$} 对应 key/value heads 数量。对于 MHA：

\[
n_{kv}=n_q.
\]

对于 MQA：

\[
n_{kv}=1.
\]

对于 GQA：

\[
1<n_{kv}<n_q.
\]

KV Cache 显存与 $n_{kv}$ 成正比。因此 GQA/MQA 对推理显存和带宽非常重要。

\textbf{Sequence 维 $S$} 对应每个请求已经缓存的 token 数。Prefill 后，$S$ 等于 prompt length；decode 每生成一个 token，$S$ 增加 1。对于多轮对话，请求的 KV Cache 可能不断增长，直到达到上下文窗口上限或触发截断、压缩、eviction。

\textbf{Head dimension $d_h$} 是每个 attention head 的向量维度。通常：

\[
H=n_qd_h.
\]

KV Cache 显存与 $d_h$ 成正比。对固定 hidden size $H$，如果使用 GQA 减少 $n_{kv}$，KV Cache 可显著下降。

\subsection{KV Cache 显存估算}
\label{subsec:kv-cache-memory-estimation}

KV Cache 显存是 serving capacity 的核心约束。对于一个 dense Decoder-only 模型，KV Cache 总显存可粗略估算为：

\[
M_{\mathrm{KV}}
2
\cdot
L
\cdot
B
\cdot
S
\cdot
n_{kv}
\cdot
d_h
\cdot
s_{\mathrm{dtype}}.
\]

其中前面的 $2$ 表示 K 和 V 两份 cache。

如果使用 tensor parallel，KV Cache 也可能按 head 或 hidden 维切分到多个 GPU 上。设 TP size 为 $P_{\mathrm{TP}}$，如果 $n_{kv}$ 能按 TP 维切分，则每 GPU KV Cache 近似为：

\[
M_{\mathrm{KV,perGPU}}
\frac{
2
\cdot
L
\cdot
B
\cdot
S
\cdot
n_{kv}
\cdot
d_h
\cdot
s_{\mathrm{dtype}}
}
{P_{\mathrm{TP}}}.
\]

如果使用 MHA，$n_{kv}=n_q$；如果使用 GQA，$n_{kv}$ 更小。以 $L=80$、$B=32$、$S=4096$、$n_{kv}=8$、$d_h=128$、BF16 为例：

\[
M_{\mathrm{KV}}
2
\cdot
80
\cdot
32
\cdot
4096
\cdot
8
\cdot
128
\cdot
2
\ \mathrm{bytes}.
\]

这个值约为：

\[
42.95\ \mathrm{GB}.
\]

这只是 KV Cache，不包括模型参数、临时 buffer、CUDA runtime、scheduler metadata 和其他开销。若 batch 更大、上下文更长，KV Cache 会迅速成为主要显存占用。

下面给出一个 KV Cache 显存估算器。

\begin{lstlisting}[language=Python,caption={KV Cache 显存估算器},label={lst:kv-cache-memory-estimator}]
def estimate_kv_cache_gb(
num_layers,
batch_size,
seq_len,
num_kv_heads,
head_dim,
dtype_bytes=2,
tensor_parallel_size=1,
):
"""
Estimate per-GPU KV cache memory in GB.

KV cache shape:
  K: [L, B, n_kv, S, d_h]
  V: [L, B, n_kv, S, d_h]
"""
total_bytes = (
    2
    * num_layers
    * batch_size
    * seq_len
    * num_kv_heads
    * head_dim
    * dtype_bytes
)

per_gpu_bytes = total_bytes / tensor_parallel_size
return per_gpu_bytes / (1024 ** 3)

gb = estimate_kv_cache_gb(
num_layers=80,
batch_size=32,
seq_len=4096,
num_kv_heads=8,
head_dim=128,
dtype_bytes=2,
tensor_parallel_size=1,
)

print(round(gb, 2), "GB")
\end{lstlisting}

KV Cache 显存估算有几个常见误区：

\begin{itemize}
\item \textbf{只看模型参数，不看 KV Cache}：对于长上下文和高并发，KV Cache 可能比参数更限制 batch size；
\item \textbf{忽略 GQA/MQA}：$n_{kv}$ 直接影响 KV Cache；
\item \textbf{忽略 TP 切分}：多卡推理时每 GPU cache 可能被切分；
\item \textbf{忽略碎片}：连续分配可能造成大量未使用空间；
\item \textbf{忽略最大长度预留}：为最大上下文预留会比实际 token 使用更多显存；
\item \textbf{忽略多轮对话增长}：同一个会话可能随着轮次持续增长。
\end{itemize}

\section{KV Cache 的性能瓶颈}
\label{sec:kv-cache-performance-bottlenecks}

KV Cache 不只是显存问题，也是性能问题。Decode 阶段每生成一个 token，都要在每一层读取历史 K/V。随着上下文变长，读取的数据量线性增长。由于每一步只处理一个新 token，计算密度不高，decode 往往受 HBM 带宽、cache layout、batch 组织和 kernel 访存效率限制。

\subsection{decode 阶段的 memory-bound 特性}
\label{subsec:decode-memory-bound}

在 decode 阶段，每层 attention 对当前 token 执行：

\[
\vect{q}*t\mat{K}*{cache}^{\mathsf{T}},
\]
\[
\softmax(\cdot)\mat{V}_{cache}.
\]

对于每个新 token、每层、每个 KV head，需要读取长度为 $S$ 的 K 和 V：

\[
\mat{K}*{cache}\in\mathbb{R}^{S\times d_h},
\quad
\mat{V}*{cache}\in\mathbb{R}^{S\times d_h}.
\]

读取量随 $S$ 线性增长：

\[
O(Sd_h).
\]

而当前 query 只有一个 token，矩阵乘法形状通常较小。相对于训练或 prefill 中的大 GEMM，decode attention 的算术强度较低：

\[
\text{Arithmetic Intensity}
\frac{\text{FLOPs}}
{\text{Bytes read/written}}.
\]

当算术强度低时，性能更受内存带宽限制，而不是 Tensor Core 峰值 FLOPs 限制。也就是说，即使 GPU 有很高的理论算力，decode 阶段也可能因为不断读取 KV Cache 而无法充分利用计算单元。

\begin{figure}[H]
\centering
\begin{tikzpicture}[
font=\small,
box/.style={draw, rounded corners=3pt, minimum width=1.65cm, minimum height=0.65cm, align=center},
mem/.style={draw, rounded corners=4pt, minimum width=2.2cm, minimum height=1.0cm, align=center, fill=orange!14},
arrow/.style={-{Latex[length=2.1mm]}, thick}
]

\node[font=\bfseries] at (0,3.25) {Decode 阶段为何容易 memory-bound};

\node[box, fill=blue!6] (q) at (-4.8,1.2) {new token\\$q_t$};
\node[mem] (kvcache) at (-1.8,1.2) {KV Cache\\history $1:t$};
\node[box, fill=green!10] (attn) at (1.3,1.2) {Attention\\small compute};
\node[box, fill=purple!10] (out) at (4.2,1.2) {next hidden};

\draw[arrow] (q) -- (attn);
\draw[arrow] (kvcache) -- node[above] {large HBM reads} (attn);
\draw[arrow] (attn) -- (out);

\node[box, fill=yellow!16] (ffn) at (1.3,-0.1) {FFN / Linear\\small batch GEMM};
\draw[arrow] (out) -- (ffn);

\node[align=left, text width=10.5cm] at (0,-1.65) {
  Decode 每步只处理新 token，计算规模小；但必须读取历史 K/V。
  上下文越长，KV Cache 读取越大，attention 越容易受 HBM 带宽限制。
};

\end{tikzpicture}
\caption{Decode 阶段的 memory-bound 特性}
\label{fig:decode-memory-bound}
\end{figure}

\subsection{长上下文显存压力}
\label{subsec:long-context-memory-pressure}

KV Cache 显存与上下文长度 $S$ 线性相关：

\[
M_{\mathrm{KV}}\propto S.
\]

如果上下文长度从 $4096$ 增加到 $32768$，KV Cache 显存也增加 $8$ 倍。长上下文模型不仅 prefill 更贵，decode 时每步读取的 cache 也更大。

设单请求 KV Cache 显存为：

\[
M_{\mathrm{KV,req}}
2L S n_{kv}d_hs_{\mathrm{dtype}}.
\]

在固定 GPU 显存预算 $M_{\mathrm{budget}}$ 下，可并发请求数近似满足：

\[
B_{\max}
\le
\frac{M_{\mathrm{budget}}}
{M_{\mathrm{KV,req}}}.
\]

这说明长上下文会直接降低可承载并发。对于在线服务，用户请求长度差异很大：有些请求只有几十 token，有些请求有几万 token。如果简单按最大长度预留 KV Cache，会浪费大量显存；如果按实际长度动态分配，又需要复杂内存管理。

长上下文还会影响调度。一个很长的请求可能占用大量 KV Cache，降低 batch 中其他请求的可用空间。服务系统必须决定是否接收、排队、截断、分块处理或拒绝该请求。这属于 admission control 的问题，后续 continuous batching 章节会继续讨论。

\subsection{batch size 与 cache 碎片}
\label{subsec:batch-size-cache-fragmentation}

推理中 batch size 不像训练那样固定。请求动态加入和退出，每个请求的 prompt length 和 output length 都不同。若使用简单连续分配，为每个请求预留最大长度，会产生两类浪费：

\begin{itemize}
\item \textbf{内部碎片}：请求实际使用长度小于预留长度，预留空间未使用；
\item \textbf{外部碎片}：已释放的空间分散，难以满足新的大请求连续分配。
\end{itemize}

例如两个请求：

\[
R_1:\ \text{最大预留 }4096,\quad \text{实际使用 }512,
\]
\[
R_2:\ \text{最大预留 }4096,\quad \text{实际使用 }3800.
\]

如果都按 $4096$ 预留，$R_1$ 浪费大量空间。如果按实际长度增长式分配，则当 $R_2$ 继续生成时，可能需要扩展连续空间；若后面空间已被其他请求占用，就需要搬移或重新分配。

这正是 PagedAttention 和 block-based allocation 的动机。把 KV Cache 切成固定大小 block，可以像操作系统分页一样管理，避免要求每个请求拥有连续大块空间。本章先讨论基本策略，下一章会深入 PagedAttention。

\subsection{多轮对话中的 cache 复用}
\label{subsec:multi-turn-cache-reuse}

多轮对话中，用户和助手的历史消息会不断追加。例如：

\[
\text{System prompt}
+
\text{User}_1
+
\text{Assistant}_1
+
\text{User}_2
+
\cdots
\]

如果每一轮都从头 prefill 整个对话历史，会重复计算大量相同前缀。KV Cache 可以在同一会话内复用历史前缀，只对新增用户消息做增量 prefill，再进入 decode。

假设前一轮已经缓存了前缀长度 $S_{\mathrm{old}}$：

\[
x_1,\ldots,x_{S_{\mathrm{old}}}.
\]

新一轮用户追加了 $\Delta S$ 个 token：

\[
x_{S_{\mathrm{old}}+1},\ldots,x_{S_{\mathrm{old}}+\Delta S}.
\]

系统只需要对新增部分做 prefill，并让新增 token attend 到旧 cache：

\[
\mat{K}*{cache}[1:S*{\mathrm{old}}+\Delta S].
\]

这样可以降低多轮对话中的重复 prefill 成本。但 cache 复用也有工程问题：

\begin{itemize}
\item 会话 cache 需要在请求之间保留，占用显存；
\item 用户可能长时间不继续对话，cache 是否保留需要 eviction 策略；
\item 多副本部署时，后续请求是否路由到同一 worker 影响复用；
\item 如果 system prompt 或采样参数变化，prefix 是否仍可复用需要判断；
\item 对话历史超出上下文窗口时，需要截断或压缩。
\end{itemize}

\begin{figure}[H]
\centering
\begin{tikzpicture}[
font=\small,
req/.style={draw, rounded corners=3pt, minimum width=1.8cm, minimum height=0.55cm, align=center},
cache/.style={draw, rounded corners=2pt, minimum width=0.65cm, minimum height=0.4cm, align=center},
arrow/.style={-{Latex[length=2.1mm]}, thick}
]

\node[font=\bfseries] at (0,3.25) {多请求 KV Cache 随生成过程增长};

\foreach \r/\y/\len/\col in {1/1.65/4/blue!10,2/0.75/7/green!10,3/-0.15/3/orange!14} {
  \node[req, fill=\col] at (-5.0,\y) {Request \r};
  \foreach \i in {1,...,\len} {
    \node[cache, fill=\col] at ({-3.5+0.45*\i},\y) {};
  }
  \draw[arrow] ({-3.5+0.45*\len+0.35},\y) -- ({-3.5+0.45*\len+1.0},\y);
  \node[cache, fill=yellow!18] at ({-3.5+0.45*\len+1.3},\y) {new};
}

\node[align=left, text width=10.5cm] at (0,-1.45) {
  不同请求的 KV Cache 长度不同，并且 decode 每步继续增长。
  Serving 系统需要动态管理 cache 空间，避免长短请求混部造成显存碎片和调度低效。
};

\end{tikzpicture}
\caption{多请求 KV Cache 的增长过程}
\label{fig:multi-request-kv-cache-growth}
\end{figure}

\section{KV Cache 管理策略}
\label{sec:kv-cache-management-strategies}

KV Cache 管理是 LLM serving 系统的核心模块之一。它不仅负责分配和释放显存，还要配合 scheduler、attention kernel、prefix cache、请求取消、长上下文、batch 重排和多 GPU 并行。

本节讨论四类基础策略：

\begin{itemize}
\item contiguous allocation；
\item block-based allocation；
\item cache eviction；
\item prefix caching。
\end{itemize}

\subsection{contiguous allocation}
\label{subsec:contiguous-allocation}

最简单的 KV Cache 管理方式是连续分配。系统为每个请求分配一段连续的 cache 空间：

\[
[\text{start},\ \text{start}+S_{\max}).
\]

每生成一个 token，就把新的 K/V 写入下一个位置。连续分配的优点是实现简单，attention kernel 读取也简单，因为一个请求的历史 K/V 在内存中是连续的。

连续分配的常见布局可以是：

\[
\texttt{cache[layer][batch_slot][kv_head][pos][head_dim]}.
\]

写入第 $t$ 个 token 时：

\[
\texttt{cache[layer][slot][head][t][:]}
\leftarrow
\texttt{new_kv}.
\]

这种方式适合固定 batch、固定最大长度、请求长度较一致的场景。但在线服务中，连续分配有明显缺点：

\begin{itemize}
\item 需要提前预留最大长度，浪费显存；
\item 请求长度不均导致内部碎片；
\item 请求释放后产生外部碎片；
\item 扩展请求长度可能需要重新分配；
\item batch slot 与 request 绑定，动态调度不灵活。
\end{itemize}

\begin{figure}[H]
\centering
\begin{tikzpicture}[
font=\small,
cell/.style={draw, minimum width=0.48cm, minimum height=0.42cm},
label/.style={font=\scriptsize}
]

\node[font=\bfseries] at (0,3.0) {Contiguous KV Cache Allocation};

\foreach \r/\y/\used/\col in {1/1.55/3/blue!10,2/0.85/6/green!10,3/0.15/2/orange!14} {
  \node[label] at (-4.8,\y) {Request \r};
  \foreach \i in {0,...,7} {
    \ifnum\i<\used
      \node[cell, fill=\col] at ({-3.8+0.52*\i},\y) {};
    \else
      \node[cell, fill=gray!16] at ({-3.8+0.52*\i},\y) {};
    \fi
  }
}

\node[draw, rounded corners=2pt, fill=blue!10, minimum width=0.55cm, minimum height=0.35cm] at (1.4,1.5) {};
\node[anchor=west] at (1.8,1.5) {used tokens};
\node[draw, rounded corners=2pt, fill=gray!16, minimum width=0.55cm, minimum height=0.35cm] at (1.4,0.9) {};
\node[anchor=west] at (1.8,0.9) {reserved but unused};

\node[align=left, text width=10.4cm] at (0,-1.35) {
  连续分配简单、访存友好，但需要预留最大长度。
  在线请求长度不均时，未使用的预留空间会造成显著显存浪费。
};

\end{tikzpicture}
\caption{连续 KV Cache 分配中的预留空间浪费}
\label{fig:contiguous-kv-cache-allocation}
\end{figure}

\subsection{block-based allocation}
\label{subsec:block-based-allocation}

Block-based allocation 将 KV Cache 切成固定大小的 blocks。每个 block 可以保存固定数量 token 的 K/V，例如 block size 为 $16$ 或 $32$ tokens。请求不再需要一段连续大空间，而是持有一个 block 列表：

\[
\text{Request }r
\rightarrow
[b_0,b_1,b_2,\ldots].
\]

当请求继续生成，需要更多空间时，从 free block pool 中申请新 block；请求结束时，把它持有的 blocks 释放回池中。

这种方式的优点是：

\begin{itemize}
\item 避免按最大长度连续预留；
\item 释放后 block 可被其他请求复用；
\item 更适合动态 batch 和长短请求混部；
\item 支持 prefix cache 的 block 级共享；
\item 为 PagedAttention 的 block table 奠定基础。
\end{itemize}

缺点是 attention kernel 读取 K/V 时不能简单假设请求的历史 token 连续存储。Kernel 需要通过 block table 把逻辑位置映射到物理 block。这会增加地址计算复杂度，但换来显存利用率提升。

\begin{figure}[H]
\centering
\begin{tikzpicture}[
font=\small,
block/.style={draw, rounded corners=2pt, minimum width=0.7cm, minimum height=0.45cm, align=center},
req/.style={draw, rounded corners=3pt, minimum width=1.55cm, minimum height=0.55cm, align=center, fill=blue!6},
arrow/.style={-{Latex[length=2.1mm]}, thick}
]

\node[font=\bfseries] at (0,3.15) {Block-based KV Cache Allocation};

\node[req] (r1) at (-5.2,1.55) {Req 1};
\node[req] (r2) at (-5.2,0.55) {Req 2};
\node[req] (r3) at (-5.2,-0.45) {Req 3};

\foreach \i/\x/\y/\col in {
  0/-2.8/1.55/blue!10,
  1/-1.85/1.55/blue!10,
  2/-0.9/1.55/blue!10,
  3/-2.8/0.55/green!10,
  4/-1.85/0.55/green!10,
  5/-0.9/0.55/green!10,
  6/0.05/0.55/green!10,
  7/-2.8/-0.45/orange!14,
  8/-1.85/-0.45/orange!14
} {
  \node[block, fill=\col] (b\i) at (\x,\y) {$b_{\i}$};
}

\draw[arrow] (r1) -- (b0);
\draw[arrow] (r2) -- (b3);
\draw[arrow] (r3) -- (b7);

\node[draw, rounded corners=3pt, fill=gray!12, minimum width=2.9cm, minimum height=0.75cm, align=center] (free) at (3.1,0.55)
  {Free block pool};

\foreach \i/\x in {9/2.25,10/3.1,11/3.95} {
  \node[block, fill=gray!20] at (\x,-0.25) {$b_{\i}$};
}

\node[align=left, text width=10.5cm] at (0,-1.95) {
  每个请求持有若干固定大小 block，而不是一段连续大空间。
  新 token 写满当前 block 后，再从 free pool 申请新 block。
};

\end{tikzpicture}
\caption{Block-based KV Cache 分配}
\label{fig:block-based-kv-cache-allocation}
\end{figure}

\subsection{cache eviction}
\label{subsec:cache-eviction}

KV Cache 占用显存，而显存是有限资源。当请求结束、取消、超时或被调度策略移除时，需要释放 cache。更复杂的是，系统可能为了 prefix caching 或多轮会话保留某些 cache，以便未来复用。此时需要 eviction 策略决定哪些 cache 可以被驱逐。

常见 eviction 依据包括：

\begin{itemize}
\item \textbf{请求是否已结束}：结束请求的 cache 可立即释放；
\item \textbf{会话是否仍活跃}：长时间无后续请求的 session cache 可释放；
\item \textbf{prefix 复用频率}：高频 system prompt 或热门 prefix 值得保留；
\item \textbf{cache 大小}：超大 cache 可能优先驱逐以释放显存；
\item \textbf{SLA 优先级}：低优先级请求的 cache 可更早释放；
\item \textbf{最近使用时间}：LRU 类策略；
\item \textbf{预计重算成本}：prefill 很贵的 prefix 更值得缓存。
\end{itemize}

可以定义一个 cache value 分数：

\[
\operatorname{score}(c)
\alpha\cdot \operatorname{reuse_prob}(c)
+
\beta\cdot \operatorname{recompute_cost}(c)
-------------------------------------------
## \gamma\cdot \operatorname{memory_size}(c)
\delta\cdot \operatorname{idle_time}(c).
\]

分数低的 cache 更适合驱逐。真实系统不一定显式使用这样的公式，但背后的权衡类似：保留 cache 可以降低未来 prefill 成本，但会占用当前显存，降低当前可服务并发。

\subsection{prefix caching}
\label{subsec:prefix-caching}

Prefix caching 指多个请求共享相同前缀时，复用该前缀的 KV Cache。典型场景包括：

\begin{itemize}
\item 所有请求共享同一个 system prompt；
\item RAG 模板中前半部分相同；
\item 多个用户请求使用相同 few-shot examples；
\item 多轮对话中复用历史上下文；
\item 批量评测中 prompt 模板一致。
\end{itemize}

假设两个请求前缀相同：

\[
x_1,\ldots,x_K.
\]

它们的 KV Cache 前 $K$ 个 token 完全相同。因此可以只保存一份 prefix KV，并让多个请求引用它。每个请求只为自己的 suffix 分配新的 cache：

\[
\text{Request A}
\text{Shared Prefix}
+
\text{Suffix A},
\]
\[
\text{Request B}
\text{Shared Prefix}
+
\text{Suffix B}.
\]

Prefix cache 的收益是减少重复 prefill 和重复显存占用。但它带来引用计数和生命周期管理问题。共享 prefix 不能在仍有请求引用时释放；如果某个请求继续生成，只能追加自己的 suffix，不能修改共享 prefix。

\begin{figure}[H]
\centering
\begin{tikzpicture}[
font=\small,
block/.style={draw, rounded corners=2pt, minimum width=0.6cm, minimum height=0.42cm, align=center},
req/.style={draw, rounded corners=3pt, minimum width=1.45cm, minimum height=0.55cm, align=center, fill=blue!6},
arrow/.style={-{Latex[length=2.1mm]}, thick}
]

\node[font=\bfseries] at (0,3.25) {Prefix Cache 复用示意};

\node[req] (ra) at (-5.2,1.25) {Req A};
\node[req] (rb) at (-5.2,0.15) {Req B};

\foreach \i/\x in {0/-2.9,1/-2.25,2/-1.6,3/-0.95} {
  \node[block, fill=green!12] (p\i) at (\x,1.85) {$p_{\i}$};
}
\draw[decorate,decoration={brace,amplitude=5pt},thick]
  (-3.25,2.25) -- (-0.6,2.25)
  node[midway,above=6pt] {shared prefix KV};

\foreach \i/\x in {0/0.4,1/1.05,2/1.7} {
  \node[block, fill=blue!10] (a\i) at (\x,1.25) {$a_{\i}$};
}
\foreach \i/\x in {0/0.4,1/1.05} {
  \node[block, fill=orange!14] (b\i) at (\x,0.15) {$b_{\i}$};
}

\draw[arrow] (ra) -- (p0);
\draw[arrow] (rb) -- (p0);

\draw[arrow] (p3) -- (a0);
\draw[arrow] (p3) -- (b0);

\node[align=left, text width=10.5cm] at (0,-1.25) {
  多个请求可以共享相同 prefix 的 KV Cache，只为各自 suffix 分配额外 cache。
  这能降低重复 prefill 和显存占用，但需要引用计数与生命周期管理。
};

\end{tikzpicture}
\caption{Prefix cache 在多个请求之间复用相同前缀}
\label{fig:prefix-cache-reuse}
\end{figure}

下面给出一个极简 prefix cache 管理伪代码。

\begin{lstlisting}[language=Python,caption={Prefix Cache 管理的极简伪代码},label={lst:prefix-cache-manager}]
class PrefixCacheEntry:
def **init**(self, blocks):
self.blocks = blocks
self.ref_count = 0
self.last_used_time = 0

class PrefixCacheManager:
def **init**(self):
self.table = {}  # prefix_hash -> PrefixCacheEntry

def lookup(self, prefix_tokens):
    key = hash_tokens(prefix_tokens)
    entry = self.table.get(key)
    if entry is None:
        return None

    entry.ref_count += 1
    entry.last_used_time = now()
    return entry.blocks

def insert(self, prefix_tokens, blocks):
    key = hash_tokens(prefix_tokens)
    entry = PrefixCacheEntry(blocks)
    entry.ref_count = 1
    entry.last_used_time = now()
    self.table[key] = entry
    return entry.blocks

def release(self, prefix_tokens):
    key = hash_tokens(prefix_tokens)
    entry = self.table.get(key)
    if entry is not None:
        entry.ref_count -= 1

def evict_if_needed(self, free_block_pool):
    # Evict unused prefixes first.
    candidates = [
        (key, entry)
        for key, entry in self.table.items()
        if entry.ref_count == 0
    \]

    # Example policy: least recently used.
    candidates.sort(key=lambda item: item[1].last_used_time)

    for key, entry in candidates:
        if free_block_pool.has_enough_free_blocks():
            break
        free_blocks(entry.blocks)
        del self.table[key]

\end{lstlisting}

真实 prefix caching 要处理更多细节，例如 tokenizer 一致性、prefix hash 冲突、block table 共享、copy-on-write、并发请求、跨 worker cache 命中、cache warmup 和安全隔离。

\section{KV Cache 与并行推理}
\label{sec:kv-cache-and-parallel-inference}

前面主要从单 GPU 或单模型副本角度讨论 KV Cache。实际 serving 中，大模型常使用 tensor parallel、pipeline parallel 或 expert parallel 进行多 GPU 推理。此时 KV Cache 也必须与并行方式匹配。

\subsection{Tensor Parallel 下的 KV Cache}
\label{subsec:kv-cache-under-tensor-parallel}

Tensor Parallel 推理中，attention heads 通常按 TP rank 切分。若 query heads 和 kv heads 都能整除 TP size，则每个 rank 保存一部分 KV heads：

\[
n_{kv,\mathrm{local}}
\frac{n_{kv}}{P_{\mathrm{TP}}}.
\]

每 GPU KV Cache 近似为：

\[
M_{\mathrm{KV,local}}
2
\cdot
L
\cdot
B
\cdot
S
\cdot
n_{kv,\mathrm{local}}
\cdot
d_h
\cdot
s_{\mathrm{dtype}}.
\]

也就是：

\[
M_{\mathrm{KV,local}}
\frac{M_{\mathrm{KV,total}}}{P_{\mathrm{TP}}}.
\]

在这种设计中，每个 rank 本地执行自己负责 heads 的 attention，读取本地 KV Cache。之后 attention output projection 通常需要 TP all-reduce 或其他通信恢复 hidden state。

如果使用 GQA，需要保证 kv heads 与 TP 的映射合理。若 $n_{kv}$ 小于 TP size，某些 TP rank 可能没有独立 kv head，或者需要复制 K/V。实现上可能有多种策略：

\begin{itemize}
\item 让 TP size 不超过 $n_{kv}$；
\item 在多个 query ranks 之间共享 kv heads；
\item 复制 K/V 到多个 ranks；
\item 使用更复杂的 attention kernel mapping。
\end{itemize}

这说明模型架构中的 GQA 参数会影响 serving 并行配置。

\subsection{Pipeline Parallel 下的 KV Cache}
\label{subsec:kv-cache-under-pipeline-parallel}

Pipeline Parallel 推理中，不同 GPU 负责不同层。每个 pipeline stage 只保存自己负责层的 KV Cache。若第 $p$ 个 stage 负责层集合 $\mathcal{L}_p$，则它保存：

\[
{\mat{K}^{(\ell)},\mat{V}^{(\ell)}\mid \ell\in\mathcal{L}_p}.
\]

每个 token 在 decode 时依次经过所有 stages。Stage 0 处理完后把 hidden state 发送给 Stage 1，直到最后 stage 输出 logits。

PP 推理中的 KV Cache 有几个问题：

\begin{itemize}
\item 每个 stage 的 cache 显存取决于它负责的层数；
\item stage 间需要传递每个 decode step 的 hidden state；
\item decode token-by-token 使 pipeline bubble 更明显；
\item 如果 batch 很小，PP 利用率可能较低；
\item pipeline stage 负载不均会影响 TPOT。
\end{itemize}

训练中的 pipeline 依靠多个 micro-batches 填充流水线；推理 decode 中，每个请求每步只有一个 token，流水线填充更加困难。因此很多推理系统更偏向 tensor parallel，而不是过深 pipeline parallel。对于非常大的模型，PP 仍可能必要，但需要结合 continuous batching 提高利用率。

\subsection{KV Cache 与多副本部署}
\label{subsec:kv-cache-and-replica-serving}

在线服务通常部署多个模型副本。一个用户会话的 KV Cache 常驻在某个 worker 上。如果下一轮请求被路由到另一个 worker，就无法直接复用该 cache，除非做跨 worker cache 迁移或共享存储。

这带来 sticky session 问题：

\[
\text{same session}
\rightarrow
\text{same worker}
\]

可以提高 cache 命中率，但可能导致负载不均。例如某些 worker 上积累了很多长会话，显存压力更高；另一些 worker 空闲。若为了负载均衡把会话迁移到其他 worker，则可能失去 KV Cache 复用，增加 prefill 成本。

因此 serving router 需要在以下目标之间权衡：

\begin{itemize}
\item cache locality；
\item worker load balance；
\item GPU memory pressure；
\item request priority；
\item tail latency；
\item session affinity。
\end{itemize}

\section{KV Cache 实现骨架}
\label{sec:kv-cache-implementation-skeleton}

下面给出一个教学版 KV Cache 类。它采用连续分配，不支持分页、prefix sharing 或动态 batch 退出，但有助于理解 cache 的基本操作：初始化、append、读取和长度维护。

\begin{lstlisting}[language=Python,caption={教学版连续 KV Cache 实现骨架},label={lst:kv-cache-class-skeleton}]
import torch

class ContiguousKVCache:
def **init**(
self,
num_layers,
batch_size,
max_seq_len,
num_kv_heads,
head_dim,
dtype=torch.bfloat16,
device="cuda",
):
self.num_layers = num_layers
self.batch_size = batch_size
self.max_seq_len = max_seq_len
self.num_kv_heads = num_kv_heads
self.head_dim = head_dim

    # K/V: [L, B, H_kv, S_max, D]
    self.k_cache = torch.empty(
        num_layers,
        batch_size,
        num_kv_heads,
        max_seq_len,
        head_dim,
        dtype=dtype,
        device=device,
    )
    self.v_cache = torch.empty_like(self.k_cache)

    # Current valid length for each batch slot.
    self.lengths = torch.zeros(
        batch_size,
        dtype=torch.long,
        device=device,
    )

def append(self, layer_idx, batch_indices, k_new, v_new):
    """
    Append K/V for one decode step.

    k_new, v_new: [B_active, H_kv, 1, D]
    batch_indices: [B_active]
    """
    for local_i, batch_i in enumerate(batch_indices):
        pos = self.lengths[batch_i].item()

        if pos >= self.max_seq_len:
            raise RuntimeError("KV cache is full")

        self.k_cache[layer_idx, batch_i, :, pos:pos+1, :] = k_new[local_i]
        self.v_cache[layer_idx, batch_i, :, pos:pos+1, :] = v_new[local_i]

def increment_lengths(self, batch_indices):
    self.lengths[batch_indices] += 1

def get_prefix(self, layer_idx, batch_i):
    """
    Return valid K/V prefix for one request.

    k: [H_kv, S, D]
    v: [H_kv, S, D]
    """
    length = self.lengths[batch_i].item()
    k = self.k_cache[layer_idx, batch_i, :, :length, :]
    v = self.v_cache[layer_idx, batch_i, :, :length, :]
    return k, v

def reset_slot(self, batch_i):
    # We do not need to clear memory immediately.
    # Updating length is enough for correctness.
    self.lengths[batch_i] = 0

\end{lstlisting}

这个实现有许多生产系统不能接受的问题：

\begin{itemize}
\item Python loop 不能用于高性能 decode；
\item batch slot 固定，动态加入退出不灵活；
\item 每个请求预留 \texttt{max\_seq\_len}，浪费显存；
\item 不支持 block table；
\item 不支持 prefix cache；
\item 不支持 paged attention kernel；
\item 不支持 tensor parallel 中的 head 切分；
\item 不支持请求取消后的 block 复用；
\item 不支持异步调度与 stream 管理。
\end{itemize}

但它清楚地展示了 KV Cache 的基本语义：每一层、每个 batch slot、每个 KV head、每个 position 都保存一条 K/V 向量。

生产级实现通常会将 append 和 attention 读取融合到 C++/CUDA/Triton kernel 中，并使用 block table 或 paged memory 避免连续分配问题。

\section{本章小结}
\label{sec:chapter12-summary}

本章系统介绍了 KV Cache 的机制、推理流程、显存估算、性能瓶颈和管理策略。

\begin{itemize}
\item \textbf{LLM 推理分为 prefill 和 decode 两个阶段}。Prefill 一次处理完整 prompt，建立 KV Cache；decode 每次生成一个 token，并持续读取和追加 KV Cache。
\item \textbf{Prefill 更像训练 forward}，计算密度较高，但长 prompt 会带来 $S^2$ attention 成本，并决定 TTFT。
\item \textbf{Decode 是 token-by-token 的强依赖过程}，每步只处理新 token，但需要读取历史所有 K/V，因此容易 memory-bound。
\item \textbf{Key 和 Value 可以缓存}，因为历史 token 在 causal decoder 中不会被未来 token 改变；未来 token 只需要用当前 query 查询历史 key，并加权历史 value。
\item \textbf{Query 通常不缓存}，因为历史 query 不会被未来 token 复用。
\item \textbf{KV Cache shape 通常为 $[L,B,n_{kv},S,d_h]$}，不同实现会根据 attention kernel 和内存访问模式选择不同 layout。
\item \textbf{KV Cache 显存公式} 为 $M_{\mathrm{KV}}=2LBSn_{kv}d_hs_{\mathrm{dtype}}$，其中 $2$ 表示 K 和 V。
\item \textbf{GQA/MQA 能显著降低 KV Cache}，因为 cache 大小与 $n_{kv}$ 成正比。
\item \textbf{Decode 阶段的主要瓶颈常是 HBM 带宽}，因为每步需要读取随上下文长度线性增长的 K/V。
\item \textbf{长上下文会线性增加 KV Cache 显存和 decode 读取量}，直接限制并发能力和成本。
\item \textbf{连续 KV Cache 分配简单但容易浪费显存}，尤其在请求长度不均和动态加入退出的在线 serving 场景中。
\item \textbf{Block-based allocation 更适合动态请求}，它把 cache 切成固定大小 blocks，为 PagedAttention 奠定基础。
\item \textbf{Cache eviction 与 prefix caching 是 serving 系统中的关键策略}，它们在显存占用、重算成本和请求延迟之间做权衡。
\item \textbf{并行推理中的 KV Cache 必须与 TP/PP 配合}：TP 常按 KV heads 切分 cache，PP 则按层切分 cache。
\end{itemize}

下一章将深入 PagedAttention 与 vLLM。我们会看到，KV Cache 最大的工程问题之一不是公式中的显存大小，而是在线请求中长度不均、动态增长和释放带来的碎片。PagedAttention 通过类似操作系统分页的 block table，把逻辑 token 位置映射到物理 cache blocks，从而显著提升显存利用率，并支撑 continuous batching。

\chapter{PagedAttention 与 vLLM}
\label{chap:paged-attention-vllm}

\section{Serving 系统中的内存碎片问题}
\label{sec:serving-memory-fragmentation}

上一章我们讨论了 KV Cache：在自回归推理中，prefill 阶段为 prompt 建立每一层的 Key/Value 缓存，decode 阶段每生成一个新 token 都读取历史 KV，并追加新的 K/V。KV Cache 的基本显存公式为：

[
M_{\mathrm{KV}}
===============

2
\cdot
L
\cdot
B
\cdot
S
\cdot
n_{kv}
\cdot
d_h
\cdot
s_{\mathrm{dtype}}.
]

这个公式描述的是“真实有效 token”需要的理论显存。但在线 serving 系统中的难点并不只是理论显存，而是\textbf{如何在动态请求中高效分配、增长、复用和释放 KV Cache}。

如果所有请求长度相同、同时开始、同时结束，那么 KV Cache 管理会非常简单。系统只需要分配一个规则的矩形张量：

[
[L,B,n_{kv},S,d_h].
]

然而真实 LLM serving 中，请求具有高度动态性：

\begin{itemize}
\item 每个请求的 prompt length 不同；
\item 每个请求的 output length 不同；
\item 请求随时到达；
\item 请求随时结束；
\item 有些请求会被取消；
\item 有些请求会进入多轮对话；
\item 有些请求共享 prefix，有些请求完全不同；
\item 有些请求只生成几个 token，有些请求生成几千 token。
\end{itemize}

因此，KV Cache 更像一个在线内存分配系统，而不是一个静态张量。PagedAttention 的核心价值，正是把 KV Cache 管理从“连续大块分配”改造成“分页式 block 管理”，从而显著减少碎片和 over-reservation。

\subsection{不同请求长度导致的浪费}
\label{subsec:waste-from-variable-request-lengths}

考虑一个简单场景：系统同时服务 $B$ 个请求，每个请求最多允许上下文长度为 $S_{\max}$。若使用最朴素的连续分配方式，系统可能为每个请求都预留 $S_{\max}$ 个 token 的 KV Cache。

第 $i$ 个请求实际使用长度为：

[
S_i=S_{\mathrm{prompt},i}+S_{\mathrm{generated},i}.
]

如果为每个请求预留 $S_{\max}$，则第 $i$ 个请求的 cache 利用率为：

[
u_i=\frac{S_i}{S_{\max}}.
]

整个 batch 的平均利用率为：

[
U
=

\frac{\sum_{i=1}^{B}S_i}
{B S_{\max}}.
]

当请求长度差异很大时，$U$ 可能非常低。例如，一个 batch 中有如下请求：

[
[128,\ 256,\ 512,\ 2048,\ 4096]
]

若都按 $4096$ 预留，则总预留 token slots 为：

[
5\times 4096=20480.
]

真实使用 token slots 为：

[
128+256+512+2048+4096=7040.
]

利用率为：

[
U=\frac{7040}{20480}\approx 34.4%.
]

也就是说，超过一半的 KV Cache 显存只是被预留，却没有存放有效 token。

这类浪费在在线服务中非常常见。很多用户请求很短，例如“解释一个概念”“改写一句话”“写一个 SQL”；也有一些请求很长，例如“总结一篇长文档”“分析代码仓库”“多轮对话历史很长”。如果系统为了兼容长请求而给所有请求预留最大长度，短请求会造成大量内部浪费。

\begin{figure}[H]
\centering
\begin{tikzpicture}[
font=\small,
cell/.style={draw, minimum width=0.42cm, minimum height=0.38cm},
label/.style={font=\scriptsize},
used/.style={fill=blue!12},
unused/.style={fill=gray!18}
]

```
\node[font=\bfseries] at (0,3.05) {连续 KV Cache 预留导致的空间浪费};

\foreach \r/\y/\usedcells in {1/1.65/2,2/1.05/4,3/0.45/6,4/-0.15/10,5/-0.75/16} {
  \node[label] at (-4.8,\y) {Req \r};
  \foreach \i in {0,...,15} {
    \ifnum\i<\usedcells
      \node[cell, used] at ({-3.9+0.45*\i},\y) {};
    \else
      \node[cell, unused] at ({-3.9+0.45*\i},\y) {};
    \fi
  }
}

\node[draw, rounded corners=2pt, fill=blue!12, minimum width=0.55cm, minimum height=0.35cm] at (3.8,1.45) {};
\node[anchor=west] at (4.2,1.45) {used KV};
\node[draw, rounded corners=2pt, fill=gray!18, minimum width=0.55cm, minimum height=0.35cm] at (3.8,0.85) {};
\node[anchor=west] at (4.2,0.85) {reserved unused};

\node[align=left, text width=10.5cm] at (0,-2.0) {
  所有请求都按最大长度预留时，短请求会占用大量未使用 slots。
  对长短请求混部的在线 serving，这种 over-reservation 会显著降低可并发请求数。
};
```

\end{tikzpicture}
\caption{传统连续 KV Cache 分配中的长度浪费}
\label{fig:traditional-contiguous-fragmentation}
\end{figure}

\subsection{over-reservation}
\label{subsec:over-reservation}

Over-reservation 指系统为了避免未来扩容困难，提前为请求预留超过当前需要的 KV Cache 空间。它通常来自两个考虑：

\begin{itemize}
\item 请求未来还会继续 decode，cache 会增长；
\item 连续分配难以原地扩展，因此一开始预留最大长度更简单。
\end{itemize}

设请求当前长度为 $S_{\mathrm{cur}}$，最大允许长度为 $S_{\max}$。连续分配通常分配：

[
A_{\mathrm{contiguous}}=S_{\max}.
]

真实当前使用：

[
U_{\mathrm{cur}}=S_{\mathrm{cur}}.
]

当前浪费：

[
W_{\mathrm{over}}
=================

S_{\max}-S_{\mathrm{cur}}.
]

如果请求最终只生成到 $S_{\mathrm{final}}$，则最终浪费为：

[
W_{\mathrm{final}}
==================

S_{\max}-S_{\mathrm{final}}.
]

这类浪费尤其影响 admission control。假设 GPU 上剩余 KV Cache token capacity 为 $C_{\mathrm{free}}$。如果系统按最大长度预留，则新请求能否进入取决于：

[
S_{\max}\le C_{\mathrm{free}}.
]

但如果按实际增长分配，进入条件可能只需要：

[
S_{\mathrm{prompt}}+\Delta \le C_{\mathrm{free}},
]

其中 $\Delta$ 是少量初始 decode 余量。连续预留会使系统过早拒绝请求，降低吞吐。

Over-reservation 的本质是用显存换实现简单性。对于小模型或低并发服务，这可能可以接受；对于长上下文、高并发、GPU 显存昂贵的 LLM serving，它会成为主要成本来源。

\begin{table}[H]
\centering
\small
\caption{连续预留与按需分配的差异}
\label{tab:over-reservation-vs-demand-allocation}
\begin{tabular}{p{0.24\textwidth}p{0.34\textwidth}p{0.32\textwidth}}
\toprule
\textbf{维度} & \textbf{连续最大长度预留} & \textbf{按需增长分配} \
\midrule
实现复杂度
& 简单
& 需要动态内存管理 \
\midrule
访存模式
& 连续，kernel 读取简单
& 需要地址映射或 block table \
\midrule
显存利用率
& 长短请求混部时较低
& 通常更高 \
\midrule
请求扩容
& 无需扩容，已预留
& 生成过程中申请新 block \
\midrule
碎片风险
& 内部浪费严重，外部碎片也可能出现
& 主要是最后一个 block 的内部碎片 \
\midrule
适用场景
& 固定长度、离线批处理、低并发
& 在线 serving、长短请求混部、高并发 \
\bottomrule
\end{tabular}
\end{table}

\subsection{internal fragmentation}
\label{subsec:internal-fragmentation}

Internal fragmentation 指已经分配给某个请求的内存块内部，有一部分没有被使用。连续预留中的内部碎片可能非常大，因为每个请求都预留到最大长度。

PagedAttention 使用固定大小 block 后，仍然存在内部碎片，但碎片被限制在最后一个 block 内。

设 block size 为 $B_{\mathrm{blk}}$，请求实际长度为 $S_i$。需要分配的 block 数为：

[
N_{\mathrm{blk},i}
==================

\left\lceil
\frac{S_i}{B_{\mathrm{blk}}}
\right\rceil.
]

分配的 token slots 为：

[
A_i
===

N_{\mathrm{blk},i}B_{\mathrm{blk}}.
]

内部碎片为：

[
W_i
===

A_i-S_i.
]

因为只可能最后一个 block 没用满，所以：

[
0\le W_i < B_{\mathrm{blk}}.
]

这意味着每个请求最多浪费 $B_{\mathrm{blk}}-1$ 个 token slots。若 block size 为 $16$，每个请求最多浪费 $15$ 个 token slots；若 block size 为 $32$，最多浪费 $31$ 个 token slots。

与按 $S_{\max}$ 连续预留相比，这个浪费上界小得多。

[
W_{\mathrm{paged}}
<
B_{\mathrm{blk}},
]
[
W_{\mathrm{contiguous}}
=======================

S_{\max}-S_i.
]

当 $S_{\max}$ 为几千或几万，而 $B_{\mathrm{blk}}$ 只有几十时，分页式分配的空间效率优势非常明显。

\begin{figure}[H]
\centering
\begin{tikzpicture}[
font=\small,
cell/.style={draw, minimum width=0.38cm, minimum height=0.34cm},
blockbox/.style={draw, rounded corners=2pt, dashed, thick},
used/.style={fill=green!12},
unused/.style={fill=gray!18}
]

```
\node[font=\bfseries] at (0,3.0) {Block-based 分配中的内部碎片只出现在最后一个 block};

\node[anchor=east] at (-4.7,1.35) {logical tokens};
\foreach \i in {0,...,18} {
  \node[cell, used] at ({-4.1+0.4*\i},1.35) {};
}

\foreach \i in {19,...,23} {
  \node[cell, unused] at ({-4.1+0.4*\i},1.35) {};
}

\draw[blockbox] (-4.32,1.08) rectangle (-1.02,1.62);
\draw[blockbox] (-0.92,1.08) rectangle (2.38,1.62);
\draw[blockbox] (2.48,1.08) rectangle (5.78,1.62);

\node[font=\scriptsize] at (-2.67,0.75) {block 0 full};
\node[font=\scriptsize] at (0.73,0.75) {block 1 full};
\node[font=\scriptsize] at (4.13,0.75) {last block partially used};

\node[draw, rounded corners=2pt, fill=green!12, minimum width=0.55cm, minimum height=0.35cm] at (-1.0,-0.15) {};
\node[anchor=west] at (-0.6,-0.15) {used};
\node[draw, rounded corners=2pt, fill=gray!18, minimum width=0.55cm, minimum height=0.35cm] at (1.2,-0.15) {};
\node[anchor=west] at (1.6,-0.15) {internal fragmentation};

\node[align=left, text width=10.4cm] at (0,-1.45) {
  请求长度不是 block size 的整数倍时，最后一个 block 会有少量未使用 slots。
  这种浪费上界由 block size 决定，而不是由最大上下文长度决定。
};
```

\end{tikzpicture}
\caption{Block-based allocation 中的内部碎片}
\label{fig:block-internal-fragmentation}
\end{figure}

\subsection{为什么传统 batch 管理低效}
\label{subsec:why-traditional-batch-management-is-inefficient}

传统 batch 管理通常假设一个 batch 内的请求同步执行。对于普通深度学习推理，这很自然：一批输入进入模型，一次 forward 后全部完成。但 LLM decode 是迭代式的，一个请求可能生成 5 个 token，另一个请求可能生成 500 个 token。如果仍按静态 batch 管理，会出现严重低效。

静态 batch 的典型问题包括：

\begin{itemize}
\item \textbf{padding waste}：为了对齐长度，短请求被 padding 到长请求长度；
\item \textbf{head-of-line blocking}：短请求已经结束，但必须等长请求完成才能释放 batch；
\item \textbf{GPU 利用率波动}：batch 内请求逐渐结束，实际活跃请求数下降；
\item \textbf{KV Cache 释放不及时}：结束请求的 cache 不能立刻被新请求复用；
\item \textbf{新请求无法及时加入}：需要等整个 batch 结束；
\item \textbf{长请求拖慢短请求}：短请求 tail latency 被长请求影响。
\end{itemize}

Continuous batching 的思想是 iteration-level scheduling：每个 decode iteration 都可以让新请求加入，让完成请求退出。PagedAttention 提供了内存管理基础，使这种动态调度可行。因为请求不再绑定到固定连续 slot，而是持有一组可动态增长和释放的 blocks。

\begin{figure}[H]
\centering
\begin{tikzpicture}[
font=\small,
cell/.style={draw, minimum width=0.58cm, minimum height=0.42cm, align=center},
active/.style={fill=blue!12},
done/.style={fill=gray!18},
wait/.style={fill=orange!14}
]

```
\node[font=\bfseries] at (0,3.05) {Static Batching 中短请求等待长请求};

\foreach \t in {0,...,9} {
  \node[font=\scriptsize] at ({-3.7+0.65*\t},2.25) {t\t};
}

\foreach \r/\y/\len in {1/1.55/3,2/0.95/5,3/0.35/9} {
  \node[font=\scriptsize] at (-4.55,\y) {Req \r};
  \foreach \t in {0,...,9} {
    \ifnum\t<\len
      \node[cell, active] at ({-3.7+0.65*\t},\y) {run};
    \else
      \node[cell, done] at ({-3.7+0.65*\t},\y) {};
    \fi
  }
}

\node[cell, wait, minimum width=1.6cm] at (3.85,0.95) {new req waits};

\node[align=left, text width=10.4cm] at (0,-1.0) {
  静态 batch 中，新请求通常要等当前 batch 结束才能加入。
  短请求虽然已经完成，但资源释放和新请求调度都被长请求阻塞。
};
```

\end{tikzpicture}
\caption{传统静态 batch 管理中的 head-of-line blocking}
\label{fig:static-batching-head-of-line-blocking}
\end{figure}

\section{PagedAttention 核心思想}
\label{sec:paged-attention-core-idea}

PagedAttention 的名字来自操作系统中的分页内存管理。它不是改变 attention 的数学公式，而是改变 KV Cache 的\textbf{物理存储方式}和 attention kernel 的\textbf{读取方式}。

标准 attention 仍然是：

[
\operatorname{Attention}(Q,K,V)
===============================

\softmax
\left(
\frac{QK^{\mathsf{T}}}{\sqrt{d_h}}
\right)V.
]

PagedAttention 改变的是：逻辑上连续的 $K,V$ 序列，不再要求物理上连续存储。每个请求的逻辑 token 序列被切成多个 logical blocks；系统从全局 physical block pool 中分配 physical blocks；block table 记录 logical block 到 physical block 的映射。Attention kernel 在读取历史 K/V 时，通过 block table 做地址翻译。

\subsection{操作系统分页类比}
\label{subsec:os-paging-analogy}

操作系统虚拟内存中，进程看到的是连续虚拟地址空间。实际物理内存可能是不连续的 page frames。页表 page table 负责把虚拟页号映射到物理页框。

PagedAttention 中，请求看到的是连续 token positions：

[
0,1,2,\ldots,S-1.
]

但这些 token 的 KV Cache 可以分散存放在不同 physical blocks 中。Block table 类似页表：

[
\text{logical block id}
\rightarrow
\text{physical block id}.
]

类比如下：

\begin{table}[H]
\centering
\small
\caption{操作系统分页与 PagedAttention 的类比}
\label{tab:os-paging-pagedattention-analogy}
\begin{tabular}{p{0.28\textwidth}p{0.30\textwidth}p{0.32\textwidth}}
\toprule
\textbf{概念} & \textbf{操作系统分页} & \textbf{PagedAttention} \
\midrule
地址空间
& 进程虚拟地址空间
& 请求的逻辑 token 序列 \
\midrule
页
& virtual page
& logical KV block \
\midrule
物理页框
& physical page frame
& physical KV block \
\midrule
页表
& page table
& block table \
\midrule
地址翻译
& virtual address $\rightarrow$ physical address
& token position $\rightarrow$ physical KV address \
\midrule
共享页
& 多进程共享只读页
& 多请求共享 prefix blocks \
\midrule
Copy-on-write
& 写共享页前复制
& 请求分叉或 beam search 时复制 block \
\bottomrule
\end{tabular}
\end{table}

这个类比非常重要，因为它解释了 PagedAttention 的设计哲学：不要要求每个请求在物理显存中连续，而是让 kernel 支持一次轻量级地址翻译。用少量寻址开销换取更高显存利用率和更灵活的调度。

\begin{figure}[H]
\centering
\begin{tikzpicture}[
font=\small,
box/.style={draw, rounded corners=3pt, minimum width=1.9cm, minimum height=0.65cm, align=center},
page/.style={draw, rounded corners=2pt, minimum width=0.68cm, minimum height=0.46cm, align=center},
arrow/.style={-{Latex[length=2.1mm]}, thick}
]

```
\node[font=\bfseries] at (0,3.25) {操作系统分页与 PagedAttention 的结构类比};

% OS side
\node[font=\bfseries] at (-3.7,2.45) {OS Paging};
\foreach \i/\x in {0/-5.0,1/-4.25,2/-3.5,3/-2.75} {
  \node[page, fill=blue!10] (vp\i) at (\x,1.75) {VP\i};
}
\node[box, fill=orange!14] (pt) at (-3.85,0.75) {Page Table};
\foreach \i/\x in {0/-5.0,1/-4.25,2/-3.5,3/-2.75} {
  \node[page, fill=green!10] (pf\i) at (\x,-0.35) {PF\i};
}
\foreach \i in {0,1,2,3} {
  \draw[arrow] (vp\i) -- (pt);
  \draw[arrow] (pt) -- (pf\i);
}

% PA side
\node[font=\bfseries] at (3.7,2.45) {PagedAttention};
\foreach \i/\x in {0/2.35,1/3.1,2/3.85,3/4.6} {
  \node[page, fill=blue!10] (lb\i) at (\x,1.75) {LB\i};
}
\node[box, fill=orange!14] (bt) at (3.5,0.75) {Block Table};
\foreach \i/\x/\pid in {0/2.35/7,1/3.1/2,2/3.85/9,3/4.6/5} {
  \node[page, fill=green!10] (pb\i) at (\x,-0.35) {PB\pid};
}
\foreach \i in {0,1,2,3} {
  \draw[arrow] (lb\i) -- (bt);
  \draw[arrow] (bt) -- (pb\i);
}

\draw[dashed, thick] (0,-0.9) -- (0,2.7);

\node[align=left, text width=10.5cm] at (0,-1.95) {
  请求逻辑上看到连续 token positions，但物理 KV blocks 可以分散存放。
  block table 扮演页表角色，把逻辑位置映射到物理显存地址。
};
```

\end{tikzpicture}
\caption{操作系统分页与 PagedAttention 的类比}
\label{fig:os-paging-vs-pagedattention}
\end{figure}

\subsection{logical block 与 physical block}
\label{subsec:logical-block-and-physical-block}

设 block size 为：

[
B_{\mathrm{blk}}.
]

请求的逻辑 token 位置 $s$ 可以分解为：

[
b_{\mathrm{logical}}
====================

\left\lfloor
\frac{s}{B_{\mathrm{blk}}}
\right\rfloor,
]
[
o
=

s\bmod B_{\mathrm{blk}}.
]

其中 $b_{\mathrm{logical}}$ 是 logical block id，$o$ 是 block 内 offset。

Block table 给出：

[
b_{\mathrm{physical}}
=====================

\operatorname{BlockTable}[r,b_{\mathrm{logical}}],
]

其中 $r$ 是 request id 或 sequence id。

最终物理地址可以抽象为：

[
\operatorname{addr}
===================

\operatorname{Base}
+
b_{\mathrm{physical}}
\cdot
\operatorname{BlockStride}
+
o
\cdot
\operatorname{TokenStride}
+
h
\cdot
\operatorname{HeadStride}
+
d
\cdot
\operatorname{DimStride}.
]

这里 $h$ 是 KV head index，$d$ 是 head dimension index。真实 kernel 会根据具体 layout 优化 stride 和向量化加载。

Logical block 和 physical block 的分离带来三个重要能力：

\begin{itemize}
\item \textbf{非连续物理存储}：请求的逻辑 blocks 可以映射到任意 physical blocks；
\item \textbf{按需增长}：请求需要更多 cache 时，只需追加新的 physical block；
\item \textbf{共享与复制}：多个请求可以引用相同 physical block，写入时再 copy-on-write。
\end{itemize}

\subsection{block table}
\label{subsec:block-table}

Block table 是 PagedAttention 的核心元数据。对于每个请求，它记录每个 logical block 对应的 physical block：

[
\operatorname{BlockTable}
\in
\mathbb{N}^{N_{\mathrm{req}}\times N_{\mathrm{max_blocks}}}.
]

例如，请求 $r$ 的 block table 为：

[
[7,\ 2,\ 9,\ 5].
]

这表示：

[
\text{logical block }0\rightarrow\text{physical block }7,
]
[
\text{logical block }1\rightarrow\text{physical block }2,
]
[
\text{logical block }2\rightarrow\text{physical block }9,
]
[
\text{logical block }3\rightarrow\text{physical block }5.
]

请求逻辑序列仍然是连续的，但 physical blocks 可以任意分散。

\begin{figure}[H]
\centering
\begin{tikzpicture}[
font=\small,
lb/.style={draw, rounded corners=2pt, minimum width=0.85cm, minimum height=0.48cm, align=center, fill=blue!10},
pb/.style={draw, rounded corners=2pt, minimum width=0.85cm, minimum height=0.48cm, align=center, fill=green!10},
table/.style={draw, rounded corners=3pt, minimum width=2.5cm, minimum height=1.3cm, align=center, fill=orange!14},
arrow/.style={-{Latex[length=2.1mm]}, thick}
]

```
\node[font=\bfseries] at (0,3.25) {PagedAttention block table 映射};

\foreach \i/\x in {0/-5.0,1/-4.0,2/-3.0,3/-2.0} {
  \node[lb] (lb\i) at (\x,1.65) {LB \i};
}

\node[table] (bt) at (-3.5,0.45) {
  Block Table\\
  \begin{tabular}{c|c}
  LB & PB\\
  0 & 7\\
  1 & 2\\
  2 & 9\\
  3 & 5
  \end{tabular}
};

\foreach \i/\x/\pid in {0/1.2/2,1/2.2/5,2/3.2/7,3/4.2/9} {
  \node[pb] (pb\pid) at (\x,1.65) {PB \pid};
}

\foreach \i in {0,1,2,3} {
  \draw[arrow] (lb\i) -- (bt);
}

\draw[arrow] (bt) -- (pb7);
\draw[arrow] (bt) -- (pb2);
\draw[arrow] (bt) -- (pb9);
\draw[arrow] (bt) -- (pb5);

\node[align=left, text width=10.4cm] at (0,-1.35) {
  block table 把请求的逻辑 block 映射到全局 physical block pool 中的任意 block。
  Attention kernel 根据该表读取逻辑上连续、物理上不连续的 KV Cache。
};
```

\end{tikzpicture}
\caption{PagedAttention 中 logical block 到 physical block 的映射}
\label{fig:pagedattention-block-table}
\end{figure}

Block table 的大小通常远小于 KV Cache 本身。若每个 block 存 $B_{\mathrm{blk}}$ 个 token，一个请求长度为 $S$，则 block table 只需要：

[
\left\lceil\frac{S}{B_{\mathrm{blk}}}\right\rceil
]

个整数。相比每个 token 的 K/V 向量，metadata 开销很小。

\subsection{attention kernel 如何读取分页 KV}
\label{subsec:how-attention-kernel-reads-paged-kv}

PagedAttention kernel 的核心是在遍历历史 token 时进行地址翻译。对于当前 query token，attention 需要遍历历史位置：

[
s=0,1,\ldots,S-1.
]

对于每个 $s$，kernel 计算：

[
b=\left\lfloor\frac{s}{B_{\mathrm{blk}}}\right\rfloor,
\quad
o=s\bmod B_{\mathrm{blk}}.
]

然后从 block table 读取 physical block id：

[
p=\operatorname{BlockTable}[r,b].
]

再从物理 KV Cache 中读取：

[
\mat{K}*{phys}[p,o,h,:],
\quad
\mat{V}*{phys}[p,o,h,:].
]

因此，PagedAttention 的读取路径可以写成：

[
s
\rightarrow
(b,o)
\rightarrow
p
\rightarrow
\operatorname{addr}(p,o,h,d)
\rightarrow
K/V.
]

下面是一个高度简化的伪代码，用于说明逻辑，不代表高性能 kernel。

\begin{lstlisting}[language=Python,caption={PagedAttention 中根据 block table 读取 KV 的伪代码},label={lst:pagedattention-read-kv-pseudocode}]
def read_kv_for_position(
request_id,
position,
kv_head,
block_size,
block_table,
k_blocks,
v_blocks,
):
logical_block = position // block_size
offset = position % block_size

```
physical_block = block_table[request_id][logical_block]

# k_blocks/v_blocks:
# [num_physical_blocks, num_kv_heads, block_size, head_dim]
k = k_blocks[physical_block, kv_head, offset, :]
v = v_blocks[physical_block, kv_head, offset, :]

return k, v
```

\end{lstlisting}

真正的 PagedAttention kernel 需要处理：

\begin{itemize}
\item 多个 requests；
\item 多个 heads；
\item 每个请求不同 context length；
\item 向量化加载 K/V；
\item block 内连续访存；
\item softmax 的数值稳定；
\item warp/block 级并行；
\item shared memory 与 register 使用；
\item GQA/MQA head mapping；
\item causal mask；
\item 最后一个 block 未满的边界；
\item tensor parallel 下的 head shard。
\end{itemize}

PagedAttention 的挑战在于：它必须在引入 block table 间接寻址的同时，保持 attention kernel 的访存和并行效率。若地址翻译开销过大，显存节省可能被 kernel 开销抵消。因此 block size、layout、kernel tiling 和调度策略必须共同设计。

\begin{figure}[H]
\centering
\begin{tikzpicture}[
font=\small,
box/.style={draw, rounded corners=3pt, minimum width=1.75cm, minimum height=0.65cm, align=center},
arrow/.style={-{Latex[length=2.1mm]}, thick}
]

```
\node[font=\bfseries] at (0,3.25) {PagedAttention kernel 的 KV 读取路径};

\node[box, fill=blue!6] (pos) at (-5.2,1.1) {logical\\position $s$};
\node[box, fill=orange!14] (split) at (-3.0,1.1) {$s\rightarrow(b,o)$};
\node[box, fill=green!10] (table) at (-0.75,1.1) {Block\\Table};
\node[box, fill=yellow!16] (phys) at (1.55,1.1) {physical\\block $p$};
\node[box, fill=purple!10] (addr) at (3.8,1.1) {KV memory\\address};
\node[box, fill=red!8] (kv) at (5.85,1.1) {$K,V$};

\draw[arrow] (pos) -- (split);
\draw[arrow] (split) -- (table);
\draw[arrow] (table) -- (phys);
\draw[arrow] (phys) -- (addr);
\draw[arrow] (addr) -- (kv);

\node[align=left, text width=10.6cm] at (0,-1.0) {
  PagedAttention 在 attention kernel 内部完成逻辑位置到物理 block 的地址翻译。
  逻辑序列连续，物理存储可以离散，从而避免大块连续 KV Cache 分配。
};
```

\end{tikzpicture}
\caption{PagedAttention kernel 通过 block table 读取分页 KV}
\label{fig:pagedattention-kernel-read-path}
\end{figure}

\subsection{copy-on-write 与共享 block}
\label{subsec:copy-on-write-shared-blocks}

PagedAttention 也天然支持共享 prefix。多个请求如果共享相同前缀，可以让它们的 block table 指向相同 physical blocks。每个 physical block 维护引用计数：

[
\operatorname{refcount}(p).
]

当多个请求共享 block 时：

[
\operatorname{refcount}(p)>1.
]

如果某个请求需要写入一个共享 block，而该写入会改变 block 内容，就不能直接原地修改，否则会影响其他请求。此时需要 copy-on-write：

[
p_{\mathrm{new}}=\operatorname{AllocateBlock}(),
]
[
\operatorname{Copy}(p_{\mathrm{old}},p_{\mathrm{new}}),
]
[
\operatorname{BlockTable}[r,b]\leftarrow p_{\mathrm{new}}.
]

在 LLM decode 中，通常只会向最后一个 block 追加 token。共享 prefix blocks 大多是只读的；分叉、beam search、parallel sampling 或 speculative decoding 可能会产生共享后再分叉的场景，此时 copy-on-write 很有用。

\begin{lstlisting}[language=Python,caption={Block 引用计数与 copy-on-write 的简化伪代码},label={lst:block-copy-on-write-pseudocode}]
def ensure_writable_block(request, logical_block_id, block_manager):
physical_block = request.block_table[logical_block_id]

```
if block_manager.ref_count[physical_block] == 1:
    return physical_block

# Copy-on-write.
new_block = block_manager.allocate_block()
copy_block(src=physical_block, dst=new_block)

block_manager.ref_count[physical_block] -= 1
block_manager.ref_count[new_block] = 1

request.block_table[logical_block_id] = new_block
return new_block
```

\end{lstlisting}

Copy-on-write 的意义不只是节省显存，也让复杂 decoding 策略更容易实现。例如 beam search 中多个 beams 共享同一个 prompt prefix，直到它们生成不同 suffix。若每个 beam 都复制完整 KV Cache，显存会迅速膨胀；若共享 prefix blocks，只复制分叉后的 blocks，就能显著降低开销。

\section{vLLM 架构}
\label{sec:vllm-architecture}

vLLM 是围绕高吞吐 LLM serving 设计的推理系统。它的核心思想之一是使用 PagedAttention 管理 KV Cache，并配合 continuous batching、scheduler、worker、block manager 等组件，让在线请求能动态加入、退出和共享 cache。

本节不追逐某个版本的具体代码结构，而从系统架构角度解释一个 vLLM 风格 serving engine 通常如何组织。

\subsection{总体架构}
\label{subsec:vllm-overall-architecture}

一个 vLLM 风格的 serving 系统可以抽象为以下组件：

\begin{itemize}
\item \textbf{API Server / Frontend}：接收 HTTP/gRPC/OpenAI-compatible 请求，处理鉴权、限流、请求解析和流式输出；
\item \textbf{Tokenizer}：将文本转为 token ids，将生成 token 转回文本；
\item \textbf{Scheduler}：决定每个 iteration 哪些请求运行 prefill，哪些请求运行 decode，哪些请求等待；
\item \textbf{Block Manager}：管理 KV Cache physical blocks 的分配、释放、引用计数和 block table；
\item \textbf{Worker / Executor}：在 GPU 上执行模型 forward、PagedAttention kernel、sampling 等；
\item \textbf{Model Runner}：封装模型权重、attention backend、KV Cache layout 和 CUDA graph 等执行细节；
\item \textbf{Output Processor}：处理采样、stop condition、流式 token 返回和请求结束逻辑。
\end{itemize}

它们的关系如下图所示。

\begin{figure}[H]
\centering
\begin{tikzpicture}[
font=\small,
comp/.style={draw, rounded corners=4pt, minimum width=2.05cm, minimum height=0.75cm, align=center},
gpu/.style={draw, rounded corners=4pt, minimum width=2.15cm, minimum height=0.85cm, align=center, fill=blue!6},
arrow/.style={-{Latex[length=2.2mm]}, thick}
]

```
\node[font=\bfseries] at (0,3.5) {vLLM 风格 Serving 架构};

\node[comp, fill=orange!12] (api) at (-5.2,1.95) {API Server\\Frontend};
\node[comp, fill=green!10] (tok) at (-2.75,1.95) {Tokenizer};
\node[comp, fill=yellow!16] (sched) at (-0.25,1.95) {Scheduler};
\node[comp, fill=purple!10] (bm) at (2.35,1.95) {Block\\Manager};
\node[gpu] (worker) at (4.95,1.95) {GPU Worker\\Model Runner};

\node[comp, fill=red!8] (out) at (-0.25,0.35) {Output\\Processor};
\node[comp, fill=blue!8] (kv) at (2.35,0.35) {KV Blocks\\Block Tables};
\node[gpu] (kernels) at (4.95,0.35) {PagedAttention\\Sampling Kernels};

\draw[arrow] (api) -- (tok);
\draw[arrow] (tok) -- (sched);
\draw[arrow] (sched) -- (bm);
\draw[arrow] (bm) -- (worker);
\draw[arrow] (worker) -- (kernels);
\draw[arrow] (bm) -- (kv);
\draw[arrow] (kv) -- (kernels);
\draw[arrow] (kernels) -- (out);
\draw[arrow] (out.west) .. controls (-2.2,0.15) and (-4.6,0.6) .. (api.south);

\node[align=left, text width=10.6cm] at (0,-1.25) {
  Scheduler 决定每个 iteration 的请求集合；Block Manager 为请求分配 KV blocks 并维护 block table；
  Worker 在 GPU 上执行模型和 PagedAttention kernel，Output Processor 处理采样结果与流式返回。
};
```

\end{tikzpicture}
\caption{vLLM 风格 serving engine 的核心组件}
\label{fig:vllm-serving-architecture}
\end{figure}

\subsection{scheduler}
\label{subsec:vllm-scheduler}

Scheduler 是 LLM serving 的大脑。它每个 iteration 都要决定：

\begin{itemize}
\item 哪些新请求可以进入系统；
\item 哪些请求执行 prefill；
\item 哪些请求执行 decode；
\item 哪些请求需要等待；
\item 哪些请求因为资源不足被延迟；
\item 哪些请求已经完成，需要释放资源；
\item 本 iteration 的 token budget 如何分配。
\end{itemize}

与传统 batch scheduler 不同，LLM scheduler 的基本调度单位不是“整个请求”，而是“本 iteration 的 token”。对于 decode 请求，每个 running request 通常只生成一个 token；对于 prefill 请求，一个 prompt 可能包含很多 token，系统可能一次处理完整 prompt，也可能使用 chunked prefill 将长 prompt 分块处理。

Scheduler 常维护几类队列：

\begin{itemize}
\item \textbf{waiting queue}：新到达但尚未分配资源的请求；
\item \textbf{running set}：已经拥有 KV blocks、正在 decode 或 prefill 的请求；
\item \textbf{finished set}：已完成，需要释放 blocks 并返回结果；
\item \textbf{paused/swapped set}：资源压力下暂时暂停或迁出的请求。
\end{itemize}

一个简化调度目标可以写成：

[
\max
\quad
\operatorname{tokens_scheduled}
]

同时满足：

[
\operatorname{KVBlocksNeeded}
\le
\operatorname{FreeBlocks},
]

[
\operatorname{BatchedTokens}
\le
T_{\max},
]

[
\operatorname{NumSequences}
\le
N_{\max}.
]

其中 $T_{\max}$ 是每 iteration 最大 token budget，$N_{\max}$ 是最大并发序列数。调度器需要在吞吐和延迟之间折中：多塞请求可以提升吞吐，但可能增加单请求延迟；优先短请求可以改善平均延迟，但可能让长请求饥饿；优先 decode 可以降低 TPOT，但可能让新请求 TTFT 变差。

\begin{lstlisting}[language=Python,caption={Iteration-level scheduler 的简化伪代码},label={lst:vllm-scheduler-pseudocode}]
def schedule_iteration(waiting, running, block_manager, token_budget, max_seqs):
scheduled_prefills = []
scheduled_decodes = []

```
remaining_tokens = token_budget
remaining_slots = max_seqs - len(running)

# 1. Schedule decode for already running requests.
for req in list(running):
    if req.is_finished():
        continue

    if not block_manager.can_append_slot(req):
        continue

    scheduled_decodes.append(req)
    remaining_tokens -= 1

    if remaining_tokens <= 0:
        break

# 2. Admit new requests if there is budget and KV memory.
while waiting and remaining_tokens > 0 and remaining_slots > 0:
    req = waiting[0]
    prefill_tokens = req.prompt_len

    if prefill_tokens > remaining_tokens:
        break

    if not block_manager.can_allocate_for_prefill(req):
        break

    waiting.pop(0)
    block_manager.allocate_for_prefill(req)
    running.add(req)

    scheduled_prefills.append(req)
    remaining_tokens -= prefill_tokens
    remaining_slots -= 1

return scheduled_prefills, scheduled_decodes
```

\end{lstlisting}

真实 scheduler 会更复杂。它会考虑 chunked prefill、prefix cache 命中、优先级、超时、beam search、多模态输入、LoRA adapter、speculative decoding、抢占和资源回收。但核心仍然是：在每个 iteration 中，根据 KV Cache 资源和 token budget 选择请求集合。

\subsection{worker}
\label{subsec:vllm-worker}

Worker 负责在 GPU 上执行实际计算。它通常接收 scheduler 产生的执行计划，包括：

\begin{itemize}
\item 本次 prefill 的 requests；
\item 本次 decode 的 requests；
\item input token ids；
\item position ids；
\item slot mapping；
\item block tables；
\item sequence lengths；
\item sampling parameters；
\item LoRA 或其他 adapter 信息。
\end{itemize}

Worker 的执行流程可以抽象为：

\begin{enumerate}
\item 根据调度计划准备输入张量；
\item 调用模型 forward；
\item attention kernel 根据 block table 读取 KV Cache；
\item 将新 token 的 K/V 写入对应 physical blocks；
\item 计算 logits；
\item 执行 sampling；
\item 返回 sampled tokens 和请求状态。
\end{enumerate}

Worker 的关键性能目标是减少 CPU/GPU 同步、减少小 kernel、减少动态 shape 造成的开销，并尽量让 GPU 一直有足够工作。

在高性能实现中，worker 可能使用 CUDA graph 固化部分执行路径，以减少 launch overhead。但 LLM serving 的动态性很强，batch size、sequence length、请求组合都在变化，CUDA graph 捕获和复用需要做 shape bucketing 或其他工程折中。

\subsection{block manager}
\label{subsec:vllm-block-manager}

Block Manager 是 PagedAttention 系统中的内存管理器。它维护全局 physical block pool，并为每个请求维护 block table。

Block Manager 的职责包括：

\begin{itemize}
\item 初始化 GPU KV block pool；
\item 为 prefill 请求分配足够 blocks；
\item 为 decode 追加 token 时分配新 block；
\item 维护 request 到 block table 的映射；
\item 维护 physical block 引用计数；
\item 请求结束时释放 blocks；
\item prefix cache 命中时共享 blocks；
\item copy-on-write；
\item 在资源不足时触发 eviction、pause 或拒绝请求。
\end{itemize}

一个极简 Block Manager 可以写成：

\begin{lstlisting}[language=Python,caption={Block Manager 的简化实现骨架},label={lst:block-manager-skeleton}]
import math
from collections import deque

class BlockManager:
def **init**(self, num_blocks, block_size):
self.num_blocks = num_blocks
self.block_size = block_size

```
    self.free_blocks = deque(range(num_blocks))
    self.ref_count = [0 for _ in range(num_blocks)]

    # request_id -> list of physical block ids
    self.block_tables = {}

def num_required_blocks(self, num_tokens):
    return math.ceil(num_tokens / self.block_size)

def can_allocate(self, num_tokens):
    return len(self.free_blocks) >= self.num_required_blocks(num_tokens)

def allocate(self, request_id, num_tokens):
    n = self.num_required_blocks(num_tokens)
    if len(self.free_blocks) < n:
        raise RuntimeError("not enough KV blocks")

    blocks = []
    for _ in range(n):
        block = self.free_blocks.popleft()
        self.ref_count[block] = 1
        blocks.append(block)

    self.block_tables[request_id] = blocks
    return blocks

def append_slot(self, request_id, current_length):
    blocks = self.block_tables[request_id]

    # Allocate a new block when current length reaches a block boundary.
    if current_length % self.block_size == 0:
        if not self.free_blocks:
            raise RuntimeError("no free block for decode")
        block = self.free_blocks.popleft()
        self.ref_count[block] = 1
        blocks.append(block)

    return blocks

def release(self, request_id):
    blocks = self.block_tables.pop(request_id, [])

    for block in blocks:
        self.ref_count[block] -= 1
        if self.ref_count[block] == 0:
            self.free_blocks.append(block)
```

\end{lstlisting}

真实 Block Manager 还要处理共享 prefix、copy-on-write、CPU swap、并发安全、GPU/CPU metadata 同步和调度器交互。它是 serving 系统中最容易影响稳定性和显存利用率的组件之一。

\subsection{continuous batching}
\label{subsec:vllm-continuous-batching}

Continuous batching 指系统在每个 decode iteration 重新组织 batch，让新请求可以加入，让完成请求可以退出。它与 PagedAttention 是互补关系：

\begin{itemize}
\item Continuous batching 解决\textbf{请求调度动态性}；
\item PagedAttention 解决\textbf{动态请求的 KV Cache 内存管理}。
\end{itemize}

如果没有 PagedAttention，continuous batching 仍然会遇到连续 KV Cache 分配困难：请求加入和退出会导致 batch slot 移动、cache 拷贝、碎片和 over-reservation。PagedAttention 让请求的 cache 独立于 batch slot，只需要更新 block table 和调度 metadata。

一个 iteration-level timeline 如下：

\begin{figure}[H]
\centering
\begin{tikzpicture}[
font=\scriptsize,
cell/.style={draw, minimum width=0.7cm, minimum height=0.42cm, align=center},
decode/.style={fill=blue!12},
prefill/.style={fill=green!12},
done/.style={fill=gray!18},
newreq/.style={fill=orange!14}
]

```
\node[font=\bfseries\small] at (0,3.0) {vLLM 风格 Continuous Batching 时间线};

\foreach \t in {0,...,8} {
  \node at ({-3.6+0.78*\t},2.35) {iter \t};
}

\foreach \r/\y in {A/1.65,B/1.1,C/0.55,D/0.0} {
  \node at (-4.45,\y) {Req \r};
}

% A decode then finish
\foreach \t in {0,1,2,3} {
  \node[cell, decode] at ({-3.6+0.78*\t},1.65) {D};
}
\foreach \t in {4,...,8} {
  \node[cell, done] at ({-3.6+0.78*\t},1.65) {};
}

% B long decode
\foreach \t in {0,...,8} {
  \node[cell, decode] at ({-3.6+0.78*\t},1.1) {D};
}

% C joins at 2
\foreach \t in {0,1} {
  \node[cell, done] at ({-3.6+0.78*\t},0.55) {};
}
\node[cell, prefill] at ({-3.6+0.78*2},0.55) {P};
\foreach \t in {3,...,8} {
  \node[cell, decode] at ({-3.6+0.78*\t},0.55) {D};
}

% D joins at 5
\foreach \t in {0,...,4} {
  \node[cell, done] at ({-3.6+0.78*\t},0.0) {};
}
\node[cell, prefill] at ({-3.6+0.78*5},0.0) {P};
\foreach \t in {6,...,8} {
  \node[cell, decode] at ({-3.6+0.78*\t},0.0) {D};
}

\node[draw, rounded corners=2pt, fill=green!12, minimum width=0.5cm, minimum height=0.32cm] at (3.8,1.5) {};
\node[anchor=west] at (4.15,1.5) {Prefill};
\node[draw, rounded corners=2pt, fill=blue!12, minimum width=0.5cm, minimum height=0.32cm] at (3.8,1.0) {};
\node[anchor=west] at (4.15,1.0) {Decode};
\node[draw, rounded corners=2pt, fill=gray!18, minimum width=0.5cm, minimum height=0.32cm] at (3.8,0.5) {};
\node[anchor=west] at (4.15,0.5) {inactive};

\node[align=left, text width=10.5cm, font=\small] at (0,-1.35) {
  每个 iteration 都可以重新组成 batch。已完成请求释放 blocks，新请求分配 blocks 后加入。
  PagedAttention 让这种动态调度不需要移动已有请求的 KV Cache。
};
```

\end{tikzpicture}
\caption{Continuous batching 中请求动态加入和退出}
\label{fig:vllm-continuous-batching-timeline}
\end{figure}

图中 P 表示 prefill，D 表示 decode。Req C 和 Req D 在不同 iteration 加入；Req A 结束后释放资源；Req B 持续 decode。每个 iteration 的实际 batch 都不同。

\subsection{prefix caching 与 speculative decoding}
\label{subsec:vllm-prefix-caching-speculative-decoding}

Prefix caching 与 PagedAttention 结合得很自然，因为共享 prefix 可以通过共享 physical blocks 实现。多个请求的 block table 前缀指向相同 blocks，refcount 增加即可。只有当请求在共享 prefix 后继续生成自己的 suffix 时，才分配新的 blocks。

Speculative decoding 则引入另一类 cache 管理问题。它通常包含 draft model 和 target model：draft model 先快速生成多个候选 token，target model 验证这些 token。若候选 token 被接受，则可以一次提交多个 token；若被拒绝，需要回滚到某个位置。

从 KV Cache 管理角度看，speculative decoding 需要处理：

\begin{itemize}
\item 草稿 token 的临时 KV；
\item 被接受 token 的 KV 提交；
\item 被拒绝 token 的 KV 回滚；
\item draft 和 target 的 cache 是否分开；
\item 多 token 提交时 block table 的更新；
\item 回滚时释放未提交 blocks。
\end{itemize}

Paged block 管理对这些场景有帮助，因为未提交 token 可能只占用少量 blocks；回滚时释放这些 blocks 即可，而不需要移动连续大 cache。

\begin{figure}[H]
\centering
\begin{tikzpicture}[
font=\small,
block/.style={draw, rounded corners=2pt, minimum width=0.55cm, minimum height=0.38cm, align=center},
arrow/.style={-{Latex[length=2.0mm]}, thick}
]

```
\node[font=\bfseries] at (0,3.15) {Speculative Decoding 中的提交与回滚};

\foreach \i/\x in {0/-4.0,1/-3.4,2/-2.8,3/-2.2} {
  \node[block, fill=green!12] (c\i) at (\x,1.4) {$c$};
}
\node[anchor=east] at (-4.5,1.4) {committed};

\foreach \i/\x in {0/-1.2,1/-0.6,2/0.0,3/0.6} {
  \node[block, fill=orange!14] (d\i) at (\x,1.4) {$d$};
}
\node[anchor=south] at (-0.3,1.85) {draft tokens};

\foreach \i/\x in {0/1.8,1/2.4} {
  \node[block, fill=blue!12] at (\x,1.4) {$a$};
}
\foreach \i/\x in {0/3.0,1/3.6} {
  \node[block, fill=gray!18] at (\x,1.4) {$r$};
}

\draw[arrow] (0.9,0.9) -- node[below] {verify} (1.55,0.9);

\node[align=center, font=\scriptsize] at (2.1,0.65) {accepted};
\node[align=center, font=\scriptsize] at (3.35,0.65) {rollback};

\node[align=left, text width=10.5cm] at (0,-0.95) {
  Speculative decoding 会产生临时 token。验证通过的 token 提交到正式 KV Cache；
  未通过的 token 对应 blocks 需要释放或回滚。分页式 block 管理让这种操作更自然。
};
```

\end{tikzpicture}
\caption{Speculative decoding 中 KV blocks 的提交与回滚}
\label{fig:speculative-decoding-kv-commit-rollback}
\end{figure}

\section{PagedAttention 的性能收益}
\label{sec:pagedattention-performance-benefits}

PagedAttention 的收益不能简单理解为“attention 算法更少 FLOPs”。它的主要价值在于\textbf{提高 KV Cache 显存利用率}，从而允许更多请求同时运行，并使 continuous batching 更有效。吞吐提升往往来自更高并发、更低内存浪费和更好的调度，而不是单个 attention step 的数学计算减少。

\subsection{显存利用率提升}
\label{subsec:pagedattention-memory-utilization}

连续分配中，每个请求分配 $S_{\max}$ token slots：

[
A_{\mathrm{contiguous},i}=S_{\max}.
]

PagedAttention 中，每个请求按 block 分配：

[
A_{\mathrm{paged},i}
====================

B_{\mathrm{blk}}
\left\lceil
\frac{S_i}{B_{\mathrm{blk}}}
\right\rceil.
]

总利用率分别为：

[
U_{\mathrm{contiguous}}
=======================

\frac{\sum_i S_i}
{\sum_i S_{\max}},
]

[
U_{\mathrm{paged}}
==================

\frac{\sum_i S_i}
{\sum_i B_{\mathrm{blk}}
\left\lceil
\frac{S_i}{B_{\mathrm{blk}}}
\right\rceil}.
]

当 $B_{\mathrm{blk}}\ll S_{\max}$ 时，PagedAttention 的利用率通常显著更高。

下面用一个简单例子说明。假设 block size 为 $16$，请求长度为：

[
[128,\ 256,\ 512,\ 2048,\ 4096].
]

PagedAttention 分配 slots 为：

[
[128,\ 256,\ 512,\ 2048,\ 4096],
]

因为这些长度恰好是 16 的倍数，利用率接近 $100%$。若长度为：

[
[130,\ 257,\ 519,\ 2050,\ 4097],
]

则分配 slots 为：

[
[144,\ 272,\ 528,\ 2064,\ 4112].
]

总使用为：

[
7053.
]

总分配为：

[
7120.
]

利用率为：

[
\frac{7053}{7120}\approx 99.1%.
]

相比连续最大长度预留的约 $34%$，差异非常大。当然真实系统中还有 metadata、block table、alignment 和其他 buffer，但核心趋势成立。

\begin{table}[H]
\centering
\small
\caption{连续预留与 PagedAttention 分配 token slots 对比}
\label{tab:contiguous-vs-paged-slots}
\begin{tabular}{rrrrr}
\toprule
\textbf{请求} & \textbf{实际长度} & \textbf{连续预留} & \textbf{Paged block size 16} & \textbf{Paged 浪费} \
\midrule
1 & 130 & 4096 & 144 & 14 \
2 & 257 & 4096 & 272 & 15 \
3 & 519 & 4096 & 528 & 9 \
4 & 2050 & 4096 & 2064 & 14 \
5 & 4097 & 4096 & 4112 & 15 \
\bottomrule
\end{tabular}
\end{table}

表中第 5 个请求超过了 $4096$，说明连续预留若固定为 $4096$ 甚至无法容纳；若把 $S_{\max}$ 提高到更大，短请求浪费会更严重。Paged allocation 则可以随请求增长追加 blocks。

\subsection{吞吐提升}
\label{subsec:pagedattention-throughput-improvement}

PagedAttention 提升吞吐的核心路径可以写成：

[
\text{更高 KV 利用率}
\rightarrow
\text{同样显存容纳更多请求}
\rightarrow
\text{decode batch 更大}
\rightarrow
\text{GPU 利用率更高}
\rightarrow
\text{tokens/s 提升}.
]

Decode 阶段每个请求每次只生成一个 token。如果 batch 太小，GEMM 和 attention kernel 难以充分利用 GPU。通过降低 KV Cache 浪费，系统可以保留更多 running requests，从而提高 decode batch size。

但吞吐提升不是无限的。Decode batch 变大后，KV Cache 读取量也增加。最终系统可能受 HBM 带宽、attention kernel、sampling、CPU scheduler 或网络输出限制。PagedAttention 解决的是 KV Cache 分配和调度瓶颈，不会消除 decode attention 本身随上下文长度增长的读取成本。

可以把单 GPU serving 吞吐粗略看作：

[
\operatorname{Throughput}
\approx
\frac{
\operatorname{ActiveRequests}
}
{
T_{\mathrm{iteration}}
}.
]

PagedAttention 增加的是可同时活跃的请求数，但 $T_{\mathrm{iteration}}$ 也可能随着 batch 和上下文增长而上升。实际收益取决于请求长度分布、模型大小、GPU 显存、block size、kernel 实现和 scheduler 策略。

\subsection{长短请求混部}
\label{subsec:mixed-long-short-requests}

长短请求混部是在线 serving 的常态。PagedAttention 对混部尤其重要，因为它避免了短请求被长请求的最大长度预留拖累。

在连续分配中，如果 batch 内允许一个长请求达到 $S_{\max}=32768$，那么短请求也可能被迫预留接近这个上限。这样会严重降低短请求并发能力。

在 paged allocation 中，短请求只分配少量 blocks；长请求随着生成或 prefill 逐步占用更多 blocks。两类请求可以更自然地共存。

这对实际服务有几个影响：

\begin{itemize}
\item 短请求不会因为系统支持长上下文而付出巨大显存预留；
\item 长请求不会要求一开始就占用所有最大上下文 cache；
\item 请求结束后 blocks 可立即回收；
\item scheduler 可以根据剩余 blocks 做 admission control；
\item prefix cache 可让多个长请求共享共同前缀；
\item chunked prefill 可减少超长 prompt 对 decode 的阻塞。
\end{itemize}

\begin{figure}[H]
\centering
\begin{tikzpicture}[
font=\small,
cell/.style={draw, minimum width=0.35cm, minimum height=0.32cm},
long/.style={fill=green!12},
short/.style={fill=blue!12},
free/.style={fill=gray!16}
]

```
\node[font=\bfseries] at (0,3.0) {PagedAttention 支持长短请求混部};

\node[anchor=east] at (-4.8,1.55) {Short reqs};
\foreach \r/\y/\len in {1/1.75/3,2/1.25/2,3/0.75/4} {
  \foreach \i in {0,...,5} {
    \ifnum\i<\len
      \node[cell, short] at ({-4.2+0.38*\i},\y) {};
    \else
      \node[cell, free] at ({-4.2+0.38*\i},\y) {};
    \fi
  }
}

\node[anchor=east] at (0.5,1.25) {Long req};
\foreach \i in {0,...,17} {
  \node[cell, long] at ({1.0+0.38*\i},1.25) {};
}

\node[align=left, text width=10.4cm] at (0,-0.8) {
  短请求只占用少量 blocks；长请求占用更多 blocks。
  二者不需要按同一个最大长度预留，因此显存利用率更高。
};
```

\end{tikzpicture}
\caption{长短请求在 paged KV Cache 中自然共存}
\label{fig:mixed-long-short-requests-paged-kv}
\end{figure}

\subsection{trade-off 与适用边界}
\label{subsec:pagedattention-tradeoffs-boundaries}

PagedAttention 不是免费的。它用间接寻址和更复杂的内存管理换取更高显存利用率。主要 trade-off 包括：

\begin{itemize}
\item \textbf{地址翻译开销}：kernel 需要通过 block table 找到 physical block；
\item \textbf{metadata 开销}：需要维护 block table、refcount、request state；
\item \textbf{kernel 复杂度}：attention 读取不再是简单连续地址；
\item \textbf{block size 选择困难}：太大浪费多，太小 metadata 和寻址开销大；
\item \textbf{最后 block 仍有碎片}：内部碎片上界为 block size；
\item \textbf{prefill 长 prompt 仍然昂贵}：PagedAttention 不改变 $S^2$ prefill attention 成本；
\item \textbf{decode 长上下文仍然 memory-bound}：它不消除读取历史 KV 的线性成本；
\item \textbf{多 GPU 下更复杂}：TP/PP、prefix cache、block table 同步需要额外设计。
\end{itemize}

Block size 是一个典型权衡。设 block size 为 $B_{\mathrm{blk}}$。较大的 block size：

\begin{itemize}
\item block table 更小；
\item kernel 内 block 处理更规整；
\item metadata 更少；
\item 但最后 block 内部碎片更大。
\end{itemize}

较小的 block size：

\begin{itemize}
\item 内部碎片更小；
\item 按需增长更精细；
\item 但 block table 更大；
\item 地址翻译更多；
\item kernel 可能更难高效向量化。
\end{itemize}

因此 PagedAttention 适合动态在线 serving、请求长度差异大、高并发、KV Cache 显存紧张的场景。对于离线批处理、请求长度固定、显存充足或 kernel 极端追求连续访存的场景，简单连续布局可能也有竞争力。

\section{PagedAttention 的工程实现细节}
\label{sec:pagedattention-engineering-details}

为了让本章更贴近工程实践，本节进一步展开 PagedAttention 落地时的几个关键问题：block size、物理内存布局、slot mapping、prefill 写入、decode 追加、以及与采样策略的交互。

\subsection{block size 选择}
\label{subsec:block-size-choice}

Block size 是 PagedAttention 的核心超参数之一。它决定了每个 physical block 能存多少 token 的 K/V。设 block size 为 $B_{\mathrm{blk}}$，每个 token 的 KV Cache 大小为：

[
M_{\mathrm{token}}
==================

2
\cdot
L
\cdot
n_{kv}
\cdot
d_h
\cdot
s_{\mathrm{dtype}}.
]

如果每个 block 覆盖一个请求的 $B_{\mathrm{blk}}$ 个 token，则每个请求级 block 的 KV 大小为：

[
M_{\mathrm{block}}
==================

B_{\mathrm{blk}}
\cdot
M_{\mathrm{token}}.
]

实际系统中，物理 KV block 可能按 layer、head 或 worker 内部 layout 切分，因此实现细节会不同。但从逻辑上看，block size 越大，单次分配粒度越粗。

选择 block size 时需要同时考虑：

\begin{itemize}
\item 平均请求长度；
\item 最后 block 内部碎片；
\item block table metadata 大小；
\item attention kernel block 内连续读取效率；
\item GPU memory alignment；
\item prefix sharing 粒度；
\item copy-on-write 粒度；
\item 调度器分配和释放开销。
\end{itemize}

一个简单的期望碎片估算是：如果请求长度对 block size 的余数近似均匀分布，则每个请求最后 block 的平均浪费约为：

[
\mathbb{E}[W]
\approx
\frac{B_{\mathrm{blk}}-1}{2}.
]

因此 block size 从 $16$ 增加到 $32$，平均每请求浪费 token slots 约增加一倍。但如果 block size 太小，kernel 每次处理的连续 token 过少，metadata 和地址翻译开销上升。

\subsection{physical block layout}
\label{subsec:physical-block-layout}

Paged KV Cache 的物理 layout 需要服务 attention kernel。一个可能的 layout 是：

[
\mat{K}*{blocks}
\in
\mathbb{R}^{
N*{\mathrm{blocks}}
\times
n_{kv}
\times
B_{\mathrm{blk}}
\times
d_h
}.
]

V Cache 同理。也可以把 layer 维放在外面：

[
[L,N_{\mathrm{blocks}},n_{kv},B_{\mathrm{blk}},d_h].
]

或者为每层单独维护 block pool。不同 layout 的取舍包括：

\begin{itemize}
\item 每层独立 block pool 管理简单，但 metadata 多；
\item 所有层合并的 block 可以减少请求级 metadata，但 block 更大；
\item head 维连续有利于某些向量化加载；
\item block 内 token 维连续有利于顺序扫描历史位置；
\item head_dim 维连续通常有利于 vectorized load。
\end{itemize}

高性能 kernel 通常会对 layout 做非常具体的优化。本书不展开某个实现的所有细节，但读者应记住：PagedAttention 的设计不是只有 block table，\textbf{block table、memory layout 和 attention kernel 必须一起设计}。

\subsection{slot mapping}
\label{subsec:slot-mapping}

在每个 iteration 中，scheduler 选出一批要执行的 tokens。每个 token 的新 K/V 需要写入某个 physical block 的某个 offset。为了让 GPU kernel 知道写到哪里，系统通常会构造 slot mapping。

对一个 token，slot mapping 记录：

[
\text{token}
\rightarrow
(\text{physical block id},\ \text{offset}).
]

对于 decode，每个请求通常追加一个 token。设请求当前长度为 $S_i$，新 token 位置为 $S_i$。则：

[
b=\left\lfloor\frac{S_i}{B_{\mathrm{blk}}}\right\rfloor,
\quad
o=S_i\bmod B_{\mathrm{blk}}.
]

若 $o=0$，说明需要申请新 block。之后从 block table 得到 physical block id：

[
p=\operatorname{BlockTable}[i,b].
]

Slot mapping 为：

[
(p,o).
]

对于 prefill，一个请求一次写入多个 token。系统可以为这些 token 批量生成 slot mapping。若 prompt 跨多个 blocks，slot mapping 会跨越多个 physical blocks。

\begin{lstlisting}[language=Python,caption={为 decode token 构造 slot mapping 的简化伪代码},label={lst:decode-slot-mapping-pseudocode}]
def build_decode_slot_mapping(requests, block_manager, block_size):
slot_mapping = []

```
for req in requests:
    pos = req.current_length
    logical_block = pos // block_size
    offset = pos % block_size

    if offset == 0:
        block_manager.append_slot(req.request_id, pos)

    physical_block = req.block_table[logical_block]
    slot_mapping.append((physical_block, offset))

return slot_mapping
```

\end{lstlisting}

Slot mapping 是连接 scheduler、block manager 和 model runner 的关键桥梁。Scheduler 只知道请求要生成 token；Block Manager 知道请求的 blocks；Model Runner 需要具体写入地址。Slot mapping 把这三者连接起来。

\subsection{prefill 与 decode 的不同路径}
\label{subsec:prefill-and-decode-different-paths}

Prefill 和 decode 对 PagedAttention 的要求不同。

Prefill 中，一个请求可能一次写入很多 token。系统需要：

\begin{itemize}
\item 为 prompt 长度申请足够 blocks；
\item 批量写入 K/V；
\item 对 prompt 内部执行 causal attention；
\item 生成最后位置 logits；
\item 处理长 prompt 的 chunking。
\end{itemize}

Decode 中，每个请求每 iteration 通常只写入一个 token。系统需要：

\begin{itemize}
\item 检查当前 block 是否有空位；
\item 必要时申请新 block；
\item 读取完整历史 blocks；
\item 写入新 token K/V；
\item 执行 sampling；
\item 更新请求长度和 stop 状态。
\end{itemize}

因此，serving engine 可能为 prefill 和 decode 使用不同的 attention kernel 或不同调度策略。Prefill 更接近大矩阵计算，decode 更受 KV 读取和 batch 组织影响。PagedAttention 主要解决两者共同的 KV Cache 存储问题，但它在 decode 阶段的收益往往更直接，因为 decode 请求持续增长、加入和退出最频繁。

\section{本章小结}
\label{sec:chapter13-summary}

本章系统介绍了 PagedAttention 与 vLLM 风格 serving 架构。核心思想可以概括为：\textbf{把 KV Cache 从连续大块分配改造成分页式 block 管理，并让 attention kernel 通过 block table 读取逻辑连续、物理离散的 K/V。}

\begin{itemize}
\item \textbf{LLM serving 中的 KV Cache 是动态内存管理问题}：请求长度不同、到达时间不同、结束时间不同，cache 会不断增长和释放。
\item \textbf{连续 KV Cache 分配会产生 over-reservation}：如果每个请求都按最大上下文长度预留，短请求会浪费大量显存。
\item \textbf{Block-based allocation 将浪费限制在最后一个 block}：若 block size 为 $B_{\mathrm{blk}}$，每个请求最多浪费 $B_{\mathrm{blk}}-1$ 个 token slots。
\item \textbf{PagedAttention 类似操作系统分页}：请求逻辑 token 序列对应 logical blocks，物理 KV Cache 对应 physical blocks，block table 类似 page table。
\item \textbf{Block table 是核心元数据}：它记录每个请求的 logical block 到 physical block 的映射，使请求无需连续显存。
\item \textbf{Attention kernel 需要通过 block table 读取 KV}：逻辑位置 $s$ 被分解为 logical block id 和 offset，再映射到 physical block 地址。
\item \textbf{共享 prefix 可以通过共享 physical blocks 实现}：多个请求的 block table 指向相同 prefix blocks，并使用 refcount 管理生命周期。
\item \textbf{Copy-on-write 支持分叉场景}：beam search、parallel sampling、speculative decoding 等场景可以共享前缀，只在写入分叉时复制 block。
\item \textbf{vLLM 风格架构由 scheduler、block manager、worker、model runner 等组件组成}：scheduler 负责 iteration-level 调度，block manager 负责 KV blocks，worker 在 GPU 上执行模型。
\item \textbf{Continuous batching 与 PagedAttention 相互配合}：前者让请求动态加入和退出，后者让这种动态调度不需要移动连续 KV Cache。
\item \textbf{PagedAttention 的吞吐收益主要来自更高显存利用率和更大有效并发}，而不是改变 attention 的数学 FLOPs。
\item \textbf{PagedAttention 也有 trade-off}：block table 间接寻址、metadata 管理、kernel 复杂度和 block size 选择都需要工程权衡。
\item \textbf{Prefill 与 decode 的需求不同}：prefill 更偏大块计算和批量写入，decode 更偏动态追加、历史 KV 读取和调度效率。
\end{itemize}

下一章将进入 FlashAttention。PagedAttention 解决的是 KV Cache 的内存管理和动态调度问题；FlashAttention 解决的是 Attention 计算中的 IO 瓶颈。我们会从标准 Attention 需要物化 $S\times S$ attention matrix 的问题出发，推导 online softmax，理解为什么把 Q/K/V 分块搬入 SRAM 计算可以显著减少 HBM 读写，并最终形成高性能 fused attention kernel。

\chapter{FlashAttention 算法}
\label{chap:flash-attention}

\section{标准 Attention 的 IO 瓶颈}
\label{sec:standard-attention-io-bottleneck}

上一章讨论了 PagedAttention。PagedAttention 解决的是在线 serving 中 KV Cache 的\textbf{内存管理问题}：请求长度不一、动态加入退出、KV Cache 增长和释放，都会造成碎片与 over-reservation。PagedAttention 通过 block table 让逻辑连续的 KV Cache 可以物理离散存储。

本章讨论另一个方向：\textbf{Attention 计算本身的 IO 瓶颈}。这就是 FlashAttention 要解决的问题。

标准 Attention 的数学公式非常简单：

\[
\operatorname{Attention}(\mat{Q},\mat{K},\mat{V})
\softmax
\left(
\frac{\mat{Q}\mat{K}^{\mathsf{T}}}{\sqrt{d_h}}
\right)
\mat{V}.
\]

如果只从数学上看，计算过程似乎只是三步：

\[
\mat{S}=\frac{\mat{Q}\mat{K}^{\mathsf{T}}}{\sqrt{d_h}},
\]
\[
\mat{P}=\softmax(\mat{S}),
\]
\[
\mat{O}=\mat{P}\mat{V}.
\]

但在 GPU 上，性能不只取决于 FLOPs，还取决于数据如何在 HBM、L2、shared memory、register 之间移动。标准 Attention 最大的问题，是它通常会显式物化巨大的中间矩阵：

\[
\mat{S}\in\mathbb{R}^{S\times S},
\quad
\mat{P}\in\mathbb{R}^{S\times S}.
\]

其中 $S$ 是序列长度。随着上下文长度增长，$S^2$ 级别的中间矩阵会带来巨大的 HBM 读写压力。FlashAttention 的核心不是改变 Attention 的数学结果，而是改变计算路径：\textbf{不把完整 attention matrix 写入 HBM，而是在片上 SRAM/register 中分块完成 softmax 和 weighted sum}。

\subsection{$QK^{\mathsf{T}}$ 的中间矩阵}
\label{subsec:qk-intermediate-matrix}

考虑单个 batch、单个 attention head。设：

\[
\mat{Q},\mat{K},\mat{V}\in\mathbb{R}^{S\times d_h}.
\]

标准 Attention 第一步计算 score matrix：

\[
\mat{S}=\mat{Q}\mat{K}^{\mathsf{T}},
\]

其中：

\[
\mat{S}\in\mathbb{R}^{S\times S}.
\]

第 $i$ 行第 $j$ 列为：

\[
S_{ij}
\langle \vect{q}_i,\vect{k}_j\rangle.
\]

对于 causal attention，还需要 mask：

\[
S_{ij}=-\infty,\quad j>i.
\]

然后对每一行做 softmax：

\[
P_{ij}
\frac{\exp(S_{ij}-m_i)}
{\sum_{j'}\exp(S_{ij'}-m_i)},
\]
其中 $m_i=\max_j S_{ij}$ 用于数值稳定。

最后输出：

\[
\vect{o}_i
\sum_j P_{ij}\vect{v}_j.
\]

问题出现在 $\mat{S}$ 和 $\mat{P}$。如果 $S=4096$，则单 head 的 attention matrix 元素个数为：

\[
4096^2=16{,}777{,}216.
\]

若使用 FP16/BF16，每个元素 2 bytes，则一个矩阵约为：

\[
32\ \mathrm{MB}.
\]

如果同时保存 score、softmax probability、dropout mask 或 backward 所需中间结果，显存和 HBM traffic 会更高。对于多 head、多 batch、多层，这个开销会迅速放大。

\[
M_{\mathrm{attn\ matrix}}
B\cdot n_h\cdot S^2\cdot s_{\mathrm{dtype}}.
\]

例如：

\[
B=8,\quad n_h=32,\quad S=4096,\quad s_{\mathrm{dtype}}=2.
\]

则仅一个 $B\times n_h\times S\times S$ 矩阵就需要：

\[
8\times 32\times 4096^2\times 2
\approx
8.59\ \mathrm{GB}.
\]

这只是单层 attention 的一个中间矩阵。训练时还要考虑 backward 保存和重读，开销更加不可接受。

\begin{figure}[H]
\centering
\begin{tikzpicture}[
font=\small,
tensor/.style={draw, rounded corners=3pt, minimum width=1.55cm, minimum height=0.75cm, align=center},
big/.style={draw, rounded corners=3pt, minimum width=2.0cm, minimum height=1.15cm, align=center},
arrow/.style={-{Latex[length=2.2mm]}, thick}
]

\node[font=\bfseries] at (0,3.35) {标准 Attention 显式物化 $S\times S$ 中间矩阵};

\node[tensor, fill=blue!8] (q) at (-5.2,1.4) {$Q$\\$S,d_h$};
\node[tensor, fill=blue!8] (k) at (-5.2,0.35) {$K$\\$S,d_h$};
\node[big, fill=orange!14] (score) at (-2.6,0.85) {$S=QK^T$\\$S,S$};
\node[big, fill=green!12] (prob) at (0.25,0.85) {$P=\softmax(S)$\\$S,S$};
\node[tensor, fill=blue!8] (v) at (2.65,0.2) {$V$\\$S,d_h$};
\node[tensor, fill=purple!10] (o) at (5.0,0.85) {$O=PV$\\$S,d_h$};

\draw[arrow] (q) -- (score);
\draw[arrow] (k) -- (score);
\draw[arrow] (score) -- node[above] {write/read HBM} (prob);
\draw[arrow] (prob) -- (o);
\draw[arrow] (v) -- (o);

\node[draw, rounded corners=3pt, fill=red!8, minimum width=2.6cm, minimum height=0.55cm, align=center] at (-1.2,-0.95)
  {large $S\times S$ traffic};

\node[align=left, text width=10.5cm] at (0,-2.05) {
  标准实现往往把 score matrix 和 softmax matrix 写入 HBM，再从 HBM 读回。
  当序列长度变大时，中间矩阵读写成为主要瓶颈。
};

\end{tikzpicture}
\caption{标准 Attention 的中间矩阵物化流程}
\label{fig:standard-attention-materialization}
\end{figure}

\subsection{HBM 与 SRAM 的带宽差异}
\label{subsec:hbm-vs-sram-bandwidth}

现代 GPU 有多级存储层次。对深度学习 kernel 来说，最重要的几层是：

\begin{itemize}
\item \textbf{HBM}：容量大，带宽高，但相比片上存储仍然慢；
\item \textbf{L2 Cache}：片上共享缓存，容量比 HBM 小得多；
\item \textbf{Shared Memory / SRAM}：每个 SM 内部可编程片上存储，延迟低、带宽高、容量有限；
\item \textbf{Registers}：线程私有，最快但数量有限。
\end{itemize}

标准 Attention 的问题是，它在 HBM 中读写了太多中间结果。即使 GPU 的 Tensor Core 对矩阵乘法很快，如果数据需要反复从 HBM 搬进搬出，整体性能仍会被内存带宽限制。

可以把一个 GPU kernel 的执行理解为：

\[
\text{HBM}
\rightarrow
\text{SRAM/register}
\rightarrow
\text{compute}
\rightarrow
\text{HBM}.
\]

如果中间矩阵 $\mat{S}$ 和 $\mat{P}$ 都写回 HBM，再读回继续计算，则数据路径变成：

\[
Q,K,V
\rightarrow
\mat{S}\ \text{write HBM}
\rightarrow
\mat{S}\ \text{read HBM}
\rightarrow
\mat{P}\ \text{write HBM}
\rightarrow
\mat{P}\ \text{read HBM}
\rightarrow
O.
\]

FlashAttention 的目标是减少这类 HBM 往返。它将 Q、K、V 按块搬入 SRAM/register，在片上计算 score、softmax 和局部输出累积，中间 attention matrix 不落地到 HBM。

\begin{figure}[H]
\centering
\begin{tikzpicture}[
font=\small,
mem/.style={draw, rounded corners=4pt, minimum width=3.0cm, minimum height=0.75cm, align=center},
arrow/.style={-{Latex[length=2.2mm]}, thick}
]

\node[font=\bfseries] at (0,3.25) {GPU 存储层次与 Attention IO};

\node[mem, fill=orange!14] (hbm) at (0,2.0) {HBM\\large capacity, high latency};
\node[mem, fill=green!12] (l2) at (0,0.95) {L2 Cache};
\node[mem, fill=blue!8] (smem) at (0,-0.1) {Shared Memory / SRAM\\small capacity, high bandwidth};
\node[mem, fill=purple!10] (reg) at (0,-1.15) {Registers\\fastest, tiny};

\draw[arrow] (hbm) -- (l2);
\draw[arrow] (l2) -- (smem);
\draw[arrow] (smem) -- (reg);

\node[align=left, text width=10.5cm] at (0,-2.35) {
  FlashAttention 的关键是把尽可能多的中间计算留在 SRAM/register 中完成，
  避免把 $S\times S$ attention matrix 写入 HBM。
};

\end{tikzpicture}
\caption{HBM 与片上 SRAM/register 的存储层次}
\label{fig:hbm-sram-memory-hierarchy}
\end{figure}

\subsection{计算复杂度与 IO 复杂度}
\label{subsec:compute-complexity-and-io-complexity}

标准 Attention 和 FlashAttention 在数学上都计算精确 Attention。它们的 FLOPs 复杂度同阶：

\[
O(S^2d_h).
\]

这意味着 FlashAttention 并不是通过把 attention 变成稀疏或近似来减少数学计算。它的收益来自 IO 复杂度下降。

为了理解这一点，需要区分：

\begin{itemize}
\item \textbf{计算复杂度}：执行多少乘加、指数、除法等运算；
\item \textbf{IO 复杂度}：数据在 HBM 与片上存储之间搬运多少次。
\end{itemize}

在 GPU 上，如果 kernel 是 compute-bound，优化 FLOPs 更重要；如果 kernel 是 memory-bound，减少 HBM 读写更重要。标准 Attention 中，$S\times S$ 中间矩阵导致 HBM traffic 很高，尤其在长序列下，IO 成本可能主导运行时间。

粗略看，标准 Attention 至少涉及：

\[
\text{read }Q,K,V,
\]
\[
\text{write }S,
\]
\[
\text{read }S,
\]
\[
\text{write }P,
\]
\[
\text{read }P,V,
\]
\[
\text{write }O.
\]

其中 $S$ 和 $P$ 都是 $S\times S$ 大矩阵。FlashAttention 则尽量只读 Q/K/V，写 O，并保存少量 softmax 统计量，例如每行最大值和归一化分母：

\[
m_i,\quad l_i.
\]

因此，它把 attention 的显存复杂度从需要保存 $O(S^2)$ 中间矩阵，降低到只需要 $O(S)$ 级别的辅助统计量和 $O(Sd_h)$ 的输出。

\begin{table}[H]
\centering
\small
\caption{标准 Attention 与 FlashAttention 的关键区别}
\label{tab:standard-vs-flash-attention}
\begin{tabular}{p{0.24\textwidth}p{0.32\textwidth}p{0.34\textwidth}}
\toprule
\textbf{维度} & \textbf{标准 Attention} & \textbf{FlashAttention} \
\midrule
数学结果
& 精确 Attention
& 精确 Attention \
\midrule
FLOPs 复杂度
& $O(S^2d_h)$
& $O(S^2d_h)$，可能有少量重计算 \
\midrule
中间矩阵
& 常物化 $S\times S$ score/probability
& 不物化完整 attention matrix \
\midrule
主要优化目标
& 简单分步执行
& 减少 HBM 读写，提升 IO efficiency \
\midrule
显存占用
& 需要保存较大中间矩阵
& 只保存输出和少量 softmax 统计量 \
\midrule
适用价值
& 短序列或简单实现
& 长序列、高 batch、训练和 prefill 场景 \
\bottomrule
\end{tabular}
\end{table}

\subsection{为什么不是 FLOPs 越少越快}
\label{subsec:why-fewer-flops-not-always-faster}

很多性能优化初学者会自然地认为：算法 FLOPs 越少，速度越快。但在 GPU 上，这并不总成立。一个 kernel 的速度常常受以下因素共同决定：

\begin{itemize}
\item HBM 读写量；
\item L2 cache 命中；
\item shared memory 使用；
\item register pressure；
\item Tensor Core 利用率；
\item warp occupancy；
\item memory coalescing；
\item kernel launch 和同步；
\item 数据 layout；
\item 是否有中间 tensor 落地。
\end{itemize}

FlashAttention 的一个重要启发是：\textbf{宁可做一些额外计算，也要避免昂贵的 HBM 读写}。尤其在 backward 中，FlashAttention 可以选择不保存完整 attention probability，而是在 backward 时重新计算需要的 score 和 softmax 块。这增加了 FLOPs，但减少了巨大的中间矩阵存储和读取。

这与 gradient checkpointing 的思想类似：用重计算换显存和 IO。区别在于 FlashAttention 的重计算发生在更细粒度的 attention kernel 内部，并且与 tiling 和 online softmax 结合。

\begin{figure}[H]
\centering
\begin{tikzpicture}[
font=\small,
axis/.style={-{Latex[length=2.2mm]}, thick},
point/.style={circle, draw, fill=blue!10, minimum size=0.18cm},
line/.style={thick}
]

\node[font=\bfseries] at (0,3.0) {Roofline 视角：Attention 可能受 IO 而非 FLOPs 限制};

\draw[axis] (-5.0,-1.2) -- (5.1,-1.2) node[right] {arithmetic intensity};
\draw[axis] (-5.0,-1.2) -- (-5.0,2.2) node[above] {performance};

\draw[line, orange!80!black] (-4.7,-1.0) -- (-1.5,1.5);
\draw[line, orange!80!black] (-1.5,1.5) -- (4.7,1.5);

\node[point, fill=red!18] (std) at (-3.2,0.15) {};
\node[anchor=west] at (-3.0,0.15) {standard attention};

\node[point, fill=green!18] (flash) at (0.1,1.2) {};
\node[anchor=west] at (0.3,1.2) {FlashAttention};

\node[align=center, font=\scriptsize] at (-3.2,-0.45) {memory-bound};
\node[align=center, font=\scriptsize] at (2.2,1.85) {compute-bound ceiling};

\node[align=left, text width=10.5cm] at (0,-2.35) {
  减少 HBM traffic 可以提高 arithmetic intensity，使 kernel 从 memory-bound 区域向更高性能区域移动。
  因此即使 FLOPs 没有减少，甚至略有增加，wall-clock 时间仍可能下降。
};

\end{tikzpicture}
\caption{从 roofline 模型理解 FlashAttention 的 IO-aware 思路}
\label{fig:attention-roofline-intuition}
\end{figure}

\section{Online Softmax}
\label{sec:online-softmax}

FlashAttention 的数学核心之一是 online softmax。Attention 的困难在于 softmax 的分母依赖整行所有元素：

\[
\softmax(x_i)
\frac{e^{x_i}}
{\sum_j e^{x_j}}.
\]

如果把 score matrix 按列分块计算，那么在只看到一个块时，并不知道整行的全局最大值和全局分母。Online softmax 解决了这个问题：它允许我们一块一块处理 score，同时维护每一行的最大值和指数和，最终得到与一次性 softmax 完全一致的结果。

\subsection{softmax 的数值稳定性}
\label{subsec:softmax-numerical-stability}

直接计算：

\[
\frac{e^{x_i}}{\sum_j e^{x_j}}
\]

容易溢出。数值稳定做法是减去最大值：

\[
m=\max_j x_j.
\]

然后：

\[
\softmax(x_i)
\frac{e^{x_i-m}}
{\sum_j e^{x_j-m}}.
\]

因为：

\[
\frac{e^{x_i-m}}{\sum_j e^{x_j-m}}
# \frac{e^{x_i}e^{-m}}{\sum_j e^{x_j}e^{-m}}
\frac{e^{x_i}}{\sum_j e^{x_j}}.
\]

减去 $m$ 不改变 softmax 结果，但能防止指数溢出。

标准 softmax 通常分两步：

\[
m=\max_j x_j,
\]
\[
l=\sum_j e^{x_j-m}.
\]

然后输出：

\[
y_i=\frac{e^{x_i-m}}{l}.
\]

如果 $x$ 是 attention score 的一行，那么 $j$ 遍历所有 key positions。FlashAttention 分块处理时，需要在多个 blocks 之间维护同样的 $m$ 和 $l$。

\subsection{分块计算中的 max 与 sum 更新}
\label{subsec:blockwise-max-sum-update}

假设一行 score 被分成多个 blocks。当前已经处理过若干 blocks，维护：

\[
m_{\mathrm{old}}
\max_{\text{old blocks}} x_j,
\]
\[
l_{\mathrm{old}}
\sum_{\text{old blocks}} e^{x_j-m_{\mathrm{old}}}.
\]

现在处理新 block，计算该 block 的局部最大值：

\[
m_{\mathrm{blk}}
\max_{\text{new block}} x_j,
\]

和局部指数和：

\[
l_{\mathrm{blk}}
\sum_{\text{new block}} e^{x_j-m_{\mathrm{blk}}}.
\]

合并 old blocks 与 new block 的全局最大值为：

\[
m_{\mathrm{new}}
\max(m_{\mathrm{old}},m_{\mathrm{blk}}).
\]

旧的指数和 $l_{\mathrm{old}}$ 是以 $m_{\mathrm{old}}$ 为基准的。要改成以 $m_{\mathrm{new}}$ 为基准，需要乘缩放因子：

\[
e^{m_{\mathrm{old}}-m_{\mathrm{new}}}.
\]

新 block 的指数和也要改成以 $m_{\mathrm{new}}$ 为基准：

\[
e^{m_{\mathrm{blk}}-m_{\mathrm{new}}}.
\]

因此：

\[
l_{\mathrm{new}}
e^{m_{\mathrm{old}}-m_{\mathrm{new}}}l_{\mathrm{old}}
+
e^{m_{\mathrm{blk}}-m_{\mathrm{new}}}l_{\mathrm{blk}}.
\]

这就是 online softmax 的核心更新公式。

\subsection{避免完整 attention matrix}
\label{subsec:avoid-full-attention-matrix}

Attention 输出不是 softmax 本身，而是：

\[
\vect{o}
\sum_j
\frac{e^{x_j-m}}{l}
\vect{v}_j.
\]

分块时，不仅要维护 $m$ 和 $l$，还要维护未归一化的输出累积。设旧输出累积为：

\[
\vect{o}_{\mathrm{old}}
\sum_{\text{old blocks}}
\frac{e^{x_j-m_{\mathrm{old}}}}
{l_{\mathrm{old}}}
\vect{v}_j.
\]

更方便的做法是维护带分母的累积量：

\[
\vect{a}_{\mathrm{old}}
# l_{\mathrm{old}}\vect{o}_{\mathrm{old}}
\sum_{\text{old blocks}}
e^{x_j-m_{\mathrm{old}}}\vect{v}_j.
\]

新 block 的累积量为：

\[
\vect{a}_{\mathrm{blk}}
\sum_{\text{new block}}
e^{x_j-m_{\mathrm{blk}}}\vect{v}_j.
\]

合并到新基准 $m_{\mathrm{new}}$ 后：

\[
\vect{a}_{\mathrm{new}}
e^{m_{\mathrm{old}}-m_{\mathrm{new}}}\vect{a}*{\mathrm{old}}
+
e^{m*{\mathrm{blk}}-m_{\mathrm{new}}}\vect{a}_{\mathrm{blk}}.
\]

最终输出：

\[
\vect{o}_{\mathrm{new}}
\frac{\vect{a}*{\mathrm{new}}}{l*{\mathrm{new}}}.
\]

这意味着我们可以一块一块处理 score 和 V，持续更新：

\[
m_i,\quad l_i,\quad \vect{a}_i,
\]

而不需要保存整行 $S$ 个 softmax probability，更不需要保存完整 $S\times S$ attention matrix。

\begin{figure}[H]
\centering
\begin{tikzpicture}[
font=\small,
block/.style={draw, rounded corners=3pt, minimum width=1.45cm, minimum height=0.65cm, align=center},
state/.style={draw, rounded corners=3pt, minimum width=1.4cm, minimum height=0.55cm, align=center, fill=green!10},
arrow/.style={-{Latex[length=2.1mm]}, thick}
]

\node[font=\bfseries] at (0,3.2) {Online Softmax 分块更新};

\node[block, fill=blue!8] (b0) at (-5.1,1.3) {score\\block 0};
\node[state] (s0) at (-3.3,1.3) {$m,l,a$};

\node[block, fill=blue!8] (b1) at (-1.7,1.3) {score\\block 1};
\node[state] (s1) at (0.1,1.3) {$m,l,a$};

\node[block, fill=blue!8] (b2) at (1.7,1.3) {score\\block 2};
\node[state] (s2) at (3.5,1.3) {$m,l,a$};

\node[block, fill=purple!10] (out) at (5.2,1.3) {$O=a/l$};

\draw[arrow] (b0) -- (s0);
\draw[arrow] (s0) -- (b1);
\draw[arrow] (b1) -- (s1);
\draw[arrow] (s1) -- (b2);
\draw[arrow] (b2) -- (s2);
\draw[arrow] (s2) -- (out);

\node[align=left, text width=10.5cm] at (0,-0.65) {
  每处理一个 score block，就更新该行的最大值 $m$、指数和 $l$ 与输出累积 $a$。
  最终结果与一次性 softmax 完全一致，但不需要物化完整 attention row。
};

\end{tikzpicture}
\caption{Online softmax 的分块归约流程}
\label{fig:online-softmax-blockwise}
\end{figure}

\subsection{online softmax 推导}
\label{subsec:online-softmax-derivation}

下面把推导整理成更系统的形式。对某一行 attention score，分块处理。第 $r$ 个 block 的元素集合为 $\mathcal{B}_r$。

局部最大值：

\[
m_r=\max_{j\in\mathcal{B}_r} x_j.
\]

局部指数和：

\[
l_r=\sum_{j\in\mathcal{B}_r}e^{x_j-m_r}.
\]

局部 value 累积：

\[
\vect{a}_r
\sum_{j\in\mathcal{B}_r}
e^{x_j-m_r}\vect{v}_j.
\]

若当前已经合并到第 $r-1$ 个 block，状态为：

\[
m^{(r-1)},\quad l^{(r-1)},\quad \vect{a}^{(r-1)}.
\]

合并第 $r$ 个 block：

\[
m^{(r)}
\max(m^{(r-1)},m_r).
\]

\[
l^{(r)}
e^{m^{(r-1)}-m^{(r)}}l^{(r-1)}
+
e^{m_r-m^{(r)}}l_r.
\]

\[
\vect{a}^{(r)}
e^{m^{(r-1)}-m^{(r)}}\vect{a}^{(r-1)}
+
e^{m_r-m^{(r)}}\vect{a}_r.
\]

处理完所有 blocks 后：

\[
\vect{o}
\frac{\vect{a}^{(R)}}{l^{(R)}}.
\]

这个公式的关键在于缩放因子。每当全局最大值变大，旧累积量必须被重新缩放，否则 softmax 分母不一致。FlashAttention 的 kernel 就是把这种数学更新嵌入到 tiled attention 计算中。

下面给出一个教学版 online softmax attention 的 Python 伪代码。它强调数学流程，不代表高性能实现。

\begin{lstlisting}[language=Python,caption={教学版 Online Softmax Attention 伪代码},label={lst:online-softmax-attention-pseudocode}]
import math
import torch

def online_attention_one_row(q, k, v, block_size):
"""
q: [d]
k: [seq_len, d]
v: [seq_len, d]
return: [d]
"""
seq_len = k.shape[0]
scale = 1.0 / math.sqrt(q.shape[0])

m = torch.tensor(-float("inf"), device=q.device)
l = torch.tensor(0.0, device=q.device)
acc = torch.zeros_like(q)

for start in range(0, seq_len, block_size):
    end = min(start + block_size, seq_len)

    k_blk = k[start:end]  # [B_blk, d]
    v_blk = v[start:end]  # [B_blk, d]

    scores = (k_blk @ q) * scale  # [B_blk]
    m_blk = torch.max(scores)
    p_blk = torch.exp(scores - m_blk)
    l_blk = torch.sum(p_blk)
    acc_blk = p_blk @ v_blk

    m_new = torch.maximum(m, m_blk)

    old_scale = torch.exp(m - m_new)
    blk_scale = torch.exp(m_blk - m_new)

    acc = old_scale * acc + blk_scale * acc_blk
    l = old_scale * l + blk_scale * l_blk
    m = m_new

return acc / l

\end{lstlisting}

实际 FlashAttention 会同时处理一个 $Q$ block 的多行，而不是一行；会在 shared memory/register 中计算 block GEMM；会进行向量化、warp-level reduction 和 Tensor Core MMA。上面代码只是展示 online softmax 的数学骨架。

\section{FlashAttention 核心算法}
\label{sec:flashattention-core-algorithm}

FlashAttention 的核心可以概括为四句话：

\begin{enumerate}
\item 将 $Q,K,V$ 沿序列维分块；
\item 将 $K,V$ blocks 搬入 SRAM；
\item 对 $Q$ blocks 计算局部 attention，并用 online softmax 更新输出；
\item 不把完整 $S\times S$ attention matrix 写入 HBM。
\end{enumerate}

它的设计目标不是减少 Attention 的数学复杂度，而是最小化 HBM 与片上存储之间的 IO。

\subsection{tiling}
\label{subsec:flashattention-tiling}

设序列长度为 $S$。将 Q 按行切成 blocks：

\[
\mat{Q}_1,\mat{Q}*2,\ldots,\mat{Q}*{T_Q},
\]

每个 Q block 大小约为：

\[
B_Q\times d_h.
\]

将 K/V 也按序列维切成 blocks：

\[
(\mat{K}_1,\mat{V}_1),
(\mat{K}*2,\mat{V}*2),
\ldots,
(\mat{K}*{T_K},\mat{V}*{T_K}),
\]

每个 K/V block 大小约为：

\[
B_K\times d_h.
\]

对于一个 Q block 和一个 K block，局部 score 为：

\[
\mat{S}_{ij}
\mat{Q}_i\mat{K}_j^{\mathsf{T}}.
\]

它的 shape 为：

\[
B_Q\times B_K.
\]

这个局部矩阵可以放在 SRAM/register 中处理，而不需要写入 HBM。对所有 K/V blocks 迭代，使用 online softmax 更新 Q block 对应的输出。

\[
\mat{O}_i
\operatorname{Attention}(\mat{Q}_i,\mat{K},\mat{V}).
\]

最终只将 $\mat{O}_i$ 写回 HBM。

\begin{figure}[H]
\centering
\begin{tikzpicture}[
font=\small,
tile/.style={draw, rounded corners=2pt, minimum width=0.9cm, minimum height=0.55cm, align=center},
mat/.style={draw, rounded corners=3pt, minimum width=1.5cm, minimum height=1.0cm, align=center},
arrow/.style={-{Latex[length=2.1mm]}, thick}
]

\node[font=\bfseries] at (0,3.35) {FlashAttention Tiling 数据流};

\node[mat, fill=blue!8] (q) at (-5.2,1.25) {$Q$};
\foreach \i/\y in {0/1.75,1/1.25,2/0.75} {
  \node[tile, fill=blue!12] at (-5.2,\y) {$Q_{\i}$};
}

\node[mat, fill=green!10] (kv) at (-2.8,1.25) {$K,V$};
\foreach \i/\y in {0/1.75,1/1.25,2/0.75} {
  \node[tile, fill=green!14] at (-2.8,\y) {$K_{\i},V_{\i}$};
}

\node[mat, fill=orange!14, minimum width=2.2cm] (sram) at (0.0,1.25) {SRAM/register\\block compute};

\node[mat, fill=purple!10] (online) at (2.9,1.25) {online\\softmax update};

\node[mat, fill=red!8] (out) at (5.2,1.25) {$O$\\write once};

\draw[arrow] (q) -- (sram);
\draw[arrow] (kv) -- (sram);
\draw[arrow] (sram) -- (online);
\draw[arrow] (online) -- (out);

\node[align=left, text width=10.5cm] at (0,-0.85) {
  $Q,K,V$ 按块搬入片上存储。局部 score、mask、softmax 和 $PV$ 累积都在片上完成。
  完整 $S\times S$ attention matrix 不写入 HBM。
};

\end{tikzpicture}
\caption{FlashAttention 的 tiling 数据流}
\label{fig:flashattention-tiling-dataflow}
\end{figure}

\subsection{SRAM 中计算 block}
\label{subsec:compute-block-in-sram}

FlashAttention 的一个典型循环结构可以抽象为：

\begin{enumerate}
\item 外层遍历 K/V blocks；
\item 内层遍历 Q blocks；
\item 将当前 K/V block 加载到 SRAM；
\item 将当前 Q block 加载到 SRAM/register；
\item 计算局部 score；
\item 应用 causal mask 或 attention mask；
\item 计算局部 softmax 统计；
\item 更新输出累积；
\item 最后写回 O。
\end{enumerate}

不同实现可能选择不同循环顺序。关键原则是：让 K/V block 被加载后尽可能复用，减少 HBM 访问；让局部 score matrix 只在片上存在。

对一个 Q block，维护每一行的：

\[
m_i,\quad l_i,\quad \vect{a}_i.
\]

每处理一个 K/V block，更新这些状态。最后输出：

\[
\vect{o}_i=\frac{\vect{a}_i}{l_i}.
\]

如果是 causal attention，当 K block 的位置超过 Q block 当前行允许关注的位置，需要 mask 掉未来位置：

\[
S_{ij}=-\infty,\quad j>i.
\]

在 block 级别，causal mask 可能出现三种情况：

\begin{itemize}
\item 当前 K block 完全在 Q block 之前：无需 mask；
\item 当前 K block 完全在 Q block 之后：整块跳过；
\item 当前 K block 与 Q block 对角线相交：需要块内 mask。
\end{itemize}

这种 block-level 判断可以避免不必要计算。

\subsection{fused kernel}
\label{subsec:flashattention-fused-kernel}

标准实现可能将 Attention 拆成多个 kernel：

\[
QK^{\mathsf{T}}\ \text{kernel},
\]
\[
mask/scale\ \text{kernel},
\]
\[
softmax\ \text{kernel},
\]
\[
dropout\ \text{kernel},
\]
\[
PV\ \text{kernel}.
\]

每个 kernel 都可能把中间结果写回 HBM，再被下一个 kernel 读入。FlashAttention 则将这些操作融合到一个或少数几个 kernel 中：

\[
QK^{\mathsf{T}}
+
mask
+
softmax
+
dropout
+
PV
\rightarrow
\text{fused attention kernel}.
\]

Kernel fusion 的收益包括：

\begin{itemize}
\item 减少 kernel launch overhead；
\item 减少 HBM 中间读写；
\item 中间值停留在 register/shared memory；
\item 更容易做跨操作优化；
\item 更低显存峰值。
\end{itemize}

但 fused kernel 也更复杂。它需要同时管理矩阵乘法、softmax、mask、dropout、数值稳定、寄存器压力、shared memory、warp 分工和边界条件。FlashAttention 的难点正是把这些操作合并后仍然保持高性能。

\begin{figure}[H]
\centering
\begin{tikzpicture}[
font=\small,
box/.style={draw, rounded corners=3pt, minimum width=1.45cm, minimum height=0.58cm, align=center},
arrow/.style={-{Latex[length=2.1mm]}, thick}
]

\node[font=\bfseries] at (0,3.1) {Kernel Fusion：从多次 HBM 往返到单个 fused kernel};

% standard
\node[font=\bfseries] at (-3.5,2.35) {Standard};
\node[box, fill=blue!8] (a1) at (-5.5,1.45) {$QK^T$};
\node[box, fill=orange!14] (a2) at (-3.8,1.45) {mask};
\node[box, fill=green!10] (a3) at (-2.1,1.45) {softmax};
\node[box, fill=purple!10] (a4) at (-0.4,1.45) {$PV$};
\draw[arrow] (a1) -- node[above,font=\scriptsize] {HBM} (a2);
\draw[arrow] (a2) -- node[above,font=\scriptsize] {HBM} (a3);
\draw[arrow] (a3) -- node[above,font=\scriptsize] {HBM} (a4);

% flash
\node[font=\bfseries] at (2.5,2.35) {FlashAttention};
\node[box, fill=red!8, minimum width=3.0cm, minimum height=0.85cm] (fused) at (3.2,1.45)
  {$QK^T$ + mask + softmax + $PV$};
\node[box, fill=yellow!16] (o) at (5.7,1.45) {$O$};
\draw[arrow] (fused) -- node[above,font=\scriptsize] {write once} (o);

\node[align=left, text width=10.5cm] at (0,-0.45) {
  标准实现的多个 kernel 之间需要通过 HBM 传递中间矩阵。
  FlashAttention 把这些步骤融合，让中间 attention block 不落地。
};

\end{tikzpicture}
\caption{FlashAttention 的 fused kernel 思路}
\label{fig:flashattention-fused-kernel}
\end{figure}

\subsection{forward pass}
\label{subsec:flashattention-forward-pass}

FlashAttention forward 的核心状态包括：

\[
\mat{O}\in\mathbb{R}^{S\times d_h},
\]
\[
\vect{m}\in\mathbb{R}^{S},
\]
\[
\vect{l}\in\mathbb{R}^{S}.
\]

其中 $\vect{m}$ 保存每一行的运行最大值，$\vect{l}$ 保存每一行的运行 softmax 分母。输出 $\mat{O}$ 在分块过程中逐步更新。

抽象伪代码如下：

\begin{lstlisting}[language=Python,caption={FlashAttention forward 的简化伪代码},label={lst:flashattention-forward-pseudocode}]
def flash_attention_forward(Q, K, V, block_q, block_k):
"""
Q, K, V: [seq_len, head_dim]
This is a mathematical sketch, not an efficient implementation.
"""
S, D = Q.shape

O = zeros([S, D])
m = full([S], -inf)
l = zeros([S])

for q_start in range(0, S, block_q):
    q_end = min(q_start + block_q, S)
    Q_blk = Q[q_start:q_end]

    O_blk = zeros([q_end - q_start, D])
    m_blk_running = full([q_end - q_start], -inf)
    l_blk_running = zeros([q_end - q_start])

    for k_start in range(0, S, block_k):
        k_end = min(k_start + block_k, S)
        K_blk = K[k_start:k_end]
        V_blk = V[k_start:k_end]

        scores = Q_blk @ K_blk.T / sqrt(D)

        # Apply causal mask if needed.
        scores = apply_causal_mask(
            scores,
            q_start=q_start,
            k_start=k_start,
        )

        m_local = row_max(scores)
        p = exp(scores - m_local[:, None])
        l_local = row_sum(p)
        acc_local = p @ V_blk

        m_new = maximum(m_blk_running, m_local)
        old_scale = exp(m_blk_running - m_new)
        new_scale = exp(m_local - m_new)

        O_blk = (
            old_scale[:, None] * O_blk
            + new_scale[:, None] * acc_local
        )
        l_blk_running = (
            old_scale * l_blk_running
            + new_scale * l_local
        )
        m_blk_running = m_new

    O_blk = O_blk / l_blk_running[:, None]

    O[q_start:q_end] = O_blk
    m[q_start:q_end] = m_blk_running
    l[q_start:q_end] = l_blk_running

return O, m, l

\end{lstlisting}

实际 kernel 不会像 Python 这样分配中间张量。它会把 block 内矩阵乘法映射到 GPU thread block、warp 和 Tensor Core MMA 上，尽量保持 Q/K/V block 的加载、计算和累积在片上完成。

\subsection{backward pass}
\label{subsec:flashattention-backward-pass}

训练时还需要 backward。标准 Attention backward 如果保存了 $\mat{P}$，可以直接使用它计算梯度。但 $\mat{P}$ 是 $S\times S$ 大矩阵，保存它会带来高显存和 IO。

FlashAttention 的 backward 通常不保存完整 $\mat{P}$，而是保存 forward 中每行的 softmax 统计量：

\[
m_i,\quad l_i.
\]

在 backward 中，对每个 block 重新计算：

\[
\mat{S}_{ij}=\mat{Q}_i\mat{K}*j^{\mathsf{T}},
\]
\[
\mat{P}*{ij}
\frac{\exp(\mat{S}_{ij}-m_i)}{l_i}.
\]

然后使用重算得到的 $\mat{P}_{ij}$ 计算：

\[
d\mat{V},
\quad
d\mat{Q},
\quad
d\mat{K}.
\]

这体现了重计算换 IO 的思想。Backward 的 FLOPs 增加了一些，但避免读取和保存巨大的 attention probability matrix。

Softmax backward 的核心公式为：

\[
dS_i
P_i
\odot
\left(
dP_i
-
\sum_j dP_{ij}P_{ij}
\right).
\]

在 Attention 中：

\[
O=PV.
\]

因此：

\[
dP=dO V^{\mathsf{T}},
\]
\[
dV=P^{\mathsf{T}}dO.
\]

再由：

\[
S=QK^{\mathsf{T}},
\]

得到：

\[
dQ=dS K,
\]
\[
dK=dS^{\mathsf{T}}Q.
\]

FlashAttention backward 通过分块重算 $P$，在片上完成这些局部梯度计算，并将 $dQ,dK,dV$ 写回或累积。

\begin{figure}[H]
\centering
\begin{tikzpicture}[
font=\small,
box/.style={draw, rounded corners=3pt, minimum width=1.9cm, minimum height=0.7cm, align=center},
arrow/.style={-{Latex[length=2.1mm]}, thick}
]

\node[font=\bfseries] at (0,3.25) {FlashAttention Backward：重算 Attention Block 而非保存完整 $P$};

\node[box, fill=blue!8] (saved) at (-5.0,1.4) {saved\\$m,l$};
\node[box, fill=green!10] (qkv) at (-5.0,0.35) {$Q,K,V$\\blocks};
\node[box, fill=orange!14] (recompute) at (-2.3,0.9) {recompute\\$S,P$ block};
\node[box, fill=purple!10] (grad) at (0.5,0.9) {local gradient\\$dQ,dK,dV$};
\node[box, fill=red!8] (acc) at (3.3,0.9) {accumulate\\write grads};

\draw[arrow] (saved) -- (recompute);
\draw[arrow] (qkv) -- (recompute);
\draw[arrow] (recompute) -- (grad);
\draw[arrow] (grad) -- (acc);

\node[align=left, text width=10.5cm] at (0,-1.2) {
  Forward 只保存每行 softmax 统计量 $m,l$。Backward 中分块重算 $S$ 和 $P$，
  用额外计算换取更低显存与更少 HBM 读取。
};

\end{tikzpicture}
\caption{FlashAttention backward 的重计算思想}
\label{fig:flashattention-backward-recompute}
\end{figure}

\subsection{FlashAttention-2 / FlashAttention-3 的优化方向}
\label{subsec:flashattention2-flashattention3}

FlashAttention 的第一版解决了 IO-aware exact attention 的核心问题，但仍然距离高度优化的 GEMM 有差距。后续版本主要围绕 GPU 并行度、work partitioning、硬件特性和低精度展开。

\textbf{FlashAttention-2} 的核心方向包括：

\begin{itemize}
\item 减少非矩阵乘法 FLOPs；
\item 改进 thread block 和 warp 之间的工作划分；
\item 对单个 head 的 attention 进行更好的并行切分；
\item 降低 shared memory 读写和跨 warp 通信；
\item 提高 occupancy 和 Tensor Core 利用率。
\end{itemize}

从系统角度看，FlashAttention-2 的启发是：减少 IO 只是第一步。一个 kernel 想接近 GEMM 性能，还必须让 GPU 有足够并行度，避免某些 warp 空闲，避免 shared memory 成为内部瓶颈。

\textbf{FlashAttention-3} 则进一步面向新硬件特性，尤其是 Hopper 架构中的异步数据搬运、Tensor Memory Accelerator、warp specialization 和 FP8 支持。其思想可以概括为：

\begin{itemize}
\item 让数据搬运和计算重叠；
\item 让不同 warp 承担不同角色，例如 producer/consumer；
\item 更细粒度地交错矩阵乘法和 softmax；
\item 利用 FP8 等低精度能力提升吞吐；
\item 在低精度下控制数值误差。
\end{itemize}

对 AI Infra 工程师来说，理解这些方向比记住某个版本的具体实现更重要。FlashAttention 的演进说明，高性能 attention kernel 不是单一算法公式，而是算法、数值稳定、GPU 存储层次、线程组织、低精度和硬件特性的共同设计。

\begin{table}[H]
\centering
\small
\caption{FlashAttention 系列的优化演进}
\label{tab:flashattention-evolution}
\begin{tabular}{p{0.22\textwidth}p{0.34\textwidth}p{0.34\textwidth}}
\toprule
\textbf{版本} & \textbf{核心问题} & \textbf{主要优化方向} \
\midrule
FlashAttention
& 标准 Attention 物化 $S\times S$ 中间矩阵，HBM IO 过高
& IO-aware tiling、online softmax、不落地 attention matrix。 \
\midrule
FlashAttention-2
& 第一版仍未充分接近 GEMM 效率
& 更好的 work partitioning、更高 occupancy、减少非 matmul FLOPs。 \
\midrule
FlashAttention-3
& 新硬件特性未充分利用
& Hopper 异步、warp specialization、TMA、FP8 与低精度数值控制。 \
\bottomrule
\end{tabular}
\end{table}

\section{工程实现视角}
\label{sec:flashattention-engineering-view}

FlashAttention 的论文公式展示了算法思想，但工程实现还要处理大量底层细节。一个高性能 attention kernel 需要同时满足：

\begin{itemize}
\item 数据从 HBM 读取次数少；
\item shared memory 容量足够；
\item register pressure 不过高；
\item Tensor Core 利用率高；
\item warp 分工合理；
\item occupancy 足够；
\item 支持 causal mask、padding mask、dropout；
\item 支持不同 head dimension；
\item 支持 MHA/MQA/GQA；
\item 支持 variable sequence length；
\item 支持训练 backward 和推理 decode/prefill；
\item 与 PyTorch、Triton、CUDA graph、serving engine 集成。
\end{itemize}

本节从几个工程维度展开。

\subsection{block size 选择}
\label{subsec:flashattention-block-size-choice}

FlashAttention 的 block size 决定每次搬入 SRAM 的 Q/K/V tile 大小。设 Q block 大小为 $B_Q$，K/V block 大小为 $B_K$。局部 score block 为：

\[
B_Q\times B_K.
\]

Block size 受以下约束：

\begin{itemize}
\item shared memory 容量；
\item register 数量；
\item head dimension $d_h$；
\item Tensor Core tile 形状；
\item warp 数量；
\item occupancy；
\item causal mask 边界；
\item sequence length；
\item dtype。
\end{itemize}

如果 block 太大：

\begin{itemize}
\item shared memory 不够；
\item register pressure 高；
\item occupancy 下降；
\item block 内同步和 reduction 变重。
\end{itemize}

如果 block 太小：

\begin{itemize}
\item 数据复用不足；
\item kernel launch 和 block 调度开销相对变高；
\item Tensor Core 利用率可能下降；
\item online softmax 更新次数增加。
\end{itemize}

因此 block size 不是越大越好，也不是越小越好。高性能库通常会根据 head dimension、dtype、GPU 架构和是否 causal 选择不同 kernel specialization。

\subsection{warp-level parallelism}
\label{subsec:warp-level-parallelism}

在 GPU 中，warp 是调度基本单位。FlashAttention kernel 需要将一个 attention tile 的工作分配给多个 warps。例如：

\begin{itemize}
\item 一些 warps 负责加载 Q/K/V；
\item 一些 warps 负责 MMA；
\item 一些 warps 负责 softmax reduction；
\item 一些 warps 负责输出累积和写回。
\end{itemize}

如果 warp 分工不好，会出现：

\begin{itemize}
\item 某些 warp 等待其他 warp；
\item shared memory 读写冲突；
\item reduction 开销过大；
\item occupancy 低；
\item Tensor Core pipeline 填不满。
\end{itemize}

FlashAttention-2 的一个重要方向就是更好的 work partitioning。特别是在序列长度较短、batch/head 数不足或单 head 工作量较大时，需要在一个 head 内部进一步并行，否则 GPU 可能吃不满。

\begin{figure}[H]
\centering
\begin{tikzpicture}[
font=\small,
warp/.style={draw, rounded corners=3pt, minimum width=1.4cm, minimum height=0.55cm, align=center},
smem/.style={draw, rounded corners=4pt, minimum width=4.2cm, minimum height=0.75cm, align=center, fill=orange!14},
arrow/.style={-{Latex[length=2.0mm]}, thick}
]

\node[font=\bfseries] at (0,3.15) {FlashAttention kernel 中的 warp-level 分工示意};

\foreach \i/\x/\task/\col in {
  0/-4.5/load Q/blue!8,
  1/-2.6/load KV/green!10,
  2/-0.7/MMA/purple!10,
  3/1.2/softmax/yellow!16,
  4/3.1/accumulate/red!8
} {
  \node[warp, fill=\col] (w\i) at (\x,1.55) {Warp \i\\\task};
}

\node[smem] (smem) at (-0.7,0.25) {Shared Memory / Registers};

\foreach \i in {0,1,2,3,4} {
  \draw[arrow] (w\i) -- (smem);
}

\node[align=left, text width=10.5cm] at (0,-1.25) {
  实际实现中 warp 分工会更加复杂。核心目标是让数据搬运、矩阵乘法、
  softmax reduction 和输出累积尽量流水化，减少等待和片上通信。
};

\end{tikzpicture}
\caption{FlashAttention 中 warp-level parallelism 的工程直觉}
\label{fig:flashattention-warp-level-parallelism}
\end{figure}

\subsection{shared memory 使用}
\label{subsec:flashattention-shared-memory-usage}

Shared memory 是 FlashAttention 的关键资源。它用于缓存 K/V blocks、部分 Q blocks、中间 score 或 softmax 相关数据。使用 shared memory 的目标是减少 HBM 访问，但 shared memory 容量有限，且会影响 occupancy。

需要关注：

\begin{itemize}
\item 每个 thread block 使用多少 shared memory；
\item 是否导致每个 SM 可驻留 block 数下降；
\item shared memory bank conflict；
\item 数据 layout 是否利于 vectorized load；
\item Q/K/V tile 是否复用充分；
\item 是否需要 double buffering。
\end{itemize}

Double buffering 是常见优化：当当前 tile 正在计算时，异步加载下一个 tile。这样可以重叠数据搬运和计算。较新的 GPU 架构提供更强的异步拷贝能力，使这类优化更重要。

\subsection{与 Triton kernel 的关系}
\label{subsec:flashattention-and-triton}

Triton 是一种用于编写高性能 GPU kernel 的语言和编译器，常用于实现矩阵乘法、softmax、attention 等算子。FlashAttention 可以用 CUDA C++ 实现，也可以用 Triton 实现。Triton 的优势是：

\begin{itemize}
\item Python 风格开发体验；
\item 更容易写出 tile-based kernel；
\item 自动处理部分底层细节；
\item 与 PyTorch 集成方便；
\item 适合快速迭代和定制 kernel。
\end{itemize}

但极致性能实现仍可能需要 CUDA C++、CUTLASS、内联 PTX 或硬件特定优化。尤其当需要利用最新架构的异步拷贝、TMA、warp specialization、FP8 Tensor Core 等特性时，底层实现复杂度会增加。

对工程团队而言，选择 Triton 还是 CUDA C++ 取决于：

\begin{itemize}
\item 性能目标；
\item 开发速度；
\item 目标 GPU 架构；
\item 需要支持的 dtype 和 head dimension；
\item 是否需要快速适配模型变体；
\item 团队对底层 CUDA 的掌握程度。
\end{itemize}

下面给出一个简化的 Triton 风格伪代码，用于展示 tiled attention kernel 的形状。它不是可直接运行的完整实现。

\begin{lstlisting}[language=Python,caption={Triton 风格 FlashAttention kernel 的结构示意},label={lst:triton-flashattention-sketch}]

# This is a structural sketch, not runnable Triton code.

@triton_jit
def flash_attention_kernel(Q, K, V, O, LSE, stride_info, BLOCK_M, BLOCK_N):
# Program id selects a Q block and a head.
q_block_id = program_id(0)
head_id = program_id(1)

# Load Q tile.
q = load_q_tile(Q, q_block_id, head_id, BLOCK_M)

# Initialize online softmax states.
m = full([BLOCK_M], -inf)
l = zeros([BLOCK_M])
acc = zeros([BLOCK_M, HEAD_DIM])

# Loop over K/V blocks.
for kv_block_id in range(0, num_kv_blocks):
    k = load_k_tile(K, kv_block_id, head_id, BLOCK_N)
    v = load_v_tile(V, kv_block_id, head_id, BLOCK_N)

    scores = dot(q, transpose(k)) * scale
    scores = apply_mask_if_needed(scores, q_block_id, kv_block_id)

    m_new = maximum(m, max(scores, axis=1))
    p = exp(scores - m_new[:, None])

    alpha = exp(m - m_new)
    acc = acc * alpha[:, None] + dot(p, v)
    l = l * alpha + sum(p, axis=1)
    m = m_new

# Normalize and write output.
acc = acc / l[:, None]
store_o_tile(O, q_block_id, head_id, acc)

# Save log-sum-exp or m/l for backward if training.
store_lse(LSE, q_block_id, head_id, m + log(l))

\end{lstlisting}

这个伪代码保留了 FlashAttention 的核心结构：tile、循环、online softmax、累积输出、不保存 $S\times S$ 矩阵。

\section{FlashAttention 与推理系统}
\label{sec:flashattention-and-serving-systems}

FlashAttention 最初常被训练和长序列 attention 场景关注，但它同样影响推理系统。推理中有 prefill 和 decode 两种阶段，它们对 attention kernel 的需求不同。

\subsection{Prefill 中的 FlashAttention}
\label{subsec:flashattention-in-prefill}

Prefill 阶段一次处理 prompt 的全部 token。对于长 prompt，attention 仍然是 $S\times S$ 结构，因此 FlashAttention 非常有价值。它可以：

\begin{itemize}
\item 降低 prefill 显存峰值；
\item 减少 attention 中间矩阵 HBM 读写；
\item 提升长 prompt 的 TTFT；
\item 支持更长上下文；
\item 与 causal mask 结合。
\end{itemize}

在 serving 中，TTFT 通常由排队时间和 prefill 时间共同决定。长 prompt 请求尤其依赖高效 prefill attention。FlashAttention 可以显著降低 prefill kernel 的 IO 开销，但它不能消除 $O(S^2)$ 的数学计算，因此超长 prompt 仍然昂贵。对于超长上下文，还需要 chunked prefill、prefix cache、sparse attention 或 context parallel 等技术配合。

\subsection{Decode 中的 FlashAttention}
\label{subsec:flashattention-in-decode}

Decode 阶段每个请求通常只处理一个新 token。此时 Q 的 sequence length 为 $1$，K/V length 为历史上下文长度 $S$。Attention 形状是：

\[
Q: 1\times d_h,
\quad
K,V: S\times d_h.
\]

这与 prefill 中的 $S\times S$ attention 不同。Decode attention 不会物化完整 $S\times S$ matrix，因为只有一行 score：

\[
1\times S.
\]

因此 decode 的瓶颈更多是 KV Cache 读取和 batch 组织，而不是 $S^2$ 中间矩阵物化。FlashAttention 的 online softmax 和 fused kernel 思想仍然有帮助，但 decode kernel 通常需要针对 single-query 或 small-query 情况专门优化。

对于 batch decode，多个请求可以组成：

\[
B_{\mathrm{decode}}\times 1
\]

的 query batch，每个请求上下文长度不同。若配合 PagedAttention，K/V 还可能是分页存储。此时 attention kernel 需要同时处理：

\begin{itemize}
\item variable context length；
\item paged KV block table；
\item GQA/MQA；
\item small Q length；
\item memory-bound KV reads；
\item sampling 后的 dynamic scheduling。
\end{itemize}

因此，serving 系统中常常同时需要 FlashAttention 风格的 prefill kernel 和 PagedAttention 风格的 decode kernel。二者解决的问题不同：

\begin{itemize}
\item FlashAttention：减少 attention 计算中的 HBM IO；
\item PagedAttention：管理动态 KV Cache 并支持 continuous batching。
\end{itemize}

\begin{table}[H]
\centering
\small
\caption{Prefill 与 Decode 对 Attention Kernel 的不同要求}
\label{tab:prefill-decode-attention-kernel}
\begin{tabular}{p{0.2\textwidth}p{0.34\textwidth}p{0.34\textwidth}}
\toprule
\textbf{维度} & \textbf{Prefill} & \textbf{Decode} \
\midrule
Q 长度
& prompt 长度 $S$
& 通常为 1 \
\midrule
K/V 长度
& prompt 长度 $S$
& 历史上下文长度 $S_{\mathrm{cache}}$ \
\midrule
主要问题
& $S\times S$ attention IO 和计算
& KV Cache 读取、动态 batch、memory-bound \
\midrule
FlashAttention 价值
& 非常高，避免物化 attention matrix
& 思想有用，但通常需要 decode-specialized kernel \
\midrule
PagedAttention 价值
& 管理写入和长 prompt blocks
& 非常高，支持动态 KV 读取与 continuous batching \
\bottomrule
\end{tabular}
\end{table}

\subsection{与 PagedAttention 的互补关系}
\label{subsec:flashattention-pagedattention-complementary}

PagedAttention 和 FlashAttention 很容易被混淆，因为它们都与 Attention 和显存有关。但二者关注点不同。

\[
\text{FlashAttention}
\Rightarrow
\text{How to compute attention with less HBM IO?}
\]

\[
\text{PagedAttention}
\Rightarrow
\text{How to store and manage KV Cache for dynamic requests?}
\]

在一个现代 LLM serving engine 中，它们可能同时存在：

\begin{itemize}
\item prefill 使用 FlashAttention 或其变体；
\item decode 使用 PagedAttention kernel 读取分页 KV Cache；
\item 长 prompt chunked prefill 中也可能结合 paged KV 写入；
\item prefix cache 使用 block table 共享 KV；
\item attention backend 根据阶段、shape、GPU 架构选择不同 kernel。
\end{itemize}

可以把关系理解为：

\[
\text{FlashAttention 是计算优化，PagedAttention 是内存管理与调度优化。}
\]

\begin{figure}[H]
\centering
\begin{tikzpicture}[
font=\small,
box/.style={draw, rounded corners=4pt, minimum width=2.35cm, minimum height=0.8cm, align=center},
arrow/.style={-{Latex[length=2.1mm]}, thick}
]

\node[font=\bfseries] at (0,3.25) {FlashAttention 与 PagedAttention 在 Serving 中的互补};

\node[box, fill=green!10] (prefill) at (-4.2,1.35) {Prefill\\long prompt};
\node[box, fill=orange!14] (flash) at (-1.35,1.35) {FlashAttention\\IO-aware compute};
\node[box, fill=blue!8] (kvwrite) at (1.45,1.35) {Write KV\\blocks};
\node[box, fill=purple!10] (decode) at (4.2,1.35) {Decode\\token loop};

\node[box, fill=yellow!16] (paged) at (1.45,0.0) {PagedAttention\\paged KV read};
\node[box, fill=red!8] (sched) at (4.2,0.0) {Continuous\\batching};

\draw[arrow] (prefill) -- (flash);
\draw[arrow] (flash) -- (kvwrite);
\draw[arrow] (kvwrite) -- (decode);
\draw[arrow] (paged) -- (decode);
\draw[arrow] (sched) -- (decode);
\draw[arrow] (kvwrite) -- (paged);

\node[align=left, text width=10.5cm] at (0,-1.55) {
  Prefill 阶段常用 FlashAttention 减少 $S\times S$ attention 的 IO；
  Decode 阶段常用 PagedAttention 管理和读取动态 KV Cache。
};

\end{tikzpicture}
\caption{FlashAttention 与 PagedAttention 在推理系统中的互补关系}
\label{fig:flashattention-pagedattention-complementary}
\end{figure}

\section{常见误区与调试方法}
\label{sec:flashattention-misconceptions-debugging}

\subsection{误区一：FlashAttention 是近似 Attention}
\label{subsec:misconception-flashattention-approximate}

FlashAttention 不是近似 attention。它计算的是与标准 attention 数学上等价的精确结果，只是改变计算顺序和内存访问方式。由于浮点数加法和指数运算顺序变化，结果可能与标准实现存在极小数值差异，但这不是近似算法意义上的误差。

这点很重要。FlashAttention 不像某些 sparse attention、linear attention 或 low-rank attention 那样改变模型的 attention pattern。它的目标是在不改变模型语义的前提下提高速度和降低显存。

\subsection{误区二：FlashAttention 降低了 $O(S^2)$ 计算复杂度}
\label{subsec:misconception-flashattention-compute-complexity}

FlashAttention 不改变 dense attention 的 $O(S^2d_h)$ 计算复杂度。对于 full attention，所有 query-key 对仍然需要计算。它降低的是 HBM IO 和中间显存，不是数学计算阶数。

因此，当序列极长时，FlashAttention 仍然要面对 $S^2$ 计算量。若需要从复杂度上降低长上下文成本，需要考虑 sparse attention、sliding window attention、state space model、linear attention、attention sink、KV compression 等其他方向。

\subsection{误区三：所有场景都一定更快}
\label{subsec:misconception-flashattention-always-faster}

FlashAttention 在许多长序列和训练/prefill 场景中非常有效，但并不意味着所有 shape 都一定更快。影响因素包括：

\begin{itemize}
\item sequence length；
\item head dimension；
\item batch size；
\item causal 或 non-causal；
\item GPU 架构；
\item dtype；
\item dropout 是否开启；
\item mask 类型；
\item 是否 variable length；
\item kernel 是否支持当前 shape；
\item 是否 decode single-query。
\end{itemize}

对于很短序列，标准实现或其他 fused attention 可能已经足够快；对于某些特殊 mask 或 layout，FlashAttention kernel 可能无法直接使用；对于 decode 阶段，专门的 paged decode kernel 可能比通用 FlashAttention 更合适。

\subsection{调试与验证}
\label{subsec:flashattention-debug-validation}

接入 FlashAttention 或替换 attention backend 时，需要验证：

\begin{itemize}
\item 输出数值与标准实现是否在 tolerance 内一致；
\item backward 梯度是否正确；
\item causal mask 是否正确；
\item padding mask 或 variable length 是否正确；
\item dropout 随机性是否符合训练要求；
\item mixed precision 是否稳定；
\item 长序列下是否 OOM；
\item profiler 中 HBM traffic 是否下降；
\item kernel 是否真正被调用，而不是 fallback；
\item 不同 GPU 架构是否走了合适 kernel。
\end{itemize}

一个最小数值验证可以写成：

\begin{lstlisting}[language=Python,caption={FlashAttention 与标准 Attention 的数值验证示意},label={lst:flashattention-correctness-check}]
import torch
import torch.nn.functional as F

def standard_attention(q, k, v, causal=True):
# q, k, v: [batch, heads, seq, dim]
scale = q.shape[-1] ** -0.5
scores = torch.matmul(q, k.transpose(-1, -2)) * scale

if causal:
    seq = q.shape[-2]
    mask = torch.triu(
        torch.ones(seq, seq, device=q.device, dtype=torch.bool),
        diagonal=1,
    )
    scores = scores.masked_fill(mask, float("-inf"))

prob = torch.softmax(scores, dim=-1)
return torch.matmul(prob, v)

def check_attention_impl(q, k, v, flash_attention_fn):
ref = standard_attention(q, k, v, causal=True)
out = flash_attention_fn(q, k, v, causal=True)

max_abs_err = (ref - out).abs().max().item()
max_rel_err = ((ref - out).abs() / ref.abs().clamp_min(1e-6)).max().item()

print("max abs err:", max_abs_err)
print("max rel err:", max_rel_err)

torch.testing.assert_close(out, ref, rtol=1e-2, atol=1e-2)

\end{lstlisting}

在 BF16/FP16 下，不应要求逐 bit 一致。应使用合理的 absolute/relative tolerance，并分别验证不同 sequence length、head dimension、causal mask、batch size 和 dtype。

\section{本章小结}
\label{sec:chapter14-summary}

本章系统介绍了 FlashAttention 的动机、online softmax 推导、核心算法、前反向计算和工程实现视角。

\begin{itemize}
\item \textbf{标准 Attention 的主要问题之一是 IO 瓶颈}：它常常显式物化 $S\times S$ score matrix 和 softmax probability matrix，导致大量 HBM 读写。
\item \textbf{FlashAttention 不改变 Attention 数学结果}，而是通过 IO-aware tiling 和 kernel fusion 减少 HBM traffic。
\item \textbf{FlashAttention 不降低 dense attention 的 FLOPs 阶数}，仍然是 $O(S^2d_h)$，但显著降低中间显存和 HBM 访问。
\item \textbf{HBM 与片上 SRAM/register 的差异是理解 FlashAttention 的关键}：把中间 attention block 留在片上计算，通常比写入 HBM 再读回更快。
\item \textbf{Online softmax 使分块 attention 成为可能}：通过维护每行的运行最大值 $m$、指数和 $l$ 和输出累积 $a$，可以分块得到与完整 softmax 等价的结果。
\item \textbf{FlashAttention forward 按 Q/K/V blocks 计算}，局部 score、mask、softmax 和 $PV$ 累积都在 SRAM/register 中完成，最终只写回输出和少量统计量。
\item \textbf{FlashAttention backward 使用重计算换 IO}：不保存完整 attention probability，而是保存 $m,l$，在 backward 中分块重算需要的 softmax block。
\item \textbf{FlashAttention-2 进一步优化 work partitioning}，提升 occupancy，减少非矩阵乘法 FLOPs 和 shared memory 开销。
\item \textbf{FlashAttention-3 面向新硬件特性}，利用异步数据搬运、warp specialization、TMA 和 FP8 等能力继续提升性能。
\item \textbf{工程实现需要综合考虑 block size、shared memory、register pressure、warp 分工、Tensor Core 利用率和 mask 支持}。
\item \textbf{在推理系统中，FlashAttention 与 PagedAttention 是互补的}：前者主要优化 attention 计算 IO，后者主要优化动态 KV Cache 管理和 continuous batching。
\item \textbf{FlashAttention 不是所有 shape 都必然更快}。短序列、特殊 mask、single-query decode 或不支持的 layout 可能需要其他 attention backend。
\end{itemize}

下一章将进入量化技术。FlashAttention 通过减少 IO 提升 attention kernel 效率，而量化则从数据表示本身入手，降低模型权重、激活和 KV Cache 的字节数。我们会系统讨论 PTQ、QAT、GPTQ、AWQ、weight-only quantization、per-channel/per-group scale，以及量化对显存、带宽、吞吐和精度的影响。

\chapter{量化技术}
\label{chap:quantization}

\section{为什么推理需要量化}
\label{sec:why-inference-needs-quantization}

前面两章分别讨论了 PagedAttention 与 FlashAttention。PagedAttention 解决 KV Cache 的动态内存管理问题，FlashAttention 解决 Attention 计算中的 HBM IO 瓶颈。本章进入推理优化的另一个核心方向：\textbf{量化 Quantization}。

量化的基本思想是：用更低 bit-width 的数值格式表示模型权重、激活或 KV Cache，从而降低显存占用、减少内存带宽压力，并在硬件支持时提升计算吞吐。

在大模型推理中，量化的价值尤其明显。以一个 dense Decoder-only 模型为例，如果参数量为 $N$，使用 FP16/BF16 存储权重，则权重显存约为：

[
M_{\mathrm{weight,fp16}} = 2N\ \mathrm{bytes}.
]

如果使用 INT8 权重量化：

[
M_{\mathrm{weight,int8}} \approx 1N\ \mathrm{bytes}.
]

如果使用 INT4 权重量化：

[
M_{\mathrm{weight,int4}} \approx 0.5N\ \mathrm{bytes}.
]

忽略 scale、zero point、alignment 和 packing metadata，仅从权重本体看，INT8 相比 FP16 约节省一半显存，INT4 约节省四分之三显存。对于 70B 参数模型，FP16/BF16 权重大约需要：

[
70\times 10^9\times 2
=====================

140\ \mathrm{GB}.
]

如果权重压到 INT4，理论权重本体约为：

[
70\times 10^9\times 0.5
=======================

35\ \mathrm{GB}.
]

这会直接改变部署形态：原本需要多张 GPU 才能装下的模型，可能可以用更少 GPU 部署；原本显存主要被权重占满的服务，可以留出更多空间给 KV Cache，从而支持更高并发或更长上下文。

\subsection{权重显存}
\label{subsec:weight-memory}

推理阶段不需要保存梯度和优化器状态，显存主要由以下部分组成：

[
M_{\mathrm{inference}}
======================

M_{\mathrm{weights}}
+
M_{\mathrm{KV}}
+
M_{\mathrm{activations}}
+
M_{\mathrm{workspace}}
+
M_{\mathrm{runtime}}.
]

其中权重显存通常是常驻显存，KV Cache 随请求增长，activation 和 workspace 随 batch、sequence length 和 kernel 实现变化。

在 batch 较小、上下文不太长的场景中，权重可能是主要显存占用；在长上下文和高并发 serving 中，KV Cache 可能成为主要显存占用。量化可以作用于不同对象：

\begin{itemize}
\item \textbf{Weight Quantization}：压缩模型权重，例如 W8A16、W4A16；
\item \textbf{Activation Quantization}：压缩中间激活，例如 W8A8；
\item \textbf{KV Cache Quantization}：压缩推理缓存，例如 FP8 KV、INT8 KV；
\item \textbf{Optimizer Quantization}：训练中压缩优化器状态，本章不重点展开。
\end{itemize}

其中推理部署最常见的是 weight-only quantization，即权重量化到 INT8/INT4，激活仍用 FP16/BF16。它实现难度相对较低，精度风险也比同时量化激活更可控。

\begin{table}[H]
\centering
\small
\caption{常见推理量化对象}
\label{tab:quantization-targets}
\begin{tabular}{p{0.22\textwidth}p{0.28\textwidth}p{0.38\textwidth}}
\toprule
\textbf{量化对象} & \textbf{常见格式} & \textbf{主要收益与风险} \
\midrule
权重
& INT8, INT4, NF4, FP8
& 降低常驻显存和权重读取带宽；风险主要是模型精度下降。 \
\midrule
激活
& INT8, FP8
& 降低 GEMM 输入带宽并启用低精度算力；风险是 activation outlier 造成误差放大。 \
\midrule
KV Cache
& INT8, FP8, INT4
& 降低长上下文和高并发的 cache 显存；风险是 decode 质量和长程依赖下降。 \
\midrule
Logits / Sampling
& 通常保持 FP16/FP32
& 采样对数值稳定敏感，通常不作为主要量化对象。 \
\bottomrule
\end{tabular}
\end{table}

\subsection{显存带宽}
\label{subsec:memory-bandwidth}

量化不仅减少显存容量，也减少显存带宽压力。推理中的 Linear 层通常形如：

[
\mat{Y}=\mat{X}\mat{W}.
]

在 decode 阶段，batch 维和 token 维可能很小，权重读取成本会非常重要。每生成一个 token，模型需要经过所有层，每层读取 QKV projection、output projection、FFN up/gate/down projection 等权重。若权重以 FP16 存储，读取字节数大约是 INT8 的 2 倍、INT4 的 4 倍。

当 kernel 是 memory-bound 时，减少权重读取字节数可以直接提升速度。设某层需要读取权重大小 $M_W$，有效 HBM 带宽为 $\beta$，则权重读取时间下界为：

[
T_{\mathrm{read}}
\ge
\frac{M_W}{\beta}.
]

如果量化后权重字节数降低为原来的 $\rho$，则读取时间下界也降低为：

[
T_{\mathrm{read,quant}}
\ge
\frac{\rho M_W}{\beta}.
]

其中：

[
\rho_{\mathrm{INT8}}\approx \frac{1}{2},
\quad
\rho_{\mathrm{INT4}}\approx \frac{1}{4}.
]

当然，实际速度不一定按这个比例提升，因为还要考虑 dequantization、packing/unpacking、Tensor Core 支持、kernel fusion 和 batch size。量化后的 GEMM 如果需要先把 INT4 解包并反量化到 FP16，再执行普通 GEMM，可能会被解包开销抵消。高性能量化推理必须让 dequantization 与 GEMM 融合。

\subsection{compute-bound 与 memory-bound}
\label{subsec:compute-bound-memory-bound-quantization}

量化是否能加速，取决于当前瓶颈。若推理主要受权重读取或 KV Cache 读取限制，量化往往收益明显；若推理主要受计算 FLOPs 限制，则需要硬件真正支持低精度计算，才能获得吞吐提升。

可以用 roofline 直觉理解。算术强度为：

[
I=\frac{\mathrm{FLOPs}}{\mathrm{Bytes}}.
]

量化减少 bytes，因此提高 $I$。如果原 kernel 处于 memory-bound 区域，提高算术强度可以提升性能；如果原 kernel 已经 compute-bound，仅减少 bytes 未必能提升很多，除非低精度格式也提高 compute throughput。

\begin{figure}[H]
\centering
\begin{tikzpicture}[
font=\small,
axis/.style={-{Latex[length=2.2mm]}, thick},
curve/.style={thick},
point/.style={circle, draw, minimum size=0.16cm}
]

```
\node[font=\bfseries] at (0,3.0) {量化通过减少 Bytes 提高算术强度};

\draw[axis] (-5.1,-1.15) -- (5.2,-1.15) node[right] {Arithmetic Intensity};
\draw[axis] (-5.1,-1.15) -- (-5.1,2.25) node[above] {Performance};

\draw[curve, orange!80!black] (-4.8,-0.95) -- (-1.4,1.55);
\draw[curve, orange!80!black] (-1.4,1.55) -- (4.8,1.55);

\node[point, fill=red!20] (fp16) at (-3.2,0.15) {};
\node[anchor=west] at (-3.0,0.15) {FP16 weights};

\node[point, fill=green!20] (int4) at (-0.8,1.05) {};
\node[anchor=west] at (-0.6,1.05) {INT4 weights};

\draw[->, thick, blue!70] (fp16) -- (int4);

\node[align=left, text width=10.5cm] at (0,-2.2) {
  对 memory-bound 推理，量化减少 HBM 读取字节数，使 kernel 向更高算术强度移动。
  但若 dequantization 开销过高，或 kernel 无法利用低精度算力，速度收益会下降。
};
```

\end{tikzpicture}
\caption{从 roofline 视角理解量化的性能收益}
\label{fig:quantization-roofline}
\end{figure}

\subsection{量化与前几章优化的关系}
\label{subsec:quantization-relation-to-previous-chapters}

量化不是孤立优化。它与前几章技术相互影响：

\begin{itemize}
\item 与 \textbf{PagedAttention} 的关系：KV Cache 量化可以进一步降低 paged block 的显存占用，提高可并发请求数；
\item 与 \textbf{FlashAttention} 的关系：FlashAttention 降低 attention 中间 IO，量化降低输入/权重/KV 字节数，两者可互补；
\item 与 \textbf{Tensor Parallel} 的关系：TP 下权重和 KV head 已经切分，量化进一步降低每 GPU shard 的显存和带宽；
\item 与 \textbf{Serving Scheduler} 的关系：量化降低单请求资源占用，改变 admission control、batch size 和成本模型；
\item 与 \textbf{模型架构} 的关系：RMSNorm、GQA、SwiGLU、MoE、RoPE 等结构会影响量化敏感性。
\end{itemize}

因此，量化不能只看模型文件大小，还要看端到端 serving 指标：

[
\text{TTFT},\quad
\text{TPOT},\quad
\text{tokens/s},\quad
\text{P99 latency},\quad
\text{quality},\quad
\text{cost/token}.
]

\section{数值表示与量化基础}
\label{sec:numerical-representation-and-quantization-basics}

量化的本质是用有限个离散值近似连续或高精度值。对深度学习而言，量化通常包含三个步骤：

\begin{enumerate}
\item 选择低精度表示范围；
\item 将高精度 tensor 映射到低精度整数或浮点值；
\item 在计算时直接使用低精度算子，或反量化回高精度参与计算。
\end{enumerate}

\subsection{FP32、FP16、BF16、FP8、INT8、INT4}
\label{subsec:number-formats}

常见数值格式可以分为浮点和整数两类。

\textbf{FP32} 是传统训练和高精度计算格式，精度和动态范围都较高，但显存和带宽开销大。

\textbf{FP16} 使用 16 bit 浮点，显存减半，Tensor Core 支持好，但指数位较少，动态范围小，训练中可能需要 loss scaling。

\textbf{BF16} 也是 16 bit 浮点，它保留与 FP32 类似的指数位宽，因此动态范围较大，但尾数位较少。大模型训练和推理中 BF16 非常常见。

\textbf{FP8} 是 8 bit 浮点格式，常见有 E4M3、E5M2 等变体。FP8 的优点是比 INT8 更像浮点，有指数动态范围；缺点是硬件和软件栈要求更高，并且需要 scale 管理。

\textbf{INT8} 是 8 bit 整数量化格式，范围通常为：

[
[-128,127]
]

或：

[
[0,255].
]

\textbf{INT4} 是 4 bit 整数量化格式，范围通常为：

[
[-8,7]
]

或：

[
[0,15].
]

INT4 的压缩率高，但误差更大，通常需要 per-group scale、精心校准和高质量量化算法。

\begin{table}[H]
\centering
\small
\caption{常见数值格式对比}
\label{tab:numeric-format-comparison-quant}
\begin{tabular}{p{0.16\textwidth}p{0.16\textwidth}p{0.24\textwidth}p{0.32\textwidth}}
\toprule
\textbf{格式} & \textbf{bit 数} & \textbf{典型用途} & \textbf{特点} \
\midrule
FP32
& 32
& 高精度训练、累积
& 动态范围和精度高，显存/带宽开销大。 \
\midrule
FP16
& 16
& 训练/推理
& Tensor Core 支持好，但动态范围较小。 \
\midrule
BF16
& 16
& 大模型训练/推理
& 动态范围接近 FP32，尾数精度低于 FP16。 \
\midrule
FP8
& 8
& 新硬件上的训练/推理
& 有指数动态范围，需要 scale 与 amax 管理。 \
\midrule
INT8
& 8
& W8A8、weight-only
& 整数表示，硬件支持较广，校准较重要。 \
\midrule
INT4
& 4
& weight-only LLM 推理
& 压缩率高，需要 group scale 和专用 kernel。 \
\bottomrule
\end{tabular}
\end{table}

\subsection{uniform quantization}
\label{subsec:uniform-quantization}

最基础的是 uniform quantization，即用等间隔离散值近似连续值。给定实数 $x$，scale 为 $s$，zero point 为 $z$，量化为整数 $q$：

[
q
=

\operatorname{clip}
\left(
\left\lfloor \frac{x}{s}\right\rceil + z,
q_{\min},
q_{\max}
\right).
]

反量化为：

[
\hat{x}
=======

s(q-z).
]

其中：

\begin{itemize}
\item $s$ 控制量化步长；
\item $z$ 控制实数 0 对应的整数位置；
\item $q_{\min},q_{\max}$ 由 bit-width 决定；
\item $\lfloor\cdot\rceil$ 表示四舍五入。
\end{itemize}

量化误差为：

[
e=x-\hat{x}.
]

Uniform quantization 的优点是简单、硬件友好；缺点是当 tensor 分布有 outlier 时，scale 被 outlier 拉大，大部分普通值的量化分辨率变差。

\subsection{symmetric 与 asymmetric quantization}
\label{subsec:symmetric-asymmetric-quantization}

\textbf{Symmetric quantization} 使用零点 $z=0$，量化区间关于 0 对称。例如 INT8 symmetric 范围为：

[
q\in[-127,127]
]

或：

[
q\in[-128,127].
]

Scale 可取：

[
s
=

\frac{\max(|x_{\min}|,|x_{\max}|)}{q_{\max}}.
]

量化为：

[
q
=

\operatorname{clip}
\left(
\left\lfloor \frac{x}{s}\right\rceil,
q_{\min},
q_{\max}
\right).
]

Symmetric quantization 的优点是计算简单，尤其适合权重，因为权重分布通常近似以 0 为中心。

\textbf{Asymmetric quantization} 使用非零 zero point，使实数范围 $[x_{\min},x_{\max}]$ 映射到整数范围 $[q_{\min},q_{\max}]$：

[
s=
\frac{x_{\max}-x_{\min}}
{q_{\max}-q_{\min}},
]

[
z=
q_{\min}
--------

\left\lfloor
\frac{x_{\min}}{s}
\right\rceil.
]

Asymmetric quantization 更适合非对称分布，例如 ReLU 后激活。但它在矩阵乘法中需要处理 zero point 修正，kernel 更复杂。

\begin{figure}[H]
\centering
\begin{tikzpicture}[
font=\small,
axis/.style={-{Latex[length=2.1mm]}, thick},
tick/.style={thick},
qpt/.style={circle, draw, minimum size=0.18cm, fill=blue!12}
]

```
\node[font=\bfseries] at (0,3.0) {Symmetric 与 Asymmetric Quantization};

% symmetric
\node[font=\bfseries] at (-3.3,2.25) {Symmetric};
\draw[axis] (-5.4,1.3) -- (-1.2,1.3);
\foreach \x/\lab in {-5.0/$-a$,-3.3/$0$,-1.6/$a$} {
  \draw[tick] (\x,1.18) -- (\x,1.42);
  \node[font=\scriptsize] at (\x,0.92) {\lab};
}
\foreach \x in {-5.0,-4.55,-4.1,-3.65,-3.3,-2.95,-2.5,-2.05,-1.6} {
  \node[qpt] at (\x,1.3) {};
}

% asymmetric
\node[font=\bfseries] at (3.0,2.25) {Asymmetric};
\draw[axis] (0.8,1.3) -- (5.4,1.3);
\foreach \x/\lab in {1.2/$x_{\min}$,2.1/$0$,5.0/$x_{\max}$} {
  \draw[tick] (\x,1.18) -- (\x,1.42);
  \node[font=\scriptsize] at (\x,0.92) {\lab};
}
\foreach \x in {1.2,1.68,2.16,2.64,3.12,3.60,4.08,4.56,5.04} {
  \node[qpt, fill=green!12] at (\x,1.3) {};
}

\node[align=left, text width=10.5cm] at (0,-0.35) {
  symmetric quantization 以 0 为中心，权重中常用；
  asymmetric quantization 可以更好覆盖非对称激活分布，但 zero point 会增加 kernel 复杂度。
};
```

\end{tikzpicture}
\caption{对称与非对称量化的区间映射}
\label{fig:symmetric-asymmetric-quantization}
\end{figure}

\subsection{per-tensor、per-channel、per-group scale}
\label{subsec:per-tensor-channel-group-scale}

Scale 的粒度决定了量化误差和 metadata 开销。

\textbf{Per-tensor quantization} 对整个 tensor 使用一个 scale：

[
s_{\mathrm{tensor}}.
]

优点是 metadata 少、实现简单；缺点是如果 tensor 不同通道分布差异大，一个 scale 难以兼顾所有通道。

\textbf{Per-channel quantization} 对每个输出通道或输入通道使用独立 scale。例如 Linear 权重：

[
\mat{W}\in\mathbb{R}^{H_{\mathrm{out}}\times H_{\mathrm{in}}}.
]

若按 output channel 量化：

[
s_i,\quad i=0,\ldots,H_{\mathrm{out}}-1.
]

第 $i$ 行权重使用 scale $s_i$。这种方式通常显著降低误差。

\textbf{Per-group quantization} 将通道或列划分为小组，每组一个 scale。例如 group size 为 $G$，则每 $G$ 个元素共享一个 scale：

[
s_{i,g}.
]

INT4 LLM weight-only quantization 中 per-group scale 很常见。它在精度和 metadata 之间折中：group 越小，量化更精细，但 scale metadata 越多，kernel 更复杂；group 越大，metadata 少，但误差增大。

[
\text{per-tensor}
\rightarrow
\text{per-channel}
\rightarrow
\text{per-group}
]

不是简单单调关系。实际选择取决于 weight layout、kernel 支持和精度目标。

\begin{table}[H]
\centering
\small
\caption{Scale 粒度对比}
\label{tab:scale-granularity}
\begin{tabular}{p{0.22\textwidth}p{0.30\textwidth}p{0.36\textwidth}}
\toprule
\textbf{粒度} & \textbf{优点} & \textbf{缺点} \
\midrule
Per-tensor
& metadata 最少，实现简单
& outlier 影响全 tensor，误差可能较大。 \
\midrule
Per-channel
& 每个通道适配自身范围，精度较好
& scale 数量增加，kernel 需支持通道 scale。 \
\midrule
Per-group
& INT4 常用，精度和 metadata 折中
& group size 需要调参，解包和 scale 应用更复杂。 \
\bottomrule
\end{tabular}
\end{table}

\subsection{dequantization}
\label{subsec:dequantization}

低精度权重可以有两种计算方式。

第一种是\textbf{反量化后计算}：

[
\hat{\mat{W}}=s(\mat{Q}_W-z),
]
[
\mat{Y}=\mat{X}\hat{\mat{W}}.
]

这种方式实现简单，但如果先把整块权重反量化到 FP16/BF16 再 GEMM，会失去显存带宽收益，并增加临时显存。

第二种是\textbf{融合反量化与 GEMM}。Kernel 从 HBM 读取 packed INT4/INT8 权重，在寄存器或片上存储中解包并乘 scale，然后参与矩阵乘法：

[
\mat{Y}
=======

\mat{X}
\left[
s(\mat{Q}_W-z)
\right].
]

高性能 weight-only 推理必须使用这种方式。也就是说，量化模型的速度不仅取决于量化算法，还取决于是否有高质量 kernel。

\begin{figure}[H]
\centering
\begin{tikzpicture}[
font=\small,
box/.style={draw, rounded corners=3pt, minimum width=2.05cm, minimum height=0.7cm, align=center},
arrow/.style={-{Latex[length=2.2mm]}, thick}
]

```
\node[font=\bfseries] at (0,3.1) {反量化路径：离线展开 vs fused dequant GEMM};

\node[box, fill=blue!8] (q1) at (-5.0,1.55) {INT4 packed\\weights};
\node[box, fill=orange!14] (dq1) at (-2.6,1.55) {dequant to\\FP16 tensor};
\node[box, fill=green!10] (gemm1) at (-0.2,1.55) {FP16 GEMM};
\draw[arrow] (q1) -- (dq1);
\draw[arrow] (dq1) -- (gemm1);

\node[box, fill=blue!8] (q2) at (-5.0,0.2) {INT4 packed\\weights};
\node[box, fill=purple!10, minimum width=3.0cm] (fused) at (-1.9,0.2) {fused unpack + scale + GEMM};
\node[box, fill=green!10] (out) at (1.3,0.2) {output};
\draw[arrow] (q2) -- (fused);
\draw[arrow] (fused) -- (out);

\node[align=left, text width=10.5cm] at (0,-1.35) {
  若先反量化完整 FP16 权重，再执行普通 GEMM，会失去许多带宽收益。
  高性能量化推理通常把 unpack、scale 和 GEMM 融合到一个 kernel 中。
};
```

\end{tikzpicture}
\caption{量化权重的两种计算路径}
\label{fig:dequantization-paths}
\end{figure}

\section{PTQ 与 QAT}
\label{sec:ptq-and-qat}

量化方法可以从训练介入程度分为两大类：

\begin{itemize}
\item \textbf{PTQ, Post-Training Quantization}：训练完成后量化，不重新训练或只做少量校准；
\item \textbf{QAT, Quantization-Aware Training}：训练或微调时模拟量化误差，让模型适应低精度。
\end{itemize}

LLM 推理部署中，PTQ 非常流行，因为重新训练大模型成本很高。GPTQ、AWQ、SmoothQuant 都属于后训练量化或校准型方法。QAT 精度潜力更高，但成本也更高，且对训练数据、训练 recipe 和工程流程要求更强。

\subsection{Post-Training Quantization}
\label{subsec:post-training-quantization}

PTQ 的基本流程是：

\begin{enumerate}
\item 加载训练好的高精度模型；
\item 准备少量校准数据；
\item 收集权重和激活统计；
\item 选择 scale、zero point、group size、clipping 策略；
\item 量化权重或激活；
\item 评估 perplexity、benchmark、下游任务和生成质量；
\item 导出量化模型和推理 kernel 所需 metadata。
\end{enumerate}

PTQ 的优点是：

\begin{itemize}
\item 不需要重新训练；
\item 成本低；
\item 部署速度快；
\item 适合已有开源模型和内部 checkpoint；
\item 可以对不同硬件生成不同量化格式。
\end{itemize}

缺点是：

\begin{itemize}
\item 低 bit 下精度可能明显下降；
\item 对校准数据敏感；
\item outlier 和敏感层难处理；
\item 不同模型架构量化难度不同；
\item 需要专门算法降低误差。
\end{itemize}

\subsection{calibration dataset}
\label{subsec:calibration-dataset}

校准数据用于估计激活范围、outlier、通道重要性和 scale。它不一定需要很大，但必须代表模型实际使用分布。如果部署场景是代码生成，用普通百科文本校准可能不理想；如果部署场景是长文档问答，短指令数据可能无法覆盖长上下文激活分布。

校准数据需要考虑：

\begin{itemize}
\item 领域是否匹配；
\item prompt 长度是否匹配；
\item 是否包含真实 chat template；
\item 是否包含 system prompt；
\item 是否覆盖代码、数学、多语言等高风险输入；
\item 是否覆盖长上下文；
\item 数据量是否足够稳定估计统计量。
\end{itemize}

校准样本数量不是越多越好。太少统计不稳，太多校准耗时高。很多 PTQ 方法使用几百到几千条样本即可得到可用结果，但具体取决于模型和 bit-width。

\subsection{Quantization-Aware Training}
\label{subsec:quantization-aware-training}

QAT 在训练或微调中引入 fake quantization。Forward 时模拟量化：

[
\hat{x}
=======

s\cdot
\operatorname{round}
\left(
\frac{x}{s}
\right),
]

但 backward 时通常使用 straight-through estimator，简称 STE：

[
\frac{\partial \hat{x}}{\partial x}\approx 1.
]

这样模型可以在训练中适应量化误差。QAT 通常比 PTQ 精度更好，尤其在极低 bit、activation quantization 或任务敏感场景中。但 QAT 需要训练资源，也可能需要原始训练数据或高质量微调数据。

QAT 适合：

\begin{itemize}
\item 对质量要求极高的生产模型；
\item 低 bit 下 PTQ 精度不够；
\item 硬件固定，需要深度适配；
\item 模型训练流程可控；
\item 有足够校准/微调数据。
\end{itemize}

\begin{lstlisting}[language=Python,caption={Fake Quantization 的简化 PyTorch 实现},label={lst:fake-quantization-pytorch}]
import torch

class FakeQuantize(torch.autograd.Function):
@staticmethod
def forward(ctx, x, scale, qmin, qmax):
q = torch.round(x / scale)
q = torch.clamp(q, qmin, qmax)
return q * scale

```
@staticmethod
def backward(ctx, grad_output):
    # Straight-through estimator.
    return grad_output, None, None, None
```

def fake_quant(x, scale, qmin=-127, qmax=127):
return FakeQuantize.apply(x, scale, qmin, qmax)
\end{lstlisting}

这个代码只是说明 STE 思想。真实 QAT 需要 observer、scale 更新、per-channel quantization、梯度稳定性、混合精度和训练 recipe 配合。

\subsection{PTQ 与 QAT 的工程选择}
\label{subsec:ptq-vs-qat-engineering-choice}

工程上常见选择是：

[
\text{先 PTQ，若质量不够，再考虑 QAT 或低秩微调补偿。}
]

原因是 PTQ 成本低，可以快速得到多个候选量化配置。例如：

[
\text{W8A16},\quad
\text{W4A16 group size 128},\quad
\text{W8A8 SmoothQuant},\quad
\text{FP8 KV}.
]

如果某个配置满足质量和性能要求，就无需进行昂贵的 QAT。如果 PTQ 在关键任务上损失明显，再考虑 QAT、LoRA-QAT、混合精度保护敏感层等方案。

\begin{table}[H]
\centering
\small
\caption{PTQ 与 QAT 对比}
\label{tab:ptq-vs-qat}
\begin{tabular}{p{0.20\textwidth}p{0.34\textwidth}p{0.34\textwidth}}
\toprule
\textbf{维度} & \textbf{PTQ} & \textbf{QAT} \
\midrule
训练成本
& 低，只需校准或重构
& 高，需要训练或微调 \
\midrule
部署速度
& 快
& 慢 \
\midrule
精度潜力
& 中到高，取决于算法和 bit-width
& 通常更高 \
\midrule
数据需求
& 少量校准数据
& 微调或训练数据 \
\midrule
工程复杂度
& 中等
& 高 \
\midrule
适合场景
& 快速部署、已有模型、多硬件格式
& 极致低 bit、质量要求极高、训练流程可控 \
\bottomrule
\end{tabular}
\end{table}

\section{Weight-only Quantization}
\label{sec:weight-only-quantization}

Weight-only quantization 是 LLM 推理中最常见的量化方式之一。它只量化权重，激活仍保持 FP16/BF16：

[
\mat{Y}
=======

\mat{X}*{\mathrm{fp16}}
\cdot
\operatorname{Dequant}(\mat{W}*{\mathrm{int4/int8}}).
]

常见形式包括：

[
\text{W8A16},
\quad
\text{W4A16}.
]

其中 W 表示 weight bit-width，A 表示 activation bit-width。

\subsection{W8A16 / W4A16}
\label{subsec:w8a16-w4a16}

W8A16 使用 INT8 权重和 FP16/BF16 激活。它通常精度较稳，压缩率约 2 倍，kernel 实现也相对容易。

W4A16 使用 INT4 权重和 FP16/BF16 激活。它压缩率约 4 倍，但更容易损失精度，需要更复杂算法和 kernel 支持。W4A16 常用于单机部署更大模型，或在显存固定时提高并发和上下文长度。

在 weight-only quantization 中，矩阵乘法可以写作：

[
\mat{Y}
=======

\mat{X}
\left(
\mat{S}\odot
\mat{Q}_W
\right),
]

其中 $\mat{Q}_W$ 是低精度整数权重，$\mat{S}$ 是 scale。根据 scale 粒度不同，$\mat{S}$ 可以按 tensor、channel 或 group broadcast。

\subsection{group-wise quantization}
\label{subsec:group-wise-weight-quantization}

INT4 常用 group-wise quantization。设权重矩阵某一行被划分为多个 group，每个 group 大小为 $G$。第 $g$ 个 group 的权重为：

[
\vect{w}_g\in\mathbb{R}^{G}.
]

Symmetric INT4 量化可取：

[
s_g=
\frac{\max_i |w_{g,i}|}{7}.
]

量化：

[
q_{g,i}
=======

\operatorname{clip}
\left(
\left\lfloor
\frac{w_{g,i}}{s_g}
\right\rceil,
-8,
7
\right).
]

反量化：

[
\hat{w}*{g,i}=s_gq*{g,i}.
]

Group size 常见选择如 32、64、128。Group 越小，scale 更精细，精度更好；但 scale metadata 更多，kernel 读取 scale 更频繁。Group 越大，metadata 少，但 outlier 更容易拉大量化范围。

\begin{figure}[H]
\centering
\begin{tikzpicture}[
font=\small,
cell/.style={draw, minimum width=0.38cm, minimum height=0.34cm},
groupbox/.style={draw, rounded corners=2pt, dashed, thick},
arrow/.style={-{Latex[length=2.0mm]}, thick}
]

```
\node[font=\bfseries] at (0,3.0) {Group-wise INT4 Weight Quantization};

\foreach \i in {0,...,23} {
  \node[cell, fill=blue!10] at ({-4.6+0.38*\i},1.45) {};
}

\draw[groupbox] (-4.82,1.17) rectangle (-1.62,1.72);
\draw[groupbox] (-1.54,1.17) rectangle (1.66,1.72);
\draw[groupbox] (1.74,1.17) rectangle (4.94,1.72);

\node[font=\scriptsize] at (-3.22,0.9) {scale $s_0$};
\node[font=\scriptsize] at (0.06,0.9) {scale $s_1$};
\node[font=\scriptsize] at (3.34,0.9) {scale $s_2$};

\node[align=left, text width=10.5cm] at (0,-0.65) {
  每个 group 使用独立 scale。group size 越小，量化误差通常越低；
  但 scale metadata 和 kernel 复杂度也更高。
};
```

\end{tikzpicture}
\caption{按 group 量化权重并保存 scale}
\label{fig:group-wise-weight-quantization}
\end{figure}

\subsection{packing 与 kernel}
\label{subsec:int4-packing-and-kernel}

INT4 权重需要 packing。一个 8-bit byte 可以存两个 4-bit 权重：

[
\text{byte}
===========

(q_0\ &\ 0xF)
|
((q_1\ &\ 0xF)\ll 4).
]

Packing 的目的包括：

\begin{itemize}
\item 降低存储大小；
\item 降低 HBM 读取字节数；
\item 对齐 kernel 的向量化加载；
\item 适配 Tensor Core 或自定义 MMA 路径。
\end{itemize}

下面是教学版 INT4 packing 代码。

\begin{lstlisting}[language=Python,caption={教学版 INT4 packing 与 unpacking},label={lst:int4-pack-unpack}]
import torch

def pack_int4(q):
"""
q: int tensor with values in [0, 15], shape [..., N]
N must be even.
return: uint8 tensor shape [..., N // 2]
"""
assert q.shape[-1] % 2 == 0
q = q.to(torch.uint8)

```
low = q[..., 0::2] & 0x0F
high = (q[..., 1::2] & 0x0F) << 4

return low | high
```

def unpack_int4(packed):
"""
packed: uint8 tensor shape [..., N_packed]
return: uint8 tensor with values in [0, 15], shape [..., 2 * N_packed]
"""
low = packed & 0x0F
high = (packed >> 4) & 0x0F

```
out = torch.empty(
    *packed.shape[:-1],
    packed.shape[-1] * 2,
    dtype=torch.uint8,
    device=packed.device,
)
out[..., 0::2] = low
out[..., 1::2] = high
return out
```

\end{lstlisting}

真实 kernel 不会先 unpack 整个权重矩阵，而是在 GEMM 内部按 tile 解包、转换、乘 scale，并与 activation 做乘加。否则 INT4 的带宽优势会被临时展开抵消。

\subsection{accuracy drop 的来源}
\label{subsec:weight-only-accuracy-drop-sources}

Weight-only quantization 的精度下降主要来自：

\begin{itemize}
\item 权重分布 outlier；
\item 通道重要性差异；
\item 量化误差在层间累积；
\item attention projection 或 FFN 某些矩阵更敏感；
\item 最前/最后几层更敏感；
\item embedding/lm head 更敏感；
\item group size 过大；
\item calibration 数据不匹配；
\item 低 bit 下 rounding error 太大。
\end{itemize}

一个常见工程策略是混合精度保护敏感部分。例如：

\begin{itemize}
\item embedding 保持 FP16；
\item lm head 保持 FP16；
\item 第一层和最后一层保持 FP16/INT8；
\item 某些 outlier channels 使用更高精度；
\item attention output 或 FFN down projection 使用更保守量化。
\end{itemize}

这类策略会牺牲部分压缩率，但常能显著改善质量。

\section{GPTQ}
\label{sec:gptq}

GPTQ 是 LLM PTQ 中非常有代表性的方法。它属于 weight-only post-training quantization，核心思想是利用近似二阶信息，在逐列量化权重时补偿误差，使量化后的层输出尽量接近原始层输出。

\subsection{二阶误差直觉}
\label{subsec:gptq-second-order-intuition}

对某个 Linear 层：

[
\mat{Y}=\mat{X}\mat{W}.
]

量化后权重为：

[
\hat{\mat{W}}=\mat{W}+\Delta\mat{W}.
]

输出误差为：

[
\Delta\mat{Y}
=============

\mat{X}\Delta\mat{W}.
]

我们希望最小化：

[
|\mat{X}\mat{W}-\mat{X}\hat{\mat{W}}|_2^2
=========================================

|\mat{X}\Delta\mat{W}|_2^2.
]

展开可得误差与输入协方差有关：

[
|\mat{X}\Delta\mat{W}|_2^2
==========================

\operatorname{tr}
\left(
\Delta\mat{W}^{\mathsf{T}}
\mat{X}^{\mathsf{T}}\mat{X}
\Delta\mat{W}
\right).
]

令：

[
\mat{H}\approx \mat{X}^{\mathsf{T}}\mat{X}.
]

则量化误差可以用 Hessian-like 矩阵 $\mat{H}$ 加权。直觉是：如果某个权重方向对输出影响大，则量化误差应该更小；如果某个方向对输出影响小，则可以承受更大误差。

\subsection{逐列量化与误差补偿}
\label{subsec:gptq-sequential-quantization}

GPTQ 的一个关键思想是逐列或分块量化权重，并在量化某一列后，把误差补偿到尚未量化的列上。

设当前要量化权重列 $\vect{w}_i$，量化后为 $\hat{\vect{w}}_i$，误差为：

[
\vect{e}_i=\vect{w}_i-\hat{\vect{w}}_i.
]

如果直接忽略误差，后续输出会偏离。GPTQ 利用 $\mat{H}^{-1}$ 的信息，把该误差对后续列的影响进行更新。抽象地：

[
\mat{W}*{j}
\leftarrow
\mat{W}*{j}
-----------

\alpha_{ij}\vect{e}_i,
\quad j>i.
]

其中 $\alpha_{ij}$ 来自 Hessian inverse 的相关项。这样后续列在量化前已经吸收部分误差补偿。

这个过程类似“带二阶信息的贪心量化”。它不是简单地对每个 group 独立 round，而是考虑输入分布和权重列之间的耦合。

\begin{figure}[H]
\centering
\begin{tikzpicture}[
font=\small,
col/.style={draw, rounded corners=2pt, minimum width=0.55cm, minimum height=1.6cm, align=center},
arrow/.style={-{Latex[length=2.1mm]}, thick}
]

```
\node[font=\bfseries] at (0,3.0) {GPTQ：逐列量化并向后续列传播误差补偿};

\foreach \i/\x/\colr in {0/-3.6/green!12,1/-2.9/blue!10,2/-2.2/blue!10,3/-1.5/blue!10,4/-0.8/blue!10} {
  \node[col, fill=\colr] (c\i) at (\x,1.25) {$w_{\i}$};
}

\node[draw, rounded corners=3pt, fill=orange!14, minimum width=1.4cm, minimum height=0.6cm, align=center] (q) at (-3.6,-0.1)
  {quantize\\$w_0$};

\draw[arrow] (c0) -- (q);

\foreach \i in {1,2,3,4} {
  \draw[arrow, red!70] (q) .. controls (-3.0,-0.85) and (-2.2,-0.85) .. (c\i.south);
}

\node[align=left, text width=10.5cm] at (0,-1.5) {
  GPTQ 不是独立量化每一列，而是在量化当前列后，将误差按近似二阶信息补偿到后续未量化列。
  这样可以降低层输出重构误差。
};
```

\end{tikzpicture}
\caption{GPTQ 的逐列量化与误差补偿直觉}
\label{fig:gptq-sequential-error-compensation}
\end{figure}

\subsection{Hessian 近似}
\label{subsec:gptq-hessian-approximation}

GPTQ 不会计算完整训练 loss 的真实 Hessian，而是对每层局部重构误差使用近似：

[
\mat{H}=\mat{X}^{\mathsf{T}}\mat{X}.
]

这里 $\mat{X}$ 是校准数据经过该层时的输入 activation。实际实现中会加入 damping，避免矩阵病态：

[
\mat{H}_{\mathrm{damped}}
=========================

\mat{H}
+
\lambda \operatorname{diag}(\mat{H}).
]

然后使用 Cholesky 分解或其他数值方法得到所需的逆相关信息。由于层维度很大，GPTQ 通常采用 blockwise 处理，避免一次操作完整巨大矩阵。

\subsection{工程流程}
\label{subsec:gptq-engineering-flow}

GPTQ 的工程流程可以概括为：

\begin{enumerate}
\item 准备校准样本；
\item 对模型逐层 forward，收集每层输入 activation；
\item 对每个 Linear 层构造 $\mat{H}\approx \mat{X}^{\mathsf{T}}\mat{X}$；
\item 对权重按 block/group 逐步量化；
\item 使用 Hessian inverse 信息进行误差补偿；
\item 导出 INT4/INT3/INT8 权重、scale、zero point 和 packing layout；
\item 使用支持 GPTQ 格式的推理 kernel 执行。
\end{enumerate}

GPTQ 的优势是低 bit 下精度较好，尤其适合 weight-only INT4。代价是量化过程比简单 min-max PTQ 更复杂，需要校准数据和矩阵分解，量化耗时更长。

\section{AWQ}
\label{sec:awq}

AWQ，Activation-aware Weight Quantization，是另一类重要的 LLM weight-only PTQ 方法。它的核心观察是：权重并非同等重要，少量 salient weight channels 对模型输出影响很大；识别这些重要通道时，应该关注 activation 分布，而不只是权重本身。

\subsection{activation-aware 的动机}
\label{subsec:awq-activation-aware-motivation}

对 Linear 层：

[
\mat{Y}=\mat{X}\mat{W}.
]

某个权重通道的量化误差对输出影响，不仅取决于权重误差本身，还取决于输入 activation 的大小。如果某个输入通道的 activation 经常很大，那么对应权重列的误差会被放大。

设第 $j$ 个输入通道 activation 平均幅度较大：

[
\mathbb{E}[|X_j|]\ \text{large}.
]

那么权重列 $\mat{W}_{j,:}$ 的量化误差会对输出产生更大影响：

[
\Delta\mat{Y}
=============

\sum_j
\mat{X}*j
\Delta\mat{W}*{j,:}.
]

因此，AWQ 使用 activation 统计来判断哪些 channels 更重要，并对这些 channels 做保护。

\subsection{salient channels}
\label{subsec:awq-salient-channels}

AWQ 的思想可以简化为：

[
\text{找出 activation 大、对输出敏感的 channels}
\rightarrow
\text{降低这些 channels 的量化误差}.
]

一种直接做法是混合精度：把重要 channels 保持高精度。但混合精度可能不硬件友好，因为 kernel 需要同时处理 INT4 和 FP16 权重，导致实现复杂。

AWQ 更偏向硬件友好方案：通过缩放变换保护 salient channels，同时保持整体权重仍可统一低 bit 量化。

\subsection{scale 变换}
\label{subsec:awq-scale-transformation}

考虑 Linear：

[
\mat{Y}=\mat{X}\mat{W}.
]

引入对角缩放矩阵 $\mat{S}$：

[
\mat{Y}
=======

(\mat{X}\mat{S}^{-1})
(\mat{S}\mat{W}).
]

这是数学等价变换：

[
\mat{X}\mat{W}
==============

\mat{X}\mat{S}^{-1}\mat{S}\mat{W}.
]

AWQ 的直觉是，适当放大某些重要 weight channels，使其在量化时相对误差更小，同时在 activation 侧用反向缩放保持等价。实际中，需要搜索或确定合适的 scale，使量化后误差最小。

这个方法的优点是：最终仍然可以使用统一的 low-bit weight-only kernel，而不必在 kernel 内混合处理不同精度的权重。

\begin{figure}[H]
\centering
\begin{tikzpicture}[
font=\small,
box/.style={draw, rounded corners=3pt, minimum width=1.8cm, minimum height=0.65cm, align=center},
arrow/.style={-{Latex[length=2.1mm]}, thick}
]

```
\node[font=\bfseries] at (0,3.05) {AWQ：通过等价缩放保护重要通道};

\node[box, fill=blue!8] (x) at (-5.0,1.3) {$X$};
\node[box, fill=orange!14] (sinv) at (-2.9,1.3) {$S^{-1}$\\on activation};
\node[box, fill=green!10] (sweight) at (-0.6,1.3) {$S W$\\scaled weight};
\node[box, fill=purple!10] (quant) at (1.7,1.3) {INT4\\quantize};
\node[box, fill=red!8] (y) at (4.2,1.3) {$Y$};

\draw[arrow] (x) -- (sinv);
\draw[arrow] (sinv) -- (sweight);
\draw[arrow] (sweight) -- (quant);
\draw[arrow] (quant) -- (y);

\node[align=left, text width=10.5cm] at (0,-0.65) {
  AWQ 使用 activation 统计识别重要通道，并通过等价缩放降低这些通道的量化误差。
  最终仍保持硬件友好的 weight-only quantization。
};
```

\end{tikzpicture}
\caption{AWQ 中 activation-aware scaling 的直觉}
\label{fig:awq-scale-transformation}
\end{figure}

\subsection{AWQ 的适用场景}
\label{subsec:awq-use-cases}

AWQ 适合：

\begin{itemize}
\item INT4 weight-only 推理；
\item 希望校准成本较低；
\item 希望保持 kernel 硬件友好；
\item 模型是 instruction-tuned LLM 或多模态 LLM；
\item 部署在边缘设备或显存紧张环境；
\item 不希望做复杂二阶重构。
\end{itemize}

与 GPTQ 相比，AWQ 更强调 activation-aware channel protection 和硬件友好性；GPTQ 更强调二阶误差补偿和逐列重构。实际中二者都常作为 INT4 量化候选，需要在目标模型和任务上评估。

\section{SmoothQuant 与 W8A8}
\label{sec:smoothquant-w8a8}

Weight-only quantization 只量化权重，激活仍是 FP16/BF16。若想进一步利用 INT8 GEMM，需要同时量化权重和激活，即 W8A8。SmoothQuant 是 LLM W8A8 PTQ 中代表性方法之一。

\subsection{activation outlier 问题}
\label{subsec:activation-outlier-problem}

LLM 激活中常存在 outlier channels。某些通道的 activation 幅度远大于其他通道。如果直接对 activation 做 per-tensor INT8 量化，scale 会被 outlier 拉大，导致大多数普通值分辨率不足。

设 activation 向量某些通道非常大：

[
|x_j|\gg |x_k|.
]

Per-tensor scale 必须覆盖最大值：

[
s=\frac{\max_j |x_j|}{127}.
]

若 $\max_j |x_j|$ 很大，普通通道的值会被量化到很粗的整数网格上，误差增大。

权重相比激活更容易量化，因为权重是固定的，可以按通道或 group 使用 scale；activation 是动态的，范围随输入变化，outlier 更难处理。

\subsection{将量化难度从 activation 转移到 weight}
\label{subsec:smoothquant-migrate-difficulty}

SmoothQuant 的核心思想是使用等价变换，将 activation outlier 的难度迁移到 weight 上。

对于 Linear：

[
\mat{Y}=\mat{X}\mat{W}.
]

引入对角矩阵 $\mat{S}$：

[
\mat{Y}
=======

(\mat{X}\mat{S}^{-1})
(\mat{S}\mat{W}).
]

如果某些 activation channel 很大，可以选择较大的 scale $s_j$，让：

[
\mat{X}'_j=\frac{\mat{X}_j}{s_j}
]

变得更平滑。对应权重：

[
\mat{W}'_j=s_j\mat{W}_j.
]

这样 activation 更容易 INT8 量化，而权重量化难度增加。由于权重是静态的，更容易通过 per-channel scale 处理。因此整体 W8A8 可行性提高。

SmoothQuant 通常用参数 $\alpha$ 控制平滑程度。若 activation 通道幅度为 $a_j$，weight 通道幅度为 $w_j$，scale 可设计为某种函数：

[
s_j
===

\frac{a_j^{\alpha}}
{w_j^{1-\alpha}}.
]

直觉上，$\alpha$ 越大，越多地平滑 activation；$\alpha$ 越小，越少改变权重。

\begin{figure}[H]
\centering
\begin{tikzpicture}[
font=\small,
box/.style={draw, rounded corners=3pt, minimum width=1.95cm, minimum height=0.7cm, align=center},
arrow/.style={-{Latex[length=2.1mm]}, thick}
]

```
\node[font=\bfseries] at (0,3.1) {SmoothQuant：平滑 Activation Outlier};

\node[box, fill=blue!8] (x) at (-5.1,1.45) {activation\\with outliers};
\node[box, fill=orange!14] (scale) at (-2.7,1.45) {$X'=XS^{-1}$\\smooth};
\node[box, fill=green!10] (w) at (-0.2,1.45) {$W'=SW$\\adjust weight};
\node[box, fill=purple!10] (w8a8) at (2.3,1.45) {W8A8\\GEMM};
\node[box, fill=red!8] (y) at (4.8,1.45) {$Y$};

\draw[arrow] (x) -- (scale);
\draw[arrow] (scale) -- (w8a8);
\draw[arrow] (w) -- (w8a8);
\draw[arrow] (w8a8) -- (y);

\node[align=left, text width=10.5cm] at (0,-0.65) {
  SmoothQuant 通过等价缩放把 activation outlier 平滑掉，
  将部分量化难度转移到静态权重上，从而支持 W8A8 PTQ。
};
```

\end{tikzpicture}
\caption{SmoothQuant 的 activation 平滑思想}
\label{fig:smoothquant-activation-smoothing}
\end{figure}

\subsection{W8A8 GEMM}
\label{subsec:w8a8-gemm}

W8A8 的目标是让权重和激活都用 INT8 表示，矩阵乘法使用 INT8 Tensor Core 或等价硬件路径：

[
\mat{Y}_{\mathrm{int32}}
========================

\mat{X}*{\mathrm{int8}}
\mat{W}*{\mathrm{int8}}.
]

累积通常使用 INT32，然后再乘 scale 转回 FP16/BF16 或继续保持低精度：

[
\hat{\mat{Y}}
=============

s_Xs_W
\mat{Y}_{\mathrm{int32}}.
]

W8A8 的潜在收益包括：

\begin{itemize}
\item 权重读取减半；
\item 激活读取减半；
\item 使用 INT8 矩阵乘加硬件；
\item 降低中间 activation 显存。
\end{itemize}

难点包括：

\begin{itemize}
\item activation 动态范围；
\item outlier；
\item scale 管理；
\item layernorm/RMSNorm 后分布变化；
\item residual add 的 scale 对齐；
\item kernel 对动态 scale 的支持；
\item 不同层量化敏感度差异。
\end{itemize}

\section{FP8 与低精度浮点}
\label{sec:fp8-low-precision-floating-point}

INT 量化使用固定线性 scale 映射到整数。FP8 则使用 8 bit 浮点格式，包含 sign、exponent 和 mantissa。相比 INT8，FP8 可以覆盖更大动态范围，适合训练和推理中的激活、权重或梯度。

\subsection{E4M3 与 E5M2}
\label{subsec:e4m3-e5m2}

常见 FP8 格式包括：

\begin{itemize}
\item \textbf{E4M3}：4 bit exponent，3 bit mantissa；
\item \textbf{E5M2}：5 bit exponent，2 bit mantissa。
\end{itemize}

E4M3 尾数更多，精度更高，但动态范围较小。E5M2 指数更多，动态范围更大，但精度较低。不同张量可能适合不同格式。例如 forward activation 和 weight 可能更偏 E4M3，gradient 可能需要更大动态范围。

\begin{figure}[H]
\centering
\begin{tikzpicture}[
font=\small,
bit/.style={draw, minimum width=0.5cm, minimum height=0.45cm, align=center}
]

```
\node[font=\bfseries] at (0,3.0) {FP8 E4M3 与 E5M2 bit layout};

\node[anchor=east] at (-4.2,1.55) {E4M3};
\node[bit, fill=red!15] at (-3.6,1.55) {S};
\foreach \i/\x in {0/-3.05,1/-2.5,2/-1.95,3/-1.4} {
  \node[bit, fill=blue!12] at (\x,1.55) {E};
}
\foreach \i/\x in {0/-0.85,1/-0.3,2/0.25} {
  \node[bit, fill=green!12] at (\x,1.55) {M};
}

\node[anchor=east] at (-4.2,0.65) {E5M2};
\node[bit, fill=red!15] at (-3.6,0.65) {S};
\foreach \i/\x in {0/-3.05,1/-2.5,2/-1.95,3/-1.4,4/-0.85} {
  \node[bit, fill=blue!12] at (\x,0.65) {E};
}
\foreach \i/\x in {0/-0.3,1/0.25} {
  \node[bit, fill=green!12] at (\x,0.65) {M};
}

\node[align=left, text width=10.4cm] at (0,-0.75) {
  E4M3 提供更多 mantissa 精度，E5M2 提供更大 exponent 动态范围。
  FP8 通常需要配合 scale、amax history 和硬件支持使用。
};
```

\end{tikzpicture}
\caption{FP8 两种常见格式的 bit 分配}
\label{fig:fp8-e4m3-e5m2}
\end{figure}

\subsection{scale 与 amax}
\label{subsec:fp8-scale-amax}

FP8 表示范围有限，需要 scale 管理。常见做法是记录 tensor 的绝对最大值：

[
a_{\max}=\max_i |x_i|.
]

根据 FP8 格式最大可表示值 $f_{\max}$，选择 scale：

[
s=
\frac{a_{\max}}{f_{\max}}.
]

量化时：

[
x_{\mathrm{fp8}}
================

\operatorname{cast}_{\mathrm{fp8}}
\left(
\frac{x}{s}
\right).
]

反量化时：

[
\hat{x}=s\cdot x_{\mathrm{fp8}}.
]

为了避免 scale 随 batch 抖动，系统可能维护 amax history，并用滑动窗口或 delayed scaling 更新 scale。这样可以减少异常 batch 导致的 scale 波动。

FP8 的工程难点在于：每个 tensor 或每类 tensor 的 scale 如何选择、何时更新、是否 per-tensor 或 per-channel、是否 delayed、如何与分布式并行和 kernel fusion 配合。

\subsection{FP8 训练与推理}
\label{subsec:fp8-training-inference}

FP8 可用于训练和推理。训练中，部分 GEMM 输入可以使用 FP8，累积通常使用更高精度，master weights 仍可能保留 BF16/FP32。推理中，权重和激活可使用 FP8，以降低带宽和提高吞吐。

FP8 相比 INT8 的优势是动态范围更自然；相比 BF16 的优势是字节数更低。它更依赖硬件支持，例如特定 GPU 架构的 FP8 Tensor Core，以及软件库对 scale 管理和 kernel 的支持。

FP8 与 INT8/INT4 的关系不是谁完全替代谁：

\begin{itemize}
\item INT4 weight-only 适合极致压缩权重；
\item INT8 W8A8 适合成熟整数硬件路径；
\item FP8 适合有硬件支持、希望兼顾动态范围和低字节数的训练/推理；
\item KV Cache 可用 FP8 或 INT8，取决于质量和 kernel 支持。
\end{itemize}

\section{KV Cache 量化}
\label{sec:kv-cache-quantization}

前面几章已经看到，KV Cache 是长上下文和高并发 serving 的关键显存瓶颈。KV Cache 量化的目标是降低：

[
M_{\mathrm{KV}}
===============

2LBSn_{kv}d_hs_{\mathrm{dtype}}.
]

如果把 BF16 KV Cache 换成 FP8 或 INT8，$s_{\mathrm{dtype}}$ 从 2 bytes 降为 1 byte，理论上 KV Cache 显存减半。若使用 INT4，理论上可进一步减半，但质量风险和 kernel 复杂度更高。

\subsection{KV Cache 为什么适合量化}
\label{subsec:why-kv-cache-quantization}

KV Cache 量化适合以下场景：

\begin{itemize}
\item 长上下文；
\item 高并发；
\item batch decode；
\item GQA/MQA 后仍然 cache 很大；
\item 显存主要被 KV Cache 占用；
\item decode memory bandwidth 成为瓶颈。
\end{itemize}

在 decode 中，每一步都读取历史 K/V。如果 K/V 字节数减半，HBM 读取压力也可能下降。尤其对于 memory-bound decode attention，KV 量化有潜在速度收益。

\subsection{K 与 V 的敏感性}
\label{subsec:k-v-sensitivity}

K 和 V 的量化误差影响不同。K 参与 attention score：

[
qk^{\mathsf{T}}.
]

K 的误差会改变 attention 分布：

[
\softmax(q(K+\Delta K)^{\mathsf{T}}).
]

V 的误差会改变被加权汇聚的内容：

[
P(V+\Delta V).
]

K 的误差可能通过 softmax 放大，因为 score 的相对差异会影响注意力权重；V 的误差则更像输出内容误差。实际系统可能对 K 和 V 使用不同 scale、不同格式，甚至只量化 V 或对 K 更保守。

\subsection{per-token 与 per-channel KV scale}
\label{subsec:kv-per-token-per-channel-scale}

KV Cache 量化需要 scale。常见粒度包括：

\begin{itemize}
\item per-tensor：整个 cache 一个 scale，简单但误差大；
\item per-layer：每层一个 scale；
\item per-head：每个 KV head 一个 scale；
\item per-token：每个 token 的 K/V 向量一个 scale；
\item per-channel：每个 head dimension 一个 scale；
\item per-block：PagedAttention block 内共享 scale。
\end{itemize}

Per-token scale 可以适应不同 token 的幅度变化，但 metadata 较多；per-channel scale 可以保护某些维度，但 kernel 读取 scale 更复杂；per-block scale 与 paged KV blocks 很自然，metadata 与 block table 可结合。

KV Cache 量化的工程设计需要同时考虑：

\begin{itemize}
\item scale metadata 显存；
\item attention kernel 读取 scale 的开销；
\item PagedAttention block layout；
\item TP 下 KV heads 切分；
\item prefix cache 共享 blocks；
\item copy-on-write；
\item 长上下文质量。
\end{itemize}

\begin{figure}[H]
\centering
\begin{tikzpicture}[
font=\small,
block/.style={draw, rounded corners=2pt, minimum width=0.7cm, minimum height=0.45cm, align=center},
scale/.style={draw, rounded corners=2pt, minimum width=0.9cm, minimum height=0.35cm, align=center, fill=orange!14}
]

```
\node[font=\bfseries] at (0,3.0) {Paged KV Cache 中的 per-block scale};

\foreach \i/\x in {0/-3.6,1/-2.7,2/-1.8,3/-0.9,4/0.0} {
  \node[block, fill=blue!10] (b\i) at (\x,1.35) {$B_{\i}$};
  \node[scale] at (\x,0.72) {$s_{\i}$};
}

\foreach \i/\x in {5/1.2,6/2.1,7/3.0} {
  \node[block, fill=green!10] (b\i) at (\x,1.35) {$B_{\i}$};
  \node[scale] at (\x,0.72) {$s_{\i}$};
}

\node[align=left, text width=10.5cm] at (0,-0.75) {
  KV Cache 量化可以为每个 physical block 保存 scale。
  这种方式与 PagedAttention 的 block 管理自然结合，但 attention kernel 必须读取并应用对应 scale。
};
```

\end{tikzpicture}
\caption{KV Cache per-block scale 与 paged blocks}
\label{fig:kv-cache-per-block-scale}
\end{figure}

\section{量化对 Serving Infra 的影响}
\label{sec:quantization-impact-on-serving-infra}

量化不是只改变模型文件格式。它会影响模型加载、GPU kernel、调度器、KV Cache、并行推理、监控和质量评估。

\subsection{模型加载与格式}
\label{subsec:quantized-model-loading-format}

量化模型通常包含：

\begin{itemize}
\item packed low-bit weights；
\item scale tensors；
\item zero point tensors；
\item group size metadata；
\item quantization method metadata；
\item layer-wise configuration；
\item kernel expected layout；
\item optional outlier/salient channel metadata。
\end{itemize}

Serving engine 加载模型时必须确保：

\begin{itemize}
\item 权重 layout 与 kernel 匹配；
\item TP 切分与 packed format 兼容；
\item scale 切分与权重 shard 对齐；
\item checkpoint 转换没有破坏 group 边界；
\item lm head、embedding 等特殊层处理正确；
\item adapter 或 LoRA 与量化权重兼容。
\end{itemize}

一个常见错误是：量化权重在单 GPU 上能跑，但切到 TP 后 group scale 对齐出错，导致输出异常。原因是 TP shard 边界可能切开 group，或 scale metadata 没有同步切分。因此，量化格式必须从一开始就考虑并行推理。

\subsection{kernel 支持}
\label{subsec:quantized-kernel-support}

量化收益高度依赖 kernel。一个量化配置在纸面上显存很省，但如果没有高性能 kernel，可能速度很慢。

Kernel 需要支持：

\begin{itemize}
\item INT4/INT8 unpack；
\item per-group/per-channel scale；
\item zero point；
\item fused dequant GEMM；
\item Tensor Parallel shard；
\item activation dtype；
\item batch GEMM shape；
\item CUDA graph 或动态 shape；
\item MoE grouped GEMM；
\item LoRA adapter 合并或旁路。
\end{itemize}

如果推理框架 fallback 到普通 FP16 GEMM，可能会先反量化权重，导致显存和速度都不理想。因此上线前必须通过 profiler 或日志确认量化 kernel 被真正调用。

\subsection{质量评估}
\label{subsec:quantization-quality-evaluation}

量化上线不能只看 perplexity。需要多层评估：

\begin{itemize}
\item 语言建模 perplexity；
\item 通用 benchmark；
\item 指令跟随；
\item 代码能力；
\item 数学推理；
\item 多语言；
\item 长上下文；
\item RAG 场景；
\item 安全策略；
\item 真实线上 prompt 回放；
\item 人工偏好评测。
\end{itemize}

低 bit 量化常见问题包括：

\begin{itemize}
\item 输出重复；
\item 格式遵循下降；
\item 长推理链更容易出错；
\item 代码生成中小语法错误增加；
\item 数学计算稳定性下降；
\item 长上下文引用准确率下降；
\item 多语言弱语言质量下降。
\end{itemize}

因此，量化模型需要按业务场景评估，而不是只看一个平均指标。

\subsection{成本模型}
\label{subsec:quantization-cost-model}

量化的最终目标通常是降低成本。可以定义单位 token 成本：

[
C_{\mathrm{token}}
==================

\frac{C_{\mathrm{GPU/hour}}}
{\mathrm{tokens/hour}}.
]

量化可能通过以下路径降低成本：

\begin{itemize}
\item 用更少 GPU 部署同一模型；
\item 同一 GPU 承载更多并发；
\item 提高 tokens/s；
\item 降低能耗；
\item 留出更多显存给 KV Cache；
\item 支持更长上下文而不扩容。
\end{itemize}

但量化也可能引入额外成本：

\begin{itemize}
\item 质量下降导致重试或人工干预；
\item kernel 不成熟导致延迟变差；
\item 多格式模型管理复杂；
\item 量化校准与评估成本；
\item 特定硬件依赖增加。
\end{itemize}

因此，工程决策应比较端到端成本，而不是只看模型大小。

\section{一个教学版量化实现}
\label{sec:toy-quantization-implementation}

下面给出一个教学版 per-group symmetric INT4 权重量化实现。它用于理解数据结构，不是生产级实现。

\begin{lstlisting}[language=Python,caption={Per-group symmetric INT4 权重量化示意},label={lst:per-group-int4-quantization}]
import torch

def quantize_int4_per_group(weight, group_size=128):
"""
weight: [out_features, in_features], FP16/BF16/FP32
return:
q_packed: packed uint8 INT4 weights
scales: [out_features, num_groups]
"""
out_features, in_features = weight.shape
assert in_features % group_size == 0
assert group_size % 2 == 0

```
num_groups = in_features // group_size

w = weight.float().view(out_features, num_groups, group_size)

max_abs = w.abs().amax(dim=-1).clamp_min(1e-8)
scales = max_abs / 7.0

q = torch.round(w / scales[..., None]).clamp(-8, 7)

# Convert signed INT4 [-8, 7] to unsigned nibble [0, 15].
q_u = (q.to(torch.int16) + 8).to(torch.uint8)

# Pack two int4 values into one byte.
low = q_u[..., 0::2] & 0x0F
high = (q_u[..., 1::2] & 0x0F) << 4
q_packed = (low | high).contiguous()

return q_packed, scales.half()
```

def dequantize_int4_per_group(q_packed, scales, group_size=128):
"""
Slow reference dequantization for correctness tests.
"""
out_features, num_groups, packed_per_group = q_packed.shape
assert packed_per_group * 2 == group_size

```
low = q_packed & 0x0F
high = (q_packed >> 4) & 0x0F

q_u = torch.empty(
    out_features,
    num_groups,
    group_size,
    dtype=torch.uint8,
    device=q_packed.device,
)
q_u[..., 0::2] = low
q_u[..., 1::2] = high

q = q_u.to(torch.int16) - 8
w = q.float() * scales.float()[..., None]

return w.view(out_features, num_groups * group_size)
```

\end{lstlisting}

量化后需要验证重构误差：

[
\frac{|\mat{W}-\hat{\mat{W}}|_F}
{|\mat{W}|_F}.
]

更重要的是验证层输出误差：

[
\frac{|\mat{X}\mat{W}-\mat{X}\hat{\mat{W}}|_F}
{|\mat{X}\mat{W}|_F}.
]

因为模型质量受输出误差影响，而不是只受权重误差影响。

\begin{lstlisting}[language=Python,caption={量化权重的重构误差与层输出误差检查},label={lst:quantization-error-check}]
def check_quantization_error(weight, x, group_size=128):
q_packed, scales = quantize_int4_per_group(weight, group_size)
w_hat = dequantize_int4_per_group(q_packed, scales, group_size)

```
weight_error = torch.norm(weight.float() - w_hat.float())
weight_error = weight_error / torch.norm(weight.float())

y_ref = x.float() @ weight.float().t()
y_hat = x.float() @ w_hat.float().t()

output_error = torch.norm(y_ref - y_hat) / torch.norm(y_ref)

return {
    "relative_weight_error": weight_error.item(),
    "relative_output_error": output_error.item(),
}
```

\end{lstlisting}

生产系统中，不能对每层都用这种慢速反量化路径执行推理。但这些 reference 实现适合单元测试：确认 packing、scale、zero point、group 切分和 TP shard 都正确。

\section{量化方案选型}
\label{sec:quantization-scheme-selection}

不同量化方案适合不同部署目标。

\subsection{显存优先}
\label{subsec:quantization-memory-first}

如果目标是尽可能降低显存，使模型能装进更少 GPU，可优先考虑：

[
\text{W4A16 GPTQ/AWQ}
]

或更低 bit weight-only。适合：

\begin{itemize}
\item 单机部署大模型；
\item 边缘设备；
\item 低并发但显存紧张；
\item 成本敏感场景；
\item 能接受少量质量损失。
\end{itemize}

需要重点验证模型质量和 kernel 是否高效。

\subsection{吞吐优先}
\label{subsec:quantization-throughput-first}

如果目标是提升吞吐，需要看硬件是否支持低精度计算。候选包括：

[
\text{W8A8},
\quad
\text{FP8},
\quad
\text{INT4 fused GEMM}.
]

如果只压缩权重但计算仍受 activation 或 compute 限制，速度收益可能有限。吞吐优先场景应使用 profiler 分析：

\begin{itemize}
\item weight read bandwidth 是否下降；
\item GEMM kernel 是否使用低精度路径；
\item dequantization 是否 fused；
\item decode batch 是否变大；
\item TTFT 和 TPOT 是否改善。
\end{itemize}

\subsection{质量优先}
\label{subsec:quantization-quality-first}

如果质量优先，可考虑：

\begin{itemize}
\item W8A16 而非 W4A16；
\item INT4 中使用更小 group size；
\item 保留敏感层为 FP16；
\item 使用 GPTQ/AWQ 而非 naive min-max；
\item 使用 QAT；
\item 对业务数据做校准；
\item 对长上下文和推理任务做专项评估。
\end{itemize}

质量优先并不意味着不能量化，而是需要保守配置和充分评估。

\subsection{推理框架兼容性}
\label{subsec:inference-framework-compatibility}

量化方案必须与推理框架兼容。选择前需要确认：

\begin{itemize}
\item 框架是否支持该量化格式；
\item 是否支持目标 GPU；
\item 是否支持 TP/PP；
\item 是否支持 PagedAttention；
\item 是否支持 LoRA；
\item 是否支持 speculative decoding；
\item 是否支持 streaming；
\item 是否支持 sharded checkpoint；
\item 是否支持导出和加载 scale metadata；
\item 是否能在生产中稳定运行。
\end{itemize}

一个在离线脚本中可用的量化模型，不一定能直接放进 serving engine。模型格式、kernel、scheduler、KV Cache 和并行策略必须一起验证。

\section{本章小结}
\label{sec:chapter15-summary}

本章系统介绍了大模型推理中的量化技术，从数值表示、PTQ/QAT、weight-only quantization，到 GPTQ、AWQ、SmoothQuant、FP8、KV Cache 量化和 serving infra 影响。

\begin{itemize}
\item \textbf{量化的核心目标} 是用更低 bit-width 表示权重、激活或 KV Cache，从而降低显存占用、减少 HBM 带宽压力，并在硬件支持时提升计算吞吐。
\item \textbf{Weight-only quantization} 是 LLM 推理中最常见的方案之一，典型形式包括 W8A16 和 W4A16。
\item \textbf{Uniform quantization} 使用 scale 和 zero point 将浮点值映射到离散整数；symmetric quantization 简单硬件友好，asymmetric quantization 更适合非对称分布。
\item \textbf{Scale 粒度很关键}：per-tensor 简单但误差大，per-channel 精度较好，per-group 是 INT4 LLM 推理中的常见折中。
\item \textbf{高性能量化推理依赖 fused dequant GEMM}。如果先把低 bit 权重完整反量化为 FP16，再执行普通 GEMM，会损失许多带宽收益。
\item \textbf{PTQ 成本低，适合快速部署已有模型}；QAT 精度潜力更高，但需要训练或微调成本。
\item \textbf{GPTQ 利用近似二阶信息进行逐列量化和误差补偿}，适合低 bit weight-only PTQ。
\item \textbf{AWQ 使用 activation 统计识别重要通道}，通过等价缩放保护 salient channels，同时保持硬件友好的 weight-only quantization。
\item \textbf{SmoothQuant 通过等价变换平滑 activation outliers}，将量化难度从动态 activation 迁移到静态 weight，从而支持 W8A8 PTQ。
\item \textbf{FP8 是低精度浮点格式}，通常需要 scale、amax history 和硬件支持，适合训练和推理中的低精度 GEMM。
\item \textbf{KV Cache 量化可以显著降低长上下文和高并发 serving 的显存压力}，但需要处理 K/V 敏感性、scale 粒度和 paged block layout。
\item \textbf{量化会影响整个 serving stack}：模型格式、kernel、TP/PP 切分、PagedAttention、scheduler、质量评估和成本模型都要一起考虑。
\item \textbf{量化方案选型必须结合目标}：显存优先可选 INT4 weight-only，吞吐优先要看低精度硬件路径，质量优先则需要保守配置、敏感层保护和充分评估。
\end{itemize}

下一章将进入推理服务调度与 Continuous Batching。我们会把前面几章的 KV Cache、PagedAttention、FlashAttention 和量化技术放到一个在线 serving engine 中，分析请求如何排队、prefill 与 decode 如何混部、chunked prefill 如何降低首 token 延迟，以及调度器如何在吞吐、延迟、公平性和显存之间做决策。

\chapter{推理服务调度与 Continuous Batching}
\label{chap:serving-scheduler-continuous-batching}

\section{在线 LLM Serving 的系统目标}
\label{sec:online-llm-serving-goals}

前几章分别讨论了 KV Cache、PagedAttention、FlashAttention 和量化。它们分别解决了推理系统中的不同问题：

\begin{itemize}
\item KV Cache 让自回归 decode 不必重复计算历史 Key/Value；
\item PagedAttention 让动态请求的 KV Cache 可以按 block 管理，减少显存碎片；
\item FlashAttention 减少 attention 计算中的 HBM IO，尤其适合 prefill 和长序列场景；
\item 量化降低权重、激活或 KV Cache 的字节数，减少显存与带宽压力。
\end{itemize}

但是，一个真正的在线 LLM serving 系统不能只依赖单个 kernel 或单个显存优化。它必须面对持续到达的用户请求，在有限 GPU 资源上同时优化吞吐、延迟、公平性、稳定性和成本。

一个 LLM serving engine 的核心问题可以表述为：

\[
\text{在有限 GPU 显存和计算资源下，如何动态选择每个 iteration 要执行哪些 token？}
\]

这里的“token”既包括新请求的 prompt tokens，也包括已经运行请求的 decode token。不同 token 的成本差异很大：prefill token 通常计算密度高，长 prompt 可能一次带来大量计算；decode token 每次只生成一个，但要读取历史 KV Cache，并且对用户的流式体验非常敏感。

因此，推理服务调度不是简单的 batch 拼接，而是一个在线资源管理问题：

\[
\text{Scheduling}
\text{Queueing}
+
\text{KV Memory Management}
+
\text{Kernel Selection}
+
\text{Admission Control}
+
\text{SLA Optimization}.
\]

\subsection{请求生命周期}
\label{subsec:serving-request-lifecycle}

一个典型 LLM 请求从进入系统到完成，大致经历以下阶段：

\begin{enumerate}
\item \textbf{接收请求}：API server 接收 prompt、采样参数、最大生成长度、streaming 标志等；
\item \textbf{分词}：tokenizer 将文本转换为 token ids；
\item \textbf{入队}：请求进入 waiting queue，等待调度器分配资源；
\item \textbf{准入控制}：系统判断是否有足够 KV Cache blocks、batch token budget 和执行 slot；
\item \textbf{prefill}：处理 prompt，建立 KV Cache，得到首 token logits；
\item \textbf{decode}：循环生成新 token，每轮读取历史 KV，追加新的 K/V；
\item \textbf{采样与停止判断}：根据 logits 采样 token，判断 EOS、stop words、max tokens 或取消；
\item \textbf{流式返回}：将生成 token 或文本片段返回给用户；
\item \textbf{资源释放}：请求结束后释放 KV Cache blocks 和调度状态。
\end{enumerate}

这个流程中的关键资源包括：

\begin{itemize}
\item GPU compute；
\item GPU HBM；
\item KV Cache blocks；
\item token budget；
\item batch slots；
\item network output bandwidth；
\item CPU scheduler 和 tokenizer 线程；
\item 多副本部署中的 worker 容量。
\end{itemize}

\begin{figure}[H]
\centering
\begin{tikzpicture}[
font=\small,
box/.style={draw, rounded corners=4pt, minimum width=1.75cm, minimum height=0.65cm, align=center},
arrow/.style={-{Latex[length=2.1mm]}, thick}
]

\node[font=\bfseries] at (0,3.3) {在线 LLM 请求生命周期};

\node[box, fill=blue!8] (api) at (-5.5,1.45) {API\\Request};
\node[box, fill=green!10] (tok) at (-3.55,1.45) {Tokenizer};
\node[box, fill=orange!14] (queue) at (-1.55,1.45) {Waiting\\Queue};
\node[box, fill=yellow!16] (admit) at (0.55,1.45) {Admission\\Control};
\node[box, fill=purple!10] (prefill) at (2.75,1.45) {Prefill};
\node[box, fill=red!8] (decode) at (4.9,1.45) {Decode\\Loop};

\node[box, fill=cyan!10] (sample) at (4.9,0.25) {Sampling\\Stop Check};
\node[box, fill=gray!15] (release) at (2.75,0.25) {Release\\KV Blocks};
\node[box, fill=blue!6] (stream) at (0.55,0.25) {Stream\\Output};

\draw[arrow] (api) -- (tok);
\draw[arrow] (tok) -- (queue);
\draw[arrow] (queue) -- (admit);
\draw[arrow] (admit) -- (prefill);
\draw[arrow] (prefill) -- (decode);
\draw[arrow] (decode) -- (sample);
\draw[arrow] (sample.west) -- (stream.east);
\draw[arrow] (sample) -- (release);
\draw[arrow] (sample.east) .. controls (6.2,0.8) and (6.2,1.45) .. (decode.east);

\node[align=left, text width=10.6cm] at (0,-1.4) {
  调度器不仅决定请求何时进入 GPU，还要在每个 decode iteration 决定哪些请求继续生成。
  请求结束、取消或超时后，KV Cache blocks 必须及时释放。
};

\end{tikzpicture}
\caption{在线 LLM serving 的请求生命周期}
\label{fig:serving-request-lifecycle}
\end{figure}

\subsection{关键指标：TTFT、TPOT、吞吐与尾延迟}
\label{subsec:serving-key-metrics}

LLM serving 的指标比普通推理更复杂。因为用户不只关心整个请求什么时候结束，也关心模型什么时候开始输出，以及输出过程中是否流畅。

\textbf{TTFT, Time To First Token} 指从请求到达系统，到第一个生成 token 返回给用户的时间。它通常包括：

\[
T_{\mathrm{TTFT}}
T_{\mathrm{queue}}
+
T_{\mathrm{tokenize}}
+
T_{\mathrm{prefill}}
+
T_{\mathrm{first_sample}}
+
T_{\mathrm{network}}.
\]

TTFT 主要受排队时间和 prefill 时间影响。长 prompt 请求的 TTFT 往往较高。

\textbf{TPOT, Time Per Output Token} 指 decode 阶段生成每个 token 的平均时间。若生成 $N_{\mathrm{out}}$ 个 tokens，decode 总时间为 $T_{\mathrm{decode}}$，则：

\[
T_{\mathrm{TPOT}}
\frac{T_{\mathrm{decode}}}{N_{\mathrm{out}}}.
\]

TPOT 影响用户看到流式输出的速度。它主要受 decode batch size、KV Cache 读取、PagedAttention kernel、量化 kernel 和 GPU 利用率影响。

\textbf{E2E Latency} 指请求从到达到完成的端到端延迟：

\[
T_{\mathrm{E2E}}
T_{\mathrm{TTFT}}
+
T_{\mathrm{decode}}.
\]

\textbf{Throughput} 可以用 tokens/s 表示：

\[
\mathrm{Throughput}
\frac{\mathrm{Total\ generated\ tokens}}
{\mathrm{Wall\ clock\ time}}.
\]

也可以区分 prefill tokens/s 和 decode tokens/s。对 serving 系统来说，仅看总 tokens/s 可能误导，因为 prefill token 和 decode token 的成本不同。

\textbf{Tail Latency} 通常看 P95、P99 或 P999。在线系统中，平均延迟不够，尾延迟才决定用户体验和 SLA。一个调度策略可能提高平均吞吐，但让长请求或低优先级请求长期饥饿，导致 P99 恶化。

\begin{table}[H]
\centering
\small
\caption{LLM Serving 常见指标}
\label{tab:serving-metrics}
\begin{tabular}{p{0.20\textwidth}p{0.38\textwidth}p{0.30\textwidth}}
\toprule
\textbf{指标} & \textbf{含义} & \textbf{主要受哪些因素影响} \
\midrule
TTFT
& 从请求到达到首 token 输出的时间
& 排队、tokenizer、prefill、chunked prefill、prefix cache。 \
\midrule
TPOT
& 每个输出 token 的平均生成时间
& decode batch、KV Cache、PagedAttention、量化、GPU 利用率。 \
\midrule
E2E Latency
& 请求完整完成时间
& TTFT、输出长度、decode 速度、stop 条件。 \
\midrule
Throughput
& 系统单位时间处理 tokens 数
& batch size、continuous batching、kernel、显存容量。 \
\midrule
Tail Latency
& P95/P99/P999 延迟
& 调度公平性、长请求、资源竞争、preemption、排队。 \
\midrule
Cost/token
& 每生成 token 的成本
& GPU 价格、利用率、量化、并发、缓存命中。 \
\bottomrule
\end{tabular}
\end{table}

\subsection{吞吐与延迟的基本矛盾}
\label{subsec:throughput-latency-tradeoff}

Serving 调度的核心矛盾是吞吐与延迟。

为了提高吞吐，系统倾向于：

\begin{itemize}
\item 增大 batch；
\item 等待更多请求凑 batch；
\item 同时处理更多 decode sequences；
\item 让 GPU 尽量满载；
\item 接收更多长请求以提高资源利用。
\end{itemize}

但这些做法可能增加延迟：

\begin{itemize}
\item 等待凑 batch 会增加 queueing delay；
\item 长 prefill 会阻塞短 decode；
\item 过大 batch 会增加单 iteration 时间；
\item 过高并发会导致 KV Cache 紧张；
\item 长请求占用 blocks 过久，影响后续请求准入。
\end{itemize}

因此，调度器不是简单地“尽量多塞请求”。它需要在每个 iteration 解决一个受约束优化问题：

\[
\max
\quad
\mathrm{utility}(\mathrm{scheduled\ tokens})
\]

满足：

\[
\mathrm{KVBlocksUsed}
\le
\mathrm{KVBlocksCapacity},
\]

\[
\mathrm{BatchedTokens}
\le
\mathrm{TokenBudget},
\]

\[
\mathrm{NumSequences}
\le
\mathrm{MaxSequences},
\]

\[
\mathrm{LatencySLA}
\le
\mathrm{Target}.
\]

不同业务对 utility 的定义不同。聊天机器人可能更重视 TTFT 和流式体验；离线批量生成可能更重视吞吐；代码补全可能要求极低 TPOT；长文档总结可能更重视长上下文 prefill 能力。

\section{Prefill 与 Decode 混合调度}
\label{sec:prefill-decode-mixed-scheduling}

LLM serving 的独特难点之一，是 prefill 和 decode 的计算形态完全不同。

Prefill 一次处理多个 prompt tokens，计算密度高，容易形成较大 GEMM，也更适合 FlashAttention。Decode 每次每个请求只处理一个 token，需要读取历史 KV Cache，通常更 memory-bound。调度器必须把这两类任务放到同一个 GPU 时间线上执行。

\subsection{prefill 的计算特征}
\label{subsec:prefill-computation-characteristics}

Prefill 输入长度可能是：

\[
S_{\mathrm{prompt}}\in[1,S_{\max}].
\]

一个 prefill 请求需要计算整个 prompt 的 Transformer forward，并写入 KV Cache。Attention 中会出现：

\[
S_{\mathrm{prompt}}\times S_{\mathrm{prompt}}
\]

的 causal attention 结构。虽然 FlashAttention 避免物化完整 attention matrix，但 dense attention 的计算复杂度仍然与：

\[
O(S_{\mathrm{prompt}}^2d_h)
\]

相关。

Prefill 的特点包括：

\begin{itemize}
\item 对长 prompt 成本很高；
\item 计算密度较高；
\item 更容易吃满 GPU compute；
\item 会一次写入大量 KV Cache；
\item 直接决定 TTFT；
\item 若不切分，长 prefill 会阻塞 decode。
\end{itemize}

\subsection{decode 的计算特征}
\label{subsec:decode-computation-characteristics}

Decode 阶段每个请求每轮通常只生成一个 token。对第 $i$ 个请求，当前上下文长度为：

\[
S_i.
\]

每层 attention 读取该请求的历史 KV Cache：

\[
K_i,V_i\in\mathbb{R}^{S_i\times n_{kv}\times d_h}.
\]

Decode 的特点包括：

\begin{itemize}
\item 单请求单步计算很小；
\item 强自回归依赖；
\item 需要持续读取 KV Cache；
\item 对 TPOT 敏感；
\item batch 中请求长度不同；
\item 请求会动态结束和加入。
\end{itemize}

若 decode batch 太小，GPU 利用率低；若 decode batch 太大，单 iteration 时间上升，流式输出变慢。调度器需要维持合适的 running sequences 数量。

\subsection{prefill 阻塞 decode 的问题}
\label{subsec:prefill-blocks-decode}

如果调度器简单地按请求到达顺序执行，一个很长的 prefill 可能占用 GPU 很长时间。在这段时间内，已经处于 decode 阶段的请求无法生成下一个 token，用户会感到流式输出卡顿。

这就是 prefill-decode interference：

\[
\text{Long Prefill}
\rightarrow
\text{Decode Delay}
\rightarrow
\text{TPOT / P99 worsens}.
\]

反过来，如果调度器总是优先 decode，新请求可能长时间无法 prefill，TTFT 变差。调度器必须在 TTFT 和 TPOT 之间折中。

\begin{figure}[H]
\centering
\begin{tikzpicture}[
font=\scriptsize,
cell/.style={draw, minimum width=0.68cm, minimum height=0.42cm, align=center},
prefill/.style={fill=green!12},
decode/.style={fill=blue!12},
wait/.style={fill=gray!18}
]

\node[font=\bfseries\small] at (0,3.0) {长 Prefill 阻塞 Decode 的时间线};

\foreach \t in {0,...,9} {
  \node at ({-3.8+0.75*\t},2.25) {t\t};
}

\node at (-4.55,1.55) {GPU};
\node at (-4.55,0.7) {Decode reqs};

\foreach \t in {0,1} {
  \node[cell, decode] at ({-3.8+0.75*\t},1.55) {D};
}
\foreach \t in {2,3,4,5,6} {
  \node[cell, prefill] at ({-3.8+0.75*\t},1.55) {Long P};
}
\foreach \t in {7,8,9} {
  \node[cell, decode] at ({-3.8+0.75*\t},1.55) {D};
}

\foreach \t in {0,1,7,8,9} {
  \node[cell, decode] at ({-3.8+0.75*\t},0.7) {run};
}
\foreach \t in {2,3,4,5,6} {
  \node[cell, wait] at ({-3.8+0.75*\t},0.7) {wait};
}

\node[align=left, text width=10.5cm, font=\small] at (0,-0.85) {
  长 prompt prefill 若一次性执行，会让已经在 decode 的请求等待多个 iteration。
  用户看到的流式输出会出现明显停顿。
};

\end{tikzpicture}
\caption{长 prefill 对 decode 流式体验的阻塞}
\label{fig:long-prefill-blocks-decode}
\end{figure}

\subsection{调度器的 token budget}
\label{subsec:scheduler-token-budget}

现代 LLM serving 调度器常用 token budget 控制每个 iteration 的总工作量。设每轮最大可调度 token 数为：

\[
T_{\max}.
\]

对于 decode 请求，每个请求通常消耗 1 个 token budget；对于 prefill 请求，消耗 prompt tokens 数或 chunk tokens 数。

如果本轮调度 $N_D$ 个 decode 请求和若干 prefill chunks，则约束为：

\[
N_D
+
\sum_{r\in\mathcal{P}}
S_{r,\mathrm{chunk}}
\le
T_{\max}.
\]

同时还要满足最大序列数约束：

\[
N_{\mathrm{seq}}
\le
N_{\max}.
\]

以及 KV block 约束：

\[
B_{\mathrm{needed}}
\le
B_{\mathrm{free}}.
\]

Token budget 的作用是防止单个 iteration 太大，造成 decode 延迟抖动。它也为 prefill 和 decode 的混合提供统一资源度量。

\begin{lstlisting}[language=Python,caption={基于 token budget 的 prefill/decode 混合调度骨架},label={lst:mixed-prefill-decode-scheduler}]
def schedule_iteration(
waiting_queue,
running_requests,
block_manager,
max_num_batched_tokens,
max_num_sequences,
):
scheduled_decode = []
scheduled_prefill = []

token_budget = max_num_batched_tokens

# Decode-first policy: protect TPOT for running requests.
for req in running_requests:
    if req.is_finished:
        continue
    if token_budget <= 0:
        break
    if not block_manager.can_append_one_token(req):
        continue

    scheduled_decode.append(req)
    token_budget -= 1

# Admit new prefill requests if budget and KV blocks allow.
while waiting_queue and token_budget > 0:
    if len(running_requests) >= max_num_sequences:
        break

    req = waiting_queue.peek()
    prefill_tokens = req.prompt_len

    if prefill_tokens > token_budget:
        break

    if not block_manager.can_allocate_prefill(req, prefill_tokens):
        break

    waiting_queue.pop()
    block_manager.allocate_prefill(req, prefill_tokens)

    scheduled_prefill.append((req, prefill_tokens))
    running_requests.add(req)
    token_budget -= prefill_tokens

return scheduled_prefill, scheduled_decode

\end{lstlisting}

这个策略优先保护 decode，因此 TPOT 较稳，但长 prompt 的 TTFT 可能变差。后面我们会讨论 chunked prefill 来缓解这个问题。

\section{Continuous Batching}
\label{sec:continuous-batching}

Continuous batching 是 LLM serving 的核心调度技术之一。它的思想是：不要等一个 batch 中所有请求都完成才处理下一批，而是在每个 decode iteration 动态重组 batch，让完成请求退出，让新请求加入。

\[
\text{Static batching}
\text{batch-level scheduling}.
\]

\[
\text{Continuous batching}
\text{iteration-level scheduling}.
\]

\subsection{静态 batching 的局限}
\label{subsec:static-batching-limitations}

静态 batching 适合传统模型推理：一批请求进来，一次 forward，所有请求都结束。但 LLM 生成长度不同。如果一个 batch 内有短请求和长请求，短请求完成后会空占 batch slot，或者必须等待长请求结束才能释放资源。

假设一个 batch 有三个请求，输出长度分别为：

\[
[8,\ 64,\ 512].
\]

如果静态执行到所有请求完成，前两个请求很早结束，但第三个请求仍然运行很久。短请求的资源利用低，后续新请求也无法及时加入。

静态 batching 的问题包括：

\begin{itemize}
\item 结束请求不能及时释放 KV Cache；
\item 新请求必须等待整个 batch 完成；
\item 长请求拖慢短请求；
\item batch 内有效请求数逐渐下降；
\item GPU 利用率随时间下降；
\item 尾延迟恶化。
\end{itemize}

\subsection{iteration-level scheduling}
\label{subsec:iteration-level-scheduling}

Continuous batching 每个 iteration 做一次调度。一个 iteration 通常对应一次 decode step，也可能混入 prefill chunks。调度器在每轮开始时：

\begin{enumerate}
\item 移除已经完成或取消的请求；
\item 释放它们的 KV blocks；
\item 为 running 请求安排 decode token；
\item 从 waiting queue 中选择新请求或 prefill chunks；
\item 检查 token budget 和 KV block budget；
\item 构造本轮模型输入；
\item 执行 GPU forward；
\item 采样并更新请求状态。
\end{enumerate}

这种方式让 batch 成为动态集合：

\[
\mathcal{B}^{(t)}
\ne
\mathcal{B}^{(t+1)}.
\]

每轮 batch 都可以变化。

\begin{figure}[H]
\centering
\begin{tikzpicture}[
font=\scriptsize,
cell/.style={draw, minimum width=0.68cm, minimum height=0.42cm, align=center},
prefill/.style={fill=green!12},
decode/.style={fill=blue!12},
inactive/.style={fill=gray!18},
finish/.style={fill=red!10}
]

\node[font=\bfseries\small] at (0,3.0) {Continuous Batching：每个 iteration 动态重组 batch};

\foreach \t in {0,...,8} {
  \node at ({-3.65+0.75*\t},2.35) {iter \t};
}

\foreach \r/\y in {A/1.65,B/1.1,C/0.55,D/0.0,E/-0.55} {
  \node at (-4.45,\y) {Req \r};
}

% A
\foreach \t in {0,1,2} {
  \node[cell, decode] at ({-3.65+0.75*\t},1.65) {D};
}
\node[cell, finish] at ({-3.65+0.75*3},1.65) {end};
\foreach \t in {4,...,8} {
  \node[cell, inactive] at ({-3.65+0.75*\t},1.65) {};
}

% B
\foreach \t in {0,...,8} {
  \node[cell, decode] at ({-3.65+0.75*\t},1.1) {D};
}

% C joins
\foreach \t in {0,1} {
  \node[cell, inactive] at ({-3.65+0.75*\t},0.55) {};
}
\node[cell, prefill] at ({-3.65+0.75*2},0.55) {P};
\foreach \t in {3,...,8} {
  \node[cell, decode] at ({-3.65+0.75*\t},0.55) {D};
}

% D joins
\foreach \t in {0,...,4} {
  \node[cell, inactive] at ({-3.65+0.75*\t},0.0) {};
}
\node[cell, prefill] at ({-3.65+0.75*5},0.0) {P};
\foreach \t in {6,7,8} {
  \node[cell, decode] at ({-3.65+0.75*\t},0.0) {D};
}

% E joins at 7
\foreach \t in {0,...,6} {
  \node[cell, inactive] at ({-3.65+0.75*\t},-0.55) {};
}
\node[cell, prefill] at ({-3.65+0.75*7},-0.55) {P};
\node[cell, decode] at ({-3.65+0.75*8},-0.55) {D};

\node[align=left, text width=10.5cm, font=\small] at (0,-1.85) {
  Continuous batching 允许请求 A 结束后立即释放资源，同时请求 C、D、E 在后续 iteration 加入。
  每个 iteration 的 batch 都是调度器重新构造的。
};

\end{tikzpicture}
\caption{Continuous batching 的动态请求集合}
\label{fig:continuous-batching-dynamic-batch}
\end{figure}

\subsection{running、waiting、swapped 队列}
\label{subsec:running-waiting-swapped-queues}

一个 serving scheduler 通常维护多个状态集合：

\begin{itemize}
\item \textbf{Waiting}：请求已到达，但尚未获得 KV blocks 或执行机会；
\item \textbf{Running}：请求已被准入，拥有 KV Cache，正在 prefill 或 decode；
\item \textbf{Finished}：请求已完成，等待输出收尾和资源释放；
\item \textbf{Swapped / Paused}：请求因为资源压力被暂停或迁出；
\item \textbf{Preempted}：请求被抢占，后续需要恢复或重算。
\end{itemize}

状态转移可以表示为：

\[
\text{Waiting}
\rightarrow
\text{Running}
\rightarrow
\text{Finished}.
\]

也可能出现：

\[
\text{Running}
\rightarrow
\text{Paused}
\rightarrow
\text{Running}.
\]

或者：

\[
\text{Running}
\rightarrow
\text{Preempted}
\rightarrow
\text{Waiting}.
\]

\begin{figure}[H]
\centering
\begin{tikzpicture}[
font=\small,
state/.style={draw, rounded corners=4pt, minimum width=2.0cm, minimum height=0.75cm, align=center},
arrow/.style={-{Latex[length=2.1mm]}, thick}
]

\node[font=\bfseries] at (0,3.15) {Serving Scheduler 请求状态机};

\node[state, fill=orange!14] (waiting) at (-4.4,1.25) {Waiting};
\node[state, fill=green!10] (running) at (-1.35,1.25) {Running};
\node[state, fill=blue!8] (finished) at (1.85,1.25) {Finished};
\node[state, fill=gray!15] (released) at (4.75,1.25) {Released};

\node[state, fill=purple!10] (paused) at (-1.35,-0.25) {Paused\\Swapped};
\node[state, fill=red!8] (preempted) at (-4.4,-0.25) {Preempted};

\draw[arrow] (waiting) -- node[above] {admit} (running);
\draw[arrow] (running) -- node[above] {stop/eos} (finished);
\draw[arrow] (finished) -- node[above] {free blocks} (released);

\draw[arrow] (running) -- node[right] {memory pressure} (paused);
\draw[arrow] (paused) -- node[right] {resume} (running);

\draw[arrow] (running) -- (preempted);
\draw[arrow] (preempted) -- node[left] {requeue} (waiting);

\node[align=left, text width=10.5cm] at (0,-1.6) {
  Running 请求持有 KV Cache blocks，是显存管理的重点。
  当资源不足时，调度器可能暂停、迁出或抢占低优先级请求。
};

\end{tikzpicture}
\caption{请求在 scheduler 中的状态转移}
\label{fig:serving-scheduler-state-machine}
\end{figure}

\subsection{continuous batching 的核心伪代码}
\label{subsec:continuous-batching-pseudocode}

下面给出一个简化 continuous batching 主循环。它展示调度器、block manager、model runner 和 output processor 的关系。

\begin{lstlisting}[language=Python,caption={Continuous batching serving loop 的简化伪代码},label={lst:continuous-batching-main-loop}]
def serving_loop(engine):
while True:
# 1. Pull new requests from frontend.
new_requests = engine.frontend.poll_requests()
for req in new_requests:
engine.waiting_queue.push(req)

    # 2. Free finished or cancelled requests.
    for req in list(engine.running_requests):
        if req.is_finished() or req.is_cancelled():
            engine.block_manager.release(req)
            engine.running_requests.remove(req)
            engine.output_processor.finalize(req)

    # 3. Build an iteration-level schedule.
    schedule = engine.scheduler.schedule(
        waiting_queue=engine.waiting_queue,
        running_requests=engine.running_requests,
        block_manager=engine.block_manager,
    )

    if schedule.is_empty():
        engine.sleep_briefly()
        continue

    # 4. Prepare model inputs and KV slot mapping.
    model_inputs = engine.input_builder.build(schedule)

    # 5. Execute prefill/decode on GPU.
    model_outputs = engine.model_runner.execute(model_inputs)

    # 6. Sample next tokens and update request states.
    sampled_tokens = engine.sampler.sample(model_outputs.logits)
    engine.output_processor.process(schedule, sampled_tokens)

    # 7. Update sequence lengths and KV metadata.
    engine.scheduler.update_after_step(schedule, sampled_tokens)

\end{lstlisting}

真实系统会使用异步队列、多线程、CUDA streams、RPC、多 GPU executor 和更复杂的错误处理。但核心循环仍然类似：每轮接收请求、释放资源、调度、执行、采样、更新状态。

\section{Chunked Prefill}
\label{sec:chunked-prefill}

Chunked prefill 是解决长 prompt 阻塞 decode 的重要技术。它把长 prompt 拆成多个较小 chunks，在多个 iterations 中逐步执行，让 decode 请求可以穿插运行。

\subsection{为什么需要 chunked prefill}
\label{subsec:why-need-chunked-prefill}

如果一个 prompt 长度为 $S_{\mathrm{prompt}}=16000$，一次性 prefill 可能占用 GPU 很长时间。期间所有 decode 请求都无法生成新 token，TPOT 和 P99 会恶化。

Chunked prefill 将 prompt 切成大小不超过 $C$ 的 chunks：

\[
S_{\mathrm{prompt}}
C_1+C_2+\cdots+C_m,
\]

其中：

\[
C_i\le C_{\max}.
\]

每个 iteration 只处理一个或几个 chunks，并保留一部分 token budget 给 decode。这样长 prompt 的 TTFT 可能略有增加，但正在运行请求的 TPOT 更稳定。

\[
\text{Without chunking:}
\quad
\text{one huge prefill blocks decode}.
\]

\[
\text{With chunking:}
\quad
\text{prefill chunks interleave with decode}.
\]

\begin{figure}[H]
\centering
\begin{tikzpicture}[
font=\scriptsize,
cell/.style={draw, minimum width=0.7cm, minimum height=0.42cm, align=center},
prefill/.style={fill=green!12},
decode/.style={fill=blue!12}
]

\node[font=\bfseries\small] at (0,3.1) {Chunked Prefill：长 Prompt 被切分并与 Decode 穿插};

\foreach \t in {0,...,9} {
  \node at ({-3.8+0.75*\t},2.35) {t\t};
}

\node at (-4.65,1.55) {No chunk};
\foreach \t in {0,1,2,3,4} {
  \node[cell, prefill] at ({-3.8+0.75*\t},1.55) {P};
}
\foreach \t in {5,6,7,8,9} {
  \node[cell, decode] at ({-3.8+0.75*\t},1.55) {D};
}

\node at (-4.65,0.55) {Chunked};
\foreach \t/\op in {0/P1,1/D,2/P2,3/D,4/P3,5/D,6/P4,7/D,8/P5,9/D} {
  \ifodd\t
    \node[cell, decode] at ({-3.8+0.75*\t},0.55) {\op};
  \else
    \node[cell, prefill] at ({-3.8+0.75*\t},0.55) {\op};
  \fi
}

\node[align=left, text width=10.5cm, font=\small] at (0,-0.85) {
  Chunked prefill 将长 prompt 拆成多个 chunks，避免一次性占满 GPU。
  这样可以保护 decode 请求的流式输出体验。
};

\end{tikzpicture}
\caption{Chunked prefill 与 decode 的交错执行}
\label{fig:chunked-prefill-interleaving}
\end{figure}

\subsection{chunk 大小的选择}
\label{subsec:chunk-size-choice}

Chunk size 是一个重要参数。若 chunk 太大：

\begin{itemize}
\item 仍然会阻塞 decode；
\item 单 iteration 时间波动大；
\item 长 prompt 对 TPOT 影响明显。
\end{itemize}

若 chunk 太小：

\begin{itemize}
\item prefill 被拆得太碎；
\item kernel 启动和调度开销增加；
\item FlashAttention 的大块计算效率下降；
\item TTFT 可能变差。
\end{itemize}

设每轮 token budget 为 $T_{\max}$，decode 请求数为 $N_D$。可用于 prefill 的 budget 为：

\[
T_{\mathrm{prefill}}
T_{\max}-N_D.
\]

如果采用 decode-first 策略，则每轮 prefill chunk size 可以限制为：

\[
C
\le
T_{\max}-N_D.
\]

系统也可以保留最小 decode budget：

\[
N_D
\le
T_{\max}-C_{\min}.
\]

实际中，chunk size 可能动态变化：当 running decode 请求少时，允许更大 prefill chunk；当 decode 请求多时，缩小 chunk 以保护 TPOT。

\subsection{chunked prefill 的 KV Cache 写入}
\label{subsec:chunked-prefill-kv-write}

Chunked prefill 不是把 prompt 独立分段计算。第 $k$ 个 chunk 必须 attend 到前面已经 prefill 的 chunks。因此，chunked prefill 需要在每个 chunk 后把 K/V 写入 KV Cache，并在下一 chunk 中读取已有 KV。

设 prompt 被切为：

\[
X^{(1)},X^{(2)},\ldots,X^{(m)}.
\]

处理第 $r$ 个 chunk 时，已有 cache 长度为：

\[
S_{\mathrm{past}}
\sum_{i=1}^{r-1}|X^{(i)}|.
\]

当前 chunk 长度为 $C_r$。该 chunk 中 token 的 attention 需要看：

\[
S_{\mathrm{past}}+C_r
\]

个位置，其中前 $S_{\mathrm{past}}$ 来自 KV Cache，当前 chunk 内部也需要 causal attention。

因此 chunked prefill 的 attention 形态介于 full prefill 和 decode 之间：

\[
Q: C_r\times d_h,
\quad
K,V: (S_{\mathrm{past}}+C_r)\times d_h.
\]

它比 decode 更大，但比一次性完整 prefill 更可控。

\begin{lstlisting}[language=Python,caption={Chunked prefill 调度和 KV 写入的伪代码},label={lst:chunked-prefill-pseudocode}]
def run_chunked_prefill(req, model_runner, block_manager, chunk_size):
prompt = req.prompt_tokens
offset = req.num_prefilled_tokens

while offset < len(prompt):
    end = min(offset + chunk_size, len(prompt))
    chunk = prompt[offset:end]

    # Allocate KV slots for this chunk.
    slot_mapping = block_manager.allocate_slots(req, num_tokens=len(chunk))

    # The model attends to past KV cache plus current chunk.
    outputs = model_runner.prefill_chunk(
        request=req,
        input_tokens=chunk,
        past_length=offset,
        slot_mapping=slot_mapping,
    )

    # KV for this chunk has been written into cache.
    offset = end
    req.num_prefilled_tokens = offset

    # Scheduler may yield here to run decode requests.
    if scheduler_should_yield():
        break

return req.num_prefilled_tokens == len(prompt)

\end{lstlisting}

在实际系统中，chunked prefill 通常由 scheduler 在每个 iteration 决定执行多少 chunk，而不是在一个 Python 循环里独占 GPU。

\subsection{对 TTFT 与 TPOT 的影响}
\label{subsec:chunked-prefill-impact-ttft-tpot}

Chunked prefill 对指标的影响是双向的。

它改善：

\begin{itemize}
\item running decode 请求的 TPOT；
\item P99 流式稳定性；
\item scheduler 对长 prompt 的可控性；
\item GPU iteration 时间波动。
\end{itemize}

它可能恶化：

\begin{itemize}
\item 长 prompt 请求自身的 TTFT；
\item prefill 总 kernel 效率；
\item 调度和 metadata 开销。
\end{itemize}

因此，chunked prefill 的策略取决于业务目标。如果系统更重视正在流式输出请求的稳定性，可以更激进地 chunk；如果系统更重视新请求尽快首 token，可以允许更大 prefill chunk。

\section{KV Cache 准入控制与显存调度}
\label{sec:kv-cache-admission-control}

LLM serving 的准入控制不同于普通服务。普通服务可能主要看 CPU queue 或 QPS；LLM serving 必须看 KV Cache 显存。一个请求能否进入 running set，取决于是否有足够 cache blocks 支持它的 prompt 和未来生成。

\subsection{block 预算}
\label{subsec:block-budget}

使用 PagedAttention 时，KV Cache 被划分为 fixed-size blocks。设 block size 为 $B_{\mathrm{blk}}$。请求当前需要缓存 $S$ 个 tokens，则所需 blocks 为：

\[
N_{\mathrm{blocks}}(S)
\left\lceil
\frac{S}{B_{\mathrm{blk}}}
\right\rceil.
\]

对于新请求，prefill 需要的初始 blocks 为：

\[
N_{\mathrm{prefill}}
\left\lceil
\frac{S_{\mathrm{prompt}}}{B_{\mathrm{blk}}}
\right\rceil.
\]

Decode 每生成一个 token，若当前长度正好到达 block 边界，就需要新增一个 block：

\[
S\bmod B_{\mathrm{blk}}=0
\Rightarrow
\text{allocate new block}.
\]

调度器需要维护：

\[
B_{\mathrm{free}},
\quad
B_{\mathrm{used}},
\quad
B_{\mathrm{reserved}}.
\]

其中 reserved blocks 可以用于为 running 请求预留未来若干 decode slots，避免运行到一半发现没有 block 可追加。

\begin{lstlisting}[language=Python,caption={KV block 准入控制的简化逻辑},label={lst:kv-block-admission-control}]
def can_admit_request(req, block_manager, block_size, lookahead_tokens):
prompt_blocks = ceil_div(req.prompt_len, block_size)

# Reserve a small number of future decode slots.
future_len = req.prompt_len + lookahead_tokens
future_blocks = ceil_div(future_len, block_size)

needed_blocks = max(prompt_blocks, future_blocks)

return block_manager.num_free_blocks() >= needed_blocks

\end{lstlisting}

Lookahead tokens 的作用是给请求一些未来 decode 空间。但预留太多会降低并发；预留太少会增加运行中 block 不足的风险。

\subsection{最大生成长度与保守预留}
\label{subsec:max-generation-length-reservation}

如果为每个请求按最大生成长度预留 KV Cache，则需要：

\[
N_{\mathrm{max}}
\left\lceil
\frac{S_{\mathrm{prompt}}+S_{\mathrm{max_new}}}
{B_{\mathrm{blk}}}
\right\rceil.
\]

这种方式安全但浪费，因为很多请求不会生成到最大长度。PagedAttention 的优势之一是可以不必一次性预留最大长度，而是按需追加。

但完全不预留未来也有风险。如果系统过度接收请求，后续 decode 时 blocks 耗尽，就必须抢占、暂停或失败。因此实际系统会采用折中：

\begin{itemize}
\item 为 prompt 立即分配；
\item 为未来少量 decode 预留；
\item 到 block 边界时再追加；
\item 使用水位线控制 admission；
\item 在内存紧张时暂停低优先级请求。
\end{itemize}

可以设置高低水位线：

\[
B_{\mathrm{free}} < B_{\mathrm{low}}
\Rightarrow
\text{stop admitting new requests}.
\]

\[
B_{\mathrm{free}} > B_{\mathrm{high}}
\Rightarrow
\text{resume admitting}.
\]

这样可以避免频繁在接收和拒绝之间抖动。

\subsection{preemption、swap 与 recompute}
\label{subsec:preemption-swap-recompute}

当 KV Cache blocks 不足时，系统有几种选择：

\begin{itemize}
\item \textbf{Reject}：拒绝新请求；
\item \textbf{Queue}：让新请求继续等待；
\item \textbf{Preempt}：抢占某些 running 请求；
\item \textbf{Swap}：将某些请求的 KV Cache 迁出到 CPU 内存；
\item \textbf{Recompute}：释放 KV Cache，未来从 prompt 重新 prefill；
\item \textbf{Evict prefix cache}：驱逐可重算的共享 prefix。
\end{itemize}

这些策略的成本不同。Swap 保留状态但需要 PCIe/CPU 内存带宽；recompute 释放显存但未来恢复时需要重新 prefill；reject 最简单但影响可用性；queue 会增加 TTFT。

可以用一个成本函数比较：

\[
C_{\mathrm{evict}}
\alpha T_{\mathrm{restore}}
+
\beta M_{\mathrm{freed}}
+
\gamma P_{\mathrm{user_impact}}.
\]

实际系统通常用启发式策略，例如：

\begin{itemize}
\item 优先驱逐已结束或长时间 idle 的 prefix cache；
\item 优先暂停低优先级请求；
\item 避免抢占接近完成的请求；
\item 避免抢占已经生成很多 token 的请求；
\item 对超长请求设置单独配额。
\end{itemize}

\begin{figure}[H]
\centering
\begin{tikzpicture}[
font=\small,
box/.style={draw, rounded corners=4pt, minimum width=2.0cm, minimum height=0.7cm, align=center},
arrow/.style={-{Latex[length=2.1mm]}, thick}
]

\node[font=\bfseries] at (0,3.1) {KV Cache 不足时的资源处理策略};

\node[box, fill=red!8] (pressure) at (0,2.0) {KV Memory\\Pressure};

\node[box, fill=orange!14] (queue) at (-5.0,0.8) {Queue\\new reqs};
\node[box, fill=yellow!16] (reject) at (-2.5,0.8) {Reject};
\node[box, fill=blue!8] (preempt) at (0,0.8) {Preempt};
\node[box, fill=green!10] (swap) at (2.5,0.8) {Swap to\\CPU};
\node[box, fill=purple!10] (recompute) at (5.0,0.8) {Release and\\Recompute};

\foreach \n in {queue,reject,preempt,swap,recompute} {
  \draw[arrow] (pressure) -- (\n);
}

\node[align=left, text width=10.5cm] at (0,-0.85) {
  显存压力下没有免费选择。不同策略在用户体验、恢复时间、系统吞吐和实现复杂度之间取舍。
};

\end{tikzpicture}
\caption{KV Cache 资源不足时的处理路径}
\label{fig:kv-memory-pressure-actions}
\end{figure}

\section{调度策略：吞吐、延迟与公平性}
\label{sec:scheduling-policies-throughput-latency-fairness}

调度器的策略决定系统行为。不同策略可能使用相同 kernel 和相同 GPU，但表现出完全不同的 TTFT、TPOT、吞吐和尾延迟。

\subsection{FCFS}
\label{subsec:fcfs-scheduling}

FCFS，First-Come First-Served，是最简单的策略。请求按到达顺序进入 running set。优点是公平、易理解、实现简单。缺点是容易被长请求阻塞。

如果前面有一个超长 prompt 请求，后面的短请求必须等待，即 head-of-line blocking。对于用户交互场景，FCFS 可能导致短请求体验很差。

\subsection{decode-first}
\label{subsec:decode-first-scheduling}

Decode-first 策略优先调度已经在 running set 中的 decode 请求，再用剩余 token budget 处理 prefill。它的目标是保护流式输出稳定性：

\[
\text{minimize TPOT}.
\]

优点：

\begin{itemize}
\item 已开始输出的请求不容易卡顿；
\item 用户感知更流畅；
\item decode latency 更稳定。
\end{itemize}

缺点：

\begin{itemize}
\item 新请求 TTFT 可能变差；
\item waiting queue 在高负载下可能积压；
\item 长 prompt 更难被准入。
\end{itemize}

Decode-first 常与 chunked prefill 配合：先保护 decode，再把剩余 budget 用于 prefill chunks。

\subsection{prefill-first}
\label{subsec:prefill-first-scheduling}

Prefill-first 策略优先处理新请求 prefill，目标是降低 TTFT。它适合希望用户尽快看到首 token的场景。缺点是可能阻塞 running decode，导致流式输出卡顿。

Prefill-first 在低负载时问题不大；在高负载、长 prompt 混入时，容易导致 TPOT 和尾延迟恶化。因此生产系统中通常不会无条件 prefill-first，而是结合 budget 或优先级。

\subsection{Shortest-Remaining-Processing-Time 直觉}
\label{subsec:srpt-intuition}

在排队论中，SRPT，Shortest Remaining Processing Time，倾向于优先处理剩余工作量最小的任务，以降低平均完成时间。LLM serving 中可以借鉴这个直觉：短请求优先可以改善平均延迟。

但 LLM 请求的剩余生成长度通常未知。用户可能设置 max tokens，但模型可能提前 EOS。系统只能估计：

\[
\hat{S}_{\mathrm{remaining}}.
\]

估计可来自：

\begin{itemize}
\item 用户 max tokens；
\item 历史请求分布；
\item prompt 类型；
\item 当前已生成长度；
\item stop sequence；
\item 业务接口类型。
\end{itemize}

短请求优先的风险是长请求饥饿。因此需要 aging 或配额机制，让等待时间较长的请求逐渐提高优先级。

\subsection{优先级与 SLA}
\label{subsec:priority-and-sla}

生产系统中，请求可能有不同优先级：

\begin{itemize}
\item 付费等级；
\item 业务类型；
\item 用户交互请求 vs 离线任务；
\item 系统内部调用；
\item 实时补全 vs 批量生成；
\item 短上下文 vs 长上下文。
\end{itemize}

调度器可以为每个请求定义 priority：

\[
\pi_r.
\]

调度 utility 可写为：

\[
U_r
\pi_r
\cdot
f(T_{\mathrm{waiting}},S_{\mathrm{prompt}},S_{\mathrm{generated}},\mathrm{SLA}).
\]

高优先级请求获得更早 prefill、更稳定 decode 或更多资源配额。但优先级系统必须防止低优先级请求永久饥饿。常见做法是引入 waiting time aging：

\[
\pi_r^{\mathrm{effective}}
\pi_r
+
\lambda T_{\mathrm{waiting}}.
\]

等待越久，有效优先级越高。

\begin{table}[H]
\centering
\small
\caption{常见 LLM Serving 调度策略对比}
\label{tab:serving-scheduling-policies}
\begin{tabular}{p{0.18\textwidth}p{0.32\textwidth}p{0.38\textwidth}}
\toprule
\textbf{策略} & \textbf{优点} & \textbf{风险} \
\midrule
FCFS
& 简单公平，易实现
& 长请求造成 head-of-line blocking。 \
\midrule
Decode-first
& TPOT 稳定，流式体验好
& 新请求 TTFT 可能变差。 \
\midrule
Prefill-first
& 新请求首 token 更快
& 长 prefill 可能阻塞 decode。 \
\midrule
Shortest-first
& 平均完成时间较好
& 长请求可能饥饿，需要 aging。 \
\midrule
Priority/SLA
& 支持业务分级和延迟目标
& 复杂度高，低优先级需防饥饿。 \
\bottomrule
\end{tabular}
\end{table}

\section{调度器与底层优化的协同}
\label{sec:scheduler-and-low-level-optimizations}

调度器不是孤立模块。它必须了解底层 kernel、量化格式、KV Cache 管理和多 GPU 并行，否则会做出看似合理但性能很差的决策。

\subsection{与 PagedAttention 的协同}
\label{subsec:scheduler-with-pagedattention}

PagedAttention 让请求的 KV Cache 不必连续存储，调度器因此可以更灵活地动态加入和退出请求。Scheduler 需要与 Block Manager 交互：

\begin{itemize}
\item prefill 前申请 blocks；
\item decode 到 block 边界时申请新 block；
\item 请求结束释放 blocks；
\item prefix cache 命中时共享 blocks；
\item copy-on-write 时申请新 block；
\item 显存压力下触发 eviction 或 preemption。
\end{itemize}

每个 iteration，调度器会构造：

\begin{itemize}
\item block table；
\item slot mapping；
\item sequence lengths；
\item request ids；
\item position ids。
\end{itemize}

这些 metadata 传给 PagedAttention kernel，用于读取和写入 KV Cache。

\subsection{与 FlashAttention 的协同}
\label{subsec:scheduler-with-flashattention}

Prefill 阶段常使用 FlashAttention。调度器影响 FlashAttention 的 shape：

\[
\text{batched prompt tokens}
\rightarrow
\text{attention kernel shape}.
\]

如果调度器把 prefill chunk 切得太小，FlashAttention 可能无法充分发挥大块计算效率；如果切得太大，又会阻塞 decode。调度器需要在 kernel efficiency 和 latency 之间折中。

对于 variable-length prompts，调度器还要生成 packed sequence metadata，使 FlashAttention 可以处理不同长度序列，避免 padding 到最大长度。

\subsection{与量化的协同}
\label{subsec:scheduler-with-quantization}

量化改变资源模型。INT4 weight-only 降低权重显存和权重读取带宽，可能允许：

\begin{itemize}
\item 更大 KV Cache budget；
\item 更高 running request 数；
\item 更长上下文；
\item 更少 GPU 副本；
\item 更低成本。
\end{itemize}

但量化 kernel 可能对 batch size、group size、alignment、TP shard 有要求。调度器若构造出不友好的 shape，量化 kernel 性能可能下降。比如某些 INT4 GEMM 对 batch 或 hidden 维对齐敏感；某些 FP8 kernel 需要特定 scale metadata；KV Cache 量化需要 attention kernel 读取 scale。

因此，调度器的资源估算不能只使用 token 数，还需要考虑 kernel backend 的 shape 约束。

\subsection{统一执行图}
\label{subsec:unified-serving-execution-graph}

一个现代 serving engine 的每轮执行可以抽象为：

\[
\text{Schedule}
\rightarrow
\text{Build Inputs}
\rightarrow
\text{Quantized Linear}
\rightarrow
\text{Attention Backend}
\rightarrow
\text{KV Write}
\rightarrow
\text{Sampling}
\rightarrow
\text{State Update}.
\]

其中 Attention Backend 可能根据阶段选择：

\[
\text{Prefill}
\rightarrow
\text{FlashAttention}.
\]

\[
\text{Decode}
\rightarrow
\text{PagedAttention}.
\]

Linear backend 可能根据模型选择：

\[
\text{FP16/BF16 GEMM},
\quad
\text{INT4 weight-only GEMM},
\quad
\text{W8A8 GEMM},
\quad
\text{FP8 GEMM}.
\]

\begin{figure}[H]
\centering
\begin{tikzpicture}[
font=\small,
box/.style={draw, rounded corners=4pt, minimum width=2.05cm, minimum height=0.7cm, align=center},
arrow/.style={-{Latex[length=2.1mm]}, thick}
]

\node[font=\bfseries] at (0,3.25) {调度器与底层优化的统一执行路径};

\node[box, fill=orange!14] (sched) at (-5.0,1.45) {Scheduler};
\node[box, fill=green!10] (inputs) at (-2.65,1.45) {Input Builder\\slot mapping};
\node[box, fill=blue!8] (linear) at (-0.2,1.45) {Quantized\\Linear};
\node[box, fill=purple!10] (attn) at (2.25,1.45) {Attention\\Backend};
\node[box, fill=red!8] (kv) at (4.7,1.45) {KV Write\\Blocks};

\node[box, fill=yellow!16] (sample) at (2.25,0.15) {Sampling};
\node[box, fill=gray!15] (update) at (-0.2,0.15) {State\\Update};

\draw[arrow] (sched) -- (inputs);
\draw[arrow] (inputs) -- (linear);
\draw[arrow] (linear) -- (attn);
\draw[arrow] (attn) -- (kv);
\draw[arrow] (kv) -- (sample);
\draw[arrow] (sample) -- (update);
\draw[arrow] (update.west) .. controls (-2.5,-0.7) and (-4.9,0.1) .. (sched.south);

\node[align=left, text width=10.5cm] at (0,-1.4) {
  调度器决定 shape 和资源分配；底层 kernel 决定每个 shape 的执行效率。
  高性能 serving 必须让 scheduler、PagedAttention、FlashAttention 和量化 kernel 协同设计。
};

\end{tikzpicture}
\caption{Serving iteration 的统一执行路径}
\label{fig:serving-unified-execution-path}
\end{figure}

\section{多 GPU 与多副本 Serving 调度}
\label{sec:multi-gpu-and-replica-serving-scheduling}

单 GPU serving 已经复杂，多 GPU 和多副本部署会引入新的调度问题。

\subsection{Replica routing}
\label{subsec:replica-routing}

生产系统通常部署多个模型副本。Router 需要把请求分配给某个 replica。简单策略包括：

\begin{itemize}
\item round-robin；
\item least outstanding requests；
\item least KV memory used；
\item least estimated latency；
\item consistent hashing by session id；
\item priority-aware routing。
\end{itemize}

LLM serving 中，least outstanding requests 不一定好，因为一个短请求和一个长请求成本差异巨大。更合理的是估计每个 worker 的负载：

\[
\operatorname{Load}
\alpha N_{\mathrm{running}}
+
\beta B_{\mathrm{KV\ used}}
+
\gamma T_{\mathrm{queued}}
+
\delta C_{\mathrm{prefill\ pending}}.
\]

Router 根据估计负载选择 replica。

\subsection{Sticky session 与 KV locality}
\label{subsec:sticky-session-and-kv-locality}

多轮对话中，KV Cache 通常驻留在某个 worker 上。如果下一轮请求路由到同一个 worker，可以复用会话 cache 或 prefix cache；如果路由到其他 worker，就需要重新 prefill 或迁移 KV Cache。

Sticky session 策略：

\[
\text{session id}
\rightarrow
\text{same worker}.
\]

优点是 cache locality 好，TTFT 低。缺点是负载可能不均。某些 worker 可能积累大量长会话，显存紧张；另一些 worker 空闲。

因此 routing 要在 cache locality 和 load balancing 之间折中：

\[
\text{RouteScore}
-\alpha\operatorname{Load}
+
\beta\operatorname{CacheHit}
----------------------------
\gamma\operatorname{KVPressure}.
\]

\subsection{Tensor Parallel Serving}
\label{subsec:tensor-parallel-serving-scheduling}

当单 GPU 放不下模型，或希望提升单请求吞吐时，会使用 tensor parallel。一个 replica 可能由多个 GPU 组成：

\[
P_{\mathrm{TP}}>1.
\]

此时一个请求必须被同一个 TP group 内的所有 ranks 共同执行。Scheduler 的资源单位不再是单 GPU，而是一个 TP group。TP group 内部需要同步执行每个 iteration，不能某个 rank 跑请求 A，另一个 rank 跑请求 B。

TP serving 的调度影响包括：

\begin{itemize}
\item KV Cache 按 heads 或 hidden shard 分布在 TP ranks；
\item block allocation 要在 TP ranks 上一致；
\item PagedAttention block table 需要每个 rank 对应本地 shard；
\item sampling logits 可能需要 gather 或只在某个 rank 完成；
\item all-reduce/all-gather 通信影响 TPOT。
\end{itemize}

\subsection{Prefill/Decode 分离部署}
\label{subsec:prefill-decode-disaggregation}

更高级的 serving 架构可能将 prefill 和 decode 分离到不同 GPU 池。原因是二者资源特征不同：

\begin{itemize}
\item prefill 更 compute-bound，适合大块计算；
\item decode 更 memory-bound，受 KV Cache 和 HBM 带宽影响；
\item prefill 延迟影响 TTFT；
\item decode 资源决定持续输出吞吐。
\end{itemize}

Prefill/decode disaggregation 的挑战是：prefill 完成后生成的 KV Cache 需要转移到 decode worker，或者使用共享存储/高速网络让 decode worker 访问。KV transfer 可能成为新瓶颈。

迁移成本约为：

\[
M_{\mathrm{KV,prompt}}
2L S_{\mathrm{prompt}} n_{kv}d_hs_{\mathrm{dtype}}.
\]

如果 prompt 很长，KV transfer 可能非常大。因此 disaggregation 需要高带宽网络、KV 压缩、分层缓存或特定路由策略支持。

\begin{figure}[H]
\centering
\begin{tikzpicture}[
font=\small,
pool/.style={draw, rounded corners=5pt, minimum width=3.2cm, minimum height=1.3cm, align=center},
arrow/.style={-{Latex[length=2.1mm]}, thick}
]

\node[font=\bfseries] at (0,3.1) {Prefill / Decode 分离部署的基本形态};

\node[pool, fill=green!10] (prefill) at (-3.2,1.25) {Prefill Pool\\compute-oriented};
\node[pool, fill=blue!8] (decode) at (3.2,1.25) {Decode Pool\\KV / memory-oriented};

\node[pool, fill=orange!14, minimum width=2.8cm] (router) at (0,-0.45) {Router\\Scheduler};

\draw[arrow] (router) -- (prefill);
\draw[arrow] (prefill) -- node[above] {KV transfer} (decode);
\draw[arrow] (decode) -- (router);

\node[align=left, text width=10.5cm] at (0,-1.85) {
  Prefill/decode 分离能针对两类工作负载分别优化资源，但 KV Cache 转移成本很高。
  是否值得取决于 prompt 长度、网络带宽、worker 负载和缓存策略。
};

\end{tikzpicture}
\caption{Prefill 与 decode 分离的 serving 架构}
\label{fig:prefill-decode-disaggregation}
\end{figure}

\section{监控、压测与调优}
\label{sec:serving-monitoring-benchmarking-tuning}

Serving 系统的调度策略必须靠数据验证。没有监控和压测，很难判断瓶颈到底在 tokenizer、scheduler、KV Cache、attention kernel、GEMM、采样、网络还是路由。

\subsection{核心监控指标}
\label{subsec:serving-monitoring-metrics}

推荐监控以下指标：

\begin{itemize}
\item QPS；
\item input tokens/s；
\item output tokens/s；
\item TTFT P50/P90/P99；
\item TPOT P50/P90/P99；
\item E2E latency；
\item waiting queue length；
\item running request count；
\item GPU utilization；
\item HBM used/free；
\item KV blocks used/free；
\item block allocation failures；
\item preemption/swap/recompute 次数；
\item prefix cache hit rate；
\item prefill/decode token ratio；
\item average prompt length；
\item average output length；
\item per-iteration batch size；
\item per-iteration token count；
\item kernel time breakdown。
\end{itemize}

这些指标应该按 worker、model、priority、tenant、prompt length bucket 分维度统计。全局平均值常常掩盖问题。例如整体 TPOT 很好，但长上下文请求 P99 很差；整体显存充足，但某些 worker 因 sticky session 显存耗尽。

\subsection{负载模型}
\label{subsec:serving-load-model}

压测时需要构造真实负载模型。至少包括：

\begin{itemize}
\item 请求到达过程，例如 Poisson、burst、trace replay；
\item prompt length 分布；
\item output length 分布；
\item 并发用户数；
\item streaming 开关；
\item 多轮会话比例；
\item prefix 共享比例；
\item 取消请求比例；
\item 不同优先级混合；
\item 长上下文请求比例。
\end{itemize}

只用固定 prompt length 和固定 output length 压测，会严重低估真实调度难度。真实流量的长尾分布才是 tail latency 和 KV Cache 压力的主要来源。

\subsection{调优参数}
\label{subsec:serving-tuning-knobs}

常见调优参数包括：

\begin{itemize}
\item \texttt{max_num_batched_tokens}；
\item \texttt{max_num_seqs}；
\item prefill chunk size；
\item KV block size；
\item GPU memory utilization target；
\item lookahead slots；
\item prefix cache capacity；
\item scheduler policy；
\item decode-first ratio；
\item request timeout；
\item max model length；
\item quantization format；
\item tensor parallel size；
\item tokenizer worker 数；
\item router load balancing policy。
\end{itemize}

调优时应一次只改少量参数，并观察 TTFT、TPOT、吞吐和 KV block 使用率的变化。一个参数可能改善吞吐但恶化尾延迟。例如增大 \texttt{max_num_batched_tokens} 可以提高 GPU 利用率，但可能让单 iteration 更长，TPOT P99 上升。

\section{本章小结}
\label{sec:chapter16-summary}

本章系统介绍了 LLM 推理服务调度与 Continuous Batching。至此，第四篇从 KV Cache、PagedAttention、FlashAttention、量化一路走到完整 serving engine 的在线调度问题。

\begin{itemize}
\item \textbf{LLM serving 是在线资源管理问题}，不仅是模型 forward。系统需要同时管理请求队列、KV Cache、token budget、GPU kernel、采样和流式输出。
\item \textbf{请求生命周期} 包括接收、分词、入队、准入、prefill、decode、采样、流式返回和资源释放。
\item \textbf{核心指标} 包括 TTFT、TPOT、E2E latency、throughput、tail latency 和 cost/token。不同业务对这些指标权重不同。
\item \textbf{Prefill 与 decode 的计算形态不同}：prefill 计算密度高、长 prompt 成本大；decode token-by-token，强依赖 KV Cache，通常更 memory-bound。
\item \textbf{长 prefill 会阻塞 decode}，导致流式输出卡顿。调度器需要在 TTFT 和 TPOT 之间折中。
\item \textbf{Continuous batching 是 iteration-level scheduling}，每个 decode iteration 都可以让完成请求退出、新请求加入，从而提高 GPU 利用率。
\item \textbf{Scheduler 通常维护 waiting、running、finished、paused 等状态集合}，并在每轮根据 token budget、KV block budget 和序列数限制构造 batch。
\item \textbf{Chunked prefill 将长 prompt 拆成多个 chunks}，让 prefill 可以与 decode 交错执行，保护 TPOT 和尾延迟。
\item \textbf{KV Cache 准入控制是 serving 的关键}。请求能否进入 running set，不只取决于 batch slot，也取决于是否有足够 KV blocks。
\item \textbf{显存压力下可以 queue、reject、preempt、swap 或 recompute}，不同策略在吞吐、延迟、实现复杂度和用户体验之间权衡。
\item \textbf{调度策略没有唯一最优}。FCFS 简单但容易被长请求阻塞；decode-first 保护流式输出；prefill-first 降低首 token；priority/SLA 策略适合多租户生产系统。
\item \textbf{调度器必须与 PagedAttention、FlashAttention 和量化 kernel 协同}。调度器决定 shape 和资源分配，底层 kernel 决定每个 shape 的执行效率。
\item \textbf{多副本部署需要 replica routing 和 sticky session 权衡}。KV locality 能降低多轮对话 TTFT，但可能造成 worker 负载不均。
\item \textbf{多 GPU serving 中，一个 TP group 是调度资源单位}，KV Cache、block table、sampling 和通信都必须跨 ranks 协调。
\item \textbf{监控与压测是调度调优的基础}。必须观察 TTFT/TPOT 分位数、KV blocks、running requests、prefix hit rate、kernel breakdown 等指标。
\end{itemize}

至此，第四篇“推理架构优化”完成。我们已经从 KV Cache 的基本机制出发，讨论了 PagedAttention 如何解决动态 cache 管理，FlashAttention 如何减少 attention IO，量化如何降低显存与带宽压力，最后回到 serving scheduler 如何把所有这些优化组合成可运行的在线推理系统。

下一篇将进入硬件与资源调度。我们会从 GPU 架构、HBM、NVLink、PCIe、RDMA、NCCL 与集群网络开始，进一步理解为什么 AI Infra 的性能最终取决于硬件拓扑、通信路径、资源隔离、调度策略和故障恢复。


% ==================================================
% 第五篇：集群网络与硬件调度 Hardware & Networking
% ==================================================
\part{集群网络与硬件调度：AI Infra 的物理世界}

\chapter{GPU 架构底座}
\label{chap:gpu-architecture}

在前面的章节中，我们已经从模型结构、分布式训练、推理服务和调度系统的角度理解了 AI Infrastructure 的核心问题：模型越来越大，序列越来越长，请求越来越多，而系统仍然必须在有限的硬件资源上完成训练与推理。到了第五篇，我们开始进入 AI Infra 的物理世界：GPU、显存、互联网络、RDMA、Kubernetes 调度，以及这些底层资源如何共同决定一个大模型系统的上限。

本章讨论 GPU 架构底座。这里的“底座”并不是指某一种具体型号的 GPU，而是指现代 AI GPU 反复出现的一组基础结构：SM、CUDA Core、Tensor Core、Register File、Shared Memory、L1/L2 Cache、HBM，以及围绕它们建立起来的性能分析方法。对于 AI Infra 工程师来说，理解 GPU 架构的目标不是记住某张产品规格表，而是能够回答下面这些工程问题：

\begin{itemize}
  \item 为什么同样是矩阵乘法，某些 kernel 可以接近峰值吞吐，而某些 kernel 只有很低利用率？
  \item 为什么 decode 阶段的 LLM 推理经常是 memory-bound，而训练中的大 GEMM 更容易是 compute-bound？
  \item 为什么 batch size、sequence length、hidden size、head dimension 的改变会影响 GPU 利用率？
  \item 为什么有些优化需要减少 HBM 读写，而不是减少理论 FLOPs？
  \item 当 profiler 显示 occupancy 不高、memory throughput 不高、SM efficiency 也不高时，应该从哪里开始排查？
\end{itemize}

如果说前几章讨论的是“模型如何被切分、调度和服务”，那么本章讨论的是“每一次矩阵乘法、每一次 attention、每一次归约最终如何落到 GPU 硬件上”。这也是后续理解 NVLink、NVSwitch、RDMA 与 AI 集群调度的前置知识。

\section{GPU 硬件结构总览}
\label{sec:gpu-hardware-overview}

GPU 的设计目标与 CPU 明显不同。CPU 强调低延迟、复杂控制流、强大的分支预测和较大的缓存层级；GPU 强调高吞吐、海量线程、数据并行以及对大规模向量和矩阵运算的持续供给。对深度学习来说，GPU 的价值并不只在于“核心多”，而在于它把大量算术单元、寄存器、片上 SRAM、显存带宽和矩阵专用计算单元组织成了适合张量计算的执行机器。

从抽象上看，一个现代数据中心 GPU 可以被分成几个层级：

\begin{enumerate}
  \item GPU 芯片整体，包含多个 GPC、TPC 或类似的高层组织单元，不同架构命名不同；
  \item 多个 SM，即 Streaming Multiprocessor，是 CUDA 程序最重要的执行单元；
  \item SM 内部包含 warp scheduler、dispatch unit、CUDA Core、Tensor Core、Load/Store Unit、Special Function Unit、Register File、Shared Memory / L1 Cache 等；
  \item 多个 SM 共享更大的 L2 Cache；
  \item GPU 通过显存控制器访问 HBM；
  \item GPU 通过 PCIe、NVLink 等互联接口与 CPU、其他 GPU 或网络设备交互。
\end{enumerate}

对于本书的目标读者来说，最重要的不是画出某一代 GPU 的精确芯片平面图，而是建立一个稳定的 mental model：一个 CUDA kernel 启动后，线程块被调度到 SM；线程块中的线程被组织成 warp；warp 中的指令由 SM 内部的调度器发射到不同执行单元；数据在 register、shared memory、L2 cache 和 HBM 之间移动；性能最终受计算吞吐、内存带宽、访存模式、并行度、同步开销与指令调度共同影响。

\begin{figure}[H]
  \centering
  \begin{tikzpicture}[
    font=\small,
    node distance=0.8cm,
    box/.style={draw, rounded corners, minimum width=2.6cm, minimum height=0.9cm, align=center},
    smallbox/.style={draw, rounded corners, minimum width=2.2cm, minimum height=0.75cm, align=center},
    arrow/.style={-{Stealth[length=2mm]}, thick}
  ]
    \node[box, minimum width=12cm, minimum height=7.3cm] (gpu) {};
    \node[anchor=north west] at ($(gpu.north west)+(0.2,-0.2)$) {GPU Device};

    \node[box, minimum width=8.5cm, minimum height=4.6cm, anchor=north] (sms) at ($(gpu.north)+(0,-0.8)$) {};
    \node[anchor=north west] at ($(sms.north west)+(0.2,-0.2)$) {多个 SM};

    \node[smallbox] (sm0) at ($(sms.center)+(-2.8,0.9)$) {SM 0};
    \node[smallbox] (sm1) at ($(sms.center)+(0,0.9)$) {SM 1};
    \node[smallbox] (sm2) at ($(sms.center)+(2.8,0.9)$) {SM 2};

    \node[smallbox] (sm3) at ($(sms.center)+(-2.8,-0.9)$) {SM 3};
    \node[smallbox] (sm4) at ($(sms.center)+(0,-0.9)$) {SM 4};
    \node[smallbox] (smn) at ($(sms.center)+(2.8,-0.9)$) {$\cdots$};

    \node[box, minimum width=9.2cm] (l2) at ($(gpu.south)+(0,2.0)$) {L2 Cache};
    \node[box, minimum width=9.2cm] (hbm) at ($(gpu.south)+(0,0.8)$) {HBM / Global Memory};

    \draw[arrow] (sms.south) -- (l2.north);
    \draw[arrow] (l2.south) -- (hbm.north);

    \node[box, minimum width=2.4cm] (pcie) at ($(gpu.east)+(-1.7,-1.0)$) {PCIe /\\NVLink};
    \draw[arrow] (hbm.east) -- (pcie.west);
  \end{tikzpicture}
  \caption{GPU 结构的工程抽象：多个 SM 共享 L2 Cache，并通过显存控制器访问 HBM。}
  \label{fig:gpu-device-overview}
\end{figure}

\subsection{SM：GPU 的基本执行单元}
\label{subsec:sm}

SM 是 Streaming Multiprocessor 的缩写。对 CUDA 程序而言，SM 可以理解为“真正执行线程块的地方”。当我们启动一个 kernel 时，会指定 grid 和 block 的形状。一个 grid 中包含多个 block，每个 block 被分配到某一个 SM 上执行。一个 SM 可以同时驻留多个 block，前提是这些 block 所需的线程数、寄存器数量、shared memory 数量和硬件限制允许。

SM 内部并不是一个简单的大核心，而是一组协同工作的子单元。典型结构包括：

\begin{itemize}
  \item \textbf{Warp Scheduler}：负责选择就绪的 warp，并将指令发射出去。
  \item \textbf{CUDA Core}：执行标量或向量化的整数、浮点运算，是传统 CUDA 计算的主要执行单元。
  \item \textbf{Tensor Core}：执行矩阵乘加相关指令，是深度学习训练与推理的关键硬件。
  \item \textbf{Load/Store Unit}：负责访存指令，包括从 global memory、shared memory、local memory 等位置读写数据。
  \item \textbf{Special Function Unit}：负责某些特殊函数或近似函数，例如指数、倒数、三角函数等，具体能力取决于架构。
  \item \textbf{Register File}：每个线程最常用、最快的私有存储空间。
  \item \textbf{Shared Memory / L1 Cache}：线程块内共享的片上存储空间，常用于 tiling、数据复用和访存合并。
\end{itemize}

从软件角度看，SM 的核心特征是“大量 warp 的交错执行”。一个 warp 遇到 global memory load 时，可能需要较长时间才能拿到数据。GPU 并不会像 CPU 那样主要依赖复杂的乱序执行和大缓存来隐藏延迟，而是切换到其他就绪 warp 执行。只要 SM 上有足够多的活跃 warp，访存延迟就可能被隐藏。

这也是 occupancy 概念的来源。Occupancy 通常表示一个 SM 上实际活跃 warp 数量与理论最大活跃 warp 数量的比例。高 occupancy 有助于隐藏延迟，但并不等于高性能。对于一些高度优化的 GEMM kernel，单个线程块可能使用大量 register 和 shared memory，导致 occupancy 不高，但每个 warp 的计算效率和数据复用非常好，最终仍然可以获得高吞吐。反过来，一个 occupancy 很高的 kernel 如果访存完全不合并、指令依赖严重或共享内存 bank conflict 很多，性能仍可能很差。

\begin{figure}[H]
  \centering
  \begin{tikzpicture}[
    font=\small,
    node distance=0.55cm,
    box/.style={draw, rounded corners, minimum width=2.4cm, minimum height=0.72cm, align=center},
    wide/.style={draw, rounded corners, minimum width=7.8cm, minimum height=0.75cm, align=center},
    arrow/.style={-{Stealth[length=2mm]}, thick}
  ]
    \node[draw, rounded corners, minimum width=10.8cm, minimum height=7.2cm] (sm) {};
    \node[anchor=north west] at ($(sm.north west)+(0.2,-0.2)$) {SM 内部抽象结构};

    \node[wide] (scheduler) at ($(sm.north)+(0,-1.0)$) {Warp Scheduler / Dispatch};
    \node[box] (cuda) at ($(scheduler.south)+(-3.2,-1.0)$) {CUDA Core};
    \node[box] (tensor) at ($(scheduler.south)+(0,-1.0)$) {Tensor Core};
    \node[box] (lsu) at ($(scheduler.south)+(3.2,-1.0)$) {Load / Store};

    \node[box] (sfu) at ($(cuda.south)+(0,-1.0)$) {SFU};
    \node[box] (reg) at ($(tensor.south)+(0,-1.0)$) {Register File};
    \node[box] (shared) at ($(lsu.south)+(0,-1.0)$) {Shared Memory\\L1 Cache};

    \node[wide] (l2) at ($(sm.south)+(0,0.75)$) {连接到 L2 Cache / HBM};

    \draw[arrow] (scheduler.south) -- (cuda.north);
    \draw[arrow] (scheduler.south) -- (tensor.north);
    \draw[arrow] (scheduler.south) -- (lsu.north);
    \draw[arrow] (cuda.south) -- (sfu.north);
    \draw[arrow] (tensor.south) -- (reg.north);
    \draw[arrow] (lsu.south) -- (shared.north);
    \draw[arrow] (shared.south) -- (l2.north);
  \end{tikzpicture}
  \caption{SM 内部的抽象结构：调度器将 warp 指令发射到计算单元、访存单元和片上存储结构。}
  \label{fig:sm-internal-abstract}
\end{figure}

\subsection{CUDA Core：通用并行计算单元}
\label{subsec:cuda-core}

CUDA Core 是 GPU 上执行常规算术指令的重要单元。早期深度学习训练中，大量矩阵乘法和卷积最终会被映射为 CUDA Core 上的浮点乘加运算。随着 Tensor Core 的出现，主流深度学习 GEMM 的峰值吞吐更多来自 Tensor Core，但 CUDA Core 仍然非常重要，因为并不是所有计算都能转化为规整的矩阵乘加。

例如，下面这些操作常常依赖 CUDA Core、Load/Store Unit、SFU 或其他非 Tensor Core 路径：

\begin{itemize}
  \item LayerNorm / RMSNorm 中的均值、方差或平方和归约；
  \item Softmax 中的 max、exp、sum、除法；
  \item 采样中的 top-k、top-p、random sampling；
  \item embedding lookup、position embedding、RoPE 变换；
  \item 激活函数，如 GELU、SiLU、ReLU；
  \item 各类 elementwise operator 和数据重排；
  \item attention 中除矩阵乘法之外的 mask、scale、概率归一化和写回。
\end{itemize}

这意味着 AI Infra 工程师不能只关注 Tensor Core 峰值。一个 Transformer block 看似由大 GEMM 主导，但在小 batch、长序列、decode 或低精度推理场景中，非 GEMM 部分可能成为瓶颈。尤其在 LLM serving 中，每一步 decode 只生成一个 token，大 GEMM 的矩阵形状可能变得“瘦长”或 batch 太小，导致 Tensor Core 利用率下降，而访存、调度和小 kernel 启动开销变得突出。

\subsection{Tensor Core：深度学习吞吐的关键}
\label{subsec:tensor-core}

Tensor Core 是为矩阵乘加设计的专用计算单元。深度学习中的大部分线性层、attention projection、FFN up/down projection、卷积、部分 batched GEMM 都可以映射到矩阵乘法：

\[
  \mat{C} = \mat{A}\mat{B} + \mat{C}.
\]

更底层地看，Tensor Core 通常执行的是小 tile 上的矩阵乘加。例如把一个大矩阵乘法拆成许多小的 $m \times n \times k$ tile，每个 tile 由 warp 或 warp group 协同执行。具体 tile 形状和支持的数据类型会随 GPU 架构变化，但抽象规律稳定：Tensor Core 通过一次指令完成一小块矩阵的乘加，并把结果累加到寄存器中的 accumulator。

在 AI 训练和推理中，Tensor Core 的价值体现在三点：

\begin{enumerate}
  \item \textbf{提升矩阵乘法吞吐}：相比使用通用 CUDA Core 做逐元素乘加，Tensor Core 可以在单位周期内完成更多乘加操作。
  \item \textbf{支持混合精度}：输入可以是 FP16、BF16、TF32、FP8 或整数低精度格式，累加通常使用更高精度或专门设计的累加路径。
  \item \textbf{推动软件栈重构}：cuBLAS、cuDNN、CUTLASS、Triton、PyTorch Inductor、FlashAttention 等优化都围绕 tile 化、数据复用和 Tensor Core 指令展开。
\end{enumerate}

不过 Tensor Core 并不会自动让所有算子变快。要有效利用 Tensor Core，矩阵形状、数据类型、内存布局、对齐方式、tile 大小和调度策略都必须合适。常见问题包括：

\begin{itemize}
  \item 矩阵维度太小，无法提供足够并行度；
  \item $M$、$N$、$K$ 维度不符合高效 tile 的对齐习惯；
  \item 输入数据类型无法走 Tensor Core 快路径；
  \item 数据布局导致额外 transpose、copy 或不合并访存；
  \item kernel 被内存带宽限制，即使 Tensor Core 很快也吃不饱。
\end{itemize}

\subsection{Register File：最快但最稀缺的存储}
\label{subsec:register-file}

Register 是每个线程私有的最快存储空间。一个 CUDA kernel 中的局部变量、临时 accumulator、索引变量和部分编译器生成的中间值都会占用寄存器。寄存器访问速度极快，但数量有限。一个 kernel 如果每个线程使用过多寄存器，会降低每个 SM 能同时驻留的 warp 或 block 数量，从而降低 occupancy。

寄存器压力是 GPU kernel 优化中的常见主题。比如在一个 tiled GEMM 中，每个线程可能需要在寄存器中维护多个 accumulator，用于保存当前 tile 的部分结果。accumulator 越多，计算复用越好，但寄存器占用越高。寄存器不够时，编译器可能将部分变量 spill 到 local memory。所谓 local memory 并不是片上本地存储，而通常位于 global memory 地址空间中，访问代价很高。一旦发生严重 register spilling，kernel 性能可能急剧下降。

因此，在性能分析时，register 的关键问题不是“越多越好”或“越少越好”，而是平衡：

\begin{itemize}
  \item 寄存器多，可以保存更多中间结果，提高计算密度和数据复用；
  \item 寄存器少，可以让更多 warp 同时驻留，提高延迟隐藏能力；
  \item 一旦寄存器使用超过阈值导致 spilling，性能通常会明显恶化；
  \item 不同 kernel 的最优点不同，必须结合 profiler 和实际输入 shape 判断。
\end{itemize}

\subsection{Shared Memory / L1 Cache：片上数据复用}
\label{subsec:shared-memory-l1}

Shared Memory 是线程块内共享的片上存储空间，延迟和带宽通常远优于 HBM。它最常见的用途是 tiling：把 global memory 中的一小块数据加载到 shared memory，然后让线程块内多个线程重复使用，从而减少 HBM 访问次数。

以矩阵乘法为例，计算 $\mat{C} = \mat{A}\mat{B}$ 时，$\mat{A}$ 的一个 tile 会被多个输出元素复用，$\mat{B}$ 的一个 tile 也会被多个输出元素复用。如果每次乘加都直接从 HBM 读取 $\mat{A}$ 和 $\mat{B}$，带宽需求会非常高。更好的方法是：

\begin{enumerate}
  \item 线程块协同从 HBM 读取 $\mat{A}$ 和 $\mat{B}$ 的 tile；
  \item 将 tile 放入 shared memory；
  \item 线程从 shared memory 读取数据，在寄存器中累加；
  \item 处理下一个 $K$ 方向 tile；
  \item 最后将结果写回 HBM。
\end{enumerate}

Shared Memory 的优势来自可控性。L1 Cache 由硬件缓存策略控制，而 Shared Memory 由程序显式管理。高性能 GEMM、FlashAttention、卷积、归约等 kernel 都会利用 shared memory 组织数据复用。

但是 shared memory 也有代价：

\begin{itemize}
  \item 容量有限，使用过多会降低每个 SM 可驻留 block 数；
  \item 需要显式同步，例如 \texttt{\_\_syncthreads()}；
  \item 访问模式不当会造成 bank conflict；
  \item 复杂 tiling 增加代码复杂度和调试难度。
\end{itemize}

对 AI Infra 工程师而言，不一定要手写所有底层 kernel，但必须知道为什么 FlashAttention 这类算法强调“不要物化完整 attention matrix”，为什么 GEMM library 需要复杂的 tile scheduler，为什么长上下文 attention 的瓶颈常常不是数学公式本身，而是 Q、K、V、score、probability、output 在不同存储层级之间移动的方式。

\subsection{HBM：大容量高带宽显存}
\label{subsec:hbm}

HBM 是现代数据中心 GPU 的主要显存。它提供高带宽和较大容量，承载模型参数、梯度、优化器状态、激活、KV Cache、临时 buffer 等。在 AI Infra 中，HBM 同时是“容量资源”和“带宽资源”。

容量维度上，HBM 决定了单卡能放下多大的模型、batch、sequence length 和 KV Cache。训练阶段，一个参数可能对应参数本身、梯度、优化器一阶矩、二阶矩以及混合精度 master weight。推理阶段，模型权重通常占固定显存，而 KV Cache 随 batch size 和上下文长度增长。前面第 12 章讨论过 KV Cache 的显存估算，本章从硬件角度补充一点：即使 HBM 容量足够，如果 decode 阶段每生成一个 token 都要反复从 HBM 读取大量历史 K/V，那么 HBM 带宽也会成为限制。

带宽维度上，HBM 决定了 memory-bound kernel 的上限。如果一个 kernel 每做少量计算就要读写大量数据，那么即使 Tensor Core 理论峰值很高，也无法发挥出来。典型例子包括：

\begin{itemize}
  \item embedding lookup：读取为主，计算很少；
  \item LayerNorm / RMSNorm：多次读写和归约，算术强度有限；
  \item decode attention：每个新 token 需要扫描历史 K/V；
  \item 小 batch GEMM：计算规模不足，无法摊薄数据搬运开销；
  \item 大量 elementwise kernel：每个元素只做一两次运算，却要读写 HBM。
\end{itemize}

因此，GPU 性能优化的一个重要原则是：不要只看 FLOPs，要看数据从哪里来、到哪里去、被复用了几次。很多“算法优化”本质上是“IO 优化”。

\section{Tensor Core 与矩阵计算}
\label{sec:tensor-core-matrix}

深度学习工作负载的核心是矩阵计算。无论是 Transformer 中的 QKV projection、attention output projection、FFN 的 up projection 和 down projection，还是 MoE 中每个 expert 的 MLP，本质上都包含大量 GEMM。Tensor Core 的出现，使得 AI GPU 的峰值算力越来越依赖矩阵乘加硬件，而不是传统标量浮点单元。

\subsection{MMA 指令}
\label{subsec:mma}

MMA 是 Matrix Multiply-Accumulate 的缩写，即矩阵乘加。抽象形式如下：

\[
  \mat{D} = \mat{A}\mat{B} + \mat{C}.
\]

其中 $\mat{A}$ 和 $\mat{B}$ 是输入 tile，$\mat{C}$ 是累加器中的旧值，$\mat{D}$ 是新的累加结果。在 GPU 的底层指令语义中，MMA 通常不是由单个线程独立完成，而是由一个 warp 或一组线程协同完成。每个线程持有 tile 的一部分 fragment，硬件将这些 fragment 送入 Tensor Core，完成小矩阵乘加，然后把结果分布式地返回给线程寄存器。

这与普通 for-loop 形式的矩阵乘法有明显差异。普通写法可能是：

\begin{lstlisting}[language=C++,caption={概念上的标量矩阵乘法},label={lst:scalar-gemm-concept}]
for (int m = 0; m < M; ++m) {
  for (int n = 0; n < N; ++n) {
    float acc = 0.0f;
    for (int k = 0; k < K; ++k) {
      acc += A[m][k] * B[k][n];
    }
    C[m][n] = acc;
  }
}
\end{lstlisting}

而高性能 GPU GEMM 会把它重写为多层 tiling：

\begin{itemize}
  \item grid 层：不同 thread block 负责 $\mat{C}$ 的不同大 tile；
  \item block 层：一个线程块协同加载 $\mat{A}$、$\mat{B}$ 的 tile；
  \item warp 层：warp 负责更小的矩阵 tile；
  \item instruction 层：MMA 指令在 Tensor Core 上执行小 tile 乘加；
  \item register 层：每个线程保存 accumulator fragment。
\end{itemize}

\begin{figure}[H]
  \centering
  \begin{tikzpicture}[
    font=\small,
    cell/.style={draw, minimum width=0.48cm, minimum height=0.48cm},
    arrow/.style={-{Stealth[length=2mm]}, thick},
    labelnode/.style={align=center}
  ]
    \node[labelnode] at (-2.2,2.1) {$\mat{A}$ tile};
    \foreach \i in {0,1,2,3} {
      \foreach \j in {0,1,2,3} {
        \node[cell] at (-3+\j*0.48,1.3-\i*0.48) {};
      }
    }

    \node[labelnode] at (1.2,2.1) {$\mat{B}$ tile};
    \foreach \i in {0,1,2,3} {
      \foreach \j in {0,1,2,3} {
        \node[cell] at (0.4+\j*0.48,1.3-\i*0.48) {};
      }
    }

    \node[draw, rounded corners, minimum width=2.6cm, minimum height=1.2cm, align=center] (tc) at (-0.1,-1.0) {Tensor Core\\MMA};

    \node[labelnode] at (3.7,2.1) {Accumulator\\fragment};
    \foreach \i in {0,1,2,3} {
      \foreach \j in {0,1,2,3} {
        \node[cell] at (3+\j*0.48,1.3-\i*0.48) {};
      }
    }

    \draw[arrow] (-2.05,-0.65) -- (tc.west);
    \draw[arrow] (1.35,-0.65) -- (tc.east);
    \draw[arrow] (tc.north east) -- (3.25,0.05);

    \node[align=center] at (-0.1,-2.2) {小 tile 乘加：$\mat{D}_{tile} = \mat{A}_{tile}\mat{B}_{tile} + \mat{C}_{tile}$};
  \end{tikzpicture}
  \caption{Tensor Core tile 计算示意：大 GEMM 被拆成许多小 tile，每个 tile 通过 MMA 指令完成矩阵乘加。}
  \label{fig:tensor-core-tile}
\end{figure}

\subsection{FP16、BF16、TF32 与 FP8}
\label{subsec:precision-formats}

AI GPU 支持多种数值格式。不同格式的目标不是单纯追求“精度越高越好”，而是在精度、吞吐、显存和带宽之间取得平衡。

\begin{table}[H]
  \centering
  \caption{常见深度学习数值格式的工程取舍}
  \label{tab:numeric-format-tradeoff}
  \begin{tabular}{p{2.0cm}p{3.2cm}p{4.2cm}p{4.2cm}}
    \toprule
    格式 & 主要用途 & 优势 & 风险或限制 \\
    \midrule
    FP32 & 高精度训练、累加、调试基线 & 动态范围和精度较好，数值稳定 & 显存和带宽开销大，Tensor Core 快路径未必最高效 \\
    TF32 & FP32 输入场景下的加速矩阵乘 & 兼顾 FP32 易用性与 Tensor Core 加速 & 精度低于完整 FP32，需关注数值敏感任务 \\
    FP16 & 混合精度训练和推理 & 显存占用低，吞吐高 & 动态范围较小，可能需要 loss scaling \\
    BF16 & 大模型训练常用格式 & 动态范围接近 FP32，训练稳定性好 & 有效尾数少于 FP16，某些细粒度数值误差更大 \\
    FP8 & 新一代低精度训练和推理 & 显存和带宽进一步降低，吞吐潜力高 & 需要良好的 scaling、校准和框架支持 \\
    INT8/INT4 & 推理量化 & 权重体积和带宽压力显著下降 & 精度损失、kernel 支持和量化策略复杂 \\
    \bottomrule
  \end{tabular}
\end{table}

训练中常见的 mixed precision 策略是：前向和反向的大矩阵乘法使用 FP16 或 BF16，部分归约和累加使用 FP32 或更高稳定性的路径，优化器状态根据框架和策略可能保留 FP32 或采用分片、低精度、offload 等方式。推理中则更加激进，常见策略包括 FP16/BF16 权重、INT8/INT4 weight-only quantization，以及部分场景下的 FP8 推理。

从硬件映射角度看，数值格式影响三个方面：

\begin{enumerate}
  \item \textbf{Tensor Core 指令路径}：不同数据类型对应不同的 MMA 指令和吞吐上限。
  \item \textbf{HBM 带宽需求}：同样元素数量下，低精度格式需要搬运的字节更少。
  \item \textbf{寄存器和 shared memory 压力}：tile 中可容纳的元素数量、对齐方式和 fragment 布局会改变。
\end{enumerate}

需要强调的是，低精度不是免费的。以量化推理为例，INT4 权重可以显著减少 HBM 读取量，但计算前可能需要 dequantize；如果 dequantize 开销、数据重排开销或 kernel 不够优化，实际收益可能低于理论显存节省比例。这也是第 15 章中反复强调“量化技术必须与 kernel 实现结合”的原因。

\subsection{Tile-based GEMM}
\label{subsec:tile-based-gemm}

Tile-based GEMM 是高性能矩阵乘法的核心思想。假设要计算：

\[
  \mat{C}_{M \times N} = \mat{A}_{M \times K}\mat{B}_{K \times N}.
\]

朴素矩阵乘法中，每个 $\mat{C}[m,n]$ 都要遍历 $K$ 维，重复读取 $\mat{A}$ 和 $\mat{B}$。如果直接从 HBM 读取，每个元素的复用很差。Tile-based GEMM 则把 $\mat{C}$ 分成多个二维 tile，每个 thread block 负责一个或多个 tile。对于每个 $\mat{C}$ tile，它沿 $K$ 方向分多轮读取 $\mat{A}$ tile 和 $\mat{B}$ tile，并在寄存器中累加结果。

其性能优势来自三个层次的数据复用：

\begin{itemize}
  \item \textbf{HBM 到 shared memory 的复用}：一个 $\mat{A}$ tile 被同一线程块中的多个输出列复用，一个 $\mat{B}$ tile 被多个输出行复用；
  \item \textbf{shared memory 到 register 的复用}：线程从 shared memory 读取 fragment，并在多个 MMA 中复用；
  \item \textbf{register accumulator 的复用}：中间结果长期保留在寄存器中，直到最终写回 HBM。
\end{itemize}

\begin{lstlisting}[language=C++,caption={高度简化的 tiled GEMM 伪代码},label={lst:tiled-gemm-pseudo}]
for each output tile C_tile assigned to a thread block:
  initialize accumulator fragments in registers

  for k_tile in range(0, K, TILE_K):
    load A_tile from HBM to shared memory
    load B_tile from HBM to shared memory
    synchronize threads

    for each mma step inside the tile:
      load A_fragment from shared memory to registers
      load B_fragment from shared memory to registers
      accumulator = mma(A_fragment, B_fragment, accumulator)

    synchronize threads

  store accumulator fragments from registers to HBM
\end{lstlisting}

实际高性能 GEMM 比这个伪代码复杂得多。它可能包含多级 pipeline、double buffering、asynchronous copy、warp specialization、split-K、persistent thread block、stream-K、epilogue fusion 等优化。但核心目标始终一致：让 Tensor Core 持续有数据可算，并尽量减少对 HBM 的重复访问。

\subsection{深度学习算子的硬件映射}
\label{subsec:operator-hardware-mapping}

Transformer 中的主要算子可以粗略映射到不同硬件瓶颈上：

\begin{table}[H]
  \centering
  \caption{Transformer 常见算子的硬件映射}
  \label{tab:transformer-operator-hardware}
  \begin{tabular}{p{3.0cm}p{3.8cm}p{4.2cm}p{3.2cm}}
    \toprule
    算子 & 主要硬件路径 & 常见瓶颈 & 优化方向 \\
    \midrule
    QKV Projection & Tensor Core GEMM & 小 batch 或维度不对齐导致利用率不足 & 合并 QKV、选择合适 layout、使用高性能 GEMM \\
    Attention Score & Tensor Core 或 CUDA Core，取决于形状和实现 & 长序列下中间矩阵和 HBM IO & FlashAttention、tiling、避免物化 score \\
    Softmax & CUDA Core / SFU / Load Store & 归约、指数函数、访存 & online softmax、融合 kernel \\
    Attention Output & Tensor Core GEMM & KV 读取、概率矩阵读取 & 与 attention 融合、减少 HBM 写回 \\
    FFN Up / Down & Tensor Core GEMM & 通常是计算主导，但受 shape 影响 & 高效 GEMM、activation fusion \\
    LayerNorm / RMSNorm & CUDA Core + Load Store & memory-bound、归约开销 & fusion、向量化读写、减少中间结果 \\
    RoPE & CUDA Core + Load Store & elementwise 访存 & 与 Q/K projection 或 attention 融合 \\
    Sampling & CUDA Core + Load Store & top-k/top-p 排序或筛选 & 专用 kernel、减少 CPU/GPU 同步 \\
    \bottomrule
  \end{tabular}
\end{table}

这里最重要的工程直觉是：不同算子对 GPU 的压力完全不同。大 GEMM 希望 Tensor Core 满载；norm 和 elementwise 希望减少 HBM 读写；softmax 希望高效归约并避免数值不稳定；attention 希望避免 $O(S^2)$ 中间矩阵写回 HBM；sampling 希望减少小 kernel 和 CPU/GPU 同步。一个端到端 LLM 系统的性能，是这些算子共同作用的结果。

\section{GPU Memory Hierarchy}
\label{sec:gpu-memory-hierarchy}

GPU 的内存层级决定了大量 AI workload 的真实性能。一个 kernel 的理论 FLOPs 很容易计算，但它是否真的快，往往取决于数据是否能停留在靠近计算单元的位置。如果每一步计算都需要从 HBM 重新读取数据，性能就会受 HBM 带宽限制；如果数据能在 register 或 shared memory 中被多次复用，则更容易接近 Tensor Core 或 CUDA Core 的计算上限。

\subsection{Register}
\label{subsec:memory-register}

Register 是每个线程私有的片上存储。它的典型用途包括：

\begin{itemize}
  \item 保存循环索引和地址计算结果；
  \item 保存输入 fragment；
  \item 保存 MMA accumulator；
  \item 保存归约中的局部结果；
  \item 保存临时变量和编译器生成的中间值。
\end{itemize}

在矩阵乘法中，register 的关键角色是保存 accumulator。假设一个线程负责输出 tile 中的多个元素，那么这些元素的部分和可以一直保存在寄存器中，直到 $K$ 维遍历完成后再写回 HBM。这样每个输出元素只写一次 HBM，而不是每次乘加都读写 HBM。

不过 register 是稀缺资源。编译器通常会在编译报告或 profiler 中显示每个线程使用的寄存器数量。若寄存器使用过多，会影响 active warps per SM；若发生 spilling，会产生 local memory 访问。优化时常见做法包括减少临时变量生命周期、调整 unroll 因子、改变 tile shape、拆分 kernel 或使用编译参数限制寄存器数。但限制寄存器并不总是好事，因为它可能降低单线程计算复用或导致更多指令。

\subsection{Shared Memory}
\label{subsec:memory-shared}

Shared Memory 是线程块内共享的片上 SRAM。它既可以作为显式管理的缓存，也可以作为线程之间交换数据的区域。典型使用场景包括：

\begin{itemize}
  \item GEMM 中缓存 $\mat{A}$ 和 $\mat{B}$ 的 tile；
  \item FlashAttention 中缓存 $\mat{Q}$、$\mat{K}$、$\mat{V}$ tile；
  \item reduction 中保存 partial sum；
  \item stencil、卷积或局部窗口计算中复用邻域数据；
  \item transpose 中重排数据以实现合并访存。
\end{itemize}

Shared Memory 的一个常见坑是 bank conflict。Shared Memory 通常被组织成多个 bank。如果同一个 warp 中多个线程访问落在同一 bank 的不同地址，访问可能被串行化，从而降低带宽。对于矩阵 tile，常见优化包括 padding、改变 layout、使用向量化 load/store、调整线程映射方式等。

另一个问题是同步。线程块内多个线程协同加载 shared memory 后，通常需要 \texttt{\_\_syncthreads()} 保证所有数据已经写入，然后才能继续计算。同步本身有成本，而且会限制线程执行的灵活性。高性能 kernel 会尽量通过 double buffering、asynchronous copy 和更细粒度的 warp-level 协同减少同步开销。

\subsection{L2 Cache}
\label{subsec:memory-l2}

L2 Cache 是多个 SM 共享的缓存层级。与 shared memory 不同，L2 由硬件管理，程序通常不能像管理 shared memory 一样精确控制它。不过，L2 对很多 AI workload 仍然非常重要：

\begin{itemize}
  \item 多个 SM 可能重复读取相同模型权重；
  \item batch 内多个请求可能共享部分 prefix 或权重访问路径；
  \item 某些算子之间存在 producer-consumer 局部性；
  \item atomic、跨 block 数据交换或全局归约可能受 L2 行为影响；
  \item GPU 间通信和 DMA 路径在某些架构下也会与 L2 一致性策略相关。
\end{itemize}

在 LLM 推理中，模型权重很大，通常无法完全留在 L2 中。但对小模型、局部层、重复访问的 scale/zero point、metadata、block table 等数据，L2 命中率仍可能显著影响性能。对于 PagedAttention 这类系统，block table 和 KV block 的访问局部性也会影响 cache 行为。

\subsection{HBM}
\label{subsec:memory-hbm}

HBM 是 GPU 的主显存，容量大、带宽高，但与片上存储相比延迟仍然高。绝大多数深度学习张量最终都位于 HBM，包括：

\begin{itemize}
  \item 模型权重；
  \item optimizer state；
  \item activation；
  \item gradient；
  \item KV Cache；
  \item logits、采样中间结果；
  \item 各类临时 workspace。
\end{itemize}

HBM 优化的核心是减少不必要读写，提高合并访存，并增加每个字节被读取后的计算量。下面给出一个非常简单但重要的算术强度例子。

对两个 FP16 向量做逐元素加法：

\[
  c_i = a_i + b_i.
\]

每个元素需要读取 $a_i$ 和 $b_i$，写出 $c_i$。如果忽略 cache，每个元素大约搬运 $2 + 2 + 2 = 6$ 字节，只做一次加法。算术强度大约是：

\[
  I \approx \frac{1\ \text{FLOP}}{6\ \text{Bytes}}.
\]

这个 kernel 几乎一定是 memory-bound。相比之下，矩阵乘法中 $\mat{A}$ 和 $\mat{B}$ 的元素可以被多次复用，算术强度高得多，更可能接近 compute-bound。

\begin{figure}[H]
  \centering
  \begin{tikzpicture}[
    font=\small,
    box/.style={draw, rounded corners, minimum width=8.5cm, minimum height=0.8cm, align=center},
    arrow/.style={-{Stealth[length=2mm]}, thick}
  ]
    \node[box] (reg) at (0,4.0) {Register：线程私有，最快，容量最小};
    \node[box] (shm) at (0,2.8) {Shared Memory / L1：线程块内共享，显式复用};
    \node[box] (l2) at (0,1.6) {L2 Cache：多个 SM 共享，硬件管理};
    \node[box] (hbm) at (0,0.4) {HBM / Global Memory：容量最大，带宽高但延迟高};
    \node[box] (host) at (0,-0.8) {Host Memory / Storage：CPU 内存、NVMe、网络存储};

    \draw[arrow] (reg.south) -- (shm.north);
    \draw[arrow] (shm.south) -- (l2.north);
    \draw[arrow] (l2.south) -- (hbm.north);
    \draw[arrow] (hbm.south) -- (host.north);

    \node[align=left] at (6.2,4.0) {快\\小\\近计算};
    \node[align=left] at (6.2,-0.8) {慢\\大\\远计算};
    \draw[arrow] (5.7,-0.45) -- (5.7,3.65);
  \end{tikzpicture}
  \caption{GPU 内存层级：越靠近计算单元，访问越快但容量越小；越远离计算单元，容量越大但访问代价越高。}
  \label{fig:gpu-memory-hierarchy-inline}
\end{figure}

\subsection{Roofline Model}
\label{subsec:roofline}

Roofline model 是理解 GPU 性能瓶颈的经典工具。它把 kernel 性能分解为两个上限：

\begin{itemize}
  \item \textbf{计算上限}：硬件峰值计算吞吐，例如某种精度下的 Tensor Core 峰值；
  \item \textbf{带宽上限}：内存带宽乘以算术强度。
\end{itemize}

算术强度定义为：

\[
  I = \frac{\text{FLOPs}}{\text{Bytes moved}}.
\]

在 roofline 图中，横轴是算术强度，纵轴是实际可达到的计算吞吐。低算术强度区域受内存带宽限制，性能大致等于：

\[
  P \leq B \cdot I,
\]

其中 $B$ 是内存带宽。高算术强度区域受计算峰值限制，性能大致满足：

\[
  P \leq P_{\text{peak}}.
\]

最终上限可以写成：

\[
  P \leq \min(P_{\text{peak}}, B \cdot I).
\]

这条式子虽然简单，却能解释大量 AI Infra 问题：

\begin{itemize}
  \item 大 GEMM 算术强度高，通常更接近 compute-bound；
  \item elementwise 算子算术强度低，通常 memory-bound；
  \item FlashAttention 的核心收益不是降低 attention 的理论 FLOPs 数量，而是减少 HBM IO；
  \item KV Cache decode 要反复读取历史 K/V，容易受 HBM 带宽限制；
  \item 量化推理减少每个权重元素的字节数，本质上提高了有效算术强度。
\end{itemize}

\begin{figure}[H]
  \centering
  \begin{tikzpicture}[
    font=\small,
    arrow/.style={-{Stealth[length=2mm]}, thick}
  ]
    \draw[arrow] (0,0) -- (9,0) node[right] {算术强度 $I$};
    \draw[arrow] (0,0) -- (0,5.5) node[above] {性能 $P$};

    \draw[thick] (0.6,0.5) -- (4.6,4.2) -- (8.4,4.2);
    \node[align=center] at (2.2,2.6) {memory-bound\\$P \approx B \cdot I$};
    \node[align=center] at (6.6,4.55) {compute-bound\\$P \approx P_{\text{peak}}$};

    \draw[dashed] (4.6,0) -- (4.6,4.2);
    \node[below] at (4.6,0) {ridge point};

    \filldraw (1.2,1.05) circle (2pt);
    \node[align=left, right] at (1.3,1.05) {LayerNorm\\Elementwise};

    \filldraw (2.6,2.35) circle (2pt);
    \node[align=left, right] at (2.7,2.35) {Decode Attention};

    \filldraw (6.4,4.05) circle (2pt);
    \node[align=left, below] at (6.4,3.8) {Large GEMM};

    \filldraw (5.2,3.9) circle (2pt);
    \node[align=left, below] at (5.2,3.55) {FFN GEMM};
  \end{tikzpicture}
  \caption{Roofline model 示意：低算术强度 kernel 受内存带宽限制，高算术强度 kernel 受计算峰值限制。}
  \label{fig:roofline-model}
\end{figure}

Roofline model 不是精确预测器。真实性能还会受到指令调度、cache 命中率、访存合并、bank conflict、同步、warp divergence、kernel launch overhead、通信和框架调度影响。但它提供了一个非常实用的第一判断：一个 kernel 到底应该优先优化计算，还是优先优化 IO？

\section{Profiling 与性能分析}
\label{sec:profiling-performance}

性能优化不能只靠猜。GPU 程序的执行路径非常复杂，框架层、编译器层、kernel 层、驱动层和硬件层都会影响最终表现。Profiler 的作用是把“感觉 GPU 没跑满”转化为可观察、可定位、可验证的问题。

在 AI Infra 工作中，常见的性能分析分为两类：

\begin{itemize}
  \item \textbf{系统级 timeline 分析}：关注一次训练 step 或一次推理请求中，CPU、GPU、通信、kernel、memory copy 的时间关系。
  \item \textbf{kernel 级指标分析}：关注单个 kernel 内部的 occupancy、warp stall、memory throughput、SM utilization、Tensor Core utilization、register usage、shared memory usage 等。
\end{itemize}

前者适合判断“时间花在哪里”，后者适合判断“这个 kernel 为什么慢”。

\subsection{Nsight Systems}
\label{subsec:nsight-systems}

Nsight Systems 适合做系统级 profiling。它展示 CPU thread、CUDA API、GPU kernel、memory copy、NCCL communication、NVTX range 等事件在时间轴上的分布。对 AI Infra 工程师来说，它特别适合回答下面这些问题：

\begin{itemize}
  \item GPU 是否存在明显空洞？
  \item CPU 是否在频繁阻塞 GPU？
  \item kernel launch 是否过多、过碎？
  \item H2D / D2H memory copy 是否出现在关键路径上？
  \item NCCL 通信是否与计算重叠？
  \item prefill 和 decode 是否被合理调度？
  \item 不同 stream 之间是否存在不必要同步？
\end{itemize}

例如，在分布式训练中，如果 Nsight Systems timeline 显示 backward 结束后才开始 all-reduce，说明通信没有与反向计算充分重叠；如果大量小 kernel 之间有明显空隙，可能是 CPU launch overhead 或框架调度开销；如果 GPU 长时间空闲而 CPU 忙于 dataloader，瓶颈可能在数据输入 pipeline，而不是模型 kernel。

在 LLM serving 中，Nsight Systems 可以帮助观察 prefill、decode、KV Cache copy、sampling、network IO 之间的关系。如果每个 token decode 都伴随大量 CPU/GPU 同步，那么即使单个 kernel 很快，端到端 TPOT 也可能很差。

\begin{figure}[H]
  \centering
  \begin{tikzpicture}[
    font=\small,
    lane/.style={draw, minimum width=11.5cm, minimum height=0.75cm, anchor=west},
    task/.style={draw, rounded corners, minimum height=0.45cm, anchor=west, align=center},
    arrow/.style={-{Stealth[length=2mm]}, thick}
  ]
    \node[lane] (cpu) at (0,3.0) {};
    \node[anchor=west] at (0.1,3.0) {CPU};

    \node[lane] (gpu0) at (0,2.0) {};
    \node[anchor=west] at (0.1,2.0) {GPU stream 0};

    \node[lane] (gpu1) at (0,1.0) {};
    \node[anchor=west] at (0.1,1.0) {GPU stream 1};

    \node[lane] (comm) at (0,0.0) {};
    \node[anchor=west] at (0.1,0.0) {NCCL / Copy};

    \node[task, minimum width=1.3cm] at (1.4,3.0) {launch};
    \node[task, minimum width=1.3cm] at (3.1,3.0) {launch};
    \node[task, minimum width=1.3cm] at (5.0,3.0) {sync};
    \node[task, minimum width=1.3cm] at (7.2,3.0) {launch};

    \node[task, minimum width=2.1cm] at (1.6,2.0) {GEMM};
    \node[task, minimum width=1.5cm] at (3.9,2.0) {Norm};
    \node[task, minimum width=2.0cm] at (7.3,2.0) {Attention};

    \node[task, minimum width=1.8cm] at (5.4,1.0) {Sampling};
    \node[task, minimum width=2.8cm] at (4.6,0.0) {AllReduce / Copy};

    \draw[arrow] (5.0,2.65) -- (5.0,2.25);
    \draw[arrow] (6.0,0.35) -- (6.0,0.75);
  \end{tikzpicture}
  \caption{系统级 timeline 的抽象示例：通过观察 CPU、GPU、通信和同步关系定位端到端瓶颈。}
  \label{fig:profiling-timeline-abstract}
\end{figure}

\subsection{Nsight Compute}
\label{subsec:nsight-compute}

Nsight Compute 适合做单 kernel 深度分析。它不像 Nsight Systems 那样关注全局时间线，而是深入一个 kernel 内部，分析其资源使用和硬件效率。常见指标包括：

\begin{itemize}
  \item \textbf{Achieved Occupancy}：实际活跃 warp 与理论最大 warp 的比例；
  \item \textbf{SM Throughput}：SM 执行单元忙碌程度；
  \item \textbf{Memory Throughput}：内存子系统带宽利用率；
  \item \textbf{L1/L2 Hit Rate}：缓存命中情况；
  \item \textbf{Warp Stall Reasons}：warp 等待的主要原因，例如等待内存、等待依赖、等待同步；
  \item \textbf{Register Usage}：每线程寄存器数量；
  \item \textbf{Shared Memory Usage}：每 block 使用的 shared memory；
  \item \textbf{Tensor Core Utilization}：矩阵乘加单元利用程度。
\end{itemize}

单独看某一个指标很容易误判。比如 occupancy 低不一定是问题，memory throughput 低也不一定说明内存不是瓶颈。如果 kernel 因为非合并访存、依赖链或分支导致无法发出足够请求，那么实际 memory throughput 可能不高，但 warp stall 却显示大量等待内存。这时瓶颈仍然可能是访存模式，而不是带宽已经打满。

一个实用的分析顺序是：

\begin{enumerate}
  \item 先从 Nsight Systems 找到耗时最长或最可疑的 kernel；
  \item 用 Nsight Compute 分析该 kernel 的资源使用；
  \item 判断它更像 compute-bound、memory-bound、latency-bound 还是 launch/sync-bound；
  \item 对应地修改 shape、layout、fusion、tiling、并行度或调度策略；
  \item 再次 profiling 验证。
\end{enumerate}

\subsection{Kernel Timeline}
\label{subsec:kernel-timeline}

Kernel timeline 是最直观的性能入口。它告诉我们一次训练 step 或推理请求中，到底执行了哪些 GPU kernel，它们的顺序是什么，耗时是多少，是否存在空隙，是否与通信重叠。

在 PyTorch 或类似框架中，一个高层算子可能对应多个底层 kernel。例如 LayerNorm 可能包含 reduction、normalize、scale、bias；attention 可能包含 QK GEMM、mask、softmax、dropout、PV GEMM；采样可能包含 logits processor、top-k、top-p、random sampling 等。如果没有 fusion，这些小 kernel 会频繁读写 HBM，并产生 kernel launch overhead。

Kernel timeline 中常见的坏味道包括：

\begin{itemize}
  \item \textbf{大量短小 kernel}：每个 kernel 只运行几微秒，launch overhead 和 HBM 往返显著；
  \item \textbf{GPU 空洞}：kernel 之间存在明显间隔，可能是 CPU 调度、同步或数据准备问题；
  \item \textbf{通信孤岛}：NCCL kernel 与计算完全不重叠；
  \item \textbf{意外 memory copy}：H2D、D2H 或 D2D copy 出现在关键路径；
  \item \textbf{同步过密}：频繁 cudaDeviceSynchronize 或隐式同步破坏流水。
\end{itemize}

对于 LLM serving，kernel timeline 还可以揭示 prefill 与 decode 的差异。Prefill 阶段通常包含较大的 GEMM 和 attention，GPU 更容易被喂饱；decode 阶段每轮只处理一个新 token，更容易出现小 GEMM、KV Cache 读取和调度开销。Continuous batching、PagedAttention、prefix caching、speculative decoding 等优化，本质上都试图改善这些 timeline 特征。

\subsection{Occupancy、Memory Throughput 与 SM Efficiency}
\label{subsec:key-metrics}

GPU 性能分析中经常出现三个指标：occupancy、memory throughput 和 SM efficiency。它们都重要，但都不能孤立解释性能。

\subsubsection{Occupancy}

Occupancy 反映 SM 上活跃 warp 的数量。影响 occupancy 的因素包括：

\begin{itemize}
  \item 每个 block 的线程数；
  \item 每个线程使用的寄存器数；
  \item 每个 block 使用的 shared memory；
  \item 架构允许的最大 block、warp、register 和 shared memory 限制。
\end{itemize}

高 occupancy 可以帮助隐藏内存延迟，但高 occupancy 不保证高吞吐。对于 compute-bound GEMM，较低 occupancy 也可能因为 Tensor Core 利用率高而表现很好。对于 latency-bound 或 memory-bound kernel，提升 occupancy 可能有效，但前提是有足够独立 warp 可以发出访存请求。

\subsubsection{Memory Throughput}

Memory throughput 表示 kernel 实际使用的内存带宽。它需要结合理论峰值和算子特征解读：

\begin{itemize}
  \item 如果 memory-bound kernel 已接近带宽上限，继续优化计算指令意义不大；
  \item 如果 memory-bound kernel 带宽很低，可能是访存不合并、请求粒度差、cache miss、warp divergence 或依赖导致；
  \item 如果 compute-bound kernel memory throughput 不高，可能是正常现象，因为瓶颈在计算；
  \item 如果大量中间结果写回 HBM，fusion 或 tiling 可能比减少 FLOPs 更有效。
\end{itemize}

\subsubsection{SM Efficiency}

SM efficiency 或 SM utilization 表示 SM 忙碌程度。它低可能有多种原因：

\begin{itemize}
  \item kernel 太小，无法覆盖所有 SM；
  \item block 数量不足；
  \item 每个 block 资源占用过高，导致活跃 warp 少；
  \item warp 经常等待内存、同步或依赖；
  \item CPU launch 间隔导致 GPU 空闲；
  \item 多 GPU 通信或数据加载让 GPU 等待。
\end{itemize}

因此，看到 SM utilization 低时，不能立刻说“需要提高 occupancy”或“需要增大 batch”。正确做法是先确定低利用率发生在单 kernel 内部，还是 kernel 之间的空隙。如果是前者，用 Nsight Compute；如果是后者，用 Nsight Systems。

\subsection{从 Profiler 判断瓶颈}
\label{subsec:diagnose-bottleneck}

下面给出一个面向 AI Infra 的简化诊断表。它不是绝对规则，但可以作为排查起点。

\begin{longtable}{p{4.0cm}p{5.0cm}p{5.0cm}}
  \caption{Profiler 现象与可能瓶颈}
  \label{tab:profiler-diagnosis} \\
  \toprule
  现象 & 可能原因 & 优先排查方向 \\
  \midrule
  \endfirsthead
  \toprule
  现象 & 可能原因 & 优先排查方向 \\
  \midrule
  \endhead
  GPU timeline 中有大空洞 & CPU 调度、dataloader、同步、通信等待 & 检查 CPU thread、CUDA API、数据输入、隐式同步 \\
  大量很短的小 kernel & 算子未融合、框架 eager 开销、elementwise 过多 & 使用 fusion、graph capture、编译器或手写 fused kernel \\
  GEMM 耗时高但 Tensor Core 利用低 & shape 不适合、dtype 不对、layout 不佳、batch 太小 & 检查矩阵维度、对齐、数据类型、是否走 Tensor Core 路径 \\
  Memory throughput 接近峰值 & memory-bound & 减少 HBM 读写、做 fusion、提高数据复用、量化 \\
  Memory throughput 不高但 warp stall memory 高 & 访存延迟、非合并访存、请求不足 & 改善 coalescing、增加并行度、调整 layout、使用 shared memory \\
  Occupancy 很低 & register/shared memory 过多，block 配置不合理 & 检查寄存器、shared memory、block size，但不要盲目追求高 occupancy \\
  L2 hit rate 异常低 & 访问局部性差、working set 太大、随机访问 & 调整数据布局、分块、重排访问、减少 metadata 随机访问 \\
  NCCL 与 backward 不重叠 & bucket 设置、调度顺序、依赖关系问题 & 调整 bucket、检查 DDP/ZeRO 配置、观察 stream 关系 \\
  Decode TPOT 高 & KV Cache 读取、batch 太小、小 GEMM、采样开销 & Continuous batching、PagedAttention、speculative decoding、kernel fusion \\
  Prefill TTFT 高 & 长 prompt attention、GEMM、调度排队 & FlashAttention、chunked prefill、prefix caching、admission control \\
  \bottomrule
\end{longtable}

为了让诊断更具体，我们考虑两个常见案例。

\subsubsection{案例一：训练中大 GEMM 没有打满}

假设某个 Transformer FFN 的 up projection 很慢，profiler 显示 GEMM kernel 占据大量时间，但 Tensor Core utilization 不高。此时可能原因包括：

\begin{itemize}
  \item 输入 dtype 是 FP32，未启用 TF32 或混合精度；
  \item $M$ 维过小，例如 micro batch size 太小且 sequence length 不大；
  \item hidden size 或 intermediate size 不满足高效对齐；
  \item tensor layout 导致额外 transpose 或 copy；
  \item 使用了非最优 backend；
  \item kernel 被其他 stream 或资源竞争影响。
\end{itemize}

排查路径可以是：确认 dtype 和 matmul precision；打印实际矩阵 shape；用单独 benchmark 测试相同 shape 的 GEMM；比较 cuBLAS、Triton 或框架编译后的 kernel；观察是否存在额外 copy；调整 micro batch 或 gradient accumulation 策略。

\subsubsection{案例二：推理 decode 随 batch 增加后仍不快}

LLM decode 阶段常见现象是：增大 batch 后吞吐上升，但 TPOT 也变差；某些 batch 规模后吞吐不再提升。这可能是因为：

\begin{itemize}
  \item KV Cache 读取带宽成为瓶颈；
  \item 请求长度差异导致 cache 访问不规则；
  \item batch 中短请求被长请求拖累；
  \item attention kernel 对 paged KV 的访问局部性不足；
  \item sampling 或 logits processing 成为 CPU/GPU 同步瓶颈；
  \item memory allocator 或 block manager 开销上升。
\end{itemize}

这时不能只看 GEMM。应该同时观察 attention kernel、KV Cache 读写、scheduler、sampling、CPU timeline 和 GPU 空洞。优化方向可能包括 PagedAttention、prefix caching、chunked prefill、continuous batching、speculative decoding、量化 KV Cache 或调度策略调整。

\section{GPU 架构对 AI Infra 设计的影响}
\label{sec:gpu-architecture-ai-infra}

理解 GPU 架构不是为了写出所有底层 kernel，而是为了在系统设计时做出正确取舍。本节从训练、推理和集群调度三个角度总结 GPU 架构的影响。

\subsection{对训练系统的影响}
\label{subsec:gpu-training-impact}

训练系统关心 step time、吞吐、显存、通信和稳定性。GPU 架构影响训练系统的方式包括：

\begin{itemize}
  \item \textbf{Tensor Core 决定大 GEMM 上限}：模型 hidden size、FFN intermediate size、head dimension 等最好与高效矩阵 tile 友好。
  \item \textbf{HBM 容量决定并行策略}：单卡放不下模型或 optimizer state 时，需要 ZeRO、TP、PP 或 offload。
  \item \textbf{HBM 带宽影响非 GEMM 算子}：norm、dropout、activation、optimizer update 都可能成为大规模训练中的隐性开销。
  \item \textbf{Register 和 shared memory 影响自定义 kernel}：FlashAttention、fused optimizer、fused norm 都需要在片上资源间权衡。
  \item \textbf{L2 和通信路径影响 overlap}：梯度同步、reduce-scatter、all-gather 与计算重叠时，数据路径和 stream 调度都很重要。
\end{itemize}

例如，在 3D 并行中，TP 会引入更多 intra-node 通信，PP 会改变 micro-batch 时间线，DP 会引入梯度同步。GPU 本身的计算速度越快，通信和调度空洞就越容易显得突出。这也是为什么大规模训练不是简单地“堆更多 GPU”，而是要让计算、显存和通信平衡。

\subsection{对推理系统的影响}
\label{subsec:gpu-serving-impact}

推理系统关心 TTFT、TPOT、吞吐、SLA 和成本。GPU 架构对推理的影响更加复杂，因为推理 workload 的形状变化非常大。

Prefill 阶段通常类似训练中的前向计算，序列长度较长，GEMM 和 attention 规模较大，更容易利用 Tensor Core。Decode 阶段每次只生成一个 token，虽然 batch 可以通过 continuous batching 增大，但每个请求的历史上下文长度不同，KV Cache 访问成为关键问题。

因此，推理优化常围绕以下方向展开：

\begin{itemize}
  \item \textbf{提高 batch 的有效性}：continuous batching 让 decode 阶段聚合更多请求，增加并行度；
  \item \textbf{降低 KV Cache IO}：PagedAttention、prefix caching、KV Cache quantization 等减少显存浪费或带宽压力；
  \item \textbf{减少小 kernel 和同步}：融合 RoPE、norm、activation、sampling 等操作；
  \item \textbf{选择合适精度}：FP16/BF16、FP8、INT8、INT4 在不同模型和 SLA 下取舍；
  \item \textbf{分离 prefill 与 decode}：针对 compute-heavy prefill 和 memory-heavy decode 使用不同资源池。
\end{itemize}

一个成熟的 LLM serving 系统，本质上是对 GPU 计算单元、HBM 容量、HBM 带宽、kernel launch、cache 管理和调度策略的综合优化。

\subsection{对集群调度的影响}
\label{subsec:gpu-scheduling-impact}

Kubernetes 或其他调度系统看到的资源往往是“几张 GPU”，但真实性能取决于更细粒度的硬件结构：

\begin{itemize}
  \item GPU 型号不同，HBM 容量、带宽、Tensor Core 能力不同；
  \item 同一节点内 GPU 拓扑不同，P2P 通信性能不同；
  \item MIG 或 GPU sharing 会改变 SM、显存、L2、带宽隔离方式；
  \item 多租户推理会产生 cache、带宽和 kernel 调度干扰；
  \item 训练作业的 TP/PP/DP rank mapping 需要匹配 GPU 拓扑。
\end{itemize}

因此，AI 集群调度不能只做数量匹配，还要做拓扑感知、带宽感知、显存感知和 workload 感知。例如，一个 8 GPU TP 作业最好被放在 NVLink/NVSwitch 拓扑良好的单节点或紧密 GPU group 中；一个 memory-bound 推理服务即使 SM utilization 不高，也可能已经耗尽 HBM 带宽；多个 decode-heavy 服务混部时，显存带宽可能比显存容量更早成为瓶颈。

\section{本章小结}
\label{sec:gpu-architecture-summary}

本章从 AI Infra 工程视角介绍了 GPU 架构底座。我们没有追求某一代 GPU 的完整硬件规格，而是建立了几个跨架构稳定的核心概念：

\begin{itemize}
  \item SM 是 CUDA kernel 的基本执行场所，warp 调度、CUDA Core、Tensor Core、Load/Store Unit、Register File 和 Shared Memory 共同决定 kernel 表现。
  \item Tensor Core 是深度学习矩阵计算的关键，但只有在合适的数据类型、矩阵形状、tile 策略和内存布局下才能发挥高吞吐。
  \item Register、Shared Memory、L2 Cache、HBM 构成了 GPU 的内存层级。性能优化经常不是减少数学计算，而是减少数据在 HBM 与片上存储之间的往返。
  \item Roofline model 提供了判断 compute-bound 与 memory-bound 的第一性工具。大 GEMM 通常算术强度高，而 norm、elementwise、decode attention、KV Cache 访问常常受内存带宽限制。
  \item Nsight Systems 适合分析端到端 timeline，Nsight Compute 适合分析单 kernel。性能优化应从 profiler 证据出发，而不是只凭经验猜测。
  \item GPU 架构直接影响训练并行策略、推理服务设计和集群调度。AI Infra 工程师需要把模型 shape、kernel 行为、显存容量、显存带宽、GPU 拓扑和调度系统放在一起思考。
\end{itemize}

下一章将继续沿着硬件路径向外扩展：当一台机器中有多张 GPU 时，它们如何通信？PCIe、NVLink 和 NVSwitch 的差异是什么？为什么单机多卡拓扑会影响 tensor parallel、pipeline parallel 和 NCCL collective 的性能？这些问题将进入第 18 章：NVLink 与 NVSwitch。

\chapter{NVLink 与 NVSwitch}
\label{chap:nvlink-nvswitch}

上一章我们从单个 GPU 的内部结构出发，讨论了 SM、Tensor Core、HBM、Shared Memory、L2 Cache 以及 profiler 如何帮助我们判断 kernel 是 compute-bound 还是 memory-bound。但是在真正的大模型训练和推理系统中，单个 GPU 很快就会遇到边界：模型参数放不下，batch 放不下，KV Cache 放不下，训练时间太长，推理吞吐不够，单卡再强也无法独自承担所有计算。

于是，多 GPU 系统成为 AI Infrastructure 的核心形态。多 GPU 系统中的关键问题不是“把几张 GPU 插到一台机器里”这么简单，而是这些 GPU 之间如何通信、通信带宽和延迟如何影响并行策略、通信拓扑如何影响 NCCL collective、以及错误的 rank mapping 如何让昂贵的 GPU 集群跑出低效性能。

本章讨论单机多卡场景中最重要的 scale-up 互联技术：PCIe、GPU Direct P2P、NVLink 与 NVSwitch。它们共同决定了一台 AI 服务器内部 GPU 之间的数据移动能力。理解这些内容，是进入下一章 RDMA / RoCE 多机通信之前的必要准备。

本章将围绕四个问题展开：

\begin{itemize}
  \item 单机多卡为什么不能只依赖 PCIe？
  \item NVLink 如何改善 GPU-to-GPU 通信？
  \item NVSwitch 如何把点对点互联扩展为接近 all-to-all 的 GPU fabric？
  \item 如何用 \texttt{nvidia-smi topo}、NCCL tests 和 profiler 诊断单机多卡通信问题？
\end{itemize}

\section{单机多卡通信问题}
\label{sec:single-node-multi-gpu}

大模型训练和推理中的多卡通信可以分为两类：一类是显式通信，例如 tensor parallel 中的 all-reduce、all-gather、reduce-scatter；另一类是隐式数据移动，例如不同 GPU 之间访问 peer memory、模型 shard 之间交换激活、流水线并行 stage 之间发送 tensor。无论是哪一类，本质上都需要 GPU 与 GPU 之间高速传输数据。

如果没有高速互联，多卡系统就会出现一种尴尬局面：每张 GPU 的 Tensor Core 都很快，但它们经常在等待数据。训练中，等待梯度同步；推理中，等待 KV Cache 或中间激活跨卡传输；模型并行中，等待前一个 shard 的输出。此时系统瓶颈不在计算，而在 GPU 之间的数据路径。

\subsection{PCIe 的局限}
\label{subsec:pcie-limitations}

PCIe 是通用 I/O 总线，CPU、GPU、NIC、NVMe 等设备都可以通过 PCIe 接入系统。它的优点是生态成熟、通用性强、设备支持广；缺点是它并不是专门为多 GPU 高密度通信设计的。

在传统服务器中，多张 GPU 通常挂在 CPU root complex 或 PCIe switch 下面。GPU 之间通信时，数据可能需要经过 PCIe switch，甚至经过 CPU root complex。这样会带来几个问题：

\begin{itemize}
  \item \textbf{带宽受限}：PCIe 的带宽相对 GPU 内部 HBM 和 Tensor Core 吞吐低得多，容易成为通信瓶颈。
  \item \textbf{拓扑不均匀}：同一台机器中，不同 GPU pair 的通信路径可能不同，有的只经过一个 switch，有的需要跨 CPU socket。
  \item \textbf{延迟较高}：相比 GPU-to-GPU 专用互联，PCIe 路径更长，协议栈更通用，延迟也更高。
  \item \textbf{与其他 I/O 竞争}：NIC、NVMe、其他加速卡可能共享 PCIe 资源，在高负载下互相影响。
  \item \textbf{跨 socket 代价}：双路 CPU 服务器中，如果 GPU 分布在不同 NUMA domain 下，跨 socket 通信可能显著变慢。
\end{itemize}

这并不是说 PCIe 不重要。事实上，PCIe 仍然是 GPU 接入主机系统的基础路径，也常用于 CPU-GPU 数据传输、GPU-NIC 连接、GPU-NVMe 数据路径等。但对于需要频繁 GPU-to-GPU 通信的大模型训练，单靠 PCIe 往往难以满足高吞吐需求。

\begin{figure}[H]
  \centering
  \begin{tikzpicture}[
    font=\small,
    box/.style={draw, rounded corners, minimum width=1.9cm, minimum height=0.75cm, align=center},
    wide/.style={draw, rounded corners, minimum width=4.8cm, minimum height=0.8cm, align=center},
    arrow/.style={-{Stealth[length=2mm]}, thick},
    link/.style={thick}
  ]
    \node[wide] (cpu) at (0,3.5) {CPU / Root Complex};
    \node[wide] (sw) at (0,2.2) {PCIe Switch};
    \node[box] (g0) at (-3,0.8) {GPU 0};
    \node[box] (g1) at (-1,0.8) {GPU 1};
    \node[box] (g2) at (1,0.8) {GPU 2};
    \node[box] (g3) at (3,0.8) {GPU 3};

    \draw[link] (cpu.south) -- (sw.north);
    \draw[link] (sw.south) -- (g0.north);
    \draw[link] (sw.south) -- (g1.north);
    \draw[link] (sw.south) -- (g2.north);
    \draw[link] (sw.south) -- (g3.north);

    \draw[arrow, dashed] (g0.south) .. controls (-2,-0.4) and (2,-0.4) .. (g3.south);
    \node[align=center] at (0,-0.75) {GPU-to-GPU 通信可能需要经过 PCIe Switch / Root Complex};

    \node[align=center] at (0,4.35) {PCIe 拓扑：通用但路径可能不均匀};
  \end{tikzpicture}
  \caption{PCIe 多 GPU 拓扑的抽象示意：不同 GPU 之间通信路径可能经过 PCIe Switch 或 CPU root complex。}
  \label{fig:pcie-gpu-topology}
\end{figure}

\subsection{GPU Direct P2P}
\label{subsec:gpu-direct-p2p}

GPU Direct P2P 允许一张 GPU 直接访问另一张 GPU 的显存地址空间，而不必把数据先拷贝到 CPU host memory 再拷贝到目标 GPU。对上层程序而言，它使 GPU-to-GPU 数据传输更直接，也让通信库可以更高效地实现 peer copy 和 collective。

没有 P2P 时，GPU 之间的数据交换可能退化为如下路径：

\[
  \text{GPU 0 HBM} \rightarrow \text{Host Memory} \rightarrow \text{GPU 1 HBM}.
\]

有 P2P 时，路径可以变成：

\[
  \text{GPU 0 HBM} \rightarrow \text{PCIe / NVLink} \rightarrow \text{GPU 1 HBM}.
\]

这看似只是少了一次 host staging，但工程影响很大。对于大模型训练中的高频通信，绕过 CPU host memory 可以降低延迟、减少 CPU 内存带宽压力、减少 CPU 参与度，并让 NCCL 等通信库更容易构造高效路径。

在 CUDA 程序中，可以通过 peer access 相关 API 检查和启用 GPU 之间的 peer 访问能力。下面是一个概念示例：

\begin{lstlisting}[language=C++,caption={检查并启用 GPU peer access 的概念代码},label={lst:cuda-peer-access}]
int can_access_peer = 0;
cudaDeviceCanAccessPeer(&can_access_peer, 0, 1);

if (can_access_peer) {
  cudaSetDevice(0);
  cudaDeviceEnablePeerAccess(1, 0);

  cudaSetDevice(1);
  cudaDeviceEnablePeerAccess(0, 0);
}
\end{lstlisting}

实际 AI 框架通常不会要求用户手写这些代码。PyTorch、NCCL、DeepSpeed、Megatron-LM、vLLM 等系统会在底层利用 CUDA runtime、driver 和 NCCL 能力。但作为 Infra 工程师，理解 P2P 的意义非常重要：当 P2P 不可用或拓扑不佳时，某些多卡通信性能会突然下降，甚至出现本该在 GPU 间完成的数据传输绕回 CPU 的情况。

\subsection{Intra-node Bandwidth}
\label{subsec:intra-node-bandwidth}

Intra-node bandwidth 指同一台服务器内部 GPU 之间的通信带宽。它对以下并行策略尤其关键：

\begin{itemize}
  \item \textbf{Tensor Parallelism}：每层 Transformer 内部都可能发生 all-reduce、all-gather 或 reduce-scatter，对单机内 GPU 带宽极其敏感。
  \item \textbf{Sequence Parallelism}：activation 沿 sequence 维切分后，需要额外 collective。
  \item \textbf{Expert Parallelism}：MoE token dispatch 和 combine 会产生 all-to-all 通信。
  \item \textbf{Pipeline Parallelism}：stage 之间要传递 activation 和 gradient，通信量通常小于 TP，但延迟仍然影响 bubble。
  \item \textbf{Inference Tensor Parallel}：推理每生成一个 token 都会触发层内跨卡通信，TPOT 对通信延迟敏感。
\end{itemize}

一个常见误区是只看“总带宽”。多 GPU 通信还要关注路径是否均匀、是否全双工、是否需要经过共享 switch、是否跨 NUMA、是否被其他设备抢占、collective 算法是否匹配拓扑。对于 tensor parallel 这种频繁、小步、层内通信，延迟和拓扑均匀性同样重要。对于大 batch 数据并行梯度同步，带宽通常更重要。

\subsection{通信拓扑对并行训练的影响}
\label{subsec:topology-impact-training}

通信拓扑会直接影响并行策略选择。假设一台机器有 8 张 GPU。如果它们通过 NVSwitch 构成高带宽 all-to-all fabric，那么 TP size = 8 可能是合理选择；如果它们只是通过 PCIe 分成两个相对独立的 4 GPU group，那么 TP size = 8 可能会跨越低带宽路径，性能反而不如 TP size = 4 加上更大的 DP 或 PP。

对训练系统来说，rank mapping 是拓扑优化的关键。所谓 rank mapping，就是把全局 rank 分配到物理 GPU 上。错误的 rank mapping 可能让通信密集的 rank 被放在通信路径最差的 GPU pair 上。例如：

\begin{itemize}
  \item TP group 内 rank 应尽量放在同一高速互联域中；
  \item PP 相邻 stage 应尽量减少跨慢链路传输；
  \item DP group 的通信量通常较大但频率低于 TP，可根据节点间网络权衡；
  \item MoE expert parallel 中 all-to-all 的 rank 分布需要结合节点内和节点间带宽。
\end{itemize}

\begin{figure}[H]
  \centering
  \begin{tikzpicture}[
    font=\small,
    gpu/.style={draw, rounded corners, minimum width=1.2cm, minimum height=0.65cm, align=center},
    group/.style={draw, rounded corners, dashed, inner sep=0.25cm},
    fast/.style={line width=1.2pt},
    slow/.style={line width=0.5pt, dashed}
  ]
    \node[gpu] (g0) at (0,2.5) {GPU0};
    \node[gpu] (g1) at (2,2.5) {GPU1};
    \node[gpu] (g2) at (0,1.2) {GPU2};
    \node[gpu] (g3) at (2,1.2) {GPU3};

    \node[gpu] (g4) at (5,2.5) {GPU4};
    \node[gpu] (g5) at (7,2.5) {GPU5};
    \node[gpu] (g6) at (5,1.2) {GPU6};
    \node[gpu] (g7) at (7,1.2) {GPU7};

    \draw[fast] (g0) -- (g1);
    \draw[fast] (g0) -- (g2);
    \draw[fast] (g1) -- (g3);
    \draw[fast] (g2) -- (g3);
    \draw[fast] (g0) -- (g3);
    \draw[fast] (g1) -- (g2);

    \draw[fast] (g4) -- (g5);
    \draw[fast] (g4) -- (g6);
    \draw[fast] (g5) -- (g7);
    \draw[fast] (g6) -- (g7);
    \draw[fast] (g4) -- (g7);
    \draw[fast] (g5) -- (g6);

    \draw[slow] (g1) -- (g4);
    \draw[slow] (g3) -- (g6);

    \node[group, fit=(g0)(g1)(g2)(g3), label=below:{高速互联域 A}] {};
    \node[group, fit=(g4)(g5)(g6)(g7), label=below:{高速互联域 B}] {};

    \node[align=center] at (3.5,3.25) {错误的 TP group 跨越慢链路会放大通信瓶颈};
  \end{tikzpicture}
  \caption{拓扑感知 rank mapping 示例：通信密集的 group 应优先放在高速互联域内。}
  \label{fig:topology-aware-rank-mapping}
\end{figure}

\section{NVLink}
\label{sec:nvlink}

NVLink 是 NVIDIA 面向 GPU-to-GPU 高速通信设计的互联技术。与 PCIe 相比，NVLink 更接近专用 GPU fabric，目标是为多 GPU 训练和推理提供更高带宽、更低延迟和更好的拓扑特性。

需要注意的是，NVLink 不是一个固定带宽的单一产品。不同 GPU 架构、不同 NVLink 代际、不同服务器形态中，link 数量、单 link 带宽、拓扑连接方式都可能不同。本书更关注它在 AI Infra 中的抽象作用：为 GPU 之间提供比普通 PCIe 更适合大规模张量通信的高速路径。

\subsection{点对点高速互联}
\label{subsec:nvlink-p2p}

最基本的 NVLink 可以理解为 GPU 与 GPU 之间的高速点对点链路。若 GPU0 和 GPU1 之间存在 NVLink 连接，那么它们之间的 P2P 传输可以走 NVLink，而不是走 PCIe。对于通信库来说，这意味着同样一次 peer copy、send/recv 或 collective chunk transfer 可以获得更高的传输能力。

\begin{figure}[H]
  \centering
  \begin{tikzpicture}[
    font=\small,
    gpu/.style={draw, rounded corners, minimum width=1.4cm, minimum height=0.75cm, align=center},
    nv/.style={line width=1.4pt},
    pcie/.style={dashed, line width=0.6pt},
    box/.style={draw, rounded corners, minimum width=3.8cm, minimum height=0.8cm, align=center}
  ]
    \node[box] (cpu) at (0,3.4) {CPU / PCIe Root};
    \node[gpu] (g0) at (-3,1.7) {GPU0};
    \node[gpu] (g1) at (-1,1.7) {GPU1};
    \node[gpu] (g2) at (1,1.7) {GPU2};
    \node[gpu] (g3) at (3,1.7) {GPU3};

    \draw[pcie] (cpu.south) -- (g0.north);
    \draw[pcie] (cpu.south) -- (g1.north);
    \draw[pcie] (cpu.south) -- (g2.north);
    \draw[pcie] (cpu.south) -- (g3.north);

    \draw[nv] (g0) -- (g1);
    \draw[nv] (g1) -- (g2);
    \draw[nv] (g2) -- (g3);
    \draw[nv] (g0) to[bend right=20] (g2);
    \draw[nv] (g1) to[bend left=20] (g3);

    \node[align=center] at (0,0.35) {实线：NVLink GPU-to-GPU 高速路径\\虚线：PCIe 到主机路径};
  \end{tikzpicture}
  \caption{NVLink 点对点互联的抽象示意：GPU 之间可通过专用高速链路直接通信。}
  \label{fig:nvlink-point-to-point}
\end{figure}

在训练系统中，NVLink 对 tensor parallel 的影响尤其明显。以 Transformer 中的 row-parallel linear 为例，局部 matmul 后需要对 partial output 做 all-reduce。如果 TP group 内 GPU 之间有高带宽 NVLink，all-reduce 的开销可以显著降低；如果 TP group 跨越 PCIe 或跨节点网络，则通信可能成为每层的主瓶颈。

\subsection{Bandwidth 与 Lane}
\label{subsec:nvlink-bandwidth-lane}

NVLink 的物理实现通常由多条 link 或 lane 组成。一个 GPU 可以通过多条 NVLink 连接到其他 GPU 或交换芯片。不同代际的 NVLink 在单 link 速率、每 GPU link 数量、总带宽、协议能力等方面不同。

在工程分析中，与其死记某个型号的数字，不如关注三个问题：

\begin{enumerate}
  \item \textbf{某两个 GPU 之间是否存在 NVLink 路径？}
  \item \textbf{这条路径是直接连接，还是需要经过 NVSwitch 或其他转发？}
  \item \textbf{这条路径的有效带宽是否被 NCCL 或应用实际利用？}
\end{enumerate}

硬件规格表给出的是理论带宽，而训练作业看到的是有效通信性能。有效通信性能还受数据大小、collective 算法、协议、并发通信数量、GPU 拓扑、NCCL channel 数量、stream 调度、GPU 负载等影响。小消息更容易受延迟影响，大消息更容易受带宽影响。Tensor parallel 中的通信往往频繁且分布在每层中，因此对延迟和带宽都敏感。

\subsection{Peer Access}
\label{subsec:nvlink-peer-access}

NVLink 常与 CUDA peer access 一起发挥作用。启用 peer access 后，GPU 可以通过统一虚拟地址空间访问其他 GPU 的显存。对用户来说，最直观的表现是 \texttt{cudaMemcpyPeerAsync} 或框架内部的 P2P copy 可以走更高效路径。

但是，peer access 是否可用并不只取决于 API 调用，还取决于硬件拓扑、驱动、系统配置、IOMMU、虚拟化、容器权限等因素。某些环境中，即使物理上有多张 GPU，P2P 也可能不可用或受限。

在实际排查中，可以使用以下方式检查：

\begin{lstlisting}[language=bash,caption={检查 GPU 拓扑和 P2P 能力的常用命令},label={lst:check-topology-p2p}]
# 查看 GPU 之间的拓扑关系
nvidia-smi topo -m

# 查看 GPU 设备基本信息
nvidia-smi -L

# 运行 CUDA samples 中的 p2pBandwidthLatencyTest
./p2pBandwidthLatencyTest
\end{lstlisting}

如果 \texttt{nvidia-smi topo -m} 显示两个 GPU 之间是 \texttt{NV}、\texttt{NVL} 或类似 NVLink 路径，通常说明它们之间存在 NVLink 连接；如果显示 \texttt{PIX}、\texttt{PXB}、\texttt{PHB}、\texttt{SYS} 等，则说明路径可能经过不同层级的 PCIe switch、host bridge 或跨 socket。具体标记含义应以当前驱动文档和设备输出为准。

\subsection{NCCL 如何利用 NVLink}
\label{subsec:nccl-nvlink}

NCCL 是深度学习中最常用的 GPU collective 通信库之一。PyTorch DDP、Megatron-LM、DeepSpeed、Horovod、vLLM 的多卡路径以及许多训练框架都会依赖 NCCL。NCCL 会检测机器拓扑，并根据 GPU 之间的连接关系选择通信路径和 collective 算法。

常见 collective 包括：

\begin{itemize}
  \item \textbf{AllReduce}：所有 rank 输入张量求和后，每个 rank 得到完整结果。数据并行梯度同步常用。
  \item \textbf{ReduceScatter}：先 reduce，再把结果切分给不同 rank。ZeRO、tensor parallel 和 sequence parallel 中常见。
  \item \textbf{AllGather}：每个 rank 持有一部分数据，最终所有 rank 收集完整数据。
  \item \textbf{AllToAll}：每个 rank 向每个其他 rank 发送不同数据。MoE token dispatch 常用。
  \item \textbf{Broadcast}：一个 rank 的数据分发到所有 rank。
  \item \textbf{Send / Recv}：点对点通信，pipeline parallel 中常见。
\end{itemize}

在 NVLink 拓扑下，NCCL 可以利用高速 GPU-to-GPU 链路构造 ring、tree 或其他拓扑感知通信计划。对于一个 8 GPU 节点，如果 GPU 之间通过 NVSwitch 形成高带宽互联，NCCL 的 all-reduce 和 all-gather 可以获得比普通 PCIe 拓扑更好的性能。

下面是一个极简的 NCCL all-reduce 调用概念示例：

\begin{lstlisting}[language=C++,caption={NCCL AllReduce 概念代码},label={lst:nccl-allreduce-concept}]
ncclAllReduce(
    sendbuff,
    recvbuff,
    count,
    ncclFloat16,
    ncclSum,
    comm,
    stream
);
\end{lstlisting}

在 PyTorch 中，用户通常通过 \texttt{torch.distributed} 间接使用 NCCL：

\begin{lstlisting}[language=Python,caption={PyTorch 中触发 NCCL AllReduce 的概念示例},label={lst:pytorch-allreduce}]
import torch
import torch.distributed as dist

x = torch.ones(1024, device="cuda")
dist.all_reduce(x, op=dist.ReduceOp.SUM)
\end{lstlisting}

上层代码看起来很简单，但底层性能高度依赖拓扑。相同的 \texttt{dist.all\_reduce}，在 PCIe、NVLink、NVSwitch、跨节点 InfiniBand / RoCE 上的性能可能完全不同。

\section{NVSwitch}
\label{sec:nvswitch}

NVLink 解决了 GPU-to-GPU 高速互联问题，但如果只依赖点对点连接，拓扑会随着 GPU 数量增加而变复杂。每张 GPU 不可能无限制地直接连接所有其他 GPU，link 数量、布线复杂度、拓扑非均匀性都会成为问题。NVSwitch 的目标，就是把 NVLink 从点对点链路扩展为更大规模、更均匀的 GPU fabric。

\subsection{All-to-All GPU Fabric}
\label{subsec:all-to-all-gpu-fabric}

NVSwitch 可以理解为面向 NVLink 的高速交换芯片。GPU 不再只通过点对点 NVLink 直接连接某几个邻居，而是连接到 NVSwitch，再由 NVSwitch 在 GPU 之间转发数据。这样可以在一个节点或一个更大的 scale-up domain 中构造更接近 all-to-all 的通信能力。

\begin{figure}[H]
  \centering
  \begin{tikzpicture}[
    font=\small,
    gpu/.style={draw, rounded corners, minimum width=1.2cm, minimum height=0.65cm, align=center},
    sw/.style={draw, rounded corners, minimum width=2.0cm, minimum height=0.8cm, align=center, fill=gray!10},
    link/.style={line width=1.1pt}
  ]
    \node[sw] (s0) at (0,2.4) {NVSwitch 0};
    \node[sw] (s1) at (4,2.4) {NVSwitch 1};

    \node[gpu] (g0) at (-1.8,0.7) {GPU0};
    \node[gpu] (g1) at (-0.6,0.7) {GPU1};
    \node[gpu] (g2) at (0.6,0.7) {GPU2};
    \node[gpu] (g3) at (1.8,0.7) {GPU3};

    \node[gpu] (g4) at (2.8,0.7) {GPU4};
    \node[gpu] (g5) at (4.0,0.7) {GPU5};
    \node[gpu] (g6) at (5.2,0.7) {GPU6};
    \node[gpu] (g7) at (6.4,0.7) {GPU7};

    \foreach \g in {g0,g1,g2,g3,g4,g5,g6,g7} {
      \draw[link] (\g.north) -- (s0.south);
      \draw[link] (\g.north) -- (s1.south);
    }

    \draw[link, dashed] (s0) -- (s1);

    \node[align=center] at (2.3,-0.3) {每张 GPU 通过 NVLink 连接到交换芯片，形成更均匀的 GPU fabric};
  \end{tikzpicture}
  \caption{NVSwitch 全互联抽象示意：GPU 通过 NVLink 接入交换芯片，由交换芯片转发 GPU-to-GPU 流量。}
  \label{fig:nvswitch-all-to-all}
\end{figure}

在这种架构中，GPU-to-GPU 通信不再强依赖某两个 GPU 是否存在直接物理链路，而是通过交换 fabric 获得更均匀的路径。这对大模型训练非常重要，因为 tensor parallel、expert parallel 和一些 all-to-all workload 希望任意 GPU pair 之间都有可预测的通信性能。

\subsection{HGX / DGX 拓扑}
\label{subsec:hgx-dgx-topology}

HGX 和 DGX 是 NVIDIA 数据中心多 GPU 平台中常见的系统形态。不同代际和型号的具体连接方式不同，但它们的共同目标是把多张 GPU、NVLink、NVSwitch、CPU、NIC、PCIe 等组件组织成适合 AI 训练与推理的节点。

从 AI Infra 角度看，关心的不是平台名称本身，而是以下几个拓扑问题：

\begin{itemize}
  \item 节点内有多少张 GPU？
  \item GPU 之间是否通过 NVSwitch 实现高带宽互联？
  \item GPU 与 NIC 的距离如何？是否支持高效 GPU Direct RDMA？
  \item GPU 与 CPU socket 的 NUMA 关系如何？
  \item 多个 NVSwitch 或交换平面之间是否存在瓶颈？
  \item 节点内拓扑是否对 NCCL 可见且被正确识别？
\end{itemize}

一个常见设计原则是：节点内的高速 GPU fabric 用于 scale-up，并优先承载 TP、EP 等通信密集型并行；节点间网络用于 scale-out，并承载 DP、PP 跨节点 stage 或更大规模的 EP 通信。换句话说，拓扑决定并行维度如何放置。

\subsection{多 GPU 全互联}
\label{subsec:multi-gpu-full-connectivity}

NVSwitch 最大的工程价值之一，是让多 GPU 节点更接近“任意 GPU 到任意 GPU 都高带宽”的抽象。这种抽象简化了上层系统设计。例如在 8 GPU NVSwitch 节点中，TP group 可以自然设置为 8，模型每层内部的 collective 都在节点内高速 fabric 上完成。

但是，“全互联”不等于通信无成本。即使 NVSwitch 提供了很高带宽，通信仍然需要时间，也仍然可能与计算争用资源。对于每层都通信的 tensor parallel，通信量和通信频率仍然需要谨慎估算。

以 tensor parallel 中的 all-reduce 为例，假设每层输出 activation 大小为 $B \times S \times H$，其中 $B$ 是 micro batch size，$S$ 是 sequence length，$H$ 是 hidden size。如果 TP size 为 $p$，某些实现中每层可能需要对局部输出做 all-reduce。通信数据量大致与 activation size 成正比。随着层数增加，这类通信在 step time 中会累积成明显开销。

\begin{figure}[H]
  \centering
  \begin{tikzpicture}[
    font=\small,
    gpu/.style={draw, rounded corners, minimum width=1.1cm, minimum height=0.6cm, align=center},
    sw/.style={draw, rounded corners, minimum width=2.3cm, minimum height=0.8cm, align=center, fill=gray!10},
    arrow/.style={-{Stealth[length=2mm]}, thick},
    link/.style={line width=1.0pt}
  ]
    \node[sw] (sw) at (0,2.4) {NVSwitch Fabric};

    \foreach \i/\x/\y in {
      0/-3/0.6,
      1/-2/0.6,
      2/-1/0.6,
      3/0/0.6,
      4/1/0.6,
      5/2/0.6,
      6/3/0.6,
      7/4/0.6
    } {
      \node[gpu] (g\i) at (\x,\y) {G\i};
      \draw[link] (g\i.north) -- (sw.south);
    }

    \draw[arrow, bend left=35] (g0.south) to node[below] {\scriptsize collective chunk} (g7.south);
    \draw[arrow, bend left=25] (g2.south) to (g5.south);
    \draw[arrow, bend right=25] (g6.south) to (g1.south);

    \node[align=center] at (0,-0.85) {NCCL collective 将大 tensor 切成多个 chunk，在 GPU fabric 中并行传输与归约};
  \end{tikzpicture}
  \caption{NVSwitch 节点中的 collective 抽象：通信库将数据切分并在 GPU fabric 中并行传输。}
  \label{fig:nvswitch-collective}
\end{figure}

\subsection{对 TP / PP / DP 的意义}
\label{subsec:nvswitch-parallelism-impact}

NVSwitch 对不同并行策略的意义不同。

\subsubsection{对 Tensor Parallelism 的意义}

TP 是最依赖节点内高速互联的并行方式之一。因为 TP 把单层内部的矩阵计算切到多张 GPU 上，每一层都可能引入通信。如果节点内有 NVSwitch，TP group 内 GPU 之间通信更均匀、更高带宽，因此可以支持更大的 TP size 或更高效的层内通信。

但是 TP size 不是越大越好。TP size 增大后，每张 GPU 上的 matmul 规模变小，Tensor Core 利用率可能下降；通信次数不一定减少，通信相对占比可能上升。因此，选择 TP size 时要同时考虑：

\begin{itemize}
  \item 单卡是否能放下 shard 后的参数和 activation；
  \item shard 后的 GEMM shape 是否仍然高效；
  \item all-reduce / all-gather / reduce-scatter 通信是否可接受；
  \item TP group 是否完全位于 NVSwitch 高速域内；
  \item 是否需要与 PP、DP、ZeRO 组合。
\end{itemize}

\subsubsection{对 Pipeline Parallelism 的意义}

PP 把模型按层切分成多个 stage。相邻 stage 之间需要传递 activation 和 gradient。相比 TP，PP 的通信频率通常较低，但通信发生在 pipeline 边界上，延迟会影响 bubble。NVSwitch 可以降低同节点 PP stage 之间的数据传输成本，但 PP 的主要瓶颈往往还包括负载均衡、micro-batch 数量、bubble、activation memory 等。

如果一个节点内有 8 张 GPU，可以把连续几个 pipeline stage 放在同一节点内，使 stage 间通信走 NVLink/NVSwitch；跨节点 PP 则需要走 RDMA 网络。通常应尽量避免高频、细粒度通信跨越低带宽路径。

\subsubsection{对 Data Parallelism 的意义}

DP 的主要通信是梯度同步。单节点内 DP 可以通过 NVLink/NVSwitch 加速 all-reduce；多节点 DP 则还要经过 NIC 和 RDMA 网络。大规模训练中，NCCL 通常会构造分层通信：节点内先利用 NVLink/NVSwitch 聚合，节点间再通过 InfiniBand / RoCE 通信，最后节点内广播或分发结果。

这类 hierarchical collective 的效率取决于节点内和节点间两层网络是否都被正确利用。如果节点内 NVLink 可用但 NCCL 没有识别拓扑，或者节点间 NIC 绑定不合理，整体 all-reduce 性能都会下降。

\begin{table}[H]
  \centering
  \caption{NVLink / NVSwitch 对不同并行策略的影响}
  \label{tab:nvlink-parallelism-impact}
  \begin{tabular}{p{3.0cm}p{4.0cm}p{4.2cm}p{3.5cm}}
    \toprule
    并行策略 & 主要通信 & 对节点内互联的依赖 & 拓扑建议 \\
    \midrule
    Tensor Parallel & all-reduce、all-gather、reduce-scatter & 极高，每层频繁通信 & TP group 尽量放在 NVSwitch / NVLink 高速域内 \\
    Pipeline Parallel & send / recv activation 与 gradient & 中等，边界通信影响 bubble & 相邻 stage 尽量放在通信路径短的位置 \\
    Data Parallel & gradient all-reduce 或 reduce-scatter & 中到高，取决于 batch 和 bucket & 节点内聚合利用 NVLink，节点间利用 RDMA \\
    ZeRO & reduce-scatter、all-gather & 中到高，取决于 stage 和 bucket & 分层通信，避免频繁跨慢链路 \\
    Expert Parallel & all-to-all token dispatch & 很高，尤其 MoE 场景 & EP group 需要高带宽、均匀拓扑 \\
    Inference TP & 每层 collective & 很高，直接影响 TPOT & 推理 TP group 应优先保持在单节点高速 fabric 内 \\
    \bottomrule
  \end{tabular}
\end{table}

\section{通信性能诊断}
\label{sec:communication-diagnosis}

多 GPU 通信问题经常表现为“GPU 利用率低”“训练扩展性差”“8 卡不比 4 卡快多少”“推理 TPOT 异常高”“NCCL 卡住或超时”。这些现象背后的原因可能是硬件拓扑、NCCL 配置、驱动、容器权限、NUMA、PCIe ACS、IB/RoCE、rank mapping、框架 bucket 设置、数据 shape 等。诊断时必须从事实出发。

\subsection{\texttt{nvidia-smi topo}}
\label{subsec:nvidia-smi-topo}

\texttt{nvidia-smi topo -m} 是查看单机 GPU 拓扑的第一步。它会输出 GPU 与 GPU、GPU 与 NIC、GPU 与 CPU affinity 之间的关系。典型输出包含 GPU pair 之间的连接标记，以及 CPU affinity、NUMA affinity 等信息。

\begin{lstlisting}[language=bash,caption={查看单机 GPU 拓扑},label={lst:nvidia-smi-topo}]
nvidia-smi topo -m
\end{lstlisting}

一个简化的输出可能类似下面这样：

\begin{lstlisting}[language=bash,caption={nvidia-smi topo -m 的简化示意输出},label={lst:nvidia-smi-topo-output}]
        GPU0    GPU1    GPU2    GPU3    NIC0    CPU Affinity
GPU0     X      NVL     NVL     SYS     PIX     0-31
GPU1    NVL      X      NVL     SYS     PIX     0-31
GPU2    NVL     NVL      X      SYS     PXB     0-31
GPU3    SYS     SYS     SYS      X      PHB     32-63
NIC0    PIX     PIX     PXB     PHB      X
\end{lstlisting}

上面的示例不是某台真实机器的标准输出，而是为了说明几个重点：

\begin{itemize}
  \item GPU0、GPU1、GPU2 之间存在较好的连接；
  \item GPU3 与其他 GPU 之间路径较差，可能跨 socket 或 host bridge；
  \item NIC0 与不同 GPU 的距离不同；
  \item 如果把 TP group 放在 GPU0--GPU3 上，可能引入慢链路；
  \item 如果跨节点通信主要使用 NIC0，那么 GPU 与 NIC 的亲和性也需要考虑。
\end{itemize}

在真实环境中，不同驱动版本、GPU 型号和系统拓扑显示的标记可能不同。关键不是记住所有缩写，而是理解它们反映的是通信路径差异。遇到性能问题时，应把 \texttt{nvidia-smi topo -m} 输出与 rank mapping、NCCL 日志和作业配置放在一起看。

\subsection{NCCL Tests}
\label{subsec:nccl-tests}

NCCL tests 是测试 GPU collective 性能的常用工具。它可以测量 all-reduce、all-gather、reduce-scatter、broadcast、all-to-all 等操作在不同消息大小、不同 GPU 数量下的带宽和延迟。常见命令如下：

\begin{lstlisting}[language=bash,caption={运行 NCCL all-reduce 性能测试的示例},label={lst:nccl-tests-allreduce}]
# 8 GPU 单机 all-reduce 测试
./build/all_reduce_perf -b 8 -e 4G -f 2 -g 8
\end{lstlisting}

其中：

\begin{itemize}
  \item \texttt{-b} 表示起始消息大小；
  \item \texttt{-e} 表示结束消息大小；
  \item \texttt{-f} 表示每轮消息大小的增长倍数；
  \item \texttt{-g} 表示每个进程使用的 GPU 数量。
\end{itemize}

多节点测试时，通常会结合 MPI 或其他 launcher：

\begin{lstlisting}[language=bash,caption={多节点 NCCL all-reduce 测试的概念示例},label={lst:nccl-tests-multinode}]
mpirun -np 16 -H node0:8,node1:8 \
  -x NCCL_DEBUG=INFO \
  -x NCCL_SOCKET_IFNAME=eth0 \
  ./build/all_reduce_perf -b 8M -e 4G -f 2 -g 1
\end{lstlisting}

单机 NVLink/NVSwitch 诊断时，应优先测试节点内 2 GPU、4 GPU、8 GPU 的 all-reduce、all-gather、reduce-scatter。若测试结果明显低于同类机器预期，则说明问题可能在硬件拓扑、驱动、NCCL、容器权限或系统配置，而不一定是训练框架问题。

\subsection{Bus Bandwidth 与 Algorithm Bandwidth}
\label{subsec:busbw-algbw}

NCCL tests 输出中经常包含 \texttt{algbw} 和 \texttt{busbw}。很多初学者会困惑：为什么同一个通信操作有两个带宽？

可以粗略理解为：

\begin{itemize}
  \item \textbf{Algorithm bandwidth}：从算法输入输出角度看到的有效数据处理速度。
  \item \textbf{Bus bandwidth}：根据 collective 类型和通信模式折算出的底层链路带宽利用估计。
\end{itemize}

不同 collective 的数据移动量与用户 tensor size 不完全相同。例如 ring all-reduce 中，每个 rank 并不是简单发送一次完整 tensor，而是分块、多步传输、归约、再分发。NCCL tests 用 bus bandwidth 帮助用户更接近地理解底层互联是否被充分利用。

对工程诊断来说，重点不是纠结某一个数字，而是比较：

\begin{itemize}
  \item 同一台机器上不同 GPU 数量的扩展曲线是否合理；
  \item 不同消息大小下，带宽是否能在大消息时趋于稳定；
  \item 单机结果是否明显优于跨节点结果；
  \item NVLink/NVSwitch 节点是否明显优于普通 PCIe 节点；
  \item 同类节点之间是否存在异常离群值。
\end{itemize}

如果 8 GPU all-reduce 的带宽远低于预期，可以进一步尝试：

\begin{lstlisting}[language=bash,caption={打开 NCCL 调试日志和拓扑信息},label={lst:nccl-debug-topo}]
export NCCL_DEBUG=INFO
export NCCL_DEBUG_SUBSYS=INIT,GRAPH,TOPO,COLL
./build/all_reduce_perf -b 8M -e 4G -f 2 -g 8
\end{lstlisting}

NCCL 日志中通常会显示拓扑识别、通信 channel、算法选择、网络接口选择等信息。虽然日志较长，但对排查拓扑误识别、P2P 不可用、NIC 选择错误非常有帮助。

\subsection{常见拓扑误配置}
\label{subsec:common-topology-misconfig}

多 GPU 通信问题经常来自配置，而不是硬件本身。下面列出一些常见问题。

\subsubsection{问题一：容器看不到完整拓扑}

在容器环境中，如果 GPU device、驱动库、\texttt{/sys} 拓扑信息或相关权限没有正确挂载，NCCL 可能无法正确识别 GPU 拓扑。表现为：

\begin{itemize}
  \item NCCL 日志中拓扑信息异常；
  \item GPU P2P 不可用；
  \item 同一节点内通信性能远低于裸机；
  \item \texttt{nvidia-smi topo -m} 在容器内外输出不同。
\end{itemize}

排查时应比较宿主机与容器内的 \texttt{nvidia-smi topo -m}、NCCL tests 和驱动库版本。

\subsubsection{问题二：Rank Mapping 不符合拓扑}

框架可能默认按 GPU index 分配 rank，但 GPU index 不一定与拓扑最优顺序一致。比如 GPU0--GPU3 是一个高速域，GPU4--GPU7 是另一个高速域，如果 TP group 被交错放置为 \{0,4,1,5\}，通信会跨越慢路径。解决方法包括：

\begin{itemize}
  \item 显式设置 \texttt{CUDA\_VISIBLE\_DEVICES} 的顺序；
  \item 在 launcher 中控制 local rank 到 GPU 的映射；
  \item 使用框架提供的 topology-aware rank mapping；
  \item 根据 \texttt{nvidia-smi topo -m} 和 NCCL tests 验证 group 划分。
\end{itemize}

\subsubsection{问题三：GPU 与 NIC 亲和性错误}

多节点训练中，GPU-to-NIC 的路径非常重要。如果某个 rank 使用的 GPU 离实际通信 NIC 很远，跨节点通信可能绕过不理想的 PCIe 或 NUMA 路径。表现为：

\begin{itemize}
  \item 单机 NCCL tests 正常，多机性能很差；
  \item 某些 rank 带宽明显低；
  \item NCCL 日志显示选择了非预期网卡；
  \item 跨 socket 流量异常。
\end{itemize}

常用排查方向包括检查 \texttt{nvidia-smi topo -m} 中 GPU 与 NIC 的关系，设置 \texttt{NCCL\_IB\_HCA}、\texttt{NCCL\_SOCKET\_IFNAME} 等环境变量，并结合下一章的 RDMA / RoCE 网络诊断方法。

\subsubsection{问题四：PCIe ACS 或 IOMMU 影响 P2P}

某些服务器 BIOS 或系统配置可能启用 PCIe ACS、IOMMU 等功能，影响 P2P 或 GPUDirect 路径。具体行为与平台有关，不能一概而论。表现可能包括 P2P bandwidth 异常、GPU Direct RDMA 不可用、NCCL fallback 到较慢路径等。排查时应参考硬件厂商、驱动和 NCCL 文档，并与平台团队确认 BIOS 设置。

\subsubsection{问题五：错误理解 GPU 利用率}

通信瓶颈下，\texttt{nvidia-smi} 看到的 GPU utilization 可能不高，但这并不意味着 GPU “没有任务”。它可能在等待通信、等待同步、等待 CPU launch，或者 NCCL kernel 本身没有表现为高 SM utilization。此时应结合 Nsight Systems timeline 判断 GPU 是否处于通信等待，而不是只看单一利用率数字。

\section{从通信视角重新理解并行策略}
\label{sec:parallelism-from-communication}

在前面的章节中，我们已经学习过 DP、TP、PP、ZeRO 和 3D 并行。本节从 NVLink/NVSwitch 的角度重新审视它们。

\subsection{为什么 TP 更偏好节点内高速互联}
\label{subsec:tp-prefers-intranode}

TP 的通信特点是高频、细粒度、层内发生。以一个 decoder-only Transformer 为例，每一层 attention 和 FFN 内部都可能有 TP collective。假设模型有 80 层，那么一次 forward 可能触发大量跨 GPU 通信。训练还要加上 backward 中的通信。

这意味着 TP group 跨越慢链路会非常昂贵。即使单次 all-reduce 数据量不算极大，层数乘以 micro-batch 数后，总通信开销也会变成 step time 的重要组成部分。因此：

\begin{itemize}
  \item 单机 8 GPU NVSwitch 节点上，TP size = 8 常常是自然选择；
  \item 普通 PCIe 节点上，TP size 可能需要更谨慎；
  \item 跨节点 TP 只有在节点间网络足够强且框架高度优化时才值得考虑；
  \item 对推理服务，TP 跨节点通常会显著增加 TPOT，应尽量避免。
\end{itemize}

\subsection{为什么 DP 可以更自然地跨节点}
\label{subsec:dp-cross-node}

DP 的通信主要发生在 backward 梯度同步阶段。虽然梯度同步数据量很大，但它不像 TP 那样每层 forward 都要通信。通过 bucket、overlap、reduce-scatter 等技术，DP 通信可以与 backward 计算重叠。因此，DP 更适合作为跨节点扩展的维度。

当然，这不意味着 DP 通信不重要。在千卡训练中，跨节点 all-reduce 仍可能成为主要瓶颈。只是从拓扑映射角度看，通常会优先把 TP 放在节点内 NVLink/NVSwitch 域，把 DP 扩展到节点间 RDMA 网络。

\subsection{为什么 MoE 对 all-to-all 拓扑敏感}
\label{subsec:moe-alltoall-topology}

MoE 模型中的 expert parallel 会产生 token dispatch 和 combine。每个 rank 上的 token 根据 router 结果被发送到不同 expert 所在 rank，这通常对应 all-to-all 或类似通信模式。All-to-all 比 all-reduce 更依赖均匀拓扑，因为每个 rank 都可能与多个其他 rank 交换不同大小的数据。

在 NVSwitch 节点内，all-to-all 的路径相对均匀；跨节点时，all-to-all 会给网络造成复杂的 east-west traffic。MoE 训练和推理中的很多系统优化，例如 expert placement、capacity factor、token dropping、grouped GEMM、hierarchical all-to-all，本质上都是为了同时处理计算不均衡和通信不均衡。

\subsection{推理系统中的 NVLink / NVSwitch}
\label{subsec:nvlink-serving}

推理系统使用 NVLink/NVSwitch 的方式与训练不同。训练更关注 step time 和 scaling efficiency，推理更关注 TTFT、TPOT、吞吐、SLA 和成本。

在 tensor parallel 推理中，每层都可能发生 collective，因此节点内高速互联直接影响 TPOT。尤其是 decode 阶段，每生成一个 token 都要跨过所有层。如果每层 collective 都增加一点延迟，最终会累积成用户可感知的响应变慢。

NVLink/NVSwitch 对推理的典型价值包括：

\begin{itemize}
  \item 支持更大模型在单节点内 TP 部署；
  \item 降低层内 collective 延迟，改善 TPOT；
  \item 支持多 GPU 共享模型 shard，提高单实例吞吐；
  \item 在多实例部署中减少跨 GPU 数据移动成本；
  \item 对 prefill/decode 分离架构提供更灵活的节点内数据路径。
\end{itemize}

但是推理也会遇到新问题。比如同一节点上部署多个模型实例时，它们可能共享 NVLink/NVSwitch fabric 和 HBM 带宽。即使每个实例单独测试都正常，混部后也可能出现干扰。因此，AI serving 调度器不能只看 GPU 显存剩余量，还要关注通信和带宽压力。

\section{实践：单机多卡通信排查流程}
\label{sec:single-node-comm-debug-workflow}

本节给出一个面向工程实践的排查流程。假设你发现一个 8 GPU 训练作业在某台机器上速度异常，比同型号机器慢很多，或者 8 卡扩展效率明显低于预期。

\subsection{第一步：确认硬件与驱动可见性}
\label{subsec:debug-step-hardware}

首先确认 GPU 是否全部正常可见：

\begin{lstlisting}[language=bash,caption={确认 GPU 可见性},label={lst:debug-gpu-visible}]
nvidia-smi
nvidia-smi -L
\end{lstlisting}

检查内容包括：

\begin{itemize}
  \item GPU 数量是否正确；
  \item 是否有 GPU 处于异常状态；
  \item ECC、温度、功耗、时钟是否异常；
  \item 是否有其他进程占用 GPU；
  \item MIG 是否被启用并改变了资源形态。
\end{itemize}

如果 GPU 基本状态异常，先不要分析 NCCL 或框架配置。

\subsection{第二步：查看拓扑}
\label{subsec:debug-step-topology}

运行：

\begin{lstlisting}[language=bash,caption={查看拓扑},label={lst:debug-topology}]
nvidia-smi topo -m
\end{lstlisting}

保存输出，并重点观察：

\begin{itemize}
  \item GPU 之间是否存在 NVLink / NVSwitch 路径；
  \item 是否有某些 GPU pair 显示为明显更差的路径；
  \item GPU 与 NIC 的亲和性；
  \item CPU affinity 与 NUMA affinity；
  \item 容器内外输出是否一致。
\end{itemize}

\subsection{第三步：运行 P2P 测试}
\label{subsec:debug-step-p2p}

如果环境中有 CUDA samples，可以运行 P2P 带宽延迟测试：

\begin{lstlisting}[language=bash,caption={运行 P2P bandwidth/latency 测试},label={lst:debug-p2p}]
./p2pBandwidthLatencyTest
\end{lstlisting}

观察 GPU pair 之间带宽是否符合拓扑预期。若某些应走 NVLink 的 pair 表现很差，可能是 P2P 未启用、拓扑不可见、驱动或系统配置问题。

\subsection{第四步：运行 NCCL 单机测试}
\label{subsec:debug-step-nccl-single}

运行 all-reduce、all-gather、reduce-scatter 测试：

\begin{lstlisting}[language=bash,caption={单机 NCCL collective 测试},label={lst:debug-nccl-single}]
export NCCL_DEBUG=INFO
./build/all_reduce_perf -b 8M -e 4G -f 2 -g 8
./build/all_gather_perf -b 8M -e 4G -f 2 -g 8
./build/reduce_scatter_perf -b 8M -e 4G -f 2 -g 8
\end{lstlisting}

如果这些基础测试异常，则说明问题在框架之下。此时应先修复 NCCL / 拓扑 / 系统问题，而不是调模型代码。

\subsection{第五步：检查 NCCL 日志}
\label{subsec:debug-step-nccl-log}

打开更详细的 NCCL 日志：

\begin{lstlisting}[language=bash,caption={NCCL 拓扑调试日志},label={lst:debug-nccl-log}]
export NCCL_DEBUG=INFO
export NCCL_DEBUG_SUBSYS=INIT,GRAPH,TOPO,COLL
export NCCL_TOPO_DUMP_FILE=/tmp/nccl_topo.xml
\end{lstlisting}

关注 NCCL 是否识别到预期的 NVLink/NVSwitch 路径，是否启用了 P2P，是否选择了合理的 channel 和算法。对于多节点，还要关注 NIC 选择和 IB/RoCE 路径。

\subsection{第六步：对照训练框架配置}
\label{subsec:debug-step-framework}

基础通信正常后，再检查训练框架：

\begin{itemize}
  \item TP size、PP size、DP size 是否符合拓扑；
  \item \texttt{CUDA\_VISIBLE\_DEVICES} 顺序是否合理；
  \item local rank 到 GPU 的映射是否正确；
  \item DDP bucket、ZeRO stage、overlap 设置是否合理；
  \item micro batch size 是否过小导致计算无法覆盖通信；
  \item 是否有意外 CPU/GPU 同步或数据拷贝。
\end{itemize}

最后用 Nsight Systems 观察真实训练 step 的 timeline，确认通信是否与计算重叠、是否存在 GPU 空洞、是否有 NCCL kernel 占据关键路径。

\section{设计原则与经验法则}
\label{sec:nvlink-design-principles}

本节总结若干实践原则。

\subsection{原则一：通信密集 group 优先放在高速域内}
\label{subsec:principle-fast-domain}

TP、EP、推理 TP 这类高频通信 group 应优先放在 NVLink/NVSwitch 域内。DP 可以更自然地跨节点扩展，PP 则要关注相邻 stage 的路径。

\subsection{原则二：不要只看 GPU 数量，要看拓扑}
\label{subsec:principle-not-just-gpu-count}

两台机器都可能有 8 张 GPU，但如果一台是 NVSwitch all-to-all，另一台是普通 PCIe 拓扑，它们对 TP 和 MoE 的支持能力完全不同。调度系统在分配作业时，应把 GPU 拓扑作为资源属性，而不是只记录 GPU count。

\subsection{原则三：先测 NCCL，再测框架}
\label{subsec:principle-test-nccl-first}

当训练扩展性异常时，先用 NCCL tests 验证基础通信，再分析框架。否则很容易在模型代码中浪费时间。基础通信不正常时，框架层调参通常无法根治问题。

\subsection{原则四：通信和计算要一起看}
\label{subsec:principle-comm-compute}

高带宽互联不是为了让通信无限快，而是为了让通信更容易被计算覆盖。优化目标不是单独最大化 NCCL 带宽，而是最小化端到端 step time 或 TPOT。因此要结合 profiler 看通信是否在关键路径上。

\subsection{原则五：推理调度也要拓扑感知}
\label{subsec:principle-serving-topology}

推理系统常常只按显存和 GPU 数量部署实例，但多卡推理实例对 NVLink/NVSwitch 依赖很强。TP 推理实例应尽量独占或稳定绑定高速互联 group，避免与其他通信密集 workload 混部。

\section{本章小结}
\label{sec:nvlink-nvswitch-summary}

本章讨论了单机多卡通信中的核心硬件路径：PCIe、GPU Direct P2P、NVLink 与 NVSwitch。

\begin{itemize}
  \item PCIe 是通用 I/O 总线，生态成熟，但对频繁 GPU-to-GPU 通信来说，带宽、延迟和拓扑均匀性都可能成为瓶颈。
  \item GPU Direct P2P 允许 GPU 之间直接访问 peer memory，避免不必要的 host staging，是高效 GPU-to-GPU 通信的基础能力。
  \item NVLink 提供面向 GPU-to-GPU 的高速点对点互联，显著改善节点内通信能力，尤其适合 tensor parallel、pipeline stage 传输和多卡推理。
  \item NVSwitch 把 NVLink 从点对点连接扩展为更均匀的 GPU fabric，使多 GPU 节点更接近 all-to-all 高带宽互联。
  \item NCCL 会根据拓扑选择 collective 算法和通信路径。相同的 \texttt{all\_reduce}，在 PCIe、NVLink、NVSwitch 和跨节点 RDMA 上性能可能完全不同。
  \item \texttt{nvidia-smi topo -m}、P2P bandwidth test、NCCL tests、NCCL debug log 和 Nsight Systems 是诊断单机多卡通信问题的核心工具。
  \item 并行策略必须拓扑感知：TP 和 EP 优先放在节点内高速互联域，DP 更适合跨节点扩展，PP 要关注相邻 stage 的通信路径。
\end{itemize}

至此，我们已经理解了一台机器内部多张 GPU 如何通过高速互联形成 scale-up 计算域。下一章将继续向外扩展：当训练和推理跨越多台服务器时，GPU 之间的数据必须经过 NIC、交换机和数据中心网络。为了让跨节点通信不吞噬 GPU 利用率，我们需要理解 RDMA、RoCE、InfiniBand、GPUDirect RDMA、PFC、ECN 以及 AI 集群网络拓扑。这将进入第 19 章：RDMA / RoCE 网络基础。

\chapter{RDMA / RoCE 网络基础}
\label{chap:rdma-roce}

上一章讨论了单机多卡通信：在一台服务器内部，GPU 可以通过 PCIe、NVLink 与 NVSwitch 交换数据。但是，当模型规模、数据规模或并发请求规模继续扩大时，单机已经不够。训练一个千亿参数模型可能需要数百到数千张 GPU；服务一个高并发大模型系统可能需要跨越多个节点部署 prefill、decode、embedding、rerank、tool calling 或多模型实例。此时，通信路径从“GPU 到 GPU”进一步扩展为：

\[
  \text{GPU} \rightarrow \text{PCIe/NVLink} \rightarrow \text{NIC} \rightarrow \text{Switch Fabric} \rightarrow \text{NIC} \rightarrow \text{GPU}.
\]

这就是本章要讨论的问题：多机 GPU 集群中的网络如何影响训练和推理性能。

在传统 Web 服务中，网络通常承载 RPC、HTTP、数据库访问、日志、对象存储等流量。对这些系统而言，毫秒级延迟往往已经足够，吞吐和可用性更重要。而在大模型训练中，GPU 每一步训练都可能需要跨节点同步数 GB 到数十 GB 的梯度、参数分片或激活张量。GPU 计算速度越快，网络越容易成为瓶颈。如果网络跟不上，GPU 就会在通信 barrier 前等待，昂贵的 Tensor Core 将处于空转状态。

因此，AI 集群网络不是普通数据中心网络的简单放大版。它必须同时满足：

\begin{itemize}
  \item 高带宽：支持大规模 all-reduce、reduce-scatter、all-gather、all-to-all；
  \item 低延迟：降低小消息、pipeline stage、decode 通信和控制路径等待；
  \item 低 CPU 开销：避免 CPU 参与数据搬运，减少内核协议栈开销；
  \item 拓扑可预测：让 collective 算法能够充分利用网络结构；
  \item 拥塞可控：避免少量热点流量拖慢整个训练作业；
  \item 可观测：能够定位链路、交换机、队列、NIC、NCCL 和作业层面的瓶颈。
\end{itemize}

RDMA、InfiniBand、RoCE 与 GPUDirect RDMA 正是为这类问题服务的关键技术。本章会从工程视角解释它们的基本概念，而不是陷入协议细节。读完本章后，你应该能够理解为什么 AI 集群常常选择 InfiniBand 或 RoCE 网络，为什么 RoCE 需要关注 PFC、ECN 和 lossless Ethernet，为什么 GPUDirect RDMA 能显著改善跨节点 GPU 通信，以及如何从 NCCL、网卡计数器和网络拓扑角度排查训练扩展性问题。

\section{多机训练为什么依赖高性能网络}
\label{sec:why-high-performance-network}

多机训练的核心矛盾是：计算被分散到多张 GPU 上，但模型更新仍然需要在逻辑上保持一致。无论使用 DP、TP、PP、ZeRO 还是 MoE，只要跨越节点，就会产生网络通信。随着 GPU 数量增加，单卡计算时间可能下降，但通信路径变长、通信规模变大、同步等待增加，最终扩展效率会受到网络限制。

\subsection{跨节点 All-Reduce}
\label{subsec:cross-node-allreduce}

数据并行中最典型的跨节点通信是梯度 all-reduce。假设有 $N$ 张 GPU，每张 GPU 持有一份完整模型副本，并处理不同 mini-batch。反向传播后，每张 GPU 得到本地梯度 $\nabla \theta_i$。为了让所有副本保持一致，需要计算平均梯度：

\[
  \nabla \theta = \frac{1}{N}\sum_{i=1}^{N}\nabla \theta_i.
\]

这通常通过 all-reduce 完成。每个 rank 输入自己的梯度 buffer，通信结束后每个 rank 都得到 reduce 后的完整结果。对于一个参数量为 $P$ 的模型，如果梯度使用 BF16 或 FP16，每个参数的梯度约为 2 字节，则一次梯度同步的逻辑数据规模大约为：

\[
  2P\ \text{Bytes}.
\]

对于十亿参数模型，这是 GB 级数据；对于百亿参数模型，梯度同步规模会更大。虽然实际系统会使用 bucket、overlap、reduce-scatter、ZeRO 等方法降低峰值显存和改善重叠，但通信总量仍然不容忽视。

在 ring all-reduce 中，每个 rank 的通信量大致为：

\[
  V_{\text{ring}} \approx 2 \cdot \frac{N-1}{N} \cdot S,
\]

其中 $S$ 是待 all-reduce tensor 的大小。随着 $N$ 增大，单 rank 通信量接近 $2S$。如果网络有效带宽为 $B$，仅从带宽角度估计通信时间下界为：

\[
  T_{\text{comm}} \gtrsim \frac{2S}{B}.
\]

这个估计忽略了延迟、协议开销、拓扑、拥塞和实现细节，但足以说明一个问题：当梯度 buffer 很大时，网络带宽直接决定同步时间。

\begin{figure}[H]
  \centering
  \begin{tikzpicture}[
    font=\small,
    nodebox/.style={draw, rounded corners, minimum width=2.0cm, minimum height=1.0cm, align=center},
    gpu/.style={draw, rounded corners, minimum width=1.0cm, minimum height=0.55cm, align=center},
    sw/.style={draw, rounded corners, minimum width=2.0cm, minimum height=0.75cm, align=center, fill=gray!10},
    arrow/.style={-{Stealth[length=2mm]}, thick}
  ]
    \node[nodebox] (n0) at (-4,1.4) {Node 0};
    \node[gpu] (g00) at (-4.55,0.75) {G0};
    \node[gpu] (g01) at (-3.45,0.75) {G1};

    \node[nodebox] (n1) at (0,1.4) {Node 1};
    \node[gpu] (g10) at (-0.55,0.75) {G2};
    \node[gpu] (g11) at (0.55,0.75) {G3};

    \node[nodebox] (n2) at (4,1.4) {Node 2};
    \node[gpu] (g20) at (3.45,0.75) {G4};
    \node[gpu] (g21) at (4.55,0.75) {G5};

    \node[sw] (leaf0) at (-2,3.1) {Leaf};
    \node[sw] (leaf1) at (2,3.1) {Leaf};
    \node[sw] (spine) at (0,4.5) {Spine};

    \draw[arrow] (n0.north) -- (leaf0.south);
    \draw[arrow] (n1.north) -- (leaf0.south);
    \draw[arrow] (n1.north) -- (leaf1.south);
    \draw[arrow] (n2.north) -- (leaf1.south);
    \draw[arrow] (leaf0.north) -- (spine.south);
    \draw[arrow] (leaf1.north) -- (spine.south);

    \draw[arrow, dashed] (g00.south) .. controls (-2.2,-0.7) and (2.2,-0.7) .. (g21.south);
    \node[align=center] at (0,-1.15) {跨节点 collective：GPU 数据经 NIC 与交换网络在节点间移动};
  \end{tikzpicture}
  \caption{多节点训练通信路径：节点内 GPU 通过本地互联到达 NIC，节点间数据经过交换网络传输。}
  \label{fig:multi-node-training-path}
\end{figure}

\subsection{网络带宽与训练扩展性}
\label{subsec:network-bandwidth-scaling}

训练扩展性的目标是：增加 GPU 数量后，单位时间完成的训练 token 或 sample 尽可能线性增长。理想情况下，$N$ 张 GPU 的吞吐是单卡的 $N$ 倍。但现实中，扩展效率总是低于 100\%。一个简化的 step time 模型可以写成：

\[
  T_{\text{step}} = T_{\text{compute}} + T_{\text{comm}} - T_{\text{overlap}} + T_{\text{bubble}} + T_{\text{misc}}.
\]

其中：

\begin{itemize}
  \item $T_{\text{compute}}$ 是 GPU 计算时间；
  \item $T_{\text{comm}}$ 是通信时间；
  \item $T_{\text{overlap}}$ 是计算与通信重叠抵消掉的部分；
  \item $T_{\text{bubble}}$ 是 pipeline bubble、同步等待等空洞；
  \item $T_{\text{misc}}$ 包括数据加载、checkpoint、CPU 调度等其他开销。
\end{itemize}

网络优化的目标不是让 $T_{\text{comm}}$ 变成 0，而是让它尽量小，并尽可能与计算重叠。当 GPU 更快、模型更大、并行维度更多时，通信在 step time 中的占比会上升。尤其在下面几类场景中，网络更容易成为瓶颈：

\begin{itemize}
  \item 使用较小 micro-batch，单步计算时间短，通信难以隐藏；
  \item TP 或 EP 跨节点，通信频率高；
  \item ZeRO-3 每层 all-gather 参数，通信分散且频繁；
  \item MoE all-to-all 产生复杂 east-west traffic；
  \item 网络存在 oversubscription，多个作业互相争用带宽；
  \item RoCE 拥塞控制不当，出现 PFC storm、丢包、重传或长尾延迟。
\end{itemize}

\subsection{Latency-Sensitive 与 Bandwidth-Sensitive}
\label{subsec:latency-vs-bandwidth-sensitive}

不同通信模式对网络的要求不同。可以粗略分为 latency-sensitive 和 bandwidth-sensitive 两类。

\textbf{Latency-sensitive} 通信的特点是消息较小、频率较高、在关键路径上等待。例如：

\begin{itemize}
  \item pipeline parallel 相邻 stage 之间的小 activation 传输；
  \item 推理 decode 阶段每层 tensor parallel collective；
  \item 控制面元数据同步；
  \item 小 batch 或小 tensor 上的 all-reduce；
  \item MoE token dispatch 中大量小粒度消息。
\end{itemize}

这类通信即使数据量不大，也会因为等待和同步影响端到端延迟。优化重点是降低 RTT、减少协议栈开销、避免排队、减少 CPU 参与和减少不必要同步。

\textbf{Bandwidth-sensitive} 通信的特点是消息大、传输时间主要由链路带宽决定。例如：

\begin{itemize}
  \item 数据并行大 bucket 梯度 all-reduce；
  \item ZeRO reduce-scatter / all-gather 大参数分片；
  \item checkpoint 分布式保存；
  \item 大 batch embedding 或 parameter server 风格通信；
  \item 大规模 MoE all-to-all。
\end{itemize}

这类通信的优化重点是提高有效带宽、增加多 rail 并行、选择合适 collective 算法、避免 oversubscription、让通信与计算重叠。

\begin{table}[H]
  \centering
  \caption{Latency-sensitive 与 Bandwidth-sensitive 通信对比}
  \label{tab:latency-bandwidth-communication}
  \begin{tabular}{p{3.2cm}p{4.6cm}p{4.6cm}p{3.0cm}}
    \toprule
    类型 & 特征 & 典型场景 & 优化重点 \\
    \midrule
    Latency-sensitive & 小消息、高频、同步等待明显 & Decode TP、PP stage、控制面同步、小 tensor collective & 降低 RTT、减少 CPU 参与、减少同步和排队 \\
    Bandwidth-sensitive & 大消息、传输时间由带宽主导 & DP all-reduce、ZeRO、checkpoint、MoE 大规模 all-to-all & 提高链路利用率、多 rail、拓扑感知 collective \\
    Mixed & 消息大小变化大，既受延迟也受带宽影响 & MoE token dispatch、动态 batch serving、参数分片加载 & 分层调度、分桶、合并小消息、拥塞控制 \\
    \bottomrule
  \end{tabular}
\end{table}

\subsection{通信瓶颈如何吞噬 GPU 利用率}
\label{subsec:communication-bottleneck-gpu-utilization}

通信瓶颈经常表现为 GPU 利用率下降，但它的根源并不在 GPU kernel 本身。一个训练 step 中，GPU 可能先执行 forward 和 backward kernel，然后等待 NCCL collective 完成，再进入下一层或下一步。如果通信无法与计算重叠，timeline 上就会出现明显空洞。

更隐蔽的情况是：NCCL kernel 正在 GPU 上运行，但 SM utilization 并不高。此时从 \texttt{nvidia-smi} 看 GPU 可能不是完全空闲，但 Tensor Core 没有执行模型计算。训练吞吐仍然下降。对于 AI Infra 工程师，关键不是追求某个单一利用率指标，而是判断 GPU 时间到底花在：

\begin{itemize}
  \item 计算 kernel；
  \item NCCL 通信 kernel；
  \item memory copy；
  \item CPU launch 等待；
  \item stream synchronization；
  \item dataloader 或存储等待；
  \item checkpoint I/O。
\end{itemize}

这需要结合 Nsight Systems、NCCL 日志、网络计数器和训练框架日志一起看。一个成熟的性能分析流程一定是跨层的：从模型 step time 到 NCCL collective，从 GPU timeline 到 NIC counters，从 rank mapping 到交换机队列。

\section{RDMA 基础}
\label{sec:rdma-basics}

RDMA 是 Remote Direct Memory Access 的缩写，即远程直接内存访问。它允许一台机器上的网卡直接读写另一台机器上的内存区域，而不需要远端 CPU 频繁参与数据拷贝，也不需要每个数据包都经过传统内核网络协议栈的完整处理路径。

RDMA 的核心价值可以概括为三点：

\begin{enumerate}
  \item \textbf{低延迟}：绕过传统内核协议栈，减少系统调用、协议处理和上下文切换；
  \item \textbf{低 CPU 开销}：数据搬运由 NIC 和 DMA 引擎完成，CPU 不需要参与每次拷贝；
  \item \textbf{高带宽}：更容易接近网卡和链路的有效带宽上限。
\end{enumerate}

在 AI 集群中，RDMA 通常不是直接由训练代码手写 verbs API，而是被 NCCL、UCX、MPI、SHARP、NVSHMEM 或框架通信层间接使用。尽管如此，理解 RDMA 的基本对象仍然非常有帮助，因为很多故障和性能问题会暴露为 QP、CQ、MR、GID、LID、MTU、PFC、ECN、RoCE mode 等配置问题。

\subsection{Kernel Bypass}
\label{subsec:kernel-bypass}

传统 TCP/IP 数据路径通常需要经过操作系统内核协议栈。发送数据时，应用把数据交给内核，内核处理 socket、TCP/IP 协议、拥塞控制、拷贝、网卡队列等；接收数据时，数据包到达网卡后，内核处理中断、协议栈、缓冲区，再把数据交给应用。这个路径通用、可靠、易用，但对高频低延迟通信来说开销较大。

RDMA 的 kernel bypass 思想是：连接建立、资源注册等控制操作仍然需要内核参与，但数据路径尽量绕过内核。应用通过用户态库把 work request 提交给网卡队列，网卡直接从已注册内存读取或写入数据，并在完成后通过 completion queue 通知应用。

\begin{figure}[H]
  \centering
  \begin{tikzpicture}[
    font=\small,
    box/.style={draw, rounded corners, minimum width=2.3cm, minimum height=0.7cm, align=center},
    wide/.style={draw, rounded corners, minimum width=3.2cm, minimum height=0.75cm, align=center},
    arrow/.style={-{Stealth[length=2mm]}, thick}
  ]
    \node[wide] (app1) at (-4,4.0) {Application};
    \node[wide] (kernel1) at (-4,2.8) {Kernel TCP/IP};
    \node[wide] (nic1) at (-4,1.6) {NIC};
    \node[wide] (net1) at (-4,0.4) {Network};

    \draw[arrow] (app1) -- (kernel1);
    \draw[arrow] (kernel1) -- (nic1);
    \draw[arrow] (nic1) -- (net1);

    \node[wide] (app2) at (3.2,4.0) {Application};
    \node[wide] (verbs) at (3.2,2.8) {User-space RDMA Lib};
    \node[wide] (nic2) at (3.2,1.6) {RNIC / HCA};
    \node[wide] (net2) at (3.2,0.4) {Network};

    \draw[arrow] (app2) -- (verbs);
    \draw[arrow] (verbs) -- (nic2);
    \draw[arrow] (nic2) -- (net2);

    \node[box, minimum width=2.4cm] at (-4,-0.8) {传统 TCP 路径};
    \node[box, minimum width=2.4cm] at (3.2,-0.8) {RDMA 数据路径};

    \draw[dashed] (-0.6,-1.3) -- (-0.6,4.5);
    \node[align=center] at (-0.6,4.85) {对比};
  \end{tikzpicture}
  \caption{传统 TCP 与 RDMA 数据路径对比：RDMA 的数据面尽量绕过内核协议栈。}
  \label{fig:tcp-vs-rdma-path}
\end{figure}

Kernel bypass 不是说内核完全不参与。RDMA 仍然需要内核驱动管理设备、权限、内存注册、连接建立和安全隔离。所谓 bypass 主要指高频数据收发路径不再为每个消息经过传统 TCP/IP 内核处理。

\subsection{Zero-Copy}
\label{subsec:zero-copy}

Zero-copy 指数据在发送和接收过程中尽量避免 CPU 参与的额外内存拷贝。传统网络通信中，数据可能经历应用缓冲区到内核缓冲区、内核缓冲区到网卡 DMA buffer 等多次移动。RDMA 中，应用先把一段内存注册给 RDMA 设备，NIC 获得访问这段内存的权限和地址映射。之后，NIC 可以直接通过 DMA 从这段内存读取或写入。

在 AI 训练中，zero-copy 的意义非常大。梯度、参数分片、activation 或 KV Cache 都可能是大 tensor。如果每次跨节点通信都需要 CPU 拷贝一份，CPU 内存带宽、缓存污染和延迟都会变成负担。RDMA 让通信库更接近“从源 buffer 直接送到目标 buffer”的理想路径。

不过 zero-copy 也不是免费的。为了让 NIC 安全、正确地访问内存，需要 memory registration。注册过程可能涉及 pin memory、建立地址转换、设置访问权限等操作，成本不低。因此高性能通信库通常会复用注册内存，避免频繁注册和注销。

\subsection{Queue Pair}
\label{subsec:queue-pair}

Queue Pair，简称 QP，是 RDMA 编程模型中的核心对象。一个 QP 通常包含：

\begin{itemize}
  \item Send Queue：应用向其中提交发送、写、读等 work request；
  \item Receive Queue：用于接收远端发送过来的消息；
  \item 与之关联的 Completion Queue：用于获取完成事件。
\end{itemize}

可以把 QP 理解为 RDMA 连接或通信端点的抽象。不同 QP 类型支持不同通信语义，例如可靠连接、非可靠连接、数据报等。AI 集群中最常见的上层用户通常不直接管理 QP，但 NCCL、UCX、MPI 等通信库会在底层创建并管理这些对象。

\begin{figure}[H]
  \centering
  \begin{tikzpicture}[
    font=\small,
    box/.style={draw, rounded corners, minimum width=2.6cm, minimum height=0.7cm, align=center},
    host/.style={draw, rounded corners, minimum width=4.0cm, minimum height=4.4cm, align=center},
    arrow/.style={-{Stealth[length=2mm]}, thick}
  ]
    \node[host] (h0) at (-3.2,1.8) {};
    \node[anchor=north west] at ($(h0.north west)+(0.2,-0.2)$) {Host A};
    \node[box] (sq0) at (-3.2,2.8) {Send Queue};
    \node[box] (rq0) at (-3.2,1.8) {Receive Queue};
    \node[box] (cq0) at (-3.2,0.8) {Completion Queue};

    \node[host] (h1) at (3.2,1.8) {};
    \node[anchor=north west] at ($(h1.north west)+(0.2,-0.2)$) {Host B};
    \node[box] (sq1) at (3.2,2.8) {Send Queue};
    \node[box] (rq1) at (3.2,1.8) {Receive Queue};
    \node[box] (cq1) at (3.2,0.8) {Completion Queue};

    \draw[arrow] (sq0.east) -- node[above] {Work Request} (rq1.west);
    \draw[arrow] (sq1.west) -- node[below] {Work Request} (rq0.east);
    \draw[arrow] (rq0.south) -- (cq0.north);
    \draw[arrow] (rq1.south) -- (cq1.north);

    \node[align=center] at (0,-0.6) {QP 将用户态提交的 work request 与网卡数据路径连接起来};
  \end{tikzpicture}
  \caption{RDMA Queue Pair 抽象：发送队列、接收队列和完成队列共同描述通信端点。}
  \label{fig:rdma-queue-pair}
\end{figure}

\subsection{Completion Queue}
\label{subsec:completion-queue}

Completion Queue，简称 CQ，用来通知应用某个 work request 已经完成。应用可以轮询 CQ，也可以使用事件通知。高性能通信库通常偏向轮询，因为轮询可以避免中断和上下文切换开销，但会消耗 CPU core。许多 AI 集群会为通信线程、NCCL proxy、MPI progress engine 等预留 CPU 资源，以免这些通信路径被业务线程或系统任务干扰。

CQ 的存在提醒我们：RDMA 降低了 CPU 数据搬运开销，但并不意味着 CPU 完全无关。连接管理、队列推进、completion 处理、错误处理、通信库调度仍然需要 CPU 或专门的硬件机制支持。如果 CPU 过载，网络性能仍然可能下降。

\subsection{Memory Registration}
\label{subsec:memory-registration}

RDMA 设备不能随意访问进程的任意虚拟内存。应用需要先注册内存区域，得到 Memory Region，简称 MR。注册后，设备获得这段内存的访问权限和地址映射，并生成 local key 或 remote key，用于本地或远端访问验证。

Memory registration 的工程意义包括：

\begin{itemize}
  \item 注册内存通常需要 pin 住物理页，防止被操作系统换出或移动；
  \item 注册和注销有成本，不应在热路径中频繁执行；
  \item 通信库通常会维护 registration cache；
  \item 容器、权限、IOMMU、驱动和内核模块都可能影响注册行为；
  \item GPUDirect RDMA 中，注册对象可能是 GPU memory，而不只是 CPU memory。
\end{itemize}

对于大模型训练，频繁分配和释放通信 buffer 可能导致 registration cache miss、内存碎片和性能抖动。因此框架通常会复用 gradient bucket、通信 workspace、参数 shard buffer 等内存区域。

\section{InfiniBand 与 RoCE}
\label{sec:infiniband-and-roce}

RDMA 是一种能力或编程模型，而 InfiniBand 和 RoCE 是实现 RDMA 网络的两类常见方式。它们在 AI 集群中都很常见。

\subsection{InfiniBand 网络}
\label{subsec:infiniband-network}

InfiniBand 是为高性能计算和低延迟通信设计的专用网络技术。它通常由 HCA、InfiniBand switch、subnet manager、专用线缆和协议栈组成。InfiniBand 的优势在于端到端设计较为一致，天然支持 RDMA，延迟低、带宽高、拥塞控制和拓扑管理能力成熟。

在 AI 训练集群中，InfiniBand 常用于跨节点 GPU 通信。NCCL 可以通过 InfiniBand HCA 在节点间传输 collective 数据，并结合 GPUDirect RDMA 实现 GPU memory 与 NIC 的直接数据路径。

InfiniBand 的典型运维关注点包括：

\begin{itemize}
  \item HCA 速率、端口状态和 link width；
  \item subnet manager 是否稳定；
  \item LID、SL、MTU、GID 等配置；
  \item switch 端口错误计数；
  \item fabric 拓扑和路由；
  \item 多 rail 绑定和 rank 到 NIC 的映射；
  \item NCCL 是否正确选择 IB 设备。
\end{itemize}

\subsection{RoCEv1 / RoCEv2}
\label{subsec:rocev1-rocev2}

RoCE 是 RDMA over Converged Ethernet，即在以太网上承载 RDMA。它的吸引力在于可以复用以太网生态：交换机、布线、运维体系、IP 网络能力等。RoCE 主要有两个版本：

\begin{itemize}
  \item \textbf{RoCEv1}：运行在二层以太网上，不能跨三层路由；
  \item \textbf{RoCEv2}：把 RDMA transport 封装在 UDP/IP 上，可以跨三层网络。
\end{itemize}

RoCEv2 是现代数据中心更常见的形态，因为它更容易与 leaf-spine、L3 routing、ECMP 等数据中心网络设计结合。但 RoCE 的挑战也很明显：传统以太网默认不是无损网络，而 RDMA 对丢包非常敏感。如果发生丢包，性能可能急剧下降。因此 RoCE 网络通常需要配置 PFC、ECN、QoS、buffer、拥塞控制和优先级映射。

\begin{figure}[H]
  \centering
  \begin{tikzpicture}[
    font=\small,
    layer/.style={draw, rounded corners, minimum width=8.8cm, minimum height=0.65cm, align=center},
    arrow/.style={-{Stealth[length=2mm]}, thick}
  ]
    \node[layer] (app) at (0,4.2) {NCCL / UCX / MPI / Application};
    \node[layer] (verbs) at (0,3.35) {RDMA Verbs / Communication Runtime};
    \node[layer] (ibtransport) at (0,2.5) {InfiniBand Transport};
    \node[layer] (udp) at (0,1.65) {UDP / IP（RoCEv2）};
    \node[layer] (eth) at (0,0.8) {Ethernet + QoS / PFC / ECN};
    \node[layer] (phy) at (0,-0.05) {NIC / Switch / Physical Link};

    \draw[arrow] (app) -- (verbs);
    \draw[arrow] (verbs) -- (ibtransport);
    \draw[arrow] (ibtransport) -- (udp);
    \draw[arrow] (udp) -- (eth);
    \draw[arrow] (eth) -- (phy);

    \node[align=center] at (5.4,2.1) {RoCEv2：\\将 RDMA 传输封装\\在 UDP/IP/Ethernet 上};
  \end{tikzpicture}
  \caption{RoCEv2 协议栈抽象：RDMA transport 通过 UDP/IP 承载在以太网 fabric 上。}
  \label{fig:roce-v2-stack}
\end{figure}

\subsection{Lossless Ethernet}
\label{subsec:lossless-ethernet}

RoCE 的关键工程目标是构造近似无损的以太网环境。所谓 lossless 并不是说永远不会出错，而是通过链路层流控、拥塞通知、队列管理和 QoS 配置，尽量避免因拥塞丢包导致 RDMA 重传或连接异常。

在普通 TCP 网络中，丢包常被用作拥塞信号，TCP 会重传并降低发送速率。但 RDMA transport 对丢包更敏感，重传和恢复成本高。对于大规模 all-reduce 或 all-to-all，少量链路丢包可能造成整个训练 step 长尾变长，因为 collective 通常需要等待最慢 rank。

Lossless Ethernet 的核心组件包括：

\begin{itemize}
  \item \textbf{PFC}：Priority Flow Control，在特定优先级上暂停上游发送，避免队列溢出；
  \item \textbf{ECN}：Explicit Congestion Notification，在不丢包的情况下标记拥塞；
  \item \textbf{DCQCN 或类似拥塞控制}：发送端根据拥塞信号调节速率；
  \item \textbf{QoS / DSCP / PCP 映射}：确保 RoCE 流量进入正确优先级队列；
  \item \textbf{交换机 buffer 配置}：避免瞬时突发导致丢包；
  \item \textbf{ECMP 与路由设计}：在多路径网络中平衡流量。
\end{itemize}

\subsection{PFC 与 ECN}
\label{subsec:pfc-ecn}

PFC 和 ECN 经常一起出现，但它们解决的问题不同。

PFC 是链路层暂停机制。当某个交换机端口上的特定优先级队列接近拥塞阈值时，可以向上游发送 pause frame，要求对方暂停该优先级流量。这样可以避免队列溢出丢包。它的优点是直接有效，缺点是可能造成 head-of-line blocking、拥塞扩散，甚至在配置不当时出现 PFC storm。

ECN 是网络层拥塞标记机制。当交换机检测到队列拥塞时，不直接丢包，而是在 IP header 中标记 ECN。接收端或协议栈把拥塞信号反馈给发送端，发送端降低速率。相比 PFC，ECN 更适合做端到端拥塞控制，但它要求网络路径上的设备都正确支持和配置。

\begin{table}[H]
  \centering
  \caption{PFC 与 ECN 的工程对比}
  \label{tab:pfc-ecn-comparison}
  \begin{tabular}{p{2.8cm}p{5.0cm}p{5.0cm}}
    \toprule
    机制 & 优势 & 风险 \\
    \midrule
    PFC & 可以快速阻止特定优先级队列丢包，适合构造无损队列 & 可能导致 head-of-line blocking、拥塞扩散、pause storm \\
    ECN & 不丢包即可反馈拥塞，支持端到端速率调节 & 需要端到端设备和主机协议栈正确配置，阈值不当会影响吞吐或延迟 \\
    QoS 映射 & 将 RoCE 流量放入专用优先级，避免与普通流量互相干扰 & 映射错误会让 RoCE 流量进入非无损队列或错误拥塞控制域 \\
    Buffer 调优 & 吸收突发流量，降低瞬时丢包概率 & 过大 buffer 可能增加排队延迟，过小 buffer 容易丢包或触发 PFC \\
    \bottomrule
  \end{tabular}
\end{table}

AI 集群中的 PFC/ECN 配置不能照搬普通业务网络。训练流量具有明显 collective 特征：all-reduce、reduce-scatter、all-gather、all-to-all 会产生同步、突发和多对多模式。错误的阈值可能在小规模测试中看不出问题，但在大规模训练中导致长尾延迟、NCCL timeout 或吞吐抖动。

\subsection{拥塞控制}
\label{subsec:congestion-control}

拥塞控制的目标是让网络在高负载下仍然稳定运行。对 RoCE 来说，拥塞控制尤其重要，因为它既要避免丢包，又要避免过度 PFC 导致拥塞扩散。

AI 集群中的拥塞来源包括：

\begin{itemize}
  \item 多个训练作业共享同一 leaf-spine 网络；
  \item MoE all-to-all 造成多对多突发；
  \item checkpoint 或数据加载与训练通信混用网络；
  \item 错误 rank mapping 导致某些链路热点；
  \item ECMP hash 不均匀，多个大流落到同一路径；
  \item oversubscription 过高；
  \item 某个节点或 NIC 异常导致重传和尾延迟。
\end{itemize}

拥塞控制不是只调一个参数。它涉及主机、NIC、交换机、QoS、队列、路由、NCCL 算法和作业调度。工程上应建立一套闭环：压测、监控、告警、定位、回归。不要只在训练作业报错后才查看交换机丢包计数；更好的做法是在集群上线前就用 NCCL tests、ib\_write\_bw、perftest、RoCE counters、交换机 telemetry 做基线测试。

\section{GPUDirect RDMA}
\label{sec:gpudirect-rdma}

RDMA 让 NIC 可以直接访问主机内存，而 GPUDirect RDMA 进一步让 NIC 可以直接访问 GPU memory。对于多节点 GPU 训练，这是非常关键的能力。没有 GPUDirect RDMA 时，跨节点通信可能需要经过 CPU host memory staging：

\[
  \text{GPU HBM} \rightarrow \text{CPU Memory} \rightarrow \text{NIC} \rightarrow \text{Network}.
\]

有 GPUDirect RDMA 时，理想路径可以变成：

\[
  \text{GPU HBM} \rightarrow \text{NIC} \rightarrow \text{Network}.
\]

接收方向同理，远端数据可以由 NIC 直接写入 GPU memory。这样可以减少 CPU 内存拷贝、降低延迟、释放 CPU 资源，并提高跨节点 GPU collective 的效率。

\subsection{GPU Memory 直接被 NIC 访问}
\label{subsec:nic-access-gpu-memory}

GPUDirect RDMA 需要 GPU、NIC、驱动、内核模块、PCIe 拓扑和系统配置共同支持。NIC 要能够通过 PCIe peer-to-peer DMA 访问 GPU memory；GPU 驱动需要暴露相应的 peer memory 支持；通信库需要正确注册 GPU buffer 并使用支持 GPUDirect 的数据路径。

从应用层看，这些细节通常被 NCCL 或 UCX 封装。例如 PyTorch 中一个跨节点 \texttt{dist.all\_reduce} 操作，底层可能通过 NCCL 使用 InfiniBand 或 RoCE，并利用 GPUDirect RDMA 直接传输 CUDA tensor。

\begin{figure}[H]
  \centering
  \begin{tikzpicture}[
    font=\small,
    box/.style={draw, rounded corners, minimum width=2.1cm, minimum height=0.75cm, align=center},
    nodebox/.style={draw, rounded corners, minimum width=5.0cm, minimum height=3.5cm, align=center},
    arrow/.style={-{Stealth[length=2mm]}, thick}
  ]
    \node[nodebox] (nodea) at (-3.3,1.8) {};
    \node[anchor=north west] at ($(nodea.north west)+(0.2,-0.2)$) {Node A};
    \node[box] (gpua) at (-4.2,2.4) {GPU HBM};
    \node[box] (nica) at (-2.4,1.2) {NIC / HCA};
    \node[box] (cpua) at (-4.2,1.2) {CPU Memory};

    \node[nodebox] (nodeb) at (3.3,1.8) {};
    \node[anchor=north west] at ($(nodeb.north west)+(0.2,-0.2)$) {Node B};
    \node[box] (gpub) at (4.2,2.4) {GPU HBM};
    \node[box] (nicb) at (2.4,1.2) {NIC / HCA};
    \node[box] (cpub) at (4.2,1.2) {CPU Memory};

    \node[box, minimum width=2.8cm] (net) at (0,1.2) {RDMA Network};

    \draw[arrow] (gpua) -- node[above left] {\scriptsize P2P DMA} (nica);
    \draw[arrow] (nica) -- (net);
    \draw[arrow] (net) -- (nicb);
    \draw[arrow] (nicb) -- node[above right] {\scriptsize P2P DMA} (gpub);

    \draw[dashed] (gpua) -- (cpua);
    \draw[dashed] (cpub) -- (gpub);

    \node[align=center] at (0,-0.3) {实线：GPUDirect RDMA 数据路径；虚线：传统 CPU staging 路径};
  \end{tikzpicture}
  \caption{GPUDirect RDMA 数据流：NIC 可以直接读写 GPU memory，减少 CPU host memory 中转。}
  \label{fig:gpudirect-rdma-dataflow}
\end{figure}

\subsection{绕过 CPU Copy}
\label{subsec:bypass-cpu-copy}

绕过 CPU copy 的价值不仅是减少一次数据拷贝。它还减少了多个隐性开销：

\begin{itemize}
  \item CPU 内存带宽不再成为跨节点 GPU 通信的中间瓶颈；
  \item CPU cache 不会被大 tensor 通信污染；
  \item CPU core 可以用于 dataloader、通信进度、调度和服务逻辑；
  \item 降低端到端延迟，尤其对小消息和频繁通信有帮助；
  \item 减少内存分配、pin memory 和 staging buffer 管理复杂度。
\end{itemize}

不过 GPUDirect RDMA 也有条件和限制。它依赖设备拓扑。如果 GPU 和 NIC 之间跨越不支持 P2P 的 PCIe switch 或平台配置不当，性能可能下降或功能不可用。多 NUMA、多 NIC、多 GPU 节点中，GPU 与 NIC 的亲和性尤其重要。上一章中 \texttt{nvidia-smi topo -m} 的 GPU-NIC 路径，在跨节点通信中会变得更加关键。

\subsection{NCCL 跨节点通信}
\label{subsec:nccl-cross-node-communication}

NCCL 在多节点环境中会自动选择可用的网络传输路径。常见环境变量包括：

\begin{itemize}
  \item \texttt{NCCL\_DEBUG}：打开调试日志；
  \item \texttt{NCCL\_DEBUG\_SUBSYS}：指定调试子系统；
  \item \texttt{NCCL\_SOCKET\_IFNAME}：指定 socket 网络接口；
  \item \texttt{NCCL\_IB\_HCA}：指定 InfiniBand / RoCE HCA；
  \item \texttt{NCCL\_IB\_GID\_INDEX}：选择 RoCE GID；
  \item \texttt{NCCL\_IB\_ROCE\_VERSION\_NUM}：指定 RoCE version；
  \item \texttt{NCCL\_NET\_GDR\_LEVEL}：影响 GPUDirect RDMA 使用策略。
\end{itemize}

下面是一个用于调试的多节点 NCCL 测试示例：

\begin{lstlisting}[language=bash,caption={多节点 NCCL 通信调试示例},label={lst:nccl-rdma-debug}]
export NCCL_DEBUG=INFO
export NCCL_DEBUG_SUBSYS=INIT,NET,GRAPH,TOPO,COLL
export NCCL_IB_HCA=mlx5_0,mlx5_1
export NCCL_SOCKET_IFNAME=eth0

mpirun -np 16 -H node0:8,node1:8 \
  ./build/all_reduce_perf -b 8M -e 8G -f 2 -g 1
\end{lstlisting}

NCCL 日志中需要重点观察：

\begin{itemize}
  \item 是否检测到 IB/RoCE 设备；
  \item 是否使用了预期 HCA；
  \item 是否启用 GPUDirect RDMA；
  \item channel 数量是否合理；
  \item ring/tree 拓扑是否符合节点结构；
  \item 是否 fallback 到 socket；
  \item 是否存在连接错误、timeout 或 retry。
\end{itemize}

如果 NCCL 意外 fallback 到 TCP socket，训练仍然可能运行，但性能会显著下降。很多“训练突然变慢”的问题，最终都能在 NCCL 初始化日志中发现线索。

\subsection{多 NIC、多 Rail 与拓扑感知}
\label{subsec:multi-nic-multirail}

高端 AI 服务器通常配备多张 NIC 或多端口 HCA。多 rail 的目标是让多个网络端口并行承载通信，提高总带宽并改善拓扑利用率。例如，一个 8 GPU 节点可能有 4 张 400Gb/s NIC，每两张 GPU 共享或靠近一张 NIC。理想情况下，rank 应该优先使用距离自己更近的 NIC，避免跨 NUMA 或跨 PCIe switch 访问远端 NIC。

多 rail 设计涉及三个映射：

\begin{enumerate}
  \item GPU 到 NIC 的物理拓扑映射；
  \item rank 到 GPU 的映射；
  \item NCCL channel 到 NIC rail 的映射。
\end{enumerate}

如果这些映射不一致，可能出现某些 NIC 过载、某些 NIC 空闲，或者 GPU 访问远端 NIC 导致 PCIe/NUMA 瓶颈。

\begin{figure}[H]
  \centering
  \begin{tikzpicture}[
    font=\small,
    gpu/.style={draw, rounded corners, minimum width=0.9cm, minimum height=0.55cm, align=center},
    nic/.style={draw, rounded corners, minimum width=1.2cm, minimum height=0.6cm, align=center, fill=gray!10},
    sw/.style={draw, rounded corners, minimum width=2.0cm, minimum height=0.7cm, align=center},
    fast/.style={line width=1.1pt},
    arrow/.style={-{Stealth[length=2mm]}, thick}
  ]
    \foreach \i/\x in {0/-3.5,1/-2.5,2/-1.5,3/-0.5,4/0.5,5/1.5,6/2.5,7/3.5} {
      \node[gpu] (g\i) at (\x,0.4) {G\i};
    }

    \node[nic] (n0) at (-3.0,2.0) {NIC0};
    \node[nic] (n1) at (-1.0,2.0) {NIC1};
    \node[nic] (n2) at (1.0,2.0) {NIC2};
    \node[nic] (n3) at (3.0,2.0) {NIC3};

    \node[sw] (fabric) at (0,3.6) {Network Fabric};

    \draw[fast] (g0.north) -- (n0.south);
    \draw[fast] (g1.north) -- (n0.south);
    \draw[fast] (g2.north) -- (n1.south);
    \draw[fast] (g3.north) -- (n1.south);
    \draw[fast] (g4.north) -- (n2.south);
    \draw[fast] (g5.north) -- (n2.south);
    \draw[fast] (g6.north) -- (n3.south);
    \draw[fast] (g7.north) -- (n3.south);

    \draw[arrow] (n0.north) -- (fabric.south);
    \draw[arrow] (n1.north) -- (fabric.south);
    \draw[arrow] (n2.north) -- (fabric.south);
    \draw[arrow] (n3.north) -- (fabric.south);

    \node[align=center] at (0,-0.75) {多 rail 节点：rank、GPU 与 NIC 的亲和性会影响跨节点通信性能};
  \end{tikzpicture}
  \caption{多 NIC / 多 rail 拓扑感知：GPU 应尽量使用拓扑距离更近的 NIC。}
  \label{fig:multi-nic-multirail}
\end{figure}

\section{AI 集群网络设计}
\label{sec:ai-cluster-network-design}

理解 RDMA 和 RoCE 之后，还需要从集群整体设计角度看网络。AI 集群网络不只是“每台机器插几张网卡”，还包括拓扑、交换机层级、oversubscription、路由、拥塞控制、监控和调度协同。

\subsection{Fat-Tree}
\label{subsec:fat-tree}

Fat-tree 是数据中心和 HPC 中常见的拓扑思想。它通过多层交换机构造高带宽 bisection bandwidth，让任意两个节点之间都有较好的通信路径。典型 leaf-spine 架构可以看作 fat-tree 的一种数据中心实现：服务器连接到 leaf switch，leaf switch 上联到 spine switch，spine 之间提供横向互联能力。

对于 AI 训练，fat-tree 的优点是拓扑规则、扩展性好、路径多，适合集体通信。缺点是成本高，尤其当目标是低 oversubscription 或 non-blocking fabric 时，需要大量高速交换机端口和光模块。

\begin{figure}[H]
  \centering
  \begin{tikzpicture}[
    font=\small,
    srv/.style={draw, rounded corners, minimum width=1.1cm, minimum height=0.55cm, align=center},
    sw/.style={draw, rounded corners, minimum width=1.5cm, minimum height=0.65cm, align=center, fill=gray!10},
    link/.style={line width=0.8pt}
  ]
    \node[sw] (s0) at (-3,4.0) {Spine0};
    \node[sw] (s1) at (-1,4.0) {Spine1};
    \node[sw] (s2) at (1,4.0) {Spine2};
    \node[sw] (s3) at (3,4.0) {Spine3};

    \node[sw] (l0) at (-3,2.2) {Leaf0};
    \node[sw] (l1) at (-1,2.2) {Leaf1};
    \node[sw] (l2) at (1,2.2) {Leaf2};
    \node[sw] (l3) at (3,2.2) {Leaf3};

    \foreach \l in {l0,l1,l2,l3} {
      \foreach \s in {s0,s1,s2,s3} {
        \draw[link] (\l.north) -- (\s.south);
      }
    }

    \foreach \i/\x in {0/-3.6,1/-2.4,2/-1.6,3/-0.4,4/0.4,5/1.6,6/2.4,7/3.6} {
      \node[srv] (n\i) at (\x,0.6) {Node\i};
    }

    \draw[link] (n0.north) -- (l0.south);
    \draw[link] (n1.north) -- (l0.south);
    \draw[link] (n2.north) -- (l1.south);
    \draw[link] (n3.north) -- (l1.south);
    \draw[link] (n4.north) -- (l2.south);
    \draw[link] (n5.north) -- (l2.south);
    \draw[link] (n6.north) -- (l3.south);
    \draw[link] (n7.north) -- (l3.south);

    \node[align=center] at (0,-0.35) {Leaf-Spine / Fat-tree：提供多路径、高横向带宽的 AI 集群网络};
  \end{tikzpicture}
  \caption{AI 集群 leaf-spine 网络拓扑示意：节点通过 leaf 接入，多 leaf 通过 spine 互联。}
  \label{fig:ai-cluster-leaf-spine}
\end{figure}

\subsection{Dragonfly}
\label{subsec:dragonfly}

Dragonfly 是另一类高性能网络拓扑，常见于大型 HPC 系统。它通过 group 内高带宽连接和 group 间全局连接，在较低直径下连接大量节点。相比 fat-tree，dragonfly 可以在某些规模下减少交换机和线缆成本，但路由、拥塞控制和流量工程更加复杂。

对于 AI 集群，dragonfly 的适用性取决于通信模式。如果 workload 以规则 collective 为主，并且通信库和路由能够充分利用 topology-aware 算法，dragonfly 可以提供良好性能。但如果多个作业混部、all-to-all 流量复杂、拥塞控制不足，热点路径可能导致长尾问题。

本书不深入展开 dragonfly 的路由算法，只强调一点：AI 集群网络拓扑必须与通信模式共同设计。大模型训练不是普通均匀流量。DP、TP、PP、EP、checkpoint、数据加载、推理请求各自产生不同流量形态，网络拓扑和调度系统必须理解这些差异。

\subsection{Oversubscription}
\label{subsec:oversubscription}

Oversubscription 指下行接入带宽大于上行汇聚带宽。例如一个 leaf switch 下接入 16 台服务器，每台 400Gb/s，总下行 6.4Tb/s，但上行到 spine 总带宽只有 3.2Tb/s，则 oversubscription ratio 为 2:1。普通数据中心为了降低成本，经常接受一定 oversubscription。但 AI 训练网络对 oversubscription 非常敏感。

如果多个节点之间频繁进行 all-reduce 或 all-to-all，跨 leaf 流量会占用上行带宽。oversubscription 会导致：

\begin{itemize}
  \item collective 带宽下降；
  \item 某些训练 step 长尾变长；
  \item 多作业混部互相影响；
  \item RoCE PFC/ECN 更容易触发；
  \item NCCL timeout 风险增加。
\end{itemize}

因此，大规模训练集群通常倾向于低 oversubscription 或 non-blocking fabric，尤其是训练网络。推理集群则可以根据 workload 做分层设计：prefill-heavy、decode-heavy、embedding-heavy、存储访问等流量可能使用不同网络或 QoS 策略。

\subsection{East-West Traffic}
\label{subsec:east-west-traffic}

AI 集群中的主要流量通常是 east-west traffic，即节点之间的横向流量，而不是 north-south traffic。传统 Web 服务中，请求可能从用户进入负载均衡，再访问服务、缓存、数据库，南北向流量明显。而分布式训练中，大多数通信发生在训练节点之间：rank 与 rank 交换梯度、参数分片、token、activation。

East-west traffic 的特点是：

\begin{itemize}
  \item 流量同步性强，多个 rank 同时发起 collective；
  \item 数据量大，容易形成突发；
  \item 通信模式规则，但受并行策略影响；
  \item all-to-all 可能造成复杂多对多流量；
  \item 单个慢 rank 会拖慢整个 job。
\end{itemize}

这也是为什么 AI 集群调度应尽量把同一个训练作业放在网络拓扑上相邻、带宽充足的节点集合中，而不是随机散布在整个数据中心。后续第 20 章讨论 Kubernetes 调度时，我们会看到 topology-aware scheduling、gang scheduling 和 queueing system 如何与网络拓扑结合。

\subsection{网络监控指标}
\label{subsec:network-monitoring-metrics}

AI 集群网络必须可观测。只监控 GPU utilization 不够，只监控交换机端口 up/down 也不够。至少需要覆盖以下层次：

\begin{longtable}{p{3.2cm}p{5.0cm}p{5.2cm}}
  \caption{AI 集群网络常见监控指标}
  \label{tab:ai-network-monitoring} \\
  \toprule
  层次 & 指标 & 说明 \\
  \midrule
  \endfirsthead
  \toprule
  层次 & 指标 & 说明 \\
  \midrule
  \endhead
  NIC & TX/RX bandwidth & 网卡发送和接收带宽，判断 rail 是否均衡 \\
  NIC & packet error / retry & 发现链路错误、丢包、重传和 RoCE 异常 \\
  NIC & RDMA counters & 包括 RDMA read/write/send、completion error 等 \\
  PCIe & GPU-NIC 拓扑和吞吐 & 判断是否跨 NUMA、跨 switch 或 P2P 不可用 \\
  Switch & port utilization & 发现热点链路和 oversubscription \\
  Switch & buffer occupancy & 判断是否存在队列积压和拥塞 \\
  Switch & PFC pause frame & 发现 PFC 触发频率、pause storm 或拥塞扩散 \\
  Switch & ECN marking & 判断拥塞标记是否合理，是否过早或过晚 \\
  NCCL & algbw / busbw & 验证 collective 有效带宽 \\
  NCCL & timeout / retry / fallback & 发现网络异常、fallback 到 socket 或连接问题 \\
  Job & step time / comm time & 判断通信是否在训练关键路径 \\
  Job & straggler rank & 找出拖慢整个作业的慢 rank \\
  \bottomrule
\end{longtable}

网络监控的关键是建立 baseline。同一类节点、同一类网卡、同一类拓扑，应有稳定的 NCCL tests 结果和端口带宽曲线。一旦某台机器或某条链路偏离 baseline，就应自动标记为可疑节点，避免调度大作业。

\section{实践：RDMA / RoCE 排查流程}
\label{sec:rdma-roce-debug-workflow}

本节给出一个常见排查流程。假设一个多节点训练作业扩展性异常，单机 8 卡正常，两机 16 卡明显变慢，或者 NCCL 偶发 timeout。

\subsection{第一步：确认 NCCL 是否走预期网络}
\label{subsec:debug-nccl-network}

首先打开 NCCL 日志：

\begin{lstlisting}[language=bash,caption={打开 NCCL 网络调试日志},label={lst:nccl-net-debug}]
export NCCL_DEBUG=INFO
export NCCL_DEBUG_SUBSYS=INIT,NET,GRAPH,TOPO,COLL
\end{lstlisting}

观察是否出现以下问题：

\begin{itemize}
  \item 没有识别到 IB/RoCE 设备；
  \item fallback 到 Socket；
  \item 使用了错误网卡；
  \item GID index 不正确；
  \item GPUDirect RDMA 未启用；
  \item 多 rail 未被利用；
  \item 某些 rank 连接失败或 timeout。
\end{itemize}

\subsection{第二步：运行基础带宽测试}
\label{subsec:debug-bandwidth-test}

在训练框架之外运行 NCCL tests：

\begin{lstlisting}[language=bash,caption={多节点 NCCL all-reduce 基线测试},label={lst:multi-node-nccl-baseline}]
mpirun -np 16 -H node0:8,node1:8 \
  -x NCCL_DEBUG=INFO \
  -x NCCL_IB_HCA=mlx5_0,mlx5_1 \
  ./build/all_reduce_perf -b 8M -e 8G -f 2 -g 1
\end{lstlisting}

如果 NCCL tests 已经异常，说明问题在通信基础设施或配置，不要先调模型。应继续检查 HCA、链路、交换机和系统配置。

\subsection{第三步：检查 HCA 与链路状态}
\label{subsec:debug-hca-link}

常见命令包括：

\begin{lstlisting}[language=bash,caption={检查 RDMA 设备和链路状态的常用命令},label={lst:rdma-link-check}]
# 查看 RDMA 设备
ibv_devinfo

# 查看网卡和端口
ibstat

# 查看 RDMA link
rdma link show

# 查看网卡统计
ethtool -S <interface>

# 查看 IP 与接口
ip addr show
\end{lstlisting}

检查内容包括：

\begin{itemize}
  \item link 是否 up；
  \item 速率是否符合预期；
  \item MTU 是否一致；
  \item 是否存在大量 error、drop、retry；
  \item RoCE 接口 IP、GID、VLAN、priority 是否正确；
  \item 多 NIC 是否全部可用。
\end{itemize}

\subsection{第四步：检查 RoCE QoS、PFC 与 ECN}
\label{subsec:debug-roce-qos}

RoCE 环境中，需要确认主机和交换机配置一致：

\begin{itemize}
  \item RoCE 流量是否映射到正确 priority；
  \item PFC 是否只对目标 priority 启用；
  \item ECN 阈值是否合理；
  \item DSCP/PCP 是否在交换机路径中被保留；
  \item 是否存在 pause frame 暴增；
  \item 是否存在 ECN marking 异常；
  \item 是否有普通业务流量混入 RoCE 无损队列。
\end{itemize}

如果 PFC pause frame 持续增加，同时训练 step 出现长尾，说明可能存在拥塞扩散或队列阻塞。此时需要结合交换机 telemetry 和作业调度信息，判断是否有多个大作业共享同一瓶颈链路。

\subsection{第五步：检查 GPU-NIC 拓扑}
\label{subsec:debug-gpu-nic-topology}

运行：

\begin{lstlisting}[language=bash,caption={检查 GPU 与 NIC 拓扑},label={lst:gpu-nic-topology}]
nvidia-smi topo -m
\end{lstlisting}

重点观察：

\begin{itemize}
  \item 每张 GPU 到每张 NIC 的路径；
  \item 是否跨 NUMA；
  \item 是否跨 PCIe switch；
  \item 多 rail 是否与 GPU 分组匹配；
  \item local rank 是否绑定到了合适 GPU。
\end{itemize}

如果某些 rank 使用远端 NIC，跨节点带宽可能明显下降。可以通过调整 \texttt{CUDA\_VISIBLE\_DEVICES}、launcher rank mapping、NCCL HCA 选择或调度器节点分配策略改善。

\subsection{第六步：结合训练 Timeline}
\label{subsec:debug-training-timeline}

最后，用 Nsight Systems 或框架 profiler 观察真实训练 step。需要判断：

\begin{itemize}
  \item NCCL kernel 是否在关键路径；
  \item backward 与 all-reduce 是否重叠；
  \item 是否有某个 rank 明显慢；
  \item 通信前后是否存在 CPU 空洞；
  \item checkpoint 或 dataloader 是否与训练通信争用网络；
  \item MoE all-to-all 是否造成突发拥塞。
\end{itemize}

只有把网络指标与训练 timeline 对齐，才能避免误判。例如 NCCL tests 正常但训练慢，可能是框架 bucket 太小、overlap 不好、rank mapping 不合理或其他 I/O 干扰；训练 timeline 中 NCCL 长尾明显，同时交换机 PFC/ECN 异常，则更可能是网络拥塞。

\section{设计原则与经验法则}
\label{sec:rdma-roce-design-principles}

\subsection{原则一：训练网络与普通业务网络分离}
\label{subsec:principle-separate-network}

大规模训练流量具有突发、大带宽、同步等待和长尾敏感特征。把训练 RDMA/RoCE 流量与普通业务流量混在同一无隔离网络中，容易互相影响。实践中应尽量使用独立训练网络，或至少通过 QoS、VLAN、VRF、优先级队列和带宽隔离明确区分。

\subsection{原则二：先做无损，再做拥塞控制}
\label{subsec:principle-lossless-congestion}

RoCE 的基础是避免丢包，但只靠 PFC 并不够。PFC 可以阻止队列溢出，却可能造成拥塞扩散。ECN 和端到端拥塞控制用于在不丢包的情况下调节发送速率。一个稳定 RoCE 网络需要 PFC、ECN、QoS、buffer 和路由共同配置。

\subsection{原则三：拓扑感知比盲目加带宽更重要}
\label{subsec:principle-topology-aware-network}

多 NIC、多 rail、多 leaf 的网络中，平均带宽不代表实际作业带宽。一个错误的 rank mapping 或 NIC 绑定可能让通信集中在少数链路。训练调度器应该理解节点、GPU、NIC、leaf、spine 的关系，把同一作业放在拓扑上更紧凑、更均衡的位置。

\subsection{原则四：用 NCCL Tests 建立集群基线}
\label{subsec:principle-nccl-baseline}

每个 AI 集群上线前都应建立 NCCL tests 基线，包括单机、双机、跨 leaf、多机规模下的 all-reduce、all-gather、reduce-scatter、all-to-all。后续每次驱动升级、交换机配置变更、NCCL 版本升级、系统镜像更新后，都应重新验证基线。

\subsection{原则五：监控要覆盖端到端路径}
\label{subsec:principle-end-to-end-monitoring}

只监控 GPU 不够，只监控交换机也不够。一次跨节点 GPU 通信跨越 GPU HBM、PCIe、NIC、交换机、对端 NIC、对端 PCIe、对端 GPU HBM。任何一段异常都可能拖慢训练。监控和告警应覆盖端到端路径，并能关联到具体 job、rank、node、NIC 和 switch port。

\section{本章小结}
\label{sec:rdma-roce-summary}

本章从多机 AI 集群的角度介绍了 RDMA、RoCE、InfiniBand、GPUDirect RDMA 与集群网络设计。

\begin{itemize}
  \item 多机训练依赖高性能网络，因为数据并行、张量并行、流水线并行、ZeRO 和 MoE 都会产生跨节点通信。
  \item 跨节点 all-reduce、reduce-scatter、all-gather 和 all-to-all 的性能直接影响训练 step time 和扩展效率。
  \item RDMA 通过 kernel bypass、zero-copy、queue pair、completion queue 和 memory registration 降低延迟与 CPU 开销。
  \item InfiniBand 是面向高性能 RDMA 的专用网络，RoCE 则把 RDMA 承载在以太网上，便于复用以太网生态。
  \item RoCE 网络需要重点关注 lossless Ethernet、PFC、ECN、QoS、拥塞控制和交换机 buffer，否则少量拥塞就可能导致训练长尾或 NCCL timeout。
  \item GPUDirect RDMA 允许 NIC 直接访问 GPU memory，减少 CPU host memory staging，是跨节点 GPU collective 的关键能力。
  \item 多 NIC、多 rail 节点中，GPU-NIC 亲和性、rank mapping 和 NCCL HCA 选择会显著影响实际带宽。
  \item AI 集群网络设计需要考虑 fat-tree、dragonfly、oversubscription、east-west traffic 和网络可观测性。
  \item 排查网络问题应从 NCCL 日志、NCCL tests、HCA 状态、RoCE QoS、GPU-NIC 拓扑和训练 timeline 多层结合入手。
\end{itemize}

至此，我们已经从单个 GPU、单机多卡互联，扩展到了多机 RDMA/RoCE 网络。下一章将进入集群资源调度层：当一个 AI 集群同时运行训练、推理、评测、数据处理和多租户任务时，Kubernetes 如何管理 GPU？Device Plugin、MIG、gang scheduling、Volcano、Kueue、拓扑感知调度和 AI 集群可观测性又如何协同工作？这些内容将进入第 20 章：Kubernetes 在 AI 集群中的调度。

\chapter{Kubernetes 在 AI 集群中的调度}
\label{chap:k8s-ai-scheduling}

前几章我们已经沿着 AI Infrastructure 的物理路径逐层展开：第 17 章讨论单个 GPU 的计算与内存层级，第 18 章讨论单机多卡的 NVLink 与 NVSwitch，第 19 章讨论跨节点的 RDMA、RoCE 与 AI 集群网络。到这里，一个大模型系统的硬件底座已经逐渐清晰：GPU 提供计算，HBM 提供显存，NVLink/NVSwitch 提供节点内高速互联，RDMA/RoCE 提供节点间高速网络。

但是，硬件本身不会自动变成可用的集群能力。一个真实 AI 集群通常同时运行多种任务：

\begin{itemize}
  \item 大规模预训练作业，可能需要几百到几千张 GPU 连续运行数周；
  \item SFT、RLHF、DPO、评测、数据清洗等中小规模离线任务；
  \item 在线 LLM 推理服务，需要低延迟、高可用和弹性伸缩；
  \item embedding、rerank、多模态模型、语音模型等不同形态的服务；
  \item 多个团队、多个项目、多个优先级、多个预算池共享同一批 GPU；
  \item 硬件故障、节点维护、驱动升级、镜像升级、网络异常和 checkpoint 恢复。
\end{itemize}

这时，一个核心问题出现了：谁来决定哪个任务运行在哪里？谁来保证一个 256 GPU 的训练作业不会只启动了 200 个 worker 就陷入死锁？谁来保证在线推理服务不会被离线训练任务挤占 GPU？谁来记录 GPU 利用率、HBM 占用、网络吞吐、队列等待时间和每 token 成本？谁来在节点故障后重启任务并恢复 checkpoint？

这就是 AI 集群调度系统的任务。

Kubernetes 已经成为云原生基础设施的事实标准之一。它擅长管理容器、Pod、服务发现、滚动升级、健康检查、资源声明和控制器模式。因此，很多 AI 平台选择以 Kubernetes 作为统一资源管理底座，再结合 NVIDIA Device Plugin、GPU Operator、Volcano、Kueue、Ray Operator、Kubeflow Training Operator、MPI Operator、模型服务框架和自研调度器，构建面向训练与推理的一体化 AI 集群平台。

本章不试图把 Kubernetes 所有概念完整讲一遍，而是聚焦 AI Infra 工程中最关键的调度问题：

\begin{itemize}
  \item 为什么 AI 集群不能只用普通 Kubernetes 调度策略？
  \item GPU 如何在 Kubernetes 中被发现、上报、分配和隔离？
  \item MIG、GPU sharing、拓扑感知调度分别解决什么问题？
  \item 分布式训练为什么需要 gang scheduling？
  \item Volcano 和 Kueue 在 AI batch 调度中分别扮演什么角色？
  \item LLM 推理服务如何在 Kubernetes 上部署、扩缩容和混部？
  \item AI 集群应该监控哪些指标，如何从利用率走向成本和 SLA？
\end{itemize}

\section{为什么 AI 集群需要调度系统}
\label{sec:why-ai-cluster-scheduler}

调度系统的本质是把有限资源分配给无限需求。普通在线服务调度关注 CPU、内存、端口、亲和性、反亲和性和服务可用性；AI 集群调度则更加复杂，因为 GPU 是昂贵、异构、拓扑敏感且经常以“成组”方式被使用的资源。

\subsection{GPU 资源昂贵}
\label{subsec:gpu-expensive}

GPU 是 AI 集群中最昂贵的资源之一。一台 8 GPU 服务器的采购成本、机架空间、供电、散热、网络端口和运维成本都很高。对于大模型公司或 AI 平台团队来说，GPU 空闲就是直接的资本浪费；GPU 被低效任务占用，则意味着更隐蔽的机会成本。

因此，AI 集群调度系统不能只回答“有没有空闲 GPU”，还要回答：

\begin{itemize}
  \item 空闲 GPU 是否属于同一节点、同一 NVSwitch 域或同一网络拓扑区域？
  \item 这些 GPU 的型号、显存容量、驱动版本、MIG 配置是否符合任务要求？
  \item 这个任务是在线推理还是离线训练？是否可以被抢占？
  \item 这个任务需要连续运行多久？是否支持 checkpoint 恢复？
  \item 把它放在这里会不会影响更高优先级的任务？
  \item 当前集群是否应该等待更多 GPU 空出来再启动大作业，而不是立即启动小作业造成碎片？
\end{itemize}

对于普通 Kubernetes scheduler 来说，一个 Pod 请求 1 个 GPU，调度器只要找一个有可分配 GPU 的节点即可。但对于 AI 训练，一个作业可能请求 8、16、64、256 张 GPU，并且这些 GPU 之间的拓扑关系会显著影响性能。如果只做局部贪心分配，很容易产生 GPU 碎片：表面上集群还有很多空闲 GPU，但没有足够连续、同构、拓扑良好的 GPU group 来启动大作业。

\subsection{多租户}
\label{subsec:multi-tenancy}

AI 集群通常被多个团队共享。例如基础模型团队、算法研究团队、推理平台团队、数据团队、评测团队和业务团队都需要 GPU。不同团队之间存在不同优先级、预算、配额和 SLA。

多租户调度需要解决几个问题：

\begin{itemize}
  \item \textbf{配额}：每个团队最多能使用多少 GPU、多少 GPU 小时、多少 HBM 或多少节点？
  \item \textbf{公平性}：当资源不足时，如何避免某个团队长期霸占集群？
  \item \textbf{优先级}：紧急线上修复、重要训练作业和普通实验如何排序？
  \item \textbf{抢占}：低优先级任务是否可以被高优先级任务抢占？
  \item \textbf{隔离}：不同租户之间的网络、存储、镜像、权限和日志如何隔离？
  \item \textbf{审计}：谁在什么时间使用了多少 GPU，产出了什么成本？
\end{itemize}

多租户不是简单地给 namespace 设置 resource quota。AI 任务经常是 batch job，有排队、等待、成组启动、运行数小时或数天、checkpoint 恢复等特征。调度系统需要理解“作业”而不仅仅是“Pod”。

\begin{figure}[H]
  \centering
  \begin{tikzpicture}[
    font=\small,
    box/.style={draw, rounded corners, minimum width=2.4cm, minimum height=0.75cm, align=center},
    pool/.style={draw, rounded corners, minimum width=9.8cm, minimum height=2.2cm, align=center},
    arrow/.style={-{Stealth[length=2mm]}, thick}
  ]
    \node[box] (teamA) at (-4,4.0) {基础模型团队};
    \node[box] (teamB) at (-1.3,4.0) {推理平台团队};
    \node[box] (teamC) at (1.4,4.0) {评测团队};
    \node[box] (teamD) at (4.1,4.0) {数据团队};

    \node[box, minimum width=8.8cm] (queue) at (0,2.7) {队列 / 配额 / 优先级 / 抢占策略};

    \node[pool] (cluster) at (0,1.0) {};
    \node at (0,1.75) {AI GPU 资源池};
    \foreach \i/\x in {0/-4.0,1/-3.0,2/-2.0,3/-1.0,4/0.0,5/1.0,6/2.0,7/3.0,8/4.0} {
      \node[draw, rounded corners, minimum width=0.65cm, minimum height=0.45cm] at (\x,0.65) {G};
    }

    \draw[arrow] (teamA.south) -- (queue.north);
    \draw[arrow] (teamB.south) -- (queue.north);
    \draw[arrow] (teamC.south) -- (queue.north);
    \draw[arrow] (teamD.south) -- (queue.north);
    \draw[arrow] (queue.south) -- (cluster.north);

    \node[align=center] at (0,-0.55) {多租户 AI 集群调度：不同团队通过队列、配额和优先级共享 GPU 资源池。};
  \end{tikzpicture}
  \caption{多租户 AI 集群资源池：调度系统需要同时考虑配额、公平性、优先级和资源碎片。}
  \label{fig:multi-tenant-ai-cluster}
\end{figure}

\subsection{训练、推理、评测混部}
\label{subsec:training-serving-eval-colocation}

AI 集群中的任务类型差异很大。预训练任务通常是长时间、大规模、通信密集型；推理服务是长期运行、延迟敏感、需要弹性伸缩；评测任务可能是短时突发、吞吐型；数据处理可能主要消耗 CPU、内存、存储和网络；微调任务规模较小但数量很多。

不同任务对调度系统的要求不同：

\begin{table}[H]
  \centering
  \caption{AI 集群中不同任务类型的调度特征}
  \label{tab:ai-workload-scheduling-characteristics}
  \begin{tabular}{p{2.8cm}p{3.5cm}p{4.2cm}p{4.0cm}}
    \toprule
    任务类型 & 资源特征 & 调度重点 & 典型风险 \\
    \midrule
    预训练 & 大规模 GPU、长时间、强通信 & gang scheduling、拓扑感知、checkpoint 恢复 & 启动碎片、通信慢、故障重启成本高 \\
    SFT / 微调 & 中小规模 GPU、任务数量多 & 队列、公平性、镜像和数据准备 & 小任务挤占大作业资源、低利用率 \\
    RLHF / DPO & 多阶段、多角色、可能混合训练和推理 & workflow 编排、actor/rollout/reference 模型资源协调 & 阶段不均衡、资源等待 \\
    在线推理 & 长期运行、延迟敏感、显存常驻 & 高可用、扩缩容、bin packing、SLA & 冷启动、尾延迟、显存碎片 \\
    离线评测 & 短时突发、吞吐型 & 快速启动、低优先级抢占、批量排队 & 影响线上服务或重要训练 \\
    数据处理 & CPU、内存、存储、网络为主 & 与 GPU 任务隔离、I/O 调度 & 抢占网络和存储带宽 \\
    \bottomrule
  \end{tabular}
\end{table}

混部可以提高资源利用率，但也会增加干扰。例如在线推理服务可能 SM utilization 不高，但 HBM 占用高且对延迟敏感；训练作业可能吞吐高但通信突发；评测任务可能频繁启动和释放 GPU，造成资源碎片。调度系统需要根据任务特征做隔离和分层，而不是把所有 GPU workload 当成同一种 Pod。

\subsection{利用率与公平性}
\label{subsec:utilization-and-fairness}

AI 集群调度的两个目标经常冲突：提高利用率和保证公平性。为了提高利用率，调度器希望尽快把空闲 GPU 分配出去；为了公平性和大作业启动，调度器又可能需要让一些小任务等待，以便积累足够连续资源给大任务。

例如，集群当前有 4 个节点，每个节点剩余 2 张 GPU，总共 8 张空闲 GPU。一个 8 GPU 单节点 TP 推理实例无法启动，因为它需要同一节点的 8 张 GPU；一个 8 GPU 训练作业如果要求 NVSwitch 单节点，也无法启动。此时从总量看资源足够，从拓扑和连续性看资源不足。这就是 GPU 碎片。

调度系统需要在以下策略之间权衡：

\begin{itemize}
  \item \textbf{bin packing}：尽量把任务塞满少数节点，减少碎片；
  \item \textbf{spread}：把任务分散到多个节点，提高可用性或降低局部干扰；
  \item \textbf{reservation}：为大作业预留资源，避免小作业打散集群；
  \item \textbf{backfill}：在不影响高优先级大作业预计启动时间的前提下，运行短小任务；
  \item \textbf{preemption}：通过抢占低优先级任务释放资源；
  \item \textbf{quota borrowing}：允许队列临时借用其他队列未使用的配额。
\end{itemize}

优秀的 AI 调度系统不是简单地追求某一时刻 GPU utilization 最高，而是要优化更长期的目标：整体吞吐、队列等待时间、SLA、训练完成时间、公平性和成本。

\section{Kubernetes 基础回顾}
\label{sec:kubernetes-basics-review}

本节快速回顾 Kubernetes 中与 AI 调度关系最密切的几个概念：Pod、Node、Scheduler、Device Plugin 和 Operator。熟悉 Kubernetes 的读者可以快速浏览。

\subsection{Pod}
\label{subsec:pod}

Pod 是 Kubernetes 中最小的调度单元。一个 Pod 可以包含一个或多个容器，这些容器共享网络 namespace、部分存储卷和生命周期。在 AI 任务中，一个 Pod 可能代表：

\begin{itemize}
  \item 一个训练 worker；
  \item 一个 parameter server；
  \item 一个 PyTorch distributed rank；
  \item 一个 Ray worker；
  \item 一个 vLLM / TensorRT-LLM / Triton model server 实例；
  \item 一个数据预处理任务；
  \item 一个评测 worker。
\end{itemize}

Pod 通过 \texttt{resources.requests} 和 \texttt{resources.limits} 声明资源需求。GPU 通常以扩展资源的形式出现，例如 \texttt{nvidia.com/gpu: 1}。一个最小 GPU Pod 示例：

\begin{lstlisting}[caption={请求 1 张 GPU 的 Pod 示例},label={lst:k8s-gpu-pod}]
apiVersion: v1
kind: Pod
metadata:
  name: gpu-demo
spec:
  restartPolicy: Never
  containers:
  - name: app
    image: nvcr.io/nvidia/pytorch:xx.xx-py3
    command: ["python", "train.py"]
    resources:
      limits:
        nvidia.com/gpu: 1
\end{lstlisting}

在 Kubernetes 中，GPU 这类设备资源通常要求 request 与 limit 相等，常见写法是只写 limits。实际行为取决于 device plugin 和 Kubernetes 版本，但工程上应避免对 GPU 做类似 CPU 那样的 overcommit，除非明确使用 GPU sharing 或 MIG 等机制。

\subsection{Node}
\label{subsec:node}

Node 是 Kubernetes 集群中的工作节点。对 AI 集群来说，一个 Node 通常是一台 GPU 服务器，包含：

\begin{itemize}
  \item 多张 GPU；
  \item CPU 和内存；
  \item 本地 NVMe；
  \item 网卡和 RDMA 设备；
  \item NVIDIA driver、CUDA runtime、container runtime；
  \item kubelet、device plugin、monitoring agent；
  \item 可能还包括 MIG 配置、拓扑标签和健康检查组件。
\end{itemize}

Node 会向 Kubernetes API Server 上报可分配资源。调度器根据 Pod 的资源请求、node selector、affinity、taint/toleration、topology spread、priority 等信息选择目标节点。

AI 集群中常见的 node label 包括：

\begin{lstlisting}[language=bash,caption={AI 集群常见 Node Label 示例},label={lst:ai-node-labels}]
accelerator=nvidia
gpu.vendor=nvidia
gpu.model=h100
gpu.count=8
gpu.memory=80gb
gpu.fabric=nvswitch
rdma=true
network.rail=4
zone=az-a
pool=training
\end{lstlisting}

这些 label 可以帮助调度器或平台层把任务放到合适资源池。例如，预训练任务只允许调度到 \texttt{pool=training} 且 \texttt{gpu.fabric=nvswitch} 的节点；在线推理任务调度到 \texttt{pool=serving}；评测任务可以使用较低优先级节点池。

\subsection{Scheduler}
\label{subsec:scheduler}

Kubernetes Scheduler 的基本职责是为 Pending Pod 选择一个 Node。它通常经历两个阶段：

\begin{enumerate}
  \item \textbf{Filter}：过滤掉不满足资源、亲和性、污点容忍、端口、卷等条件的节点；
  \item \textbf{Score}：对可行节点打分，选择最优节点。
\end{enumerate}

默认 scheduler 对通用 workload 很有效，但对 AI batch workload 有几个不足：

\begin{itemize}
  \item 它以 Pod 为基本单位，不天然理解一个分布式训练 Job 需要所有 worker 一起启动；
  \item 它不一定理解 GPU 拓扑、NVLink、NIC 亲和性、MIG profile 等细节；
  \item 它不提供复杂队列、配额借用、gang scheduling 和 backfill 能力；
  \item 它对长时间运行的 batch job 与在线 service 的混部策略有限；
  \item 它不直接处理 checkpoint-aware preemption 和作业级恢复。
\end{itemize}

因此，AI 集群常常在默认 scheduler 之外引入 Volcano、Kueue、自研 scheduler extender、scheduler plugin 或更上层的 batch queueing system。

\subsection{Device Plugin}
\label{subsec:device-plugin}

Device Plugin 是 Kubernetes 管理 GPU、NIC、FPGA 等特殊硬件资源的扩展机制。设备厂商或平台团队通过 device plugin 向 kubelet 注册资源类型、上报设备列表，并在 Pod 被分配设备时执行必要的环境变量、设备文件、挂载和 runtime 配置。

以 NVIDIA GPU 为例，NVIDIA device plugin 会向 Kubernetes 暴露类似 \texttt{nvidia.com/gpu} 的资源。Pod 请求该资源后，kubelet 会调用 device plugin 分配具体 GPU，并把对应设备注入容器。

Device Plugin 的重要性在于：Kubernetes 本身并不知道 GPU 内部细节。它需要 device plugin 告诉它节点上有哪些设备、哪些设备健康、如何把设备分配给容器。对于 AI Infra 工程师来说，device plugin 是 GPU 从“物理硬件”变成“Kubernetes 可调度资源”的关键桥梁。

\subsection{Operator}
\label{subsec:operator}

Operator 是 Kubernetes 的一种控制器模式，用自定义资源 CRD 和控制循环管理复杂应用。AI 集群中常见 operator 包括：

\begin{itemize}
  \item NVIDIA GPU Operator：安装和管理 driver、device plugin、DCGM exporter、MIG manager 等 GPU 组件；
  \item Training Operator / Kubeflow Operator：管理 PyTorchJob、TFJob、MPIJob 等分布式训练任务；
  \item Ray Operator：管理 Ray cluster 和 Ray job；
  \item Model serving operator：管理模型部署、版本、扩缩容和流量切换；
  \item Volcano / Kueue 相关控制器：管理队列、作业准入和调度。
\end{itemize}

Operator 的价值是把复杂运维动作自动化。例如 GPU 节点加入集群后，GPU Operator 可以自动部署驱动相关组件和监控 exporter；用户提交一个 PyTorchJob 后，训练 operator 可以创建 master 和 worker Pod，配置 rendezvous，监控状态，并在失败时重试。

\begin{figure}[H]
  \centering
  \begin{tikzpicture}[
    font=\small,
    box/.style={draw, rounded corners, minimum width=2.2cm, minimum height=0.75cm, align=center},
    wide/.style={draw, rounded corners, minimum width=8.8cm, minimum height=0.85cm, align=center},
    nodebox/.style={draw, rounded corners, minimum width=8.8cm, minimum height=2.1cm, align=center},
    arrow/.style={-{Stealth[length=2mm]}, thick}
  ]
    \node[wide] (api) at (0,4.7) {API Server};
    \node[box] (sched) at (-3.2,3.4) {Scheduler};
    \node[box] (ctrl) at (0,3.4) {Controller\\Manager};
    \node[box] (op) at (3.2,3.4) {Operators};

    \node[nodebox] (node) at (0,1.55) {};
    \node[anchor=north west] at ($(node.north west)+(0.2,-0.2)$) {GPU Node};
    \node[box] (kubelet) at (-3.0,1.4) {kubelet};
    \node[box] (dp) at (0,1.4) {Device\\Plugin};
    \node[box] (pod) at (3.0,1.4) {GPU Pods};

    \node[box, minimum width=7.5cm] (runtime) at (0,0.35) {Container Runtime / NVIDIA Runtime / Driver};

    \draw[arrow] (sched) -- (api);
    \draw[arrow] (ctrl) -- (api);
    \draw[arrow] (op) -- (api);
    \draw[arrow] (api) -- (kubelet);
    \draw[arrow] (dp) -- (kubelet);
    \draw[arrow] (kubelet) -- (pod);
    \draw[arrow] (pod) -- (runtime);

    \node[align=center] at (0,-0.65) {Kubernetes 控制面通过 kubelet、device plugin 和 operator 管理 GPU 节点与 AI workload。};
  \end{tikzpicture}
  \caption{Kubernetes 控制面与 GPU 节点组件关系图。}
  \label{fig:k8s-control-plane-gpu-node}
\end{figure}

\section{GPU 资源管理}
\label{sec:gpu-resource-management}

GPU 资源管理是 Kubernetes AI 平台的核心。它不仅包括“让 Pod 能请求 GPU”，还包括设备健康检查、驱动管理、MIG、GPU sharing、拓扑感知、资源隔离、监控和故障处理。

\subsection{NVIDIA Device Plugin}
\label{subsec:nvidia-device-plugin}

NVIDIA Device Plugin 的主要职责是把节点上的 NVIDIA GPU 暴露给 Kubernetes。它通常以 DaemonSet 形式部署在 GPU 节点上，与 kubelet 通信。工作流程可以简化为：

\begin{enumerate}
  \item device plugin 启动后发现节点上的 GPU；
  \item 向 kubelet 注册资源名，例如 \texttt{nvidia.com/gpu}；
  \item 周期性上报设备健康状态；
  \item 当 Pod 请求 GPU 时，kubelet 调用 device plugin 的 Allocate 接口；
  \item device plugin 返回容器需要的设备文件、环境变量和挂载信息；
  \item 容器启动后只能看到被分配的 GPU。
\end{enumerate}

从平台角度看，需要关注以下问题：

\begin{itemize}
  \item device plugin 版本是否与驱动、container runtime、Kubernetes 版本兼容；
  \item GPU 健康状态是否能正确反映到 Kubernetes；
  \item GPU reset、Xid error、ECC error 后是否自动隔离节点；
  \item 容器中 \texttt{CUDA\_VISIBLE\_DEVICES} 是否符合预期；
  \item MIG 模式下资源名和 profile 是否正确上报；
  \item 与 GPU Operator、DCGM exporter、MIG manager 是否协同。
\end{itemize}

\subsection{GPU Request / Limit}
\label{subsec:gpu-request-limit}

用户通常通过 Pod spec 声明 GPU 需求：

\begin{lstlisting}[caption={在 Kubernetes 中请求 GPU 资源},label={lst:k8s-gpu-resource-request}]
resources:
  limits:
    nvidia.com/gpu: 4
\end{lstlisting}

这里的 \texttt{4} 表示该容器需要 4 张 GPU。Kubernetes 会把该 Pod 调度到一个可分配至少 4 张 GPU 的节点上。对于分布式训练，通常每个 Pod 请求 1 张或多张 GPU，然后由训练框架根据 rank、world size 和 local rank 建立通信。

一个 PyTorch 分布式训练 worker 的 Pod 可能类似：

\begin{lstlisting}[caption={分布式训练 Worker Pod 的资源声明示例},label={lst:pytorch-worker-gpu-request}]
apiVersion: v1
kind: Pod
metadata:
  name: train-worker-0
spec:
  restartPolicy: Never
  containers:
  - name: worker
    image: my-registry/pytorch-train:latest
    command: ["bash", "-lc", "torchrun --nproc_per_node=8 train.py"]
    resources:
      limits:
        nvidia.com/gpu: 8
        cpu: "64"
        memory: "512Gi"
\end{lstlisting}

需要注意，Kubernetes 原生资源声明并不表达 GPU 拓扑。例如请求 8 张 GPU，只表示需要同一节点上 8 个 GPU 设备，并不说明这些 GPU 是否通过 NVSwitch 互联，也不说明 GPU 与 NIC 的关系。要表达这些约束，需要结合 node label、node affinity、scheduler plugin、device plugin 拓扑信息或自研调度层。

\subsection{MIG}
\label{subsec:mig}

MIG，即 Multi-Instance GPU，是 NVIDIA 在部分数据中心 GPU 上提供的硬件级 GPU 切分能力。它可以把一张物理 GPU 切分成多个 GPU instance，每个 instance 拥有独立的一部分计算资源、显存、缓存和带宽隔离能力。MIG 的目标不是提高单个大训练任务性能，而是提高小任务、小模型推理和多租户场景下的资源利用率与隔离性。

例如，一张大显存 GPU 如果直接分配给一个只需要少量显存和算力的小模型，会造成浪费。通过 MIG，可以把它切分成多个较小 profile，让多个推理实例或轻量任务共享同一张物理 GPU，同时获得比普通进程共享更强的硬件隔离。

\begin{figure}[H]
  \centering
  \begin{tikzpicture}[
    font=\small,
    gpu/.style={draw, rounded corners, minimum width=7.8cm, minimum height=2.5cm, align=center},
    slice/.style={draw, rounded corners, minimum width=1.2cm, minimum height=1.7cm, align=center, fill=gray!10},
    arrow/.style={-{Stealth[length=2mm]}, thick}
  ]
    \node[gpu] (full) at (0,3.2) {};
    \node at (0,4.15) {Physical GPU};
    \node at (0,3.2) {完整 GPU：SM / HBM / L2 / Copy Engine};

    \node[gpu] (mig) at (0,0.7) {};
    \node at (0,1.65) {MIG Mode};
    \foreach \i/\x/\txt in {
      0/-3.0/{GI 0},
      1/-1.8/{GI 1},
      2/-0.6/{GI 2},
      3/0.6/{GI 3},
      4/1.8/{GI 4},
      5/3.0/{GI 5}
    } {
      \node[slice] at (\x,0.65) {\txt\\SM\\HBM};
    }

    \draw[arrow] (full.south) -- (mig.north);
    \node[align=center] at (0,-1.0) {MIG 将一张物理 GPU 切分为多个具有硬件隔离的 GPU instance。};
  \end{tikzpicture}
  \caption{MIG GPU 切分示意图：适合小模型推理、多租户隔离和提高碎片资源利用率。}
  \label{fig:mig-gpu-partition}
\end{figure}

MIG 在 Kubernetes 中通常由 GPU Operator 和 MIG Manager 管理。平台可以为节点设置不同 MIG strategy，例如单一 profile 或混合 profile。调度时，MIG instance 会以不同扩展资源的形式暴露给 Kubernetes，例如某种 \texttt{nvidia.com/mig-...} 资源名。

MIG 的使用需要权衡：

\begin{itemize}
  \item \textbf{优点}：硬件隔离更强，适合多租户小模型推理，提高大 GPU 的切分利用率；
  \item \textbf{缺点}：切分后无法灵活合并，profile 选择不当会造成新的碎片；
  \item \textbf{限制}：并非所有 GPU 支持，且某些多卡通信、NVLink、profiling 或共享场景会有额外约束；
  \item \textbf{运维复杂度}：需要管理 MIG 模式切换、profile 配置、节点 drain、资源上报和监控。
\end{itemize}

实践中，MIG 更适合稳定的推理资源池，而不是频繁变化的大规模训练资源池。对于训练，通常更希望保留完整 GPU，以获得完整算力、显存和互联能力。

\subsection{GPU Sharing}
\label{subsec:gpu-sharing}

GPU sharing 指多个工作负载共享同一张 GPU。它与 MIG 不同，MIG 是硬件切分，GPU sharing 更多是软件或 runtime 层的共享。常见方式包括：

\begin{itemize}
  \item 多个进程直接共享同一 GPU；
  \item NVIDIA MPS；
  \item time-slicing；
  \item framework-level batching；
  \item serving runtime 内部的多模型复用；
  \item MIG + sharing 的组合。
\end{itemize}

GPU sharing 的价值在于提高小任务利用率。例如某些 notebook、开发调试、小模型推理、embedding 服务或低 QPS 模型不需要独占整张 GPU。共享可以减少浪费。

但 GPU sharing 风险也很大：

\begin{itemize}
  \item 显存 OOM 会相互影响；
  \item SM、HBM 带宽、L2 cache、copy engine 竞争导致性能抖动；
  \item 在线推理尾延迟变差；
  \item 监控和计费复杂；
  \item 故障隔离弱于 MIG；
  \item 用户容易误以为自己独占 GPU。
\end{itemize}

因此，GPU sharing 更适合低优先级、开发测试或可容忍抖动的推理场景。对于强 SLA 在线推理和大规模训练，应谨慎使用。

\subsection{Topology-Aware Scheduling}
\label{subsec:topology-aware-scheduling}

拓扑感知调度是 AI 集群调度的关键能力。它需要理解 GPU 与 GPU、GPU 与 NIC、GPU 与 CPU NUMA、节点与交换机之间的关系。前两章已经解释过，TP、EP、RDMA 和 GPUDirect RDMA 都强依赖拓扑。

在 Kubernetes 中，可以通过多种方式实现拓扑感知：

\begin{itemize}
  \item 节点标签：例如 \texttt{gpu.fabric=nvswitch}、\texttt{rdma=true}；
  \item node affinity：限制 Pod 只能调度到满足拓扑条件的节点；
  \item pod affinity / anti-affinity：控制 Pod 之间靠近或分散；
  \item topology spread constraints：控制副本在 zone、rack、node 间分布；
  \item device plugin 上报拓扑信息；
  \item scheduler extender 或 scheduler plugin；
  \item 上层 batch scheduler 进行作业级节点选择；
  \item 自研资源画像系统记录 GPU/NIC/NVLink/NVSwitch 拓扑。
\end{itemize}

\begin{figure}[H]
  \centering
  \begin{tikzpicture}[
    font=\small,
    gpu/.style={draw, rounded corners, minimum width=0.9cm, minimum height=0.5cm, align=center},
    group/.style={draw, rounded corners, dashed, inner sep=0.22cm},
    task/.style={draw, rounded corners, minimum width=2.0cm, minimum height=0.65cm, align=center, fill=gray!10},
    link/.style={line width=0.8pt},
    arrow/.style={-{Stealth[length=2mm]}, thick}
  ]
    \foreach \i/\x in {0/-3.5,1/-2.5,2/-1.5,3/-0.5,4/0.7,5/1.7,6/2.7,7/3.7} {
      \node[gpu] (g\i) at (\x,0.6) {G\i};
    }

    \node[group, fit=(g0)(g1)(g2)(g3), label=below:{NVSwitch Group A}] {};
    \node[group, fit=(g4)(g5)(g6)(g7), label=below:{NVSwitch Group B}] {};

    \node[task] (job) at (0,3.1) {TP=4 推理实例};

    \draw[arrow] (job.south) -- (-2.0,1.25);
    \node[align=center] at (0,2.25) {拓扑感知调度选择同一高速互联 group 内的 GPU};

    \foreach \a/\b in {0/1,1/2,2/3,4/5,5/6,6/7} {
      \draw[link] (g\a) -- (g\b);
    }
  \end{tikzpicture}
  \caption{拓扑感知 GPU 分配示意：通信密集任务应优先绑定同一高速互联域内的 GPU。}
  \label{fig:topology-aware-gpu-allocation}
\end{figure}

拓扑感知调度尤其适用于：

\begin{itemize}
  \item TP 推理实例；
  \item 单节点 8 GPU 训练；
  \item 多节点训练中的 GPU-NIC 亲和；
  \item MoE expert parallel；
  \item 多 rail RDMA；
  \item 同机多模型混部。
\end{itemize}

\section{分布式训练作业调度}
\label{sec:distributed-training-scheduling}

分布式训练与普通 Pod 调度最大的区别是：它通常需要一组 Pod 共同启动、共同运行、共同失败或共同成功。一个 64 GPU PyTorch 训练作业如果只启动了 50 个 worker，剩下 14 个 worker 处于 Pending，已经启动的 worker 可能会一直等待 rendezvous，浪费 GPU。这个问题就是 gang scheduling 的动机。

\subsection{Gang Scheduling}
\label{subsec:gang-scheduling}

Gang scheduling，也称 coscheduling，要求一个作业的一组任务要么一起调度成功，要么都不运行。对于分布式训练来说，这非常重要。一个 PyTorchJob、MPIJob 或 Ray cluster 往往有最小可运行副本数。如果资源不够，调度器应该让整个 job 等待，而不是只启动一部分 Pod。

\begin{figure}[H]
  \centering
  \begin{tikzpicture}[
    font=\small,
    pod/.style={draw, rounded corners, minimum width=0.85cm, minimum height=0.45cm, align=center},
    node/.style={draw, rounded corners, minimum width=2.0cm, minimum height=0.7cm, align=center},
    arrow/.style={-{Stealth[length=2mm]}, thick}
  ]
    \node at (0,4.0) {Job 需要 8 个 worker，minAvailable = 8};

    \foreach \i/\x in {0/-3.5,1/-2.5,2/-1.5,3/-0.5,4/0.5,5/1.5,6/2.5,7/3.5} {
      \node[pod] (p\i) at (\x,3.2) {W\i};
    }

    \node[node] (n0) at (-3,1.7) {Node A\\4 GPU};
    \node[node] (n1) at (-0.5,1.7) {Node B\\2 GPU};
    \node[node] (n2) at (2.0,1.7) {Node C\\1 GPU};

    \draw[arrow, dashed] (p0.south) -- (n0.north);
    \draw[arrow, dashed] (p1.south) -- (n0.north);
    \draw[arrow, dashed] (p2.south) -- (n0.north);
    \draw[arrow, dashed] (p3.south) -- (n0.north);
    \draw[arrow, dashed] (p4.south) -- (n1.north);
    \draw[arrow, dashed] (p5.south) -- (n1.north);
    \draw[arrow, dashed] (p6.south) -- (n2.north);

    \node[draw, rounded corners, minimum width=8.5cm, minimum height=0.7cm, align=center, fill=gray!10] at (0,0.35) {资源不足 8 个 worker：整个 Job 等待，避免部分启动浪费 GPU};
  \end{tikzpicture}
  \caption{Gang scheduling 等待与启动流程：资源不足时不应只启动部分训练 worker。}
  \label{fig:gang-scheduling}
\end{figure}

Gang scheduling 需要调度器理解 job 级语义，而不是只看单个 Pod。它通常与队列、配额、优先级和抢占结合使用。例如一个高优先级 128 GPU 训练 job 等待时，调度器可能需要抢占一些低优先级任务，释放足够资源后一次性启动整个 job。

\subsection{Volcano}
\label{subsec:volcano}

Volcano 是面向云原生 batch 和高性能计算 workload 的调度系统。它与 Kubernetes 兼容，并提供队列、gang scheduling、优先级、抢占、资源回收、backfill、DRF、公平调度等能力。对 AI 训练来说，Volcano 常用于调度 PyTorch、TensorFlow、MPI、Spark、Ray 等计算密集型任务。

Volcano 的核心价值在于它把调度对象从单个 Pod 提升到 Job 和 Queue。一个 VolcanoJob 可以包含多个 task，每个 task 可以有副本数、资源需求和最小可用数量。调度器根据队列和作业约束决定是否启动。

一个简化 VolcanoJob 示例：

\begin{lstlisting}[caption={VolcanoJob 的概念示例},label={lst:volcano-job-example}]
apiVersion: batch.volcano.sh/v1alpha1
kind: Job
metadata:
  name: pretrain-job
spec:
  minAvailable: 8
  schedulerName: volcano
  queue: foundation-model
  tasks:
  - name: worker
    replicas: 8
    template:
      spec:
        restartPolicy: Never
        containers:
        - name: trainer
          image: my-registry/pretrain:latest
          command: ["bash", "-lc", "torchrun --nnodes=1 --nproc_per_node=8 train.py"]
          resources:
            limits:
              nvidia.com/gpu: 1
\end{lstlisting}

这个示例表达了一个重要语义：作业至少需要 8 个 worker 才能运行。如果只调度了 3 个 worker，它没有意义。Volcano 的 gang plugin 可以帮助避免这种部分启动问题。

Volcano 适合的场景包括：

\begin{itemize}
  \item 大规模分布式训练；
  \item 多队列、多租户 batch 集群；
  \item 需要 gang scheduling 的 MPI/PyTorch/Ray workload；
  \item 需要公平调度、优先级、抢占和 backfill 的场景；
  \item HPC 与 AI 混合计算集群。
\end{itemize}

\subsection{Kueue}
\label{subsec:kueue}

Kueue 是 Kubernetes 原生的作业排队系统，面向 batch、HPC、AI/ML 等 workload。它的核心思想不是替代所有调度逻辑，而是在 Pod 真正创建或运行之前，对作业进行准入控制：根据队列、配额、资源可用性和优先级判断作业什么时候可以进入集群运行。

Kueue 中常见概念包括：

\begin{itemize}
  \item \textbf{Workload}：代表一个待调度的作业抽象；
  \item \textbf{LocalQueue}：namespace 内用户提交作业的队列；
  \item \textbf{ClusterQueue}：跨 namespace 的资源配额和队列管理对象；
  \item \textbf{ResourceFlavor}：不同类型资源，例如不同 GPU 型号、不同节点池；
  \item \textbf{Admission}：当资源和配额满足时，Kueue 接纳作业运行；
  \item \textbf{Preemption}：在策略允许时，抢占低优先级 workload。
\end{itemize}

Kueue 的一个典型价值是把“排队”和“运行”分开。用户可以提交大量训练或评测任务，但只有被接纳的任务才真正创建大量 Pod 消耗资源。这对于控制集群压力、实现配额、公平性和多租户非常重要。

一个简化的 Kueue 资源概念如下：

\begin{lstlisting}[caption={Kueue ClusterQueue 与 ResourceFlavor 的概念示例},label={lst:kueue-clusterqueue-example}]
apiVersion: kueue.x-k8s.io/v1beta1
kind: ResourceFlavor
metadata:
  name: h100-nvswitch
---
apiVersion: kueue.x-k8s.io/v1beta1
kind: ClusterQueue
metadata:
  name: foundation-model-queue
spec:
  namespaceSelector: {}
  resourceGroups:
  - coveredResources: ["cpu", "memory", "nvidia.com/gpu"]
    flavors:
    - name: h100-nvswitch
      resources:
      - name: nvidia.com/gpu
        nominalQuota: 256
      - name: cpu
        nominalQuota: 2048
      - name: memory
        nominalQuota: 16Ti
\end{lstlisting}

Kueue 适合与 Job、RayJob、MPIJob、PyTorchJob 等作业控制器结合，作为统一准入和配额层。它尤其适合多团队共享 batch 集群，希望按照资源 flavor 和队列控制不同团队使用资源的场景。

\subsection{Priority 与 Preemption}
\label{subsec:priority-preemption}

优先级和抢占是调度系统必须谨慎设计的能力。没有优先级，紧急任务无法快速获得资源；没有抢占，高优先级任务可能长期等待；但如果抢占过于激进，低优先级任务可能永远无法完成，集群也会频繁浪费已经完成的训练进度。

AI 任务的抢占策略需要结合 checkpoint。对于支持 checkpoint 的训练任务，可以在抢占前发送信号，让任务保存状态；对于不支持 checkpoint 的任务，抢占意味着从头开始，成本很高。对于在线推理服务，抢占必须考虑流量迁移、连接 draining 和 SLA。

常见策略包括：

\begin{itemize}
  \item 高优先级线上服务不被抢占；
  \item 低优先级实验任务可被抢占；
  \item 预训练任务只有在 checkpoint 完成后才允许安全抢占；
  \item 抢占前设置 grace period；
  \item 被抢占任务重新进入队列；
  \item 对频繁被抢占的任务提供 aging 或补偿机制；
  \item 抢占应尽量释放连续资源，而不是随机杀 Pod。
\end{itemize}

\subsection{Checkpoint-Based Rescheduling}
\label{subsec:checkpoint-based-rescheduling}

大规模训练不可能假设节点永远不坏。GPU Xid、ECC error、网络抖动、节点重启、驱动问题、存储故障、作业代码异常都可能导致训练中断。因此，调度系统必须与 checkpoint 机制结合。

Checkpoint-based rescheduling 的流程通常是：

\begin{enumerate}
  \item 训练任务周期性保存 distributed checkpoint；
  \item 调度系统监控 Pod、Node 和 Job 状态；
  \item 节点故障或抢占发生时，Job 标记为失败或待恢复；
  \item 调度器重新为 Job 分配资源；
  \item 训练任务从最近 checkpoint 恢复；
  \item 平台记录失败原因、恢复耗时和损失的训练步数。
\end{enumerate}

对大模型训练来说，checkpoint 本身也是重 I/O 操作。如果所有大作业同时 checkpoint，会冲击存储和网络。调度系统可以参与 checkpoint policy，例如错峰 checkpoint、限制并发 checkpoint、根据作业优先级调整 checkpoint 间隔。

\section{推理服务编排}
\label{sec:serving-orchestration}

训练作业通常是 batch workload，而推理服务是 online workload。它们都使用 GPU，但调度目标完全不同。训练追求吞吐和完成时间，推理追求可用性、延迟、吞吐和成本。Kubernetes 原生适合管理长期运行服务，因此很多 LLM Serving 系统会部署在 Kubernetes 上，但 GPU 推理服务仍然有大量特殊问题。

\subsection{Model Server Deployment}
\label{subsec:model-server-deployment}

一个 LLM 推理服务通常包含：

\begin{itemize}
  \item 前端 API gateway；
  \item tokenizer / request processor；
  \item model server，例如 vLLM、TensorRT-LLM、Triton、Text Generation Inference 或自研 runtime；
  \item scheduler，用于 continuous batching、prefill/decode 调度和 admission control；
  \item KV Cache manager；
  \item metrics exporter；
  \item model artifact loader；
  \item 日志、trace 和计费组件。
\end{itemize}

在 Kubernetes 中，最简单的部署方式是使用 Deployment 或 StatefulSet 管理模型服务 Pod，用 Service 暴露访问入口。对于需要多 GPU tensor parallel 的模型，一个 Pod 可能请求 2、4 或 8 张 GPU，并在容器内部启动多进程 runtime。

\begin{lstlisting}[caption={LLM Model Server Deployment 的概念示例},label={lst:llm-serving-deployment}]
apiVersion: apps/v1
kind: Deployment
metadata:
  name: llama-serving
spec:
  replicas: 2
  selector:
    matchLabels:
      app: llama-serving
  template:
    metadata:
      labels:
        app: llama-serving
    spec:
      nodeSelector:
        pool: serving
        gpu.fabric: nvswitch
      containers:
      - name: server
        image: my-registry/llm-serving:latest
        command: ["bash", "-lc", "python -m server --model /models/llama --tensor-parallel-size 4"]
        ports:
        - containerPort: 8000
        resources:
          limits:
            nvidia.com/gpu: 4
            cpu: "32"
            memory: "256Gi"
\end{lstlisting}

需要注意，Deployment 的滚动升级语义对大模型服务未必总是合适。模型加载可能需要数分钟，显存初始化和 CUDA graph capture 也需要时间。如果滚动策略设置不当，可能造成可用副本不足或冷启动抖动。对于大模型服务，通常需要自定义 readiness probe、预热流程、流量切换和灰度发布策略。

\subsection{Autoscaling}
\label{subsec:autoscaling}

普通 Web 服务常用 CPU utilization 或 QPS 做 autoscaling，但 LLM 推理服务的扩缩容指标更复杂。常见指标包括：

\begin{itemize}
  \item QPS；
  \item active requests；
  \item waiting queue length；
  \item TTFT；
  \item TPOT；
  \item request latency P50/P95/P99；
  \item GPU utilization；
  \item HBM usage；
  \item KV Cache usage；
  \item tokens per second；
  \item prefill queue length 与 decode queue length；
  \item SLA violation rate。
\end{itemize}

LLM serving 的扩缩容难点在于冷启动成本高。一个大模型实例启动需要拉取镜像、加载权重、初始化 CUDA context、构建通信组、预热 kernel，甚至加载量化权重和构建 CUDA graph。这个过程可能远慢于普通服务。因此，autoscaling 不能只在流量升高后才启动新实例，还需要预测、预留和 warm pool。

推理扩缩容的常见策略包括：

\begin{itemize}
  \item 保留最小 warm replicas；
  \item 根据 queue length 和 token throughput 扩容；
  \item 对 prefill-heavy 和 decode-heavy 资源池分别扩缩容；
  \item 使用 admission control 防止队列无限增长；
  \item 在流量低谷缩容，但保留热点模型；
  \item 对多模型服务做模型级别的容量规划。
\end{itemize}

\subsection{GPU Bin Packing}
\label{subsec:gpu-bin-packing}

推理服务通常比训练更容易产生 GPU 碎片。原因包括：

\begin{itemize}
  \item 不同模型大小不同，显存需求不同；
  \item 不同模型 QPS 不同，SM 和 HBM 带宽使用率不同；
  \item 一些模型需要 1 张 GPU，一些需要 4 张或 8 张 GPU；
  \item LoRA adapter、多租户模型和 embedding/rerank 服务共存；
  \item 在线服务不能随意迁移，否则会影响 SLA。
\end{itemize}

Bin packing 的目标是把模型实例合理放置到 GPU 节点上，提高资源利用率并减少碎片。例如，一个 8 GPU 节点可以放置两个 TP=4 的模型实例，也可以放置一个 TP=8 的大模型实例，还可以放置多个 MIG 小实例。不同放置方式对显存、通信、带宽和可用性影响不同。

推理 bin packing 需要考虑：

\begin{itemize}
  \item GPU 显存；
  \item HBM 带宽；
  \item SM utilization；
  \item KV Cache 峰值；
  \item 模型权重常驻；
  \item TP group 拓扑；
  \item 实例之间的干扰；
  \item 负载均衡和故障域。
\end{itemize}

如果只按显存打包，可能把多个 memory-bound decode 服务放在同一节点，导致 HBM 带宽争用和尾延迟上升。如果只按 SM utilization 打包，也可能忽略 KV Cache 和网络通信。因此，高质量推理调度需要 workload profiling。

\subsection{Cold Start}
\label{subsec:cold-start}

大模型推理的冷启动是一个严重问题。普通微服务冷启动可能是秒级，而大模型服务可能需要几十秒到数分钟。冷启动包括：

\begin{itemize}
  \item 镜像拉取；
  \item 模型权重从对象存储或分布式文件系统加载；
  \item 权重反序列化和分片；
  \item GPU 显存分配；
  \item CUDA context 初始化；
  \item NCCL communicator 初始化；
  \item kernel autotuning；
  \item CUDA graph capture；
  \item KV Cache block 初始化；
  \item readiness probe 通过。
\end{itemize}

降低冷启动的策略包括：

\begin{itemize}
  \item 使用本地 NVMe 缓存模型权重；
  \item 预拉取镜像；
  \item 使用 warm pool；
  \item 模型分层加载；
  \item 权重格式优化；
  \item 避免频繁 scale-to-zero；
  \item 将热点模型常驻；
  \item 使用灰度流量切换，避免未预热实例接收真实请求。
\end{itemize}

\subsection{多模型共存}
\label{subsec:multi-model-coexistence}

一个生产集群通常同时服务多个模型：大模型、小模型、embedding、rerank、多模态、代码模型、专用领域模型、实验版本和灰度版本。多模型共存带来几个调度问题：

\begin{itemize}
  \item 热点模型需要更多副本，冷门模型是否常驻？
  \item 多个小模型是否合并到同一 serving runtime？
  \item LoRA adapter 是否动态加载？
  \item 模型之间如何共享 KV Cache 或 prefix cache？
  \item 不同模型的 SLA 是否不同？
  \item 模型升级时如何做流量切换和回滚？
\end{itemize}

\begin{figure}[H]
  \centering
  \begin{tikzpicture}[
    font=\small,
    box/.style={draw, rounded corners, minimum width=2.1cm, minimum height=0.7cm, align=center},
    wide/.style={draw, rounded corners, minimum width=8.5cm, minimum height=0.75cm, align=center},
    arrow/.style={-{Stealth[length=2mm]}, thick}
  ]
    \node[wide] (client) at (0,5.0) {Clients / Applications};
    \node[wide] (gw) at (0,4.0) {API Gateway / Auth / Rate Limit};
    \node[wide] (router) at (0,3.0) {Model Router / Admission Control};

    \node[box] (s0) at (-3.2,1.8) {Prefill Pool\\GPU Pods};
    \node[box] (s1) at (0,1.8) {Decode Pool\\GPU Pods};
    \node[box] (s2) at (3.2,1.8) {Embedding /\\Rerank};

    \node[wide] (obs) at (0,0.6) {Metrics / Logs / Traces / Cost Accounting};
    \node[wide] (store) at (0,-0.55) {Model Registry / Object Storage / NVMe Cache};

    \draw[arrow] (client) -- (gw);
    \draw[arrow] (gw) -- (router);
    \draw[arrow] (router) -- (s0);
    \draw[arrow] (router) -- (s1);
    \draw[arrow] (router) -- (s2);
    \draw[arrow] (s0) -- (obs);
    \draw[arrow] (s1) -- (obs);
    \draw[arrow] (s2) -- (obs);
    \draw[arrow] (store) -- (s0);
    \draw[arrow] (store) -- (s1);
    \draw[arrow] (store) -- (s2);
  \end{tikzpicture}
  \caption{LLM Serving on Kubernetes 架构图：模型路由、GPU Pod、观测系统和模型存储共同构成在线服务平台。}
  \label{fig:llm-serving-on-k8s}
\end{figure}

\section{AI 集群可观测性}
\label{sec:ai-cluster-observability}

没有可观测性，就没有真正的调度优化。AI 集群的可观测性必须跨越 Kubernetes、GPU、网络、存储、作业和业务指标。只看 GPU utilization 会误导决策；只看 Pod 状态也无法解释训练慢在哪里。

\subsection{GPU Utilization}
\label{subsec:gpu-utilization}

GPU utilization 是最常见指标，但也是最容易被误解的指标。它通常表示一段采样时间内 GPU 是否有活跃 work，而不是 Tensor Core 是否高效运行。一个 NCCL kernel、一个 memory-bound kernel、一个低效小 kernel 都可能让 GPU utilization 看起来不低，但业务吞吐并不好。

因此，GPU utilization 应与以下指标一起看：

\begin{itemize}
  \item SM utilization；
  \item Tensor Core utilization；
  \item HBM bandwidth；
  \item GPU memory usage；
  \item kernel timeline；
  \item NCCL communication time；
  \item tokens per second 或 samples per second；
  \item cost per token 或 cost per step。
\end{itemize}

\subsection{SM Occupancy}
\label{subsec:sm-occupancy-monitoring}

SM occupancy 或 SM active 反映计算单元活跃程度。训练中的大 GEMM 通常希望 SM 和 Tensor Core 利用率较高；推理 decode 阶段则可能因为 batch 小、KV Cache 读取和小 GEMM 导致 SM 利用率不高。

调度系统可以利用 SM 指标判断混部机会。例如一个服务显存占用高但 SM 使用低，理论上可以与另一个计算型任务混部；但如果它对延迟敏感，混部可能导致 SLA 抖动。因此，资源利用率指标必须与 workload 类型和 SLA 结合。

\subsection{HBM Usage}
\label{subsec:hbm-usage-monitoring}

HBM usage 是训练和推理都非常关键的指标。训练中，HBM 被参数、梯度、优化器状态、activation 和临时 buffer 占用；推理中，HBM 被模型权重、KV Cache、workspace 和 runtime buffer 占用。

推理服务尤其需要监控 KV Cache 使用率。如果 KV Cache block 使用率接近上限，新的请求可能被拒绝、排队或触发 eviction。调度系统可以根据 KV Cache 压力而不是只根据 QPS 做扩容。

\subsection{Network Throughput}
\label{subsec:network-throughput-monitoring}

网络指标包括节点内和节点间两类：

\begin{itemize}
  \item 节点内：NVLink/NVSwitch 带宽、PCIe 带宽、GPU P2P；
  \item 节点间：NIC TX/RX、RDMA counters、PFC、ECN、丢包、重传、交换机端口利用率。
\end{itemize}

对于训练作业，需要把 NCCL collective 时间与网络指标关联起来。如果某个 job step time 变长，同时某些 leaf-spine 链路利用率接近上限，说明可能存在网络热点。如果 NCCL timeout 与 PFC pause storm 同时出现，说明 RoCE 拥塞控制需要排查。

\subsection{Job Queue Time}
\label{subsec:job-queue-time}

队列等待时间是 AI 平台用户最关心的指标之一。一个训练任务真正的完成时间不是运行时间，而是：

\[
  T_{\text{completion}} = T_{\text{queue}} + T_{\text{startup}} + T_{\text{run}} + T_{\text{recovery}}.
\]

其中 $T_{\text{queue}}$ 可能非常长。大作业经常因为资源碎片等待，小作业则可能因为低优先级等待。平台应监控：

\begin{itemize}
  \item 不同队列的平均等待时间；
  \item P95/P99 等待时间；
  \item 不同 GPU flavor 的等待时间；
  \item 大作业与小作业等待差异；
  \item 被抢占次数；
  \item backfill 成功率；
  \item 资源碎片率。
\end{itemize}

队列时间可以反向指导容量规划。如果 H100 队列长期拥塞，而 A100 队列空闲，说明资源结构与需求不匹配；如果 8 GPU 单节点作业长期等待，但单卡任务很快启动，说明 bin packing 或 reservation 策略需要调整。

\subsection{Cost Per Token / Cost Per Step}
\label{subsec:cost-per-token-step}

最终，AI Infra 的目标不是让某个指标好看，而是降低单位业务产出的成本。训练中常用 cost per step、cost per token、tokens per GPU-hour；推理中常用 cost per 1K tokens、cost per request、GPU-hour per SLA bucket。

这些指标把硬件利用率、调度效率和业务结果连接起来。例如：

\begin{itemize}
  \item 一个推理服务 GPU utilization 很高，但 P99 延迟也很高，可能需要扩容；
  \item 一个训练作业 GPU utilization 不低，但 tokens per GPU-hour 低，可能是通信瓶颈；
  \item 一个队列 GPU 使用率高，但大量作业被抢占，实际有效训练 token 低；
  \item 一个量化模型成本低，但质量下降导致重试率上升，端到端成本未必低。
\end{itemize}

\begin{figure}[H]
  \centering
  \begin{tikzpicture}[
    font=\small,
    layer/.style={draw, rounded corners, minimum width=9.0cm, minimum height=0.7cm, align=center},
    arrow/.style={-{Stealth[length=2mm]}, thick}
  ]
    \node[layer] (biz) at (0,5.0) {业务层：QPS、Latency、SLA、Cost per Token、用户体验};
    \node[layer] (job) at (0,4.0) {作业层：Queue Time、Run Time、Checkpoint、Failure、Preemption};
    \node[layer] (k8s) at (0,3.0) {Kubernetes 层：Pod、Node、Scheduler、Quota、Events};
    \node[layer] (gpu) at (0,2.0) {GPU 层：Utilization、SM、HBM、Tensor Core、Xid、ECC};
    \node[layer] (net) at (0,1.0) {网络层：NVLink、RDMA、NIC、PFC、ECN、NCCL Bandwidth};
    \node[layer] (storage) at (0,0.0) {存储层：Dataset、Checkpoint、Model Loading、NVMe Cache};

    \draw[arrow] (storage) -- (net);
    \draw[arrow] (net) -- (gpu);
    \draw[arrow] (gpu) -- (k8s);
    \draw[arrow] (k8s) -- (job);
    \draw[arrow] (job) -- (biz);
  \end{tikzpicture}
  \caption{AI Infra 监控指标分层图：从硬件资源到业务成本需要跨层关联。}
  \label{fig:ai-infra-monitoring-layers}
\end{figure}

\section{从 Kubernetes 到 AI 平台}
\label{sec:from-kubernetes-to-ai-platform}

Kubernetes 是强大的基础设施底座，但一个完整 AI 平台通常还需要更上层的能力：

\begin{itemize}
  \item 训练任务提交系统；
  \item 模型注册与版本管理；
  \item 数据集管理；
  \item 实验追踪；
  \item 分布式 checkpoint 管理；
  \item 镜像构建与依赖管理；
  \item 权限、审计和成本分摊；
  \item 自动调参和超参搜索；
  \item 模型评测流水线；
  \item 推理流量治理；
  \item 多集群和多区域调度；
  \item 故障自动诊断。
\end{itemize}

因此，不要把“在 Kubernetes 上跑 GPU Pod”误认为“已经有了 AI 平台”。真正的平台需要把用户体验、资源效率、可靠性和成本闭环起来。对用户来说，他们希望提交一个训练配置或模型部署配置，而不是手写复杂 YAML；对平台来说，需要把这些高级意图翻译成队列、配额、Pod、Service、PVC、Secret、ConfigMap、Job、监控和告警。

一个成熟 AI 平台通常会形成如下分层：

\begin{enumerate}
  \item \textbf{硬件层}：GPU、NIC、NVMe、交换机、机架和供电；
  \item \textbf{系统层}：驱动、CUDA、container runtime、RDMA、存储客户端；
  \item \textbf{Kubernetes 层}：Pod、Node、Service、Scheduler、Device Plugin、Operator；
  \item \textbf{Batch / Serving 层}：Volcano、Kueue、Training Operator、Serving Runtime；
  \item \textbf{平台层}：任务提交、模型管理、数据管理、权限、计费、可观测性；
  \item \textbf{业务层}：训练产出、推理 SLA、模型质量、成本收益。
\end{enumerate}

AI Infra 工程师的价值就在于能跨越这些层次定位问题。例如用户反馈“训练慢”，原因可能是模型代码、GEMM shape、NCCL all-reduce、RoCE 拥塞、GPU-NIC 亲和、Pod 被错误调度、存储读取慢或 checkpoint 阻塞。只懂其中一层，很难完成端到端优化。

\section{本章小结}
\label{sec:k8s-ai-scheduling-summary}

本章介绍了 Kubernetes 在 AI 集群中的调度角色。我们从 GPU 资源昂贵、多租户、训练推理混部、利用率与公平性的矛盾出发，解释了为什么 AI 集群需要比普通 Kubernetes 更强的调度能力。

本章核心要点如下：

\begin{itemize}
  \item AI 集群调度不是简单分配 GPU 数量，而是要同时考虑 GPU 型号、显存、拓扑、网络、队列、优先级、SLA 和成本。
  \item Kubernetes 的 Pod、Node、Scheduler、Device Plugin 和 Operator 构成了 GPU 集群管理的基础。
  \item NVIDIA Device Plugin 把 GPU 暴露为 Kubernetes 扩展资源，GPU Operator 可以进一步管理驱动、监控、MIG 等组件。
  \item MIG 适合小模型推理和多租户隔离，GPU sharing 适合低优先级或可容忍抖动的共享场景，但二者都需要谨慎处理性能隔离。
  \item 拓扑感知调度对 TP、EP、多 rail RDMA 和多卡推理至关重要，不能只按 GPU count 调度。
  \item 分布式训练需要 gang scheduling，避免 worker 部分启动造成 GPU 浪费。
  \item Volcano 提供面向 batch/HPC/AI workload 的队列、gang、优先级、抢占和 backfill 等能力。
  \item Kueue 提供 Kubernetes 原生作业队列、配额和准入控制，适合多租户 batch/AI 集群。
  \item 推理服务编排关注 Deployment、autoscaling、bin packing、cold start、多模型共存和 SLA。
  \item AI 集群可观测性必须跨越业务、作业、Kubernetes、GPU、网络和存储层，最终落到 cost per token 与 cost per step。
\end{itemize}

至此，本书第五篇完成了从 GPU 硬件到集群调度的完整路径：单卡计算、单机互联、多机网络、Kubernetes 调度。结合前面几篇内容，我们已经从数学、CUDA、Transformer、分布式训练、推理优化一路走到了 AI Infrastructure 的完整系统视角。后续附录将补充符号约定、CUDA 与分布式环境搭建、AI 集群上线检查清单以及参考资料索引，帮助读者把本书内容进一步落到工程实践中。


% ==================================================
% 附录
% ==================================================
\appendix

\part{附录}

% appendices/appendix_a_symbols.tex

\chapter{数学符号与张量维度约定}

在大模型工程中，许多问题表面上看是系统问题，实际根源却来自张量维度、矩阵乘法方向、并行切分方式或内存布局理解不一致。例如，同样是一个 attention tensor，在论文中可能写作 $[B, H, S, D_h]$，在框架实现中可能变成 $[S, B, H, D_h]$，在 CUDA kernel 中又可能被重排为适合访存合并的连续布局。本附录统一本书使用的数学符号、维度命名和公式排版规则，方便读者在阅读后续章节时快速对齐语义。

\section{常用数学符号}

\subsection{标量、向量、矩阵与张量}

本书使用如下约定：

\begin{longtable}{p{0.22\textwidth}p{0.68\textwidth}}
\toprule
符号 & 含义 \
\midrule
$a, b, c$ & 标量，通常表示单个数值，例如 loss、学习率、温度参数等。 \
$\vect{x}, \vect{y}, \vect{z}$ & 向量，通常表示一个 token 的 hidden state、一个参数向量或一个概率分布。 \
$\mat{A}, \mat{B}, \mat{W}$ & 矩阵，通常表示线性层权重、二维激活或批量矩阵乘法中的单个矩阵。 \
$\tensor{X}, \tensor{Y}, \tensor{Z}$ & 高维张量，通常表示 batch、sequence、head、hidden 等多个维度组成的数据。 \
$\theta$ & 模型参数集合。 \
$\nabla_{\theta}\mathcal{L}$ & loss 对模型参数的梯度。 \
$\mathcal{L}$ & 损失函数。 \
$\eta$ & 学习率。 \
\bottomrule
\end{longtable}

在具体实现中，一个变量到底是向量、矩阵还是高维张量，取决于上下文。例如，Transformer 中的 hidden states 在单个 token 视角下是向量，在单个序列视角下是矩阵，在 batch 视角下是三维张量。

\subsection{索引与集合}

本书常用如下索引：

\begin{longtable}{p{0.22\textwidth}p{0.68\textwidth}}
\toprule
符号 & 含义 \
\midrule
$i, j, k$ & 通用索引，常用于矩阵元素、token 位置或循环变量。 \
$b$ & batch 中的样本索引。 \
$s, t$ & sequence 中的 token 位置索引。 \
$h$ & attention head 索引。 \
$l$ & Transformer 层索引。 \
$r$ & 分布式训练中的 rank 索引。 \
$N$ & 通用数量，例如样本数、节点数、GPU 数等。 \
$V$ & 词表大小 vocabulary size。 \
\bottomrule
\end{longtable}

若存在歧义，本书会在公式前显式说明索引含义。例如：
[
\tensor{X}_{b,s,:}
]
表示第 $b$ 个样本、第 $s$ 个 token 对应的 hidden vector。

\section{Transformer 中的维度命名}

Transformer 工程实现中最容易出错的地方之一，是不同维度含义相近但用途不同。尤其在多头注意力、KV Cache、张量并行和推理服务章节中，维度命名必须保持一致。

\subsection{基础维度}

本书默认使用如下维度符号：

\begin{longtable}{p{0.2\textwidth}p{0.24\textwidth}p{0.46\textwidth}}
\toprule
符号 & 名称 & 含义 \
\midrule
$B$ & batch size & 一个 step 或一次推理迭代中同时处理的样本数量。 \
$S$ & sequence length & 输入序列长度，也可表示 prefill 阶段上下文长度。 \
$T$ & generated length & 自回归生成阶段生成的 token 数。 \
$D$ & hidden size & 每个 token 的隐藏状态维度。 \
$H$ & number of heads & attention head 数量。 \
$D_h$ & head dimension & 单个 attention head 的维度，通常满足 $D = H \cdot D_h$。 \
$D_{\mathrm{ff}}$ & FFN intermediate size & Feed-Forward Network 中间层维度。 \
$V$ & vocabulary size & tokenizer 对应的词表大小。 \
$L$ & number of layers & Transformer block 层数。 \
\bottomrule
\end{longtable}

典型 Decoder-only LLM 的主要张量形状如下：

\begin{longtable}{p{0.28\textwidth}p{0.28\textwidth}p{0.34\textwidth}}
\toprule
对象 & 形状 & 说明 \
\midrule
Token IDs & $[B, S]$ & 输入 token 编号。 \
Embedding Output & $[B, S, D]$ & token embedding 后的隐藏状态。 \
Q / K / V & $[B, S, H, D_h]$ & 多头注意力中的 query、key、value。 \
Attention Scores & $[B, H, S, S]$ & prefill 阶段完整 self-attention 分数矩阵。 \
Attention Output & $[B, S, H, D_h]$ & 每个 head 的 attention 输出。 \
Merged Attention Output & $[B, S, D]$ & 多头合并后的输出。 \
FFN Intermediate & $[B, S, D_{\mathrm{ff}}]$ & FFN 上投影后的中间激活。 \
Logits & $[B, S, V]$ & 每个 token 位置上的词表 logits。 \
\bottomrule
\end{longtable}

\subsection{MHA、MQA 与 GQA 中的 head 维度}

标准 Multi-Head Attention 中，Query、Key、Value 通常具有相同数量的 heads：
[
H_q = H_k = H_v = H.
]

在 Multi-Query Attention 中，多个 query heads 共享一组 key/value heads：
[
H_q = H,\qquad H_k = H_v = 1.
]

在 Grouped-Query Attention 中，query heads 被分组，每组共享一个 key/value head：
[
H_q = H,\qquad H_k = H_v = H_{\mathrm{kv}},\qquad H_{\mathrm{kv}} < H.
]

因此，KV Cache 的大小与 $H_{\mathrm{kv}}$ 直接相关，而不是只与 $H_q$ 相关。推理系统中估算显存时，应使用：
[
\mathrm{KVCacheBytes}
=====================

2 \cdot B \cdot S \cdot L \cdot H_{\mathrm{kv}} \cdot D_h \cdot \mathrm{bytes_per_element}.
]
其中系数 $2$ 分别对应 Key Cache 与 Value Cache。

\section{batch、sequence、hidden、head 维度约定}

\subsection{默认逻辑布局}

除非特别说明，本书在数学推导中默认使用 batch-first 布局：
[
\tensor{X} \in \mathbb{R}^{B \times S \times D}.
]

这与许多现代 PyTorch 模型实现保持一致，便于读者理解 batch、sequence 与 hidden 的关系。在讲解注意力时，本书通常将 Q/K/V 写作：
[
\tensor{Q}, \tensor{K}, \tensor{V}
\in
\mathbb{R}^{B \times S \times H \times D_h}.
]

实际工程中，为了适配 kernel 或提高访存效率，常见布局还包括：

\begin{itemize}
\item $[S, B, D]$：早期序列模型和部分框架内部实现常用。
\item $[B, H, S, D_h]$：attention kernel 中常用，方便按 head 并行。
\item $[B, S, H \cdot D_h]$：linear projection 前后的合并 hidden 表示。
\item $[L, B, H_{\mathrm{kv}}, S, D_h]$：KV Cache 的一种逻辑表达方式。
\end{itemize}

本书会在涉及 kernel 或内存布局时明确区分“逻辑形状”和“物理布局”。逻辑形状描述语义，物理布局描述内存中连续存储的顺序。

\subsection{矩阵乘法方向}

线性层通常写作：
[
\mat{Y} = \mat{X}\mat{W} + \vect{b},
]
其中：
[
\mat{X} \in \mathbb{R}^{N \times D_{\mathrm{in}}},
\qquad
\mat{W} \in \mathbb{R}^{D_{\mathrm{in}} \times D_{\mathrm{out}}},
\qquad
\mat{Y} \in \mathbb{R}^{N \times D_{\mathrm{out}}}.
]

在 PyTorch 的 \texttt{nn.Linear(in_features, out_features)} 中，权重参数的存储形状通常是：
[
[D_{\mathrm{out}}, D_{\mathrm{in}}],
]
前向计算在语义上等价于：
[
\mat{Y} = \mat{X}\mat{W}^{\top} + \vect{b}.
]

因此，阅读代码时不要只根据权重 tensor 的 shape 判断数学公式方向，而要结合框架 API 的约定。

\section{本书公式排版规范}

\subsection{公式与代码之间的对应关系}

本书尽量保持公式与代码之间的一一对应。例如，Scaled Dot-Product Attention 写作：
[
\Attention(\mat{Q}, \mat{K}, \mat{V})
=====================================

\softmax\left(
\frac{\mat{Q}\mat{K}^{\top}}{\sqrt{D_h}}
\right)\mat{V}.
]

对应到伪代码时，通常写为：

\begin{lstlisting}[language=Python,caption={Scaled Dot-Product Attention 的基本形式},label={lst:appendix-attention-basic}]
scores = q @ k.transpose(-2, -1)
scores = scores / math.sqrt(head_dim)
probs = torch.softmax(scores, dim=-1)
out = probs @ v
\end{lstlisting}

这里的 \texttt{transpose(-2, -1)} 对应公式中的 $\mat{K}^{\top}$，\texttt{dim=-1} 表示沿 key sequence 维度做 softmax。

\subsection{复杂度记号}

本书使用 $\mathcal{O}(\cdot)$ 表示渐进复杂度。若讨论 Transformer attention 的计算复杂度，常写作：
[
\mathcal{O}(B \cdot H \cdot S^2 \cdot D_h).
]

若讨论 KV Cache 显存复杂度，常写作：
[
\mathcal{O}(B \cdot L \cdot H_{\mathrm{kv}} \cdot S \cdot D_h).
]

需要注意，复杂度表达式只表达增长趋势，不等价于真实运行时间。真实性能还受到 kernel fusion、访存模式、HBM 带宽、Tensor Core 利用率、通信拓扑和调度策略影响。

\subsection{单位约定}

本书在系统性能章节中常用如下单位：

\begin{longtable}{p{0.22\textwidth}p{0.68\textwidth}}
\toprule
单位 & 含义 \
\midrule
FLOPs & 浮点运算次数。 \
TFLOPS & 每秒万亿次浮点运算。 \
GB / GiB & 十进制 GB 与二进制 GiB。工程排查中需注意二者差异。 \
GB/s & 每秒传输的字节数，常用于显存带宽、PCIe 带宽、网络带宽。 \
tokens/s & 推理或训练中每秒处理的 token 数。 \
steps/s & 训练中每秒完成的 optimizer step 数。 \
TTFT & Time To First Token，推理请求首 token 延迟。 \
TPOT & Time Per Output Token，生成阶段单 token 平均耗时。 \
\bottomrule
\end{longtable}

\section{附录 A 小结}

本附录定义了本书使用的主要数学符号、Transformer 维度命名、张量布局约定与公式排版规范。读者在阅读训练并行、推理优化、GPU kernel 和集群调度章节时，可以随时回到本附录对齐变量含义。尤其在分析 KV Cache、Tensor Parallelism、FlashAttention 与 Continuous Batching 时，明确维度语义往往比记住某个公式更重要。


% appendices/appendix_b_cuda_environment.tex

\chapter{CUDA 与分布式训练环境搭建}

AI Infrastructure 工程的第一道门槛，通常不是模型结构，而是环境。一个看似简单的 \texttt{import torch}，背后涉及 NVIDIA Driver、CUDA Runtime、cuDNN、NCCL、PyTorch wheel、GPU 架构、容器运行时、网络驱动和权限配置。环境配置错误会导致大量隐蔽问题：训练一启动就崩溃、NCCL 卡死、GPU 不可见、性能只有预期的一半、不同节点行为不一致等。

本附录给出一套面向单机多卡与多机训练的环境搭建和排查方法。它不是某个固定版本的安装手册，而是一套可迁移的检查框架。

\section{CUDA Toolkit}

\subsection{CUDA Toolkit 的角色}

CUDA Toolkit 包含编译 CUDA 程序所需的编译器、头文件、库文件和工具链。常见组件包括：

\begin{itemize}
\item \texttt{nvcc}：CUDA C++ 编译器。
\item CUDA Runtime：运行 CUDA 程序所需的运行时库。
\item cuBLAS：GPU 上的 BLAS 矩阵运算库。
\item cuFFT、cuRAND、cuSPARSE：常用数学库。
\item Nsight Systems / Nsight Compute：性能分析工具。
\end{itemize}

对于只运行 PyTorch 的用户，很多预编译 wheel 已经内置了相应 CUDA Runtime，不一定需要系统安装完整 CUDA Toolkit。但如果要编译自定义 CUDA Extension、Triton kernel、DeepSpeed fused op 或 Megatron-LM 中的扩展模块，通常需要安装匹配的 Toolkit。

\subsection{检查 CUDA Toolkit}

可以使用以下命令检查 \texttt{nvcc}：

\begin{lstlisting}[language=bash,caption={检查 CUDA Toolkit},label={lst:check-cuda-toolkit}]
which nvcc
nvcc --version
\end{lstlisting}

也可以检查 CUDA 安装目录：

\begin{lstlisting}[language=bash,caption={检查 CUDA 安装目录},label={lst:check-cuda-path}]
ls -l /usr/local/cuda
ls -l /usr/local/cuda/bin
ls -l /usr/local/cuda/lib64
\end{lstlisting}

常见环境变量如下：

\begin{lstlisting}[language=bash,caption={CUDA 常见环境变量},label={lst:cuda-env-vars}]
export CUDA_HOME=/usr/local/cuda
export PATH=$CUDA_HOME/bin:$PATH
export LD_LIBRARY_PATH=$CUDA_HOME/lib64:$LD_LIBRARY_PATH
\end{lstlisting}

\section{NVIDIA Driver}

\subsection{Driver 与 CUDA 的关系}

NVIDIA Driver 负责让操作系统识别和管理 GPU。CUDA 应用最终需要通过 Driver 与 GPU 交互。需要特别注意：

\begin{itemize}
\item Driver 安装在宿主机层面。
\item CUDA Toolkit 可以安装在宿主机，也可以随容器镜像提供。
\item 容器中的 CUDA Runtime 必须与宿主机 Driver 兼容。
\item 多节点训练时，不同节点的 Driver 版本最好保持一致。
\end{itemize}

在容器化环境中，一个常见误区是以为容器内安装了 CUDA 就可以直接使用 GPU。实际上，容器仍然依赖宿主机 Driver 以及 NVIDIA Container Toolkit 将 GPU 设备和驱动库挂载进容器。

\subsection{检查 Driver 与 GPU 状态}

最常用的检查命令是：

\begin{lstlisting}[language=bash,caption={检查 NVIDIA Driver 与 GPU 状态},label={lst:nvidia-smi-basic}]
nvidia-smi
\end{lstlisting}

重点观察以下信息：

\begin{itemize}
\item Driver Version：宿主机驱动版本。
\item CUDA Version：Driver 支持的最高 CUDA 运行时版本。
\item GPU Name：GPU 型号。
\item Memory-Usage：显存占用。
\item GPU-Util：GPU 利用率。
\item Processes：当前占用 GPU 的进程。
\end{itemize}

检查 GPU 拓扑：

\begin{lstlisting}[language=bash,caption={检查 GPU 拓扑},label={lst:nvidia-smi-topo}]
nvidia-smi topo -m
\end{lstlisting}

该命令对单机多卡训练尤其重要。它可以帮助判断 GPU 之间通过 PCIe、NVLink 还是 NVSwitch 互联，也可以观察 GPU 与 NIC 的亲和关系。

\section{PyTorch CUDA 环境}

\subsection{检查 PyTorch 是否识别 CUDA}

在 Python 中执行：

\begin{lstlisting}[language=Python,caption={检查 PyTorch CUDA 可用性},label={lst:pytorch-cuda-check}]
import torch

print("torch version:", torch.**version**)
print("cuda available:", torch.cuda.is_available())
print("cuda version:", torch.version.cuda)
print("device count:", torch.cuda.device_count())

for i in range(torch.cuda.device_count()):
print(i, torch.cuda.get_device_name(i))
\end{lstlisting}

若 \texttt{torch.cuda.is_available()} 返回 \texttt{False}，常见原因包括：

\begin{itemize}
\item 安装的是 CPU-only 版本 PyTorch。
\item 宿主机 Driver 未正确安装。
\item 容器未正确挂载 GPU。
\item \texttt{CUDA_VISIBLE_DEVICES} 设置错误。
\item Driver 与 CUDA Runtime 不兼容。
\end{itemize}

\subsection{简单矩阵乘法测试}

安装完成后，不应只检查 CUDA 是否可用，还应做一次实际计算：

\begin{lstlisting}[language=Python,caption={PyTorch GPU 矩阵乘法测试},label={lst:pytorch-matmul-test}]
import torch
import time

device = "cuda"
n = 8192

a = torch.randn(n, n, device=device, dtype=torch.float16)
b = torch.randn(n, n, device=device, dtype=torch.float16)

torch.cuda.synchronize()
start = time.time()

c = a @ b

torch.cuda.synchronize()
end = time.time()

print("elapsed:", end - start)
print("result:", c[0, 0].item())
\end{lstlisting}

该测试不能替代专业 benchmark，但可以快速确认 GPU 计算链路是否基本正常。

\section{NCCL 环境变量}

\subsection{NCCL 的角色}

NCCL 是 NVIDIA Collective Communications Library，负责多 GPU 和多节点之间的集合通信。DDP、Tensor Parallelism、Pipeline Parallelism、ZeRO 等训练策略都会依赖 NCCL 完成 AllReduce、AllGather、ReduceScatter、Broadcast 等通信操作。

当分布式训练卡住时，问题经常不在模型，而在 NCCL 初始化、网络接口选择、拓扑发现或防火墙配置。

\subsection{常用 NCCL 环境变量}

\begin{longtable}{p{0.32\textwidth}p{0.58\textwidth}}
\toprule
环境变量 & 用途 \
\midrule
\texttt{NCCL_DEBUG} & 打印 NCCL 日志，常设为 \texttt{INFO}。 \
\texttt{NCCL_DEBUG_SUBSYS} & 指定日志子系统，如 \texttt{INIT}、\texttt{COLL}、\texttt{NET}。 \
\texttt{NCCL_SOCKET_IFNAME} & 指定用于通信的网卡接口。 \
\texttt{NCCL_IB_HCA} & 指定 InfiniBand 或 RoCE HCA。 \
\texttt{NCCL_IB_DISABLE} & 是否禁用 IB/RDMA。 \
\texttt{NCCL_P2P_DISABLE} & 是否禁用 GPU P2P。 \
\texttt{NCCL_SHM_DISABLE} & 是否禁用共享内存通信。 \
\texttt{NCCL_NET_GDR_LEVEL} & 控制 GPUDirect RDMA 使用策略。 \
\bottomrule
\end{longtable}

调试时可以先打开详细日志：

\begin{lstlisting}[language=bash,caption={NCCL 调试日志},label={lst:nccl-debug-env}]
export NCCL_DEBUG=INFO
export NCCL_DEBUG_SUBSYS=INIT,NET,COLL
\end{lstlisting}

指定网络接口示例：

\begin{lstlisting}[language=bash,caption={指定 NCCL 网络接口},label={lst:nccl-ifname}]
export NCCL_SOCKET_IFNAME=eth0
\end{lstlisting}

实际生产环境中，接口名可能是 \texttt{bond0}、\texttt{ib0}、\texttt{ens5f0} 等，应根据机器网络配置确定。

\subsection{NCCL Tests}

建议在运行真实训练任务前，先使用 NCCL Tests 检查通信性能：

\begin{lstlisting}[language=bash,caption={NCCL all_reduce 性能测试示例},label={lst:nccl-tests}]
./build/all_reduce_perf -b 8 -e 1G -f 2 -g 8
\end{lstlisting}

多节点场景中，可以通过 \texttt{mpirun}、\texttt{torchrun} 或集群调度系统启动。重点观察：

\begin{itemize}
\item 通信是否能正常初始化。
\item bus bandwidth 是否接近预期。
\item 单机内与跨节点带宽是否存在异常差距。
\item 是否出现 timeout、connection refused、unhandled system error 等错误。
\end{itemize}

\section{DeepSpeed / Megatron-LM 安装}

\subsection{DeepSpeed 安装检查}

DeepSpeed 常用于 ZeRO、offload、混合精度训练和大模型训练优化。安装后可执行：

\begin{lstlisting}[language=bash,caption={DeepSpeed 环境检查},label={lst:deepspeed-env-report}]
ds_report
\end{lstlisting}

该命令会显示当前环境中各类 fused op 是否可用。若某些 op 未编译成功，DeepSpeed 仍可能运行，但性能可能明显下降。

常见问题包括：

\begin{itemize}
\item \texttt{CUDA_HOME} 未设置。
\item \texttt{nvcc} 不存在。
\item PyTorch CUDA 版本与系统 CUDA Toolkit 不匹配。
\item 编译器版本过旧。
\item 容器中缺少构建工具链。
\end{itemize}

\subsection{Megatron-LM 环境要点}

Megatron-LM 通常依赖如下能力：

\begin{itemize}
\item 多进程分布式启动。
\item NCCL 正常工作。
\item 可选 fused kernels。
\item 数据集预处理工具。
\item 与 Tensor Parallelism / Pipeline Parallelism 匹配的 rank mapping。
\end{itemize}

启动前应确认：

\begin{lstlisting}[language=bash,caption={分布式启动基础环境变量},label={lst:distributed-env-vars}]
export MASTER_ADDR=127.0.0.1
export MASTER_PORT=29500
export WORLD_SIZE=8
export RANK=0
export LOCAL_RANK=0
\end{lstlisting}

在真实多节点任务中，这些变量通常由 \texttt{torchrun}、MPI、Slurm、Kubernetes Operator 或训练平台自动注入。

\section{常见环境问题排查}

\subsection{问题一：GPU 不可见}

现象：

\begin{itemize}
\item \texttt{nvidia-smi} 看不到 GPU。
\item \texttt{torch.cuda.is_available()} 返回 \texttt{False}。
\item 容器内没有 \texttt{/dev/nvidia*} 设备。
\end{itemize}

排查路径：

\begin{enumerate}
\item 在宿主机执行 \texttt{nvidia-smi}。
\item 检查 Driver 是否安装成功。
\item 检查容器启动参数是否包含 GPU 访问权限。
\item 检查 Kubernetes 中 NVIDIA Device Plugin 是否正常。
\item 检查 \texttt{CUDA_VISIBLE_DEVICES} 是否被错误设置为空。
\end{enumerate}

\subsection{问题二：PyTorch CUDA 版本混乱}

现象：

\begin{itemize}
\item 编译 extension 失败。
\item 运行时报找不到 \texttt{libcudart}、\texttt{libcublas} 或 \texttt{libnccl}。
\item 安装多个 CUDA 版本后行为不一致。
\end{itemize}

排查路径：

\begin{lstlisting}[language=bash,caption={检查 CUDA 相关路径},label={lst:check-cuda-paths}]
which python
python -c "import torch; print(torch.**version**, torch.version.cuda)"
which nvcc
nvcc --version
echo $CUDA_HOME
echo $LD_LIBRARY_PATH
\end{lstlisting}

原则是：运行时、编译时、框架 wheel 与宿主机 Driver 之间要形成兼容组合。不要在同一环境中无意识混用多个 CUDA 路径。

\subsection{问题三：NCCL 初始化卡住}

现象：

\begin{itemize}
\item 训练日志停在 initializing process group。
\item 只有部分 rank 打印日志。
\item 多节点任务长时间无输出。
\end{itemize}

排查路径：

\begin{enumerate}
\item 确认所有节点能互相解析和访问 \texttt{MASTER_ADDR}。
\item 确认 \texttt{MASTER_PORT} 未被防火墙或安全组阻断。
\item 打开 \texttt{NCCL_DEBUG=INFO}。
\item 指定正确的 \texttt{NCCL_SOCKET_IFNAME}。
\item 使用 NCCL Tests 单独验证通信。
\item 检查节点时间、DNS、容器网络和多网卡路由。
\end{enumerate}

\subsection{问题四：训练能跑但性能很差}

可能原因包括：

\begin{itemize}
\item GPU 拓扑选择不合理，跨 NUMA 或跨 PCIe switch 通信。
\item NCCL 没有使用 NVLink、NVSwitch 或 RDMA。
\item DataLoader 成为瓶颈。
\item 混合精度未启用 Tensor Core。
\item batch size 太小，GPU 无法充分利用。
\item kernel launch 太碎，缺少 fusion。
\item 日志、checkpoint 或评测过于频繁。
\end{itemize}

建议先分别测量：

\begin{itemize}
\item 单卡模型吞吐。
\item 单机多卡 NCCL 吞吐。
\item 多节点 NCCL 吞吐。
\item 数据读取吞吐。
\item checkpoint 写入吞吐。
\end{itemize}

不要在没有 baseline 的情况下直接调大规模训练任务，否则很难判断瓶颈来自模型、通信、存储还是调度系统。

\section{附录 B 小结}

CUDA 与分布式训练环境搭建的核心不是背安装命令，而是理解依赖链路：Driver 负责 GPU 访问，CUDA Runtime 支撑程序运行，PyTorch 绑定具体 CUDA 能力，NCCL 负责通信，容器和调度系统负责把这些能力正确暴露给任务。遇到问题时，应按“硬件可见性、运行时兼容性、通信连通性、性能基准”四层逐步排查。


% appendices/appendix_c_cluster_checklist.tex

\chapter{AI 集群上线检查清单}

AI 集群上线不是把 GPU 服务器接入网络、安装驱动、部署 Kubernetes 就结束。对于大模型训练和推理来说，任何一个薄弱环节都会被放大：一块 GPU ECC 异常可能导致训练中断，一条链路带宽不足可能拖垮 AllReduce，一组错误的节点标签可能导致拓扑感知调度失效，一个过慢的共享存储可能让数百张 GPU 等待数据。

本附录给出一份面向生产环境的 AI 集群上线检查清单。读者可以将其改造成内部 Runbook、验收表或故障排查流程。

\section{GPU 健康检查}

\subsection{基础状态}

每台 GPU 节点上线前，应至少检查：

\begin{itemize}
\item GPU 是否全部可见。
\item Driver 版本是否符合集群标准。
\item GPU 型号、显存容量是否与资产记录一致。
\item 是否存在 Xid 错误。
\item 是否存在 ECC 错误。
\item 温度、功耗、风扇状态是否正常。
\item Persistence Mode 是否按需开启。
\end{itemize}

基础命令：

\begin{lstlisting}[language=bash,caption={GPU 基础健康检查},label={lst:gpu-health-basic}]
nvidia-smi
nvidia-smi -q
nvidia-smi -L
\end{lstlisting}

如果集群中存在 MIG，应额外检查 MIG 切分状态：

\begin{lstlisting}[language=bash,caption={MIG 状态检查},label={lst:mig-check}]
nvidia-smi mig -lgi
nvidia-smi mig -lci
\end{lstlisting}

\subsection{错误与告警}

重点关注如下问题：

\begin{itemize}
\item \textbf{Xid Error}：通常表示 GPU driver、硬件、PCIe 或应用异常。
\item \textbf{ECC Error}：可能表示显存可靠性问题。
\item \textbf{GPU Fallen Off Bus}：常见于硬件连接、PCIe、供电或驱动异常。
\item \textbf{Thermal Throttling}：散热不足会导致性能波动。
\item \textbf{Power Limit Throttling}：功耗限制可能使 GPU 无法达到预期频率。
\end{itemize}

生产集群中，不建议让存在持续 Xid 或 ECC 异常的节点进入训练资源池。可以先将节点标记为不可调度，完成硬件诊断后再恢复。

\section{网络连通性检查}

\subsection{基础网络}

多节点训练对网络极其敏感。上线前应检查：

\begin{itemize}
\item 节点之间 hostname 是否可解析。
\item 管理网络、存储网络、训练网络是否按设计隔离。
\item MTU 是否一致。
\item 多网卡路由是否正确。
\item 防火墙、安全组、ACL 是否允许训练通信端口。
\item RDMA / RoCE / InfiniBand 设备是否可见。
\end{itemize}

基础命令：

\begin{lstlisting}[language=bash,caption={网络基础检查},label={lst:network-basic-check}]
hostname
ip addr
ip route
ping <peer-node>
ssh <peer-node>
\end{lstlisting}

检查 MTU：

\begin{lstlisting}[language=bash,caption={MTU 检查},label={lst:mtu-check}]
ip link show
\end{lstlisting}

\subsection{RDMA / InfiniBand / RoCE 检查}

若集群使用 RDMA 网络，应检查 HCA 状态：

\begin{lstlisting}[language=bash,caption={RDMA 设备检查},label={lst:rdma-device-check}]
ibv_devices
ibv_devinfo
rdma link
\end{lstlisting}

对于 InfiniBand 网络，还可以检查：

\begin{lstlisting}[language=bash,caption={InfiniBand 状态检查},label={lst:ib-check}]
ibstat
ibstatus
\end{lstlisting}

对于 RoCE 网络，应额外关注：

\begin{itemize}
\item PFC 是否按优先级正确启用。
\item ECN 标记与拥塞控制是否配置。
\item 交换机 buffer 与队列策略是否符合设计。
\item lossless Ethernet 配置是否在所有链路上保持一致。
\end{itemize}

\section{NCCL 性能检查}

\subsection{单机多卡通信}

单机多卡应先检查 GPU 拓扑：

\begin{lstlisting}[language=bash,caption={单机 GPU 拓扑检查},label={lst:single-node-topo-check}]
nvidia-smi topo -m
\end{lstlisting}

然后运行 NCCL Tests：

\begin{lstlisting}[language=bash,caption={单机 NCCL all_reduce 测试},label={lst:single-node-nccl-test}]
./build/all_reduce_perf -b 8 -e 4G -f 2 -g 8
\end{lstlisting}

检查重点：

\begin{itemize}
\item 是否使用预期的 NVLink / NVSwitch / PCIe 路径。
\item 带宽是否与同型号机器 baseline 接近。
\item 不同 GPU 组合之间是否存在明显性能差异。
\end{itemize}

\subsection{多节点通信}

多节点测试应从小规模开始：

\begin{enumerate}
\item 2 节点，每节点 1 GPU。
\item 2 节点，每节点 8 GPU。
\item 多节点完整规模。
\end{enumerate}

这样可以逐步区分问题来自单机拓扑、跨节点网络还是大规模集合通信。

多节点性能验收指标包括：

\begin{itemize}
\item AllReduce 带宽。
\item AllGather 带宽。
\item ReduceScatter 带宽。
\item 小消息延迟。
\item 大消息稳定吞吐。
\item 长时间压测中的错误率。
\end{itemize}

\section{存储吞吐检查}

\subsection{训练数据读取}

大模型训练通常需要持续读取大量 tokenized dataset。存储瓶颈会导致 GPU 空转。上线前应验证：

\begin{itemize}
\item 本地 NVMe 读取吞吐。
\item 共享文件系统读取吞吐。
\item 对象存储下载吞吐。
\item 多节点并发读取下的吞吐。
\item 小文件场景下的 metadata 性能。
\end{itemize}

示例命令：

\begin{lstlisting}[language=bash,caption={简单顺序读取测试},label={lst:storage-read-test}]
dd if=/path/to/large_file of=/dev/null bs=1G count=16 iflag=direct
\end{lstlisting}

更严谨的测试应使用 \texttt{fio}：

\begin{lstlisting}[language=bash,caption={fio 顺序读取测试示例},label={lst:fio-read-test}]
fio --name=readtest 
--filename=/path/to/testfile 
--rw=read 
--bs=1M 
--size=64G 
--numjobs=4 
--iodepth=16 
--direct=1
\end{lstlisting}

\subsection{Checkpoint 写入}

训练任务不仅读取数据，也会周期性写入 checkpoint。Checkpoint 写入慢会造成训练 step 抖动，甚至导致大规模任务同时阻塞。

应检查：

\begin{itemize}
\item 单节点写入吞吐。
\item 多节点并发写入吞吐。
\item checkpoint 文件数量对 metadata 服务的压力。
\item 是否支持异步 checkpoint。
\item checkpoint 保留策略是否明确。
\end{itemize}

\section{K8s 节点标签与污点检查}

\subsection{节点标签}

AI 集群通常依赖 Kubernetes label 描述节点能力，例如 GPU 型号、GPU 数量、网络能力、拓扑域、存储类型等。上线前应检查标签是否完整且一致。

示例：

\begin{lstlisting}[language=bash,caption={查看 Kubernetes 节点标签},label={lst:k8s-node-labels}]
kubectl get nodes --show-labels
kubectl describe node <node-name>
\end{lstlisting}

常见标签类型包括：

\begin{itemize}
\item GPU 型号：如 A100、H100、L40S 等。
\item GPU 数量：如 8-GPU 节点、4-GPU 节点。
\item 网络类型：如 IB、RoCE、Ethernet。
\item 拓扑域：如 rack、pod、zone。
\item 存储能力：如 local-nvme、shared-fs。
\item 用途分组：如 training、serving、eval。
\end{itemize}

\subsection{污点与容忍}

为了避免普通 Pod 被调度到 GPU 节点，生产集群中常使用 taint：

\begin{lstlisting}[language=bash,caption={查看节点污点},label={lst:k8s-taints}]
kubectl describe node <node-name> | grep -i taints
\end{lstlisting}

常见策略是：

\begin{itemize}
\item GPU 节点默认添加 taint。
\item 只有声明 toleration 的训练或推理 Pod 可以调度。
\item 特殊高性能节点单独划分资源池。
\item 故障节点添加维护 taint 或直接 cordon。
\end{itemize}

\section{监控与告警检查}

\subsection{核心监控指标}

AI 集群至少应覆盖以下指标：

\begin{longtable}{p{0.3\textwidth}p{0.6\textwidth}}
\toprule
类别 & 指标 \
\midrule
GPU & utilization、memory used、temperature、power、ECC、Xid。 \
CPU & utilization、load average、NUMA、context switch。 \
Memory & used、available、page cache、OOM。 \
Network & throughput、packet drop、retransmission、RDMA error、PFC pause。 \
Storage & read/write throughput、IOPS、latency、metadata ops。 \
Kubernetes & pending pods、failed pods、node not ready、evictions。 \
Training & step time、tokens/s、loss、gradient norm、checkpoint time。 \
Serving & QPS、TTFT、TPOT、batch size、KV Cache usage、request error rate。 \
\bottomrule
\end{longtable}

\subsection{告警原则}

告警不应只覆盖“机器挂了”，还应覆盖“性能正在异常下降”。例如：

\begin{itemize}
\item GPU 利用率长期低于阈值但任务仍在运行。
\item 某节点网络吞吐明显低于同组节点。
\item checkpoint 写入耗时持续升高。
\item 推理服务 TTFT 或 TPOT 超过 SLA。
\item Pending GPU Pod 数量持续增长。
\item 某类 GPU 错误在短时间内频繁出现。
\end{itemize}

\section{作业失败排查流程}

\subsection{训练任务失败}

训练任务失败时，建议按以下顺序排查：

\begin{enumerate}
\item 查看调度系统事件：任务是否成功分配资源。
\item 查看 Pod / Job 状态：是否 OOM、Evicted、ImagePullBackOff。
\item 查看 rank 0 日志：是否有模型、数据或 checkpoint 错误。
\item 查看所有 rank 日志：是否只有部分 rank 失败。
\item 查看 NCCL 日志：是否通信初始化或集合通信失败。
\item 查看节点日志：是否出现 GPU Xid、网络错误或磁盘错误。
\item 复现小规模任务：判断是否与规模相关。
\end{enumerate}

\subsection{推理服务异常}

推理服务异常时，建议关注：

\begin{itemize}
\item 请求是否进入队列但长时间未调度。
\item prefill 是否成为瓶颈。
\item decode 阶段 TPOT 是否异常。
\item KV Cache 是否耗尽。
\item batch size 是否长期过小或过大。
\item 是否存在热点模型或热点租户。
\item autoscaling 是否滞后。
\end{itemize}

\section{上线验收表}

下面是一份简化版上线验收表，可作为集群交付前的检查模板：

\begin{longtable}{p{0.28\textwidth}p{0.42\textwidth}p{0.18\textwidth}}
\toprule
检查项 & 验收标准 & 状态 \
\midrule
GPU 可见性 & 所有节点 GPU 数量、型号与资产记录一致 & 待确认 \
Driver 版本 & 所有 GPU 节点版本一致或符合兼容矩阵 & 待确认 \
GPU 健康 & 无持续 Xid / ECC / 掉卡异常 & 待确认 \
单机 NCCL & 单机 AllReduce 带宽达到 baseline & 待确认 \
跨节点 NCCL & 多节点 AllReduce 带宽达到 baseline & 待确认 \
RDMA 状态 & HCA 正常，端口 active，错误计数正常 & 待确认 \
存储读取 & 数据读取吞吐满足训练需求 & 待确认 \
Checkpoint 写入 & 并发写入不阻塞训练 & 待确认 \
K8s 标签 & GPU 型号、网络、拓扑、用途标签完整 & 待确认 \
K8s 污点 & GPU 节点隔离策略生效 & 待确认 \
监控指标 & GPU、网络、存储、训练、推理指标可见 & 待确认 \
告警规则 & 故障与性能退化均有告警 & 待确认 \
\bottomrule
\end{longtable}

\section{附录 C 小结}

AI 集群的上线验收，本质上是一次端到端系统验证：硬件是否健康，网络是否达到预期，存储是否能支撑数据流，调度系统是否正确理解资源，监控是否能及时暴露问题。只有当这些基础能力稳定后，训练框架、推理引擎和上层平台才能发挥真正价值。


% appendices/appendix_d_references.tex

\chapter{参考资料与延伸阅读}

AI Infrastructure 是一个跨越算法、系统、硬件、网络和平台工程的领域。单靠一本书无法覆盖所有细节。本附录按照主题列出建议继续阅读的论文、系统、文档和开源项目。阅读这些资料时，建议不要只关注结论，而要关注它们解决的问题、隐含的硬件假设、系统瓶颈和工程取舍。

\section{Transformer 与 LLM 论文}

\subsection{基础架构}

\begin{itemize}
\item Vaswani et al., \textit{Attention Is All You Need}. Transformer 架构的奠基论文，建议重点阅读 Scaled Dot-Product Attention、Multi-Head Attention 和 Position-wise FFN 的设计动机。
\item Radford et al., \textit{Improving Language Understanding by Generative Pre-Training}. GPT 系列早期工作，展示了预训练语言模型的基本范式。
\item Brown et al., \textit{Language Models are Few-Shot Learners}. GPT-3 论文，建议关注规模扩展、in-context learning 和训练数据混合方式。
\item Touvron et al., \textit{LLaMA: Open and Efficient Foundation Language Models}. 理解现代开源 Decoder-only LLM 的重要参考。
\item Touvron et al., \textit{Llama 2: Open Foundation and Fine-Tuned Chat Models}. 建议结合第 6 章阅读，关注预训练、SFT 和对齐流程。
\end{itemize}

\subsection{位置编码与结构改进}

\begin{itemize}
\item Su et al., \textit{RoFormer: Enhanced Transformer with Rotary Position Embedding}. RoPE 的代表性论文。
\item Shazeer, \textit{GLU Variants Improve Transformer}. 理解 SwiGLU、GEGLU 等 FFN 变体的参考。
\item Shazeer, \textit{Fast Transformer Decoding: One Write-Head is All You Need}. Multi-Query Attention 的重要参考。
\item Ainslie et al., \textit{GQA: Training Generalized Multi-Query Transformer Models from Multi-Head Checkpoints}. Grouped-Query Attention 的重要参考。
\end{itemize}

\subsection{规模定律与训练范式}

\begin{itemize}
\item Kaplan et al., \textit{Scaling Laws for Neural Language Models}. 早期语言模型规模定律研究。
\item Hoffmann et al., \textit{Training Compute-Optimal Large Language Models}. Chinchilla Scaling Law 的代表性工作。
\item Ouyang et al., \textit{Training language models to follow instructions with human feedback}. InstructGPT 与 RLHF 的重要参考。
\end{itemize}

\section{分布式训练论文}

\subsection{数据并行与集合通信}

\begin{itemize}
\item Li et al., \textit{PyTorch Distributed: Experiences on Accelerating Data Parallel Training}. 理解 PyTorch DDP 工程设计的重要资料。
\item Sergeev and Del Balso, \textit{Horovod: Fast and Easy Distributed Deep Learning in TensorFlow}. Ring AllReduce 训练实践的重要系统。
\item NVIDIA NCCL Documentation. 理解 AllReduce、AllGather、ReduceScatter 等通信原语及其性能调优。
\end{itemize}

\subsection{模型并行}

\begin{itemize}
\item Shoeybi et al., \textit{Megatron-LM: Training Multi-Billion Parameter Language Models Using Model Parallelism}. Tensor Parallelism 和大模型训练的重要论文。
\item Narayanan et al., \textit{Efficient Large-Scale Language Model Training on GPU Clusters Using Megatron-LM}. 3D 并行、流水线调度和大规模训练实践的重要参考。
\item Huang et al., \textit{GPipe: Efficient Training of Giant Neural Networks using Pipeline Parallelism}. Pipeline Parallelism 的经典论文。
\item Narayanan et al., \textit{PipeDream: Generalized Pipeline Parallelism for DNN Training}. 理解流水线并行调度的重要参考。
\end{itemize}

\subsection{ZeRO 与内存优化}

\begin{itemize}
\item Rajbhandari et al., \textit{ZeRO: Memory Optimizations Toward Training Trillion Parameter Models}. ZeRO Stage 1/2/3 的核心论文。
\item Rajbhandari et al., \textit{ZeRO-Offload: Democratizing Billion-Scale Model Training}. CPU offload 的重要参考。
\item Ren et al., \textit{ZeRO-Infinity: Breaking the GPU Memory Wall for Extreme Scale Deep Learning}. CPU / NVMe offload 与更大模型训练的参考。
\item Chen et al., \textit{Training Deep Nets with Sublinear Memory Cost}. Gradient Checkpointing 的经典工作。
\end{itemize}

\section{Serving Infra 论文与系统}

\subsection{KV Cache 与推理调度}

\begin{itemize}
\item Kwon et al., \textit{Efficient Memory Management for Large Language Model Serving with PagedAttention}. vLLM 与 PagedAttention 的核心论文。
\item Yu et al., \textit{Orca: A Distributed Serving System for Transformer-Based Generative Models}. Continuous Batching 和 iteration-level scheduling 的重要系统。
\item NVIDIA FasterTransformer / TensorRT-LLM Documentation. 理解高性能 LLM 推理 kernel、engine 构建和部署优化的重要资料。
\end{itemize}

\subsection{Attention 优化}

\begin{itemize}
\item Dao et al., \textit{FlashAttention: Fast and Memory-Efficient Exact Attention with IO-Awareness}. FlashAttention 的核心论文。
\item Dao, \textit{FlashAttention-2: Faster Attention with Better Parallelism and Work Partitioning}. 理解更高效 attention kernel 并行划分的参考。
\item Shah et al., \textit{FlashAttention-3: Fast and Accurate Attention with Asynchrony and Low-precision}. 理解新一代 GPU 上 attention 优化方向的参考。
\end{itemize}

\subsection{量化与压缩}

\begin{itemize}
\item Frantar et al., \textit{GPTQ: Accurate Post-Training Quantization for Generative Pre-trained Transformers}. GPTQ 的代表性论文。
\item Lin et al., \textit{AWQ: Activation-aware Weight Quantization for LLM Compression and Acceleration}. AWQ 的代表性论文。
\item Dettmers et al., \textit{LLM.int8(): 8-bit Matrix Multiplication for Transformers at Scale}. 大模型 INT8 推理的重要参考。
\item Dettmers et al., \textit{QLoRA: Efficient Finetuning of Quantized LLMs}. 量化与参数高效微调结合的重要参考。
\end{itemize}

\section{CUDA 与 GPU 优化资料}

\subsection{官方文档}

\begin{itemize}
\item NVIDIA CUDA C++ Programming Guide. 学习 CUDA 编程模型、线程层次、内存模型和同步机制的基础资料。
\item NVIDIA CUDA C++ Best Practices Guide. 学习 coalescing、shared memory、occupancy、异步拷贝等优化技巧。
\item NVIDIA Nsight Systems Documentation. 分析端到端 timeline、CPU-GPU 协作和通信重叠。
\item NVIDIA Nsight Compute Documentation. 分析单个 CUDA kernel 的 occupancy、memory throughput、warp stall 和 Tensor Core 利用率。
\end{itemize}

\subsection{矩阵计算与 Kernel 优化}

\begin{itemize}
\item CUTLASS Documentation. 理解 tile-based GEMM、Tensor Core MMA 和高性能矩阵乘法模板的重要资料。
\item Triton Documentation. 学习使用 Python-like DSL 编写高性能 GPU kernel。
\item cuBLAS Documentation. 理解 GEMM、batched GEMM 和 mixed precision 矩阵计算。
\item cuDNN Documentation. 理解深度学习基础算子的高性能实现。
\end{itemize}

\section{网络与集群调度资料}

\subsection{高性能网络}

\begin{itemize}
\item InfiniBand Architecture Specification. 理解 InfiniBand 网络、queue pair、completion queue 和 RDMA 语义的基础资料。
\item NVIDIA GPUDirect RDMA Documentation. 理解 GPU memory 与 NIC 直接通信的关键资料。
\item RoCE Deployment Guides. 理解 RoCEv2、PFC、ECN、lossless Ethernet 与拥塞控制。
\item NCCL Tuning Guide. 理解多 NIC、多 rail、拓扑感知通信和环境变量调优。
\end{itemize}

\subsection{Kubernetes 与 AI 调度}

\begin{itemize}
\item Kubernetes Documentation. Pod、Node、Scheduler、Device Plugin、Operator 等基础概念。
\item NVIDIA Kubernetes Device Plugin. GPU 资源暴露、调度与容器运行时集成的重要项目。
\item Volcano Documentation. Gang Scheduling 和 AI/HPC batch workload 调度的重要系统。
\item Kueue Documentation. Kubernetes 原生队列管理、配额和批作业调度的重要项目。
\item Kubeflow Training Operator Documentation. 在 Kubernetes 上运行分布式训练任务的重要参考。
\end{itemize}

\section{开源项目索引}

\subsection{训练框架与并行训练}

\begin{longtable}{p{0.26\textwidth}p{0.64\textwidth}}
\toprule
项目 & 关注点 \
\midrule
PyTorch Distributed & DDP、ProcessGroup、分布式通信基础。 \
DeepSpeed & ZeRO、offload、混合精度、大规模训练优化。 \
Megatron-LM & Tensor Parallelism、Pipeline Parallelism、3D Parallelism。 \
Hugging Face Accelerate & 多设备训练启动与配置抽象。 \
Colossal-AI & 并行训练、内存优化和大模型训练工具链。 \
FairScale & FSDP、checkpointing 等分布式训练能力。 \
\bottomrule
\end{longtable}

\subsection{推理服务}

\begin{longtable}{p{0.26\textwidth}p{0.64\textwidth}}
\toprule
项目 & 关注点 \
\midrule
vLLM & PagedAttention、Continuous Batching、高吞吐 LLM Serving。 \
TensorRT-LLM & NVIDIA GPU 上的高性能 LLM 推理优化。 \
Text Generation Inference & 面向生产部署的文本生成推理服务。 \
SGLang & LLM serving runtime、structured generation 与调度优化。 \
llama.cpp & CPU / GPU 本地推理、量化格式和轻量部署。 \
ONNX Runtime & 通用模型推理优化和跨硬件执行。 \
\bottomrule
\end{longtable}

\subsection{可观测性与运维}

\begin{longtable}{p{0.26\textwidth}p{0.64\textwidth}}
\toprule
项目 & 关注点 \
\midrule
Prometheus & 指标采集与时序存储。 \
Grafana & 监控面板与可视化。 \
DCGM Exporter & NVIDIA GPU 指标采集。 \
OpenTelemetry & 分布式 tracing、metrics、logs 统一观测。 \
Loki & 日志聚合。 \
Jaeger & 分布式链路追踪。 \
\bottomrule
\end{longtable}

\section{阅读建议}

面对大量资料，建议按问题驱动阅读，而不是按论文发布时间线机械阅读：

\begin{enumerate}
\item 如果目标是理解模型结构，先读 Transformer、GPT、LLaMA、RoPE、GQA 等资料。
\item 如果目标是训练大模型，先读 DDP、Megatron-LM、ZeRO、GPipe、NCCL 文档。
\item 如果目标是优化推理服务，先读 KV Cache、PagedAttention、FlashAttention、Continuous Batching 和量化论文。
\item 如果目标是做 GPU kernel 优化，先读 CUDA Programming Guide、Best Practices Guide、CUTLASS 和 Triton 文档。
\item 如果目标是建设 AI 集群，先读 NCCL、GPUDirect RDMA、Kubernetes Device Plugin、Volcano、Kueue 和监控系统资料。
\end{enumerate}

读论文时可以重点问五个问题：

\begin{itemize}
\item 它解决的瓶颈是什么：计算、显存、通信、调度还是工程复杂度？
\item 它假设的硬件环境是什么：单卡、多卡、NVLink、RDMA、HBM 容量？
\item 它优化的是训练还是推理，是吞吐还是延迟？
\item 它引入了哪些新的 trade-off？
\item 它在真实生产系统中落地时还缺什么？
\end{itemize}

\section{附录 D 小结}

AI Infrastructure 的知识更新很快，但底层问题相对稳定：如何让计算更满、显存更省、通信更少、调度更准、故障更可控。持续阅读论文和开源系统时，应始终把注意力放在这些长期问题上，而不是只追逐某个短期流行框架或参数配置。


\backmatter

\bibliographystyle{plain}
\bibliography{references}

\end{document}